% 日本語 LuaLaTeX。2026-11-25 / v0.44（相対モジュライ原稿の現行構成）
% エンジン、出力先、SyncTeX は repo ルートの .latexmkrc で指定する。
\RequirePackage{fix-cm}
\documentclass[a4paper,11pt]{ltjsarticle}
\usepackage{amsmath,amssymb,amsthm,mathtools,mathrsfs}
% rsfs の小さな添字も、代替サイズではなく設計フォントの拡縮で組む。
\DeclareFontFamily{U}{rsfs}{\skewchar\font127}
\DeclareFontShape{U}{rsfs}{m}{n}{<-6> rsfs5 <6-8> rsfs7 <8-> rsfs10}{}
\usepackage[margin=25mm]{geometry}
\usepackage{booktabs,enumitem}
\usepackage[unicode,pdfusetitle,colorlinks=true,
 linkcolor=blue,citecolor=blue,urlcolor=blue]{hyperref}
\numberwithin{equation}{section}
\theoremstyle{plain}
\newtheorem{theorem}{定理}[section]
\newtheorem{proposition}[theorem]{命題}
\newtheorem{lemma}[theorem]{補題}
\newtheorem{corollary}[theorem]{系}
\theoremstyle{definition}
\newtheorem{definition}[theorem]{定義}
\newtheorem{example}[theorem]{例}
\newtheorem{problem}[theorem]{問題}
\theoremstyle{remark}
\newtheorem{remark}[theorem]{注意}
\renewcommand{\proofname}{証明}
\newcommand{\kk}{k}
\newcommand{\OO}{\mathcal O}
\newcommand{\cC}{\mathcal C}
\newcommand{\PP}{\mathbb P}
\newcommand{\HH}{\mathcal H_2}
\newcommand{\HHstack}{\mathfrak H_2}
\DeclareMathOperator{\Hom}{Hom}
\DeclareMathOperator{\End}{End}
\DeclareMathOperator{\Ext}{Ext}
\DeclareMathOperator{\coh}{coh}
\DeclareMathOperator{\qgr}{qgr}
\DeclareMathOperator{\gldim}{gldim}
\DeclareMathOperator{\Aut}{Aut}
\DeclareMathOperator{\Pic}{Pic}
\DeclareMathOperator{\Sym}{Sym}
\DeclareMathOperator{\Gr}{Gr}
\DeclareMathOperator{\GL}{GL}
\DeclareMathOperator{\PGL}{PGL}
\DeclareMathOperator{\rank}{rank}
\DeclareMathOperator{\ch}{ch}
\DeclareMathOperator{\td}{td}
\DeclareMathOperator{\length}{length}
\title{構造対象付き非可換二次曲面の内在的モジュライ：\\代数性と noncubic 境界}
\author{}
\date{2026年11月25日\quad v0.44}

\begin{document}
\maketitle
\begin{abstract}
構造対象 $\mathcal O$ のみを指定した強い非可換二次曲面の
内在的な相対モジュライを考える。反 Serre の適合平方根と
正則正規二次切断を持つ AS4 被覆は局所存在条件であり、
対象の追加標識にはしない。
体上の pointed exact 同値分類を別論文の入力として、
局所有限型基底の fppf site 上で $\mathfrak H_2$ の
有限型・滑らか・既約な純次元三の Artin stack 表示と、
滑らかな八十三次元の framed 実現を扱う。
楕円型の構造対象忘却のファイバー、および
被約 noncubic 境界の滑らかな余次元二の閉埋込み、
二形式表示商との相対同値、慣性、
good/coarse moduli と一致根対の形式局所構造を記述する。

非可算な標数零代数閉体上では、標識を忘れた線形 exact 同値類を幾何点上で
正確に分類する射の標的は、準コンパクトな対角を持ち得ない。
有限型 atlas 上の構成可能集合の分割と楕円型の可算無限ファイバーからこれを示す。
\end{abstract}
\tableofcontents

\section{Introduction}\label{sec:moduli-intro}

\subsection{対象と主結果}
体上のpointed exact同値類の分類と、非被約基底を含む相対族・
自己同値・降下の分類は異なる。本稿は
\cite{QuadricIntrinsic2026} の体上の強い非可換二次曲面
$(\mathcal C,\OO)$の分類を入力に、構造対象$\OO$だけを保持する
内在的な相対$\HHstack$を扱う。適合するexactな反Serre平方根、
同じ根のAS4被覆と根の同型はfppf局所での存在条件であり、
標識には加えない。一般のNoetherian基底上で同じ素朴なAbel公理を
課した関手との一致は、局所有限型site上の主結果とは別問題である。

第一の主結果は定理~\ref{thm:global-h2-smooth-finite-type} と
系~\ref{cor:global-xas-smooth-scheme} である。
$\HHstack$は有限型・smooth・既約・純次元三のArtinスタックであり、
対角はaffine有限表示、そのframed実現$X_{\mathrm{as}}$は
smooth既約八十三次元スキームとなる。これは幾何点の分類だけでなく、
有限型atlas上の実際の相対Abel族と指定pointed同型の持上げから得る。

第二に、定理~\ref{thm:noncubic-boundary-smooth-codim-two} と
定理~\ref{thm:noncubic-intrinsic-presentation-equivalence} は、
cubic開の補集合の被約構造$\mathfrak N$がsmooth純次元一で
余次元二の正則閉埋込みをなし、二つの非退化二元二次形式の
表示商$[P/G]$と、非被約test上の相対groupoidとして同値であることを示す。
さらに境界のgood moduli $\mathbb A^1_\chi$、
DM開の粗モジュライと全$\HHstack$の粗モジュライ代数空間の不存在を
区別する（第~\ref{sec:noncubic-boundary-comparison}節）。
一致根対では有限成分を含む全慣性作用を計算し、
定理~\ref{thm:noncubic-coincident-formal-pair} で
境界を含む形式局所商を記述する。この共通étale近傍上の
good moduliが三次元でsmoothであることは、
$\HHstack$のZariski開全体へのgood moduliの降下を意味しない。
全$\HHstack$の準射影的good moduliは、本稿の相対開認識を入力とする
別稿C\cite[\texttt{thm:qa-global-h2-quasiprojective-good-moduli}]{QuiverAvatarModuli2026}
で、明示半不変式の飽和被覆から構成される。

第三に、定理~\ref{thm:unmarked-diagonal} は構造対象を忘れる際の障害を与える。
非可算な標数零代数閉体上で、二点の像の同型と基礎 Abel 圏の線形 exact 同値が
同値になるなら、分類射の標的の対角は準コンパクトでない。
有限型 atlas 上で、像が指定対象と同型になる点と、各 pointed 同型類を表す点が
ともに構成可能になることを使う。これは上記の pointed 粗モジュライの問題とは異なる。

\subsection{既知入力と本稿の追加部分}
Artin--Zhang \cite{AZ1994} の再構成、
Van den Bergh \cite{VdB2011} のcubic分類、
Lowen--Van den Bergh \cite{LowenVanDenBergh2005,LowenVanDenBergh2006}、
De Deken--Lowen \cite{DeDekenLowen2011} の変形・helix持上げ、
Polishchuk--Positselski \cite{PolishchukPositselski2005} の
Koszul変形を既知入力として用いる。

特にBelmans \cite[Proposition 31, Corollary 32, Table 2]{Belmans2019} は
cubic AS正則$\mathbb Z$-algebraに対応する
非可換二次曲面のHochschild cohomologyを扱い、
$HH^{\ge3}=0$ と $\dim HH^2-\dim HH^1=3$、
および退化型の$HH^1,HH^2$を求めた。
四頂点なので、Belmans の式の差 $\#Q_0-1$ は三になる。
これは一点での変形空間と無限小自己同型の次元差である。
定理~\ref{thm:global-h2-smooth-finite-type} を得るには、
非被約な基底にわたる族と指定同型を持ち上げ、それらを有限型の
atlas 上で実現する必要がある。さらに二形式の根対が一致する
noncubic 点では、境界に横断する二方向を含めた全慣性の作用を
定理~\ref{thm:noncubic-coincident-formal-pair} で求める。

証明は相対Koszul・Noetherian性と有限foundationの変形、
適合Serre根と同じ根のAS4被覆を準備し、
cubic/noncubicの有限型効果性を別々に確立してから
smooth atlasを得る。境界では非被約正規形、相対tilting、
反Serre再構成によって二形式表示との同値を示し、
慣性・good/coarse moduli・形式局所構造を求める。
Van den Bergh \cite[arXiv v3, Theorem 5.5.10]{VdB2011} は添字付き cubic
代数の同型を四つ組から復元する。Kitamura との共同研究の Okawa による報告
\cite[Theorems 1, 3]{OkawaKitamura2022OWR} は cubic 型の無標識同値の
$D_\infty$ 軌道と line bundle の分類を述べる。本稿は二つの内在的な
foundation から非共鳴楕円型の pointed 同型類を数え、さらに通常の代数スタックの
対角と Chevalley の定理 \cite[Tags 026N, 04XQ, 0ECX, 04YV]{StacksSmoothAtlas}
を用いて分類射への障害を導く。
ここで報告の Theorem 3 が数えるのは一つの $\operatorname{qgr}$ の中の
line bundle の同型類である。本稿の忘却で数えるのは、そのような対象を
構造対象とする組を圏の自己同値で割った同型類である。
滑らかな楕円型で $\delta$ を非捩れ、$\eta\notin\mathbb Z\delta$ と取ると、
この商を取った後にも可算無限個が残る
（定理~\ref{thm:elliptic-marking-lattice-fiber}）。
この事実が定理~\ref{thm:unmarked-diagonal} に必要な入力である。

% !TEX root = ../nc-quadric-moduli.tex
% 2026-11-22 / v1.9（相対定義とflatness剛性補題のB移管）
\section{内在的な相対対象と基底変更}
\label{sec:moduli-relative}

本節では $k$ を標数零の代数閉体とし、まず $k$ 上局所有限型の
スキームの big fppf site を用いる。この範囲は非被約スキームと
有限型の局所 Artin 環を含む。任意の noetherian $k$-代数に対する
連続拡張や極限保存性を、本節の構成だけから結論しない。
とくに finite type の基底での局所 noetherian 性と、任意の
noetherian 基底変更での strong noetherian 性を区別する。

\subsection{Abel 圏の基底変更を指定する}

局所有限型の affine スキーム $U=\operatorname{Spec}A$ 上で、
$A$-線形 noetherian Abel 圏を $\mathcal C_U$ とし、
$\mathcal A_U=\operatorname{Ind}\mathcal C_U$ と置く。
$A\to A'$ に対し $(\mathcal A_U)_{A'}$ は
$\mathcal A_U$ 内の $A'$-加群対象の圏、すなわち
$A'$ の作用を備えた対象とその作用を保つ射の圏である。
この構成において引戻しは
$A'\otimes_A-$ であり、一般には右完全、$A'$ が $A$ 上平坦なら完全である。
本稿でいう圏の基底変更は、Hom 加群を単に tensor して対象集合を固定する
線形圏の基底変更ではない。
この区別と基本的な随伴は
\cite[\S4]{LowenVanDenBergh2006}、任意の有限型 affine 基底に必要な
加群対象の定式化と flat descent は
\cite[\S7.1, Definition 7.3, Remarks 7.4--7.5]{AHLHExistence}
を参照する。後者は Artin--Zhang の abstract Hilbert schemes の
B2.2, B3.16, B6.3, C8.6 を明記して用いている。

\begin{definition}[基底変更を備える Abel 圏の族]
\label{def:relative-abelian-system}
$T$ 上の Abel 圏の族とは、$T$ 上の locally finite type affine
$U$ に $\mathcal C_U$ を割り当てる pseudofunctor であって、次を満たすものとする。
\begin{enumerate}
\item 各 $\mathcal C_U$ は $A$-線形で本質的に小さい noetherian Abel 圏であり、
  その Ind 圏は局所 noetherian Grothendieck 圏である。
\item affine の間の射 $U'=\operatorname{Spec}A'\to U$ について
  \[
    \mathcal A_{U'}\simeq(\mathcal A_U)_{A'}
  \]
  という基底変更同値があり、引戻しは $A'\otimes_A-$ に対応する。
  これらの比較同値は合成と恒等射に対して coherent である。
\item 対象、射、完全列について fppf descent が成立する。
  とくに $\mathcal C$ は対象と全ての射を保持する stack である。
\item $\operatorname{Hom}_{\mathcal C_U}(E,F)$ は有限 $A$-加群であり、
  affine flat base change と可換する。
\end{enumerate}
射はこの系の間の、基底変更と可換する線形完全同値とする。
\end{definition}

この定義は圏の局所表示を指定するものではない。
基底変更同値は、族という語の意味を固定するためのデータであり、
射はその coherent な同型を許す。
有限型 affine 基底変更での局所 noetherian 性は
Artin--Zhang B6.3 の範囲内である。
任意の noetherian 基底変更へ、この結論を無条件に拡張しない。

$\mathcal C_U$ の対象 $E$ が $A$ 上平坦とは、
$-\otimes_AE:\operatorname{Mod}A\to\mathcal A_U$ が完全であることをいう。
本稿では
\[
 \mathcal D_U=D(\mathcal A_U)^c
\]
を標準の導来 enhancement における compact 対象の圏とする。
この圏は $\mathcal C$ から決まり、別の dg enhancement を標識に加えない。
以下の導来条件は $T$ の fppf 局所で課す。
すなわち fppf 被覆 $T_i\to T$ が存在して、その各 $T_i$ への
全ての本節の affine 基底変更 $U$ 上で、$D(\mathcal A_U)$ は
compactly generated であり、$\mathcal D_U$ は $D^b(\mathcal C_U)$ に含まれ、
相対 smooth proper な $\operatorname{Perf}A$-線形圏となり、導来基底変更
\[
 \mathcal D_U\otimes_{\operatorname{Perf}A}\operatorname{Perf}A'
       \simeq\mathcal D_{U'}
 \tag{RP-BC}\label{eq:relative-perfect-basechange}
\]
と整合することを要求する。
$\mathcal O$ はこれらの局所 chart 上でこの compact 部分に属するものとする。
chart の選択は標識に含めない。
本節の stack 性の証明には、貼合せ後の大域的な
$D(\operatorname{Ind}\mathcal C_U)$ が局所導来圏の descent と一致するという
追加主張を用いない。定義そのものが fppf 局所条件であることを用いる。
特異な基底上で $\mathcal D_U=D^b(\mathcal C_U)$ としたり、
compact 部分だけに bounded heart を課したりしない。
また全ての底上平坦 coherent 対象が compact であることを
今回の定義には追加しない。この一致は別の相対認識問題である。
体上では、以下の foundation による議論で
$\mathcal D_k=D^b(\mathcal C_k)$ が従う。

以下の Hom、双対、導来 tensor はこの標準 enhancement で解釈する。
(\ref{eq:relative-perfect-basechange}) は導来圏の基底変更であり、
heart の通常の引戻しが全ての基底変更で完全であるという主張ではない。

\begin{definition}[相対 Serre 条件]
\label{def:relative-serre-condition}
上の族が相対 Serre 条件を満たすとは、前記の局所 chart の各 affine $U$ 上で
次が成立し、任意の本節の基底変更と整合することをいう。
\begin{enumerate}
\item compact 対象 $E,F$ の間の $\mathbf R\operatorname{Hom}(E,F)$
  は $A$ 上 perfect で、導来基底変更と可換する。
\item $\mathcal D_U$ に自己同値 $S_U$ が存在し、
  functorial な同型
  \[
   \mathbf R\operatorname{Hom}(E,F)^\vee_A
       \simeq\mathbf R\operatorname{Hom}(F,S_UE)
  \]
  を与える。
\end{enumerate}
Serre 関手および duality の選択自体は標識に含めない。
二つの選択の比較は、表現する双対 functor の Yoneda 一意性による。
\end{definition}

\subsection{相対的な強い\texorpdfstring{$\mathcal H_2$}{H2}条件}

相対 quadratic AS4 被覆とは、局所 affine 基底 $A$ 上の次数付き環
$\mathcal B=T_A(V)/(I)$ であって、$\operatorname{rk}V=4$、
$I\subset V^{\otimes2}$ が階数六の部分束、
各次数が有限局所自由、左右 noetherian、
全幾何 fiber が型 $(1,4,6,4,1)$ の AS 正則環となるものをいう。
加えて、標準の左右 Koszul 複体は全次数で完全であり、
基底変更後にも同じ完全列を与えるものとする。
ここで幾何 fiber の定義には、記述された有限個の係数を含む体からの
体拡大で作る次数付き環を用いる。
被覆の二次の正則正規元 $g$ は、左右乗法の単射性と正規性に加え、
$\mathcal B/(g)$ の次数別局所自由性を要求する。
従って切断の完全列は任意の本節の基底変更と両立する。
固定した相対切断環への被覆の有限判定は、分類主論文 \cite[相対AS被覆の有限判定]{QuadricIntrinsic2026} を参照する。
平方根を選ばず、その存在だけを要求する部分関手の代数性は別問題である。

\begin{definition}[有限型基底上の内在的 $\mathfrak H_2$]
\label{def:relative-h2}
$\mathfrak H_2(T)$ の対象は、定義~\ref{def:relative-abelian-system} の族
$\mathcal C$ と対象 $\mathcal O\in\mathcal C_T$ の対で、以下を満たすものとする。
\begin{enumerate}[label=(RQ\arabic*)]
\item $\mathcal O$ は基底上平坦で、fppf 局所 chart の compact 導来部分は
  (\ref{eq:relative-perfect-basechange}) と相対 Serre 条件を満たす。
  これらの chart の全ての affine 基底変更上で
  $\mathbf R\operatorname{Hom}(\mathcal O,\mathcal O)\simeq A[0]$。
\item 各幾何 fiber の対 $(\mathcal C_{\bar t},\mathcal O_{\bar t})$ は、
  分類主論文 \cite{QuadricIntrinsic2026} で定める体上の $\mathcal H_2$ に属する。
  特に標識は同論文の Gabriel--Krull 商の意味で rank one torsion-free であり、
  Serre 条件は fiber の全 $D^b(\mathcal C_{\bar t})$ に成立する。
  fiber 圏は module-object 基底変更、または下記の全軌道環の
  $\qgr$ によって計算し、その一致を要求する。
\item $T$ の fppf 局所で、系全体の線形完全自己同値 $s$ が存在する。
  前記の chart と共通細分を取った局所 compact 導来部分で
  $s^2\simeq S^{-1}[2]$ であり、全ての affine 基底変更上で
  $(\mathcal O,s)$ は Artin--Zhang ample である。
  ここで生成条件は有限直和からの全射、Hom の条件は各全射に対する
  十分大きなシフトでの全射性として課す。
\item 同じ局所的な $s$ について、前記の相対 AS4 環 $\mathcal B$ と
  二次の正則正規元 $g$ が存在し、全次数で
  \[
   \bigoplus_{n\ge0}\operatorname{Hom}(\mathcal O,s^n\mathcal O)
     \simeq\mathcal B/(g)
  \]
  となる。この同型は任意の本節の基底変更と可換する。
\end{enumerate}
射は pointed exact 同値 $(F,u)$、二射は
$u'\eta_{\mathcal O}=u$ を満たす自然同型とする。
$s$、根同型、被覆、切断元、軌道環の表示のいずれも対象の標識ではない。
\end{definition}

この定義は「全幾何 fiber が $\mathcal H_2$ に属する」ことだけを課す定義より強い。
fppf 局所に存在すべき平方根と AS 被覆は、非被約底の変形方向を含めて要求される。
体上の分類と fiberwise な判定だけで、この相対条件を確認したとは扱わない。

\begin{lemma}[Abel flatness と標識付き二射の剛性]
\label{lem:relative-h2-flat-center}
定義~\ref{def:relative-h2} の局所族について、次が成り立つ。
\begin{enumerate}
\item $\mathcal C_U$ は Lowen--Van den Bergh の意味で $A$ 上 flat な Abel 圏である。
\item $\operatorname{Aut}(\operatorname{Id}_{\mathcal C_U})=A^\times$。
  特に pointed 一射の二自己同型は自明である。
\end{enumerate}
\end{lemma}
\begin{proof}
$A$-線形完全自己同値は tensor と可換するから、
$s^{-n}\mathcal O$ は $A$ 上平坦である。
AZ の生成条件により任意の coherent 対象はこれらの有限直和の商である。
従って任意の有限提示 $A$-加群 $M$ について
$\operatorname{Tor}^A_1(M,-)$ は co-effaceable となる。
これが \cite[Proposition 3.1, Definition 3.2]{LowenVanDenBergh2006}
の flatness 条件である。Ind 圏についても同論文 Proposition 3.6 を適用できる。

自然自己同型 $\eta$ の $s^{-n}\mathcal O$ 上の成分は、相対例外性により
単元 $a_n\in A^\times$ の scalar である。
正則二次切断の完全列により
\[
 \operatorname{Hom}(s^{-n}\mathcal O,\mathcal O)
       \simeq (\mathcal B/(g))_n
\]
は階数 $(n+1)^2$ の有限局所自由 $A$-加群、従って忠実である。
自然性は $(a_n-a_0)(\mathcal B/(g))_n=0$ を与えるので $a_n=a_0$。
AZ の全射と自然性から全ての対象上で $\eta=a_0\operatorname{id}$ となる。
pointed 二射では $a_0=1$ である。
この結論は fppf 局所で確認できるため、根を大域に選んでいない族にも成り立つ。
\end{proof}


% !TEX root = ../nc-quadric-moduli.tex
% 2026-11-22 / v1.30（全17証明を相対モジュライ論文Bへ移管）
\section{三次平坦性から全次数平坦性への相対的な認識}
\label{sec:relative-koszul}

ここでは三次までの関係式の平坦性と、全次数の平坦性との間を閉じる。
用いるのは Polishchuk--Positselski の局所環上の Koszul 変形原理
\cite[Chapter 6, Theorem 7.3]{PolishchukPositselski2005} である。
この定理は被約基底や一径数の形式基底に限らず、
任意の Noetherian 局所環に適用できる。
従って本節の結論は双対数上の一次変形だけでなく、
任意の Artin 基底と非被約 Noetherian 基底を含む。
一方、得られるのは次数付き代数とその線型解消の相対的な結論であり、
相対 Abel heart、Serre 関手、平方根、AS 被覆の認識は別に必要である。

\begin{proposition}[全 $\HH$ の切断環と反 Serre 環の Koszul 性]
\label{prop:all-h2-antiserre-koszul}
$k$ を本稿の基礎体とし、$(\mathcal C,\mathcal O)\in\HH$ とする。
任意の適合平方根 $s$ に対して
\[
 B=B(\mathcal C,\mathcal O,s),\qquad
 R=R^{\mathrm{as}}(\mathcal C,\mathcal O)
\]
はともに Koszul である。
従ってこの性質に cubic/noncubic の型別仮定は不要である。
\end{proposition}
\begin{proof}
定義より $B\simeq A/(g)$ であり、$A$ は連結 Koszul 環、
$g\in A_2$ は左右正則正規元である。
正則正規二次商が Koszul 性を保つことは
\cite[Chapter 2, Corollary 6.3]{PolishchukPositselski2005}
に含まれる。中心性を必要としないことを明確にするため、
この場合の環変換スペクトル系列による証明を記す。

右 $A$-加群としての分解
\[
 0\longrightarrow A(-2)\xrightarrow{\,g\cdot\,}A
 \longrightarrow B\longrightarrow0
\]
から、$\operatorname{Tor}^A_q(B,k)$ は
$q=0$ で $k$、$q=1$ で $k(-2)$、それ以外で零である。
それぞれ一つの内部次数に集中するので、
正次数の $B$ の作用は零である。
環変換スペクトル系列は
\[
 E^2_{p,q}
 =\operatorname{Tor}^B_p
       \bigl(k,\operatorname{Tor}^A_q(B,k)\bigr)
 \ \Longrightarrow\ \operatorname{Tor}^A_{p+q}(k,k)
\]
となり、$q=0,1$ の二行だけを持つ。
$B$ が Koszul でないと仮定し、
$\operatorname{Tor}^B_p(k,k)$ が内部次数 $p$ に集中しない
最小の $p$ を取る。
連結次数付き環の極小解消では内部次数はホモロジー次数以上なので、
内部次数 $j>p$ の非零元が存在する。
この $E^2_{p,0}$ の元に入る微分はなく、
出る可能性のある微分は
\[
 d_2:E^2_{p,0}\longrightarrow E^2_{p-2,1}
\]
だけである。
しかし最小性により右辺は内部次数 $(p-2)+2=p$ に集中する。
従って内部次数 $j$ の元は終端まで残り、
$A$ の Koszul 性に反する。これで $B$ は Koszul となる。
反対環にも同じ議論が適用できる。

$s^2\simeq S^{-1}[2]$ により $R\simeq B^{(2)}$ である。
Koszul 環の Veronese 部分環は再び Koszul である
\cite[Chapter 3, Proposition 2.2(i)]{PolishchukPositselski2005}。
別の具体表示で Zhang twist を用いる場合にも、
次数付き加群圏の同値が線型自由解消を移すので結論は変わらない。
\end{proof}

\begin{lemma}[Noetherian 基底上の三次平坦性判定]
\label{lem:relative-koszul-degree-three}
$S$ を Noetherian スキーム、$\mathcal V$ を有限階数の局所自由層、
$\mathcal J\subset\mathcal V^{\otimes2}$ を局所直和因子とし、
\[
 \mathcal R=T_{\mathcal O_S}(\mathcal V)/(\mathcal J)
\]
と置く。次を仮定する。
\begin{enumerate}
\item $\mathcal R_3$ は有限局所自由である。
\item 全ての幾何ファイバー $\mathcal R_{\bar s}$ は Koszul である。
\end{enumerate}
このとき全ての $\mathcal R_n$ は有限局所自由である。
さらに、$n\ge2$ に対して
\[
 \mathcal K_n=
 \bigcap_{i=1}^{n-1}
 \mathcal V^{\otimes(i-1)}\otimes\mathcal J
       \otimes\mathcal V^{\otimes(n-i-1)}
 \ \subset\mathcal V^{\otimes n},
 \qquad
 \mathcal K_1=\mathcal V,\quad \mathcal K_0=\mathcal O_S
\]
は局所直和因子であり、任意の底変更と両立する。
左右の Koszul 複体は $\mathcal O_S$ の線型解消となり、
各内部次数で局所分裂完全である。
従ってこれらの解消は任意の底変更後にも完全である。
\end{lemma}
\begin{proof}
まず任意の点 $s\in S$ の局所環 $D=\mathcal O_{S,s}$ で考える。
$V=\mathcal V_s$、$J=\mathcal J_s$ とする。
各次数の $R=T_D(V)/(J)$ は有限生成 $D$-加群であり、
$R_1,R_2,R_3$ は自由である。
Koszul 性は体の拡大に関して忠実平坦に検出されるため、
幾何ファイバーの仮定から $R\otimes_D\kappa(s)$ も Koszul である。
従って
\cite[Chapter 6, Theorem 7.3]{PolishchukPositselski2005}
を適用して全ての $R_n$ が自由となる。
各 $\mathcal R_n$ は有限表示なので、全ての茎での自由性から
有限局所自由性が従う。

解消と底変更についても同定理の証明を用いる。
$W=V^{\otimes n}$ の部分加群
\[
 J_i=V^{\otimes(i-1)}\otimes J
       \otimes V^{\otimes(n-i-1)}\quad(1\le i<n)
\]
を考える。
$W/J_i$ は $R_2$ と $V$ のテンソル積である。
$|i-j|=1$ ならば $W/(J_i+J_j)$ は $R_3$ と $V$ の
テンソル積であり、$|i-j|>1$ ならば二つの $R_2$ と $V$ の
テンソル積である。従ってこれらは全て自由である。
剰余体上では Koszul 性によって関係部分空間の格子が分配的である。
\cite[Chapter 6, Proposition 7.2]{PolishchukPositselski2005}
より、$J_i$ が生成する格子は $D$ 上でも分配的であり、
その全ての商は自由である。
同章 Proposition 7.1 により $W$ は有限個の自由直和因子へ分解でき、
各 $J_i$ はそのうちのいくつかの和となる。
この分解において和と交叉は任意のテンソル積で保存される。
これにより $K_n=\bigcap_iJ_i$ は直和因子となり、
任意の底変更と両立する。
構成は有限表示の層の間のものなので、
各茎の分解を近傍へ延長して同じ結論を $S$ 上で得る。

例えば左 Koszul 複体は
\[
 \cdots\longrightarrow
 \mathcal R\otimes\mathcal K_2
 \longrightarrow \mathcal R\otimes\mathcal K_1
 \longrightarrow\mathcal R\longrightarrow\mathcal O_S
 \longrightarrow0
\]
であり、$\mathcal K_i$ の生成次数を $i$ とする。
各内部次数 $m$ では項
$\mathcal R_{m-i}\otimes\mathcal K_i$ からなる有限な局所自由複体である。
その幾何ファイバーは Koszul 環の解消の対応する次数成分なので完全である。
有限な局所自由複体の完全性は幾何ファイバーで検出されるから、
この複体は完全である。
アフィン開上の有界な完全射影複体は分裂するため、
これは局所分裂完全であり、任意の底変更後も完全である。
右 Koszul 複体についても同じ証明を適用する。
\end{proof}


\begin{lemma}[Noetherian 特殊ファイバーを持つ局所変形]
\label{lem:local-commutative-fibre-noetherian}
$(R,\mathfrak m,\kappa)$ を Noetherian 局所環とし、
$C=\bigoplus_{d\ge0}C_d$ を $C_0=R$ の
非負次数付き結合的 $R$-代数とする。
$R$ は中心に入り、全ての $C_d$ は有限自由 $R$-加群とする。
特殊ファイバー $C_\kappa=C\otimes_R\kappa$ が
左右 Noetherian なら、$C$ も左右 Noetherian である。
$C_\kappa$ の可換性は仮定しない。
\end{lemma}
\begin{proof}
まず $R$ が完備な場合を示す。
各斉次次数に $\mathfrak m$-進フィルトレーションを入れる。
有限自由性から、環として
\[
 \operatorname{gr}_{\mathfrak m} C
 \simeq \operatorname{gr}_{\mathfrak m}R
               \otimes_\kappa C_\kappa
 \tag{Complete-gr}\label{eq:complete-local-gr}
\]
となる。右辺は一般には非可換だが、
$\operatorname{gr}_{\mathfrak m}R$ は有限生成の可換
$\kappa$-代数であり、$C_\kappa$ に対して中心に入る。
従って中心変数の Hilbert 基底定理と商への移行により、
右辺は左右 Noetherian である。

任意の斉次右イデアル $I$ に
$I\cap\mathfrak m^pC$ による誘導フィルトレーションを入れる。
$\operatorname{gr}_{\mathfrak m}I$ は上の右 Noetherian 環の
二重斉次右イデアルだから、有限個の二重斉次初期記号で生成される。
それらを $I$ の斉次元 $x_1,\ldots,x_r$ に持ち上げる。
内部次数 $d$ の $x\in I$ について、
初期記号を $x_i$ の初期記号の右線形結合で表し、
係数を $C_{d-\deg x_i}$ に持ち上げて差を取ると
$\mathfrak m$-位数が厳密に増える。
繰り返すことで係数の $\mathfrak m$-位数が無限に増す級数を得る。
各固定次数 $C_{d-\deg x_i}$ は有限自由、
$R$ は完備だから係数級数が収束し、
$\bigcap_p\mathfrak m^pC_d=0$ より
$x=\sum_i x_i b_i$ を得る。
したがって全ての斉次右イデアルは有限生成である。
この議論は各\emph{固定次数}で完備性を使うのであり、
無限直和 $C$ 自体が $\mathfrak m$-進完備であると
仮定していない。

任意の右イデアル $J$ について、
内部次数の増加フィルトレーション
$F_nC=\bigoplus_{d\le n}C_d$ を使う。
$\operatorname{gr}_FJ$ は $C$ の斉次右イデアルで
有限生成だから、その初期記号を $J$ に持ち上げ、
最高内部次数を減らす有限帰納法により $J$ も有限生成となる。
これで右 Noetherian 性を得る。
左イデアルについて全く同じ議論を行えばよい。

一般の Noetherian 局所 $R$ では完備化 $\widehat R$ を取る。
$C_{\widehat R}=C\otimes_R\widehat R$ は全次数有限自由で、
特殊ファイバーも $C_\kappa$ だから、完備の場合より
左右 Noetherian である。
$R\to\widehat R$ は忠実平坦であり、
$C\to C_{\widehat R}$ は中心的な忠実平坦環拡大である。
任意の右イデアルの昇鎖を $C_{\widehat R}$ へ延長すると停止する。
忠実平坦性により元の昇鎖も停止する。
左イデアルについても同じ議論が成り立つ。
\end{proof}




\begin{lemma}[Noetherian 局所基底での相対 \texorpdfstring{$\chi_q$}{chi-q} 持上げ]
\label{lem:local-chi-specialization}
$(R,\mathfrak m,\kappa)$ を有限 Krull 次元の
Noetherian 局所環、$B$ を $B_0=R$ の非負次数付き
右 Noetherian $R$-代数とする。
$R$ は中心に入り、各 $B_n$ は有限自由と仮定する。
自明右加群 $R=B/B_{\ge1}$ は
$P_\bullet\to R$ という斉次射影解消を持ち、
各 $P_j$ は有限個の有限階数斉次自由 $B$-加群の直和であり、
任意の $R$-代数への基底変更後にもこの解消は完全と仮定する。
特殊ファイバー $B_\kappa$ が右 $\chi_q$、
すなわち全ての有限生成次数付き右 $B_\kappa$-加群 $N$ に対し
\[
 \underline{\operatorname{Ext}}_{B_\kappa}^{j}(\kappa,N)
 \quad(0\le j\le q)
\]
が全内部次数を合わせて有限次元であると仮定する。
このとき全ての有限生成次数付き右 $B$-加群 $M$ に対して
\[
 \underline{\operatorname{Ext}}_B^j(R,M)
 \quad(0\le j\le q)
\]
は全内部次数を合わせて有限生成 $R$-加群となる。
左側には反対環の仮定の下で同じ結論が成立する。
\end{lemma}
\begin{proof}
$\dim R$ に関する帰納法で証明する。
最初に、各固定 $j,d$ について
$\underline{\operatorname{Ext}}_B^j(R,M)_d$ は
有限 $R$-加群であることに注意する。
実際、有限生成次数付き $M$ は下に有界で各 $M_n$ は
有限 $R$-加群であり、
$\underline{\operatorname{Hom}}_B(P_j,M)_d$ と
$\underline{\operatorname{Hom}}_B(P_{j+1},M)_d$ は
有限 $R$-加群なので、その部分商も有限である。
同じ理由で固定した $j$ の Ext は内部次数が十分負なら零である。
従って、全内部次数を合わせた有限生成性を示すには
内部次数が十分正で Ext が零であることを示せばよい。

$\dim R=0$ のとき $R$ は Artin 局所環で、
$\mathfrak m^a=0$ となる。
$M$ の有限フィルトレーション
$\mathfrak m^iM/\mathfrak m^{i+1}M$ の各商は
有限生成 $B_\kappa$-加群である。
$\mathfrak m$ が零に作用する任意の $N$ に対して、
$P_\bullet$ の各項が有限斉次自由で、剰余体への
基底変更後にも完全だから、
\[
 \underline{\operatorname{Hom}}_B(P_\bullet,N)
 \cong
 \underline{\operatorname{Hom}}_{B_\kappa}
       (P_\bullet\otimes_R\kappa,N)
\]
が成立する。
特殊ファイバーの $\chi_q$ と長完全 Ext 列によって
$M$ も相対 $\chi_q$ を満たす。

$\dim R>0$ とし、局所パラメータ
$t\in\mathfrak m$ を選んで
$\dim R/(t)<\dim R$ とする。
$t$ は $R$ 上の非零因子である必要はない。
$B$ が右 Noetherian なので、
$t$-power torsion 部分
\[
 T=\{m\in M:\ t^a m=0\text{ for some }a\}
\]
は有限生成であり、ある $b$ について $t^bT=0$。
$T$ の $t$-進有限フィルトレーションの各商は
有限生成次数付き $B/tB$-加群である。
$B/tB$ の全次数は $R/(t)$ 上有限自由、
$P_\bullet/tP_\bullet\to R/(t)$ は
任意の基底変更に完全な相対射影解消であり、
特殊ファイバーは元と同じ $B_\kappa$ である。
また $B/tB$ は右 Noetherian。
従って帰納法により、$T$ に対する
$\underline{\operatorname{Ext}}_B^j(R,T)$
は $0\le j\le q$ で有限生成である。
ここでも $tN=0$ なら
$\underline{\operatorname{Ext}}_B^j(R,N)
\simeq\underline{\operatorname{Ext}}_{B/tB}^j(R/(t),N)$
を相対射影解消によって用いている。

次に $M'=M/T$ と置くと、$t:M'\to M'$ は単射である。
よって
\[
 0\longrightarrow M'\xrightarrow{\,t\,}M'
  \longrightarrow M'/tM'\longrightarrow0
\]
は次数付き短完全列である。
帰納法により、$M'/tM'$ に対する
$\underline{\operatorname{Ext}}_{B/tB}^j(R/(t),M'/tM')$
は全次数を合わせて有限生成となる。
Ext の長完全列により
\[
 \frac{\underline{\operatorname{Ext}}_B^j(R,M')}
      {t\,\underline{\operatorname{Ext}}_B^j(R,M')}
 \ \hookrightarrow\
 \underline{\operatorname{Ext}}_B^j(R,M'/tM')
 \qquad(0\le j\le q)
\]
がある。
右辺は十分大きい内部次数で零である。
左辺の各内部次数成分は前述の有限 $R$-加群なので、
Nakayama の補題から
$\underline{\operatorname{Ext}}_B^j(R,M')_d=0$
が全ての十分大きい $d$ で成立する。
従って $M'$ も相対 $\chi_q$ を満たす。
最後に $0\to T\to M\to M'\to0$ の長完全 Ext 列から
$M$ に対する結論を得る。
左側についても反対環に同じ帰納法を適用する。
\end{proof}



\begin{lemma}[有限型の中心基底変更に対する相対 \texorpdfstring{$\chi_q$}{chi-q}]
\label{lem:chi-finite-type-basechange}
$R$ を Noetherian 可換環、$B_0=R$ を満たす
右 Noetherian な非負次数付き $R$-代数を $B$ とする。
$R$ は中心に入り、各 $B_n$ は有限 projective とする。
自明右加群 $R$ は、各項が有限個の
$E\otimes_R B(-a)$（$E$ は有限 projective $R$-加群）
の直和である射影解消 $P_\bullet$ を持ち、
この解消は任意の可換基底変更後も完全と仮定する。
$B$ が相対 $\chi_q$ を満たすなら、任意の可換有限型
$R$-代数 $R'$ について $B'=B\otimes_RR'$ も
相対 $\chi_q$ を満たす。
従って、任意の有限生成次数付き右 $B'$-加群 $M$ に対し
\[
 \underline{\operatorname{Ext}}_{B'}^j(R',M)
 \quad(0\le j\le q)
\]
は全内部次数を合わせて有限生成 $R'$-加群である。
$R'$ の局所性・被約性・$R$ 上の平坦性は要求しない。
左側にも反対環について同じ結論が成立する。
\end{lemma}
\begin{proof}
まず内部次数零の中心変数 $t$ 一つを加える。
$B[t]$ は右 Noetherian であり、$P_\bullet[t]$ は
$R[t]$ の射影解消である。
有限生成次数付き $B[t]$-加群 $M$ の斉次生成元
$m_1,\ldots,m_r$ を取り、内部次数とは別に
\[
 F_pM=\sum_{a=1}^r\sum_{0\le b\le p}m_a t^bB
 \quad(p\ge0),\qquad F_pM=0\quad(p<0)
\]
という補助フィルトレーションを入れる。
$N=\operatorname{gr}_F M$ は、$B$ の補助次数を零、
$t$ の補助次数を一とした $B[t]$ 上の有限生成二重次数付き加群である。
各補助次数成分 $N_p$ は有限生成次数付き $B$-加群である。

$N$ の $t$-power torsion は Noetherian 性から有限生成であり、
ある $e$ について $t^e$ で殺される。
その斉次生成元の補助次数は有限個なので、この torsion は
十分大きい補助次数で零である。
また $N$ の有限生成性により
$t:N_p\to N_{p+1}$ は十分大きい $p$ で全射である。
従って、ある $p_0$ 以後この射は同型となる。
$N_0,\ldots,N_{p_0}$ に対する $\chi_q$ から、
全ての $p\ge0$ に共通する整数 $d_0$ を選んで
\[
 \underline{\operatorname{Ext}}_B^j(R,N_p)_d=0
 \quad(0\le j\le q,\ d\ge d_0)
 \tag{Chi-aux}\label{eq:chi-auxiliary-uniform}
\]
とできる。ここで $t$ による同型は内部次数を保つ。

固定内部次数 $d$ の複体
\[
 C_d^\bullet=
 \underline{\operatorname{Hom}}_{B[t]}(P_\bullet[t],M)_d
\]
に $F_\bullet M$ からの補助フィルトレーションを入れる。
各 $P_j$ の係数は有限 projective なので
\[
 \operatorname{gr}_{F,p} C_d^\bullet
 \simeq\underline{\operatorname{Hom}}_B(P_\bullet,N_p)_d.
\]
このフィルトレーションは下に有界で、各元は有限の補助次数に属する。
$d\ge d_0$ の $j$-cocycle の最高補助次数の記号は、
(\ref{eq:chi-auxiliary-uniform}) により境界である。
その原像を持ち上げて境界を引くと補助次数が下がる。
これを有限回繰り返せば cocycle は境界となる。
$j=0$ では同じ議論により cocycle 自体が零となる。
従って $\underline{\operatorname{Ext}}_{B[t]}^j(R[t],M)_d=0$
が $0\le j\le q$ と $d\ge d_0$ で成立する。

一方、$M$ は内部次数について下に有界であり、
各固定 $j$ の $P_j$ は有限個の内部次数 shift しか持たない。
従って上の Ext は内部次数について下にも有界である。
各内部次数成分は、$R[t]$ 上有限な Hom 複体の部分商なので有限である。
これで全内部次数を合わせた $R[t]$-有限性が従う。

変数の数による帰納法で $R[\mathbf t]$ への変更を処理する。
最後に $R'=R[\mathbf t]/I$ と表示する。
有限生成 $B'$-加群 $M$ は $B[\mathbf t]$ 上でも有限生成であり、
解消の普遍的完全性から
\[
 \underline{\operatorname{Hom}}_{B'}
       (P_\bullet\otimes_RR',M)
 \simeq
 \underline{\operatorname{Hom}}_{B[\mathbf t]}
       (P_\bullet[\mathbf t],M)
\]
となる。右辺のコホモロジーは $R[\mathbf t]$ 上有限で、
$I$ が零に作用するので $R'$ 上有限である。
\end{proof}

\begin{lemma}[相対 \texorpdfstring{$\chi_1$}{chi1} から AZ (A2) へ]
\label{lem:relative-chi1-implies-AZ-A2}
$R$ を Noetherian 可換環、
$B=\bigoplus_{n\ge0}B_n$ を $B_0=R$ の
右 Noetherian な非負次数付き環とする。
$R$ は中心に入り、各 $B_n$ は有限 projective $R$-加群とする。
全ての有限生成次数付き右 $B$-加群 $M$ について
\[
 \underline{\operatorname{Ext}}_B^j(R,M)
 \quad(j=0,1)
\]
が全内部次数を合わせて有限生成 $R$-加群
であると仮定する（相対 $\chi_1$）。
このとき $\operatorname{qgr}B$ 上で
$(\pi B,(1))$ は Artin--Zhang の
Hom 全射条件 (A2) を満たす。
すなわち任意の全射
$f:\mathcal M\twoheadrightarrow\mathcal N$ について、
十分大きい全ての $d$ で
\[
 \operatorname{Hom}_{\operatorname{qgr}B}
       (\pi B(-d),\mathcal M)
 \longrightarrow
 \operatorname{Hom}_{\operatorname{qgr}B}
       (\pi B(-d),\mathcal N)
\]
は全射となる。
$\operatorname{Hom}$ の $R$-有限性や
任意基底変更への $\chi_1$ の安定性は結論しない。
\end{lemma}
\begin{proof}
有限生成次数付き加群 $M$ の torsion 部分を $\tau M$ と書く。
$B$ は右 Noetherian なので $\tau M$ は有限生成であり、
各斉次生成元は十分大きい正次数部分に殺される。
ゆえに $\tau M$ は高内部次数で零である。
$\pi M\simeq\pi(M/\tau M)$ なので、
以下 $M$ は torsion-free と仮定してよい。

$F_n=B/B_{\ge n}$ と置く。
$F_n$ は右自明加群 $B_i(-i)$
（$0\le i<n$）を商とする有限フィルトレーションを持つ。
$B_i$ は有限 projective $R$-加群なので
\[
 \underline{\operatorname{Ext}}_B^1
       (B_i(-i),M)_d
 \simeq B_i^\vee\otimes_R
      \underline{\operatorname{Ext}}_B^1(R,M)_{d+i}.
\]
相対 $\chi_1$ より右辺は $d\gg0$ で消え、
その消滅下限は $i\ge0$ に依存しない。
フィルトレーションの長完全 Ext 列から、
ある $d_0(M)$ について
\[
 \underline{\operatorname{Ext}}_B^1(F_n,M)_d=0
 \qquad(n\ge1,\ d\ge d_0(M))
 \tag{AZ-trunc}\label{eq:relative-AZ-truncation}
\]
が成り立つ。下限の $n$ に関する一様性が重要である。

$0\to B_{\ge n}\to B\to F_n\to0$ に
$\underline{\operatorname{Hom}}_B(-,M)$ を適用する。
$M$ は torsion-free だから
$\underline{\operatorname{Hom}}_B(F_n,M)=0$ であり、
(\ref{eq:relative-AZ-truncation}) より
\[
 M_d\xrightarrow{\sim}
  \underline{\operatorname{Hom}}_B(B_{\ge n},M)_d
 \qquad(d\ge d_0(M))
\]
を得る。
Gabriel 商の標準的な Hom の有向極限表示から
\[
 \operatorname{Hom}_{\operatorname{QGr}B}
      (\pi B(-d),\pi M)
 \simeq
 \varinjlim_n
     \underline{\operatorname{Hom}}_B(B_{\ge n},M)_d
 \simeq M_d
 \quad(d\gg0).
 \tag{AZ-sat}\label{eq:relative-AZ-saturation}
\]
もとの $M$ に対しては
$(M/\tau M)_d=M_d$（$d\gg0$）と解釈する。

任意の全射 $f:\pi M\twoheadrightarrow\pi N$ を取る。
Gabriel 商の射の表示により、
十分大きい $n$ では $f$ は次数付き射
$M_{\ge n}\to N/\tau N$ によって表される。
その余核は有限生成 torsion 加群なので
高内部次数で零になり、
$M_d\to(N/\tau N)_d$ は $d\gg0$ で全射となる。
(\ref{eq:relative-AZ-saturation}) を両辺に適用すると、
求める Hom の全射性が従う。
\end{proof}

\begin{lemma}[相対 tails の Hom 有限性と平坦基底変更]
\label{lem:chi1-qgr-hom-finite-flat}
補題~\ref{lem:relative-chi1-implies-AZ-A2} の仮定を満たし、
さらに $B$ は次数一で生成されるとする。
このとき任意の $\mathcal M,\mathcal N\in\operatorname{qgr}B$ に対して
$\operatorname{Hom}(\mathcal M,\mathcal N)$ は有限 $R$-加群である。
任意の平坦な有限型可換 $R$-代数 $R'$ に対して自然な射
\[
 R'\otimes_R\operatorname{Hom}_{\operatorname{qgr}B}
      (\mathcal M,\mathcal N)
 \xrightarrow{\ \sim\ }
 \operatorname{Hom}_{\operatorname{qgr}(B\otimes_RR')}
      (\mathcal M_{R'},\mathcal N_{R'})
 \tag{Hom-flat-BC}\label{eq:qgr-hom-flat-basechange}
\]
は同型である。対象自身の $R$-平坦性は仮定しない。
\end{lemma}
\begin{proof}
有限生成次数付き加群 $M$ に対して
$\overline M=M/\tau M$ と置く。
$B$ は右 Noetherian なので $\overline M$ も有限生成である。
固定した内部次数 $d$ に対し、完全列
\[
 0\longrightarrow B_{\ge n+1}\longrightarrow B_{\ge n}
   \longrightarrow B_n(-n)\longrightarrow0
\]
を用いる。$B_n$ は有限 projective なので、$j=0,1$ において
\[
 \underline{\operatorname{Ext}}_B^j(B_n(-n),\overline M)_d
 \simeq B_n^\vee\otimes_R
       \underline{\operatorname{Ext}}_B^j(R,\overline M)_{n+d}.
\]
相対 $\chi_1$ により、右辺は $n+d$ が十分大きければ零である。
よって
\[
 \underline{\operatorname{Hom}}_B(B_{\ge n},\overline M)_d
 \longrightarrow
 \underline{\operatorname{Hom}}_B(B_{\ge n+1},\overline M)_d
\]
は $n\gg0$ で同型となる。
Gabriel 商の Hom 表示より、この安定値は
$\operatorname{Hom}_{\operatorname{QGr}B}(\pi B,\pi M(d))$ である。
$B_{\ge n}$ は有限表示であり、各 $M_a$ は有限 $R$-加群なので、
固定内部次数の上の Hom は有限 $R$-加群となる。
従って構造対象の任意の捻りからの Hom は有限である。

次に $R\to R'$ を平坦な有限型変更とする。
各 $F_n=B/B_{\ge n}$ は有限表示だから、
\[
 \tau M=\varinjlim_n
        \underline{\operatorname{Hom}}_B(F_n,M)
 \ \subset M
\]
という評価による同定は平坦基底変更と可換する。
従って $R'\otimes_R\overline M\simeq
M_{R'}/\tau(M_{R'})$。
同様に $B_{\ge n}$ の有限表示性によって
各内部次数の tail Hom は平坦基底変更と可換する。
フィルター付き極限と tensor を交換すると、
構造対象の捻りからの Hom について
(\ref{eq:qgr-hom-flat-basechange}) が従う。
この基底変更の部分は、$R'$ 上で新たな安定次数を選ぶ必要がない。

最後に有限生成次数付き加群 $N$ の有限斉次自由表示
$F_1\to F_0\to N\to0$ を取る。
$\pi$ の完全性から
\[
 0\longrightarrow\operatorname{Hom}(\pi N,\pi M)
 \longrightarrow\operatorname{Hom}(\pi F_0,\pi M)
 \longrightarrow\operatorname{Hom}(\pi F_1,\pi M)
\]
は完全である。右二項は構造対象の捻りからの Hom の有限直和であり、
$R$-有限かつ平坦基底変更と可換する。
$R$ の Noetherian 性と tensor の完全性から、左項も同じ性質を持つ。
全ての qgr 対象はこのような $M,N$ で表示できる。
\end{proof}

\subsection{局所基底からの有限型基底変更と tails}

可換点の局所環における Noetherian 性は、その局所環自身に限らず、
有限型基底変更後の代数にも保たれる。ここではこの事実と、
$\operatorname{QGr}$ の module-object 基底変更を分離して証明する。
内在的 $\mathfrak H_2$ の相対 heart の認識には、
さらに $\chi_1$、Serre、AZ の比較が必要である。

\begin{lemma}[局所 Noetherian 環の有限型基底変更]
\label{lem:local-finite-type-strong-noetherian}
$R$ を Noetherian 可換環、$B$ を $R$ が中心に入る
左右 Noetherian $R$-代数とする。
任意の可換有限型 $R$-代数 $R'$ について
$B_{R'}=B\otimes_RR'$ は左右 Noetherian である。
$R'$ の冪零元、非既約性、非平坦性に制限を課さない。
また $R'$ の局所化へさらに基底変更しても結論は保たれる。
\end{lemma}
\begin{proof}
$R'=R[t_1,\ldots,t_s]/I$ と表示する。
$B[t_1,\ldots,t_s]$ は、$t_i$ が中心と可換する
通常の多項式環である。
非可換係数環に対しても中心変数の
Hilbert 基底定理は左右それぞれに成立するので、
$B[t_1,\ldots,t_s]$ は左右 Noetherian である。
$B_{R'}$ はこの環の中心イデアル $IB[t_1,\ldots,t_s]$
による商だから左右 Noetherian。
中心元での局所化についても同じ有限生成性の議論が使える。
\end{proof}

\begin{lemma}[Gabriel 局所化と module-object 基底変更]
\label{lem:qgr-module-object-basechange}
$R$ を可換環、$B$ を $B_0=R$ の
非負次数付き $R$-代数で、$R$ は中心に入るものとする。
さらに $B$ は次数一で生成され、$B_1$ は有限 $R$-加群と仮定する。
$R\to R'$ を任意の可換 $R$-代数とし、
$B'=B\otimes_RR'$ と書く。
次数付き右加群全体の Grothendieck 圏を
$\operatorname{Gr}B$、次数 torsion 加群での
Gabriel 商を $\operatorname{QGr}B$ とする。
このとき基底変更に自然で、次数移動を保つ同値
\[
 \operatorname{QGr}(B')
   \simeq \bigl(\operatorname{QGr}B\bigr)_{R'}
 \tag{QGr-BC}\label{eq:qgr-module-objects-basechange}
\]
が存在する。右辺は $\operatorname{QGr}B$ 内の
$R'$-加群対象を表す。
この同値は構造対象 $\pi B'$ を
$R'\otimes_R\pi B$ に対応させる。
両辺の torsion は、それぞれの正次数部分による
通常の torsion として定義する。

さらに $R$ が Noetherian、$B$ が右 Noetherian、
各 $B_n$ が有限 $R$-加群で、
$R'$ が有限型 $R$-代数なら、
両辺は locally Noetherian となり、
その Noetherian 対象を取ることで
\[
 \operatorname{qgr}(B')
   \simeq
    \bigl(\operatorname{QGr}B\bigr)_{R'}^{\mathrm{noeth}}
 \tag{qgr-BC}\label{eq:qgr-noetherian-basechange}
\]
を得る。右辺の Noetherian は加群対象としての意味である。
\end{lemma}
\begin{proof}
$B_{\ge n}$ は $B_n$ から生成される有限生成右イデアルである。
従って次数 torsion は部分対象・商・直和だけでなく拡大にも閉じている。
実際、拡大の商で $x$ の像が $B_{\ge n}$ に殺されるなら、
$xB_{\ge n}$ は有限生成で torsion な部分加群に含まれる。
その有限生成元を同時に殺す $B_{\ge m}$ を取れば
$xB_{\ge n+m}=0$ となる。
よってここで使う torsion は localizing 部分圏をなし、
Gabriel 商とその右随伴が存在する。
同じ生成条件は全ての基底変更後も保たれる。

$B'$ 上の次数付き右加群は、
次数付き右 $B$-加群と、その $B$-作用と可換する
$R'$-作用の組と同じである。したがって
$\operatorname{Gr}(B')\simeq(\operatorname{Gr}B)_{R'}$。
次数 torsion の定義
\[
 \forall m\ \exists n,\quad m B_{\ge n}=0
\]
は、この $R'$-作用を加えても変わらない。
よって torsion 部分圏も対応する。

$Q:\operatorname{Gr}B\to\operatorname{QGr}B$ を
Gabriel 商とし、完全忠実な右随伴を $\omega$ とする。
両関手は $R$-線形だから、
$R'$-加群対象上にも $Q_{R'}$ と
$\omega_{R'}$ を誘導し、同じ随伴を与える。
その counit は元の counit と同じく同型であり、
unit の核と余核は torsion 加群対象である。
従って $Q_{R'}$ は torsion 加群対象を核とする
Gabriel 局所化であり、(\ref{eq:qgr-module-objects-basechange})
が従う。構造対象・次数移動との整合性は構成から明らかである。

有限型 $R'$ については
補題~\ref{lem:local-finite-type-strong-noetherian}
から $B'$ は右 Noetherian であり、各次数は有限 $R'$-加群である。
したがって $\operatorname{Gr}(B')$ は
有限生成次数付き加群を Noetherian 対象とする locally
Noetherian Grothendieck 圏である。
その torsion 商の Noetherian 対象は
$\operatorname{qgr}(B')$ に一致し、
(\ref{eq:qgr-noetherian-basechange}) を得る。
\end{proof}

\begin{corollary}[可換点での相対 tails の基底変更と AZ 生成条件]
\label{cor:quadric-local-tails-basechange-generation}
分類論文で証明された可換点の局所 Noetherian AS4 被覆~\cite[可換点での局所 AS4 被覆]{QuadricIntrinsic2026} の
可換点 $u$ と局所環 $R=\mathcal O_{U^\circ,u}$ を取る。
任意の可換有限型 $R$-代数 $R'$ に対して
$B_{R'}$ は左右 Noetherian であり、
\[
 \operatorname{QGr}(B_{R'})
  \simeq\bigl(\operatorname{QGr}B_R\bigr)_{R'}
\]
は module-object 基底変更を与える。
さらに $\mathcal O_{R'}=\pi B_{R'}$ と次数移動 $(1)$ について
Artin--Zhang の豊富性の生成条件 (A1)、
すなわち任意の $\mathcal M\in\operatorname{qgr}B_{R'}$ と
任意の整数 $N$ に対して、ある $d_i\ge N$ で
\[
 \bigoplus_{i=1}^r\mathcal O_{R'}(-d_i)
     \twoheadrightarrow\mathcal M
\]
となることが成立する。
この系は AZ の Hom 全射条件 (A2)、$\chi_1$、
canonical compact 圏との同値、相対 Serre 条件を
結論しない。
\end{corollary}
\begin{proof}
Noetherian 性と module-object 基底変更は前二補題による。
$\mathcal M$ を有限生成次数付き $B_{R'}$-加群 $M$ の
torsion 商で表示する。
$B_{R'}$ が右 Noetherian なので、
任意の $N$ に対して truncation $M_{\ge N}$ は有限生成である。
$M/M_{\ge N}$ は内部次数が上に有界なので torsion であり、
$\pi M_{\ge N}\simeq\pi M$。
$M_{\ge N}$ の有限個の斉次生成元の次数を $d_i\ge N$
とすると、次数付き全射
$\bigoplus_i B_{R'}(-d_i)\twoheadrightarrow M_{\ge N}$
から所要の全射を得る。
(A2) はこの生成操作だけでは従わない。
既知の Artin--Zhang 理論
\cite[Theorem 4.5]{AZ1994} では適切な
$\chi_1$ が (A2) を得るための十分条件となるが、
今回の相対系の全基底変更で $\chi_1$ を確認したわけではない。
\end{proof}


\begin{corollary}[局所相対 tails の fppf 降下と構造対象の平坦性]
\label{cor:quadric-local-tails-fppf-descent}
系~\ref{cor:quadric-local-tails-basechange-generation} の
局所環 $R$ と切断環 $B_R$ を取る。
$R$ 上有限型の affine scheme に対する
$T=\operatorname{Spec}R'$ への割当て
\[
   T\longmapsto\operatorname{QGr}(B_R\otimes_RR')
\]
は、基底変更を module-object で定義した
fppf 降下を満たす。Noetherian 対象だけの
$\operatorname{qgr}(B_R\otimes_RR')$ についても
対象と射の有効な fppf 降下が成立する。
構造対象 $\mathcal O_{R'}=\pi B_{R'}$ は
$R'$ 上平坦であり、任意の有限型基底変更で
この平坦性が保たれる。

この結果は Grothendieck 圏とその Noetherian heart の
基底変更・降下・構造対象の平坦性についての主張であり、
Hom の有限性、$\operatorname{Ext}$ の消滅、
$\chi_1$/AZ (A2)、canonical compact 導来圏との比較、
相対 Serre 条件や強い $\mathfrak H_2$ の認識を含まない。
\end{corollary}
\begin{proof}
$R'\to R''$ を faithfully flat かつ有限表示とする。
補題~\ref{lem:qgr-module-object-basechange} より
\[
 \operatorname{QGr}(B_{R''})
 \simeq\bigl(\operatorname{QGr}(B_{R'})\bigr)_{R''}.
\]
$R'$-線形 Grothendieck 圏 $\mathcal G=
\operatorname{QGr}(B_{R'})$ における
$R''$-加群対象への自由化
$F=R''\otimes_{R'}-:\mathcal G\to\mathcal G_{R''}$
は、$R''/R'$ の忠実平坦性により完全かつ保守的である。
$F$ は忘却関手を右随伴に持ち、
完全な左随伴に対する Barr--Beck の降下判定から、
$\mathcal G$ は $\mathcal G_{R''}$ 上の
通常の Čech cocycle を備える対象の圏に同値である。
これは上記の相対 $\operatorname{QGr}$ の fppf 降下である。
複数の affine 被覆は有限個をまとめた
faithfully flat な積代数で扱える。

$R''$ は $R'$ 上有限型なので、
前補題により $B_{R''}$ は Noetherian。
Noetherian 対象の基底変更は、有限生成次数付き加群の
tensor により Noetherian のままである。
逆に、降下した対象 $\mathcal M$ の引戻しが Noetherian なら、
$\mathcal M$ の任意の部分対象の昇鎖を引き戻すと停止する。
$F$ は完全かつ保守的なので元の昇鎖も停止する。
よって Noetherian 対象の降下も有効である。

$B_{R',n}$ は全て有限局所自由 $R'$-加群であるから、
$R'$-加群 $N$ に対する
$N\mapsto N\otimes_{R'}B_{R'}$
は次数ごとに完全である。
Gabriel 商 $\pi$ も完全なので
$N\mapsto N\otimes_{R'}\pi B_{R'}$
は完全である。これが構造対象の相対平坦性であり、
同じ議論が全ての基底変更後にも成立する。
\end{proof}


\begin{lemma}[相対深さ三と構造環の全次数 saturation]
\label{lem:relative-depth-three-saturation}
$S$ を Noetherian スキームとし、
$B=\bigoplus_{n\ge0}B_n$ を $B_0=\mathcal O_S$ の
次数一生成の次数付き $\mathcal O_S$-代数とする。
$B_n$ は全て有限局所自由で、
自明右加群 $\mathcal O_S=B/B_{\ge1}$ が、
各項が有限局所自由係数をもつ有限斉次自由加群である
相対射影解消 $P_\bullet$ を持つとする。
この解消は内部次数ごとに局所分裂完全で、
任意の底変更後も完全であると仮定する。
全ての幾何ファイバーで
\[
  \underline{\operatorname{Ext}}_{B_{\bar s}}^j
        (k(\bar s),B_{\bar s})=0\qquad(j=0,1,2)
  \tag{Depth-3}\label{eq:relative-fibre-depth-three}
\]
であるとする。このとき任意の affine $S$-基底変更
$T=\operatorname{Spec}R'$ について、全内部次数で
\[
 \underline{\operatorname{Ext}}_{B_T}^j(R',B_T)=0
 \quad(j=0,1,2).
 \tag{Rel-depth}\label{eq:relative-depth-three}
\]
さらに自然な次数付き射
\[
B_T\ \xrightarrow{\ \sim\ }
 \bigoplus_{d\in\mathbb Z}
 \operatorname{Hom}_{\operatorname{QGr}B_T}
      (\pi B_T,\pi B_T(d))
 \tag{Saturation}\label{eq:relative-full-saturation}
\]
は全次数で同型で、任意の基底変更後にも成立する。
$B_T$ の Noetherian 性は仮定しない。
\end{lemma}
\begin{proof}
固定内部次数 $d$ の cochain 複体
$C_d^\bullet=\underline{\operatorname{Hom}}_B(P_\bullet,B)_d$
の各項は有限局所自由であり、全幾何ファイバーへの
基底変更と可換する。幾何ファイバーでは次数 $0,1,2$
で完全なので、先頭の微分から順に最大階数の単元 minor を消去すると
$C_d^\bullet$ は同じ三次数で局所分裂完全となる。
これは $d$ ごとに示せばよく、
無限個の $d$ に共通する自明化は不要である。
その局所分裂完全性は任意の底変更で保たれ、
(\ref{eq:relative-depth-three}) を得る。

以下 $T$ 上で考え、$F_n=B_T/(B_T)_{\ge n}$ と置く。
$F_n$ は $0\le i<n$ において
$(B_T)_i(-i)$ を商に持つ有限フィルトレーションを持つ。
各係数 $(B_T)_i$ は有限 projective なので、
(\ref{eq:relative-depth-three}) の $j=0,1$
と長完全 Ext 列から
\[
 \underline{\operatorname{Hom}}_{B_T}(F_n,B_T)
 =\underline{\operatorname{Ext}}_{B_T}^1(F_n,B_T)=0
\]
となる。
$0\to(B_T)_{\ge n}\to B_T\to F_n\to0$ によって
\[
 B_T\simeq
 \underline{\operatorname{Hom}}_{B_T}((B_T)_{\ge n},B_T)
 \quad(n\ge1)
\]
は全次数で同型となる。
とくに $\underline{\operatorname{Hom}}_{B_T}(R',B_T)=0$ と
次数一生成性から、$x(B_T)_{\ge n}=0$ なら
$x(B_T)_{\ge n-1}=0$ を帰納し、$B_T$ の
次数 torsion は零である。
したがって Gabriel 商の標準 Hom 極限表示から
\[
 \bigoplus_d\operatorname{Hom}_{\operatorname{QGr}B_T}
          (\pi B_T,\pi B_T(d))
 \simeq\varinjlim_n
 \underline{\operatorname{Hom}}_{B_T}((B_T)_{\ge n},B_T)
 \simeq B_T.
\]
これは次数付き積と可換し、自然な射の逆を与える。
全て次数別の分裂完全性から得たので任意の底変更とも可換する。
\end{proof}

\begin{corollary}[任意の係数加群に対する切断と一次コホモロジー]
\label{cor:relative-depth-three-coefficients}
補題~\ref{lem:relative-depth-three-saturation} の仮定の下で、
任意の affine 底変更 $T=\operatorname{Spec}R'\to S$ と
任意の $R'$-加群 $N$ を取る。
$D=N\otimes_{R'}B_T$、$\mathcal O_T=\pi B_T$ と書くと、
全内部次数で
\[
 \underline{\operatorname{Ext}}_{B_T}^{j}(R',D)=0
 \qquad(0\le j\le2)
\]
であり、全ての $d\in\mathbb Z$ に対して
\begin{align*}
 \operatorname{Hom}_{\operatorname{QGr}B_T}
    (\mathcal O_T,\pi D(d))
       &\simeq N\otimes_{R'}(B_T)_d,\\
 \operatorname{Ext}^{1}_{\operatorname{QGr}B_T}
    (\mathcal O_T,\pi D(d))&=0.
\end{align*}
最初の同型は自然な射が与え、任意の基底変更に整合する。
$R'$、$B_T$ の Noetherian 性や $N$ の有限生成性・平坦性を仮定しない。
\end{corollary}
\begin{proof}
$B=B_T$、$R=R'$ と略記する。
補題の証明で得た複体
$\underline{\operatorname{Hom}}_B(P_\bullet,B)_d$
は次数 $0,1,2$ で局所分裂完全であり、
各 $P_i$ の係数は有限 projective である。従って
\[
 \underline{\operatorname{Hom}}_B(P_\bullet,D)_d
 \simeq N\otimes_R
       \underline{\operatorname{Hom}}_B(P_\bullet,B)_d
\]
も同じ三次数で完全となる。
$F_n=B/B_{\ge n}$ の有限 projective 係数を持つ
自明加群によるフィルトレーションから、
$\underline{\operatorname{Ext}}_B^j(F_n,D)=0$（$0\le j\le2$）である。
補題と同じ Hom 極限表示によって $D$ は saturated となり、
主張の Hom 同型を得る。
また $0\to B_{\ge n}\to B\to F_n\to0$ から
\[
 \underline{\operatorname{Ext}}_B^1(B_{\ge n},D)=0
 \qquad(n\ge1)
\]
も従う。

任意の拡大
$0\to\pi D(d)\to\mathcal E\to\pi B\to0$ を取る。
商の右随伴 $\omega$ の左完全性と saturation により、
$0\to D(d)\to\omega\mathcal E\to B$ は左完全である。
最後の射の像を $J\subset B$ とすると、
$\pi\omega\simeq\mathrm{id}$ と $\pi$ の完全性から $B/J$ は torsion。
従って $1+J$ を殺す tail を取れば、ある $n$ で $B_{\ge n}\subset J$ となる。
$\omega\mathcal E\twoheadrightarrow J$ をこの包含に沿って引き戻した列
\[
 0\longrightarrow D(d)\longrightarrow E_n
       \longrightarrow B_{\ge n}\longrightarrow0
\]
は上の Ext 消滅により分裂する。
完全関手 $\pi$ はこの pullback を保ち、
$\pi B_{\ge n}\simeq\pi B$ であるから、元の拡大も分裂する。
以上は Noetherian 性を用いない直接の Yoneda 拡大の議論である。
Hom 同型の自然性と基底変更との整合性は、
自然な unit と変更後も完全な複体による構成から従う。
\end{proof}

\begin{corollary}[全幾何ファイバーが $\HH$ 型である関係式族]
\label{cor:h2-geometric-family-koszul}
$S$ を Noetherian $k$-スキームとし、
$S\to X_3$ を分類論文にある有限次数反Serre関係式空間の命題~\cite[有限次数の関係式空間]{QuadricIntrinsic2026}
の関係式空間への射とする。
その二次環 $\mathcal R$ の全幾何ファイバーが
$\HH$ 対象の反 Serre 環に次数付き同型であると仮定する。
このとき全ての次数で
\[
 \operatorname{rank}\mathcal R_n=(2n+1)^2
\]
であり、$\mathcal R$ は全次数平坦で相対的に Koszul である。
同じ結論は商スタック $[X_3/\operatorname{PGL}_9]$ への射と、
同命題の正規化された普遍次数付き環にも成り立つ。

特に $\HH$ 点を閉ファイバーとする任意の局所 Artin 基底上の
$X_3$ の変形は、全次数平坦な二次代数の変形を与える。
この結論は Artin 基底の冪零次数を制限しない。
\end{corollary}
\begin{proof}
$X_3$ の定義によって次数一、二、三の成分はそれぞれ階数
$9,25,49$ の局所自由層である。
命題~\ref{prop:all-h2-antiserre-koszul} と
補題~\ref{lem:relative-koszul-degree-three} を適用すると、
全次数の局所自由性と相対 Koszul 性を得る。
階数は幾何ファイバー上の Hilbert 関数から決まる。
商スタックの場合は fppf 局所に基底を選び、
同じ議論を適用すればよい。
次数 $n$ の成分への線束の $n$ 乗による twist は
局所自由性と Koszul 複体の完全性を保つ。
これらの性質は忠実平坦に降下するので結論が従う。
局所 Artin 基底には一つの位相的点しかなく、
その幾何ファイバーに命題を適用すれば最後の主張を得る。
\end{proof}

\begin{corollary}[既存の二次 jet の全次数平坦性]
\label{cor:noncubic-second-order-all-degrees}
分類論文のnoncubic二次 jet の命題~\cite[noncubic二次 jet の持上げ]{QuadricIntrinsic2026} で得られた
$k[t]/(t^3)$ 上の各関係式族 $R_t$ は、
三次成分だけでなく全ての次数で自由であり、
\[
 \operatorname{rank}_{k[t]/(t^3)}(R_t)_n=(2n+1)^2
\]
を満たす。左右の相対 Koszul 解消も存在する。
\end{corollary}
\begin{proof}
同命題の関係部分束は二次成分の自由性を保ち、
Schur complement の等式により三次成分は自由階数 $49$ である。
閉ファイバーは既存の noncubic 反 Serre 環であるから、
直前の系を $S=\operatorname{Spec}k[t]/(t^3)$ に適用する。
\end{proof}

\begin{remark}[全次数平坦性と内在的相対モジュライの区別]
本節により、既存の二次 jet に対して
「四次以上の平坦性も別に検証する」という課題は不要になる。
また、任意の固定次数の Artin 持上げが $X_3$ 内で得られれば、
その持上げの全次数平坦性は自動的である。
一方、$X_3$ 内で次の Artin 段階へ持ち上がること自体は証明していない。
Artin 基底上の noetherian Abel heart は次の小節で構成する。
その次数移動を canonical 反 Serre 関手と同定し、適合平方根、
Artin--Zhang 豊富性、AS 正則被覆を得ることは、この平坦性の結論に含めない。
各幾何ファイバーが $\HH$ 型である条件の有限表示性も、
局所環ごとの Koszul 変形原理からは従わない。
\end{remark}

\subsection{Artin 基底上の Abel 圏への実現}

\begin{lemma}[Artin 基底上での noetherian 性]
\label{lem:artin-noetherian-lift}
$(A,\mathfrak m,k)$ を可換局所 Artin 環とし、
$R$ を $A$ が中心に作用する平坦な $A$-代数とする。
$\bar R=R/\mathfrak mR$ が左右 noetherian ならば、$R$ も左右 noetherian である。
\end{lemma}
\begin{proof}
$\mathfrak m^N=0$ となる $N$ を選ぶ。
平坦性により有限フィルターの随伴次数付き環は
\[
 \operatorname{gr}_{\mathfrak mR}R
 \simeq\operatorname{gr}_{\mathfrak m}A\otimes_k\bar R
\]
である。これは $\bar R$ 上で左右とも有限生成なので左右 noetherian である。
$R$ の右イデアルの昇鎖にこのフィルターを誘導すると、
随伴次数付き右イデアルの昇鎖は停止する。
同じ随伴次数付きイデアルを持つ包含は、有限フィルターの各次数を
順に持ち上げることで等号となる。従って元の昇鎖も停止する。
左イデアルについても同じである。
\end{proof}

\begin{proposition}[Artin 関係式族からの flat Abel 変形]
\label{prop:artin-relations-flat-abelian}
$(A,\mathfrak m,k)$ を可換局所 Artin $k$-代数とし、
$\operatorname{Spec}A\to X_3$ の閉点が $\mathcal H_2$ の反 Serre 環
$\bar R$ を表すとする。対応する二次環を $R$ とする。
このとき $R$ は全次数平坦で左右 noetherian であり、
$\operatorname{qgr}R$ は Lowen--Van den Bergh の意味で
$\operatorname{qgr}\bar R$ の flat Abel 変形である。
構造対象 $\pi R$ は $A$ 上平坦で、閉ファイバーは $\pi\bar R$ である。
従って既存の全次数平坦な二次 jet は、標識対象を持つ flat Abel 変形を与える。
ただし $\pi R(1)$ が相対逆 Serre 移動であること、
適合平方根・相対 AS 被覆の存在はこの命題に含めない。
\end{proposition}
\begin{proof}
全次数平坦性は系~\ref{cor:h2-geometric-family-koszul} による。
$\bar R$ は既存の反 Serre 再構成により左右 noetherian なので、
補題~\ref{lem:artin-noetherian-lift} を適用できる。

次数付き右加群の Grothendieck 圏を $\operatorname{Gr}R$、
次数 torsion 加群の局所化部分圏を $\operatorname{Tors}R$ と書く。
ここで torsion とは各斉次元が十分高次数の元で消されることである。
$R$ は右 noetherian で次数一生成なので、この部分圏は
部分対象・商・拡大・任意直和に閉じる。
従って Gabriel 商
\[
 \pi:\operatorname{Gr}R\longrightarrow\operatorname{QGr}R
 :=\operatorname{Gr}R/\operatorname{Tors}R
\]
は $A$-線形完全関手で、完全充満な右随伴 $\omega$ を持つ。
有限生成加群の像は noetherian であり、$\pi R(i)$ が生成系となるので、
$\operatorname{QGr}R$ は局所 noetherian で、その noetherian 部分が
$\operatorname{qgr}R$ である。

$\pi$ は完全で有限直和を保つため、有限提示 $A$-加群との tensor と可換する。
$R(i)$ は $A$ 上平坦なので、$\pi R(i)$ も $A$ 上平坦である。
$\operatorname{qgr}R$ の各対象はこれらの有限直和の商である。
従って任意の有限提示 $A$-加群 $M$ に対して
$\operatorname{Tor}^A_1(M,-)$ は co-effaceable であり、
\cite[Proposition 3.1, Definition 3.2]{LowenVanDenBergh2006}
により $\operatorname{qgr}R$ は flat Abel 圏である。
特に $\pi R$ は底上平坦である。

閉ファイバーを確認する。
$(\operatorname{QGr}R)_k$ は $\mathfrak m$ の作用が零である対象の full subcategory とする。
$\bar R$-加群を $R$-加群とみなす完全関手は torsion を保ち、
\[
 i:\operatorname{QGr}\bar R\longrightarrow\operatorname{QGr}R
\]
を誘導する。$\bar R$-加群の部分加群と商は依然 $\bar R$-加群であり、
torsion の条件も変わらない。
Serre 商の Hom を、始域の torsion 余核を持つ部分対象と
終域の torsion 部分対象による商の Hom の余極限で計算すると、
両圏の Hom は一致する。従って $i$ は完全充満である。
またその像では $\mathfrak m$ の作用が零である。

逆に $X\in(\operatorname{QGr}R)_k$ を取る。
$\omega$ は $A$-線形なので、$\omega X$ も $\mathfrak m$ で消され、
$\bar R$-加群である。$\pi\omega X\simeq X$ から $X$ は $i$ の像に入る。
よって $i$ は
\[
 \operatorname{QGr}\bar R\simeq(\operatorname{QGr}R)_k
\]
を与える。noetherian 対象へ制限すると
$\operatorname{qgr}\bar R\simeq(\operatorname{qgr}R)_k$ となる。
この制限を明示的に見るには、$X$ を $\pi R(i)$ の有限直和で覆い、
$\omega X$ の有限生成部分加群を選べばよい。
それは $\mathfrak m$ で消され、さらに商を取ることで $X$ の有限生成
$\bar R$-加群による表示を得る。
$k\otimes_A\pi R\simeq\pi\bar R$ も $\pi$ の完全性から従う。
これで flat Abel 変形の定義を満たす。
\end{proof}

% !TeX root = ../nc-quadric-surface-moduli.tex
% 2026-11-21 / v1.34（論文Bへ証明移管）
\section{次数別完備化による Noetherian 性の貼合せ}
\label{sec:formal-noetherian-patching}

\paragraph{主張の範囲.}
以下は可換 Noetherian 基底上の非可換次数付き代数と
正の向きの $\mathbb Z$ 代数に対する、証明済みの環論的補題である。
閉部分上とその開補集合上で実際の代数族が Noetherian であることを仮定する。
幾何学的ファイバーの Noetherian 性だけを仮定する主張ではない。
各次数成分の基底上の有限性を用いるが、平坦性、射影性、二次表示、
成長次数、導来同値は仮定しない。新規性は未確認である。

\subsection{中央変数と次数付き有限生成}

\begin{lemma}[中央変数に対する Hilbert 基底論法]
\label{lem:patch-central-variable}
$A$ を環、$M$ を右 Noetherian $A$ 加群とする。
$t$ を中央変数とすれば $M[t]$ は右 Noetherian $A[t]$ 加群である。
内部次数を持つ $A,M$ について、Noetherian 性を次数付き部分加群だけに対して
仮定・結論する形も成立する。この場合 $t$ の内部次数を $0$ とする。
同じ結論は、冪等元 $e_i$ を持つ $\mathbb Z$ 代数 $C$ と
右 $C$ 加群の圏にも成立する。この場合、$t$ は頂点を変えず、
$C[t]_{ij}=C_{ij}[t]$ と定める。
\end{lemma}

\begin{proof}
$N\subset M[t]$ を部分加群とする。$n\geq0$ に対して
$L_n\subset M$ を、$N$ に属し $t$ 次数が $n$ 以下である多項式の
$t^n$ の係数全体とする。これは $A$ 部分加群であり、
$t$ 倍により $L_n\subset L_{n+1}$ である。
$M$ が Noetherian なので、この列はある $n_0$ から安定する。
$0\leq n\leq n_0$ に対して $L_n$ の有限生成元を取り、
それぞれを $N$ の $t$ 次数 $n$ 以下の多項式へ持ち上げる。
この有限個の多項式は $N$ を生成する。
実際、$v\in N$ の最高次係数を消すため、$\deg_t v\leq n_0$ のときは
その次数の持上げを用い、$\deg_t v>n_0$ のときは $n_0$ 次の持上げの
適当な $t$ 冪倍を用いる。差の次数が下がるので、次数に関する帰納法が使える。

内部次数を持つ場合、次数付き $N$ から作る $L_n$ は次数付き部分加群なので、
同じ証明が使える。
$\mathbb Z$ 代数の場合も $L_n$ は $M$ の右部分加群である。
有限生成元を各頂点に同次に取り、多項式の持上げも同じ頂点内で選べば、
まったく同じ帰納法が成り立つ。変数を一つずつ付加することで、
有限個の中央変数についても結論を得る。
\end{proof}

\begin{lemma}[次数付き Noetherian 性から通常の Noetherian 性へ]
\label{lem:patch-graded-to-ungraded}
$A=\bigoplus_{d\geq0}A_d$ を非負次数付き代数とし、
$M$ を下に有界な次数付き右 $A$ 加群とする。
$M$ の次数付き部分加群がすべて有限生成なら、$M$ は
次数を忘れても Noetherian である。
\end{lemma}

\begin{proof}
$N\subset M$ を次数付きとは限らない部分加群とする。
内部次数による増大フィルトレーションを入れると、
$\operatorname{gr} N\subset\operatorname{gr} M=M$ は次数付き部分加群である。
その有限個の同次生成元を $u_1,\ldots,u_s\in N$ に持ち上げる。
任意の $x\in N$ の最高内部次数項を $u_\ell$ の右倍で消せる。
差の最高内部次数は真に下がる。$M$ が下に有界であるから
この操作は有限回で終了し、$u_1,\ldots,u_s$ が $N$ を生成する。
\end{proof}

\subsection{次数ごとの完備性による持上げ}

\begin{lemma}[完備基底上の持上げ]
\label{lem:patch-complete-lifting}
$R$ を可換 Noetherian 環、$I=(f_1,\ldots,f_r)\subset R$ をイデアルとし、
$R$ は $I$ 進完備かつ分離的であると仮定する。
$A=\bigoplus_{d\geq0}A_d$ を中央 $R$ 代数とし、各 $A_d$ は有限 $R$ 加群とする。
このとき $A/IA$ が右 Noetherian なら $A$ も右 Noetherian である。

より一般に、$M$ を各成分が有限 $R$ 加群である次数付き右 $A$ 加群とする。
$M/IM$ が次数付き右 $A/IA$ 加群として Noetherian なら、
$M$ も次数付き右 $A$ 加群として Noetherian である。
$M$ が下に有界なら、次数を忘れても Noetherian である。
\end{lemma}

\begin{proof}
各有限 $R$ 加群 $A_d,M_d$ は $I$ 進完備かつ分離的である。
これは有限加群の完備化が $-\otimes_R\widehat R_I$ と一致することと
$R=\widehat R_I$ から従う
\cite[Tag 00MA]{StacksCompletionPatching}。

$\overline A=A/IA$、$\overline M=M/IM$ と置く。
内部次数と $I$ 進次数を持つ対象に対し、自然な全射
\begin{equation}\label{eq:patch-gr-surjections}
 \overline A[t_1,\ldots,t_r]\twoheadrightarrow\operatorname{gr}_I A,
 \qquad
 \overline M[t_1,\ldots,t_r]\twoheadrightarrow\operatorname{gr}_I M
\end{equation}
がある。$t_\alpha$ は内部次数 $0$、$I$ 進次数 $1$ であり、
$f_\alpha$ の初項へ写す。これらが全射である理由は、
$I^p A_d$ と $I^p M_d$ が、次数 $p$ の $f_\alpha$ の単項式を
係数とする和で生成されることである。
ここでは全射だけを用い、平坦性を仮定しないので同型とは主張しない。
補題~\ref{lem:patch-central-variable} より、
$\operatorname{gr}_I M$ は右 $\operatorname{gr}_I A$ 加群として Noetherian である。

$N\subset M$ を次数付き部分加群とする。
$F^pN=N\cap I^pM$ による誘導フィルトレーションを入れれば
$\operatorname{gr}_F N\subset\operatorname{gr}_I M$ は二重次数付き部分加群である。
有限個の二重同次生成元を取り、それぞれを
\[
 u_\ell\in N_{d_\ell}\cap I^{a_\ell}M_{d_\ell}
 \qquad (1\leq\ell\leq s)
\]
に持ち上げる。これらが $N$ を生成することを示す。

$x\in N_d$ を固定する。$I$ 進次数 $p=0,1,\ldots$ ごとに初項を消せるので、
$a_\ell\leq p$ のとき
\[
 c_{\ell,p}\in I^{p-a_\ell}A_{d-d_\ell},
\]
$a_\ell>p$ のとき $c_{\ell,p}=0$ を選んで、
\begin{equation}\label{eq:patch-approximation}
 x-\sum_{p=0}^{q}\sum_{\ell=1}^s u_\ell c_{\ell,p}
 \in N_d\cap I^{q+1}M_d
 \qquad(q\geq0)
\end{equation}
とできる。存在しない負次数成分は $0$ とする。
固定した $\ell$ について $p-a_\ell\to\infty$ であるから、
$c_\ell=\sum_{p\geq0}c_{\ell,p}$ は有限 $R$ 加群
$A_{d-d_\ell}$ 内で収束する。$M_d$ の分離性と
積の $I$ 進連続性から、\eqref{eq:patch-approximation} の極限は
\[
 x=\sum_{\ell=1}^s u_\ell c_\ell
\]
である。したがって $N$ は有限生成である。
この証明は $N$ 全体の完備性を仮定せず、
固定した内部次数の係数加群の完備性だけを使っている。

$M=A$ とすれば、$A$ のすべての次数付き右イデアルが有限生成である。
補題~\ref{lem:patch-graded-to-ungraded} により通常の右 Noetherian 性を得る。
一般の下に有界な $M$ も同様である。
\end{proof}

\begin{remark}[全体の完備化との区別]
一般の $R$ に対して以下で使う代数は
\[
 A^{\mathrm{deg}\wedge}
 :=A\otimes_R\widehat R_I
 =\bigoplus_{d\geq0}(A_d\otimes_R\widehat R_I)
\]
である。これは各内部次数を個別に完備化した次数付き代数であり、
無限個の内部次数にまたがる元を含み得る環
$\varprojlim_n A/I^nA$ と同一視しない。
補題~\ref{lem:patch-complete-lifting} の収束は一つの内部次数内で行われるため、
$A^{\mathrm{deg}\wedge}$ 全体が $I$ 進完備である必要はない。
\end{remark}

\subsection{閉部分と開補集合からの貼合せ}

\begin{theorem}[Noetherian 性の形式的貼合せ]
\label{thm:patch-noetherian}
$R$ を可換 Noetherian 環、$I=(f_1,\ldots,f_r)\subset R$ とする。
$A=\bigoplus_{d\geq0}A_d$ を各 $A_d$ が有限 $R$ 加群である中央 $R$ 代数とする。
次を仮定する。
\begin{enumerate}
 \item $A/IA$ は右 Noetherian である。
 \item 各 $\alpha$ について $A\otimes_R R_{f_\alpha}$ は右 Noetherian である。
\end{enumerate}
このとき $A$ は右 Noetherian である。左についても同じ結論が成り立つ。

より一般に、各成分が有限 $R$ 加群である次数付き右 $A$ 加群 $M$ について、
$M/IM$ とすべての $M_{f_\alpha}$ がそれぞれ次数付き Noetherian なら、
$M$ も次数付き Noetherian である。
$M$ が下に有界なら通常の右加群としても Noetherian である。
\end{theorem}

\begin{proof}
$\widehat R=\widehat R_I$ と置く。
$\widehat R$ は Noetherian かつ $I\widehat R$ 進完備であり、
$\widehat R/I\widehat R=R/I$ である
\cite[Tags 031C, 05GH]{StacksCompletionPatching}。
$A^{\mathrm{deg}\wedge}/IA^{\mathrm{deg}\wedge}=A/IA$ なので、
補題~\ref{lem:patch-complete-lifting} により
$A^{\mathrm{deg}\wedge}$ は右 Noetherian である。

有限直積
\begin{equation}\label{eq:patch-faithful-cover}
 R^{\ast}=\widehat R_I\times\prod_{\alpha=1}^r R_{f_\alpha}
\end{equation}
を取る。完備化の平坦性 \cite[Tag 00MB]{StacksCompletionPatching} と局所化の平坦性により
$R\to R^{\ast}$ は平坦である。
$\operatorname{Spec}\widehat R_I\to\operatorname{Spec} R$ の像は、
$\widehat R_I/I\widehat R_I=R/I$ より $V(I)$ を含む。
残りは $D(f_\alpha)$ で覆われるから、スペクトル上で全射である。
したがって $R\to R^{\ast}$ は忠実平坦である
\cite[Tag 00HQ]{StacksCompletionPatching}。
なお $I=R$ の場合は開被覆だけを用い、$I=0$ の場合は仮定 (1) で結論が出る。

$A\otimes_R R^{\ast}$ は右 Noetherian 環の有限直積である。
$J\subset A$ を任意の右イデアルとすると、平坦性により
$J\otimes_R R^{\ast}$ は $A\otimes_R R^{\ast}$ の右イデアルであり、有限生成である。
その有限個の生成元をテンソルの有限和として書き、
現れる有限個の $J$ の元で生成される右イデアルを $J_0$ とする。
すると
\[
 (J/J_0)\otimes_R R^{\ast}=0.
\]
忠実平坦性より $J=J_0$ である。したがって $A$ は右 Noetherian である。
左の場合は反対代数に適用する。

一般の $M$ については、補題~\ref{lem:patch-complete-lifting} の加群版で
$M^{\mathrm{deg}\wedge}$ の次数付き Noetherian 性を得て、
任意の次数付き部分加群 $N\subset M$ に上と同じ有限生成の降下を適用する。
最後の主張は補題~\ref{lem:patch-graded-to-ungraded} による。
\end{proof}

\begin{corollary}[閉集合だけを指定した形]
$Z\subset\operatorname{Spec} R$ を閉集合とし、$I$ をその被約閉部分の定義イデアルとする。
$A/IA$ が右 Noetherian であり、$\operatorname{Spec} R\setminus Z$ を覆う有限個の
主開集合上で $A$ が右 Noetherian なら、$A$ は右 Noetherian である。
\end{corollary}

\begin{proof}
有限個の主開集合を $D(g_1),\ldots,D(g_s)$ と書く。
$J=(g_1,\ldots,g_s)$ とすれば $\sqrt J=I$ である。
$R$ が Noetherian なので、ある $m$ について $I^m\subset J$ である。
$A/JA$ は中央の冪零イデアル $IA/JA$ を持ち、その商は $A/IA$ である。
各 $(I^pA+JA)/(I^{p+1}A+JA)$ は有限個の中央係数を用いて
$A/IA$ の有限直和の商となるので、右 Noetherian である。
この有限フィルトレーションにより $A/JA$ は右 Noetherian である。
定理~\ref{thm:patch-noetherian} を $J$ に適用する。
\end{proof}

\subsection{Z 代数の各 row への適用}

\begin{definition}
$C$ を、$C_{ij}=0$ for $j<i$、恒等元 $e_i\in C_{ii}$、
積 $C_{ij}\otimes_R C_{jk}\to C_{ik}$ を持つ中央 $R$ 線形 $\mathbb Z$ 代数とする。
右加群は $M=\bigoplus_jM_j$ と作用 $M_j\otimes_R C_{jk}\to M_k$ からなり、
部分加群はこの頂点分解を保つものとする。
$P_i=e_iC$、すなわち $(P_i)_j=C_{ij}$ と置く。
$C$ が右 Noetherian であるとは、すべての $P_i$ が
この右加群圏で Noetherian であることをいう。
これは有限生成右 $C$ 加群がすべて Noetherian であることと同値である。
\end{definition}

\begin{theorem}[$\mathbb Z$ 代数の形式的貼合せ]
\label{thm:patch-z-algebra}
$R$ と $I=(f_1,\ldots,f_r)$ は定理~\ref{thm:patch-noetherian} の通りとし、
$C$ を各 $C_{ij}$ が有限 $R$ 加群である中央 $R$ 線形 $\mathbb Z$ 代数とする。
$C/IC$ と各 $C_{f_\alpha}$ が右 Noetherian なら、$C$ は右 Noetherian である。
左 Noetherian 性、すなわち各 column $Ce_j$ の Noetherian 性についても同様である。

各 $M_j$ が有限 $R$ 加群である右 $C$ 加群 $M$ に対しても、
$M/IM$ と各 $M_{f_\alpha}$ の Noetherian 性から $M$ の Noetherian 性が従う。
\end{theorem}

\begin{proof}
まず $R$ が $I$ 進完備な場合を考える。
補題~\ref{lem:patch-central-variable} の $\mathbb Z$ 代数版により、各 row
\[
 \overline P_i[t_1,\ldots,t_r]
 \quad\text{は}\quad
 \overline C[t_1,\ldots,t_r]
 \text{ 上 Noetherian である},
\]
ただし $\overline C=C/IC$、$\overline P_i=e_i\overline C$ である。
全射
\[
 \overline C[t_1,\ldots,t_r]\twoheadrightarrow\operatorname{gr}_I C,
 \qquad
 \overline P_i[t_1,\ldots,t_r]\twoheadrightarrow\operatorname{gr}_I P_i
\]
より、$\operatorname{gr}_I P_i$ は右 Noetherian 加群である。

$N\subset P_i$ を右部分加群とし、
$\operatorname{gr}_F N$ の有限生成元を $I$ 進次数と頂点について同次に選ぶ。
それらを
\[
 u_\ell\in N_{j_\ell}\cap I^{a_\ell}C_{i,j_\ell}
\]
に持ち上げる。固定した $x\in N_j$ の $I$ 進初項を逐次消す係数は
\[
 c_{\ell,p}\in I^{p-a_\ell}C_{j_\ell,j}
 \quad(p\geq a_\ell)
\]
に取れる。各 $C_{j_\ell,j}$ は有限 $R$ 加群なので完備であり、
係数の級数が収束する。$C_{ij}$ の分離性から
$x=\sum_\ell u_\ell c_\ell$ を得る。
よって $P_i$ のすべての右部分加群が有限生成である。
ここで $i$ ごとに有限生成元を取っており、$i$ に関する一様な生成次数の上界は不要である。

一般の $R$ では、各 $C_{ij}$ を $\widehat R_I$ に基底変換して
上の完備基底上の結論を得る。
次に \eqref{eq:patch-faithful-cover} を使い、各 $N\subset P_i$ の有限生成を
忠実平坦降下する。テンソルの生成元を展開する際に各元を頂点ごとに分ければ、
有限個の頂点同次生成元を降下できる。
したがってすべての row は Noetherian である。

一般の $M$ も、$\overline M[t_1,\ldots,t_r]\twoheadrightarrow\operatorname{gr}_I M$ を用いて同様に証明できる。
左の場合は各 column に同じ議論を適用すればよい。
\end{proof}

\begin{remark}[普通の環の Noetherian 性との区別]
無限個の直交冪等元を持つ $\bigoplus_{i,j}C_{ij}$ 全体を
単位的な Noetherian 環と呼ぶ主張ではない。
必要なのは各 row、各 column、およびその有限直和の商の Noetherian 性である。
正規族による商 $C/(g)$ を用いる場合も、閉部分上・各主開集合上で得る入力は
この row ごとの Noetherian 性でよい。
正規族の交換関係と正則性から $C/(g)$ の結果を $C$ へ持ち上げるための補題を、
次の定理~\ref{thm:patch-normal-family} で与える。
\end{remark}

\subsection{正則正規線族からの row Noetherian 性の持上げ}

\begin{theorem}[正則正規族による $\mathbb Z$ 代数の拡大]
\label{thm:patch-normal-family}
$R$ を可換環とし、$C$ を正の向きの中央 $R$ 線形 $\mathbb Z$ 代数とし、$C_{ii}=R$ とする。
整数 $\delta>0$ と可逆 $R$ 加群
$G_i\subset C_{i,i+\delta}$ $(i\in\mathbb Z)$ が与えられ、ある両側イデアル $K\subset C$ に対して
乗法が全ての $i,j$ で同型
\begin{equation}\label{eq:patch-normal-family-input}
 G_i\otimes_R C_{i+\delta,j}
       \xrightarrow{\sim}K_{ij}
       \xleftarrow{\sim}C_{i,j-\delta}\otimes_R G_{j-\delta}
\end{equation}
を与えると仮定する。
$\overline C=C/K$ が右 Noetherian なら $C$ も右 Noetherian である。
左についても同様である。
特に有限生成右加群、有限生成左加群はそれぞれ Noetherian になる。

この定理では $G_i$ の周期性、無限個の $G_i$ を同時に自明化する基底の開被覆、
各 $C_{ij}$ の $R$ 上の有限性を仮定しない。
\end{theorem}

\begin{proof}
\emph{(a) 頂点の線束 twist による正規族の自明化.}
各 $\delta$ を法とする剰余類で一つの $L_i$ を $R$ と定め、
全ての整数 $i$ について再帰的に
\begin{equation}\label{eq:patch-vertex-lines}
 L_{i+\delta}=G_i^{-1}\otimes_R L_i
\end{equation}
と定める。各固定した $i$ の $L_i$ は有限個の可逆加群の tensor 積なので、
可逆 $R$ 加群として定義される。
次の $\mathbb Z$ 代数を取る：
\[
 C'_{ij}=L_i^{-1}\otimes_R C_{ij}\otimes_R L_j.
\]
右加群の圏の完全同値は、各頂点で
$M_j\mapsto M_j\otimes_R L_j$ とすることで得られる。
この同値は $P_i$ を $L_i\otimes_R P'_i$ へ送るため、
各 row の Noetherian 性を保つ。
また
\[
 G'_i=L_i^{-1}\otimes_R G_i\otimes_R L_{i+\delta}\simeq R.
\]
これは基底環を変更する操作ではなく、頂点ごとの線束 twist である。
無限個の線束の同時 Zariski 自明化を仮定していない。

従って $C'$ を $C$ と書き直し、全ての $G_i$ に生成元
$g_i\in C_{i,i+\delta}$ がある場合を証明すればよい。
\eqref{eq:patch-normal-family-input} により、左からの $g_i$ 倍と
右からの $g_j$ 倍はそれぞれ単射で、同じ正規イデアルを与える。

\emph{(b) 添字を $\delta$ 移す自己同型.}
$a\in C_{ij}$ に対して、正規性と左からの正則性により一意な
$\sigma(a)\in C_{i+\delta,j+\delta}$ があり、
\begin{equation}\label{eq:patch-normal-shift}
 g_i\sigma(a)=a g_j
\end{equation}
を満たす。
左右の同型 \eqref{eq:patch-normal-family-input} から各 $\sigma:C_{ij}\to C_{i+\delta,j+\delta}$ は同型である。
左から $g_i$ を消去すると
$\sigma(ab)=\sigma(a)\sigma(b)$ が従い、$\sigma(e_i)=e_{i+\delta}$ である。
従って $\sigma$ は頂点を $\delta$ だけ移す $R$ 線形自己同型である。
さらに $a=g_i$ を \eqref{eq:patch-normal-shift} に代入して左から消去すれば
$\sigma(g_i)=g_{i+\delta}$ となる。
従って $\sigma$ は $K$ を保ち、$\overline C$ の自己同型も誘導する。

\emph{(c) 正の skew polynomial $\mathbb Z$ 代数の row Hilbert 基底論法.}
$D=\overline C[t;\sigma]$ を
\[
 D_{ij}=\bigoplus_{p\geq0}\overline C_{i,j-p\delta}\,t^p
\]
と定める。ここで $a t^p$ は、係数 $a\in\overline C_{i,j-p\delta}$ に
頂点 $j-p\delta$ から $j$ への形式的な $p$ 個の $t$ を付した記号である。
積は
\[
 (a t^p)(b t^q)=a\sigma^{-p}(b)t^{p+q}
\]
で定める。$\overline C$ が正の向きなので、固定した $i,j$ について
和に現れる $p$ は有限個である。

$N\subset e_iD$ を右部分加群とする。
各 $n\geq0$ に対して $L_n\subset e_i\overline C$ を、
頂点 $\ell$ の成分が
\[
 (L_n)_\ell=
 \left\{a\in\overline C_{i\ell}:
 \begin{array}{l}
 N_{\ell+n\delta}\text{ のある元が }t\text{ 次数 }n\text{ 以下であり、}\\
 \text{その }t^n\text{ の係数が }a\text{ である}
 \end{array}\right\}
\]
となる部分加群として定める。
これは実際に $\overline C$ 部分加群である。
$a\in(L_n)_\ell$ と $c\in\overline C_{\ell m}$ に対し、
$a$ の持上げに右から $\sigma^n(c)$ を掛ければ、最高係数は $ac$ となるためである。
また右から $t$ を掛けることで $L_n\subset L_{n+1}$ を得る。
$e_i\overline C$ は Noetherian なので、この昇鎖はある $n_0$ から安定する。

$0\leq n\leq n_0$ の各 $L_n$ の有限個の頂点同次生成元を取り、
$N$ の $t$ 次数 $n$ 以下の元へ持ち上げる。
これら有限個の元が $N$ を生成する。
実際、$x\in N_j$ の最高 $t$ 次数を $d$ とする。
$d\leq n_0$ なら $L_d$ の生成元を用い、
$d>n_0$ なら $L_{n_0}$ の持上げを右から $t^{d-n_0}$ 倍して用いる。
最高係数の表示に現れる $c\in\overline C_{\ell,j-d\delta}$ に対し
右から掛ける係数を $\sigma^d(c)$ とすれば、最高係数が正しく消える。
差の $t$ 次数が下がるので、有限回の次数帰納法が完了する。
よって $D$ は各 row について右 Noetherian である。

\emph{(d) $K$ 進次数付き代数.}
$g_i^{(p)}=g_i g_{i+\delta}\cdots g_{i+(p-1)\delta}$ と置く。
正規性から
\[
 K^p_{ij}=g_i^{(p)}C_{i+p\delta,j}
          =C_{i,j-p\delta}g_{j-p\delta}^{(p)}.
\]
両側の正則性により、固定した $p$ の右辺を $K^{p+1}_{ij}$ で割ることは
係数を $K_{i,j-p\delta}$ で割ることと一致する。
従って
\[
 \overline C_{i,j-p\delta}\xrightarrow{\sim}
 K^p_{ij}/K^{p+1}_{ij},
 \qquad
 \overline a\longmapsto
 \overline{a g_{j-p\delta}^{(p)}}
\]
を得る。式 \eqref{eq:patch-normal-shift} が積の整合性を与えるので、
\[
 \operatorname{gr}_K C\simeq\overline C[t;\sigma]=D
\]
である。

\emph{(e) 元の row への有限回の持上げ.}
$N\subset P_i$ に $F^pN=N\cap e_iK^p$ を入れる。
$\operatorname{gr}_F N\subset e_i\operatorname{gr}_K C$ は右有限生成である。
その有限個の $K$ 進次数と頂点についての同次生成元を
$u_1,\ldots,u_s\in N$ に持ち上げる。
固定した $x\in N_j$ の初項をこれらの右倍で逐次消去する。
正の向きと $\delta>0$ から
\[
 K^p_{ij}=0\qquad(p\delta>j-i)
\]
なので、この操作は有限回で終了し、$x\in\sum_\ell u_\ell C$ となる。
従って $P_i$ の全ての部分加群が有限生成であり、$C$ は右 Noetherian である。

左の主張は、添字を反転した反対 $\mathbb Z$ 代数
$C^{\mathrm{op,ref}}_{ij}=C_{-j,-i}$ に右の結果を適用して得る。
この変換は正の向きと、左右両方の
\eqref{eq:patch-normal-family-input} を保つ。
\end{proof}

\begin{remark}[偶数環による証明との関係]
二周期の族では、偶数部分から得る通常の次数付き環で
正則正規な次数二の線束を扱う証明も可能である。
定理~\ref{thm:patch-normal-family} はその周期性を仮定しないので、
根を選んでいない geometric cubic foundation の族にも適用できる。
必要な幾何的入力は、相対曲線商 $\overline C$ の row Noetherian 性と、
正則正規族の普遍的同型 \eqref{eq:patch-normal-family-input} である。
この二つをファイバーだけでなく実際の族として証明する必要は残る。
\end{remark}

\subsection{適用時に確認すべき入力}

閉部分 $Z\subset U=\operatorname{Spec} R$ と開補集合 $U\setminus Z$ に分ける適用では、
必要な入力は次の有限個の主張である。
\begin{enumerate}
 \item $R$ は Noetherian であり、すべての構造成分 $A_d$ または $C_{ij}$ は有限 $R$ 加群である。
 \item 閉部分を定義する実際の商環 $R/I$ 上で、$A/IA$ または $C/IC$ が Noetherian である。
 \item 開補集合を覆う有限個の主開集合 $D(f_\alpha)$ 上で、代数全体または各 row が Noetherian である。
\end{enumerate}
この三条件から結論は上の定理で従う。
Noetherian 性を各点の局所環で証明した後、任意のイデアルの生成元が
自動的に一点の近傍で生成を続ける、と推論する必要はない。
形式的貼合せは、閉部分に沿う完備化という一つの環上で全ての次数を扱い、
有限個の開局所化と合わせた忠実平坦被覆で有限生成を降下する。


% !TeX root = ../nc-quadric-moduli.tex
% 2026-11-21 / v1.32（論文Bへ証明移管）
\section{相対 cubic tails の導来認識と効果性}
\label{sec:relative-cubic-tails}

\subsection{三次解消による相対 tails の compact 認識}

以下では $R$ を任意の可換環とする。$R$ の正則性も
Noetherian 性も、まず仮定しない。
$C$ は正の向きの $R$-線形 $\mathbb Z$-代数、すなわち
$C_{ij}=0$（$j<i$）、$C_{ii}=R$ であり、積を
$C_{ij}\otimes_R C_{jk}\to C_{ik}$ と書くものとする。
右加群 $P_i=e_iC$ の成分は $(P_i)_j=C_{ij}$ である。
$S_i$ は頂点 $i$ にだけ $R$ を置く右加群とする。
次数 torsion は、各元が十分遠い頂点への全ての射で
消されるという意味で用いる。

\begin{theorem}[相対 cubic tails の四項 tilting]
\label{thm:relative-cubic-tails-tilting}
$C$ は隣接成分で生成され、全ての $C_{ij}$ は
有限 projective $R$-加群であるとする。
有限 projective $R$-加群 $V_i,R_i$ と可逆 $R$-加群 $W_i$ により
\[
 0\longrightarrow W_i\otimes_RP_{i+4}
 \longrightarrow R_i\otimes_RP_{i+3}
 \longrightarrow V_i\otimes_RP_{i+1}
 \longrightarrow P_i\longrightarrow S_i\longrightarrow0
 \tag{CT-res}\label{eq:cubic-tails-simple-resolution}
\]
が全ての $i$ で完全であると仮定する。さらに
\[
 \operatorname{Ext}^q_{\operatorname{Gr}C}(S_l,P_j)=0
 \quad\text{unless }q=3\text{ and }j=l+4
 \tag{CT-dual}\label{eq:cubic-tails-relative-dual}
\]
とし、$\mathcal L_l=\operatorname{Ext}^3(S_l,P_{l+4})$ は
可逆 $R$-加群であるとする。
$\mathcal O_i=\pi P_i\in\operatorname{QGr}C$ と置く。
このとき次が成立する。
\begin{enumerate}
\item $\mathbf R\omega$ の cohomological dimension は高々二であり、
  $D(\operatorname{QGr}C)$ は $\mathcal O_i$（$i\in\mathbb Z$）で
  compactly generated される。
\item 任意の $i,j$ について
  \[
  \operatorname{Ext}^q_{\operatorname{QGr}C}
       (\mathcal O_i,\mathcal O_j)
   \simeq
   \begin{cases}
     C_{ji},&q=0,\\
     \operatorname{Hom}_R(C_{i,j-4},\mathcal L_{j-4}),&q=2,\ j\ge i+4,\\
     0,&\text{otherwise}.
   \end{cases}
  \tag{CT-coh}\label{eq:cubic-tails-line-cohomology}
  \]
  とくに任意の連続四項の直和
  $T_a=\bigoplus_{i=a}^{a+3}\mathcal O_i$ は compact tilting generator
  である。その通常の自己準同型環を $\Gamma_a$ とすると
  \[
    D(\operatorname{QGr}C)\simeq D(\operatorname{Mod}\Gamma_a),\qquad
    D(\operatorname{QGr}C)^c\simeq\operatorname{Perf}\Gamma_a.
    \tag{CT-Morita}\label{eq:cubic-tails-canonical-morita}
  \]
\end{enumerate}
全ての仮定と有限 projective 係数の解消が任意の $R$-基底変更で
保たれる場合、これらの compact 比較もその基底変更と整合する。
本定理は $C$ の Noetherian 性や coherent heart の存在を結論しない。
\end{theorem}
\begin{proof}
証明は Bondal--Van den Bergh
\cite[\S4.1, Lemma 4.2.2]{BondalVanDenBergh2003}
の torsion 局所化の方法を、有限 projective 係数と
頂点付き加群に対して記述したものである。
同文献の体上の定理を、そのまま任意の $R$ へ適用するのではない。

\emph{有限 truncation と torsion の導来関手。}
$n>i$ に対し
\[
 F_{i,n}=P_i/(P_i)_{\ge n},\qquad J_{i,n}=(P_i)_{\ge n}
\]
と置く。$F_{i,n}$ は $C_{il}\otimes_RS_l$（$i\le l<n$）を
逐次商に持つ有限フィルトレーションを持つ。
係数 $C_{il}$ は有限 projective なので、
(\ref{eq:cubic-tails-simple-resolution}) と horseshoe 構成から
$F_{i,n}$ は有限生成 projective による長さ三の解消を持つ。
従って $F_{i,n}$ は perfect で projective dimension は高々三、
$J_{i,n}$ も perfect で projective dimension は高々二である。
これらからの全ての Ext は filtered colimit と可換する。

隣接成分の有限生成性により次数 torsion 加群は localizing subcategory
$\operatorname{Tors}C$ をなす。
torsion radical $\tau$ は各頂点で
\[
  (\tau M)_i=\varinjlim_{n>i}\operatorname{Hom}_C(F_{i,n},M),\qquad
  (R^q\tau M)_i=\varinjlim_{n>i}\operatorname{Ext}^q_C(F_{i,n},M)
  \tag{CT-tors}\label{eq:cubic-tails-torsion-colimit}
\]
により与えられる。後の等式は injective 解消に前の等式を適用し、
filtered colimit の完全性を使えば従う。
従って $\operatorname{cd}\tau\le3$ であり、$R^q\tau$ は
filtered colimit と可換する。

ここで torsion 加群 $N$ に対し $R^q\tau N=0$（$q>0$）を確認する。
$N$ は有限個の斉次元が生成する torsion 部分加群の filtered colimit
であり、それぞれは頂点添字について上に有界である。
ゆえに $N_j=0$（$j\ge b$）の場合を示せばよい。
$n\ge b$ なら $J_{i,n}$ は頂点 $j\ge n$ にのみ台を持ち、
その自由解消は全項を $P_j$（$j\ge n$）の直和として選べる。
従って $\operatorname{Ext}^p_C(J_{i,n},N)=0$ は全ての $p\ge0$
で成立する。完全列
$0\to J_{i,n}\to P_i\to F_{i,n}\to0$ を用いると
$\operatorname{Ext}^{q}_C(F_{i,n},N)=0$（$q>0$）となり、
(\ref{eq:cubic-tails-torsion-colimit}) により主張が従う。
この議論は $N$ の $R$ 上平坦性を必要としない。

\emph{section functor と unbounded 導来圏。}
$\pi:\operatorname{Gr}C\to\operatorname{QGr}C$ の完全な商関手と
その完全忠実な右随伴 $\omega$ を取り、$Q=\omega\pi$ と置く。
Gabriel 商の Hom 表示に、直前の torsion 加群上の消滅を使うと
\[
 (QM)_i=\varinjlim_n\operatorname{Hom}_C(J_{i,n},M),\qquad
 (R^qQM)_i=\varinjlim_n\operatorname{Ext}^q_C(J_{i,n},M)
 \tag{CT-sec}\label{eq:cubic-tails-section-colimit}
\]
となる。右辺は $q>2$ で零で、filtered colimit と可換する。
この同じ消滅により、$\pi$ は injective 加群を
$\omega$-acyclic な対象へ送る。具体的には $I$ injective に対し
$\pi I$ の injective 解消へ $\omega$ を施すと、得られる複体の
homology は torsion である。この複体へ $RQ$ を適用し、
$RQ$ の torsion 上の消滅と $Q\omega=\omega$ を使えば
元の複体が完全であることを得る。
従って
\[
  R^qQ=R^q\omega\circ\pi,\qquad \operatorname{cd}\omega\le2.
\]
これが \cite[Lemmas 4.1.4--4.1.7]{BondalVanDenBergh2003}
の相対版である。

Grothendieck 圏の K-injective 解消で unbounded $\mathbf R\omega$
を取る。$\pi$ が完全なので $\omega$ は K-injective 複体を
K-injective 複体へ送り、
$\pi\mathbf R\omega\simeq\operatorname{Id}$ となる。
さらに $\omega$ の cohomological dimension が高々二であり、
各 $R^q\omega$ が直和と可換するため、unbounded
$\mathbf R\omega$ も直和と可換する。
ここでは有限 cohomological dimension を使う acyclic-resolution の
構成、または有限幅 $0\le p\le2$ の hypercohomology spectral sequence
を使える。この段階で無限個の cohomological degree の非一様な
極限交換を仮定していない。
同様に $\mathbf R\tau$ も直和と可換し、自然な三角
\[
  \mathbf R\tau X\longrightarrow X\longrightarrow
        \mathbf R\omega\pi X\longrightarrow
 \tag{CT-loc}\label{eq:cubic-tails-derived-localization}
\]
を得る。

$P_i$ は $D(\operatorname{Gr}C)$ の compact projective generator の
一つである。随伴と $\mathbf R\omega$ の直和整合性から
$\mathcal O_i=\pi P_i$ は $D(\operatorname{QGr}C)$ で compact。
全ての $\mathcal O_i$ に右直交する $X$ は、随伴により
$\mathbf R\omega X$ が全ての $P_i$ に右直交するので零である。
$\pi\mathbf R\omega X=X$ より $X=0$ となり、生成性が従う。

\emph{四項窓の Hom と生成性。}
(\ref{eq:cubic-tails-relative-dual}) と $F_{i,n}$ の有限フィルトレーション
から、$\operatorname{Ext}^q_C(F_{i,n},P_j)=0$ は $q=0,1,2$ で成立する。
$q=3$ では $l=j-4$ の逐次商だけが寄与する。従って $n>j-4$ で
\[
 \operatorname{Ext}^3_C(F_{i,n},P_j)\simeq
 \begin{cases}
  \operatorname{Hom}_R(C_{i,j-4},\mathcal L_{j-4}),&j-4\ge i,\\
  0,&j-4<i.
 \end{cases}
\]
これは truncation の遷移写像に関して安定し、自然な同型である。
$0\to J_{i,n}\to P_i\to F_{i,n}\to0$ と
(\ref{eq:cubic-tails-section-colimit}) を合わせると
(\ref{eq:cubic-tails-line-cohomology}) を得る。
とくに四項窓の higher Ext は全て零である。

(\ref{eq:cubic-tails-simple-resolution}) の $S_i$ は torsion なので、
商圏では
\[
 0\to W_i\otimes_R\mathcal O_{i+4}
 \to R_i\otimes_R\mathcal O_{i+3}
 \to V_i\otimes_R\mathcal O_{i+1}
 \to\mathcal O_i\to0
\]
となる。$W_i$ が可逆で、他の係数も有限 projective であることから、
四項窓が生成する thick subcategory は両方向への帰納法で
全ての $\mathcal O_i$ を含む。
従って $T_a$ は compact tilting generator であり、
導来 Morita 定理 \cite[\S\S4.3, 8.2]{Keller1994}
により (\ref{eq:cubic-tails-canonical-morita}) を得る。

最後に導来 torsion kernel が $\operatorname{loc}\{S_i\}$ に一致することを確認する。
任意の $R$-加群を頂点 $i$ に置いた加群は、$R$-自由解消を用いると
$S_i$ の直和と cones から得られる。
各 torsion 加群は、有限個の斉次元が生成する有界な部分加群の
filtered colimit なので、有限の頂点フィルトレーションを経て
$\operatorname{loc}\{S_i\}$ に属する。
また $\mathbf R\tau P_j$ は上の計算から、次数三にある torsion 加群である。
$\mathbf R\tau$ は直和と可換し、$P_j$ が $D(\operatorname{Gr}C)$ を生成するため、
その像全体は $\operatorname{loc}\{S_i\}$ に含まれる。
$\pi X=0$ なら局所化三角から $X\simeq\mathbf R\tau X$ であり、
逆の包含も従う。ここでは無界複体を単に bounded truncation の極限と
同一視する議論を用いていない。

基底変更で解消と dual 条件が保たれる場合、
$-\otimes_R^{\mathbf L}R'$ による graded 加群の導来拡大は
$P_i$ と $S_i$ をそれぞれ変更後の $P_i'$ と $S_i'$ へ送る。
これは係数の projective 性と解消の普遍完全性による。
従って同関手は、直前に同定した kernel による局所化へ降りる。
$T_a$ とその自己準同型環の全ての成分は通常の tensor で引き戻されるので、
この導来左随伴は Morita 表示で
$-\otimes_{\Gamma_a}^{\mathbf L}(\Gamma_a\otimes_RR')$ となる。
心上の零次は graded tensor と商関手から得る通常の module-object 引戻しである。
これが compact 比較と実際の導来基底変更の整合性である。
\end{proof}

\begin{lemma}[有限型 heart の一様な cohomology]
\label{lem:relative-cubic-coherent-cohomology}
前定理で $R$ を Noetherian とする。
さらに $\operatorname{QGr}C$ が locally Noetherian であり、
$\mathcal O_i$ が Noetherian な生成系をなすと仮定する。
任意の Noetherian 対象 $M$ について
\[
 \operatorname{Ext}^q(\mathcal O_i,M)
 \text{ は有限 }R\text{-加群},\qquad
 \operatorname{Ext}^{q}(\mathcal O_i,M)=0
 \quad(q>0,\ i\gg0).
 \tag{CT-finite}\label{eq:cubic-tails-coherent-finiteness}
\]
任意の Noetherian 対象 $M,N$ の Hom も有限 $R$-加群である。
同じ仮定を満たす平坦基底変更の下では、これらの Hom は基底変更と可換する。
cofinal な偶数添字部分列による生成条件も成立するなら、
その二項移動は Artin--Zhang ample である。
\end{lemma}
\begin{proof}
$M$ へ有限直和の $\mathcal O_j$ から全射を取り、
その核へも同様の全射を取る。核が Noetherian なので繰り返せる。
\[
 0\to K_{r+1}\to E_r\to K_r\to0
 \quad(r=0,1,2),\qquad K_0=M
\]
とし、各 $E_r$ は有限個の $\mathcal O_j$ の直和とする。
前定理から $\operatorname{Ext}^{>2}(\mathcal O_i,-)=0$ であり、
各 $\operatorname{Ext}^q(\mathcal O_i,E_r)$ は有限 projective である。
長完全列を $q=2,1,0$ の順に用いれば、
$\operatorname{Ext}^q(\mathcal O_i,M)$ は有限 $R$-加群の
有限回の部分加群・商・拡大として得られる。
ここで部分加群の有限性に $R$ の Noetherian 性を使う。
また式 (\ref{eq:cubic-tails-line-cohomology}) により、
有限個の $E_0,E_1$ に現れる全ての添字より $i$ を大きく選べば、
$\operatorname{Ext}^2(\mathcal O_i,E_0)$ と
$\operatorname{Ext}^2(\mathcal O_i,E_1)$ は同時に零となる。
同じ長完全列から $M$ に対する全ての正次数の Ext が消える。
これは有限個の presentation の添字から得た一つの消滅下限であり、
stalk ごとの下限を貼り合わせているのではない。

$\operatorname{Hom}(M,N)$ は、有限直和 $E_0\twoheadrightarrow M$
により $\operatorname{Hom}(E_0,N)$ の部分加群になるので有限である。
さらに Morita 同値の像
$\mathbf R\operatorname{Hom}(T_a,M)$ は cohomology が
次数 $0,1,2$ にある $D^b(\operatorname{mod}\Gamma_a)$ の対象となる。
$\Gamma_a$ は有限 $R$-加群なので左右 Noetherian。
平坦 $R\to R'$ に対して、第一変数の上に有界な有限生成 projective
$\Gamma_a$-解消と、第二変数の有界な有限加群複体を選ぶ。
各 total degree に有限個しか項が現れない Hom 複体は
$R'$ への平坦 tensor と可換し、cohomology も可換する。
前定理の基底変更整合性により、これが元の Hom の平坦基底変更である。

最後に任意の全射の核に (\ref{eq:cubic-tails-coherent-finiteness}) の
$q=1$ を適用すれば、十分大きな $i$ に対する Hom の全射性、
すなわち AZ (A2) を得る。仮定した cofinal な生成性が (A1) である。
\end{proof}

\subsection{右軌道規約と既存 foundation への適用}

ここで既存の $D_{ij}=\operatorname{Hom}(L_i,L_j)$ と右軌道環
$B$ の規約を明示的に合わせる。正の向きの右加群に使う
$\mathbb Z$-代数は
\[
 C_{ij}=D_{-j,-i}=\operatorname{Hom}(L_{-j},L_{-i})\quad(i\le j)
 \tag{CT-right}\label{eq:cubic-tails-right-convention}
\]
であり、積は射の合成 $f\otimes g\mapsto f\circ g$ である。
これは $D$ の反対代数を添字反転したもの。
もとの $D$ の右 tails と、添字を変えず同一視してはならない。
この規約では $\pi P_i$ が $L_{-i}$ に対応し、
$i=-3,-2,-1,0$ の四項窓が元の foundation を与える。

\begin{proposition}[幾何的 cubic 開の canonical Grothendieck 導来比較]
\label{prop:quadric-canonical-cubic-tails}
命題~\cite[\texttt{prop:quadric-uniform-cubic-hom-open}]{QuadricIntrinsic2026} の有限型開を用いる。
任意の affine 基底変更 $T=\operatorname{Spec}R\to U^\circ$ について
自然な導来 Morita 比較
\[
 D(\operatorname{QGr}B_T)\simeq D(\operatorname{Mod}\Lambda_T),\qquad
 D(\operatorname{QGr}B_T)^c\simeq\operatorname{Perf}\Lambda_T
 \tag{QT-compact}\label{eq:quadric-tails-compact-comparison}
\]
を得る。この比較は構造対象を $F_0$ に、次数移動 $(1)$ を
既存の $\mathsf s$ に対応させ、任意の導来基底変更と整合する。
特に compact 部分は相対 smooth proper であり、
その canonical Serre 関手について
$(1)^2\simeq S^{-1}[2]$ が成立する。
これは Noetherian 性を仮定しない Grothendieck tails の結論である。
Noetherian 基底では、後の定理~\cite[\texttt{thm:quadric-finite-type-effectivity}]{QuadricIntrinsic2026} が
$B_T$ の左右 Noetherian 性を証明するので、
$\operatorname{QGr}B_T=\operatorname{Ind}(\operatorname{qgr}B_T)$
も成立する。
\end{proposition}
\begin{proof}
式 (\ref{eq:cubic-tails-right-convention}) の $C$ を取る。
$D$ の左右の cubic 標準解消は、全幾何 fiber の AS 正則性と
全成分の局所自由性により、内部次数ごとに局所分裂完全である。
左側についても、補題~\cite[\texttt{lem:relative-geometric-cubic-flatness}]{QuadricIntrinsic2026} と
同じ単元 minor の消去を用いればよい。
従って $C$ は前定理の右標準解消を任意基底変更で持つ。
その解消に $\operatorname{Hom}(-,P_j)$ を適用した四項の複体も
各項が有限局所自由である。
全幾何 fiber の AS 双対条件から、局所分裂完全性と
唯一の rank-one cohomology を得る。
よって式 (\ref{eq:cubic-tails-relative-dual}) も任意基底変更で成立し、
前定理が適用できる。

$C$ の偶数頂点部分への制限は
$\operatorname{QGr}C\simeq\operatorname{QGr}B$ を与える。
実際、
\[
 C_{2a,2b}=\operatorname{Hom}(\mathsf s^{-b}F_0,\mathsf s^{-a}F_0)
 \xrightarrow{\ \mathsf s^b\ } B_{b-a}
\]
の下で合成は右軌道積
$x*y=\mathsf s^{\deg y}(x)\circ y$ に一致する。
偶数頂点への restriction は完全で、左随伴 induction の unit は同型。
counit の核と余核は偶数成分が零であるため torsion である。
induction も torsion を torsion へ送る。これは十分長い path が
必ず偶数頂点を通ることから従う。従って Gabriel 商で両者は互いに逆となる。

四項窓 $i=-3,\ldots,0$ の Hom 公式は
\[
 \operatorname{Hom}(\pi P_i,\pi P_j)
     =C_{ji}=\operatorname{Hom}(L_{-i},L_{-j})
\]
であり、積も合成そのものなので、その自己準同型環は
$\Lambda_T=\operatorname{End}(F_0\oplus\cdots\oplus F_3)$ となる。
これで (\ref{eq:quadric-tails-compact-comparison}) を得る。

比較が全ての $\pi P_i$ を $L_{-i}$ へ送ることも確認できる。
向きが許す全 Hom が次数零に集中しているため、
対応する有向 dg 部分圏の非正次数 truncation と $H^0$ への
準同型の zigzag を通して、$P_i\mapsto L_{-i}$ を
直和を保つ導来 Morita functor として実現する。
$C$ の標準解消の像は、非退化な二つの pairing を用いて
命題~\cite[\texttt{prop:quadric-uniform-cubic-hom-open}]{QuadricIntrinsic2026} の
相対 Euler 複体へ可逆線束を tensor したものになるので零である。
具体的には $a=-i-4$ と置くと、その像は
\[
 W_a\otimes L_a\longrightarrow R_{a+1}\otimes L_{a+1}
 \longrightarrow V_{a+3}\otimes L_{a+3}\longrightarrow L_{a+4}
\]
である。重なりの二つの pairing
$W_a\otimes V_a^\vee\simeq R_{a+1}$、
$W_a\otimes R_a^\vee\simeq V_{a+3}$ により、
これは $a$ から始まる Euler 複体に $W_a$ を tensor したものと一致する。
前定理の証明で導来 torsion kernel を $\operatorname{loc}\{S_i\}$ と
同定したので、この functor は kernel 全体を殺して tails へ降りる。
四項 tilting generator 上で恒等な比較を与えるため、
得た functor は前記の Morita 同値そのものである。

偶数頂点の二項移動は $L_{-i}$ の $\mathsf s$ による移動に一致する。
従って $B$ の $(1)$ は $\mathsf s$ に対応する。
全構成は有限 projective 係数と既存の相対 Morita 根からなり、
任意の導来基底変更と整合する。
元の $\Lambda$ の相対 smooth proper 性と
$\mathsf s^2\simeq S^{-1}[2]$ を移せば、残りの主張を得る。
\end{proof}

\begin{corollary}[Noetherian な心上の Hom と豊富性]
\label{cor:quadric-noetherian-base-recognition}
上の affine 基底 $R$ を Noetherian とし、$B_T$ を左右 Noetherian
と仮定する。このとき
$\operatorname{QGr}B_T=\operatorname{Ind}(\operatorname{qgr}B_T)$ であり、
前命題はこの heart に固有の canonical compact 導来圏の比較となる。
一般の coherent 対象間の Hom は有限で、平坦基底変更と可換する。
$\pi B_T$ は基底上平坦かつ exceptional、$(\pi B_T,(1))$ は AZ ample である。
これらの主張は左右 Noetherian 性が保たれる全有限型基底変更で成立する。
特に系~\cite[\texttt{cor:quadric-point-local-noetherian-as4}]{QuadricIntrinsic2026} の
任意の Noetherian 局所基底 $R$ と、その任意の有限型可換 $R$-代数で
成立する。後の定理~\cite[\texttt{thm:quadric-finite-type-effectivity}]{QuadricIntrinsic2026} により、
単一の有限型開全体と任意の Noetherian 基底変更にも適用できる。
\end{corollary}
\begin{proof}
Noetherian 性による heart と module-object 基底変更は
補題~\cite[\texttt{lem:qgr-module-object-basechange}]{QuadricIntrinsic2026} を用いる。
偶数 restriction の下で各 $\pi P_i$ は有限生成 $B_T$-加群から得られる。
従って Noetherian な生成系をなし、
補題~\ref{lem:relative-cubic-coherent-cohomology} が Hom の有限性、
平坦基底変更、一様な AZ (A2) を与える。
十分大きな偶数 truncation の有限生成性が AZ (A1) を与える。
構造対象の平坦性は $B_n$ の局所自由性による。
構造対象の全 derived End が $R[0]$ であることは
式 (\ref{eq:cubic-tails-line-cohomology})、または foundation 比較による。
有限型可換基底変更での Noetherian 性は
補題~\cite[\texttt{lem:local-finite-type-strong-noetherian}]{QuadricIntrinsic2026} により保たれる。
その全てに前命題と cohomology 補題を適用すればよい。
\end{proof}

\begin{lemma}[相対 tails の一様 cohomology から全ての $\chi_q$ へ]
\label{lem:relative-cubic-global-chi}
系~\ref{cor:quadric-noetherian-base-recognition} の Noetherian 基底と
右 Noetherian 環 $B=B_T$ において、任意の有限生成次数付き右加群 $M$ と
任意の $q\ge0$ に対して
\[
   \underline{\operatorname{Ext}}^q_B(R,M)
   \quad\text{は有限 }R\text{-加群である}.
\]
従って右相対 $\chi_q$ が全ての $q$ で成立する。
反対環にも同じ cubic 議論を適用すれば左側も成立する。
この結論は、左右 Noetherian 性がある同じ affine 基底全体での主張である。
\end{lemma}
\begin{proof}
$\tau_B$ を次数 torsion radical、$Q_B=\omega_B\pi_B$ を section functor
とする。$B$ が右 Noetherian なので、有限 truncation
$B/B_{\ge n}$ とその kernel $B_{\ge n}$ は有限生成 graded projective
からなる解消を持つ。解消の長さの有限性は不要で、各項の有限生成性を使う。
従って各 $q$ の Ext は filtered colimit と可換する。
高い tail $B_{\ge n}$ から上に有界な target への Ext が消えることから、
前定理と同じ torsion 加群上の acyclicity と section の極限公式を得る。
偶数 restriction と前定理により
$\operatorname{Ext}^{>2}_{\operatorname{QGr}B}
(\pi B(d),-)=0$ は全ての $d$ で成立するので、
$\operatorname{cd}\omega_B\le2$、$\operatorname{cd}\tau_B\le3$ となる。
従って unbounded 導来局所化三角も成立する。

$H^b_{\mathrm{tor}}(M)=R^b\tau_B M$ と置く。
まず $H^0_{\mathrm{tor}}(M)$ は有限生成 torsion 加群なので、
上に有界であり、全体として有限 $R$-加群である。
$b\ge2$ については
\[
 H^b_{\mathrm{tor}}(M)_d
 \simeq\operatorname{Ext}^{b-1}_{\operatorname{QGr}B}
                 (\pi B,\pi M(d))
\]
である。よって前補題から各成分は有限 $R$-加群であり、
$b>3$ では零、$b=2,3$ でも $d\gg0$ では零となる。

$b=1$ には自然な完全列
\[
 0\to\tau_B M\to M\to Q_BM\to H^1_{\mathrm{tor}}(M)\to0
\]
を用いる。各 $(Q_BM)_d$ は Hom の有限性から有限 $R$-加群である。
さらに有限 free presentation $F_1\to F_0\to M\to0$ を取り、
$K=\ker(F_0\to M)$ と $K_1=\ker(F_1\to K)$ と置く。
両者は有限生成なので、既に得た一様 Serre 消滅から、十分大きな $d$ で
\[
 (Q_BF_1)_d\longrightarrow(Q_BF_0)_d\longrightarrow(Q_BM)_d
 \longrightarrow0
\]
が完全である。構造環の saturation により $Q_BF_r=F_r$ であるから、
$M_d\xrightarrow{\sim}(Q_BM)_d$ が $d\gg0$ で従う。
したがって $H^1_{\mathrm{tor}}(M)$ も内部次数について上に有界である。
以上により、有限個の $b=0,1,2,3$ 全てに共通する上界 $N$ を取れる。

$\pi_BR=0$ と導来随伴から
\[
 \mathbf R\underline{\operatorname{Hom}}_B(R,\mathbf R\omega_B\pi_BM)=0.
\]
導来局所化三角を使うと
\[
 \mathbf R\underline{\operatorname{Hom}}_B(R,M)
 \simeq
 \mathbf R\underline{\operatorname{Hom}}_B(R,\mathbf R\tau_BM)
\]
となる。その有限幅 hypercohomology spectral sequence
\[
 E_2^{a,b}=
 \underline{\operatorname{Ext}}^a_B(R,H^b_{\mathrm{tor}}(M))
 \ \Longrightarrow\ 
 \underline{\operatorname{Ext}}^{a+b}_B(R,M)
\]
を固定 $q=a+b$ で考える。
$R$ の graded free 解消の最初の $q+1$ 項は有限個の shift からなる。
$H^b_{\mathrm{tor}}(M)$ の上界とこれら有限個の shift から、
この対角線の $E_2^{a,b}$ は全て内部次数 $d\gg0$ で零である。
従って目的の Ext も上に有界となる。
一方、$M$ が下に有界であり、同じ有限個の解消項を用いることから
$\underline{\operatorname{Ext}}^q_B(R,M)$ は下にも有界である。
各成分は有限 $R$-加群の複体の cohomology なので有限である。
内部次数の有限区間にしか成分がないことと合わせ、全体の有限性を得る。
\end{proof}

\subsection{相対内在的条件への昇格}

\begin{theorem}[Noetherian な基底上での相対条件の全確認]
\label{thm:quadric-noetherian-relative-promotion}
命題~\ref{prop:quadric-canonical-cubic-tails} の幾何的 cubic 開
$U^\circ$ への affine 射 $\operatorname{Spec}R\to U^\circ$ を取り、
$R$ は Noetherian 可換環、$B_R$ は左右 Noetherian 環であると仮定する。
任意の可換有限型 $R$-代数 $R'$ に対して
\[
 \mathcal C_{R'}=\operatorname{qgr}(B_R\otimes_RR'),
 \qquad\mathcal O_{R'}=\pi(B_R\otimes_RR')
\]
と置く。この系は、有限型 $R$-基底の site 上で
定義~\cite[\texttt{def:relative-abelian-system}]{QuadricIntrinsic2026} および
定義~\cite[\texttt{def:relative-h2}]{QuadricIntrinsic2026} の全ての条件を満たす。
その局所的な平方根は次数移動 $(1)$ であり、
canonical compact 導来圏は
\[
 D(\operatorname{Ind}\mathcal C_{R'})^c
     \simeq\operatorname{Perf}(\Lambda_R\otimes_RR')
\]
である。この比較は module-object 引戻しの導来関手による
任意の有限型基底変更と整合する。

従って $\operatorname{Spec}R$ 自身が元の局所有限型 $k$-site に属する場合、
$(\mathcal C_R,\mathcal O_R)$ は実際に
$\mathfrak H_2(R)$ の対象を与える。
また系~\cite[\texttt{cor:quadric-point-local-noetherian-as4}]{QuadricIntrinsic2026} の任意の
Noetherian 局所基底 $R$ は本定理の仮定を満たす。
その $R$ が元の site に属さない場合にも、有限型 $R$-基底の系に関する
上の全条件は成立する。ただしこの記述だけで
$\mathfrak H_2$ の一般基底への極限保存的な延長を定義したとはしない。
\end{theorem}
\begin{proof}
$R'$ を任意の可換有限型 $R$-代数とし、$B'=B_R\otimes_RR'$ と略記する。
補題~\cite[\texttt{lem:local-finite-type-strong-noetherian}]{QuadricIntrinsic2026} により
$B'$ は左右 Noetherian であり、その各次数は有限 projective $R'$-加群である。
従って $\operatorname{QGr}B'$ は locally Noetherian で、
その Noetherian 部分圏は $\mathcal C_{R'}$、
その Ind 圏は $\operatorname{QGr}B'$ である。
補題~\cite[\texttt{lem:qgr-module-object-basechange}]{QuadricIntrinsic2026} により
module-object 基底変更が成立し、tensor の結合則が coherent な比較を与える。
faithfully flat な有限表示拡大に沿う引戻しは完全かつ保守的であるから、
系~\cite[\texttt{cor:quadric-local-tails-fppf-descent}]{QuadricIntrinsic2026} の証明により
対象・射・完全列が有効に降下する。
補題~\ref{lem:relative-cubic-coherent-cohomology} により、
全ての coherent 対象間の Hom は $R'$ 上有限で平坦基底変更と可換する。
これで Abel 圏の系の全条件を得る。

構造対象の $R'$-平坦性は $B'_n$ の平坦性と $\pi$ の完全性による。
命題~\ref{prop:quadric-canonical-cubic-tails} の四項 tilting generator は
有限生成 $B'$-加群の像からなるため、その各項は $\mathcal C_{R'}$ に属する。
従ってその有限 cones と直和因子である compact 対象は
$D^b(\mathcal C_{R'})$ に含まれる。
同命題が compact generation と canonical compact 圏の比較を与える。
有限 foundation の相対 smooth proper 性と canonical Serre 双対を
この同値によって移すと、相対 Serre 条件が得られる。
構造対象は foundation の一項なので
$\mathbf R\operatorname{Hom}(\mathcal O_{R'},\mathcal O_{R'})\simeq R'[0]$。

ここで非平坦基底変更に対する導来整合性も確認する。
各 tilting summand は、全成分が $R'$-平坦な $C$ の表現可能加群の像なので
module-object の意味で $R'$ 上平坦である。
よって通常の引戻しは、その summand 上で導来引戻しを計算する。
任意の $R'\to R''$ について、tilting generator とその自己準同型環は
それぞれ変更後の generator と $\Lambda_{R'}\otimes_{R'}R''$ に送られる。
従って導来 Morita 表示における実際の引戻しは
$-\otimes_{\Lambda_{R'}}^{\mathbf L}\Lambda_{R''}$ である。
compact 部分へ制限すれば (\cite[\texttt{eq:relative-perfect-basechange}]{QuadricIntrinsic2026}) となる。
heart 上の通常の引戻しが非平坦基底変更でも完全であることは使っていない。

次数移動 $(1)$ は全ての $\operatorname{QGr}B'$ 上の完全自己同値で、
Noetherian 部分と全ての module-object 基底変更を保つ。
命題~\ref{prop:quadric-canonical-cubic-tails} の比較は
$\pi P_i$ を $L_{-i}$ へ送り、二項移動 $i\mapsto i-2$ の
対象と全ての矢を既存の $\mathsf s$ に対応させる。
従って既存の双加群同型
$\mathsf s^2\simeq S_\Lambda^{-1}[2]$ は、
この heart の canonical compact 圏上で
$(1)^2\simeq S^{-1}[2]$ を与える。
十分大きい偶数 truncation の有限生成性が AZ (A1)、
補題~\ref{lem:relative-cubic-coherent-cohomology} が AZ (A2) を与える。

構造対象の全軌道環は
系~\cite[\texttt{cor:quadric-all-basechange-saturated-sections}]{QuadricIntrinsic2026} により、
全次数・積・任意基底変更に関して $B'$ そのものである。
既存の homological AS4 被覆を $R'$ へ引き戻したものは、
$B'$ の Noetherian 性と正則正規元による議論により左右 Noetherian となる。
各次数の局所自由性、左右 Koszul 解消の完全性、
正則正規二次切断と商の平坦性は任意基底変更で保たれる。
従って同じ次数移動に対して (RQ4) が成立する。

最後に (RQ2) を確認する。幾何点 $\bar t$ で
右規約に合わせた $C_{\bar t}$ は三次三次元 AS 正則
$\mathbb Z$-代数である。
定理~\cite[\texttt{thm:all-admissible-inverse}]{QuadricIntrinsic2026} をこの $C_{\bar t}$ に適用すると
$(\operatorname{qgr}C_{\bar t},\pi e_0C_{\bar t})\in\HH$ である。
偶数頂点への restriction は構造対象を保つ exact 同値
\[
 (\operatorname{qgr}C_{\bar t},\pi e_0C_{\bar t})
   \simeq(\operatorname{qgr}B_{\bar t},\pi B_{\bar t})
\]
を与える。従って Gabriel--Krull rank one torsion-free 条件と、
fiber の全 $D^b$ における Hom 有限性・Serre 条件も満たされる。
この fiber は module-object 基底変更から得るものと、
全軌道環の qgr から得るものが一致する。
ここでは compact 部分の相対 Serre 条件のみから
fiber の全 $D^b$ の条件を推論していない。
以上で全ての条件を確認した。
\end{proof}

% !TeX root = ../nc-quadric-surface-moduli.tex
% 2026-11-21 / v1.3
\section{有限 foundation と反 Serre 表示のスタック}
\label{sec:foundation-stack}

この節の有限代数スタックは直ちに構成できる。
一方、$\HH$ の相対モジュライを表す目標スタックは $\HHstack$ と書く。
$\HHstack$ の代数性およびこの有限代数スタックとの同定には、末尾に記す
相対存在・開性・降下の証明がさらに必要である。

\subsection{二つの線形写像によるパラメータ}

$P,Q$ を二次元、$W_1,W_2,W_3$ をそれぞれ二、四、六次元のベクトル空間とする。
箙 $1\overset{2}{\longrightarrow}2\overset{2}{\longrightarrow}3$ の道代数を $H$ と書く。
隣接矢空間は $P,Q$ であり、$H_{13}=P\otimes Q$ である。
次元ベクトル $(2,4,6)$ の $H$-表現は
\[
 F:P\otimes W_1\longrightarrow W_2,\qquad
 G:Q\otimes W_2\longrightarrow W_3
\]
の組で与えられる。従ってそのパラメータ空間は
\[
 \mathcal R_{\mathrm{fnd}}
 =\Hom(P\otimes W_1,W_2)\oplus\Hom(Q\otimes W_2,W_3)
 \simeq\mathbb A^{64}.
\]
表現 $M=(W_1,W_2,W_3;F,G)$ に対応する一頂点拡大を $\Lambda_M$ とする。
具体的には頂点を $0,1,2,3$ とし、
\[
 (\Lambda_M)_{0i}=W_i\ (1\le i\le3),\qquad
 (\Lambda_M)_{ij}=H_{ij}\ (1\le i\le j\le3),\qquad
 (\Lambda_M)_{00}=k
\]
と置く。積は $H$ の道の積と表現 $M$ の作用による。
これは二十四次元の有向代数である。
箙の関係式の有限型スタックの一般的構成は
Abdelgadir--Okawa--Ueda \cite[Definition 3.5, Proposition 3.6]{AOU2014} にある。
ここでは全ての合成則を明示し、結合律による方程式も消去できる。

群
\[
 \mathcal G=\GL(W_1)\times\GL(W_2)\times\GL(W_3)
             \times\GL(P)\times\GL(Q)
\]
を、$g=(g_1,g_2,g_3,a,b)$ に対し
\[
 (F,G)\longmapsto
 \bigl(g_2F(a^{-1}\otimes g_1^{-1}),\
       g_3G(b^{-1}\otimes g_2^{-1})\bigr)
\]
で作用させる。中心トーラス
\[
 K=\left\{
 \left(\lambda_1 I,\lambda_2 I,\lambda_3 I,
       \frac{\lambda_2}{\lambda_1}I,
       \frac{\lambda_3}{\lambda_2}I\right)
 \right\}\simeq\mathbb G_m^3
\]
は全空間で自明に作用する。

\begin{proposition}[有限 foundation の商スタック]
\label{prop:foundation-stack}
商
\[
 \mathfrak F^{\mathrm{lin}}=[\mathcal R_{\mathrm{fnd}}/\mathcal G],\qquad
 \mathfrak F=[\mathcal R_{\mathrm{fnd}}/(\mathcal G/K)]
\]
は有限型の滑らかな Artin スタックである。
前者は上の形の有向代数を頂点とともに記述し、後者は中心 $K$ の rigidification である。
$\dim\mathfrak F=3$ である。
全ての $\HH$ 対象の定理~\cite[\texttt{thm:normal-helix}]{QuadricIntrinsic2026} の foundation は、
その End 代数を通じて $\mathfrak F^{\mathrm{lin}}$ の幾何点を与える。
\end{proposition}

\begin{proof}
$\mathcal R_{\mathrm{fnd}}$ は滑らかな六十四次元アフィン空間であり、
$\dim\mathcal G=4+16+36+4+4=64$ である。
$K$ は滑らかな中心トーラスなので $\mathcal G/K$ も滑らかな代数群である。
従って二つの商は滑らかな Artin スタックであり\cite[Tag 0DLS]{StacksModuli}、
$\dim\mathfrak F=64-(64-3)=3$ である。

頂点を保つ $H$ の自己同型は $P,Q$ 上の基底変更で尽くされる。
したがって $\mathcal G$ の作用は、頂点を保つ $\Lambda_M$ の同型を正確に与える。
対角単元 $\sum_{i=0}^3\lambda_i e_i$ による内部自己同型は、共通倍率で割り
$\lambda_0=1$ とすれば上の $K$ となる。
これは各 foundation 対象の比較同型のスカラーによる変更を、全体の自然スカラーで割ったものに対応する。
従って $K$ の除去は、$E_0$ の標識だけを残して、
$E_1,E_2,E_3$ を保つ比較同型の選択を忘れる rigidification である。
幾何点では、foundation の各項を同型類として保つ Morita 同値の自己同型群が
$\operatorname{Aut}_{\{e_i\}}(\Lambda_M)/K$ となる。
任意の底上では、これらの対象の存在と線束 twist の降下を区別する必要がある。

$\HH$ から得た四項 $(E_0,E_1,E_2,E_3)$ では
\[
 \dim\Hom(E_i,E_j)=
 \begin{pmatrix}
 1&2&4&6\\0&1&2&4\\0&0&1&2\\0&0&0&1
 \end{pmatrix}_{ij}
\]
であり、$\Hom(E_1,E_2)\otimes\Hom(E_2,E_3)\to\Hom(E_1,E_3)$ は同型である。
後者は系~\cite[\texttt{cor:normal-cubic-defect}]{QuadricIntrinsic2026} の奇数頂点間の積の計算による。
よって $E_1,E_2,E_3$ の End 代数が $H$ となり、
$\Hom(E_0,E_i)$ が上の表現 $M$ を与える。
\end{proof}

\begin{remark}
$\mathfrak F$ は foundation の線形化を除いた有限代数のパラメータスタックである。
任意の点の $\operatorname{Perf}(\Lambda_M)$ に、$\HH$ の heart、
相対 Serre 根、豊富性、AS 正則被覆が存在するとは主張していない。
また foundation を保たない全ての pointed 同値を、この群 $\mathcal G/K$ だけで割ったとも主張しない。
一般に有限代数の導来圏の同値と元の heart の exact 同値は区別する。
本稿の体上では構造対象を保つ導来同値の exact 性を
系~\cite[\texttt{cor:pointed-derived-exact}]{QuadricIntrinsic2026} で証明しているが、族での底変更には別途確認を要する。
\end{remark}

\begin{lemma}[有限 foundation 代数の全 Artin 変形]
\label{lem:foundation-artin-deformation}
$R$ を剰余体 $k$ の可換局所 Artin $k$-代数とする。
$\Lambda_M$ の有限平坦な結合代数変形 $\Lambda_R$ は、
四つの頂点冪等元の持ち上げと各頂点成分の基底を選べば、
$\mathcal R_{\mathrm{fnd}}(R)$ の点から得られる。
固定した頂点冪等元系の下での基底選択の変更と代数同型は
$\mathcal G(R)$ の作用で与えられる。
従ってこの有限表示は、通常の有限代数としての全ての Artin 局所変形を捉える。
\end{lemma}
\begin{proof}
$\mathfrak m_R\Lambda_R$ は冪零イデアルなので、剰余代数の完全直交冪等元系
$(e_0,e_1,e_2,e_3)$ を $\Lambda_R$ へ持ち上げられる。
頂点成分 $(\Lambda_R)_{ij}$ は $\Lambda_R$ の $R$-加群としての直和因子であり、
有限平坦である。局所性から自由で、その階数は剰余体上の Hom 次元に等しい。
従って逆向き成分は零、対角成分は単位 $e_i$ によって $R$ と同型である。

積
\[
 (\Lambda_R)_{12}\otimes_R(\Lambda_R)_{23}
 \longrightarrow (\Lambda_R)_{13}
\]
は階数四の自由加群間の射で、剰余体上では同型である。
Nakayama の補題、または行列式の可逆性によりこれは同型となる。
よって頂点 $1,2,3$ の部分代数は、隣接矢空間の基底を選べば
$H\otimes_kR$ である。
残る積は
\[
 F_R:P_R\otimes_R(W_1)_R\longrightarrow(W_2)_R,
 \qquad
 G_R:Q_R\otimes_R(W_2)_R\longrightarrow(W_3)_R
\]
で決まる。$(0,1,3)$ の積は結合律によりこの二写像の合成に強制される。
各五空間の基底変更は定義された $\mathcal G(R)$ の作用に一致する。
なお冪等元系の別の持ち上げは、剰余体上で恒等となる内部共役により結ばれる。
\end{proof}

\begin{remark}[Abel 圏の変形との比較に必要な範囲]
この補題は有限代数の変形についての主張である。
geometric ample helix に対する De Deken--Lowen
\cite[\S4.1, Theorems 4.5 and 5.15]{DeDekenLowen2011} の変形比較を用いる場合には、
その定理が要求する線形性・平坦性・helix の geometric 性と ampleness、
および対象とする冪零底拡大の仮定を確認する必要がある。
本稿の全 $\HH$ の Artin 変形については、これらの確認と
Abel 圏の変形との比較を定理~\ref{thm:artin-foundation-lift} で行う。
この比較によっても、指定された反 Serre 平方根が変形先へ延長することや、
何らかの相対平方根が存在することは自動的には従わない。
従ってこの補題から相対 $\HH$ locus の開性は結論しない。
\end{remark}

\subsection{Hochschild 障害の消滅}

\begin{proposition}
\label{prop:foundation-hochschild}
任意の $M\in\mathcal R_{\mathrm{fnd}}$ に対し
\[
 \operatorname{HH}^n(\Lambda_M)=0\quad(n\ge3),\qquad
 \dim\operatorname{HH}^2(\Lambda_M)
 -\dim\operatorname{HH}^1(\Lambda_M)=3.
\]
従ってこれらの代数の通常の結合代数変形について、三次 Hochschild コホモロジーにある障害は消える。
これは相対 $\HH$ 条件の無障害性とは別の主張である。
\end{proposition}

\begin{remark}[cubic 型の既知計算との関係]
cubic AS 正則 $\mathbb Z$-algebra に由来する非可換二次曲面について、
$HH^{\geq3}=0$ および
$\dim HH^2-\dim HH^1=3$ は
Belmans \cite[Proposition 31, Corollary 32]{Belmans2019}
の既知結果である。本命題の直接証明は、通常の cubic 関係を仮定せず、
有限 foundation の所定の Hom 次元、四頂点の有向性、
$A_{12}\otimes A_{23}\xrightarrow{\sim}A_{13}$ という積の同型を用いる。
従って normal foundation の仮定を満たす noncubic 型にも適用できる。
既知の cubic 計算と、この直接証明の適用範囲を区別する。
\end{remark}

\begin{proof}
$A=\Lambda_M$、$E=\bigoplus_{i=0}^3ke_i\simeq k^4$ とする。
$E$ は分離代数なので、$E$ に相対的な正規化 bar 複体が $\operatorname{HH}^*(A)$ を計算する。
有向性からそのコチェイン項は
\[
 C^n=\bigoplus_{i_0<\cdots<i_n}
 \Hom_k\bigl(A_{i_0i_1}\otimes\cdots\otimes A_{i_{n-1}i_n},
             A_{i_0i_n}\bigr)
\]
であり、$C^0=k^4$、$C^n=0$ ($n>3$) である。
特に
\[
 C^3=\Hom(A_{01}\otimes A_{12}\otimes A_{23},A_{03}).
\]
$C^2$ の $(0,1,3)$ 成分
$\Hom(A_{01}\otimes A_{13},A_{03})$ から $C^3$ への微分を取る。
他の成分を零とすれば、bar 微分の四項のうち残るものは
\[
 (a,b,c)\longmapsto\varphi(a,bc)
\]
だけである。$A_{12}\otimes A_{23}\xrightarrow\sim A_{13}$ なので
この写像は二つの四十八次元空間の同型である。
従って $C^2\to C^3$ は全射で、$\operatorname{HH}^{\ge3}(A)=0$ となる。

各コチェイン項の次元は
\[
 (\dim C^0,\dim C^1,\dim C^2,\dim C^3)
 =(4,80,128,48).
\]
ここで $\dim C^1=2^2+4^2+6^2+2^2+4^2+2^2=80$、
$\dim C^2=16+48+48+16=128$ である。
中心の元は全 $e_i$ と可換なので対角であり、
各 $A_{0i}=W_i\ne0$ と可換であるため四つの対角係数は全て等しい。
従って $\operatorname{HH}^0(A)=Z(A)=k$ である。
Euler 標数を取ると
\[
 1-\dim\operatorname{HH}^1(A)+\dim\operatorname{HH}^2(A)
 =4-80+128-48=4
\]
となり、第二式が従う。
\end{proof}



\subsection{cubic 表示を与える明示的な開部分}

\begin{proposition}
\label{prop:cubic-foundation-open}
$F$ が同型、$G$ が全射となる開部分で $W_2\simeq P\otimes W_1$ と同定し、
$R=\ker G\subset Q\otimes P\otimes W_1$ と置く。
この階数二の部分束の包含を、因子を並べ替えたテンソル
\[
 w_R\in W_1\otimes P\otimes Q\otimes R^*
\]
とみなす。四因子を順に $V_0,V_1,V_2,V_3$ と書き、添字を四を法として取る。
全ての隣接対 $i,i+1$ と非零 $\phi_i\in V_i^*$、
$\phi_{i+1}\in V_{i+1}^*$ に対し
\[
 (\phi_i\otimes\phi_{i+1})\mathbin{\lrcorner}w_R\ne0
 \quad\text{in }V_{i+2}\otimes V_{i+3}
\]
という geometric 条件を課すと、$\operatorname{Gr}(2,Q\otimes P\otimes W_1)$ の
不変開部分 $\operatorname{Gr}(2,8)_{\mathrm{geo}}$ を得る。
対応する foundation の剛体化した表示スタックは
\[
 \mathfrak F_{\mathrm{cub}}
 \simeq[\operatorname{Gr}(2,8)_{\mathrm{geo}}/
         (\operatorname{PGL}_2)^3]
\]
であり、滑らかな有限型 Artin スタック、次元三である。
その幾何点は geometric quintuple を与え、本稿の cubic 逆構成の範囲に入る。
\end{proposition}
\begin{proof}
条件は geometric quintuple の定義そのものである
\cite[Definition 5.3.1]{VdB2011}（\cite[Equation (2.19)]{OU2014} も参照）。
一つの隣接対について縮約が零となる locus を、
二つの双対射影束の積と Grassmannian の上で考える。
これは閉部分で、その Grassmannian への射影は proper である。
従って像は閉であり、四つの像の合併を除けば所要の開部分となる。
この構成は任意の底変更と整合する開条件である。

$F$ の同型を使って $W_2$ の基底を除き、$G$ の商表示を使って
$W_3\simeq(Q\otimes P\otimes W_1)/R$ とすれば、残るデータは
$R\in\operatorname{Gr}(2,8)$ だけである。
作用群は $\operatorname{GL}(W_1)\times\operatorname{GL}(P)\times\operatorname{GL}(Q)$。
三つの中心スカラーは Grassmannian 上で自明であり、
もとの $K$ の剛体化はこの三中心を除くことに一致する。
従って表示した商を得る。次元は $2(8-2)-3\cdot3=3$ である。
幾何点の AS 正則環への対応は \cite[Lemma 5.4.12, Theorem 5.5.10]{VdB2011} と
定理~\cite[\texttt{thm:all-admissible-inverse}]{QuadricIntrinsic2026} による。
\end{proof}

\begin{example}[階数条件だけでは geometric 性を保証しない]
基底を $e_0,e_1$ として
$R=\langle e_0\otimes e_0\otimes e_0,
          e_1\otimes e_1\otimes e_1\rangle$ を取る。
$R$ は二次元だが、対応する四重テンソルは
$e_{0000}+e_{1111}$ となり、最初の二因子への
$e_0^*\otimes e_1^*$ の縮約が零である。
従って $F$ の可逆性と $G$ の全射性だけを cubic の許容条件にはできない。
\end{example}

\subsection{二つの ruling と foundation の内在性}

\begin{proposition}[ruling 対象の内在的認識]
\label{prop:intrinsic-ruling-pair}
$(\mathcal C,\mathcal O)\in\HH$、$t=S^{-1}[2]$ とする。
Gabriel--Krull rank one torsion-free exceptional 対象 $F$ で
\[
 \chi(\mathcal O,t^nF)=(2n+1)(2n+2)\qquad(n\gg0)
\]
を満たすものの同型類は、正確に二つである。
任意の適合表示の二つの次数零生成 MCM 型の tails 像が、この二類を与える。
従ってその非順序対は $(\mathcal C,\mathcal O)$ から内在的に定まる。
\end{proposition}

\begin{proof}
適合表示を一つ選び、二つの MCM 型の tails 像を $X,Y$ と書く。
定理~\cite[\texttt{thm:normal-helix}]{QuadricIntrinsic2026} の評価列により、次数付き加群 $Z=X,Y$ と
$p(Z)$ は自由加群に埋め込まれる。
$B$ の商斜体へ移せば両者の正の rank の和は二なので、各 rank は一である。
従って $Z(-1)\hookrightarrow B^2$ の非零成分 $Z(-1)\to B$ は単射である。
実際、商斜体へ移した核は零であり、$Z$ は自由加群の部分加群なので
この核はもともと零である。
また $Z$ の非零斉次元は単射 $B(-d)\hookrightarrow Z$ を与える。
tails と Gabriel--Krull 商 $q:\mathcal C\to\mathcal C/\mathcal C_1$ に移すと
\[
 q\pi B(-d)\hookrightarrow q\pi Z\hookrightarrow q\pi B(1)
\]
となり、外側は非零単純対象なので $q\pi Z$ も非零単純である。
$\pi Z\hookrightarrow\pi B(1)$ によって torsion-free 性も従う。
exceptionality は同定理の geometric helix から従い、
Hilbert 級数 $2/(1-z)^3$ から指定の Euler 多項式を得る。

命題~\cite[\texttt{prop:first-orbit-object}]{QuadricIntrinsic2026} と同じ順序で
$K_0(\mathcal C)$ を $K_0(\mathbb P^1\times\mathbb P^1)$ と同定する。
$[F]$ の Chern 指標を $r+aH_1+bH_2+cH_1H_2$ と書くと
\[
 \chi(\mathcal O,t^nF)
 =r(2n+1)^2+(a+b)(2n+1)+c,\qquad
 \chi(F,F)=r^2+2rc-2ab.
\]
係数比較と例外性から $r=1$、$a+b=1$、$c=0$、$ab=0$ となる。
従って可能な類は $(a,b)=(1,0),(0,1)$ の二つだけである。

同じ $K_0$ 類を持つ二つの rank one torsion-free exceptional 対象 $L,F$ が
この多項式を満たすなら $\chi(L,F)=1$ である。
もし $\Ext^2(L,F)\ne0$ なら Serre 双対性から $F\to t^{-1}L$ が非零であり、
rank one torsion-free 性により単射となる。
十分大きい $n$ で Hom 次元は Euler 特性と一致する。
この点は命題~\cite[\texttt{prop:first-orbit-object}]{QuadricIntrinsic2026} の証明の自由被覆の議論による。
具体的には任意の $T$ を $\mathcal O$ のシフトの有限直和 $P$ で覆い、
核も同様に $P'$ で覆う。自由シフト間の $\Ext^1$ の消滅から
$\Ext^1(t^{-n}\mathcal O,T)$ は $\Ext^2(t^{-n}\mathcal O,P')$ の部分商であり、
後者は Serre 双対性と負次数切断の消滅で零となる。
$\Ext^2(t^{-n}\mathcal O,T)$ についても $P$ で同じ議論が使える。
従って上の単射は
\[
 (2n+1)(2n+2)\le(2n-1)(2n)
\]
を与えて矛盾する。従って $\Hom(L,F)\ne0$ であり、同じ理由で単射である。
その余核 $T$ は $K_0$ 類が零なので Hilbert 多項式が零である。
ここで有限生成 $B$-加群を $A$-加群と見ると有限自由分解を持つため
Hilbert 関数は最終的に多項式であり、有限次元加群を除いてもその多項式は変わらない。
従って零 Hilbert 多項式を持つ加群は有限次元であり、$T=0$ となる。
よって各 $K_0$ 類に高々一つの同型類しかない。
$X,Y$ は互いに非同型なので二つの類を尽くす。
\end{proof}

\begin{corollary}[指定構造対象から始まる二つの foundation]
\label{cor:two-intrinsic-foundations}
$M=s\mathcal O$ を命題~\cite[\texttt{prop:first-orbit-object}]{QuadricIntrinsic2026} の内在的対象とする。
定理~\cite[\texttt{thm:normal-helix}]{QuadricIntrinsic2026} で作る、添字零が $\mathcal O$ である foundation は、
各項の同型まで正確に二つである。
それぞれ最初の三項が $(\mathcal O,X,M)$ または $(\mathcal O,Y,M)$ となり、
第四項はこれらの三項から内在的に決まる。
pointed 自己同値群はこの二元集合に作用する。
\end{corollary}

\begin{proof}
初めの三項の内在性は直前の命題と第一軌道対象の認識による。
一つの選択 $E_1=X$ または $Y$ を固定する。
既存の full foundation の完全性から
\[
 {}^{\perp}\langle\mathcal O,E_1,M\rangle\simeq D^b(\operatorname{mod}k).
\]
この左直交の exceptional 生成元は、第四項の cohomological shift に限られる。
条件 $\operatorname{RHom}(M,E_3)\simeq k^2[0]$ はその shift を一意に固定する。
従って第四項も一意である。
\end{proof}

\subsection{二つの短完全列による ruling の相対交換}
\label{subsec:foundation-ruling-involution}

二写像表示において、$F:P\otimes W_1\to W_2$ の随伴を
$\alpha_F:W_1\to P^*\otimes W_2$ と書く。
開部分
\[
 U_0=\{(F,G)\in\mathcal R_{\mathrm{fnd}}\mid
           G:Q\otimes W_2\twoheadrightarrow W_3,
           \ \alpha_F:W_1\hookrightarrow P^*\otimes W_2
           \text{ は部分束の包含}\}
\]
を取る。具体的には $G$ の階数六、$\alpha_F$ の階数二という開条件である。
$U_0$ 上では
\[
 R=\ker G,\qquad C=\operatorname{coker}\alpha_F
\]
は階数二、六のベクトル束であり、任意の底変更と両立する。
新しい二写像データを
\begin{align*}
 W'_1&=R,& W'_2&=W_2,& W'_3&=C,&
 P'&=Q^*,& Q'&=P^*,\\
 F'&:Q^*\otimes R\longrightarrow W_2,&&&&&
 G'&:P^*\otimes W_2\twoheadrightarrow C
\end{align*}
で定める。$F'$ は $R\hookrightarrow Q\otimes W_2$ の縮約であり、
$G'$ は商写像である。

\begin{proposition}[相対 ruling 交換の代数性と整合性]
\label{prop:relative-ruling-involution}
上の構成は $[U_0/\mathcal G]$ 上の代数的対合 $\tau$ を与える。
中心 $K\simeq\mathbb G_m^3$ への作用は
\[
 (\lambda_1,\lambda_2,\lambda_3)
 \longmapsto (\lambda_3,\lambda_2,\lambda_1)
\]
であり、従って $\tau$ は rigidification
\[
 \mathfrak F_0=[U_0/(\mathcal G/K)]
\]
へ降下する。二回適用と恒等の自然同型は $S_2$-作用の整合条件を満たす。
従って
\[
 \mathfrak Q_{\mathrm{fnd}}=[\mathfrak F_0/S_2]
\]
は有限型・滑らか・次元三の Artin スタックであり、
$\mathfrak F_0\to\mathfrak Q_{\mathrm{fnd}}$ は representable な有限 étale
二重被覆である。
\end{proposition}

\begin{proof}
商スタック $[U_0/\mathcal G]$ は、底上の五つの所定階数のベクトル束と
二写像 $(F,G)$ の群oidを記述する。
核・余核は $U_0$ 上でベクトル束であり、構成は底変更と同型に関して自然である。
新しい $G'$ は全射であり、$\alpha_{F'}$ は元の包含
$R\hookrightarrow Q\otimes W_2$ と一致する。
従って新しいデータも $U_0$ に属する。

二回適用すると
\[
 W''_1=\ker G'=\operatorname{im}\alpha_F\simeq W_1,
 \qquad
 W''_3=\operatorname{coker}\alpha_{F'}
      =(Q\otimes W_2)/R\simeq W_3,
\]
かつ $P''\simeq P^{**}\simeq P$、$Q''\simeq Q^{**}\simeq Q$ である。
これらの標準同型の下で、$F''$ は $\alpha_F$ の縮約なので $F$ に等しく、
$G''$ は $R$ による商なので $G$ に等しい。

対合の整合性も単なる点ごとの同型に留まらない。
同じデータを、共通の階数四の束 $W_2$ を用いる二つの短完全列
\[
 \begin{split}
 0&\longrightarrow W_1\xrightarrow{\alpha_F}
       P^*\otimes W_2\longrightarrow C\longrightarrow0,\\
 0&\longrightarrow R\longrightarrow
       Q\otimes W_2\xrightarrow G W_3\longrightarrow0
 \end{split}
\]
で記述する。この表示において $\tau$ は二つの列の交換である。
より形式的には、二つの階数二の束 $V_1=P^*$、$V_2=Q$ を導入すれば、
両列はともに $0\to A_i\to V_i\otimes W_2\to B_i\to0$
という同じ型であり、$\tau$ は添字 $1,2$ の交換に一致する。
従って二回の交換は恒等で、標準双対同型を通じた元の表示でも
$\tau(\tau^2\Rightarrow\operatorname{Id})
 = (\tau^2\Rightarrow\operatorname{Id})\tau$
という必要な整合性が成り立つ。

同型 $g=(g_1,g_2,g_3,a,b)$ に対して、$R$ 上の誘導射は
$(b\otimes g_2)|_R$、$C$ 上の誘導射は
$a^{-\mathsf T}\otimes g_2$ からの商写像である。
$P'$ と $Q'$ 上の射はそれぞれ $b^{-\mathsf T}$ と $a^{-\mathsf T}$ である。
これを $K$ の元へ適用すると、$R,W_2,C$ 上のスカラーは順に
$\lambda_3,\lambda_2,\lambda_1$ となる。
よって $\tau$ は $K$ を保ち、rigidification の普遍性によって
$\mathfrak F_0$ に降下し、整合性も降下する。

$\mathfrak F_0$ は既存の滑らかな商スタックの開部分なので滑らかで次元三である。
有限 étale 定群 $S_2$ による商は Artin スタックであり、
商への射は $S_2$-torsor である。
従って主張した有限型性・滑らかさ・被覆の性質が従う。
ここでは $S_2$ は商スタック自体に作用する。
固定した基底空間 $U_0$ 上の線形作用を選んだとは限らないので、
この商を無条件に一つの半直積群による線形商と書き換える必要はない。
\end{proof}

\begin{proposition}[mutation と $\HH$ の二つの ruling との一致]
\label{prop:ruling-involution-mutation}
$U_0$ 上の普遍有限代数の perfect 圏において、
標準 foundation を $(\mathcal O,X,M,Z)$ と書く。
相対 cone を用いて
\[
 Y=\operatorname{Cone}(Q\otimes M\longrightarrow Z)[-1],
 \qquad
 Z'=\operatorname{Cone}(X\longrightarrow P^*\otimes M)
\]
と定めると、$(\mathcal O,Y,M,Z')$ は full strong exceptional foundation である。
その End 代数の二写像データは $\tau(F,G)$ と一致する。
$\HH$ の全ての normal foundation は $U_0$ に属し、
この変換は二つの normalized ruling の交換と一致する。
\end{proposition}

\begin{proof}
cone は $\operatorname{Perf}(\Lambda)$ の標準 dg 増強で取る。
初めの foundation の相対 derived Hom は、前向きには指定されたベクトル束の次数零成分、
逆向きには零であり、各対角は構造層である。
$X,M,Z$ の間の合成
\[
 \operatorname{RHom}(X,M)\otimes\operatorname{RHom}(M,Z)
 \xrightarrow\sim\operatorname{RHom}(X,Z)
\]
が同型なので、第一の三角から
$\operatorname{RHom}(X,Y)=\operatorname{RHom}(Y,X)=0$ を得る。
二つの cone は exceptional pair の mutation であるため、$Y$ と $Z'$ は相対 exceptional である。
同じ二つの三角に Hom を適用すると
\[
 \operatorname{RHom}(Y,M)=Q^*[0],\qquad
 \operatorname{RHom}(M,Z')=P^*[0],\qquad
 \operatorname{RHom}(Y,Z')=(Q^*\otimes P^*)[0]
\]
を得て、最後の同型は $M$ を経由する合成そのものである。
逆向きの三つの derived Hom は零である。

$\mathcal O$ からの derived Hom は
\[
 \operatorname{RHom}(\mathcal O,Y)
      =\operatorname{Cone}(G)[-1]=R[0],\qquad
 \operatorname{RHom}(\mathcal O,Z')
      =\operatorname{Cone}(\alpha_F)=C[0]
\]
である。ここで全射性と部分束条件を使った。
新三項から $\mathcal O$ への derived Hom は、元の逆向き消滅から零である。
従って新 foundation は strong exceptional であり、
二つの三角から $X,Z$ を復元できるので full でもある。
上の Hom 同型の下で合成は縮約 $F'$ と商写像 $G'$ になり、
End 代数は定義した二写像データと一致する。

次に $\HH$ の normal foundation を取る。
記法を $X=\pi X_0$、$M=\pi B(1)$、$Z=\pi X_1(1)$ とする。
$G$ は次数零生成 MCM 加群 $X_1$ の自由被覆を次数一で取った写像なので全射である。
$P=\Hom(X,M)$ は二次元であり、その非零元は rank one torsion-free 対象間の
非零射であるため単射である。
従って coevaluation $X\to P^*\otimes M$ は単射であり、
$\Hom(\mathcal O,-)$ によって $\alpha_F$ の単射性を得る。
よって全ての normal foundation は $U_0$ に入る。

最後に $X_1=pc(X_0)$ なので $p(X_1)=c(X_0)$ である。
評価列を次数一へ移した
\[
 0\longrightarrow\pi c(X_0)
  \longrightarrow \pi B(1)^2
  \longrightarrow\pi X_1(1)\longrightarrow0
\]
により $Y\simeq\pi c(X_0)$ を得る。
また $p(X_0)$ の評価列を同様に移すと
\[
 0\longrightarrow\pi X_0
  \longrightarrow\pi B(1)^2
  \longrightarrow\pi p(X_0)(1)\longrightarrow0
\]
となる。この中射は、$\Hom(\pi X_0,\pi B(1))$ の基底を選べば
上記 coevaluation と一致する。
従って $Z'\simeq\pi p(X_0)(1)$ であり、これは初めの ruling を
$c(X_0)$ に替えて normal-helix の構成を行った第四項そのものである。
\end{proof}

\begin{corollary}[$\HH$ の幾何点群oidの全充満な実現]
\label{cor:h2-point-groupoid-foundation}
$k$ を代数閉体とする。
対象を $\HH$ の $(\mathcal C,\mathcal O)$、射を標識同型付き exact 同値の
自然同型類とする群oidを $\mathcal H_2(k)$ と書く。
二つの normal foundation から得る非順序データは、全充満な関手
\[
 \mathcal H_2(k)\longrightarrow\mathfrak Q_{\mathrm{fnd}}(k)
\]
を与える。従って対象の同値類だけでなく、全ての慣性群も一致する。
これは $\mathfrak Q_{\mathrm{fnd}}(k)$ の full subgroupoid としての実現である。
\end{corollary}

\begin{proof}
命題~\cite[\texttt{prop:first-orbit-object}]{QuadricIntrinsic2026} と
命題~\ref{prop:intrinsic-ruling-pair} により、pointed exact 同値は
$M=s\mathcal O$ と二つの ruling の非順序対を保つ。
第四項も最初の三項から一意に定まるので、
pointed exact 同値は二つの normal foundation を置換する。
必要なら一方に $\tau$ を作用させれば、各項を保つ同値となる。
比較同型を各項で選ぶと、これは頂点を保つ End 代数の同型を与える。
比較同型の変更は各項のスカラーによる内部自己同型であり、
全体の自然スカラーを除けばちょうど $K$ である。
従って $\mathfrak F_0$ と $S_2$ 商への射を得る。

逆に、像の二点の間の $\mathfrak Q_{\mathrm{fnd}}$ の射は、
必要に応じた ruling 交換の後の頂点保存 End 代数同型として表せる。
tilting による標準導来 Morita 同値と上の mutation は、
$\mathcal O$ を標識同型付きで保存する導来同値を与える。
系~\cite[\texttt{cor:pointed-derived-exact}]{QuadricIntrinsic2026} によりこれは元の heart の間で exact である。

全充満性の一意性については、exact 同値の標準 dg 延長を用いる。
tilting 対象の像とその End 代数への作用は、この延長を標準導来 Morita 同値の
自然同型まで決定する。各項の比較同型を変えることはちょうど前述の $K$ であり、
$\mathcal O$ 上の比較同型を保つ自然二自己同型は自明である。
従って射の間の対応は単射でもある。

慣性を明記すると、$x\in\mathfrak F_0(k)$ の像の自己同型群は
\[
 \coprod_{\epsilon\in S_2}
       \operatorname{Isom}_{\mathfrak F_0}(x,\tau^\epsilon x)
\]
に $S_2$-作用から来る合成を入れた群である。
恒等置換の成分は
$\operatorname{Aut}_{\{e_i\}}(\Lambda_x)/K$ であり、
交換成分は二つの ruling を実際に交換する pointed 自己同値がある場合にだけ非空となる。
従って各点に余分な $S_2$ 慣性を一律に付加しているわけではない。
これは全ての pointed 自己同値と正確に対応する。
\end{proof}

\begin{remark}[この相対構成の射程]
この対合は任意の底変更と両立し、二重被覆は有限代数の表示スタック上で定義される。
ただし $\mathfrak Q_{\mathrm{fnd}}$ の全点が $\HH$ であるとは主張していない。
内在的な族からの $M,X,Y$ の局所的存在・底変更・線束 twist を除く降下、
相対 heart、反 Serre 平方根、AZ 豊富性、二次 AS 正則被覆の条件による
$\HH$ 部分関手の有限表示、および相対再構成の証明は別に必要である。
例外対象の $\Ext^1$ と $\Ext^2$ の消滅や有限代数の
$\operatorname{HH}^3$ の消滅だけから、これらの相対条件は結論しない。
\end{remark}

\subsection{反 Serre 環の全次数を二次関係から回復する補題}
\label{subsec:relative-quadratic-veronese}

これまでの九生成元の関係 Grassmann 射は、
二次成分だけを記録するものとして構成した。
二次環の Veronese が再び二次であるという一般事実は
Polishchuk--Positselski~\cite[Chapter 3, Proposition 2.2]{PolishchukPositselski2005}
に既知である。ここでは任意の可換底と twist に適合する
直接証明を記し、有限局所自由な反 Serre 族へ応用する。
以下では、$\HH$ の反 Serre 環がその二次関係で
\emph{全次数にわたり}復元できることを、
非被約基底にも有効な初等的議論で示す。

\begin{lemma}[二次環の偶数 Veronese は相対的に二次]
\label{lem:relative-quadratic-veronese}
$A$ を可換環、$V$ を有限生成射影 $A$-加群、
$I\subset V^{\otimes2}$ を部分加群とし、
$B=T_A(V)/(I)$ とする。
$E=B_2$、$J=\ker(E\otimes_AE\to B_4)$ と置くと、
自然な射は次数付き環の同型
\[
 \frac{T_A(E)}{(J)}
    \xrightarrow{\ \sim\ } B^{(2)}
    :=\bigoplus_{n\ge0}B_{2n}
 \tag{QV}\label{eq:relative-quadratic-veronese}
\]
を与える。さらに次数を保つ $B$ の自己同型 $\alpha$ によって
$a\star b=a\,\alpha^n(b)$（$a\in B_{2n}$）と定義した
twist は、同じ $E$ からの二次環
\[
 T_A(E)/(J_\alpha),\qquad
 J_\alpha=(\operatorname{id}_E\otimes\alpha^{-1})(J)
\]
に同型である。この二次表示は任意の底変更に両立する
（ただし $J$ そのものの底変更には $B_4$ の平坦性を要する）。
\end{lemma}
\begin{proof}
$n\ge2$ に対し、$V^{\otimes2n}$ で隣接二文字ずつを
$n$ 個の block に分ける。
各 block の内部にある $I$ を零にした商は、
tensor 積の右完全性により
\[
 (V^{\otimes2}/I)^{\otimes n}=E^{\otimes n}
\]
である。$B_{2n}$ を得るために追加で零にすべき二次関係は
block の境界を跨ぐ位置の $I$ だけである。
一つの境界を跨ぐ関係の像は、隣接する二つの
$E$-block の間の
\[
 J=\ker(E^{\otimes2}\longrightarrow B_4)
\]
そのものである。実際、$B_4$ は四文字における
左・右の block 内部の $I$ と、
中央二文字にかかる $I$ を全て零にした商である。
従って
\[
 \ker(E^{\otimes n}\to B_{2n})
 =\sum_{j=0}^{n-2}
 E^{\otimes j}\otimes J\otimes
 E^{\otimes(n-j-2)}
\]
となり、式~(\ref{eq:relative-quadratic-veronese}) が従う。
この議論は純粋に tensor 積の商であり、体上の次元計算や
平坦性を使わない。

twist の次数 $n$ における積写像は、
$E^{\otimes n}$ に
\[
 a_1\otimes\cdots\otimes a_n
 \longmapsto
 a_1\otimes\alpha(a_2)\otimes\cdots
      \otimes\alpha^{n-1}(a_n)
\]
を作用させた後の通常の乗法である。
$\alpha\otimes\alpha$ は $J$ を保つので、
上の核の各隣接成分は
$(\operatorname{id}\otimes\alpha^{-1})(J)$ に移る。
よって twist の関係も二次で尽くされる。
底変更する場合には $B$、その Veronese と twist の定義は
tensor 積から直接底変更する。$J$ をその核と同一視して
基底変更するためには、$B_4$ が平坦であることを確認すればよい。
\end{proof}

\begin{corollary}[$\HH$ の反 Serre 環は二次]
\label{cor:h2-intrinsic-antiserre-quadratic}
代数閉体上の任意の $(\mathcal C,\mathcal O)\in\HH$ に対し、
$t=S^{-1}[2]$ の反 Serre 軌道環
\[
 R^{\mathrm{as}}(\mathcal C,\mathcal O)
 =\bigoplus_{n\ge0}\Hom(\mathcal O,t^n\mathcal O)
\]
は、次数一生成であり、二次関係だけで表示される。
生成元は九次元、二次関係は五十六次元で、
\[
 \dim R^{\mathrm{as}}_n=(2n+1)^2.
\]
これは適合平方根の選択に依存しない。
\end{corollary}
\begin{proof}
$\HH$ の定義により、選べる平方根 $s$ の
全切断環 $B(\mathcal C,\mathcal O,s)$ は二次生成の
正則正規二次商であり、Hilbert 関数は $(n+1)^2$ である。
$s^2\simeq t$ を通じてその偶数 Veronese は
内在的反 Serre 軌道環に同型となる。
直前の補題を適用すれば二次表示となる。
第一・第二次数の次元は $3^2=9$、$5^2=25$ なので、
二次関係の次元は $9^2-25=56$ である。
\end{proof}

\begin{lemma}[体上で二次なら全平坦族でも二次]
\label{lem:relative-quadratic-from-fibers}
$T$ を局所 noetherian スキームとし、
$\mathcal R=\bigoplus_{n\ge0}\mathcal R_n$ を
$\mathcal R_0=\mathcal O_T$、
各 $\mathcal R_n$ が有限局所自由である
次数付き代数とする。
全ての幾何ファイバーが次数一生成かつ二次関係で
表示されると仮定する。
このとき $\mathcal R$ 自身も、任意の底変更と両立する
次数一生成・二次関係の相対表示を持つ。
特に $\operatorname{rank}\mathcal R_1=9$、
$\operatorname{rank}\mathcal R_2=25$ の場合は
\[
 \mathcal R\simeq
 T_{\mathcal O_T}(\mathcal R_1)/(\mathcal J),
 \qquad
 \mathcal J=\ker(\mathcal R_1^{\otimes2}\to\mathcal R_2),
 \quad\operatorname{rank}\mathcal J=56.
 \tag{flat-Q}\label{eq:relative-quadratic-from-fibers}
\]
\end{lemma}
\begin{proof}
各幾何点で積
$\mathcal R_1^{\otimes n}\to\mathcal R_n$ が全射なので、
有限局所自由な標的に対する余核の Nakayama 補題により
すべての $n$ で全射である。
特に二次積は全射であり、$\mathcal J$ は有限局所自由な
局所直和因子で、任意の底変更と整合する。
$\mathcal Q=T(\mathcal R_1)/(\mathcal J)$ から
$\mathcal R$ への次数別全射を取る。
幾何点上では双方の次数 $n$ が
仮定により同型である。局所 noetherian 性により、
核 $\mathcal K_n$ は有限生成であり、
$\mathcal R_n$ が平坦なため
\[
 \mathcal K_n\otimes k(x)
 =\ker(\mathcal Q_n\otimes k(x)\to
                 \mathcal R_n\otimes k(x))=0.
\]
Nakayama 補題から $\mathcal K_n=0$ となる。
この同型と $\mathcal R_n$ の局所自由性により
全次数で任意の底変更と可換する。
\end{proof}

\subsection{cubic 上の相対反 Serre 環と九生成元の表示}
\label{subsec:cubic-relative-antiserre}

有限 foundation の表示スタックは、単に各体上の点だけでなく
相対 smooth proper dg 圏を持つ。
これを利用して、平方根を選ばずに
cubic 開部分の反 Serre 環を族として定義する。
以下の構成は noncubic 表示における
命題~\ref{prop:noncubic-antiserre-descent} と共通の標的を持つ。

\begin{lemma}[相対逆双対化と軌道 Hom]
\label{lem:foundation-relative-antiserre-perfect}
$T$ を任意の基底スキームとし、$U_0$ 上の普遍 foundation 代数を
$T$ に引き戻して $\Lambda_T$ と書く。
$\operatorname{Perf}(\Lambda_T)$ の相対 Serre 関手を $S_T$ とすれば
$\mathsf t_T=S_T^{-1}[2]$ は底変更と両立する自己同値である。
標準の頂点零対象 $\mathcal O_T$ に対して
\[
 \mathscr V_n(T)=
 \mathbf R\operatorname{Hom}_{\Lambda_T}
   \bigl(\mathcal O_T,\mathsf t_T^n\mathcal O_T\bigr)
 \qquad(n\ge0)
 \tag{F-AS}\label{eq:foundation-relative-antiserre}
\]
は底上 perfect な複体であり、任意の底変更と両立する。
それらには合成による斉次積がある。
\end{lemma}
\begin{proof}
補題~\ref{lem:foundation-relative-smooth-proper} によって
$\Lambda_T$ は底上有限局所自由で、対角双加群は
有限局所射影項からなる長さ三以下の相対 bar 解消を持つ。
従って $\operatorname{Perf}(\Lambda_T)$ は相対 smooth proper である。
相対 Serre 双加群は
$D_T\Lambda_T=\mathbf R\mathcal Hom_T(\Lambda_T,\mathcal O_T)$、
その逆は相対逆双対化双加群で表される。
いずれも対角双加群の perfect 解消から形成でき、
双加群として perfect であり、任意の底変更と両立する。
よって $\mathsf t_T$ とその反復は相対 perfect 対象を保つ。
その間の導来 Hom は底上 perfect であり、
相対 smooth proper 性によって任意の底変更と可換する。
積は、次数 $m,n$ の射 $a,b$ に対して
$\mathsf t_T^n(a)\circ b$ を取るもので、結合律は
自己同値の合成の結合律から従う。
\end{proof}

\begin{proposition}[cubic 開部分の反 Serre Grassmann 射]
\label{prop:cubic-relative-antiserre-grassmann}
幾何点が cubic な $\HH$ 対象を表す開部分
$\mathfrak C\subset\mathfrak Q_{\mathrm{fnd}}$ 上で、
(\ref{eq:foundation-relative-antiserre}) は
全次数で次数零に集中するベクトル束 $\mathscr R_n$ を与え、
\[
 \operatorname{rank}\mathscr R_n=(2n+1)^2 .
\]
積は任意の底変更と可換し、$\mathscr R$ は次数一で生成される。
従って階数 $56$ の関係部分束
\[
 \mathscr J=\ker\bigl(
 \mathscr R_1\otimes\mathscr R_1\longrightarrow\mathscr R_2
 \bigr)
\]
が定まり、標準九次元空間 $U$ に対して
\[
 \Psi_{\mathrm{cub}}:
 \mathfrak C\longrightarrow
 \left[\operatorname{Gr}(56,U\otimes U)/\operatorname{PGL}(U)\right]
 \tag{C-Gr}\label{eq:cubic-antiserre-grassmann}
\]
という有限型スタック間の射が得られる。
この射は幾何点で、元の pointed $\HH$ 対象の
内在的反 Serre 環の二次関係空間を与える。
\end{proposition}
\begin{proof}
まず $\mathfrak F_0$ の線型化した表示
$[U_0/\mathcal G]$ で作る。
任意の幾何点で cubic 逆構成から得られる
$(\mathcal C,\mathcal O)\in\HH$ を用いると、
tilting 同値は $\mathcal O_T$ と構造対象を同定し、
Serre 関手を保つ。
$\mathsf t=S^{-1}[2]$ はこの同値の下で $s^2$ に対応する。
既存の切断環と Serre 双対の計算から、
$n\ge0$ について
\[
 \operatorname{Ext}^{i}
  (\mathcal O,\mathsf t^n\mathcal O)=0\quad(i\ne0),
 \qquad
 \dim\operatorname{Hom}
  (\mathcal O,\mathsf t^n\mathcal O)=(2n+1)^2.
\]
具体的には切断環 $B(\mathcal C,\mathcal O,s)$ は
$(1-z^2)/(1-z)^4$ の Hilbert 級数を持ち、
$\operatorname{Ext}^2$ は Serre 双対と負次数切断の消滅で零、
$\operatorname{Ext}^1$ は quadratic AS 被覆からの
次数付き局所双対・飽和 section の消滅によって零である。
この結果は根の選択に依存しない。

補題~\ref{lem:foundation-relative-antiserre-perfect} の
perfect 複体は、全幾何点で cohomological amplitude $[0,0]$ である。
perfect 複体の局所自由分解における単元 minor の消去を行えば、
各 $\mathscr V_n$ は指定階数のベクトル束に等しく、
任意の底変更と両立する。
各幾何点の反 Serre 環は
一次生成の $B(\mathcal C,\mathcal O,s)$ の偶数 Veronese なので、
$\mathscr R_1^{\otimes m}\to\mathscr R_m$ は幾何点で全射となる。
有限局所自由加群への写像の全射性は幾何点で判定できるから
族でも全射であり、特に二次関係は局所直和因子となる。

構成は標識付き dg 圏、相対 Serre 関手、合成だけで決まる。
$K$ は各基底変換において
$\Lambda$ の頂点冪等元に沿う内部自己同型で作用し、
$\mathcal O_T$ 上では頂点零のスカラーを一に規格化できる。
従ってその作用は標識付き Morita 自然同型を通じて
$\mathscr R_n$ 上で自明となる。
また命題~\ref{prop:ruling-involution-mutation} の
相対 mutation は $\mathcal O_T$ を保つ Morita 同値であり、
Serre 関手の自然性により $\mathscr R_n$ と積を保つ。
相対 ruling 交換の coherent な $S_2$ 作用の下で
この同定も合成と整合する。
従って $\mathscr R_n$、$\mathscr J$ とその
射影枠束は $\mathfrak C$ へ降下する。
$\operatorname{PGL}(U)$-同変な Grassmannian の射として
(\ref{eq:cubic-antiserre-grassmann}) を得る。
\end{proof}

\begin{corollary}[cubic chart の全次数二次再構成]
\label{cor:cubic-relative-full-quadratic-presentation}
命題~\ref{prop:cubic-relative-antiserre-grassmann} の
相対反 Serre 環について
\[
 \mathscr R\simeq
 T_{\mathcal O_{\mathfrak C}}(\mathscr R_1)/(\mathscr J),
 \qquad
 \mathscr J=\ker(\mathscr R_1^{\otimes2}\to\mathscr R_2)
 \tag{C-quad}\label{eq:cubic-relative-full-quadratic}
\]
が全次数で成立し、任意の底変更に両立する。
従って $\Psi_{\mathrm{cub}}$ の二次関係は
cubic 族の反 Serre 環の全積を復元する。
\end{corollary}
\begin{proof}
全幾何ファイバーの反 Serre 環は
系~\ref{cor:h2-intrinsic-antiserre-quadratic} により二次である。
$\mathscr R_n$ は命題~\ref{prop:cubic-relative-antiserre-grassmann} により
全次数で有限局所自由なので、
補題~\ref{lem:relative-quadratic-from-fibers} を
cubic の滑らかな有限型表示上に適用できる。
構成は基底変更と ruling 交換に整合するので
$\mathfrak C$ へ降下する。
\end{proof}

\begin{remark}[二つの chart の射と内在的な相対認識]
命題~\ref{prop:cubic-relative-antiserre-grassmann} と
命題~\ref{prop:noncubic-antiserre-grassmann} により、
cubic と noncubic の双方に、同じ九生成元・
階数56の二次関係 Grassmann 商への具体的な射がある。
両者の幾何点像は、内在的な反 Serre 環の関係空間によって
pointed 同値類を区別する。
しかし二つの chart で定義された相対 functor の
重なりと合成を内在的な $\HHstack$ の
有限型スタックと同一視する定理は未証明である。
また foundation の arbitrary point について
(\ref{eq:foundation-relative-antiserre}) が次数零に
集中するとは主張しない。
\end{remark}

\subsection{cubic 開部分の稠密性と noncubic への特殊化}
\label{subsec:cubic-noncubic-components}

以下では $k$ を標数零の代数閉体とする。
$\mathfrak Q=\mathfrak Q_{\mathrm{fnd}}$ と略記し、
$q:\mathfrak F_0\to\mathfrak Q$ を相対 ruling 交換による商とする。
$\mathfrak F_{\mathrm{cub}}\subset\mathfrak F_0$ は
命題~\ref{prop:cubic-foundation-open} の geometric 開部分である。
$|\mathfrak Q|$ の点は任意の代数閉拡大体上の幾何点を同値で割ったものとし、
全ての代数閉拡大体上の $\mathcal H_2$ 対象の像の合併を $Y$ と書く。
これは像の合併としての定義であり、$\mathcal H_2$ 公理一般の任意底変更に関する
安定性を別の入力として仮定していない。
$Y$ には $|\mathfrak Q|$ からの部分空間位相を入れる。
これは内在的な相対 $\mathcal H_2$ スタックの構成を先取りする定義ではない。

\begin{proposition}[実現点集合の既約性]
\label{prop:h2-component-density}
$\mathfrak Q$ は既約である。
$\mathfrak C=q(\mathfrak F_{\mathrm{cub}})$ は非空開部分スタックであり、
全ての幾何点が cubic な $\mathcal H_2$ 対象を表す。
特に
\[
 |\mathfrak C|\subset Y\subset|\mathfrak Q|,
 \qquad
 \overline{|\mathfrak C|}^{\,Y}=Y.
\]
従って $Y$ は既約かつ連結であり、その既約成分と連結成分はともに一つである。
全 noncubic 実現点は、この同じ既約空間における cubic 開部分の閉包に属する。
\end{proposition}

\begin{proof}
$U_0$ は $\mathbb A^{64}$ の非空開部分なので既約である。
商への連続全射\cite[Tags 04XL, 0DR5]{StacksModuli}
\[
 |U_0|\longrightarrow|\mathfrak F_0|
       \longrightarrow|\mathfrak Q|
\]
により、$\mathfrak F_0$ と $\mathfrak Q$ も既約となる。

$\mathfrak F_{\mathrm{cub}}$ は非空である。例えば可換二次曲面の foundation
\[
 (\mathcal O,\mathcal O(1,0),\mathcal O(1,1),\mathcal O(2,1))
\]
はこの部分に入る。実際、$A=H^0(\mathcal O(1,0))$、$B=H^0(\mathcal O(0,1))$ と置けば、
$W_1=A$、$P=B$、$W_2=A\otimes B$、$Q=A$、
$W_3=\operatorname{Sym}^2A\otimes B$ である。
$F$ は因子交換の同型、$G$ は対称積の写像で、
$R=\ker G=\bigwedge^2A\otimes B$ となる。
対応する四重テンソルは、因子 $0,2$ の非退化交代テンソルと
因子 $1,3$ の非退化 pairing の積である。
従って任意の隣接二因子の非零縮約は消えない。

$q$ は有限 \'etale、従って開写像であるから、その像 $\mathfrak C$ は非空開部分となる。
幾何点の cubic 逆構成と、ruling 交換の mutation による記述によって
$|\mathfrak C|\subset Y$ である。
既約空間 $|\mathfrak Q|$ の非空開部分 $|\mathfrak C|$ は既約かつ稠密なので、
\[
 \overline{|\mathfrak C|}^{\,Y}
 =Y\cap\overline{|\mathfrak C|}^{\,|\mathfrak Q|}=Y.
\]
既約部分空間の閉包は既約であり、主張が従う。
ここでは単に $Y$ が既約空間の部分集合であることを使うのではなく、
$Y$ がその非空開部分を含むことを使っている。
\end{proof}

\begin{proposition}[各 noncubic 点への一次元の特殊化]
\label{prop:h2-line-degeneration}
$k$ 上の任意の $\mathcal H_2$ 対象と、その normal foundation の
基底を選んだ表示 $x=(F_0,G_0)\in U_0(k)$ を取る。
$0,1\in T\subset\mathbb A^1_k$ という非空開部分と、
階数二十四の有限自由 $\mathcal O_T$-代数 $\Lambda_T$ を取り、次を満たせる。
\begin{enumerate}
\item $t=0$ の代数は指定された foundation の End 代数 $\Lambda_x$ である。
\item $t\ne0$ の全幾何ファイバーは geometric cubic foundation の代数である。
\item 頂点零を標識とした $\operatorname{Perf}(\Lambda_T)$ は、
      底変更と両立する smooth proper な相対 dg 圏の族である。
      全幾何ファイバーは $\mathcal H_2$ 対象の標識付き導来圏と同値である。
\end{enumerate}
特に $x$ が noncubic なら、これは任意に指定した noncubic 対象を
特殊ファイバーに持ち、他の全ファイバーが cubic となる foundation の族である。
\end{proposition}

\begin{proof}
$U_{\mathrm{cub}}\subset U_0$ を $\mathfrak F_{\mathrm{cub}}$ の逆像とし、
$y=(F_1,G_1)\in U_{\mathrm{cub}}(k)$ を一つ選ぶ。
これは前命題の可換二次曲面の foundation の表示に固定してよく、
$x$ ごとに別の cubic 対象を探索する必要はない。
アフィン空間内の直線
\[
 \ell(t)=(1-t)x+ty,
 \quad
 F_t=(1-t)F_0+tF_1,
 \quad
 G_t=(1-t)G_0+tG_1
\]
を取る。$\ell^{-1}(U_{\mathrm{cub}})$ は $1$ を含む非空開部分なので、
その補集合 $Z$ は有限である。
\[
 T=\mathbb A^1_k\setminus(Z\setminus\{0\})
\]
と置けば、$0,1\in T$ であり、$T\setminus\{0\}$ は $U_{\mathrm{cub}}$ に写る。
$t=0$ の像も $U_0$ に属するので $\ell(T)\subset U_0$ である。
この直線に普遍二写像表示の代数を引き戻して $\Lambda_T$ を得る。
その加群は四つの対角成分、三つの $W_i$、二つの矢空間 $P,Q$ と
$P\otimes Q$ の直和であり、全体で階数二十四の自由加群である。
積は $F_t,G_t$ とその合成からなり、任意の底変更と両立する。
最初の二つの主張と、全幾何ファイバーの $\mathcal H_2$ 実現はこれで従う。

相対 proper 性は $\Lambda_T$ の有限自由性による。
相対 smooth 性には、次の補題を用いる。
頂点零の標準 projective を標識対象とすれば、その自己 Hom は $\mathcal O_T$ であり、
幾何ファイバー上の tilting 同値は指定された $\mathcal O$ をこの対象へ移す。
従ってこれは標識付き相対 dg 圏の族である。
必要なら $\operatorname{Spec}k[t]_{(t)}$ や
$\operatorname{Spec}k[[t]]$ へ底変更して、一般ファイバーが cubic、
特殊ファイバーが指定対象という一径数退化にできる。
\end{proof}

\begin{lemma}[有向有限代数の相対 smooth proper 性]
\label{lem:foundation-relative-smooth-proper}
任意の底スキーム上の二写像表示から得られる普遍 foundation 代数 $\Lambda$ は、
底上で perfect であり、相対対角双加群も perfect である。
この記述は任意の底変更と両立する。
\end{lemma}
\begin{proof}
$E=\bigoplus_{i=0}^3\mathcal O e_i$ とし、$J$ を非対角成分の直和とする。
有向性により $\Lambda=E\oplus J$ は増大された $E$-代数である。
$E$ に相対的な正規化 bar 分解は、次数 $n$ に
\[
 \Lambda\otimes_E J^{\otimes_E n}\otimes_E\Lambda
\]
を持つ。増大 bar 分解の単位挿入による収縮と正規化により、
この複体は対角双加群 $\Lambda$ を解消する。
非零なテンソルは頂点の狭義単調列に対応するため、$n>3$ の項は零である。
各成分 $J_{ij}$ は有限局所自由であり、上の項は
\[
 \Lambda e_i\otimes_{\mathcal O} V\otimes_{\mathcal O}e_j\Lambda
\]
という形の有限直和である。ここで $V$ は有限局所自由である。
従って各項は局所的に有限生成 projective な
$\Lambda\otimes_{\mathcal O}\Lambda^{\mathrm{op}}$-加群であり、
対角は perfect となる。
底上の perfect 性は $\Lambda$ の有限局所自由性による。
全ての項と微分は局所自由な成分と積から作られるため、底変更と両立する。
\end{proof}

\begin{proposition}[cubic から noncubic への反 Serre 環の実際の特殊化]
\label{prop:line-degeneration-antiserre-algebras}
命題~\ref{prop:h2-line-degeneration} の一次元族
$\Lambda_T$ に対して、(\ref{eq:foundation-relative-antiserre}) の複体は
全次数で有限局所自由な次数零の加群 $\mathscr R_{T,n}$ となり、
\[
 \operatorname{rank}\mathscr R_{T,n}=(2n+1)^2
 \qquad(n\ge0).
\]
$\mathscr R_T=\bigoplus_{n\ge0}\mathscr R_{T,n}$ は
次数一生成の、任意底変更と両立する相対次数付き代数であり、
その二次関係は階数 $56$ の部分束をなす。
従って曲線からの射
\[
 T\longrightarrow
 \left[\operatorname{Gr}(56,U\otimes U)/\operatorname{PGL}(U)\right]
 \tag{D-Gr}\label{eq:line-antiserre-grassmann}
\]
が存在し、$t=0$ の値は指定した noncubic 対象の
\emph{内在的反 Serre 環の二次関係}、
$t\ne0$ の値は cubic 対象のそれである。
特に noncubic 点への特殊化は foundation の dg 圏に限らず、
反 Serre の全次数付き環の族としても実現される。
\end{proposition}
\begin{proof}
補題~\ref{lem:foundation-relative-antiserre-perfect} を $T$ に適用する。
命題~\ref{prop:h2-line-degeneration} によって
$T$ の全幾何ファイバーは $\HH$ 対象の標識付き導来圏である。
従って各 $n$ の導来 Hom の幾何ファイバーは、
命題~\ref{prop:cubic-relative-antiserre-grassmann} の証明で
用いた $\HH$ の反 Serre 軌道の Ext 消滅から、
次数零の $(2n+1)^2$ 次元空間に集中する。
perfect 複体の局所分裂により
$\mathscr V_n(T)$ は同階数の有限局所自由加群になり、
任意の底変更と可換する。
各ファイバーの反 Serre 環は一次生成環の偶数 Veronese なので、
$\mathscr R_{T,1}^{\otimes m}\to\mathscr R_{T,m}$ は幾何点で全射、
従って $T$ 上でも全射である。
$\mathscr R_{T,1}$、$\mathscr R_{T,2}$ の階数がそれぞれ九、二十五であるため
二次関係は階数五十六の部分束である。
射影枠束から (\ref{eq:line-antiserre-grassmann}) を得る。
各ファイバーでの内在的反 Serre 環との一致は
Serre 関手の標準 Morita 自然性による。
\end{proof}

\begin{corollary}[特殊化族の全反 Serre 環の二次提示]
\label{cor:line-full-quadratic-degeneration}
命題~\ref{prop:line-degeneration-antiserre-algebras} の
$T$ 上で、自然な射は全次数の同型
\[
 T_{\mathcal O_T}(\mathscr R_{T,1})/
 \ker(\mathscr R_{T,1}^{\otimes2}\to\mathscr R_{T,2})
   \xrightarrow{\ \sim\ }\mathscr R_T
 \tag{D-quad}\label{eq:line-full-quadratic}
\]
を与える。従って cubic から任意の noncubic 点への
反 Serre 環の退化は、二次関係だけの
\emph{完全な}相対提示によって記述される。
\end{corollary}
\begin{proof}
$\mathscr R_{T,n}$ は各次数で有限局所自由で、
全幾何ファイバーは $\HH$ の反 Serre 環である。
系~\ref{cor:h2-intrinsic-antiserre-quadratic} と
補題~\ref{lem:relative-quadratic-from-fibers} の適用により従う。
\end{proof}

\begin{proposition}[反 Serre 関係の構成可能な実現集合と既約な代数的包絡]
\label{prop:h2-constructible-relations-and-irreducible-closure}
$k$ を標数零の代数閉体とし、九次元空間 $U$ を固定する。
\[
 \mathfrak G_9=
 \bigl[\operatorname{Gr}(56,U\otimes U)/\operatorname{PGL}(U)\bigr]
\]
と書く。すべての代数閉拡大体の $\HH$ 対象から得られる
反 Serre 二次関係の幾何点像を
$\mathcal Y_{\mathrm{as}}\subset|\mathfrak G_9|$ とする。
このとき以下が成立する。
\begin{enumerate}
\item $\mathcal Y_{\mathrm{as}}$ は
  $\Psi_{\mathrm{cub}}(\mathfrak C)$ と
  $\Psi_{\mathrm{nc}}(\mathfrak N^{\mathrm{pres}})$
  の幾何点像の和に等しく、有限型の構成可能集合である。
\item その部分空間位相は既約であり、
  \[
   \overline{\Psi_{\mathrm{cub}}(|\mathfrak C|)}
           ^{\,\mathcal Y_{\mathrm{as}}}
     =\mathcal Y_{\mathrm{as}} .
  \]
\item cubic の相対反 Serre 環の一次数束を
  $U_{\mathrm{cub}}$ 上で枠付けた
  $\operatorname{GL}(U)$-torsor から
  $\operatorname{Gr}(56,U\otimes U)$ への射の
  幾何点像の Zariski 閉包に、被約構造を入れたものを
  $Z_2$ と置く。
  $Z_2$ は有限型の既約な
  $\operatorname{PGL}(U)$-不変閉部分スキームであり、
  \[
    \mathfrak Z_2=[Z_2/\operatorname{PGL}(U)]
       \hookrightarrow\mathfrak G_9
       \tag{H2-envelope}\label{eq:h2-irreducible-envelope}
  \]
  は有限型既約 Artin スタックである。
  $\mathcal Y_{\mathrm{as}}\subset|\mathfrak Z_2|$ であり、
  cubic の実現像は $|\mathfrak Z_2|$ において稠密である。
\end{enumerate}
ただし $\mathfrak Z_2$ はあくまで
\emph{関係式の既約な代数的包絡}であり、
余分な幾何点や冪零方向を持つ可能性がある。
これを内在的な $\HHstack$ と同一視したとは主張しない。
\end{proposition}
\begin{proof}
全 $\HH$ 幾何点の分類によれば、対象は
cubic の逆構成か noncubic の二形式表示のいずれかから得られる。
双方の反 Serre 環は
系~\ref{cor:h2-intrinsic-antiserre-quadratic} により二次である。
従って二次関係は対象を完全に決め、
二つの表示射の幾何点像の和が
$\mathcal Y_{\mathrm{as}}$ となる。
$\mathfrak C$ と $\mathfrak N^{\mathrm{pres}}$ は有限型である。
$\mathfrak G_9$ の Grassmannian 平滑表示で
これらの射を引き戻し、有限型スキームの atlas を選ぶと、
その Grassmannian への像は有限型スキームの射の像である。
Chevalley の定理
\cite[Tag 054K]{StacksModuli} によりこれらの像は構成可能であり、
商の位相でも構成可能である。

$U_{\mathrm{cub}}$ は既約な $\mathbb A^{64}$ の非空開であり、
その階数九ベクトル束の枠束も、幾何学的既約な
$\operatorname{GL}_9$ をファイバーとするため既約である。
その Grassmannian への像は既約なので、
被約閉包 $Z_2$ も既約である。
枠の変更による $\operatorname{PGL}(U)$ 作用は
像を保つため、その閉包も不変である。
よって (\ref{eq:h2-irreducible-envelope}) は
有限型既約 Artin スタックを与える。

任意の noncubic 幾何点 $x$ に対し、
命題~\ref{prop:line-degeneration-antiserre-algebras} と
系~\ref{cor:line-full-quadratic-degeneration} により、
$\operatorname{Spec}k[[t]]$ から $\mathfrak G_9$ への射が存在し、
一般点が cubic、特殊点が $x$ の反 Serre 関係である。
階数九ベクトル束をこの DVR 上で自明化すれば、
Grassmannian への射に持ち上がる。
その一般点は $Z_2$ に入るので、閉性から特殊点も
$Z_2$ に入る。従って
$\mathcal Y_{\mathrm{as}}\subset|\mathfrak Z_2|$ である。

cubic の幾何点像は既約であり、
任意の noncubic 幾何点はその閉包に入る。
そのため $\mathcal Y_{\mathrm{as}}$ は
既約集合の閉包として既約である。
$Z_2$ の定義から cubic の像は
$|\mathfrak Z_2|$ に稠密である。
\end{proof}

\begin{remark}[内在的な相対モジュライへの移行条件]
命題~\ref{prop:h2-component-density} は、実現点集合の部分空間位相についての
無条件の結論である。
命題~\ref{prop:h2-line-degeneration} は、全幾何ファイバーが所要の
標識付き導来圏である実際の smooth proper dg 圏の族を与える。
これらに非被約基底上の $\mathcal H_2$ locus の構成は必要ない。
一方、この族に相対 heart、相対 Serre 平方根、AZ 豊富性、
二次 AS 正則被覆が局所的に存在するという追加主張は、ここでは証明していない。

内在的なモジュライスタック $\mathfrak H_2$ に対して、
$|\mathfrak H_2|\to Y$ が同相であると示せれば、同じ成分数が直ちに従う。
また、全ての代数閉拡大体に対して、本節の普遍 cubic 族と各対象への直線族が
$\mathfrak H_2$ の族として認められることを示しても十分である。
実際、普遍 cubic 族の像は既約であり、各直線族の一般点がその像に入るため、
その閉包は全ての対象を含む。従って $\mathfrak H_2$ は既約かつ連結となる。
この移行条件を確認せずに、相対的な根の存在まで要求するモジュライの
成分数を確定したとは主張しない。
\end{remark}

\subsection{反 Serre 環の商と標識に由来する慣性}
\label{subsec:antiserre-stack-candidate}

ここでは $k$ 上の pointed 圏の射を
$(F,u):(\mathcal C,\mathcal O)\to(\mathcal C',\mathcal O')$、
$u:F\mathcal O\xrightarrow\sim\mathcal O'$ とし、
二つの射の間の二射を $u'\eta_{\mathcal O}=u$ を満たす自然同型
$\eta:F\Rightarrow F'$ とする。
以下の自己同型群は、この二groupoidの一射を二射で割った群を意味する。

\begin{lemma}[標識付き圏の二射と次数スカラー]
\label{lem:pointed-degree-scalar}
$(\mathcal C,\mathcal O)\in\HH$ とし、
$R=\bigoplus_{n\ge0}\Hom(\mathcal O,t^n\mathcal O)$ とする。
$\operatorname{Aut}(\operatorname{Id}_{\mathcal C})=k^\times$ であり、
pointed 一射の二自己同型は自明である。
また完全列
\[
 1\longrightarrow\mathbb G_m(k)
 \xrightarrow{\lambda\mapsto\xi_\lambda}
 \operatorname{Aut}_{\mathrm{gr}}(R)
 \longrightarrow
 \operatorname{Aut}_{\mathrm{pt}}(\mathcal C,\mathcal O)
 \longrightarrow1,
 \qquad \xi_\lambda|_{R_n}=\lambda^n\operatorname{id}
\]
がある。従って graded algebra の商をそのまま用いると、
pointed 圏の自己同型群にない次数スカラーが慣性として一つ余る。
一方、標識対象自身のスカラーをさらに慣性へ加える必要はない。
\end{lemma}

\begin{proof}
自然自己同型 $\eta$ の各 $t^{-n}\mathcal O$ 上の成分は、例外性により
スカラー $\lambda_n$ である。
$\Hom(t^{-n}\mathcal O,\mathcal O)=R_n\ne0$ と自然性から
$\lambda_n=\lambda_0$ となる。
反 Serre 軌道の豊富性により全ての対象はこれらの有限直和の商なので、
$\eta$ は全ての対象で同じスカラーである。
pointed 二射は標識上で恒等であるため、このスカラーは一に限られる。
特に $(\operatorname{Id},\lambda\operatorname{id}_{\mathcal O})$ と
$(\operatorname{Id},\operatorname{id}_{\mathcal O})$ は互いに二同型である。

graded 自己同型は AZ 再構成を通じて pointed exact 自己同値を与える。
群準同型の規約として、$\phi$ には右加群の作用を $\phi^{-1}$ で twist する関手を対応させる。
逆に pointed exact 自己同値は Serre 関手と可換するので、
その可換同型を選ぶことで $R$ の graded 自己同型を与える。
可換同型の異なる選択は、前段の自然自己同型の計算により一つのスカラーだけ異なり、
$R_n$ 上ではその $n$ 乗だけ異なる。これで全射性を得る。

核を計算する。graded 自己同型 $\phi$ の誘導する pointed 関手が恒等と二同型なら、
各自由シフト $\pi R(i)$ 上のその二同型は、
基準同型 $b\mapsto\phi(b)$ のスカラー $c_i$ 倍である。
隣接シフト間の Hom に対する自然性は
$c_{i+1}\phi_1=c_i\operatorname{id}_{R_1}$ を与える。
従って $\phi_1$ はスカラーで、$R$ の一次生成性により $\phi=\xi_\lambda$ である。
逆に $\xi_\lambda^{-1}$ による加群 twist は、次数 $n$ 上で $\lambda^n$ を掛ける写像により
恒等関手と自然同型であり、次数零の標識も保つ。
\end{proof}

九次元空間 $U$ を固定し、
\[
 \operatorname{Gr}\bigl(56,U\otimes U\bigr)
\]
上の関係部分束 $J$ から二次環 $R_J=T(U)/(J)$ を作る。
反 Serre 環として実現される部分を記号的に $X_{\mathrm{as}}$ と書く。
ここで $X_{\mathrm{as}}$ はまだ単なる幾何点集合以上の意味には構成していない。
$\GL(U)$ の中心は Grassmannian に自明に作用するため、
実現部分が有限表示の不変部分スキームとして構成できれば
\[
 [X_{\mathrm{as}}/\GL(U)]\longrightarrow
 [X_{\mathrm{as}}/\PGL(U)]
\]
は次数スカラーを除く $\mathbb G_m$-rigidification となる。
補題の幾何点慣性の計算から、pointed 圏のモジュライに対応すべきものは右辺である。
ただしこの式だけでは相対 $\HH$ モジュライの構成を完了しない。

\begin{proposition}[九生成元・二十五次元二次成分の gerbe は中立]
\label{prop:antiserre-gerbe-neutral}
$Y=\operatorname{Gr}(56,U\otimes U)$ とし、
$\mathcal J\subset U\otimes U$ を普遍部分束とする。
$[Y/\GL(U)]$ 上の普遍環の初めの二成分を
\[
 \mathcal R_1=U,\qquad
 \mathcal R_2=(U\otimes U)/\mathcal J
\]
と書く。その階数はそれぞれ九、二十五である。
線束
\[
 \mathcal L=(\det\mathcal R_1)^{-11}
                 \otimes(\det\mathcal R_2)^2
\]
は中心 $\mathbb G_m\subset\GL(U)$ に関して重み一である。
従って $\mathbb G_m$-gerbe
\[
 [Y/\GL_9]\longrightarrow[Y/\PGL_9]
\]
は中立であり、選んだ $\mathcal L$ によって
$[Y/\GL_9]\simeq[Y/\PGL_9]\times B\mathbb G_m$ と同値になる。
$\mathcal R_n\otimes\mathcal L^{-n}$ は $[Y/\PGL_9]$ 上へ降下し、
その全体は通常の準連接次数付き代数をなす。
特にこの商に残る独立の Brauer 障害はない。
同じ結論は任意の $\PGL_9$-不変部分への制限に成り立つ。
\end{proposition}

\begin{proof}
中心の元 $\lambda I$ は $\mathcal R_1$ に $\lambda$、
$\mathcal R_2$ に $\lambda^2$ で作用する。
従って $\det\mathcal R_1$ と $\det\mathcal R_2$ の重みは九と五十であり、
\[
 -11\cdot9+2\cdot50=1
\]
から $\mathcal L$ の重みは一である。
$\mathcal V=\mathcal R_1\otimes\mathcal L^{-1}$ の中心作用は自明なので、
$\GL_9$-同変降下により $[Y/\PGL_9]$ 上の階数九のベクトル束となる。
その射影枠束は普遍 $\PGL_9$-torsor と一致するので、
通常の枠束はこの torsor の $\GL_9$ への持ち上げを与える。
これが gerbe の中立化である。

同じ事実を局所的に見ると、$\PGL_9$-遷移関数の持ち上げ
$g_{ij}$ による三重重なりの障害が $\lambda_{ijk}$ であるとき、
$\mathcal L$ の遷移係数 $\ell_{ij}$ は重み一のため
$\ell_{ij}\ell_{jk}\ell_{ki}=\lambda_{ijk}$ を満たす。
従って $\ell_{ij}^{-1}g_{ij}$ は真の $\GL_9$-cocycle となる。
一般にこのような重み一の線束を持つ $\mathbb G_m$-gerbe は、
線束の枠束を取ることにより底の $B\mathbb G_m$ との積に同値となる。

$\mathcal R_n$ に対する中心の重みは $n$ なので
$\mathcal R_n\otimes\mathcal L^{-n}$ の重みは零である。
積もこの twist と両立する。
具体的には $\mathcal J\otimes\mathcal L^{-2}\subset
\mathcal V\otimes\mathcal V$ が部分束として降下し、
\[
 T(\mathcal V)/(\mathcal J\otimes\mathcal L^{-2})
\]
が求める代数である。
三次以上の成分の平坦性は Grassmannian 全体では主張せず、
別に定める相対 Hilbert 条件の locus で課す。
\end{proof}


\subsection{定義を弱めて得られる二つの厳密な相対モジュライ}
\label{subsec:weak-moduli-stack-choices}

$\mathfrak H_2$ の内在的相対定義と、既に得た
有限型表示を混同しないために、
\emph{実際にモジュライスタックとして証明できる弱い定義}
を明記する。第一の定義は、体上での $\HH$ の
群oidと慣性を保つ代わりに型の間の特殊化を捨てる。
第二の定義は、特殊化を保持するために
許容対象の条件を三次 Hilbert 関数まで弱める。

\begin{proposition}[型を固定した表示族のモジュライ]
\label{prop:split-h2-presentation-moduli}
$\mathfrak C=q(\mathfrak F_{\mathrm{cub}})$ を
命題~\ref{prop:h2-component-density} の cubic 開部分、
$\mathfrak N^{\mathrm{pres}}$ を
命題~\ref{prop:noncubic-presentation-stack} の
noncubic 表示スタックとする。
\[
 \mathfrak H_2^{\mathrm{split,pres}}
 :=\mathfrak C\ \amalg\ \mathfrak N^{\mathrm{pres}}
 \tag{split-H2}\label{eq:split-h2-presentation-stack}
\]
は有限型 Artin スタックであり、
各成分は既約かつ連結で、
その既約成分数と連結成分数はいずれも二である。
任意の代数閉拡大体 $K/k$ について、
その幾何点群oidは元の $\HH(K)$ の
pointed exact 同値類と全ての慣性を再現する。
ここでこのスタックが分類するものは、
\emph{cubic または noncubic のどちらかの表示を
族全体で選び、型を局所定数とする表示族}である。
\end{proposition}
\begin{proof}
$\mathfrak C$ は $\mathfrak Q_{\mathrm{fnd}}$ の
非空開部分スタックであり、滑らか・有限型で既約である。
$\mathfrak N^{\mathrm{pres}}=[P/G]$ は
非空な既約開 $P\subset\mathbb A^6$ と
滑らかな有限型群 $G$ の商なので、
同じく滑らか・有限型・既約である。
非交和は二つの開閉な既約・連結成分を持つ。
幾何点群oidについては、
系~\ref{cor:h2-point-groupoid-foundation} の全充満性、
命題~\ref{prop:cubic-foundation-open} の cubic 逆構成、
命題~\ref{prop:noncubic-presentation-stack} の
同値類と慣性の計算、そして
定理~\cite[\texttt{thm:complete-h2-classification}]{QuadricIntrinsic2026} の全被覆・非交差性を
合わせればよい。

この表示族の定義では、連結な基底上で
cubic と noncubic を別々の開閉な場合に分けている。
従って命題~\ref{prop:h2-line-degeneration} で構成した
cubic 一般点・noncubic 特殊点の族を、
元の自然な意味での相対 $\HH$ 族として許容するかどうかに
関係なく、このスタックには入れられない。
\end{proof}

\begin{proposition}[三次Hilbert値までの相対二次提示のスタック]
\label{prop:truncated-antiserre-moduli-stack}
$U$ を九次元の $k$-線型空間、
$Y=\operatorname{Gr}(56,U\otimes U)$ とする。
$Y$ 上の普遍関係部分束を $\mathcal J$、
$\mathcal E=(U\otimes U)/\mathcal J$ とし、
\[
 \mathcal L_3:
 U\otimes\mathcal J\longrightarrow
 \mathcal E\otimes U
\]
を $(1\otimes r)$ に関係商を作用させる自然な射とする。
この階数が正確に $176$ である
スキーム論的な locally closed rank stratum を
\[
 X_3=\{J\in Y:\operatorname{rank}(\mathcal L_3)=176\}
 \tag{H2-X3}\label{eq:truncated-antiserre-locus}
\]
と置く。
すると
\[
 \boxed{\quad
 \mathfrak H_2^{(3),\mathrm{pres}}
 = [X_3/\operatorname{PGL}(U)]
 \quad}
 \tag{H2-trunc}\label{eq:truncated-antiserre-stack}
\]
は有限型 Artin スタックであり、
次の意味で\emph{三次までの二次環の相対提示}
のモジュライスタックである：
$\operatorname{GL}(U)$ 版は、
階数九の生成ベクトル束 $\mathcal V$ と
階数五十六の局所直和因子
$\mathcal J_T\subset\mathcal V^{\otimes2}$ で
二次環 $R=T(\mathcal V)/(\mathcal J_T)$ の
$\operatorname{rank}R_2=25$、
$R_3$ が階数四十九の有限局所自由加群となる
族を分類する。
$\operatorname{PGL}(U)$ 版では
次数一の共通スカラーを rigidify する。

$\mathfrak H_2^{(3),\mathrm{pres}}$ 上には、
各次数での共通線束 twist を規格化した
\emph{通常の普遍次数付き代数}
$\mathcal R^{\mathrm{norm}}$ が存在し、
その第零・一・二・三次数は
階数 $1,9,25,49$ の局所自由加群となる。
$\mathfrak C$ と $\mathfrak N^{\mathrm{pres}}$
からの既存の反 Serre 関係射は
$[X_3/\operatorname{PGL}(U)]$ を経由し、
命題~\ref{prop:line-degeneration-antiserre-algebras} の
全ての特殊化族もここを経由する。
しかし $n\ge4$ の平坦性、相対 Serre・AZ・AS 被覆、
内在的 $\HH$ の条件は $X_3$ の定義には含めない。
\end{proposition}
\begin{proof}
$\mathcal E$ は階数 $81-56=25$ である。
二次環の三次成分は
\[
 R_3\simeq
 \operatorname{coker}\bigl(
 U\otimes\mathcal J\to\mathcal E\otimes U\bigr)
\]
であるから、$\operatorname{rank}\mathcal L_3=176$
という条件は、三次成分が階数
$25\cdot9-176=49$ の局所自由加群であることと同値である。
行列の階数が所定値である条件は
大きい小行列式の消滅と少なくとも一つの
所定次数小行列式の可逆性で定める
局所閉 stratum であり、その上で像・余核は部分束になる。
この構成は基底変更と両立し、
$\operatorname{PGL}(U)$-不変な有限型スキーム $X_3$ を与える。
$\operatorname{PGL}(U)$ は滑らかな有限型群なので、
商は有限型 Artin スタックとなる
\cite[Tag 0DLS]{StacksModuli}。

$\operatorname{GL}(U)$ 版の群oid解釈は
生成元の枠束と関係部分束の降下そのものである。
中心 $\mathbb G_m$ は $X_3$ に自明に作用するが、
全次数 $R_n$ には重み $n$ で作用する。
命題~\ref{prop:antiserre-gerbe-neutral} の
線束
\[
 \mathcal L=
 (\det\mathcal R_1)^{-11}
       \otimes(\det\mathcal R_2)^2
\]
は中心重み $-11\cdot9+2\cdot(2\cdot25)=1$ を持つ。
従って全次数で
$\mathcal R_n^{\mathrm{norm}}=
 \mathcal R_n\otimes\mathcal L^{-n}$
は中心重み零になり、
環構造ごと $\operatorname{PGL}(U)$ 商へ降下する。
これは任意底変更と整合する
普遍二次環である。
高次成分が局所自由であるとはここでは主張しない。

cubic・noncubic の既存の相対反 Serre 環と
それらを結ぶ特殊化族では、
全次数で階数 $(2n+1)^2$ の局所自由性が
既に証明されている。従って三次階数は四十九であり、
その関係射は $X_3$ を経由する。
\end{proof}

\begin{proposition}[二根対一致点での一次変形の相違]
\label{prop:truncated-versus-foundation-tangent}
$k$ を標数零の代数閉体とし、
$f_1=f_2=uv$ の noncubic 対象に対応する点を $x$ とする。
この点で $\mathfrak H_2^{(3),\mathrm{pres}}$
の接空間の同型類（すなわち接複体の $H^0$）は
五次元である。一方、同じ点を
$\mathfrak Q_{\mathrm{fnd}}$ に実現したときの
接空間の同型類は四次元である。
従って、両 Artin スタックのこの点における
形式近傍は同値でない。
\end{proposition}
\begin{proof}
命題~\ref{prop:noncubic-third-degree-tangent} により
$\mathcal D_{3,\ge49}$ の点 $x$ での
Grassmannian 内の一次接空間次元は $84$ である。
$\mathcal L_3$ の階数がここでちょうど $176$ なので、
$X_3$ はこの閉 determinantal locus の
$\operatorname{rank}\mathcal L_3\ge176$ という開部分であり、
両者の接空間は一致する。
$\operatorname{PGL}_9$ 軌道の接次元は $79$ である。
従って商スタックの接複体の $[-1,0]$ 切断は
\[
 [\mathfrak{pgl}_9\longrightarrow T_xX_3]
\]
であり、その $H^0$ の次元は $84-79=5$ となる。

一方 $\mathfrak Q_{\mathrm{fnd}}$ は
命題~\ref{prop:relative-ruling-involution} による
有限 \'etale 商の構成と
命題~\ref{prop:foundation-stack} から、
滑らかで相対次元三である。
ここで慣性の Lie 次元を幾何点の集合だけから結論してはならない。
foundation を一つ選び、その表示の安定化群を $H$ とする。
$H$ は滑らかな線形群 $\mathcal G/K$ の閉部分群なので有限型である。
固定有限代数の相対逆双対化双加群と構造対象 $e_0\Lambda$ から
反 Serre 環の次数一・二の積を作る操作は有限完全複体の操作であり、
任意の係数環への底変更と整合する。
従ってこの操作は代数群の射
\[
 H\longrightarrow\operatorname{Stab}_{\operatorname{PGL}_9}(J)
\]
を与える。系~\ref{cor:h2-point-groupoid-foundation} と
補題~\ref{lem:pointed-degree-scalar} により、$k$ 点では単射であり、
その像は二つの内在的 ruling のうち選んだ一つを保つ部分群である。
従って像の指数は高々二である。
標数零の有限型群スキームは滑らかであり、代数群の射の像は閉部分群である。
よって核は自明で、像と標的は同じ Lie 次元を持つ。
後者は軌道の接次元の計算から $80-79=1$ である。
さらに $\mathfrak F_0\to\mathfrak Q_{\mathrm{fnd}}$ は有限 \'etale なので、
両者の慣性 Lie 次元は等しい。
滑らかな Artin スタックでは
\[
 \dim H^0(T_x\mathfrak Q_{\mathrm{fnd}})
 =\dim\mathfrak Q_{\mathrm{fnd}}
  +\dim H^{-1}(T_x\mathfrak Q_{\mathrm{fnd}})
 =3+1=4.
\]
したがって二つのスタックの接複体は
$H^0$ の次元が異なり、形式近傍は同値でない。
\end{proof}

\begin{remark}[三つのスタックの役割と内在的 $\HHstack$]
$\mathfrak H_2^{\mathrm{split,pres}}$ は
全幾何点・慣性を正しく含むが
cubic から noncubic への特殊化を許さず、
見かけ上二つの連結成分を持つ。
一方 $\mathfrak H_2^{(3),\mathrm{pres}}$ は
両型を結ぶ既知の相対反 Serre 環の族を含むが、
三次以外の全 Hilbert 条件を課さないため
余分な対象と非被約変形を含む可能性がある。
$\mathfrak Q_{\mathrm{fnd}}$ も
実際の有限代数のモジュライであるが、
全点が内在的 $\HH$ 対象となるわけではない。
これらのいずれも内在的 $\HHstack$ の
有限型 Artin 表示を証明するものではない。
ただし三次 locus $X_3$ は、
未構成の許容関係式 locus $X_{\mathrm{as}}$
を含む明示的有限型の代数的環境を与える。
\end{remark}

\begin{remark}[有限型 locus と相対再構成の不足]
まず、幾何点の実現集合やその被約スキーム像を取るだけでは、
非被約底上の変形を指定したことにならない。
相対 Hilbert 関数・full section algebra 条件、AZ 豊富性、相対 Serre 同定、
平方根の存在と二次 AS 正則被覆の条件を、
有限表示の部分関手として記述する必要がある。
体上で成立する有限回の同値判定も、全底上で一様な次数の
Isom 関手や降下の表現可能性をそのまま与えるものではない。

直前の命題により、$\PGL_9$-torsor の Brauer 障害そのものは解消される。
しかし線束で次数ごとに twist した環の shift は、元の shift に
底からの線束をテンソルしたものに対応する。
従って相対 tails 圏の構成、相対 Serre 関手との同定、
その $\Hom(\mathcal O,t^n\mathcal O)$ の底変更を確認する仕事はなお残る。
gerbe の中立性は、この相対再構成を代用する主張ではない。

従って有限 foundation の明示商スタックと相対 $S_2$ 商、
全 $\HH$ の幾何点群oidの全充満な実現、その Hochschild 計算、
および反 Serre 表示で除くべき慣性の同定が得られる。
これらを全 $\HH$ の有限型 Artin スタックと同一視する相対定理は残る。
\end{remark}

\subsection{選んだ Serre 根の一次変形障害}
\label{subsec:chosen-root-deformation-obstruction}

例外対象の持上げと、自己同値の持上げは異なる問題である。
後者には、自己同値が変形の Hochschild 類を保つという必要条件がある。
以下では $k=\mathbb C$、$R=k[\varepsilon]/(\varepsilon^2)$ とする。
この例は標数零の例であり、Toda の定理の基礎体の仮定に合わせている。

\begin{lemma}[自己同値を持ち上げるための必要条件]
\label{lem:autoequivalence-deformation-obstruction}
$X$ を滑らかな射影多様体、$u\in\operatorname{HH}^2(X)$ を一次変形の類とし、
$\Phi$ を Fourier--Mukai 核 $P$ を持つ自己同値とする。
$P$ が同一の変形 $X_u$ の左右の作用を持つ perfect kernel に持ち上がるなら、
\[
 \Phi_*u=u
\]
が必要である。従って、同じ変形への $R$-線形 dg/Morita 自己同値としての
持上げにもこの等式が必要である。
\end{lemma}

\begin{proof}
核 $P$ の持上げ障害は、左右から作用する変形類の差
\[
 \operatorname{ob}(P;u,u)
 =u\circ P-P\circ u\in\operatorname{Ext}^2(P,P)
\]
である。ここで左・右作用の規約は、以下の $u-\Phi_*u$ との同定に合わせる。
この障害の解釈には Lowen~\cite[Theorems 1.1 and 4.8, Corollary 4.9]{Lowen2008}
の characteristic morphism の定理を用いる。
$P$ は同値の核なので、左右の Hochschild 作用はいずれも同型であり、
その比較が $\Phi_*$ である
\cite[Proposition 5.6 and Lemma 5.8, arXiv v3]{Toda2009}。
従って一方の同型 $\operatorname{Ext}^2(P,P)\simeq\operatorname{HH}^2(X)$ の下で
障害は $u-\Phi_*u$ に対応する。持上げがあれば障害は零なので主張を得る。

Toda~\cite[Theorems 1.1 and 4.7, arXiv v3]{Toda2009} は
$v=\Phi_*u$ のとき $X_u$ と $X_v$ の間へ同値を持ち上げる十分条件を述べる。
本補題では上の障害計算によって必要性を示した。
\end{proof}

\begin{proposition}[可換 quadric 上の根でも持上げ障害が生じる]
\label{prop:chosen-root-does-not-lift}
$Q=\mathbb P^1\times\mathbb P^1$、$L=\mathcal O_Q(1,1)$ とし、
標準アフィン座標を $(z,w)$ とする。対合と自己同値
\[
 f(z,w)=(-z,w),\qquad s=(-\otimes L)\circ f^*
\]
を取る。この $s$ は $\operatorname{coh}Q$ の exact 自己同値であり、
\[
 s^2\simeq-\otimes\mathcal O_Q(2,2)=S^{-1}[2]
\]
を満たす。しかし $u$ を大域 bivector
$\pi=\partial_z\wedge\partial_w$ の Hochschild 類とすると、
類 $u$ の一次変形へ、この選んだ $s$ は $R$-線形 dg/Morita 自己同値として
持ち上がらない。
\end{proposition}

\begin{proof}
$f^*L\simeq L$ と $f^2=\operatorname{id}$ から平方の式を得る。
また $s^{2n}\mathcal O_Q=\mathcal O_Q(2n,2n)$ なので、
閉ファイバーの $(\mathcal O_Q,s)$ は Artin--Zhang ample である。
Segre 埋込みにより $f$ は $\mathbb P^3$ の線形対合へ延びる。
切断環は可換 quadric の切断環の Zhang twist であり、
多項式環の同じ twist を型 $(1,4,6,4,1)$ の AS 正則被覆に取れる。
従ってこれは元の $\mathcal H_2$ 対象の適合根である。
$\partial_z$ は $\mathbb P^1$ 上の大域ベクトル場であり、
無限遠の座標 $\zeta=z^{-1}$ では $-\zeta^2\partial_\zeta$ となる。
従って $\pi$ は $Q$ 全体の非零大域切断である。

$T_Q\simeq\mathcal O_Q(2,0)\oplus\mathcal O_Q(0,2)$ なので
$H^1(Q,T_Q)=H^2(Q,\mathcal O_Q)=0$ である。HKR 同型から
\[
 \operatorname{HH}^2(Q)
 \simeq H^0(Q,\textstyle\bigwedge^2 T_Q)
 \simeq H^0(Q,\mathcal O_Q(2,2))
\]
を得る。線束による tensor は、その線束が自明な各開集合では恒等関手と
同型であるため、局所の bivector 成分には恒等に作用する。
上の HKR 分解には他の成分がなく、局所制限への写像が単射なので、
$-\otimes L$ は $\operatorname{HH}^2(Q)$ 全体にも恒等に作用する。
一方 $f^*\pi=-\pi$ であり、従って
\[
 s_*u=-u\ne u.
\]
補題~\ref{lem:autoequivalence-deformation-obstruction} より持上げは存在しない。

対象となる一次変形は具体的にも構成できる。
$\mathcal A_\pi=\mathcal O_Q\oplus\varepsilon\mathcal O_Q$ に
\[
 a\star b=ab+\frac{\varepsilon}{2}\pi(da,db)
\]
を入れると、$\pi$ の二重導分性と $\varepsilon^2=0$ により積は結合的である。
これは Toda~\cite[\S\S3--4, arXiv v3]{Toda2009} の
$(\alpha,\beta,\gamma)=(0,0,\pi)$ に対する $R$-flat な代数層であり、
同所の Ore 局所化と coherent 加群の構成によって、
その加群圏は $\operatorname{coh}Q$ の類 $u$ の一次変形を与える。
\end{proof}

\begin{remark}[結論の範囲]
障害を持つのは、ここで自己同値として\emph{選んだ}根 $s$ である。
根同型 $s^2\simeq S^{-1}[2]$ を選ぶかどうかに先立ち、自己同値そのものが
持ち上がらない。この例は、変形後に\emph{別の} Serre 根が存在する可能性を
排除しない。また全 $\mathcal H_2$ の相対モジュライの非存在や、
その許容部分の非開性を主張していない。
ここで構成した abelian 一次変形が、相対的な根・豊富性・AS 正則被覆まで備える
$\mathcal H_2$ の族であることも、この命題の結論には含めていない。

従って、幾何学的点で平方根が存在することと例外対象の無障害性だけから、
選んだ根の忘却射が形式的に滑らか、あるいは \emph{\'etale} であるとは結論できない。
局所 Artin 環 $R$ は代数閉剰余体を持つ strictly henselian 環なので、
点を持つ \emph{\'etale} 近傍には $R$ 上の切断がある。
従ってこの障害を単なる \emph{\'etale} 基底変更で消すこともできない。
\end{remark}

% !TeX root = ../nc-quadric-moduli.tex
% 2026-11-21 / v1.11
\section{Artin 基底での foundation と反 Serre 偏極}
\label{sec:artin-foundation}

この節の基底は、剰余体 $k$ を持つ有限次元可換局所 Artin
$k$-代数 $(A,\mathfrak m,k)$ に限る。
非被約な一般有限型基底への代数化は、この節の結論には含めない。
$(\mathcal C,\mathcal O)\in\HH$ とし、
$\mathcal G=\operatorname{Ind}\mathcal C$ と置く。
flat Abel 変形は Lowen--Van den Bergh の意味で取り、
Grothendieck 圏で記すときは $\mathcal G_A$ とする。

\begin{theorem}[Artin 変形への foundation の持上げ]
\label{thm:artin-foundation-lift}
定理~\cite[\texttt{thm:normal-helix}]{QuadricIntrinsic2026} の一方の helix $(E_i)_{i\in\mathbb Z}$ を選び、
$E_0=\mathcal O$、$T=\bigoplus_{i=0}^{3}E_i$、
$\Lambda=\operatorname{End}(T)$ と置く。
任意の flat Grothendieck 変形 $\mathcal G_A$ において、次が成立する。
\begin{enumerate}
\item 各 $E_i$ は flat な noetherian 対象 $E_{i,A}$ に持ち上がり、
  その持上げの同型類は一意である。同型そのものの一意性は主張しない。
\item 持上げた列は geometric helix であり、
  $T_A=\bigoplus_{i=0}^{3}E_{i,A}$ は $D(\mathcal G_A)$ の compact tilting generator である。
\item $\Lambda_A=\operatorname{End}(T_A)$ は階数二十四の自由 $A$-代数で、
  $\Lambda_A\otimes_A k\simeq\Lambda$。
  これは補題~\ref{lem:foundation-artin-deformation} の有限 foundation 表示に入る。
\item $D(\mathcal G_A)\simeq D(\operatorname{Mod}\Lambda_A)$、
  $D(\mathcal G_A)^c\simeq\operatorname{Perf}\Lambda_A$ である。
  compact 部分は $A$ 上 smooth proper であり、相対 Serre 関手を持つ。
\end{enumerate}
さらに、閉ファイバーを同定した Abel 変形と有限 foundation 代数の線形変形は、
De Deken--Lowen の変形比較により対応する。
この対応は適合平方根または AS4 被覆の持上げを主張しない。
\end{theorem}
\begin{proof}
既存の helix が \cite{DeDekenLowen2011} の定理の仮定を満たすことを確認する。
$\mathcal G$ は局所 noetherian であり、$\mathcal C$ の対象は finitely presented である。
\cite[Proposition 6.34]{LowenVanDenBergh2006} と大域次元二により、
各 $E_i$ からの導来 Hom は直和と可換し、有限 cohomological amplitude を持つ。
従って $E_i$ は $D(\mathcal G)$ で compact である。
四連続項が $D^b(\mathcal C)$ を生成し、$\mathcal C$ が $D(\mathcal G)$ を
localizing に生成するため、四連続項は compact generators になる。
helix の mutation 公式は完全例外列と Serre 双対の通常の公式であり、
周期移動 $E_{i+4}=S^{-1}E_i[2]$ と一致する。
従ってこれは \cite[Definition 5.12]{DeDekenLowen2011} の geometric $(4,3)$-helix である。
基礎体上で全対象と全 Hom 空間は平坦である。

また定理~\cite[\texttt{thm:normal-helix}]{QuadricIntrinsic2026} の豊富性と正次数 Ext の最終的消滅は、
同論文 Corollary 3.16 の生成条件および tails-projectivity を与える。
付随する $\mathbb Z$-代数は右 noetherian なので、同論文 \S3.3 の
$L_{\mathrm{tails}}=T_{\mathrm{tails}}$ が成立する。
従って helix は $\mathbb Z$-generating である。
さらに geometric 性は、同論文 Theorem 4.5 の
$\operatorname{Ext}^{1,2}(E_i,X\otimes_k E_j)=0$（$i\le j$、$X$ 有限次元）を与える。

ここで $A\twoheadrightarrow k$ は coherent な可換環の冪零拡大であり、
$k$ は $A$ 上有限表示だから、同論文 \S4.1 の基底条件も満たす。
Theorem 5.13(1)--(4) により、flat 対象としての helix の一意な持上げ、
compact generation、strong 性、前向き Hom の平坦性と特殊ファイバーとの一致を得る。
有限次元の特殊ファイバーを持つ平坦 $A$-加群は有限自由である。
よって $\Lambda_A$ の頂点成分は元と同じ階数を持ち、全階数は二十四となる。

各 $E_{i,A}$ は有限な $\mathfrak m$-進フィルターを持つ。
その次数付き部分は平坦性により
$(\mathfrak m^r/\mathfrak m^{r+1})\otimes_k E_i$ であり、
いずれも noetherian なので $E_{i,A}$ も noetherian である。
全変形圏の局所 noetherian 性は、Theorem 4.5 と Proposition 4.6(5) による
noetherian $\mathbb Z$-代数の変形表示からも従う。
導来 Morita により表示された導来同値を得る。
有限 foundation の相対 bar 解消
（補題~\ref{lem:foundation-relative-smooth-proper}）が相対 smooth proper 性を与える。
最後の変形比較は、いま確認した仮定の下での
\cite[Theorem 5.15]{DeDekenLowen2011} である。
\end{proof}

\begin{lemma}[Artin 基底での exactness の認識]
\label{lem:artin-morita-heart-exact}
定理~\ref{thm:artin-foundation-lift} の表示の下で、$P$ を
$A$-線形の可逆 perfect $\Lambda_A$-双加群とし、
$F_A=-\otimes_{\Lambda_A}^{\mathbf L}P$ を $D(\mathcal G_A)$ の自己同値とする。
その特殊ファイバー $F_0$ が $\mathcal G$ の完全自己同値から誘導されるならば、
$F_A$ は $\mathcal G_A$ の完全自己同値から誘導される。
この自己同値は noetherian 対象および $A$ 上平坦な対象を保つ。
\end{lemma}
\begin{proof}
$i:\mathcal G\to\mathcal G_A$ を $\mathfrak m$ で消される対象の完全な包含とする。
導来 Morita 表示で $D(i)$ は $\Lambda_A\to\Lambda$ に沿う restriction of scalars になる。
実際、flat な $T_A$ に対する導来随伴により
\[
 \mathbf R\operatorname{Hom}_{\mathcal G_A}(T_A,iX)
 \simeq \mathbf R\operatorname{Hom}_{\mathcal G}(T,X)
\]
であり、$\Lambda_A$ の作用はその還元を通る。
非有界の場合も、\cite[\S5.2]{DeDekenLowen2011} の導来随伴、または
compact generator 上の比較を直和を保つ exact functor へ延長することにより成立する。

$P_0=k\otimes_A^{\mathbf L}P$ と置く。
$\Lambda_A$ は $A$ 上平坦で、双加群の左右の $A$-作用は一致するので、
任意の $M\in D(\operatorname{Mod}\Lambda)$ に対し
\[
 (\operatorname{Res}M)\otimes_{\Lambda_A}^{\mathbf L}P
 \simeq
 \operatorname{Res}\bigl(M\otimes_\Lambda^{\mathbf L}P_0\bigr).
\]
これは $\Lambda\otimes_{\Lambda_A}^{\mathbf L}P
\simeq k\otimes_A^{\mathbf L}P$ と結合則から従う。
従って $F_A i\simeq iF_0$ であり、逆関手についても同じである。

任意の $M\in\mathcal G_A$ のフィルター
$M\supset\mathfrak mM\supset\cdots\supset\mathfrak m^NM=0$
の各商は $i\mathcal G$ に属する。
$F_0$ は heart-exact なので各商の $F_A$ による像は heart に入る。
heart は拡大に閉じるから $F_A M$ も heart に入り、逆関手にも同じ議論が適用できる。
従って制限は Abel 圏の完全自己同値である。
Abel 同値は noetherian 性を保ち、$A$-線形完全関手は有限表示 $A$-加群との
tensor と可換するので、基底上の平坦性も保つ。
\end{proof}

\begin{lemma}[Artin 基底での豊富性の持上げ]
\label{lem:artin-ample-lift}
$F_A$ を前補題の完全自己同値とし、$\mathcal O_A$ を $\mathcal O$ の flat な持上げとする。
$(\mathcal O,F_0)$ が AZ ample であり、全ての $N\in\mathcal C$ について
\[
 \operatorname{Ext}^{j}_{\mathcal C}(F_0^{-n}\mathcal O,N)=0
 \quad(j>0,\ n\gg0)
\]
と仮定する。このとき $(\mathcal O_A,F_A)$ は
$\mathcal C_A=(\mathcal G_A)^{\mathrm{noeth}}$ 上 AZ ample である。
\end{lemma}
\begin{proof}
$L_{n,A}=F_A^{-n}\mathcal O_A$ は flat である。
$N\in\mathcal C$ に対する導来随伴は
\[
 \mathbf R\operatorname{Hom}_{\mathcal G_A}(L_{n,A},iN)
 \simeq\mathbf R\operatorname{Hom}_{\mathcal G}(F_0^{-n}\mathcal O,N)
\]
を与える。任意の noetherian $M\in\mathcal C_A$ をその有限な
$\mathfrak m$-進フィルターで分け、長完全列を使うと、
$\operatorname{Ext}^{j}(L_{n,A},M)=0$（$j>0,n\gg0$）を得る。
この消滅が AZ の Hom 全射条件を与える。

生成条件では $\bar M=M/\mathfrak mM$ に対して十分負のシフトからの有限全射を取り、
$\operatorname{Ext}^{1}(L_{n,A},\mathfrak mM)=0$ によって各射を $M$ へ持ち上げる。
持上げた写像の余核 $Q$ は $Q/\mathfrak mQ=0$ なので、
$Q=\mathfrak mQ=\cdots=\mathfrak m^NQ=0$ である。
従って持上げも全射となる。
\end{proof}

\begin{proposition}[全 Abel 変形の canonical 反 Serre 環]
\label{prop:artin-canonical-antiserre}
任意の flat pointed Abel 変形 $(\mathcal C_A,\mathcal O_A)$ について、
relative compact 圏の反 Serre 関手 $t_A=S_A^{-1}[2]$ は
$\mathcal C_A$ の完全自己同値に延長し、$(\mathcal O_A,t_A)$ は AZ ample である。
さらに
\[
 R_A^{\mathrm{can}}
 =\bigoplus_{n\ge0}\operatorname{Hom}(\mathcal O_A,t_A^n\mathcal O_A)
\]
は全次数自由、階数 $(2n+1)^2$ の二次環である。
生成元の枠と次数線束の標準的な扱いによって $[X_3/\operatorname{PGL}_9]$ の
Artin 変形を与え、
\[
 (\mathcal C_A,\mathcal O_A)
 \simeq(\operatorname{qgr}R_A^{\mathrm{can}},\pi R_A^{\mathrm{can}}).
\]
この構成は Artin 基底変更と両立する。
\end{proposition}
\begin{proof}
有限 foundation 表示の相対 Serre 双加群とその逆は perfect で、
有限な相対 bar 解消によって基底変更と可換する。
特殊ファイバーの $t_0=S^{-1}[2]$ は完全自己同値なので、
補題~\ref{lem:artin-morita-heart-exact} が $t_A$ の heart-exactness を与える。
系~\cite[\texttt{cor:pointed-derived-exact}]{QuadricIntrinsic2026} の証明の Ext 消滅と
補題~\ref{lem:artin-ample-lift} により AZ ample となる。

$t_A^n\mathcal O_A$ は $E_{4n}=t_0^n\mathcal O$ の flat な持上げである。
従って持上げの一意性により $E_{4n,A}$ と同型であり、
\cite[Theorem 5.13(4)]{DeDekenLowen2011} は
$\mathbf R\operatorname{Hom}(\mathcal O_A,t_A^n\mathcal O_A)$ が
次数零に集中する階数 $(2n+1)^2$ の自由加群であることを与える。
これらは compact 対象間の Hom だから、有限 foundation の perfect 複体として
計算でき、任意の Artin 基底変更に対し導来基底変更する。
次数零の自由性から通常の Hom の基底変更も従う。

特殊ファイバーは元の二次反 Serre 環である。
全次数平坦性と二次性から、補題~\ref{lem:relative-quadratic-from-fibers} により
$R_A^{\mathrm{can}}$ 自身が二次である。
次数一、二、三の階数は $9,25,49$ なので $X_3$ の条件を満たす。
noetherian 性は補題~\ref{lem:artin-noetherian-lift} による。
一般の noetherian $M,N$ に対する Hom 有限性には、$M$ の平坦性を必要としない。
$N$ を $\mathfrak m$-進フィルターで分け、その商を $iN_0$ と書けば、通常の随伴
\[
 \operatorname{Hom}_{\mathcal G_A}(M,iN_0)
 \simeq\operatorname{Hom}_{\mathcal G}(M/\mathfrak mM,N_0)
\]
の右辺は有限次元である。Hom の左完全性と有限フィルターから
$\operatorname{Hom}(M,N)$ の $A$ 上有限性が従う。
相対 AZ 再構成を適用して pointed 同値を得る。
Artin 局所基底上の線束は自由なので、次数スカラーの rigidification によって
基底の選択に依存しない商スタックの点を得る。
\end{proof}

\begin{proposition}[Artin 関係式族は表示された全軌道環である]
\label{prop:artin-presentation-full-orbit}
命題~\ref{prop:artin-relations-flat-abelian} の $R_A$ について、
$\mathcal O_A=\pi R_A$、$F_A=(1)$ と置くと、
$(\mathcal O_A,F_A)$ は AZ ample であり、自然な次数付き射
\[
 R_A\longrightarrow
 \bigoplus_{n\ge0}\operatorname{Hom}(\mathcal O_A,F_A^n\mathcal O_A)
\]
は全次数で同型である。
\end{proposition}
\begin{proof}
$F_A$ は次数付き加群圏とその商の完全自己同値であり、
特殊ファイバーは $t_0$ である。
補題~\ref{lem:artin-ample-lift} により AZ ample となる。
$\mathcal O_A$ は flat な $E_0$ の持上げなので compact であり、
$F_A$ は直和を保つ導来自己同値を与えるから、全ての $F_A^n\mathcal O_A$ も compact である。
従って両者の導来 Hom は finite foundation の表示で perfect $A$-複体となる。
その特殊ファイバーは $\operatorname{Hom}(\mathcal O,t_0^n\mathcal O)[0]$ なので、
導来 Nakayama により自身も階数 $(2n+1)^2$ の自由加群の次数零への集中となる。
上の自然射の次数 $n$ 成分は同階数の自由加群間の射であり、
特殊ファイバーで元の反 Serre 全軌道環の同型になる。
従って Nakayama により同型である。
\end{proof}

\begin{corollary}[形式的な偏極忘却の分裂]
\label{cor:artin-polarization-retraction}
$x\in[X_3/\operatorname{PGL}_9](k)$ を $\HH$ の反 Serre 環の点とする。
Artin 基底上で、関係式から pointed Abel 圏へ移る変形写像
\[
 \mathrm{Def}_{[X_3/\operatorname{PGL}_9],x}
 \longrightarrow\mathrm{Def}_{\mathrm{ab},(\mathcal C,\mathcal O)}
\]
は、canonical 反 Serre 環を取る右逆を持つ。
従って一次変形空間上の写像は分裂全射である。さらに標準的な分裂完全列
\[
 0\longrightarrow\operatorname{HH}^{1}(\Lambda)
 \longrightarrow T_x[X_3/\operatorname{PGL}_9]
 \longrightarrow\operatorname{HH}^{2}(\Lambda)
 \longrightarrow0
 \tag{Artin-pol}\label{eq:artin-polarization-tangent-split}
\]
を得る。ここで中央項は接複体の $H^0$、すなわち閉ファイバーの同定を保つ
一次変形の同型類のベクトル空間である。
二根対一致点ではその核は一次元で、
命題~\ref{prop:antiserre-polarization-tangent} の Zhang 方向が生成する。
\end{corollary}
\begin{proof}
左から右の写像は命題~\ref{prop:artin-relations-flat-abelian} により定義される。
命題~\ref{prop:artin-canonical-antiserre} が右逆を与え、合成が恒等になることは
AZ 再構成による。
ここで変形は閉ファイバーの同定を持つものとし、一次変形空間はその同型類を取る。
構造対象の持上げの同型類は例外性から一意なので、pointed 条件による
追加の一次変形方向はない。
De Deken--Lowen の変形比較と既存の有限 foundation 表示によって、
右辺の一次変形空間は $\operatorname{HH}^{2}(\Lambda)$ である。

核を計算するため $A=k[\varepsilon]/(\varepsilon^2)$ とし、
Abel 変形の自明化を選ぶ。
命題~\ref{prop:artin-presentation-full-orbit} により、関係式はこの一定圏上の
$t_0$ の完全自己同値としての持上げ $F_A$ の全軌道環である。
canonical な持上げ $t_A$ と比較すると、
$U_A=F_A t_A^{-1}$ は特殊ファイバーで恒等となる完全自己同値である。
$U_A\mathcal O_A$ は例外対象 $\mathcal O$ の flat な持上げなので
$\mathcal O_A$ と同型であり、pointed 同値として扱える。
逆にそのような $U_A$ から $F_A=U_A t_A$ の全軌道環を作れる。
これは前命題と同じ compact Hom・豊富性・二次性の議論で $X_3$ の点を与える。

この $U_A$ の Morita 双加群 $P$ を考える。
右加群として $P\otimes_A^{\mathbf L}k\simeq\Lambda$ なので、
underlying perfect $A$-複体は次数零の自由加群に集中する。
特殊ファイバーの単位元を持ち上げて作る右加群射
$\Lambda_A\to P$ は同階数自由 $A$-加群間の射で、還元が同型なので同型である。
従って左作用は $\alpha=1+\varepsilon\delta$ という代数自己同型で与えられ、
積との整合は $\delta\in\operatorname{Der}_k(\Lambda)$ と同値である。
右加群同型を $1+\varepsilon a$ だけ変えると $\delta$ は内部導分だけ変わる。
よって同値類は $\operatorname{HH}^{1}(\Lambda)$ となる。
頂点代数 $k^4$ が分離であるため、内部導分を調整して
$\delta(e_i)=0$ とでき、構造対象 $e_0\Lambda_A$ の標識を保てる。
また $Z(\Lambda)=k$ から $Z(\Lambda_A)=A$ なので、標識の scalar の選択は
自然自己同型で除かれ、新たな一次方向を加えない。

Abel 変形の自明化の変更は、特殊ファイバーで恒等となる自己同値による
$F_A$ の共役である。
canonical $t_A$ は相対 Serre 双対の一意性によって全ての $A$-線形自己同値と可換する。
従ってこれは $U_A$ の共役であり、一次の near-identity 群の共役は自明なので、
$\operatorname{HH}^{1}(\Lambda)$ の元を変えない。
逆に全軌道環の同型は AZ 再構成によって、この共役に対応する。
次数比較同型の scalar は $\operatorname{PGL}_9$ において既に除いている。
これで核の同定と完全性が従う。

二根対一致点では命題~\ref{prop:truncated-versus-foundation-tangent} の計算より、
中央項は五次元、右辺は四次元であり、核は一次元である。
既存の Zhang 方向は次数付き環として非自明で、pointed Abel 圏として一定だから、
この核の非零元を与える。
\end{proof}

\begin{corollary}[Abel 変形の無障害性と三基準点の接次元]
\label{cor:artin-abelian-unobstructed}
$\HH$ の任意の対象の pointed Abel 変形は、Artin 基底の小拡大に対して無障害である。
また
\[
 \dim T_x[X_3/\operatorname{PGL}_9]
 =3+2\dim\operatorname{HH}^{1}(\Lambda).
\]
既存の noncubic 三基準点では、Abel 変形の接次元は $3,3,4$、
偏極を保持した関係式変形の接次元は $3,3,5$ となる。
\end{corollary}
\begin{proof}
De Deken--Lowen の比較と命題~\ref{prop:foundation-hochschild} の
$\operatorname{HH}^{3}(\Lambda)=0$ により、小拡大に対する有限代数の結合律障害は消える。
例外対象 $\mathcal O$ の持上げも無障害かつ同型類一意である。
同命題の $\dim\operatorname{HH}^{2}(\Lambda)-\dim\operatorname{HH}^{1}(\Lambda)=3$ と
(\ref{eq:artin-polarization-tangent-split}) が次元公式を与える。
三点では、既存の pointed 自己同値群と有限 foundation の慣性の比較から、
$\dim\operatorname{HH}^{1}(\Lambda)=0,0,1$ である。
\end{proof}

\begin{corollary}[$\operatorname{HH}^{1}=0$ の点での全 Artin 偏極認識]
\label{cor:artin-rigid-polarization}
$\operatorname{HH}^{1}(\Lambda)=0$ と仮定する。
$x$ を閉ファイバーとする任意の Artin $X_3$-変形について、
表示された shift $F_A=(1)$ は、その pointed Abel 圏の canonical な反 Serre 関手
$t_A=S_A^{-1}[2]$ と自然同型である。
従って偏極忘却は、Artin 変形の同型類の関手上で canonical 反 Serre 環の構成と互いに逆になる。
特に noncubic の二つの根対が異なる全ての点で、この結論が成立する。
根を一つ共有する場合も含む。
\end{corollary}
\begin{proof}
前系と同様に $U_A=F_A t_A^{-1}$ を考える。
その Morita 双加群は右 $\Lambda_A$-加群として $\Lambda_A$ と同型であり、
左作用は $\mathfrak m$ を法として恒等な代数自己同型 $\alpha$ で与えられる。
この議論は任意の冪零次数について成立する。

$\alpha\equiv1\pmod{\mathfrak m^r}$ と仮定すると、先頭項
\[
 \alpha-1:\Lambda\longrightarrow
 (\mathfrak m^r/\mathfrak m^{r+1})\otimes_k\Lambda
\]
は $k$-導分である。
$\operatorname{HH}^{1}(\Lambda)=0$ により各係数の導分は内部導分なので、
ある $x_r\in\mathfrak m^r\Lambda_A$ に対して
$\operatorname{Ad}(1+x_r)^{-1}\alpha\equiv1\pmod{\mathfrak m^{r+1}}$
とできる。有限回繰り返せば $\alpha$ は内部自己同型となる。
従って $U_A$ は恒等 functor と自然同型である。
標識上の scalar の差は $Z(\Lambda_A)=A$ によって調整できる。
ここで中心の等式も、$Z(\Lambda)=k$ を使って各 $\mathfrak m$-次数を
順に scalar として取り除くことで得られる。
よって $F_A\simeq t_A$ となる。

noncubic で二つの根対が異なれば、
命題~\ref{prop:noncubic-presentation-stack} の pointed 慣性群は有限である。
標数零ではその Lie 代数は零であり、有限 foundation の慣性との比較により
$\operatorname{HH}^{1}(\Lambda)=0$ となる。
これは一般配置だけでなく、一根共有および有限な追加対称性を持つ配置にも当てはまる。
\end{proof}

\begin{corollary}[非一致根対での関係式空間の滑らかさ]
\label{cor:artin-noncubic-relations-smooth}
$x\in[X_3/\operatorname{PGL}_9](k)$ が、二つの根対の異なる noncubic 対象の
反 Serre 環を表すとする。根を一つ共有する場合も含める。
このとき $[X_3/\operatorname{PGL}_9]$ は $x$ で滑らか、局所次元は三である。
生成元の枠を選んだ対応点 $J\in X_3(k)$ では、$X_3$ は滑らかで局所次元は八十三である。
\end{corollary}
\begin{proof}
この場合の pointed 慣性群は有限なので、
$\operatorname{HH}^{1}(\Lambda)=0$ である。
系~\ref{cor:artin-rigid-polarization} と
系~\ref{cor:artin-abelian-unobstructed} により、
任意の Artin 関係式変形は canonical に偏極され、その underlying Abel 変形は
任意の小拡大に沿って持ち上がる。

枠を付けた $J$ の任意の変形 $J_{A_0}$ と小拡大 $A\twoheadrightarrow A_0$ を取る。
その Abel 変形を $A$ へ持ち上げ、canonical 反 Serre 環を作る。
$A_0$ への制限は元の環と次数スカラーを除いて同型である。
次数一の自由加群の基底を、その同型に合わせて $A$ へ持ち上げれば、
$J_{A_0}$ 自身を制限に持つ $X_3(A)$ の点を得る。
基底の持上げに障害がないことは $\operatorname{GL}_9$ の滑らかさ、
または可逆行列を局所 Artin 環の小拡大へ持ち上げることから従う。
従って $X_3$ は $J$ で形式的に滑らかであり、有限型なので滑らかである。

系~\ref{cor:artin-abelian-unobstructed} により商スタックの $H^0$ 接空間は三次元である。
慣性群が有限で標数零であることから軌道接空間は八十次元であり、
$\dim T_JX_3=3+80=83$ を得る。
滑らかな $\operatorname{PGL}_9$ による商の局所次元は $83-80=3$ である。
\end{proof}

この滑らかさは $X_3$ とその商についての結論である。
平方根と AS4 被覆の相対存在条件を課す $X_{\mathrm{as}}$ の開性や滑らかさを
結論してはいない。

\begin{remark}
この分裂は Artin 変形の比較である。
canonical な反 Serre 偏極の存在は、適合平方根と AS4 被覆の存在とは異なる。
従って全 Abel 変形が $\mathfrak H_2$ の変形であることや、
$X_{\mathrm{as}}$ の代数性・開性を結論していない。
また、四つの foundation 対象の像だけを比べても偏極は認識できない。
それらは例外対象として一意に持ち上がるため、Zhang 方向でも各像は同型になる。
関手を比較するには、有限 foundation 代数の射への作用との整合も必要である。
\end{remark}

\subsection{Artin 上の平方根・被覆の相対族への昇格}
\label{subsec:artin-relative-promotion}

Artin 基底で対象を作っただけでは、定義~\ref{def:relative-h2} が要求する
全有限型基底変更上の圏の系を確認したことにはならない。
次の命題では、既に平方根と AS 被覆が与えられた場合について、その比較を閉じる。
任意の Abel 圏に対する非平坦基底変更の導来降下定理は用いない。

\begin{proposition}[Artin 上の適合表示から相対 $\mathfrak H_2$ へ]
\label{prop:artin-witness-relative-promotion}
$(A,\mathfrak m,k)$ を有限次元可換局所 Artin $k$-代数とする。
$(\mathcal C_A,\mathcal O_A)$ を $(\mathcal C,\mathcal O)\in\HH$ の
flat pointed Abel 変形とし、次のデータが与えられたと仮定する。
\begin{enumerate}
\item 適合する体上の平方根 $s_0$ の $A$-線形 Morita 持上げ $s_A$。
  有限 foundation の Morita 表示において
  $s_A^2\simeq S_A^{-1}[2]$ という双加群の同型を持つ。
\item 全軌道環の同型
  \[
   B_A:=\bigoplus_{n\ge0}
     \operatorname{Hom}(\mathcal O_A,s_A^n\mathcal O_A)
   \simeq\widetilde B_A/(g_A).
  \]
  ここで $\widetilde B_A$ は各次数有限自由な相対 quadratic AS4 環であり、
  $g_A$ は二次の左右正則正規元、商も全次数 $A$-平坦である。
\end{enumerate}
このとき $s_A$ は heart 上完全で $(\mathcal O_A,s_A)$ は AZ ample である。
さらに、有限型 $A$-代数 $A'$ に対して
\[
 \mathcal C_{A'}=
 \bigl((\operatorname{Ind}\mathcal C_A)_{A'}\bigr)^{\mathrm{noeth}}
 \simeq \operatorname{qgr}(B_A\otimes_A A')
\]
と置き、$\mathcal O_{A'}=A'\otimes_A\mathcal O_A$ とすれば、
これらは定義~\ref{def:relative-h2} の相対族を与える。
従って $(\mathcal C_A,\mathcal O_A)$ は $\mathfrak H_2(A)$ の対象となる。
平方根と被覆はこの相対族の存在を証明するための補助データであり、標識に加えない。
\end{proposition}
\begin{proof}
補題~\ref{lem:artin-morita-heart-exact} により $s_A$ は heart-exact である。
$s_0^2=t_0$ と反 Serre の Ext 消滅により、偶数・奇数の部分列を分ければ
$\operatorname{Ext}^{>0}(s_0^{-n}\mathcal O,N)=0$（$n\gg0$）となる。
補題~\ref{lem:artin-ample-lift} が AZ ample 性を与え、
AZ 再構成により $\mathcal C_A\simeq\operatorname{qgr}B_A$ としてよい。

\emph{Abel 圏の基底変更。}
$A'$ は有限型なので、$B_{A'}=B_A\otimes_A A'$ は
$B_A$ に有限個の中心変数を加えた環の商である。
$B_A$ は Artin 持上げにより左右 noetherian だから、$B_{A'}$ も左右 noetherian である。
同じことは被覆 $\widetilde B_A$ にも成り立つ。
$B_{A'}$ の各次数は $A'$ 上有限自由である。
従って $\operatorname{QGr}B_{A'}$ は局所 noetherian で、
$\pi B_{A'}(i)$ を平坦な生成系として持ち、flat Abel 圏になる。

$\operatorname{Gr}B_A$ における $A'$-module object は
$\operatorname{Gr}B_{A'}$ の対象と同じである。
torsion 局所化とその完全充満な右随伴に作用準同型を移すことにより、
\[
 (\operatorname{QGr}B_A)_{A'}
 \simeq \operatorname{QGr}B_{A'}
\]
を得る。noetherian 部分へ制限して所要の基底変更になる。
この構成は tensor の結合則により合成と整合する。

\emph{有限 foundation と canonical compact 圏。}
$I=\mathfrak mA'$、$K=A'/I$ と置く。$I$ は冪零であり、$K$ は有限型 $k$-代数である。
閉じた圏の族は定数族 $\operatorname{QGr}(B_0\otimes_k K)$ である。
命題~\ref{prop:relative-h2-field-recovery} の定数族の構成により、
これは $K$ 上の geometric helix と compact tilting foundation $T_K$ を持つ。
その各前向き Hom は $K$ 上有限自由である。
\cite[Theorem 5.13]{DeDekenLowen2011} を coherent な環の冪零拡大
$A'\twoheadrightarrow K$ に適用できる。

Artin 基底上で構成した $T_A$ を通常の module-object tensor で移した
$T_{A'}=A'\otimes_A T_A$ は $A'$ 上平坦である。
各 summand は $\pi B_A(i)$ の有限直和による有限表示を持つので、
その基底変更は noetherian である。
その $K$ への還元は $T_K$ であり、De Deken--Lowen の持上げの一意性により、
$T_{A'}$ は同定理の compact tilting foundation と一致する。

その End 加群は $A'$ 上平坦で、$K$ への還元は
$\Lambda\otimes_k K$ である。
冪零イデアルを法として有限自由となる平坦加群は有限自由である。
実際、基底を持ち上げた有限自由加群からの射は冪零 Nakayama により全射で、
平坦性によりその核の還元も零なので核自身が零となる。
従って $\operatorname{End}(T_{A'})$ は階数二十四の自由 $A'$-加群である。
自然な代数準同型
\[
 \Lambda_A\otimes_A A'\longrightarrow\operatorname{End}(T_{A'})
 \tag{Artin-BC}\label{eq:artin-promotion-end-basechange}
\]
は $I$ を法として同型であり、両辺が同階数自由だから同型である。
高次自己 Ext は持上げた foundation の strong 性により零である。
よって canonical な導来圏について
\[
 D(\operatorname{QGr}B_{A'})^c
 \simeq\operatorname{Perf}(\Lambda_A\otimes_A A')
\]
を得る。有限な相対 bar 解消から smooth proper 性と相対 Serre 双対を得る。
この表示と (\ref{eq:artin-promotion-end-basechange}) は自然なので、
さらに任意の有限型基底変更を取っても Perf の基底変更と一致する。
これが (RP-BC) を与える。

\emph{平方根の双加群と任意基底変更。}
$s_A$ は Ind 圏に colimit を保つ完全自己同値として延長し、
$A'$-作用にも適用することにより $s_{A'}$ を与える。
その Morita 双加群は $P_A=\mathbf R\operatorname{Hom}(T_A,s_AT_A)$ である。
$s_AT_A$ は compact なので、$\operatorname{add}(T_A)$ の有界複体で代表できる。
この複体の各項と唯一の cohomology $s_AT_A$ は $A$ 上平坦である。
従って任意の $A'$ への通常の tensor は、その導来基底変更を計算し、
唯一の cohomology $A'\otimes_A s_AT_A$ を持つ。
実際、複体の最上次数から下に向かって cycle と boundary の平坦性を示し、
次数零では $0\to B^0\to Z^0\to s_AT_A\to0$ を用いれば、
全ての該当完全列が任意の tensor で保たれる。
特定の strict augmentation の選択を必要としない。
基底変更後の複体は $\operatorname{add}(T_{A'})$ に属する。
End の基底変更と strong 性を使って Hom を計算すると、自然な双加群の射
\[
 P_A\otimes_A^{\mathbf L}A'
 \longrightarrow
 \mathbf R\operatorname{Hom}(T_{A'},s_{A'}T_{A'})
\]
は quasi-isomorphism である。
左作用との整合は、$T_A$ の各自己準同型と $s_A$ によるその像に対する自然性による。
相対 Serre 双加群も有限 foundation の双対と bar 解消から基底変更する。
従って $s_A^2\simeq S_A^{-1}[2]$ は
$s_{A'}^2\simeq S_{A'}^{-1}[2]$ を与える。

\emph{豊富性、軌道 Hom と例外性。}
定数族 $(\mathcal C_K,\mathcal O_K,s_K)$ では、
全ての noetherian $N$ に対して
$\operatorname{Ext}^{>0}(s_K^{-n}\mathcal O_K,N)=0$（$n\gg0$）となる。
これは体上の Ext 消滅と、定数族における有限段の被覆を使う
命題~\ref{prop:relative-h2-field-recovery} の証明による。
$I$-進有限フィルターを使い、補題~\ref{lem:artin-ample-lift} と同じ議論を
$A'\twoheadrightarrow K$ に適用すれば $(\mathcal O_{A'},s_{A'})$ は AZ ample である。
ここでは剰余環が体である必要はなく、剰余圏の Ext 消滅と生成条件を使っている。

構造対象は foundation の一つなので、その相対例外性は
(\ref{eq:artin-promotion-end-basechange}) から従う。
$s_A^n\mathcal O_A$ は、体上の normal helix の $E_{2n}$ の flat lift である。
定理~\ref{thm:artin-foundation-lift} の geometric Hom の平坦性により、
$\mathbf R\operatorname{Hom}(\mathcal O_A,s_A^n\mathcal O_A)$ は
次数零の自由加群 $B_{A,n}$ に集中する。
同じ有界 $\operatorname{add}(T_A)$ 複体による議論を
$s_A^n\mathcal O_A$ に適用すると、全 $n\ge0$ で
\[
 \operatorname{Hom}(\mathcal O_{A'},s_{A'}^n\mathcal O_{A'})
 \simeq B_{A,n}\otimes_A A'
\]
を得る。各項は自由であり、積と任意の有限型基底変更に整合する。
被覆の各次数と切断の商は平坦なので、正則正規切断の完全列は基底変更で保たれる。
左右の相対 Koszul 解消も各内部次数で分裂完全だから、任意基底変更で保たれる。
従って $\widetilde B_A\otimes_A A'$ と $g_A\otimes1$ が (RQ4) の被覆を与える。
全幾何ファイバーは元の体上の対象の体拡大なので (RQ2) も満たす。

\emph{有限性と fppf descent。}
一般の noetherian $M,N$ について、$N$ を $I$-進フィルターで分けると、
その商への Hom は通常の随伴によって
$\operatorname{Hom}_{\mathcal C_K}(M/IM,N_0)$ となる。
定数族の Hom 有限性から $\operatorname{Hom}_{\mathcal C_{A'}}(M,N)$ は有限 $A'$-加群である。
flat な基底変更については、noetherian 対象が finitely presented であることと
Lazard の定理を使う既存の module-object の議論により Hom が基底変更する。

最後に、faithfully flat な有限表示の基底変更に沿う tensor は完全で保守的である。
保守性は有限提示生成対象からの Hom と flat tensor の交換によって検出できる。
従って通常の module-object の随伴と Beck の降下により、対象と射は fppf 降下する。
完全性も忠実平坦に検出される。
noetherian 性の降下は部分対象の昇鎖を完全かつ保守的な tensor 後に調べればよい。
上昇方向では、有限型中心拡大上の有限生成加群表示が noetherian 性を保つ。
これにより affine 上の系は fppf 局所に貼り合い、
定義~\ref{def:relative-abelian-system} の系を与える。
全条件を確認したので定義~\ref{def:relative-h2} に属する。
\end{proof}

\begin{remark}
本命題は平方根・AS 被覆の\emph{存在}を新たに作っていない。
Artin 基底上でそれらを構成したとき、任意の有限型基底変更を含む相対族の条件を
追加の未証明仮定として残さずに確認するための命題である。
一般の $X_3$-変形に対して AS 被覆が存在するという主張は含まない。
\end{remark}

% !TeX root = ../nc-quadric-moduli.tex
% 2026-11-21 / v1.34
\section{Serre 平方根の Artin 持上げ}
\label{sec:artin-serre-roots}
\subsection{Artin 基底での Serre 根の平方関係}

以下で $k$ は標数零の体、$(A,\mathfrak m,k)$ は可換局所 Artin
$k$-代数とする。$\Lambda_A$ は有限自由 $A$-代数で、相対的に smooth、
$\Lambda=\Lambda_A\otimes_Ak$ の中心は $k$ と仮定する。
$\operatorname{DPic}_A(\Lambda_A)$ は可逆で、$\Lambda_A^e$ 上 perfect な双加群の同型類の群であり、
積は導来 tensor とする。相対 Serre 関手の逆に $[2]$ を合成した元を
$t_A$ と書く。これは任意の $A$-線形 Morita 自己同値と可換である。
ここでは平方根との同型そのものを標識にしない。

\begin{lemma}[恒等に還元される Morita 自己同値]
\label{lem:artin-infinitesimal-dpic}
還元準同型の核
\[
 K_A=\ker\bigl(\operatorname{DPic}_A(\Lambda_A)
           \longrightarrow\operatorname{DPic}_k(\Lambda)\bigr)
\]
は、次の nilpotent Lie 代数の Baker--Campbell--Hausdorff 群である：
\[
 \frac{\{\delta\in\operatorname{Der}_A(\Lambda_A):
                    \delta(\Lambda_A)\subset\mathfrak m\Lambda_A\}}
      {\operatorname{ad}(\mathfrak m\Lambda_A)}.
\]
特に $K_A$ の任意の元は一意な平方根を持つ。
\end{lemma}
\begin{proof}
まず $P\in K_A$ を右 $\Lambda_A$-加群として考える。
$\Lambda_A$ は有限次元 $k$-代数なので、$P$ は有限生成 projective
右加群の有界 minimal 複体で表せる。
$\mathfrak m\Lambda_A$ は Jacobson radical に含まれ、
還元後もこの複体は minimal である。
還元は次数零の $\Lambda$ と同型なので、minimal 複体の一意性と
nilpotent Nakayama により、元の複体も次数零の projective 加群となる。
projective 加群の nilpotent ideal に沿う持上げの一意性から、これは
右 $\Lambda_A$-加群 $\Lambda_A$ と同型である。
従って双加群 $P$ は環の自己同型 $\alpha$ による twist である。
還元 $P_k\simeq\Lambda$ は $\alpha_k$ が inner であることを意味するので、
単元を持ち上げて共役することで $\alpha\equiv1\pmod{\mathfrak m}$ とできる。

このような自己同型全体を $\operatorname{Aut}_1(\Lambda_A)$ と書く。
二つの同型が同じ双加群を与える差は inner である。
その共役単元 $u$ の還元は $Z(\Lambda)^\times=k^\times$ に属するから、
スカラーで割れば $u\in1+\mathfrak m\Lambda_A$ とできる。従って
\[
 K_A=\operatorname{Aut}_1(\Lambda_A)
           /\operatorname{Ad}(1+\mathfrak m\Lambda_A).
\]

近恒等自己同型について $\alpha-1$ は nilpotent な $A$-線形作用素である。
有限和で定義される $\alpha^x=(1+(\alpha-1))^x$ は、$x$ の多項式として
任意の非負整数 $x$ で積を保つ。標数零なので多項式恒等式により
任意の $x\in k$ で積を保つ。その $x=0$ での微分
$\log\alpha$ は $A$-導分となる。
逆に、像が $\mathfrak m\Lambda_A$ に含まれる導分の有限指数和は
自己同型である。これらは互いに逆で、積は BCH 積に対応する。
また
\[
 \log(\operatorname{Ad}(u))=\operatorname{ad}(\log u),
 \qquad u\in1+\mathfrak m\Lambda_A.
\]
inner 導分は Lie ideal なので、上の商も nilpotent Lie 代数の
BCH 群である。$\exp(v)^2=\exp(2v)$ と $2\in k^\times$ より最後の主張を得る。
\end{proof}

\begin{proposition}[平方関係の補正]
\label{prop:artin-serre-square-correction}
$p_0\in\operatorname{DPic}_k(\Lambda)$ が $p_0^2=t_k$ を満たすとする。
$p_0$ の Morita 自己同値としての持上げ $p\in\operatorname{DPic}_A(\Lambda_A)$
が存在すれば、同じ還元を持つ $s$ で $s^2=t_A$ となるものが存在する。
さらに商 $A\twoheadrightarrow A'$ 上で $p_{A'}^2=t_{A'}$ が既に成立する場合、
$s_{A'}=p_{A'}$ となるように補正できる。
\end{proposition}
\begin{proof}
$d=p^2t_A^{-1}\in K_A$ とする。
Serre 関手の自然性より $p$ と $t_A$ は可換なので、$d$ も $p$ と可換する。
補題の一意な平方根を用いて $e=d^{-1/2}\in K_A$ と置く。
$pep^{-1}$ も $d^{-1}$ の平方根なので、一意性より $pe=ep$ である。
従って $s=ep$ は
\[
 s^2=e^2p^2=d^{-1}p^2=t_A
\]
を満たす。追加条件の下では $e_{A'}^2=1$ であり、$e_{A'}$ は
$K_{A'}$ に属する。同補題の一意性により $e_{A'}=1$ である。
\end{proof}

この命題は underlying Morita 自己同値の持上げが存在することを仮定する。
それ自体の障害を消したという主張ではない。また一般の基底上で
根忘却射が滑らかであるとは主張していない。

\begin{lemma}[合成と持上げ障害]
\label{lem:root-obstruction-composition}
$A\twoheadrightarrow A'$ を核 $I$ が $\mathfrak m I=0$ を満たす
small extension とする。$p'\in\operatorname{DPic}_{A'}(\Lambda_{A'})$
の閉 fiber を $p_0$ とし、$p'^2=t_{A'}$ と仮定する。
その持上げ障害を、一方の Morita 同定によって
\[
 b\in\operatorname{HH}^2(\Lambda)\otimes_kI
\]
と書くと、
\[
 b+(p_0)_*b=0.
\]
特に $(p_0)_*$ が $\operatorname{HH}^2(\Lambda)$ 上恒等なら、
$p'$ は $A$ 上の Serre 平方根へ持ち上がる。
\end{lemma}
\begin{proof}
$p'$ を表す有限 projective $\Lambda_{A'}$-双加群の有界複体を選ぶ。
相対 smoothness により可逆双加群は双加群として perfect である。
nilpotent ideal に沿う projective 加群の持上げによって、その graded
複体と微分を $\Lambda_A^e$ 上へ持ち上げる。
持ち上げた微分 $d$ の $d^2$ は $I$ に値を取り、そのコホモロジー類が
通常の perfect 複体の持上げ障害である。
$\mathfrak m I=0$ により、この群は
$\operatorname{Ext}^2_{\Lambda^e}(p_0,p_0)\otimes_kI$ と同定できる。
類が零なら次数一の項で $d$ を補正して $d^2=0$ とできる。

合成を表す tensor 複体の微分の平方は
\[
 (d\otimes1+1\otimes d)^2=d^2\otimes1+1\otimes d^2
\]
である（複体の tensor の符号規約を含む）。
具体的には、曲率類を $c_{p_0}:p_0\to p_0[2]$ とし、
$b=c_{p_0}\otimes p_0^{-1}$、
$(p_0)_*a=p_0\otimes a\otimes p_0^{-1}$ と同定する。
左右の Morita 同定の比較がこの $(p_0)_*$ なので、
合成 $p'^2$ の障害は $b+(p_0)_*b$ に対応する。
ところが $p'^2\simeq t_{A'}$ には canonical な持上げ $t_A$ があるため、
この障害は零である。

最後の仮定の下では $2b=0$ だから $b=0$ である。
得られる perfect 双加群の持上げは、閉 fiber で可逆なので可逆である。
実際、その左右随伴との評価・余評価の cone は閉 fiber で零であり、
perfect 複体の nilpotent derived Nakayama により零である。
命題~\ref{prop:artin-serre-square-correction} を適用して平方関係を補正する。
この補正は $A'$ 上の既存の根を保つ。
\end{proof}

\begin{corollary}[Hochschild 二次に自明に作用する根]
\label{cor:hhtrivial-root-artin}
$p_0^2=t_k$ とし、$(p_0)_*$ が $\operatorname{HH}^2(\Lambda)$ 上恒等とする。
任意の有限自由かつ相対 smooth な局所 Artin 変形 $\Lambda_A$ 上で、
$p_0$ は $s_A^2=t_A$ を満たす Morita 自己同値へ持ち上がる。
\end{corollary}
\begin{proof}
$A$ を small extension の有限列に分解し、直前の補題を帰納的に適用する。
各段の障害に作用するのは閉 fiber の $(p_0)_*$ であり、
途中の持上げの選択によってこの恒等作用は変わらない。
\end{proof}

\subsection{可換 quadric における根の存在と四つの一次条件}

以下は $Q=\mathbb P^1\times\mathbb P^1$ の full strong foundation を用いた
有限自由相対 smooth Artin Morita 変形についての結論である。
その閉 fiber の中心は $H^0(Q,\mathcal O_Q)=k$ なので、上の議論が適用できる。
Abel 圏の変形への適用には、その foundation の相対持上げと
heart との比較を別途用いる。

\begin{proposition}[恒等幾何型の根は全 Artin Morita 変形に延長する]
\label{prop:quadric-root-all-artin}
$s_0=(-)\otimes\mathcal O_Q(1,1)$ は、$Q$ の任意の上記 Artin Morita 変形へ
相対 Serre 平方根として延長する。
\end{proposition}
\begin{proof}
$s_0^2=(-)\otimes\mathcal O_Q(2,2)=S_Q^{-1}[2]$ である。
また $H^1(T_Q)=H^2(\mathcal O_Q)=0$ なので
\[
 \operatorname{HH}^2(Q)=H^0(Q,\textstyle\bigwedge^2T_Q)
                     =H^0(Q,\mathcal O_Q(2,2)).
\]
線束 tensor は局所的に恒等関手であり、この次数では他の HKR 成分がない。
従って局所 bivector への制限の単射性により、$s_0$ はこの九次元空間に
恒等に作用する。系~\ref{cor:hhtrivial-root-artin} を適用する。
\end{proof}

\begin{corollary}[可換点の全 Abel 変形における exact ample な根]
\label{cor:quadric-exact-root-all-artin}
$(\operatorname{coh}Q,\mathcal O_Q)$ の任意の flat pointed Abel 変形
$(\mathcal C_A,\mathcal O_A)$ には、$s_0=(-)\otimes\mathcal O_Q(1,1)$ を還元とし、
$s_A^2\simeq S_A^{-1}[2]$ を満たす完全自己同値 $s_A$ が存在する。
$(\mathcal O_A,s_A)$ は AZ ample で、その全軌道環は
各次数の階数 $(n+1)^2$ の自由 $A$-加群からなる二次 noetherian 環である。
適合する AS4 被覆は定理~\ref{thm:quadric-all-artin-h2} で構成する。
\end{corollary}
\begin{proof}
定理~\ref{thm:artin-foundation-lift} により有限 foundation の変形に移し、
前命題で Morita 平方根を持ち上げる。
補題~\ref{lem:artin-morita-heart-exact} が heart 上の完全性を与える。
特殊ファイバーの $s_0$ は豊富な線束による tensor なので、
Serre 消滅と補題~\ref{lem:artin-ample-lift} から相対 AZ 豊富性が従う。
$s_A^n\mathcal O_A$ は normal helix の $E_{2n}$ の一意な flat lift であり、
De Deken--Lowen の geometric Hom の平坦性から全軌道の階数と基底変更を得る。
二次性は補題~\ref{lem:relative-quadratic-from-fibers}、
noetherian 性は補題~\ref{lem:artin-noetherian-lift} による。
\end{proof}

\begin{lemma}[近恒等 Morita 群の次数商]
\label{lem:artin-dpic-associated-graded}
補題~\ref{lem:artin-infinitesimal-dpic} の設定で、$r\geq1$ に対し
\[
 K_A^r=\ker\bigl(\operatorname{DPic}_A(\Lambda_A)
       \longrightarrow
       \operatorname{DPic}_{A/\mathfrak m^r}
                    (\Lambda_A/\mathfrak m^r\Lambda_A)\bigr)
\]
と置く。$K_A^1=K_A$ であり、$\mathfrak m^N=0$ なら $K_A^N=1$ である。
先頭の導分類は加法群の単射
\[
 \ell_r:K_A^r/K_A^{r+1}\hookrightarrow
 \operatorname{HH}^1(\Lambda)\otimes_k
                  (\mathfrak m^r/\mathfrak m^{r+1})
\]
を与える。任意の $p\in\operatorname{DPic}_A(\Lambda_A)$ による共役は
この filtration を保ち、$\ell_r$ の像上では閉 fiber $p_0$ の
Hochschild 一次への作用と一致する。
\end{lemma}
\begin{proof}
まず任意の商 $B=A/\mathfrak m^r$ に対して
$Z(\Lambda_A\otimes_AB)=B$ である。
実際、中心元の閉 fiber は $Z(\Lambda)=k$ に属する。
そのスカラーを引き、$\mathfrak m$-進次数ごとに同じ議論を繰り返す。
各次数の中心条件は
$\Lambda\otimes_k\operatorname{gr}_{\mathfrak m}^jB$ 上の条件であり、
その各係数は $Z(\Lambda)=k$ に属するので、その次数のスカラーを引ける。
有限回で終わるため中心元は $B$ の元である。

$u\in K_A^r$ を取る。補題~\ref{lem:artin-infinitesimal-dpic} により
$u$ は環の自己同型 $\alpha$ による twist で表せる。
$A/\mathfrak m^r$ 上で $u$ は恒等なので、そこでの $\alpha$ は inner である。
その共役単元を $\Lambda_A$ へ持ち上げて補正すれば
$\alpha\equiv1\pmod{\mathfrak m^r}$ とできる。
従って $\alpha-1$ の先頭部分は
\[
 \delta_\alpha:\Lambda\longrightarrow
       \Lambda\otimes_k(\mathfrak m^r/\mathfrak m^{r+1})
\]
という $k$-導分を与える。
同じ類を表す別の自己同型との差が $\operatorname{Ad}(v)$ なら、
$v$ の $A/\mathfrak m^r$ 上の像は中心単元である。
上で証明した中心の等式により、スカラー単元で割って
$v\in1+\mathfrak m^r\Lambda_A$ とできる。
従って $\delta_\alpha$ の変化は inner 導分であり、$\ell_r$ は定義される。
逆に先頭導分が inner なら、その係数を $\mathfrak m^r\Lambda_A$ に持ち上げ、
対応する $1+x$ による共役で補正すると
$\alpha\equiv1\pmod{\mathfrak m^{r+1}}$ となる。
従って核はちょうど $K_A^{r+1}$ である。
積の先頭部分は導分の和なので、これは加法群の単射である。

最後に、この先頭導分類は恒等双加群の一次変形類である。
可逆双加群で左右から tensor して共役する操作は、その一次類を
Hochschild 一次上の Morita 作用で移す。
係数 $\mathfrak m^r/\mathfrak m^{r+1}$ は $\mathfrak m$ で零になるので、
この作用には $p$ の閉 fiber $p_0$ だけが現れる。
これが必要な同変性である。
\end{proof}

\begin{proposition}[Hochschild 一次に自明に作用する Serre 根の一意性]
\label{prop:artin-serre-root-uniqueness}
$s_0^2=t_k$ とし、$(s_0)_*$ が
$\operatorname{HH}^1(\Lambda)$ 上恒等であるとする。
$s_0$ を還元とし $s_A^2=t_A$ を満たす
$\operatorname{DPic}_A(\Lambda_A)$ の元は、存在すれば一意である。
また、そのような $s_A$ は $K_A$ の全ての元と可換である。
ここでの一意性は Morita 自己同値の同型類についての主張であり、
平方との同型そのものを一意に選ぶという意味ではない。
\end{proposition}
\begin{proof}
$s=s_A$ を一つ固定し、$\sigma=\operatorname{Ad}(s)|_{K_A}$ と置く。
前補題により $\sigma$ は各次数商に恒等に作用する。
$u\in K_A$ が $u\sigma(u)=1$ を満たすなら、
$u\in K_A^r$ から
\[
 0=\ell_r(u)+\ell_r(\sigma(u))=2\ell_r(u)
\]
を得る。標数零なので $u\in K_A^{r+1}$ であり、帰納的に $u=1$ となる。
別の根を $s'=us$ と書くと、$s'^2=s^2$ はちょうど
$u\sigma(u)=1$ を意味するから $s'=s$ である。

さらに $s^2=t_A$ は中心元なので $\sigma^2=1$ である。
任意の $g\in K_A$ に対し $v=\sigma(g)g^{-1}$ と置くと
$\sigma(v)v=1$、従って $v\sigma(v)=1$ である。
同じ帰納法により $v=1$、すなわち $\sigma(g)=g$ となる。
\end{proof}

\begin{corollary}[可換 quadric の恒等幾何型の根の一意性]
\label{cor:quadric-root-unique-all-artin}
$Q=\mathbb P^1\times\mathbb P^1$ に対し、
$s_0=(-)\otimes\mathcal O_Q(1,1)$ の相対 Serre 平方根としての持上げは、
命題~\ref{prop:quadric-root-all-artin} の任意の Artin Morita 変形上で
Morita 自己同値の同型類として一意である。
\end{corollary}
\begin{proof}
存在は命題~\ref{prop:quadric-root-all-artin} による。
$H^1(Q,\mathcal O_Q)=0$ なので
$\operatorname{HH}^1(Q)=H^0(Q,T_Q)$ である。
線束 tensor は局所的に恒等関手であり、稠密開集合へのベクトル場の制限は
単射なので、$s_0$ はこの Hochschild 一次に恒等に作用する。
前命題を適用する。
\end{proof}

\begin{proposition}[四つの根の一次持上げ条件]
\label{prop:quadric-four-root-tangents}
$k=\mathbb C$ とし、$Q$ の一次変形類
$u\in\operatorname{HH}^2(Q)$ を取る。
定理~\cite[\texttt{thm:commutative-four-roots}]{QuadricIntrinsic2026} の四代表 $s_\varphi$ の underlying
Morita 自己同値が同じ一次変形へ延長するための必要十分条件は
\[
 (s_\varphi)_*u=u.
\]
この条件が成立すれば、平方関係も満たすように持上げを選べる。
固定部分空間の次元は、同定理の順に $9,3,5,3$ である。
\end{proposition}
\begin{proof}
必要性は補題~\ref{lem:autoequivalence-deformation-obstruction}、十分性は
Toda~\cite[Theorem 4.7, arXiv v3]{Toda2009} による。
得られた持上げの平方関係は
命題~\ref{prop:artin-serre-square-correction} で補正できる。

$V=H^0(\mathbb P^1,T_{\mathbb P^1})$ と置くと
$\operatorname{HH}^2(Q)=V\otimes V$ である。
対合 $z\mapsto-z$ は $\partial_z,z\partial_z,z^2\partial_z$ に対し
固有値 $-1,1,-1$ を持つので、$V$ の正負固有空間の次元は $1,2$ である。
従って一因子対合の固定空間は次元 $1\cdot3=3$、
二因子対合では $1\cdot1+2\cdot2=5$ となる。
因子交換は bivector の wedge の順序を逆転するので、$V\otimes V$ には
単なる転置ではなく $v\otimes w\mapsto-w\otimes v$ と作用する。
その固定空間は $\bigwedge^2V$、次元 $3$ である。
恒等幾何型では全空間が固定され、次元は $9$ である。
\end{proof}

従って、可換 quadric で選んだ非自明な根の持上げ障害から、
「いずれの根も存在しない」一次方向を得ることはできない。
恒等幾何型の根は全九次元の一次方向を覆い、さらに上述の有限 Artin
Morita 変形の全段階へ持ち上がる。
この根に適合する二次 AS 正則被覆は次節で構成する。


\subsection{可換点の近傍における Morita 平方根の代数化}
\label{subsec:quadric-etale-morita-root}

本項では Abel heart の代数化に先立ち、有限 foundation の
\emph{Morita 自己同値としての}平方根を代数化する。
根を持つ有限代数の族から相対 $\mathfrak H_2$ 族が得られるとは、ここでは主張しない。

\begin{proposition}[標準根を持つ \'etale な foundation 近傍]
\label{prop:quadric-etale-morita-root}
$k$ を標数零の代数閉体とし、$S$ を有限 foundation スタックの
滑らかな有限型 affine chart とする。
$\Lambda$ を $S$ 上の普遍有限自由有向代数、$x\in S(k)$ を可換 quadric
$Q=\mathbf P^1_k\times\mathbf P^1_k$ の foundation に対応する点とする。
次の性質を用いる。
\begin{enumerate}
\item $\operatorname{Perf}(\Lambda)$ は相対 smooth proper であり、
  任意の基底変更後も $\operatorname{HH}^{<0}(\Lambda)=0$、
  $\operatorname{HH}^{0}(\Lambda)=\mathcal O_S$ である。
\item 相対 inverse Serre $t=S_{\Lambda}^{-1}[2]$ は可逆 perfect 双加群で表され、
  任意の基底変更と可換する。
\end{enumerate}
これらは補題~\ref{lem:foundation-relative-smooth-proper}、
補題~\ref{lem:foundation-relative-antiserre-perfect} と以下の普遍中心の計算により成立する。

$\mathcal G\to S$ を $\operatorname{Perf}(\Lambda)$ の Morita 自己同値の
\'etale 局所同型類の層とする。これは $S$ 上局所有限表示の代数空間群である。
$\tau\in\mathcal G(S)$ を $t$ の類とし、平方写像を
$\mathrm{sq}:\mathcal G\to\mathcal G$, $p\mapsto p^2$ と置く。
この写像は一般には群準同型ではないが、代数空間の射である。
平方根空間
\[
 \mathcal Z=\mathcal G\mathbin{\times}_{\mathrm{sq},\mathcal G,\tau}S
\]
は $S$ 上局所有限表示の代数空間となる。

$p_0$ を $s_0=(-)\otimes\mathcal O_Q(1,1)$ の Morita 類とする。
$\mathcal Z\to S$ は $p_0$ において \'etale である。
さらに $x$ を含む開近傍を必要に応じて縮めると、その \'etale 被覆
$U\to S$ 上に可逆 perfect $\Lambda_U$-双加群 $P$ と実際の同型
\[
 P\otimes^{\mathbf L}_{\Lambda_U}P\simeq t_U
\]
が存在し、選んだ $u\in U(k)$ における $P$ の類は $p_0$ である。
\end{proposition}

\begin{proof}
任意の基底変更後も有向代数の対角成分は基底環である。
中心元は頂点の冪等元と可換するので対角であり、隣接頂点間の正階数の自由加群と
可換する条件から各対角スカラーは等しい。従って中心は基底環そのものである。
負の Hochschild 次数の消滅は通常の代数の自己 Ext の消滅による。

$S=\operatorname{Spec}R$ として
$T=\operatorname{Perf}(\Lambda)$ を $R$-線形 saturated dg 圏とみなす。
\cite[著者公開版 April 2007, Corollary 3.24]{TV2007} は、
基礎環が体でなくても、$\operatorname{HH}^{<0}(T)=0$ と
$R\simeq\operatorname{HH}^{0}(T)$ の下で
$\operatorname{aut}(T)$ が局所有限表示の代数空間で表されることを述べる。
これを $R$ 上で適用すれば $\mathcal G$ が得られる。
ここで層化は essential であり、基底上の線束による tensor twist は
局所的には同じ類になる。任意の基底上の実際の双加群の同型類を、
層化せずそのままこの代数空間と同一視しない。

同論文 Lemma 3.25 は、perfect 双加群が Morita 同値を与える条件を
二つの評価射の quasi-isomorphism として表し、その条件が開であることを示す。
可逆双加群の自己 Ext は Hochschild cohomology と同一視されるので、
仮定 (1) より全て単純であり、負次数自己 Ext は零である。
従って Corollary 3.22 により、実際の可逆双加群のスタックから
$\mathcal G$ への射は \'etale 局所に自明な $\mathbf G_m$-gerbe である。
$S$ は Noetherian で、$\mathcal G$ はその上局所有限表示だから、
$S$ 上の射 $\mathrm{sq}$ も局所有限表示である。
従って $\mathcal Z$ の代数性と局所有限表示性は上のファイバー積から従う。

\'Etale 性を $p_0$ で調べる。
剰余体 $k$ の局所 Artin $k$-代数 $A$ と、閉点を $x$ に送る
$a:\operatorname{Spec}A\to S$ を任意に取る。
$\Lambda_A$ は $\Lambda_x$ の有限平坦な代数変形である。
命題~\ref{prop:quadric-root-all-artin} と
系~\ref{cor:quadric-root-unique-all-artin} により、
$p_0$ を還元に持つ $p_A\in\operatorname{DPic}_A(\Lambda_A)$ で
$p_A^2=t_A$ を満たすものが存在し、その同型類は一意である。
これらの命題は有限代数の Morita 変形についての結果であり、
$\Lambda_A$ に相対 $\mathfrak H_2$ heart が先に与えられていることを要求しない。
存在には $(s_0)_*=1$ on $\operatorname{HH}^2(Q)$、
一意性には $(s_0)_*=1$ on $\operatorname{HH}^1(Q)$ を用いる。

局所 Artin $A$ 上では線束は自明であり、$A$ は strictly henselian なので
上の \'etale gerbe は中立である。
従ってこの存在・一意性は、閉点を固定した変形関手の同型
\[
 \operatorname{Def}_{\mathcal Z,p_0}(A)
 \xrightarrow{\ \sim\ }
 \operatorname{Def}_{S,x}(A)
\]
を意味する。特に平方との比較同型そのものを標識として数えてはいない。

$p_0$ を持ち上げる \'etale なスキーム chart $V\to\mathcal Z$ と
$v\in V(k)$ を取る。\'Etale な chart は閉点を固定した Artin 変形関手を変えないから、
自然な準同型
\[
 \widehat{\mathcal O}_{S,x}\longrightarrow
 \widehat{\mathcal O}_{V,v}
\]
は同型である。
$S$ は滑らかなので右辺は $k$ 上の形式冪級数環である。
従って $V$ は $v$ で正則、したがって完全体 $k$ 上滑らかであり、
$V\to S$ の接写像は同型である。
滑らかな $k$-スキーム間の Jacobian 判定により $V\to S$ は $v$ で \'etale となる。
\'Etale 局所性から $\mathcal Z\to S$ も $p_0$ で \'etale である。
その \'etale locus の像は $x$ を含む開部分である。

この locus 上で前記の $\mathbf G_m$-gerbe を \'etale 局所に中立化し、
実際の可逆 perfect 双加群 $P$ を取る。
$[P]^2=\tau$ は、まだ大域的な平方比較の選択ではない。
定理~\ref{thm:bimodule-line-comparison} により比較線束
\[
 L=H^0\mathbf R\mathcal Hom_{\Lambda_U^e}
       (P\otimes^{\mathbf L}_{\Lambda_U}P,t_U)
\]
と評価同型
$(P\otimes^{\mathbf L}_{\Lambda_U}P)\otimes L\simeq t_U$ を得る。
さらに Zariski 被覆で $L$ を自明化すると、主張した実際の平方同型が得られる。
これらの細分はいずれも \'etale であり、必要なら有限個の affine chart に制限できる。
\end{proof}

\begin{proposition}[標準根の有限窓と相対 cubic tensor]
\label{prop:quadric-etale-finite-window}
命題~\ref{prop:quadric-etale-morita-root} の可換点 $u$ を含む
$U$ をさらに Zariski 開に縮める。
普遍 foundation の四つの projective 対象を $F_0,\ldots,F_3$ とし、
閉ファイバーで順に
\[
 \mathcal O_Q,\quad \mathcal O_Q(0,1),\quad
 \mathcal O_Q(1,1),\quad \mathcal O_Q(1,2)
\]
に対応させる。perfect Morita 根 $P$ が与える自己同値を
$\mathsf s$ と書く。このとき次が成立する。
\begin{enumerate}
\item $U$ 上で
  \[
     \mathsf sF_0\simeq F_2,\qquad
     \mathsf sF_1\simeq F_3
  \]
  となる同型を取れる。各同型の選択は単元倍を除いて一意であり、
  一度選べば任意の $U$-基底変更に対して同型のままである。
\item 全整数 $m$ に対して
  \[
    L_{2m}=\mathsf s^m F_0,\qquad
    L_{2m+1}=\mathsf s^m F_1
  \]
  と置くと、任意の連続四項 $(L_i,L_{i+1},L_{i+2},L_{i+3})$
  は相対的な full strong exceptional collection である。
  その derived Hom は任意の基底変更と可換し、
  次数零の有限局所自由加群となる。
\item $D_{ij}=\operatorname{Hom}(L_i,L_j)$ と置く。
  $1\leq j-i\leq4$ に対して higher Ext は消え、
  \[
    \operatorname{rk}D_{i,i+r}=2,4,6,9
       \quad(r=1,2,3,4).
  \]
  特に $V_i=D_{i,i+1}$ は階数二であり、
  $V_i\otimes V_{i+1}\xrightarrow{\sim}D_{i,i+2}$、
  $V_i\otimes V_{i+1}\otimes V_{i+2}\twoheadrightarrow
     D_{i,i+3}$ が成り立つ。
  後者の核 $R_i$ は階数二の直和因子である。
  さらに四重積は $D_{i,i+4}$ に全射で、
  \[
    W_i=(R_i\otimes V_{i+3})\cap
              (V_i\otimes R_{i+1})
       \subset V_i\otimes V_{i+1}\otimes V_{i+2}\otimes V_{i+3}
  \]
  は階数一の直和因子である。
  $W_i$ が誘導する $R_i$ と $V_{i+3}^{\vee}$、
  および $V_i^{\vee}$ と $R_{i+1}$ の二つの pairing は非退化である。
  これらの階数・核・pairing は任意の $U$-基底変更と可換する。
\end{enumerate}
これは全ての \emph{有限窓} における主張であり、
$D_{ij}$ の任意の長さでの次数零集中・三次提示、
相対 tails heart、AZ 豊富性や相対 $\mathfrak H_2$ 実現は結論しない。
\end{proposition}
\begin{proof}
まず $K_a=\mathbf R\operatorname{Hom}_{\Lambda_U}
       (\mathsf sF_a,F_{a+2})$（$a=0,1$）を取る。
$\Lambda_U$ の相対 smooth proper 性と $P$ の perfect 性から、
これは $U$ 上 perfect で、任意の導来基底変更と可換する。
$u$ で $\mathsf sF_a$ と $F_{a+2}$ は同じ線束であり、
$K_a\otimes^{\mathbf L}k(u)\simeq k(u)[0]$ である。
perfect 複体の局所的な有限自由表示で単元成分を消去すると、
$u$ の近傍で $K_a$ は次数零の線束になる。
その線束の閉ファイバーで恒等同型に対応する非零元を
局所的な切断に持ち上げ、対応する射
$\mathsf sF_a\to F_{a+2}$ を得る。
射の cone は基底上 perfect で閉ファイバーでは零だから、
perfect Nakayama によりさらに縮めれば零となる。
従って射は同型であり、全ての基底変更でも同型である。
$\operatorname{End}(F_{a+2})=\mathcal O_U$ なので、
二つの同型は単元倍だけ異なる。

偶数添字から始まる四項は
$\mathsf s^m(F_0,F_1,F_2,F_3)$ と同型である。
元の四項は foundation の projective 対象であり、
全ての基底変更上で full strong exceptional だから (2) は
偶数窓では従う。
奇数窓は $\mathsf s^m(F_1,F_2,F_3,tF_0)$ である
（$\mathsf s^2=t=S^{-1}[2]$）。
full exceptional collection の Serre rotation により、
$(F_1,F_2,F_3,S^{-1}F_0)$ は相対的に full exceptional である。
最後の項を $[2]$ だけ移しても例外性と生成性は保たれる。
$u$ 上ではこの奇数窓が
\[
  \bigl(\mathcal O_Q(0,1),\mathcal O_Q(1,1),
        \mathcal O_Q(1,2),\mathcal O_Q(2,2)\bigr)
\]
となり、K\"unneth 公式から全ての higher Ext と
逆向きの全 Ext が消える。
関係する有限個の相対 Hom 複体が perfect であることから、
もう一度縮めればそれぞれ次数零の所定の階数の束、
または零となる。よって奇数窓も full strong で、
この性質は任意の基底変更で保たれる。

(3) は $\mathsf s$ による二項移動で $i=0,1$ に帰着する。
閉ファイバーで隣接の bidegree は $(0,1),(1,0)$ が交互に現れ、
距離四は $(2,2)$ だから、K\"unneth 公式によって
前記の四つの階数と高次消滅を得る。
関連する有限個の perfect Hom 複体を同時に縮めれば、
全ての $D_{i,i+r}$ が次数零の有限局所自由加群となる。
二重・三重・四重積の全射性は閉ファイバーの
$\mathbf P^1\times\mathbf P^1$ の斉次座標環の積の全射性から従い、
有限局所自由加群の射の最大 minor を可逆化すれば
基底変更後にも成立する。二重積は同階数なので同型である。
三重積の核 $R_i$ は階数二の直和因子となる。

四重積の核は階数 $16-9=7$ である。
$R_i\otimes V_{i+3}$ と $V_i\otimes R_{i+1}$
はいずれもその核に含まれる。閉ファイバーでは
可換 quadric の cubic 解消により、この二つの和は
七次元で四重積の核全体と一致する。
従って七次 minor を可逆化した開近傍では、
和から核への射は全射であり、その核すなわち $W_i$ は
階数 $4+4-7=1$ の直和因子となる。
閉ファイバーの標準 cubic syzygy は二つの pairing について
非退化なので、その行列式が単元となる開近傍まで縮めれば
相対的にも非退化である。
全て有限個の条件を二つの parity について課しただけであるから、
$\mathsf s$ による移動によって全整数 $i$ に同時に成立する。
\end{proof}


\begin{theorem}[可換点の局所環上の全次数 cubic 提示]
\label{thm:quadric-local-all-degree-cubic}
命題~\ref{prop:quadric-etale-finite-window} の $u\in U(k)$ において
$A=\mathcal O_{U,u}$ と置き、根 $\mathsf s$ と全列 $L_i$ を
$\operatorname{Spec}A$ に引き戻す。以下は体 $k$ 上の可換点を
閉点とする局所環 $A$ 上で成立する。
\begin{enumerate}
\item 全ての整数 $i\le j$ について、canonical な導来 Hom は
  任意の $A$-代数への導来基底変更と可換し、
  \[
    \mathbf R\operatorname{Hom}(L_i,L_j)
       \simeq A^{\,h_{j-i}}[0],\qquad
    h_{2m}=(m+1)^2,\quad
    h_{2m+1}=(m+1)(m+2)\quad(m\ge0)
  \]
  となる。等号は局所自由加群としての存在であり、
  基底の標準的な選択を意味しない。
\item $D_{ij}=\operatorname{Hom}(L_i,L_j)$（$i\le j$）とし、
  $V_i=D_{i,i+1}$、
  \[
    R_i=\ker\bigl(
      V_i\otimes_A V_{i+1}\otimes_A V_{i+2}
             \longrightarrow D_{i,i+3}\bigr)
  \]
  と置くと、$V_i$ は階数二、$R_i$ は階数二の直和因子であり、
  $D=\bigoplus_{i\le j}D_{ij}$ は全内部次数で
  \[
     D \simeq T_{\mathbb Z,A}(V_i)/(R_i)
     \tag{Loc-cub}\label{eq:quadric-local-cubic-presentation}
  \]
  と表示される。左辺の積は射の合成である。
  この表示、次数別の有限自由性、及び相対 Hom との同定は
  任意の $A$-代数への基底変更と整合する。
\item 右軌道積を用いた偶数環
  \[
    B_A=\bigoplus_{n\ge0}\operatorname{Hom}
                (L_0,\mathsf s^nL_0)
  \]
  は全次数で階数 $(n+1)^2$ の自由 $A$-加群であり、
  \[
      B_A\simeq T_A(B_{A,1})/(J_A),\qquad
      \operatorname{rk}B_{A,1}=4,\quad
      \operatorname{rk}J_A=7.
  \]
  この二次環は相対 Koszul である。
  その左右 Koszul 複体は各内部次数で分裂完全であり、
  任意の $A$-代数への基底変更後も完全である。
\item $P_i=e_iD$ を右表現可能加群とし、
  $S_i$ を次数 $i$ の成分が $A$、他が零の頂点加群
  （体上の単純加群の相対版）とすると、
  \[
  0\longrightarrow W_i\otimes_A P_{i+4}
   \longrightarrow R_i\otimes_A P_{i+3}
   \longrightarrow V_i\otimes_A P_{i+1}
   \longrightarrow P_i\longrightarrow S_i\longrightarrow0
  \tag{Loc-res}\label{eq:quadric-local-cubic-resolution}
  \]
  は相対的な標準 cubic 解消である。
  ここで $W_i$ は命題~\ref{prop:quadric-etale-finite-window}
  の階数一の重なり束である。
  各内部次数は $A$-加群として分裂完全で、
  任意の $A$-代数へ tensor しても完全である。
\end{enumerate}
ここで $\operatorname{Spec}A$ は一般に $k$ 上有限型ではない。
本定理は $B_A$ の左右 noetherian 性、相対 tails heart、
Serre・AZ 条件、相対 AS4 被覆の存在、
あるいは $\mathfrak H_2$ の代数性を主張しない。
\end{theorem}
\begin{proof}
$\Lambda_U$ の相対 smooth proper 性と $\mathsf s$ の perfect
Morita 表示から、任意の $i,j$ について
$\mathbf R\operatorname{Hom}(L_i,L_j)$ は基底上 perfect であり、
導来基底変更と可換する。
閉点で $L_{2a}=\mathcal O_Q(a,a)$、
$L_{2a+1}=\mathcal O_Q(a,a+1)$ である。
$i\le j$ のとき、その差は非負の bidegree だから、
K\"unneth 公式より導来 Hom は次数零に集中し、
次元は表示した $h_{j-i}$ となる。
局所環上の perfect 複体の有限自由表示から
単元係数を消去すると、剰余体への導来 tensor が
次数零だけにある場合、元の複体はその階数の自由加群を
次数零に置いたものに同型である。
これは各 $(i,j)$ に個別に適用でき、同じ局所環 $A$ を
さらに縮める必要はない。従って (1) が従う。

任意の $i<j$ について、隣接 $V_a$ の tensor 積から
$D_{ij}$ への合成射は有限自由加群の射である。
閉点では $\mathbb P^1\times\mathbb P^1$ の線束切断の積であり、
全射であるから Nakayama により $A$ 上でも全射である。
三重積の核 $R_i$ は階数二の直和因子になる。
三次関係による tensor $\mathbb Z$-代数の商を $D'$ とすると、
合成により全内部次数で $D'_{ij}\twoheadrightarrow D_{ij}$
を得る。各 $D'_{ij}$ は有限生成 $A$-加群である。

閉点における $D'$ の同定は直接確認できる。
偶数・奇数の隣接矢空間は交互に二次元空間 $V,U$ であり、
三次関係はそれぞれ
$\bigwedge^2V\otimes U$ と
$\bigwedge^2U\otimes V$（因子の並べ方は隣接順）
に対応する。
従って三次関係は同じ parity の隣り合う出現を
交換する関係であり、これらの交換が全ての置換を生成する。
ゆえに長さ $n$ の商は二色それぞれの対称冪の tensor 積で、
$D_{ij}\otimes_Ak$ に一致する。
各 $D'_{ij}\to D_{ij}$ は剰余体上同型であり、
終域が自由なので核との短完全列は $A$-加群として分裂する。
核は有限生成かつ剰余体上零だから Nakayama により零となる。
これで (2) の全次数と任意基底変更での主張を得る。

(1) により $B_{A,n}$ は全て自由で階数 $(n+1)^2$、
特殊ファイバーは
$\bigoplus_{n\ge0}H^0(Q,\mathcal O_Q(n,n))$
という Segre 環である。
$B_{A,1}^{\otimes n}\to B_{A,n}$ は特殊ファイバーで全射だから
$A$ 上も全射であり、$J_A$ は階数 $16-9=7$ の直和因子となる。
二次関係による商から $B_A$ への全射は
剰余体上、Segre 環の既知の二次提示と一致する。
(2) と同様、各次数で自由な終域へ分裂する全射の核に
Nakayama を適用して二次提示を得る。
Segre 環は二次 Koszul 環であり、$B_{A,3}$ は自由なので、
補題~\ref{lem:relative-koszul-degree-three} の証明にある
Polishchuk--Positselski の局所 Koszul 変形原理
\cite[Chapter 6, Theorem 7.3 and Propositions 7.1--7.2]
{PolishchukPositselski2005}
を局所環 $A$ とその閉ファイバーに適用する。
関係空間の格子は自由な直和因子に分解され、
左右の各内部次数の Koszul 複体は有限自由で分裂完全となる。
この性質は任意の $A$-代数への tensor で保たれる。

最後に (2) による全次数の提示と
命題~\ref{prop:quadric-etale-finite-window} の $W_i$ を使って
(\ref{eq:quadric-local-cubic-resolution}) の微分を定義する。
$W_i\subset R_i\otimes V_{i+3}$ という包含と、
$R_i$ の関係式から、これは複体になる。
各内部次数では (1) と (2) により有限自由 $A$-加群の
有界複体である。その閉ファイバーは
$\mathbb P^1\times\mathbb P^1$ の標準 cubic
$\mathbb Z$-代数の解消と一致して完全である
\cite[\S4.2]{VdB2011}。
局所環上の有限自由複体が剰余体への導来 tensor
の下で完全なら、単元成分を消去して
元の複体も分裂完全になる。
内部次数ごとに適用して (4) を得る。
無限個の内部次数について一つの有限型開を
選んだという主張ではない。
\end{proof}

\begin{corollary}[各有限次数の開近傍への拡張]
\label{cor:quadric-each-finite-degree-open}
任意の整数 $N\ge0$ について、$u$ を含む Zariski 開部分
$U_N\subset U$ が存在し、上の導来 Hom の次数零集中と
階数、cubic 提示の次数 $N$ までの同定、及び
$B$ の二次提示の次数 $N$ までの同定が成立し、
任意の $U_N$-基底変更と可換する。
この系の有限次数の議論だけからは、
$U_N$ を $N$ に依存せずに選べるとは結論しない。
同一開上の全次数の主張は後の
命題~\ref{prop:quadric-uniform-cubic-hom-open} で別に証明する。
\end{corollary}
\begin{proof}
$\mathsf s$ による二項移動により、必要な derived Hom と積・関係の
条件は始点の二つの parity 及び $0\le j-i\le N$
の有限個の次数へ帰着する。
定理~\ref{thm:quadric-local-all-degree-cubic} の局所環上の同型は、
有限表示加群の射の核・余核の消滅と perfect 複体の自由表示で
記述できるので、有限個の分母を可逆にすることで
同時に $u$ の開近傍へ延長できる。
この開近傍で局所自由性・分裂性を保てば任意基底変更にも保たれる。
\end{proof}


\begin{proposition}[幾何的 cubic ファイバーの有限型開近傍]
\label{prop:quadric-geometric-quintuple-open}
命題~\ref{prop:quadric-etale-finite-window} の可換点
$u\in U(k)$ を含む Zariski 開部分
$U_{\mathrm{geo}}\subset U$ があり、全ての幾何点
$\bar v\to U_{\mathrm{geo}}$ に対して
\[
 D^{\mathrm{pres}}_{\bar v}
       =T_{\mathbb Z,k(\bar v)}(V_{i,\bar v})/
                                  (R_{i,\bar v})
\]
は幾何的 cubic AS 正則 $\mathbb Z$-代数である。
同じ性質は任意の非被約な有限型基底変更後も幾何ファイバーごとに成り立つ。
ただし $D^{\mathrm{pres}}$ の全次数の相対局所自由性、
noetherian 性、$D^{\mathrm{pres}}_{ij}$ と
$\operatorname{Hom}(L_i,L_j)$ の全次数の一致、
相対 tails や $\mathfrak H_2$ の効果性は結論しない。
\end{proposition}
\begin{proof}
命題~\ref{prop:quadric-etale-finite-window} により
$W_i\subset V_i\otimes V_{i+1}\otimes V_{i+2}\otimes V_{i+3}$
は線束であり、$R_i$ と $V_{i+3}^{\vee}$、
$V_i^{\vee}$ と $R_{i+1}$ への pairing は非退化である。
根 $\mathsf s$ の二項移動により、$i=0,1$ の二つを調べればよい。

$i=0,1$ と巡回する四つの隣接因子対をそれぞれ固定する。
$W_i$ の元に対する二つの非零 covector の縮約が零になる条件は、
該当する二つの双対射影束の積
$\mathbf P(V_{a}^{\vee})\times_U\mathbf P(V_{b}^{\vee})$
上の vector bundle の切断の零条件として閉である。
この射影束の積から $U$ への射は proper なので、
零集合の像も閉である。
それらの有限個の像を除けば、全幾何ファイバーで
Van den Bergh の geometric quintuple の縮約非零条件が成り立つ。
可換 $Q$ の $W_i$ は二つの二次元 symplectic
coevaluation の tensor であり、該当する全縮約は非零である。
よって得られる開部分は $u$ を含む。

非退化 pairing から $W_i$ の縮約で $R_i,R_{i+1}$ を回復する。
$W_i\subset V_i\otimes R_{i+1}$ と
$W_{i+1}\subset R_{i+1}\otimes V_{i+4}$ の
coevaluation を合成して $V_i\simeq V_{i+4}$ とすれば、
二つの $W$ はこの同型による巡回で一致する
（補題~\ref{lem:relative-cubic-tensors-artin} と同じ線形代数）。
従って全関係は一つの geometric quintuple の標準再構成であり、
体上の分類 \cite[arXiv v3, Theorem 5.5.10]{VdB2011} により
各幾何ファイバーの cubic AS 正則性が従う。
体上の定理を相対 flatness の定理に読み替えてはいない。
\end{proof}


\begin{lemma}[幾何的 cubic 表示の相対全次数平坦性]
\label{lem:relative-geometric-cubic-flatness}
$S$ を局所 Noetherian スキームとし、
$V_i$ を階数二の局所自由層、
$R_i\subset V_i\otimes V_{i+1}\otimes V_{i+2}$
を階数二の直和因子とする。
$D=T_{\mathbb Z,\mathcal O_S}(V_i)/(R_i)$ と置く。
全幾何ファイバーが三次元 cubic AS 正則
$\mathbb Z$-代数であり、標準解消に現れる
\[
 W_i=(R_i\otimes V_{i+3})\cap(V_i\otimes R_{i+1})
\]
が $S$ 上の階数一の直和因子で、全幾何ファイバーの
syzygy と一致すると仮定する。
このとき全ての $i\le j$ について $D_{ij}$ は
有限局所自由で、
\[
 \operatorname{rk}D_{i,i+2m}=(m+1)^2,\qquad
 \operatorname{rk}D_{i,i+2m+1}=(m+1)(m+2).
\]
さらに右表現可能加群 $P_i=e_iD$ に対する
\[
 0\longrightarrow W_i\otimes P_{i+4}
 \longrightarrow R_i\otimes P_{i+3}
 \longrightarrow V_i\otimes P_{i+1}
 \longrightarrow P_i\longrightarrow S_i\longrightarrow0
 \tag{Geo-res}\label{eq:relative-geometric-cubic-resolution}
\]
は各内部次数で局所分裂完全であり、
任意の $S$-基底変更後にも完全である。
$S_i$ は頂点 $i$ に $\mathcal O_S$ を置く加群とする。
この結論は全次数に対して\emph{同じ} $S$ 上で成立する。
左右 Noetherian 性はここでは主張しない。
\end{lemma}
\begin{proof}
各 $D_{ij}$ は有限個の tensor 加群の商なので有限表示である。
$n=j-i$ に関して帰納する。
$n=0,1,2,3$ では階数は $1,2,4,6$ で、
関係の直和因子性から局所自由である。
$n\ge4$ とし、$D_{i+a,j}$（$a=1,3,4$）の
局所自由性と基底変更整合性が既知とする。
三次関係による商の定義から
\[
 W_i\otimes D_{i+4,j}\xrightarrow{d_2}
 R_i\otimes D_{i+3,j}\xrightarrow{d_1}
 V_i\otimes D_{i+1,j}\longrightarrow D_{ij}\longrightarrow0
 \tag{Geo-$n$}\label{eq:relative-geometric-cubic-degree}
\]
は複体で、右端は実際の余核表示である。
全幾何ファイバーでは体上の cubic AS 正則代数の
標準解消が完全である
\cite[arXiv v3, Theorem 5.5.7(1)]{VdB2011}。
左三項は帰納法で局所自由なので、
幾何点で $d_2$ が単射となる最大 minor を可逆化すれば
$d_2$ は直和因子への単射である。
その余核から右の局所自由層への誘導写像も
全幾何点で単射であり、同様に直和因子への単射となる。
従って複体は局所分裂完全で、
最後の余核 $D_{ij}$ も有限局所自由である。
幾何ファイバーの Hilbert 関数から階数公式が従う。
全 $n$ についてこの帰納法を行えば、同じ $S$ 上で
全内部次数の分裂完全性が得られる。
局所分裂完全列は任意の tensor 底変更で完全だから、
(\ref{eq:relative-geometric-cubic-resolution}) も
任意底変更と可換する。
\end{proof}

\begin{proposition}[可換点の有限型開近傍における全軌道 Hom]
\label{prop:quadric-uniform-cubic-hom-open}
命題~\ref{prop:quadric-geometric-quintuple-open} の
$U_{\mathrm{geo}}$ を可換点 $u$ のまわりで
\emph{一度だけ} Zariski 開に縮めると、
全ての $i\le j$ について、相対 perfect 圏で
\[
 D^{\mathrm{pres}}_{ij}[0]
 \xrightarrow{\ \sim\ }
 \mathbf R\operatorname{Hom}(L_i,L_j)
 \tag{Geo-Hom}\label{eq:quadric-uniform-geometric-hom}
\]
が成立する。左辺は
$D^{\mathrm{pres}}=T_{\mathbb Z}(V_i)/(R_i)$ の成分である。
従って両辺は任意の基底変更と可換し、
全次数で階数 $h_{j-i}$ の局所自由層となる。
とくに、$\mathsf s$ の二項移動から作る偶数環
\[
 B=\bigoplus_{n\ge0}
   \operatorname{Hom}(L_0,\mathsf s^nL_0)
\]
は全次数で局所自由かつ任意底変更と整合し、
\[
 \operatorname{rk}B_n=(n+1)^2,\qquad
 B=T(B_1)/(J),\quad
 \operatorname{rk}B_1=4,\quad\operatorname{rk}J=7.
\]
この結果は左右 Noetherian 性、相対 tails heart、
AZ 豊富性、相対 AS4 被覆、内在的 $\mathfrak H_2$
の効果性を含まない。
\end{proposition}
\begin{proof}
補題~\ref{lem:relative-geometric-cubic-flatness} により
$D^{\mathrm{pres}}$ の全成分とその標準解消は
同じ $U_{\mathrm{geo}}$ 上で局所分裂完全である。

四項の相対 Euler 複体
\[
 \mathcal E_i^\bullet:\quad
 L_i\longrightarrow V_i^\vee\otimes L_{i+1}
 \longrightarrow R_i^\vee\otimes L_{i+3}
 \longrightarrow W_i^\vee\otimes L_{i+4}
 \tag{Geo-Euler}\label{eq:quadric-relative-geometric-euler}
\]
を作る。矢は合成による coevaluation と
$R_i,W_i$ の包含から定まり、関係とその重なりにより
連続する合成が零になる。四項窓の相対 Hom は
次数零に集中しているため、この有限複体を
perfect enhancement 上で coherent に実現できる。
可換点では二つの parity の各複体は
$\mathbb P^1\times\mathbb P^1$ の二つの Euler 完全列を
継いだものであり、totalization は零である。
perfect cone の零性は開条件なので、
$i=0,1$ の二つの複体が完全になるよう
$u$ の一つの開近傍へ縮める。
$\mathsf s$ の二項移動により、同じ開上で
全ての $i$ の Euler 複体が完全である。

(\ref{eq:quadric-uniform-geometric-hom}) を
$n=j-i$ に関して帰納する。
$n\le3$ は四項窓と積の定義から成立する。
$n\ge4$ では (\ref{eq:quadric-relative-geometric-euler}) に
$\mathbf R\operatorname{Hom}(-,L_j)$ を適用すると、
$\mathbf R\operatorname{Hom}(L_i,L_j)$ は
\[
 W_i\otimes\mathbf R\operatorname{Hom}(L_{i+4},L_j)
 \longrightarrow
 R_i\otimes\mathbf R\operatorname{Hom}(L_{i+3},L_j)
 \longrightarrow
 V_i\otimes\mathbf R\operatorname{Hom}(L_{i+1},L_j)
\]
の導来余核となる。
帰納法により右辺はすべてより短い距離の
$D^{\mathrm{pres}}$ の成分であり、
その複体は (\ref{eq:relative-geometric-cubic-degree})
と同じである。補題による局所分裂完全性から
その導来余核は $D^{\mathrm{pres}}_{ij}[0]$ で、
合成射からの自然な比較と一致する。
これで全 $n$ に対して主張が従う。

偶数環の次数 $n$ は二項周期で $D_{0,2n}$ に一致する。
三次表示の二次 Veronese の相対的な二次提示に
補題~\ref{lem:relative-quadratic-veronese} を適用すれば、
$B$ は二次関係で表示される。二次関係の核の階数は
$4^2-9=7$ であり、局所直和因子である。
全ての比較は有限局所自由層の分裂表示から得られるため、
任意底変更と整合する。
\end{proof}


\begin{corollary}[有限型開近傍の偶数環の相対 Koszul 性]
\label{cor:quadric-uniform-even-koszul}
命題~\ref{prop:quadric-uniform-cubic-hom-open} の開近傍上で、
偶数環 $B$ は相対的に二次 Koszul である。
その左右の Koszul 複体は各内部次数で局所分裂完全で、
任意の基底変更後も完全である。
本系だけで $B$ の Noetherian 性は結論しない。
\end{corollary}
\begin{proof}
幾何点 $\bar v$ を取る。
その cubic $\mathbb Z$-代数は相対的標準解消
(\ref{eq:relative-geometric-cubic-resolution}) の
体上への制限を持ち、二項周期がある。
奇数頂点の標準解消を偶数頂点へ制限すると、
単純頂点加群が消え、奇数頂点の制限加群の
線形自由分解を繰り返し継いで得られる。
偶数頂点の標準解消を同様に制限し、
得られた奇数制限加群の線形分解を代入すると、
$B_{\bar v}$ の次数零剰余体の線形分解ができる。
これは命題~\cite[\texttt{prop:weighted-even-koszul}]{QuadricIntrinsic2026} の証明の
$d_i=0$ の場合と同じ構成であり、
そこで Noetherian 性を使うのは Noetherian 性の結論の
証明だけである。
従って全ての幾何ファイバーの $B_{\bar v}$ は Koszul である。

命題~\ref{prop:quadric-uniform-cubic-hom-open} により
$B$ は二次表示を持ち、$B_3$ は局所自由である。
Noetherian な有限型開上で
補題~\ref{lem:relative-koszul-degree-three} を適用する。
全次数の局所自由性、左右 Koszul 複体の
局所分裂完全性と任意底変更への整合性が従う。
\end{proof}


\begin{proposition}[可換点の有限型近傍での相対 AS4 被覆]
\label{prop:quadric-zariski-homological-as-cover}
命題~\ref{prop:quadric-uniform-cubic-hom-open} の
可換点 $u\in U(k)$ を含む有限型 Zariski 開部分
$U^\circ\subset U$ が存在し、同じ開の軌道環
\[
 B=\bigoplus_{n\ge0}\operatorname{Hom}
       (L_0,\mathsf s^nL_0)
       =T_{\mathcal O_{U^\circ}}(B_1)/(K),
 \qquad \operatorname{rk}K=7
\]
に対し、階数六の直和因子 $H\subset K$ を選べる。
$A=T(B_1)/(H)$、$G=K/H$ とすると、次が成立する。
\begin{enumerate}
\item $A_n$ は全次数で有限局所自由で
  $\operatorname{rk}A_n=\binom{n+3}{3}$。
  各幾何ファイバーは型 $(1,4,6,4,1)$ の quadratic
  Koszul AS 正則環である。左右 Koszul 解消および
  相対双加群 Koszul 解消は各内部次数で局所分裂完全で、
  任意の底変更後も完全である。
\item $G$ の局所生成元 $g$ は $A$ の次数二の正則正規元であり、
  $A/(g)\simeq B$ は全次数の基底変更と整合する。
  $g$ の左右乗法は全次数で直和因子への単射である。
\item $A^!$ は最高次数四の相対 Frobenius 二次環であり、
  その階数列は $(1,4,6,4,1)$。
  相対 AS Gorenstein の左右双対性を持つ。
\end{enumerate}
これらの homological 条件に加え、任意の Noetherian 底変更
$T\to U^\circ$ で $A_T,B_T$ は左右 Noetherian である。
従って $A\to B$ は定義~\ref{def:relative-h2} に必要な相対 AS4 被覆となる。
相対 tails heart、AZ 豊富性、canonical compact 圏との比較は
第\ref{sec:relative-cubic-tails}節で別に証明する。
\end{proposition}
\begin{proof}
命題~\ref{prop:quadric-uniform-cubic-hom-open} により、同じ有限型開上で
全次数の geometric cubic 表示、標準解消、二周期、偶数軌道環との同定がある。
右軌道積の規約には (\ref{eq:cubic-tails-right-convention}) を用い、
$C_{ij}=D_{-j,-i}$、$\beta_C=\mathsf s^{-1}$ の偶数環へ
定理~\ref{thm:universal-cubic-hull} を適用する。
同定理の線束値写像
\[
 \Phi:S_0\longrightarrow W_0,\qquad
 S_0=R_0\otimes V_3+V_0\otimes R_1
\]
の核が求める $H$ である。$H$ は $W_0$ の局所生成元に依存せず、
任意の底変更と整合する。必要なら $W_0$ を自明化する affine 近傍へ制限する。
全次数の平坦性、正則正規二次切断、Koszul Frobenius 性および
普遍 Noetherian 性は同定理から従う。
この構成には、Artin 厚みから等式を延長する操作は必要ない。
\end{proof}


\begin{corollary}[幾何的 cubic 点の局所基底での Noetherian AS4 被覆]
\label{cor:quadric-point-local-noetherian-as4}
命題~\ref{prop:quadric-zariski-homological-as-cover} の
幾何的 cubic 開 $U^\circ$ の任意の点 $v$ と、
閉点を $v$ へ送る任意の Noetherian 局所
$\mathcal O_{U^\circ,v}$-代数 $R$ を取る。
このとき $B_R$ と homological 被覆 $A_R$ は左右 Noetherian である。
従って $A_R\twoheadrightarrow B_R$ は
型 $(1,4,6,4,1)$ の相対 Noetherian quadratic AS4 被覆であり、
正則正規二次切断と相対 Koszul/AS 双対性は
全ての基底変更と整合する。
特に $R=\mathcal O_{U^\circ,v}$ 自体とその完備化に適用できる。
有限型の可換 $R$-代数への変更後も左右 Noetherian である。
従って $\operatorname{qgr}B_R$ は Noetherian Abel 圏であり、
$\pi B_R$ は構造対象、次数移動は exact 自己同値である。
単一の有限型開全体の Noetherian 性は、この局所判定からは従わない。
canonical compact 比較と Noetherian な基底上の全相対条件は、
第\ref{sec:relative-cubic-tails}節で別に証明する。
\end{corollary}
\begin{proof}
$B_R$ は全次数で有限自由である。
任意の幾何的 cubic ファイバーは
Van den Bergh 型三次元 AS 正則 $\mathbb Z$-代数の
偶数 Veronese である。
体上の三次元 cubic AS 正則 $\mathbb Z$-代数の
左右 Noetherian 性\cite[Corollary 5.5.9(1)]{VdB2011} と、偶数頂点への制限で
Noetherian 性が保たれる事実
（命題~\cite[\texttt{prop:weighted-even-koszul}]{QuadricIntrinsic2026} の証明の
偶数イデアルの有限生成性の部分）により、
幾何的特殊ファイバー $B_{\bar v}$ は左右 Noetherian。
体拡大の忠実平坦性により $B_R$ の特殊ファイバーも
左右 Noetherian である。
そこで補題~\cite[\texttt{lem:local-commutative-fibre-noetherian}]{QuadricIntrinsic2026}
の強化した結論を適用する。
$A_R$ は正則正規二次元 $g$ で $A_R/(g)=B_R$
となるから、定理~\ref{thm:relative-finite-as-cover} の
正規元による Noetherian 性の議論を適用して
$A_R$ も左右 Noetherian となる。
他の相対 AS4 条件は既に homological 被覆から従う。
全次数の局所自由性・幾何的 cubic 条件は
任意の局所基底変更で保たれる。
有限型の可換 $R$-代数への変更後の Noetherian 性は
補題~\ref{lem:local-finite-type-strong-noetherian}
を $B_R,A_R$ へ適用して得る。
\end{proof}




\begin{corollary}[幾何的 cubic 点での全相対 \texorpdfstring{$\chi$}{chi}]
\label{cor:quadric-local-chi}
命題~\ref{prop:quadric-zariski-homological-as-cover} の
幾何的 cubic 開 $U^\circ$ の任意の点 $v$ と、
$\operatorname{Spec}R\to U^\circ$ を取る。
ここで $(R,\mathfrak m,\kappa)$ は有限 Krull 次元の
Noetherian 局所環で、その閉点は $v$ に写るとする。
このとき切断環 $B_R$ は左右の相対 $\chi_q$ を
全ての $q\ge0$ について満たす。
すなわち任意の有限生成次数付き右・左 $B_R$-加群
$M$ に対して
\[
 \underline{\operatorname{Ext}}_{B_R}^j(R,M)
 \quad(j\ge0)
\]
は全内部次数を合わせて有限生成 $R$-加群となる。
この結論を非局所基底上での一様な
$\chi$ や AZ (A2) の証明と同一視しない。
\end{corollary}
\begin{proof}
系~\ref{cor:quadric-point-local-noetherian-as4} により
$B_R$ は左右 Noetherian である。
系~\ref{cor:quadric-uniform-even-koszul} により
自明加群 $R$ の左右相対 Koszul 解消は
各項が有限斉次自由であり、任意の底変更後にも完全となる。

$\bar\kappa$ を $\kappa$ の代数閉包とする。
幾何的ファイバー $B_{\bar\kappa}$ は
命題~\ref{prop:quadric-zariski-homological-as-cover} により
四生成の quadratic AS 正則環
$A_{\bar\kappa}$ の正則正規二次切断による商である。
また $B_{\bar\kappa}$ は geometric cubic
$\mathbb Z$-代数の偶数環であり、体上の Noetherian 性は
系~\ref{cor:quadric-point-local-noetherian-as4} の証明で
確認した。正則正規元による拡張から
$A_{\bar\kappa}$ も左右 Noetherian である。
従って Smith--Van den Bergh
\cite[\S2.8]{SVdB2013}
の既知の $\chi$ 性の結論が適用され、
$B_{\bar\kappa}$ は左右 $\chi_q$ を全ての $q$ で満たす。
ここで引用するのは体上の既知結果であり、
相対基底への拡張は引用に含まれない。

各有限生成 $B_\kappa$-加群の Ext は
有限斉次自由 Koszul 解消により計算でき、
$\kappa\to\bar\kappa$ の平坦体拡大と可換する。
ゆえに $B_{\bar\kappa}$ の $\chi_q$ から
$B_\kappa$ の左右 $\chi_q$ が従う。
補題~\ref{lem:local-chi-specialization} を
$R$ 上で任意の $q$ に適用すればよい。
\end{proof}

\begin{corollary}[幾何的 cubic 点の Noetherian 局所基底での AZ 豊富性]
\label{cor:quadric-local-AZ-ampleness}
前系の任意の Noetherian 局所基底 $R$ では
\[
 (\pi B_R,(1))
\]
は $\operatorname{qgr}B_R$ 上で
Artin--Zhang ample、すなわち (A1), (A2) を満たす。
有限型基底変更と canonical compact・Serre の比較については、
定理~\ref{thm:quadric-local-finite-type-abelian-system} と
第\ref{sec:relative-cubic-tails}節で確認する。
$\mathfrak H_2$ 全体の代数性は未証明である。
\end{corollary}
\begin{proof}
系~\ref{cor:quadric-local-chi} から右相対 $\chi_1$ があり、
各 $B_{R,n}$ は有限自由である。
補題~\ref{lem:relative-chi1-implies-AZ-A2} により (A2) を得る。
(A1) は系~\ref{cor:quadric-local-tails-basechange-generation}
を局所環 $\mathcal O_{U^\circ,v}$ とその基底変更 $R$
に適用して得る。ただし $R$ が有限型でないときも、
証明で用いた $B_R$ の左右 Noetherian 性と
高次数 truncation の有限生成性が成立するので
同じ (A1) の直接証明が使える。
\end{proof}

\begin{theorem}[局所基底上の全有限型 tails と Hom]
\label{thm:quadric-local-finite-type-abelian-system}
系~\ref{cor:quadric-local-chi} の有限 Krull 次元 Noetherian 局所環
$R$ を固定する。任意の可換有限型 $R$-代数 $R'$ に対し、
$B'=B_R\otimes_RR'$ は左右 Noetherian かつ左右相対 $\chi_q$
（全ての $q\ge0$）を満たし、
$(\pi B',(1))$ は $\operatorname{qgr}B'$ 上 AZ ample である。
さらに
\[
 \operatorname{Spec}R'\longmapsto\operatorname{qgr}B'
 \tag{Local-tails}\label{eq:local-finite-type-tails-system}
\]
は Noetherian Abel 圏の系であり、module-object 基底変更、
対象・射・完全列の fppf 降下、全ての対象間の Hom の $R'$-有限性と
平坦な有限型基底変更との整合性を持つ。
構造対象 $\pi B'$ は基底上平坦で、$A_R\otimes_RR'\twoheadrightarrow B'$
は同じ次数移動に対する相対 Noetherian AS4 被覆である。

$R'$ が局所環であることも $R$ 上平坦であることも要求しない。
canonical compact 圏と Serre 平方根の比較、ならびに有限 Krull 次元の
仮定を除いた昇格は、定理~\ref{thm:quadric-noetherian-relative-promotion}
で証明する。単一の有限型開全体の Noetherian 性は
定理~\ref{thm:relative-geometric-cubic-noetherian} によって別に得られる。
\end{theorem}
\begin{proof}
$B_R$ の Noetherian 性は
系~\ref{cor:quadric-point-local-noetherian-as4}、
全相対 $\chi_q$ は系~\ref{cor:quadric-local-chi} による。
相対 Koszul 解消の各項は有限斉次自由で任意の基底変更後も完全だから、
補題~\ref{lem:chi-finite-type-basechange} が全ての $q$ に適用できる。
$B',A_R\otimes_RR'$ の Noetherian 性は
補題~\ref{lem:local-finite-type-strong-noetherian} による。
従って高次数 truncation の有限生成性から AZ (A1)、
補題~\ref{lem:relative-chi1-implies-AZ-A2} から AZ (A2) を得る。

各 $B'_n$ は有限 projective $R'$-加群なので、
補題~\ref{lem:chi1-qgr-hom-finite-flat} により任意の二対象間の
Hom は有限で平坦基底変更と可換する。
module-object 比較は補題~\ref{lem:qgr-module-object-basechange}、
fppf 降下と構造対象の平坦性は
系~\ref{cor:quadric-local-tails-fppf-descent} の証明を
この $R$ に適用して得る。完全列の降下は、
faithfully flat な拡大が完全かつ保守的であることによる。
AS4 被覆の残りの条件は既存の homological 被覆の
普遍的完全性から全ての基底変更後にも成立する。
\end{proof}

\begin{corollary}[有限型族の全 stalk における局所 AZ 豊富性]
\label{cor:quadric-local-AZ-on-all-stalks}
$T$ を $U^\circ$ 上局所有限型な Noetherian スキームとし、
$\mathcal B_T$ を上の切断環の引戻しとする。
任意の点 $t\in T$ に対して
\[
 B_t=\mathcal B_T\otimes_{\mathcal O_T}
                 \mathcal O_{T,t}
\]
は左右 Noetherian、左右の相対 $\chi_q$ を全 $q$ で満たし、
$(\pi B_t,(1))$ は $\operatorname{qgr}B_t$ 上 AZ ample である。
特に $T$ の非被約点及び非閉点においても成立する。

これは \emph{各 stalk の}結論であり、$T$ 上で全
内部次数に共通する $\chi$ 消滅の下限や、
単一の affine chart 上の切断環の Noetherian 性、
relative AZ (A2) の一様性、
canonical heart の認識を結論しない。
\end{corollary}
\begin{proof}
任意の $t\in T$ に対して
$\mathcal O_{T,t}$ は有限 Krull 次元の
Noetherian 局所環であり、その閉点の像は
$U^\circ$ の一点である。
系~\ref{cor:quadric-point-local-noetherian-as4} と
直前の二つの系を
$R=\mathcal O_{T,t}$ に適用すればよい。
\end{proof}



\begin{corollary}[任意基底変更での相対 quadric の完全切断環]
\label{cor:quadric-all-basechange-saturated-sections}
命題~\ref{prop:quadric-zariski-homological-as-cover} の
幾何的 cubic 開 $U^\circ$ の任意の affine 底変更
$T=\operatorname{Spec}R'\to U^\circ$ に対して、
\[
 \underline{\operatorname{Ext}}_{B_T}^j(R',B_T)=0
 \qquad(j=0,1,2)
\]
が全内部次数で成立する。さらに $\mathcal O_T=\pi B_T$ に対して
\[
 \operatorname{Hom}_{\operatorname{QGr}B_T}
       (\mathcal O_T,\mathcal O_T(n))
 \simeq (B_T)_n\qquad(n\in\mathbb Z)
 \tag{Quadric-H0}\label{eq:quadric-relative-full-sections}
\]
は自然な同型であり、任意の底変更と可換する。
$n<0$ では $(B_T)_n=0$ とする。
特に構造対象の軌道切断環が全次数で元の $B_T$ を回復する。
さらに任意の $R'$-加群 $N$ について
\begin{align*}
 \operatorname{Hom}_{\operatorname{QGr}B_T}
   (\mathcal O_T,N\otimes_{R'}\mathcal O_T(n))
   &\simeq N\otimes_{R'}(B_T)_n,\\
 \operatorname{Ext}^1_{\operatorname{QGr}B_T}
   (\mathcal O_T,N\otimes_{R'}\mathcal O_T(n))&=0
 \qquad(n\in\mathbb Z).
\end{align*}
非被約・非平坦な底変更も含むが、
一般 coherent 対象の Hom 有限性には Noetherian 性を別に用いる。
Noetherian 性なしの canonical compact 比較は
命題~\ref{prop:quadric-canonical-cubic-tails} で証明する。
\end{corollary}
\begin{proof}
各幾何ファイバーでは
命題~\ref{prop:quadric-zariski-homological-as-cover}
によって $B_{\bar t}$ は四次元 quadratic AS 正則環
$A_{\bar t}$ の正則正規二次切断商である。
体上の Gorenstein 降下
\cite[\S2.8]{SVdB2013}
より $B_{\bar t}$ は深さ三で、
自明加群からの $\operatorname{Ext}^j$ は
$j=0,1,2$ で零となる。
系~\ref{cor:quadric-uniform-even-koszul} による
左右相対 Koszul 解消は全内部次数で局所分裂完全で
任意の底変更と可換する。
従って補題~\ref{lem:relative-depth-three-saturation}
と系~\ref{cor:relative-depth-three-coefficients}
を右・左両側に適用すれば結論を得る。
\end{proof}

\begin{remark}[有限型相対実現への依存関係]
この節は、標準 Morita 根、四項窓、全次数の相対 cubic 解消と Hom、
相対 Koszul 性、Zariski 局所 homological AS4 被覆を構成する。
局所 Noetherian 判定と構造対象の切断計算も得られるが、
これらを一点ごとの結果だけから有限型開全体へ移してはならない。

第\ref{sec:relative-cubic-tails}節は普遍 cubic 解消から
canonical な Grothendieck 導来比較と全 helix の根の認識を証明し、
左右 Noetherian 性から全相対条件を導く。
第\ref{sec:relative-cubic-noetherian}節は相対曲線商・linear 部分・形式的貼合せを用いて
任意の Noetherian 基底変更後の左右 Noetherian 性を証明する。
これらを合わせるのが定理~\ref{thm:quadric-finite-type-effectivity} である。
Artin 基底上の heart exactness を冪零でない基底へ外挿する必要はない。
この有限型相対族の存在と $\mathfrak H_2$ 全体の代数性は区別する。
\end{remark}

% !TeX root = ../nc-quadric-surface-moduli.tex
% 2026-11-21 / v1.34
\section{可換 quadric の全 Artin 変形と相対 AS4 被覆}
\label{sec:quadric-artin-hull}

以下では $(A,\mathfrak m,k)$ を標数零の体上の有限次元可換局所 Artin 環とする。
本節は一般の相対 canonical hull 定理~\ref{thm:universal-cubic-hull} を
可換 quadric の Artin 変形へ適用する。
主眼は、任意の pointed Abel 変形から、同じ構造対象を持つ exact ample な
Serre 平方根と相対 AS4 被覆を回復することである。
二周期の $\mathbb Z$-代数を用い、cubic 代数の一周期性は仮定しない。

\paragraph{入力を有限自由な tensor データとして固定する。}
$D$ を可換 quadric の cubic $\mathbb Z$-代数の $A$-平坦変形とし、
$V_i=D_{i,i+1}$、$R_i\subset V_i\otimes V_{i+1}\otimes V_{i+2}$ を
それぞれ階数二の一次生成元・三次関係の自由加群とする。
三次の標準解消は相対完全であり、
\[
 W_i=(R_i\otimes V_{i+3})\cap(V_i\otimes R_{i+1})
\]
は階数一で、任意の生成元 $w_i$ は両方の表示で非退化とする。
この非退化性は閉ファイバーでの非退化性と Nakayama から従う。
次数二の周期 $\beta:D\to D(2)$ が与えられたとし、
その偶数部分の graded 化を $B$ とする。
特に $B_1$ は階数四、$B_n$ は階数 $(n+1)^2$ である。

$B$ は二次環で相対 Koszul、二次双対 $E=B^!$ の各次数は有限自由で、
その階数は $1,4,7,8,8,\ldots$ と仮定する。
これらは後述の helix からの入力では既存の相対 Koszul 補題により成立する。
双対の積は tensor 部分加群の包含を双対化する規約で書く。
全ての $i$ に対して
\[
 S_i=R_i\otimes V_{i+3}+V_i\otimes R_{i+1}
 \subset V_i\otimes V_{i+1}\otimes V_{i+2}\otimes V_{i+3}
\]
と置く。これは階数七であり、偶数 $i$ では $\check B$ の二次関係加群である。

\begin{lemma}[平坦な cubic 変形の syzygy と四周期]
\label{lem:relative-cubic-tensors-artin}
$D$ の各成分が有限自由で、閉ファイバーが可換 quadric の cubic
$\mathbb Z$-代数であるとする。
このとき上記の三次提示、$R_i,W_i$ の自由性と底変更、非退化性、
三次の標準解消の相対完全性が成立する。
また生成元 $w_i\in W_i$ を選ぶと、$w_i$ の巡回から
$\theta_i:V_i\simeq V_{i+4}$ が定まり、$D\simeq D(4)$ を誘導する。
\end{lemma}
\begin{proof}
閉ファイバーで次数一生成なので、各次数の乗法の全射性は Nakayama で持ち上がる。
次数二の乗法は階数四の自由加群間で還元が同型だから同型である。
次数三では、階数八の自由加群から階数六の自由加群への全射の核が
$R_i$ であり、これは階数二で還元と可換する。
三次関係で割った tensor $\mathbb Z$-代数を $C$ とすると $C\twoheadrightarrow D$ がある。
各成分の終域 $D_{ij}$ は自由なので、核の還元は還元した射の核と一致する。
$C\otimes_Ak\to D\otimes_Ak$ は同型であるから、有限な核は Nakayama により零となる。
従って $D$ はこの三次提示そのものである。

次数四の関係加群は、階数十六から階数九への自由加群間の全射の核であり、階数七である。
三次提示から
\[
 (R_i\otimes V_{i+3})\oplus(V_i\otimes R_{i+1})
 \longrightarrow\ker(T_{i,i+4}\to D_{i,i+4})
\]
は全射なので、その核は階数一で還元と可換する。
二つの包含の差をこの射とすれば、核は交わり $W_i$ と同定される。
閉ファイバーで $w_i$ の二つの $2\times2$ pairing 行列は可逆なので、
$A$ 上でも行列式は単元であり、非退化である。

これらの包含から得る標準 cubic 複体
\[
 0\to W_i\otimes P_{i+4}\to R_i\otimes P_{i+3}
 \to V_i\otimes P_{i+1}\to P_i\to S_i^{\mathrm{simple}}\to0
\]
は複体になる。各内部次数で有限自由 $A$-加群の有界複体であり、
閉ファイバーで完全である。
単元 minor による消去、または局所環上の最小複体を使えば、
元の複体は $A$-加群として分裂完全である。従って任意の底変更後にも完全である。

四周期は非退化 tensor から直接得られる。
$w_i=\sum_a u_a\otimes r_a$ を $V_i\otimes R_{i+1}$ 内で書き、
同じ $R_{i+1}$ の基底を使って $w_{i+1}=\sum_a r_a\otimes u'_a$ と書く。
$\theta_i(u_a)=u'_a$ と定めれば、これは基底によらない同型である。
従って $w_{i+1}$ は $w_i$ の最初の因子を最後へ巡回し、
その因子に $\theta_i$ を作用させた tensor である。
四回反復すると $\bigotimes_{a=0}^3\theta_{i+a}$ は $W_i$ を $W_{i+4}$ へ移す。
最後の因子による縮約は $R_i$ を回復するから、
$\theta_i\otimes\theta_{i+1}\otimes\theta_{i+2}$ は $R_i$ を $R_{i+4}$ へ移す。
従って $\theta$ は三次提示を保ち、$D\simeq D(4)$ を与える。
ここでは $\theta=\beta^2$ は仮定も結論もしていない。
\end{proof}

\begin{lemma}[canonical functional の相対的な正規性]
\label{lem:relative-canonical-hull-functional}
$w_i$ が定める二つの coevaluation に対する評価写像を
$R_i\otimes V_{i+3}$ と $V_i\otimes R_{i+1}$ 上に取る。
両者は $W_i$ 上で一致し、全射
\[
 \phi_i:S_i\longrightarrow A,\qquad \phi_i(w_i)=2
\]
を定める。その核 $H_i$ は階数六の直和因子である。
$\phi_i$ の張る直線は $w_i$ の単元倍によらず、$\beta$ と整合する。
さらに $E$ の二次元に対応するこれらの元は、正規化列をなす。
従って $\beta$ に沿って graded 化すれば、$E_2$ に正規元
$\phi$ が得られる。
\end{lemma}
\begin{proof}
定理~\ref{thm:universal-cubic-hull} の線束値写像
$\Phi_i:S_i\to W_i$ を $W_i=Aw_i$ で自明化したものが $\phi_i$ である。
二つの逆 coevaluation は $W_i$ 上でともに $2\,\mathrm{id}$ となるので貼り合う。
正規性は同定理の証明にある、次数三の二つの縮約の tensor 恒等式による。
そこでは四周期 $\theta$ と二周期 $\beta$ を区別しており、
$\theta=\beta^2$ を仮定していない。
\end{proof}

\begin{proposition}[正規 functional からの相対 AS4 被覆]
\label{prop:relative-artin-hull-from-functional}
前補題の入力を満たすとき、$B$ は全次数平坦な二次 AS4 被覆
$\widetilde B\twoheadrightarrow B$ を持ち、その核は二次の左右正則正規元で生成される。
\end{proposition}
\begin{proof}
補題~\ref{lem:relative-cubic-tensors-artin} により、
定理~\ref{thm:universal-cubic-hull} の全入力が成立する。
同定理を適用し、可逆加群 $G$ を局所 Artin 基底上で自明化すればよい。
全次数の平坦性、左右正則正規性、相対 Koszul Frobenius 性、
任意の底変更との整合、および Noetherian 性も同時に従う。
\end{proof}

\begin{remark}[可換ファイバーでの正規化]
標準二周期を持つ Segre 環では、$\phi_0$ は六つの交換関係を零に送り、
$q=x_{00}x_{11}-x_{01}x_{10}$ を $1$ に送る。
これは \cite[(5.3)]{VdB2011} の
\[
 w=u_0v_0u_1v_1-u_1v_0u_0v_1-u_0v_1u_1v_0+u_1v_1u_0v_0
\]
を逆 coevaluation の定義へ代入して確かめられる。
従って $E_k/(\phi_0)$ は四変数の外積代数であり、
その双対 hull は四変数多項式環となる。
\end{remark}

\begin{theorem}[可換 quadric の全 Artin 変形の相対実現]
\label{thm:quadric-all-artin-h2}
$Q=\mathbb P^1_k\times\mathbb P^1_k$ とし、
$(\mathcal C_A,\mathcal O_A)$ を $(\operatorname{coh}Q,\mathcal O_Q)$ の
任意の flat pointed Abel 変形とする。
この変形には $s_0=(-)\otimes\mathcal O_Q(1,1)$ を還元とする
exact ample な相対 Serre 平方根 $s_A$ が存在し、その全軌道環 $B_A$ は
\[
 B_A\simeq\widetilde B_A/(g),\qquad
 \operatorname{rank}_A(B_A)_n=(n+1)^2
\]
と表示できる。ここで $\widetilde B_A$ は型 $(1,4,6,4,1)$ の相対 quadratic
AS 正則環であり、$g$ は二次の左右正則正規元である。
従って $(\mathcal C_A,\mathcal O_A)$ は、任意の有限型 $A$-代数への基底変更を
備えた定義~\ref{def:relative-h2} の相対 $\mathfrak H_2$ 族へ昇格する。
根と被覆は存在条件であり、新しい標識として加えない。
\end{theorem}
\begin{proof}
系~\ref{cor:quadric-exact-root-all-artin} による $s_A$ と、
$\mathcal O_Q(0,1)$ の flat な持上げ $M_A$ を取る。
構造対象と $M_A$ の存在・同型類の一意性は
定理~\ref{thm:artin-foundation-lift} の geometric helix の持上げによる。
\[
 L_{2n,A}=s_A^n\mathcal O_A,\qquad L_{2n+1,A}=s_A^nM_A
\]
と置く。各 $L_{i,A}$ は対応する可換線束の一意な flat 持上げであり、
全列は De Deken--Lowen による持ち上げた geometric helix と一致する。
従って $D_{ij}=\operatorname{Hom}(L_{-j,A},L_{-i,A})$ の各成分は有限自由であり、
補題~\ref{lem:relative-cubic-tensors-artin} により相対 cubic 提示と解消、
$V_i,R_i,W_i$ の自由性・非退化性を得る。
この cubic 変形との比較は \cite[Theorem 8.1.2, \S8.3]{VdB2011} とも整合する。

負添字の Hom 規約により $s_A^{-1}$ が
$D_{ij}\to D_{i+2,j+2}$ を与え、これを $\beta$ とする。
偶数部分の graded 化の $n$ 次成分
$D_{0,2n}=\operatorname{Hom}(s_A^{-n}\mathcal O_A,\mathcal O_A)$ に
$s_A^n$ を作用させると、$s_A$ の全軌道環 $(B_A)_n$ と同定される。
その積は $a*b=s_A^{\deg b}(a)\circ b$ という本稿の右規約になる。

この全軌道環は全次数自由であり、閉ファイバーでの次数一生成と二次提示は
次数ごとの Nakayama で持ち上がる。
三次も自由なので、補題~\ref{lem:relative-koszul-degree-three} により
$B_A$ は相対 Koszul、$E=B_A^!$ の全次数は自由である。
閉ファイバーの Hilbert 級数から階数は $1,4,7,8,8,\ldots$ となる。
従って補題~\ref{lem:relative-canonical-hull-functional} と
命題~\ref{prop:relative-artin-hull-from-functional} の全入力が成立する。
これらによって相対 AS4 被覆と正則正規な切断を得る。
最後に命題~\ref{prop:artin-witness-relative-promotion} を適用し、
全有限型基底変更上の相対 $\mathfrak H_2$ 族を得る。
\end{proof}

\begin{corollary}[可換点での形式変形問題]
\label{cor:quadric-formal-deformation-equivalence}
$Q$ の相対 $\mathfrak H_2$ 変形から flat pointed Abel 変形への忘却は、
局所 Artin $k$-代数上の変形 groupoid の同値を与える。
この変形問題は障害を持たず、その接空間は
\[
 \operatorname{HH}^2(Q)=H^0(Q,\mathcal O_Q(2,2)),\qquad \dim=9
\]
である。無限小自己同型の空間は $\operatorname{HH}^1(Q)=H^0(Q,T_Q)$、次元六である。
\end{corollary}
\begin{proof}
本質的全射性は前定理による。
基底 $A$ 上の pointed exact 同値は Ind 圏へ延長し、有限型 $A$-代数 $A'$ の
作用を移すことにより module-object 圏上の同値を与える。
相対族の基底変更同値により、その延長は自然同型を除いて一意である。
根と被覆は標識ではないので追加の射の条件を生まない。
従って忘却は全忠実であり、底変更に自然な同値である。

Abel 変形の接空間・無限小自己同型・障害空間はそれぞれ
Hochschild 次数二・一・三である
\cite[Theorem 3.1, Remark 3.5]{LowenVanDenBergh2005}。
構造対象の持上げには $\operatorname{Ext}^{1,2}(\mathcal O_Q,\mathcal O_Q)=0$ を用いる。
HKR と $Q$ 上のコホモロジーから $\operatorname{HH}^3(Q)=0$ である。
具体的には $H^1(\bigwedge^2T_Q)=H^2(T_Q)=0$ であり、
残る全次数三の成分も零である。
これにより任意の小拡大に対して変形を持ち上げられる。
表示した接空間と無限小自己同型も HKR による。
\end{proof}

この定理は可換点における全 Artin 基底の認識を与える。
一般有限型基底への代数化、一般の非可換点での同様の認識、
内在的 $\mathfrak H_2$ 全体の代数性は別に証明する必要がある。

\subsection{cubic 点での条件付き Artin 認識}

\begin{proposition}[適合 cubic helix と Hochschild 自明作用からの全 Artin 認識]
\label{prop:cubic-artin-hhtrivial-conditional}
標数零の代数閉体 $k$ 上の $(\mathcal C,\mathcal O)\in\mathcal H_2$ について、
次のデータと条件を仮定する。
\begin{enumerate}
\item $E_0=\mathcal O$ を含む ample geometric four-helix $(E_i)_{i\in\mathbb Z}$ があり、
  $D_{ij}=\operatorname{Hom}(E_{-j},E_{-i})$ は標準的な三次三次元 AS 正則
  $\mathbb Z$-代数である。
\item exact な Serre 平方根 $s_0^2\simeq S^{-1}[2]$ があり、
  $(\mathcal O,s_0)$ は Artin--Zhang ample で、
  全ての $i$ について $s_0E_i\simeq E_{i+2}$ である。
  ここでは対象ごとの同型で十分であり、別途 coherent な周期の標識は要求しない。
\item $(s_0)_*$ は $\operatorname{HH}^2(\mathcal C)$ 上恒等に作用する。
\end{enumerate}
このとき、任意の局所 Artin $k$-代数 $A$ 上の flat pointed Abel 変形
$(\mathcal C_A,\mathcal O_A)$ は、$s_0$ の持上げを用いる全軌道環の
相対 quadratic AS4 被覆を持つ。
従って全有限型基底変更を含む相対 $\mathfrak H_2$ 族へ昇格する。
この点での相対 $\mathfrak H_2$ 変形と flat pointed Abel 変形の groupoid は、
局所 Artin 基底上で同値となる。

さらに $(s_0)_*$ が $\operatorname{HH}^1(\mathcal C)$ 上も恒等なら、
この還元を持つ Serre 平方根の持上げは Morita 自己同値の同型類として一意である。
\end{proposition}
\begin{proof}
有限 foundation の持上げと相対 smooth proper 性は
定理~\ref{thm:artin-foundation-lift} による。
小拡大における根の障害 $b$ は $b+(s_0)_*b=0$ を満たすので、
Hochschild 二次の仮定により零となる。
近恒等核の一意な二可除性で平方関係を補正し、全 Artin 基底上の $s_A$ を得る。
補題~\ref{lem:artin-morita-heart-exact} と
補題~\ref{lem:artin-ample-lift} によってこれは exact ample である。

$M_A$ を $E_1$ の flat な持上げとし、
$L_{2n,A}=s_A^n\mathcal O_A$、$L_{2n+1,A}=s_A^nM_A$ と置く。
条件 (2) と例外対象の持上げの一意性により、$L_{i,A}$ は
$E_i$ の geometric helix としての持上げと同型である。
従って $D_{ij,A}=\operatorname{Hom}(L_{-j,A},L_{-i,A})$ の前向き成分は有限自由で、
閉ファイバーは元の cubic $\mathbb Z$-代数となる。
$s_A^{-1}$ が実際の次数二の周期を与える。

補題~\ref{lem:relative-cubic-tensors-artin} の証明は、閉ファイバーが可換であることを使わない。
必要なのは cubic AS 解消と成分の階数 $1,2,4,6,9,\ldots$ である。
同じ Nakayama の議論から、三次提示、階数二の $R_i$、階数一の $W_i$、
非退化な coevaluation、四周期 $\theta$、全内部次数で相対完全な cubic 解消を得る。
偶数部分の graded 化 $B_A$ は $s_A$ の全軌道環であり、
全次数自由、相対二次 Koszul、その双対 $E_A=B_A^!$ の階数列は
$1,4,7,8,8,\ldots$ となる。

正規 functional $\phi_i$ の構成と tensor 恒等式も、そのまま成立する。
とくに次数三の直和分解
\[
 (E_3^*)_i=(R_i\otimes R_{i+3})
       \oplus(V_i\otimes W_{i+1}\otimes V_{i+5})
\]
は可換点専用の恒等式ではない。
体上では右 cubic 解消の次数六の部分が
\[
 W_i\otimes V_{i+4}\otimes V_{i+5}
 \hookrightarrow R_i\otimes D_{i+3,i+6}
\]
を与える。
$D_{i+3,i+6}$ は三重 tensor を $R_{i+3}$ で割ったものなので、
これは $R_iR_{i+3}\cap W_iV_{i+4}V_{i+5}=0$ を意味する。
一方、$W_{i+1}\subset R_{i+1}V_{i+4}$ を用いると、
$R_iR_{i+3}$ と $V_iW_{i+1}V_{i+5}$ の交わりは
$W_iV_{i+4}V_{i+5}$ に入る。従って交わりは零である。
これは \cite[Lemmas 2.3--2.4]{BauwensVanDenBergh1998} と同じ
解消による議論を、右加群の規約で添字付きにしたものである。
二つの四次元部分空間と $\dim E_3=8$ から体上の分解を得て、
自由加群間の Nakayama により $A$ 上へ持ち上げる。
この分解上の二つの縮約は既存の係数計算で一致し、
$\beta$ と整合する正規元 $\phi\in(E_A)_2$ が得られる。

可換点の証明で変更する箇所は、特殊ファイバーの最後の同定だけである。
一般の cubic 点では $E_k/(\phi_0)$ は外積代数とは限らない。
その代わり、\cite[Theorem 6.1, Lemmas 6.4--6.5]{VdB2011} の
canonical hull の二次双対であり、Koszul Frobenius、階数列 $(1,4,6,4,1)$ である。
$\phi_0$ は左右正則正規である。
評価 functional との一致は
\cite[(2.11), Proposition 3.2, Remark 3.3]{BauwensVanDenBergh1998} の
canonical choice を、上の添字付き評価で表したものである。
この体上の入力だけを使い、$\phi$ の乗法を各次数の有限自由加群間で持ち上げる。
従って $P_A=E_A/(\phi)$ は全次数自由となる。
相対 Koszul 補題と Frobenius pairing の持上げから
$\widetilde B_A=P_A^!$ は相対 AS4 となる。

$\widetilde B_A\twoheadrightarrow B_A$ の次数別核は自由であり、
二次生成元 $g$ の左右乗法は、閉ファイバーの正則正規切断の同型を持ち上げる。
よって全次数で同型となり、$g$ は左右正則正規である。
閉ファイバーの canonical hull は noetherian なので、
$\mathfrak m$-進有限フィルターにより $\widetilde B_A$ も左右 noetherian である。
命題~\ref{prop:artin-witness-relative-promotion} を適用すれば相対 $\mathfrak H_2$ 族となる。
忘却の全忠実性は module-object 系の定義、
平方根の追加の一意性は命題~\ref{prop:artin-serre-root-uniqueness} による。
\end{proof}

\begin{remark}[Hochschild 作用条件は対象ごとの周期性から従わない]
\label{rem:cubic-helix-root-action}
この命題は、全ての cubic 点または全ての平方根について
$(s_0)_*=1$ on $\operatorname{HH}^2$ と主張していない。
$s_0^2=S^{-1}[2]$ から得られるのは作用が対合であることまでである。
実際、可換 $Q$ 上で $\varphi(z,w)=(-z,w)$ とすると
$s_0=(-\otimes\mathcal O_Q(1,1))\circ\varphi^*$ は Serre 平方根であり、
標準 helix の各線束を対象ごとに二項先へ送る。
しかし $\partial_z\wedge\partial_w$ を負に送るので、
$\operatorname{HH}^2$ 上恒等ではない。
従って cubic helix との対象ごとの適合性から作用条件を除くことはできない。
一般 cubic 点に前命題を適用するには、適合する根で Hochschild 二次に
恒等に作用するものが存在する範囲を別に確定する必要がある。
\end{remark}

% !TeX root = ../nc-quadric-moduli.tex
% 2026-11-21 / v1.12
\section{可換 quadric の全一次方向を覆う二次 AS 被覆}
\label{sec:quadric-first-order-as}

$k$ は標数零、$A=k[\varepsilon]/(\varepsilon^2)$、
$P=k[x_0,x_1,x_2,x_3]$、$q=x_0x_3-x_1x_2$ とする。
$\epsilon_{ijkl}$ は $\epsilon_{0123}=1$ で定まる交代記号である。

\begin{proposition}[一次の Jacobian 被覆]
\label{prop:quadric-first-order-jacobian-cover}
$f\in P_2/kq$ を取り、$q_i=\partial_iq$、$f_i=\partial_if$ と書く。
$P$ 上の交代二重導分 $\gamma_f$ を
\[
 \gamma_f(x_i,x_j)=\sum_{k,l=0}^3\epsilon_{ijkl}q_kf_l
\]
で定める。$f$ の $kq$ による変更はこの式を変えない。
$\mathcal A_f=P\otimes_kA$ に積
\[
 a\star b=ab+\frac{\varepsilon}{2}\gamma_f(a,b)
\]
を入れると、これは次数別有限自由の二次 AS4 被覆となり、
ベクトル空間の同一視で指定した $q\in(\mathcal A_f)_2$ は中心正則である。
商 $\mathcal B_f=\mathcal A_f/(q)$ も全次数有限自由であり、
階数は $(n+1)^2$ である。
これらの平坦性・相対 Koszul 完全性・相対 Gorenstein 双対は底変更と両立する。
任意の Noetherian $A$-代数への底変更後にも左右 noetherian である。
\end{proposition}
\begin{proof}
二重導分の Leibniz 則により
\[
 \gamma_f(a,b)c+\gamma_f(ab,c)
 =a\gamma_f(b,c)+\gamma_f(a,bc).
\]
従って $\varepsilon^2=0$ を使えば積は結合的である。
この一次の結合性には高次の量子化定理を用いていない。
交代記号の縮約から $\gamma_f(q,-)=0$ なので
$q\star a=qa=a\star q$ である。
従って $q$ のイデアルは通常の $qP\otimes A$ に一致し、
商は $(P/(q))\otimes A$ と同じ次数別加群を持つ。
通常の $q$ の乗法は単射なので、左右正則性も従う。

二次提示は
\[
 x_ix_j-x_jx_i=\varepsilon\gamma_f(x_i,x_j)
\]
で与えられる。右辺の二次式を任意の固定順序で持ち上げてよい。
その選択の差には再び $\varepsilon$ が付くため、全体の差は零となる。
この関係で単項式を順序付けると、主項の inversion は減少し、
補正項には $\varepsilon$ が付く。補正項の再順序付けは
$\varepsilon^2=0$ によって可換計算になる。
従って各次数は通常の順序付き単項式で張られる。
明示的な $\star$ 積の中でその像は、通常の単項式から
$1+\varepsilon M_n$ という可逆行列で得られるので独立でもある。
これで二次提示と全次数 PBW 基底を得る。

閉 fiber は四変数多項式環である。
三次の平坦性と閉 fiber の Koszul 性から
補題~\ref{lem:relative-koszul-degree-three} を適用し、
左右の相対 Koszul 解消と、その全ての syzygy 束の底変更を得る。
二次双対の階数は $(1,4,6,4,1)$、五次以降は零であり、
Frobenius pairing は閉 fiber の外積の完全 pairing を持ち上げるので完全である。
従って標準 Koszul 複体の双対は最高次数の $A(4)$ だけをコホモロジーに持ち、
相対 AS4 条件を満たす。

最後に $\varepsilon\mathcal A_f$ の平方は零で、
商は noetherian な $P$、このイデアルは左右 $P$-加群として有限生成である。
よって左右 noetherian である。
Noetherian な $A$-代数 $D$ へ底変更しても、商は
$(D/\varepsilon D)[x_0,x_1,x_2,x_3]$、nilpotent ideal の係数は
有限生成 $D/\varepsilon D$-加群 $\varepsilon D$ なので同じ議論が通る。
商 $\mathcal B_f$ についても従う。
\end{proof}

ここで $q$ は次数別加群 $P\otimes A$ 内の元である。
非可換の積記号を使って $x_0\star x_3-x_1\star x_2$ と書いた元とは、
一般に $\varepsilon$ の二次補正だけ異なる。

\begin{lemma}[全 bivector 方向の回復]
\label{lem:quadric-jacobian-bivectors}
$Q=\operatorname{Proj}(P/(q))$ の $x_0\ne0$ の chart で
$z=x_2/x_0$、$w=x_1/x_0$ と置く。
$\mathcal B_f$ の次数零局所化の一次交換子は
\[
 \{z,w\}=2\frac{f}{x_0^2}
        =2f(1,w,z,zw)
\]
である。従って $f\mapsto\gamma_f|_Q$ は
$P_2/kq\simeq H^0(Q,\bigwedge^2T_Q)$ という線形同型を与える。
\end{lemma}
\begin{proof}
Leibniz 則で分母を処理し、$x_0x_3=x_1x_2$ を使うと
\[
 \{z,w\}
 =\frac{x_0f_0+x_1f_1+x_2f_2+x_3f_3}{x_0^2}
 =\frac{2f}{x_0^2}.
\]
最後は $f$ の斉次次数が二であることによる Euler の式である。
Segre の二次切断制限 $P_2\to H^0(Q,\mathcal O_Q(2,2))$ は
全射で核が $kq$ なので主張を得る。
\end{proof}

\begin{proposition}[Abel 変形の Hochschild 類と係数の規約]
\label{prop:quadric-jacobian-abelian-class}
$\operatorname{qgr}\mathcal B_f$ を閉 fiber $\operatorname{coh}Q$ の
pointed flat Abel 一次変形とみなしたとき、その Hochschild 類は、
各標準 affine chart 上で二コサイクル $\gamma_f/2$ により表される。
HKR の規約を
\[
 v\wedge w\longmapsto
 \left[(a,b)\longmapsto
   \frac{v(a)w(b)-w(a)v(b)}2\right]
\]
と固定すると、対応する bivector は
\[
 \pi_f=2f(1,w,z,zw)\,\partial_z\wedge\partial_w
 \quad (x_0\ne0)
\]
である。従って任意の指定された一次類
$u\leftrightarrow h(z,w)\partial_z\wedge\partial_w$
に対して、$f$ の制限を $h/2$ と取ればこの類を実現できる。
\end{proposition}
\begin{proof}
$s=x_i$ と置く。$s$ の $\star$-冪は通常の冪であり、
任意の $n\ge0$ と $b$ に対して
\[
 \begin{aligned}
 s^{n+1}\star b
  &=\bigl(sb+\varepsilon(n+\tfrac12)\gamma_f(s,b)\bigr)\star s^n,\\
 b\star s^{n+1}
  &=s^n\star\bigl(sb-\varepsilon(n+\tfrac12)\gamma_f(s,b)\bigr).
 \end{aligned}
\]
これは二重導分の Leibniz 則を展開した恒等式である。
また閉 fiber で $s$ は正則であり、$\varepsilon$-filtration により
$\mathcal B_f$ でも左右正則である。
従って $\{s^n:n\ge0\}$ は左右の homogeneous Ore 集合である。

二重導分を通常の可換局所化へ延長すると、局所化の積も
$a\star b=ab+\varepsilon\gamma_f(a,b)/2$ で与えられる。
例えば $s^{-1}$ はそのまま $\star$-逆元であり、
\[
 s^{-1}\star b=s^{-1}b-\frac{\varepsilon}{2}s^{-2}\gamma_f(s,b),
 \qquad
 b\star s^{-1}=s^{-1}b+\frac{\varepsilon}{2}s^{-2}\gamma_f(s,b).
\]
上の Ore 式とこの逆元によって、
$\mathcal B_f[s^{-1}]$ は
$(P/(q))[s^{-1}]\otimes A$ 上のこの積に同型となる。
次数零を取れば
\[
 D_i=(\mathcal B_f[x_i^{-1}])_0
      =\mathcal O_Q(D_+(x_i))\otimes A
\]
に同じ積を入れた環を得る。

Ore 局所化は次数付き加群圏の exact localization であり、
局所化環は次数一の単元を持つため、その次数付き加群圏は
$\operatorname{Mod}D_i$ と同値である。
torsion 加群はこの局所化で零になるので、
$\operatorname{QGr}\mathcal B_f$ から $\operatorname{Mod}D_i$ への
exact localization が得られる。その閉 fiber は
$\operatorname{QCoh}Q$ の $D_+(x_i)$ への通常の制限である。
nilpotent 変形と局所化の対応
\cite[Theorem 7.1, Proposition 8.12]{LowenVanDenBergh2006}
により、これは元の Abel 変形の局所化である。
Hochschild 類の制限はこの局所化の変形類と一致する。
この制限と scheme の Hochschild cohomology の比較には
\cite[\S\S7.7--7.9, author preprint]{LowenVanDenBergh2005}
を用いる。

局所環の変形積の二コサイクルは定義から $\gamma_f/2$ である。
この規約では一次交換子 $[a,b]_\star/\varepsilon=\gamma_f(a,b)$ が
HKR bivector に一致する。
前の補題の affine 計算を適用すれば表示した $\pi_f$ を得る。
$Q$ では $H^1(T_Q)=H^2(\mathcal O_Q)=0$ なので、
$\operatorname{HH}^2(Q)=H^0(\bigwedge^2T_Q)$ の稠密 affine chart への
制限は単射である。従って局所計算は大域の類を一意に決定する。
最後の係数 $1/2$ は、$f\mapsto\pi_f$ に現れる Euler 因子 $2$ を相殺する。
\end{proof}

\begin{lemma}[近恒等の補正と右 Zhang twist の規約]
\label{lem:quadric-first-order-root-twist}
$\delta\in\mathfrak{sl}_2\oplus\mathfrak{sl}_2$ を Segre 表示の $P_1$ に作用させ、
$\alpha=1+\varepsilon\delta$ と置く。
これは $\mathcal A_f$ と $\mathcal B_f$ の次数付き自己同型で、$\alpha(q)=q$ である。
$\mathcal C=\operatorname{qgr}\mathcal B_f$、$F=(1)$ とし、
その全軌道環が $\mathcal B_f$ に一致すると仮定する。
右加群に作用を $\alpha^{-1}$ で twist する自己同値を $G_\alpha$ とすれば、
標識 $G_\alpha(\pi\mathcal B_f)\simeq\pi\mathcal B_f$ を
$b\mapsto\alpha(b)$ で取ることにより、
$G_\alpha$ の全軌道環上の作用は $\alpha$ となる。
従って $s=G_\alpha F$ の右規約の全軌道環は
\[
 a*_{s}b=\alpha^{\deg b}(a)*_Fb
\]
という右 Zhang twist $\mathcal B_f^\alpha$ であり、
$\mathcal A_f^\alpha$ はその相対 AS4 被覆を与える。
\end{lemma}
\begin{proof}
$\delta$ は通常の $P$ の次数零導分であり、Segre 方程式を固定する。
$\alpha$ が $\star$ 積を保つ式の一次項には、通常の積の導分条件だけが現れ、
$\gamma_f$ への追加の作用には $\varepsilon^2$ が付く。
従って自己同型性と $\alpha(q)=q$ が従う。
$G_\alpha$ は $F$ と可換し、上記の標識で次数 $n$ の射に $\alpha$ を作用させる。
右の軌道積は $a*_sb=s^{\deg b}(a)\circ b$ なので表示した式を得る。
被覆と切断の Zhang twist は全次数加群・Koszul 解消・Frobenius 双対を保ち、
中心正則 $q$ は正則正規元となる。
\end{proof}


\subsection{可換 quadric の全一次変形に対する適合表示}
\label{subsec:quadric-first-order-bridge}

$k$ は標数零、$A=k[\varepsilon]/(\varepsilon^2)$ とする。
$Q=\mathbf P^1_k\times\mathbf P^1_k$、$L=\mathcal O_Q(1,1)$ と置き、
命題~\ref{prop:quadric-first-order-jacobian-cover} の
$\mathcal B_f=\mathcal A_f/(q)$ を用いる。
以下の最初の補題では、変形された圏が $\mathfrak H_2$ に属するとは仮定しない。

\begin{lemma}[商圏の平坦性と全次数の飽和]
\label{lem:quadric-first-order-flat-saturated}
$\mathcal G_A=\operatorname{QGr}\mathcal B_f$、
$\mathcal C_A=\operatorname{qgr}\mathcal B_f$、
$\mathcal O_A=\pi\mathcal B_f$、$F=(1)$ と置く。
このとき $(\mathcal C_A,\mathcal O_A)$ は
$(\operatorname{coh}Q,\mathcal O_Q)$ の flat pointed Abel 変形であり、
$\mathcal O_A$ は $A$ 上平坦である。
自然な射は全次数で同型
\[
 \mathcal B_f\xrightarrow{\ \sim\ }
 \bigoplus_{n\ge0}\operatorname{Hom}_{\mathcal C_A}
       (\mathcal O_A,F^n\mathcal O_A)
\]
となる。また $\mathcal O_A$ は相対例外的である。
\end{lemma}
\begin{proof}
$\mathcal B_f$ は右 noetherian、次数一生成、全次数 $A$-自由である。
従って $\operatorname{QGr}\mathcal B_f$ は局所 noetherian な Gabriel 商であり、
その noetherian 部分は $\operatorname{qgr}\mathcal B_f$ である。
完全な商関手 $\pi$ は有限提示 $A$-加群との tensor と可換するので、
各 $\pi\mathcal B_f(i)$ は $A$ 上平坦である。
各 noetherian 対象はこれらの有限直和の商だから、
任意の有限提示 $A$-加群 $M$ に対して
$\operatorname{Tor}^A_1(M,-)$ は co-effaceable である。
従って \cite[Proposition 3.1, Definition 3.2]{LowenVanDenBergh2006}
の意味で $\mathcal C_A$ は flat Abel 圏である。
この議論は命題~\ref{prop:artin-relations-flat-abelian} の証明の一般部分であり、
反 Serre の次数移動または $\mathfrak H_2$ 所属を使わない。

閉ファイバーも同命題の証明で確認できる。
具体的には $\mathcal B_0=P/(q)$-加群を $\mathcal B_f$-加群とみなす関手は
torsion を保ち、商圏上で完全充満な関手を与える。
その像は $\varepsilon$ で消される対象全体である。
最後の主張には、商関手の $A$-線形右随伴 $\omega$ を使えばよい。
$\varepsilon X=0$ なら $\varepsilon\omega X=0$ であり、
$\pi\omega X\simeq X$ である。
従って閉ファイバーは $\operatorname{QCoh}Q$、noetherian 部分は
$\operatorname{coh}Q$ となり、$\mathcal O_A$ の閉ファイバーは $\mathcal O_Q$ である。

全軌道環の一致は foundation を使う前に証明できる。
$i:\operatorname{QCoh}Q\to\mathcal G_A$ を閉ファイバーの完全な包含とする。
flat な始域 $\mathcal O_A$ に対する change-of-rings により
\[
 \operatorname{Ext}^j_{\mathcal G_A}(\mathcal O_A,iN)
 \simeq\operatorname{Ext}^j_Q(\mathcal O_Q,N)\qquad(j\ge0).
\]
実際、\cite[Proposition 4.7]{LowenVanDenBergh2006} を反対圏へ適用すると、
$E_2^{p,q}=\operatorname{Ext}^p_Q(\operatorname{Tor}^A_q(k,\mathcal O_A),N)$
から左辺へのスペクトル系列を得る。
$\mathcal O_A$ の平坦性によって $q=0$ 以外は零となる。
$F^n\mathcal O_A$ も平坦なので、$n\ge0$ に対して完全列
\[
 0\longrightarrow iL^n\xrightarrow{j}F^n\mathcal O_A
   \xrightarrow{r}iL^n\longrightarrow0
\]
があり、$\varepsilon$ の作用は $j\circ r$ に一致する。
$H^{>0}(Q,L^n)=0$ を使って Hom と Ext を計算すれば
\[
 0\longrightarrow V_n\longrightarrow H_n\longrightarrow V_n\longrightarrow0,
 \qquad V_n=H^0(Q,L^n),\quad
 H_n=\operatorname{Hom}(\mathcal O_A,F^n\mathcal O_A),
\]
を得て、$H_n$ 上でも $\varepsilon=j\circ r$ である。
従って $\ker(\varepsilon:H_n\to H_n)=\varepsilon H_n$、
$H_n/\varepsilon H_n=V_n$ となり、$H_n$ は階数 $(n+1)^2$ の自由 $A$-加群である。
自然な射 $(\mathcal B_f)_n\to H_n$ は閉ファイバーで通常の全切断環の同型になる。
両辺が同階数自由なので、この射は同型である。
自然な射は右軌道積と整合するため、次数付き環の同型となる。
同じ完全列で高次 Ext は零となる。
特に $n=0$ から $\operatorname{End}(\mathcal O_A)=A$ と
$\operatorname{Ext}^{>0}(\mathcal O_A,\mathcal O_A)=0$ を得る。
\end{proof}

\begin{lemma}[平方根補正の全近恒等方向は次数付き自己同型から来る]
\label{lem:quadric-first-order-correction-geometric}
$\mathcal G_A$ の finite foundation 表示を取り、
その近恒等な derived Picard 群の核を $K_A$ と書く。
命題~\ref{prop:artin-serre-square-correction} の補正元は、
ある $\delta\in\mathfrak{sl}_2\oplus\mathfrak{sl}_2$ と
$\alpha=1+\varepsilon\delta$ に対する $G_\alpha$ で表される。
ここで $G_\alpha$ は右加群の作用を $\alpha^{-1}$ で twist する関手である。
\end{lemma}
\begin{proof}
前補題は flat Abel 変形であることを既に与えたから、
定理~\ref{thm:artin-foundation-lift} を適用できる。
補題~\ref{lem:artin-infinitesimal-dpic} を一次基底に適用すると
\[
 K_A\simeq\varepsilon\operatorname{HH}^1(Q)
 =\varepsilon H^0(Q,T_Q)
 =\varepsilon(\mathfrak{sl}_2\oplus\mathfrak{sl}_2)
\]
を得る。第二の等号には $H^1(Q,\mathcal O_Q)=0$ を使う。
Segre の四次元一次切断空間上の traceless な作用は $q$ を固定し、
補題~\ref{lem:quadric-first-order-root-twist} により
$\alpha=1+\varepsilon\delta$ は $\mathcal A_f$ と $\mathcal B_f$ の自己同型になる。

対応する $K_A$ の元が $+\varepsilon\delta$ であることを規約込みで確認する。
右作用を $\alpha^{-1}$ で変えた tensor 双加群
${}_1R_{\alpha^{-1}}$ は、写像 $r\mapsto\alpha(r)$ により
${}_\alpha R_1$ と同型である。
従って近恒等双加群を左作用の自己同型で記述したとき、その一次導分は $+\delta$ である。
標準 affine chart への Ore 局所化でも同じ式が成立し、
Hochschild 一次類は関数への導分 $\delta$ に制限される。
ここでは特殊ファイバー上の恒等関手との標準的同定を固定している。
$H^1(Q,\mathcal O_Q)=0$ によりこの局所ベクトル場は大域の
$\operatorname{HH}^1$ 類を決定する。
また閉ファイバーの大域中心は $k$ なので、特殊ファイバーの恒等関手との
同定を忘れてもこの近恒等核の記述は変わらない。
従って $\delta\mapsto[G_\alpha]$ は表示した $K_A$ の同型であり、
必要な補正元はこの形で実現される。
\end{proof}

\begin{theorem}[可換 quadric の全一次方向は相対 $\mathfrak H_2$ に入る]
\label{thm:quadric-all-first-order-h2}
$u\in\operatorname{HH}^2(Q)$ を任意に取る。
この一次類を持つ $(\operatorname{coh}Q,\mathcal O_Q)$ の flat pointed Abel 変形には、
$s_0=(-)\otimes\mathcal O_Q(1,1)$ の持上げ $s_A$ と相対 AS4 被覆が存在し、
定義~\ref{def:relative-h2} の全有限型基底変更を含む相対族を構成できる。
従って、この点での $\mathfrak H_2$ から pointed Abel 変形への一次忘却写像は全射である。
\end{theorem}
\begin{proof}
命題~\ref{prop:quadric-jacobian-abelian-class} の HKR 規約で
$u=h(z,w)\partial_z\wedge\partial_w$ と表し、
$f(1,w,z,zw)=h(z,w)/2$ を選ぶ。
前の平坦性・飽和補題と同命題により、
$(\operatorname{qgr}\mathcal B_f,\pi\mathcal B_f)$ はちょうど一次類 $u$ の変形である。
この選択は全九次元の $\operatorname{HH}^2(Q)$ を覆う。

$F=(1)$ は完全自己同値で、閉ファイバーは $s_0$ である。
その Morita 核は compact generator の像から得られる可逆 perfect 双加群である。
$t_A=S_A^{-1}[2]$ と置くと、$d=F^2t_A^{-1}$ は $K_A$ に属する。
命題~\ref{prop:artin-serre-square-correction} に従い
$e=d^{-1/2}$ と取る。
前補題によって $e=[G_\alpha]$ となる $\alpha$ を選べるので、
\[
 s_A=G_\alpha F,\qquad s_A^2\simeq t_A
\]
を得る。$G_\alpha$ は $F$ と可換する。

補題~\ref{lem:quadric-first-order-root-twist} と全軌道環の一致により、
$s_A$ の全軌道環は右 Zhang twist $\mathcal B_f^\alpha$ であり、
\[
 \mathcal B_f^\alpha\simeq\mathcal A_f^\alpha/(q)
\]
が相対 AS4 被覆となる。
各次数は自由で、$q$ は左右正則正規である。
従って命題~\ref{prop:artin-witness-relative-promotion} の仮定が全て成立し、
任意の有限型 $A$-代数への基底変更を含む相対 $\mathfrak H_2$ 族を得る。

最後に、flat Abel 一次変形の同型類はその Hochschild 二次類で分類される
\cite[Theorem 3.1 and Remark 3.5, author preprint]{LowenVanDenBergh2005}。
$\operatorname{Ext}^{1,2}_Q(\mathcal O_Q,\mathcal O_Q)=0$ なので、
構造対象の flat な持上げは存在し、その同型類は一意である。
従って同じ $u$ を持つ任意の pointed Abel 一次変形に以上の表示を移せる。
これは定理~\ref{thm:quadric-all-artin-h2} の一次の場合に対する具体的構成である。
\end{proof}


\begin{corollary}[可換点の内在的な接空間]
\label{cor:quadric-intrinsic-tangent}
可換対象 $x=(\operatorname{coh}Q,\mathcal O_Q)$ において、
一次変形の忘却により自然な同定
\[
 T_x\mathfrak H_2\simeq\operatorname{HH}^2(Q)
      \simeq H^0(Q,\mathcal O_Q(2,2))
\]
を得る。従って内在的接空間は九次元である。
反 Serre 関係式を表す点を $\bar x\in[X_3/\operatorname{PGL}_9]$ と書けば、
\[
 \dim H^0T_{\bar x}[X_3/\operatorname{PGL}_9]=15,
 \qquad \dim\ker(T_{\bar x}\to\operatorname{HH}^2(Q))=6.
\]
後者の六次元は canonical 反 Serre 偏極からのずれであり、
$\mathfrak H_2$ の追加変形方向ではない。
\end{corollary}
\begin{proof}
全射性は前定理による。
基底 $A$ 上の pointed exact 同値は Ind 圏へ延長し、
各有限型 $A'$-作用へ適用することで module-object 圏の同値を与える。
相対族の定義における基底変更同値により、この延長は自然同型を除いて一意である。
根と被覆を標識に加えていないので、同じ pointed Abel 一次変形を持つ
二つの相対 $\mathfrak H_2$ 変形は同型である。これが単射性を与える。

線形構造との整合も確認できる。
有限次元 $k$-空間 $I$ と二乗零拡大 $A=k\oplus I$ に対し、基底 $\epsilon_a$ を選んで
$\gamma=\sum_a\epsilon_a\gamma_{f_a}$ とし、$ab+\gamma(a,b)/2$ を用いる。
$I^2=0$ なので PBW・中心切断・相対 AS4 の証明は同じである。
飽和の証明では $0\to I\otimes_kL^n\to F^n\mathcal O_A\to L^n\to0$ を用い、
軌道 Hom は $A$ 上自由となる。
近恒等核は $I\otimes_k\operatorname{HH}^1(Q)$ であり、
補正は $1+\sum_a\epsilon_a\delta_a$ で表される。
従って全ての二乗零拡大で同じ全射・単射を得て、$I$ の線形写像と可換する。
このことが表示した自然な接空間の同定を与える。
最後の二式は $\dim\operatorname{HH}^1(Q)=6$ と
系~\ref{cor:artin-polarization-retraction} の分裂完全列から従う。
\end{proof}

全 Artin 変形の認識は系~\ref{cor:quadric-formal-deformation-equivalence} で得られる。
これらの形式的結果だけから内在的 $\mathfrak H_2$ 全体の代数性は従わない。

% !TeX root = ../nc-quadric-moduli.tex
% 2026-11-22 / v0.2（代数性・framed実現の正本をB内で相互参照）
\section{相対反 Serre 再構成と三次表示の認識}
\label{sec:moduli-antiserre-recognition}

本節の相対命題は第~\ref{sec:moduli-relative}節の公理の下で扱う。
体上の分類、normal helix、Veronese と foundation の入力は
分類論文 \cite{QuadricIntrinsic2026} に依存する。
相対効果性と全体の有限型代数性は、ここでの全忠実性だけからは導かない。

\begin{proposition}[相対的な反 Serre 再構成]
\label{prop:relative-h2-antiserre-reconstruction}
定義~\ref{def:relative-h2} の族について、相対反 Serre 関手の heart 上の代表
$t=S^{-1}[2]$ は局所的な $s^2$ と一致し、選択を除いて降下する。
局所的な軌道環
\[
 R=\bigoplus_{n\ge0}\operatorname{Hom}(\mathcal O,t^n\mathcal O)
\]
は、次数スカラーの比較を許して、基底変更と整合する。
各 $R_n$ は階数 $(2n+1)^2$ の有限局所自由加群であり、
\[
 R\simeq T(R_1)/(J),\qquad
 J=\ker(R_1^{\otimes2}\to R_2),\qquad\operatorname{rk}J=56.
\]
さらに局所 affine 基底で
\[
 (\mathcal C_U,\mathcal O,t)\simeq(\qgr R,\pi R,(1))
\]
が成立する。
\end{proposition}
\begin{proof}
根を選んだ fppf chart では $t=s^2$ と置ける。
compact 導来部分の二つの代表の比較は Serre 双対の一意性により決まる。
この比較は AZ 生成対象上で定義され、その全ての射と可換するため、
二段の AZ 表示を用いて全 coherent heart へ一意に延長される。
具体的には対象を十分負の軌道対象の有限直和の余核として表す。
二つの表示はさらに十分負の表示で比較でき、AZ の Hom 全射条件が
その比較に必要な lift を与える。余核上の比較はこれにより表示に依存しない。

同じ議論と右規約の軌道積により
$R\simeq(\mathcal B/(g))^{(2)}$。
正則切断の完全列と $\mathcal B$ の Koszul 型から階数公式が従う。
相対二次 Veronese の補題~\cite[相対二次 Veronese の補題]{QuadricIntrinsic2026}
により二次提示と基底変更整合性が従う。

分類論文の ample square 補題~\cite[ample square の補題]{QuadricIntrinsic2026} の証明は、
有限局所自由な $B_1$ に対する評価射を用いれば $A$ 上でもそのまま成立する。
従って $(\mathcal O,t)$ は AZ ample である。
有限 Hom 条件と $\operatorname{End}(\mathcal O)=A$ を用いて
Artin--Zhang の再構成定理
\cite[Theorem 4.5(1)]{AZ1994}
（一般の End 環に対する形は \cite[Theorem 7.3]{OACU2026} にも再録）
を適用すると上の同値を得る。
基底変更との整合性は、逆向きの関手を自由 graded 加群のシフトから
$t$ の軌道対象への写像として作り、有限表示の余核へ延長することで確認できる。
通常の基底変更は右完全なのでこの構成と可換する。
全ての対象の低次数の $\Gamma_*$ が非平坦底変更と可換すると仮定する必要はない。
\end{proof}

\begin{proposition}[fppf stack 性と三次表示への全忠実性]
\label{prop:relative-h2-stack-recognition}
定義~\ref{def:relative-h2} の二 groupoid は、二射を除くと
fppf stack in groupoids となる。
相対反 Serre 関係式は全忠実な射
\[
 \Phi:\mathfrak H_2\longrightarrow[X_3/\operatorname{PGL}_9]
\]
を与える。
従ってその像は、同型で飽和した全部分 stack として定義される。
この主張は像の有限表示性・局所閉性・代数性を含まない。
\end{proposition}
\begin{proof}
圏の系は category-valued stack とその基底変更比較を貼り合わせることで
有効に降下する。Abel 性と完全性は完全な fppf 引戻しで検出され、
finite presentation は module-object descent により降下する。
noetherian 性も、昇鎖を忠実平坦に引き戻して検出できる。
有限型 affine base change に対する Ind 圏の局所 noetherian 性と
module-object 比較は、同じ descent と基底変更の随伴から得る。
Hom の有限性は有限加群の忠実平坦降下で確認できる。
compact generation、compact 導来部分、相対 Serre、perfect Hom および
それらの基底変更は、定義上 fppf 局所 chart の存在条件である。
従って、貼合せた系も同じ局所 chart を使ってこれらを満たす。
ここで局所導来圏の descent の heart が同じであることだけから
大域的な unbounded 導来圏の同値を結論していない。
根と AS 被覆は初めから fppf 局所存在する性質であるので、
根同士を重なりで同一視したり、被覆に descent datum を追加したりする必要はない。
pointed 一射の二自己同型は前補題で消える。
したがって一射の同型類による descent cocycle には一意な二射の lift があり、
三重の整合性も一意性から従う。これで stack in groupoids を得る。

前命題の $R_1$ の枠を選べば $J$ が $\operatorname{Gr}(56,81)$ の点を与える。
全次数の Hilbert 階数から三次関係行列の階数は 176 なので、点は $X_3$ に入る。
比較同型の変更は $R_n$ に単元の $n$ 乗を掛けるだけである。
従って射影枠と関係式は大域に降下し、上の $\Phi$ が定まる。

pointed exact 同値はcompact 導来部分と Serre duality を保つので、
局所的に反 Serre 環の graded 同型を誘導する。
この同型の異なる lift は次数スカラーだけ異なる。
逆に二つの族の像の間の $\operatorname{PGL}_9$ 同型は fppf 局所で
$\operatorname{GL}_9$ に持ち上がり、二次提示により全 graded 環の同型となる。
AZ 再構成はこれを pointed exact 同値へ送る。
異なる lift の次数スカラーは graded 加群の各次数に scalar を掛ける自然同型で消える。
標識上恒等のこの二射は前補題により一意なので、局所同値は降下する。
同じ核計算により射の忠実性も成り立つ。
\end{proof}

\begin{remark}[普遍環の次数別線束と認識条件]
\label{rem:relative-h2-line-twist}
$[X_3/\operatorname{PGL}_9]$ 上の普遍環は
$L=(\det R_1)^{-11}(\det R_2)^2$ を用いて
$R_n^{\mathrm{norm}}=R_n\otimes L^{-n}$ と正規化している。
この正規化は pointed $\qgr$ 圏を変えないが、その次数移動を
元の移動と同じにするわけではない。
自然な同値の下で
\[
 (1)_{R^{\mathrm{norm}}}\simeq t\otimes L^{-1}.
\]
従って普遍環に $(1)=S^{-1}[2]$ を大域に課して像を定義してはならない。
相対反 Serre 環との次数別線束 twist を許すか、fppf 局所でこの条件を課す。
平方根の存在条件自体は、常に真の反 Serre 関手 $t$ に対して課す。
\end{remark}

\begin{proposition}[体上の対象の回復]
\label{prop:relative-h2-field-recovery}
$k$ 上の体としての値を取る評価 functor は
\[
 \mathfrak H_2(k)\simeq\mathcal H_2(k)
\]
という pointed groupoid の同値を与える。
右辺では一射を標識に適合する自然同型で割る。
この主張は任意 noetherian 基底への延長を含まない。
\end{proposition}
\begin{proof}
左から右は (RQ2) の幾何 fiber 条件による。
特に元の定義における全 $D^b(\mathcal C)$ 上の Serre 条件を、
compact 部分だけの相対 Serre 条件から推論していない。
ここで $k$ は代数閉体なので、非空の有限型 fppf chart は $k$-点を持つ。
従って局所存在として課した平方根、被覆、相対導来条件はその点へ引き戻せる。
非代数閉体上で fppf 局所に根を持つという条件と、その体上で根を持つという
条件を同一視する主張ではない。
逆に $(\mathcal C,\mathcal O)\in\mathcal H_2(k)$ を取り、
$\mathcal A=\operatorname{Ind}\mathcal C$ とする。
有限型 affine $U=\operatorname{Spec}A$ 上で
\[
 \mathcal C_U=(\mathcal A_A)^{\mathrm{noeth}}
\]
と置く。Artin--Zhang の基底変更と flat descent により
これは定義~\ref{def:relative-abelian-system} の系を与える。
局所 noetherian 性について使うのは有限型基底変更の B6.3 である。
$\mathcal O_U=A\otimes_k\mathcal O$ は $A$ 上平坦である。

分類論文の normal helix 定理~\cite[normal helix の定理]{QuadricIntrinsic2026} の full strong exceptional foundation の和を
$T_0$、その有限次元 End 環を $\Lambda$ とする。
locally noetherian Grothendieck 圏は locally coherent であり、
その noetherian 対象は finitely presented である。
\cite[Proposition 6.34]{LowenVanDenBergh2006} により、
これらの対象からの全ての $\operatorname{Ext}^i$ は filtered colimit と可換する。
noetherian 両端を持つ有限段の Yoneda 拡大は、順に noetherian 部分対象を
取ることで $\mathcal C$ の中の拡大で代表できるので、
$\operatorname{Ext}^i_{\mathcal C}$ と $\operatorname{Ext}^i_{\mathcal A}$ は一致する。
$\mathcal C$ の homological dimension は二であるから、任意の
$E\in\mathcal C$ について $\operatorname{Ext}^i_{\mathcal A}(E,-)=0$（$i>2$）となる。
この有限 cohomological amplitude と直和との可換性を導来 Hom に適用すれば、
$E$、特に $T_0$ は $D(\mathcal A)$ で compact である。
$\operatorname{thick}(T_0)=D^b(\mathcal C)$ であり、
$\mathcal C$ の対象は $D(\mathcal A)$ を localizing に生成するので、
$T_0$ は compact generator となる。

$A$ は $k$ 上平坦であるため、忘却 functor
$\mathcal A_A\to\mathcal A$ は injective を保ち、
その左随伴は $A\otimes_k-$ である。
従って $T_A=A\otimes_kT_0$ に対し
\[
 \mathbf R\operatorname{Hom}_{\mathcal A_A}(T_A,M)
 \simeq\mathbf R\operatorname{Hom}_{\mathcal A}(T_0,\operatorname{forget}M)
\]
となる。忘却は直和を保ち保守的なので $T_A$ も compact generator である。
また $\mathbf R\operatorname{End}(T_A)=\Lambda\otimes_kA$ は次数零に集中する。
導来 Morita \cite[\S4.3, \S8.2]{Keller1994} により
\[
 D(\mathcal A_A)^c\simeq\operatorname{Perf}(\Lambda\otimes_kA).
\]
有限自由 foundation の相対 bar 解消
（分類論文の相対 foundation の補題~\cite[相対 foundation の補題]{QuadricIntrinsic2026}）により、
右辺は相対 smooth proper であり、任意の本節の基底変更と可換する。
その対象は $T_A$ の有限 cones と直和因子から作られるので、
左辺は $D^b(\mathcal C_U)$ に含まれる。
Serre 関手は有限 perfect 双加群 $D\Lambda$ の基底変更で与えられる。
元の完全自己同値 $s$ は Ind 圏へ colimit を保つ完全自己同値として延長し、
作用準同型にも $s$ を適用することで加群対象の完全自己同値 $s_A$ を与える。
Morita 表示で $s$ を表す双加群は
$\mathbf R\operatorname{Hom}(T_0,sT_0)$ であり、上述の随伴によって
$s_A$ の双加群はその $A$ への tensor である。
従って元の自然同型 $s^2\simeq S^{-1}[2]$ もこの双加群同型の tensor として
延長され、$s_A^2\simeq S_A^{-1}[2]$ を与える。

相対 AZ 条件と Hom 有限性も確認する。
任意の $M\in\mathcal C_U$ は $A\otimes_kN$、$N\in\mathcal C$、の
有限直和の商である。
実際、忘却した $M$ は noetherian 部分対象の filtered colimit であり、
随伴によるそれらからの像は $M$ を生成する。
$M$ の有限生成性から有限個で足りる。
この全射の核も noetherian なので、同じ形の被覆を必要な有限段だけ取れる。
分類論文の pointed derived--exact 比較~\cite[pointed derived--exact 比較]{QuadricIntrinsic2026} の証明は
$\operatorname{Ext}^{i}(s^{-n}\mathcal O,N)=0$（$i>0,n\gg0$）を与える。
$E=s^{-n}\mathcal O$ に対し
\[
 \operatorname{Ext}^{i}_{\mathcal A_A}(A\otimes E,A\otimes N)
       =A\otimes_k\operatorname{Ext}^{i}_{\mathcal A}(E,N),
 \qquad\operatorname{Ext}^{i}_{\mathcal A_A}(A\otimes E,-)=0\ (i>2).
\]
二段の被覆に dimension shifting を適用すれば、任意の $M$ について
正次数 Ext の最終的消滅が従う。
これが AZ の Hom 全射条件を与える。
生成条件は $N$ の生成条件を tensor して得られる。
Hom の $A$ 上有限性は次の三段の被覆で確認する。
$N\in\mathcal C_U$ について $K_0=N$ とし、
$F_i=A\otimes N_i$、$N_i\in\mathcal C$、による完全列
$0\to K_{i+1}\to F_i\to K_i\to0$ を $i=0,1,2$ で取る。
固定した $E=A\otimes M_0$ に対し、
$\operatorname{Ext}^3(E,-)=0$ より
$\operatorname{Ext}^2(E,K_2)$ は $\operatorname{Ext}^2(E,F_2)$ の商であり有限である。
残る二列の長完全列により、順に
$\operatorname{Ext}^1(E,K_1)$、$\operatorname{Hom}(E,N)$ も有限となる。
一般の $M\in\mathcal C_U$ を $A\otimes M_0$ の商として表示すれば
$\operatorname{Hom}(M,N)\subset\operatorname{Hom}(A\otimes M_0,N)$ なので、
$A$ の noetherian 性から所要の有限性を得る。

flat base change $A\to A'$ については、$A'$ を有限自由 $A$-加群の
filtered colimit として表示する（Lazard の定理）。
$M$ は $\mathcal A_A$ の finitely presented 対象であるため、随伴と
有限直和との可換性により
\[
 \begin{aligned}
 \operatorname{Hom}_{\mathcal A_{A'}}(A'\otimes M,A'\otimes N)
  &\simeq\operatorname{Hom}_{\mathcal A_A}(M,A'\otimes N)\\
  &\simeq A'\otimes_A\operatorname{Hom}_{\mathcal A_A}(M,N).
 \end{aligned}
\]
この議論は一般の定義~\ref{def:relative-abelian-system} においても成立する。
相対例外性と軌道 Hom の全次数の基底変更は
$A\otimes_k\operatorname{Hom}(\mathcal O,s^n\mathcal O)$ から直接得る。

被覆を選ぶ際には、体上の全分類の逆構成が与える適合表示を取れる。
その有限な関係式、Koszul 複体と正則切断を $A$ へ tensor する。
左右 noetherian 性は、有限型 $A$ を中心多項式環の商として書くことで保たれる。
幾何体への拡大については、許容四つ組または非退化二形式という
逆構成データが体拡大に安定であることと、同じ逆構成定理を拡大体上に適用する。
cubic 側は分類論文の許容周期データの構成~\cite[許容周期データの構成]{QuadricIntrinsic2026} の標準 hull と
選んだ二周期の scalar extension が、拡大体上の同じ構成になることを用いる。
noncubic 側では明示された Ore 関係式と中心二次切断を tensor すればよい。
また $\operatorname{Ind}\mathcal C\simeq\operatorname{QGr}B$ の
module-object 基底変更は $\operatorname{QGr}(B\otimes_k\Omega)$ と一致する。
これは $\operatorname{Gr}B$ 内の $\Omega$-作用が
$\operatorname{Gr}(B\otimes_k\Omega)$ の構造と同じであり、
torsion 局所化とその右随伴に同じ作用を移すことで確認できる。
上述の逆構成により $B\otimes_k\Omega$ が noetherian であるため、
noetherian 部分を取れば幾何 fiber の $\qgr$ との一致も従う。
ここで arbitrary noetherian $A$ に対する strong noetherian 性を用いていない。
全ての幾何 fiber で元の $\mathcal H_2$ 条件が成立する。

最後に、系は module-object 基底変更で定まるため、$k$ 上の pointed exact
同値は各有限型 affine へ一意に延長する。
逆に系の同値は $k$ 上での値で決まり、その自然同型も同様である。
これで全充満性と本質的全射性を得る。
\end{proof}


\subsection{対角と代数性の判定}

\begin{proposition}[相対認識から対角と代数性判定への還元]
\label{prop:relative-h2-algebraicity-reduction}
$\Phi:\mathfrak H_2\to\mathfrak P=[X_3/\operatorname{PGL}_9]$
が全忠実な fppf スタックの射であるとする。
$X_{\mathrm{as}}=X_3\times_{\mathfrak P}\mathfrak H_2$ と置く。
\begin{enumerate}
\item $\mathfrak H_2$ の対角は affine かつ finite presentation である。
\item $X_{\mathrm{as}}$ は $X_3$ の $\operatorname{PGL}_9$-不変 fppf 部分層であり、
  fppf スタックとして $\mathfrak H_2\simeq
  [X_{\mathrm{as}}/\operatorname{PGL}_9]$ である。
\item 本節の有限型基底の site 上で、$\mathfrak H_2$ が
  $k$ 上有限型の Artin スタックにより表現されるための必要十分条件は、
  $X_{\mathrm{as}}$ が $k$ 上有限型の代数空間により表現されることである。
\end{enumerate}
ここでは $X_{\mathrm{as}}\to X_3$ が開または局所閉 immersion であるとは結論しない。
\end{proposition}
\begin{proof}
$X_3$ は有限型の分離 $k$-スキームである。
$\mathfrak P$ の対角は affine かつ finite presentation である。
実際、その $X_3\times X_3$ 上への引戻しは
\[
 \{(g,x,y):gx=y\}\longrightarrow X_3\times X_3
\]
であり、$\operatorname{PGL}_9\times X_3\times X_3$ の閉部分スキームからの射である。
有限型 $k$-スキームは noetherian なので、この閉条件は有限表示である。
affineness と finite presentation は fppf に降下する。
全忠実性により、任意の二つの $\mathfrak H_2$ 族の Isom 層は
それらの $\mathfrak P$ 像の Isom 層と一致する。
従って $\Delta_{\mathfrak H_2}$ は $\Delta_{\mathfrak P}$ の引戻しであり、(1) を得る。

全忠実性から $X_{\mathrm{as}}$ は部分層である。
射 $X_{\mathrm{as}}\to\mathfrak H_2$ は
$X_3\to\mathfrak P$ の引戻しなので、$\operatorname{PGL}_9$-torsor である。
torsor の有効な降下が (2) を与える。

$X_{\mathrm{as}}$ が有限型代数空間で表されれば、滑らかな有限型群による商は
有限型 Artin スタックであり、所要の site 上の関手を表す。
逆に $\mathfrak H_2$ が有限型 Artin スタックで表されれば、
$\Phi$ はその有限型 smooth chart 上の射と降下から代数スタックの射に延長する。
この引戻しは Artin スタックとなり、
かつ部分層なので全ての慣性が自明である。
慣性自明な Artin スタックは代数空間である。
さらに $X_{\mathrm{as}}\to\mathfrak H_2$ は滑らかな全射で有限表示、
実際には有限型の affine 群の torsor なので quasi-compact でもある。
有限型性はこの射に沿って昇降するため最後の主張が従う。
\end{proof}


本節は内在的な fppf スタックと代数性の判定を構成した。
この判定に必要な $X_{\mathrm{as}}$ の表現可能性は、
\emph{本稿B}の定理~\ref{thm:global-h2-smooth-finite-type} と
系~\ref{cor:global-xas-smooth-scheme} で証明する。
同空間は有限型 smooth スキームとなるが、
本稿の命題~\ref{prop:global-h2-not-ambient-open} により、
一致根対の近傍では $X_3$ の開部分ではない。
既約包絡の幾何点集合や incidence の像との同一視は用いない。
任意 Noetherian 基底上の素朴な Abel 公理との一致、および既存の foundation 表示との
相対同値は、この代数性判定とは別に検討する。

% !TEX root = ../nc-quadric-moduli.tex
% 2026-11-22 / v1.9（相対AS被覆の有限判定とfppf像の反例をBへ移管）
\section{固定した相対切断環の二次 AS 被覆}
\label{sec:relative-as-cover}

平方根を固定した相対軌道環が既に得られている場合には、
その AS 被覆の存在を有限表示の問題へ帰着できる。
以下では被覆を選ぶ空間と、被覆の存在だけを要求する部分関手を区別する。
特に前者の代数性だけから後者の代数性を結論しない。

\begin{theorem}[二次被覆の相対有限判定]
\label{thm:relative-finite-as-cover}
$S$ を Noetherian スキーム、$\mathcal V$ を階数四の局所自由層、
$\mathcal K\subset\mathcal V^{\otimes2}$ を階数七の局所直和因子とする。
二次環
\[
 \mathcal B=T_{\mathcal O_S}(\mathcal V)/(\mathcal K)
\]
について、各 $\mathcal B_n$ が有限局所自由で階数 $(n+1)^2$ を持ち、
全ての代数閉体からの射 $\bar s\to S$ に対して
$\mathcal B_{\bar s}$ が左右 noetherian な Koszul 環であると仮定する。
さらに $S$ の各アフィン開上で $\mathcal B$ は左右 noetherian とする。
この条件の任意の底変更への保存は仮定しない。

$G=\operatorname{Gr}_S(6,\mathcal K)$ とし、
普遍部分束を $\mathcal J$、商直線束を
$\mathcal L=\mathcal K/\mathcal J$ とする。
$G$ 上の二次環を
\[
 \mathcal A=T(\mathcal V)/\mathcal J,\qquad
 \mathcal D=\mathcal A^!=T(\mathcal V^*)/(\mathcal J^\perp)
\]
と書く。
次の条件を満たす局所閉部分スキーム $E_{\mathcal B}\subset G$ が存在し、
$S$ 上有限表示である。
\begin{enumerate}
\item $\mathcal A_3,\mathcal A_4$ が階数 $20,35$ の局所自由層である。
\item $\mathcal D_3,\mathcal D_4,\mathcal D_5$ がそれぞれ
階数 $4,1,0$ の局所自由層である。
次数零、一、二の階数は既に $1,4,6$ である。
\item 次数三の左右の積の像が、層として
\[
 \operatorname{im}(\mathcal L\otimes\mathcal V\to\mathcal A_3)
 =\operatorname{im}(\mathcal V\otimes\mathcal L\to\mathcal A_3)
\]
である。
\item $i=1,2$ に対する左右の積
\[
 \mathcal L\otimes\mathcal A_i\longrightarrow\mathcal A_{i+2},
 \qquad
 \mathcal A_i\otimes\mathcal L\longrightarrow\mathcal A_{i+2}
\]
が局所直和因子への単射である。
\item $i=1,2,3$ に対する積
$\mathcal D_i\otimes\mathcal D_{4-i}\to\mathcal D_4$
が完全な双線型組合せである。
\end{enumerate}

任意の Noetherian $S$-スキーム $T$ と
$\mathcal J_T\subset\mathcal K_T$ について、
$T\to G$ が $E_{\mathcal B}$ を経由することは、次と同値である。
\begin{enumerate}[label=(\alph*)]
\item $\mathcal A_T=T(\mathcal V_T)/(\mathcal J_T)$ の全次数が
有限局所自由であり、全ての幾何ファイバーが型 $(1,4,6,4,1)$ の
AS 正則環である。
\item $\mathcal L_T\subset(\mathcal A_T)_2$ は、局所生成元を選ぶと
左右正則正規元で生成される。
その正則性は任意の底変更後にも保たれ、
$\mathcal A_T/(\mathcal L_T)\simeq\mathcal B_T$ である。
\end{enumerate}
このとき $\operatorname{rank}(\mathcal A_T)_n=\binom{n+3}{3}$ であり、
左右の相対 Koszul 解消および双加群 Koszul 解消は長さ四で、
任意の底変更と両立する。
また、左右それぞれの相対 AS Gorenstein 条件は
\[
 \underline{\operatorname{Ext}}_{\mathcal A_T}^i
       (\mathcal O_T,\mathcal A_T)
 \simeq
 \begin{cases}
  (\mathcal D_T)_4\otimes\mathcal O_T(4),&i=4,\\
  0,&i\ne4
 \end{cases}
\]
となる。最高次数の直線束は局所的に自明化できる。
ここでの長さ四は相対解消の長さであり、
特異な基底上の環全体の大域次元を四と主張するものではない。

各アフィン開上で $\mathcal A_T$ が左右 noetherian であることは、
$\mathcal B_T$ が左右 noetherian であることと同値である。
従って $\mathcal B$ が全ての Noetherian 底変更の後にも左右 noetherian
であるという追加仮定の下では、$E_{\mathcal B}$ は
固定された $\mathcal B$ への全ての相対 noetherian AS 被覆を表現する。
この追加仮定を置かない場合は、$\mathcal B_T$ が左右 noetherian となる
テスト基底に限って同じ結論を用いる。
\end{theorem}
\begin{proof}
まず有限表示の部分スキームを構成する。
任意の固定次数 $n$ で、$\mathcal A_n$ と $\mathcal D_n$ は
有限階数の局所自由層の間の射の余核として表示される。
有限行列の余核が指定階数の局所自由層となる条件は、
より大きな minor の消滅と所定の minor が局所的に単位である条件で
表される。これは余核の局所自由性を普遍的に表す
\emph{scheme-theoretic な}局所閉階数条件である。
従って (1), (2) を課した部分スキームは局所閉である。

その上で (4) は最大階数 minor による開条件であり、
左右の次数三の像は階数四の部分束になる。
二つの部分束の一致は、一方から他方の商への射が零となる
閉条件で表される。同じ階数の部分束間の包含は一致を意味するので、
これが (3) を正確に表す。
(5) は局所自由層間の同型を要求する開条件である。
従って全条件の交叉は所要の $E_{\mathcal B}$ を与える。
特に、幾何点で minor が零というだけで被約構造を取っているのではなく、
これらの射と minor 自体の消滅を課している。
この構成は非被約なテスト基底も含めて普遍的である。

条件 (1)--(5) を満たす $T$ を取る。
各幾何ファイバーに分類論文の有限 ambient 判定~\cite[二次 AS 被覆の有限判定]{QuadricIntrinsic2026} を適用すると、
$(\mathcal A_T)_{\bar t}$ は AS 正則で、
$\mathcal L_{\bar t}$ の生成元は正則正規となる。
特に全てのファイバーは Koszul である。
$\mathcal A_T$ の次数一、二、三は階数 $4,10,20$ の局所自由層なので、
分類論文の相対 Koszul 三次判定~\cite[相対 Koszul 三次判定]{QuadricIntrinsic2026} により
全次数が局所自由であり、相対 Koszul 複体も得られる。
階数はファイバー上の AS Hilbert 級数から決まる。

局所的に $\mathcal L_T$ を $g$ で生成する。
(3) は $g\mathcal V_T=\mathcal V_Tg$ を意味するため、
次数一生成性を使って生成元を一つずつ移すことにより
$g\mathcal A_T=\mathcal A_Tg$ を得る。
ここでは全ての等式が基底上で成立し、底変更後も同じである。
各次数の左右の $g$ 乗法は、既に局所自由と分かった層の間の射であり、
全ての幾何ファイバーで単射である。
従って局所直和因子への単射となり、任意の底変更後も単射である。
これで全次数の左右正則性を得る。
商の同一視は $\mathcal J_T\subset\mathcal K_T$ の表示から直接従う。

$\mathcal D_5=0$ と一次生成性から $\mathcal D_i=0$ ($i>4$) である。
相対 Koszul 解消の係数は $(\mathcal D_i)^*$ なので長さ四となる。
完全な組合せ (5) によって、この解消の $\mathcal A_T$-双対は、
最高次数の直線束、次数四の移動、および Frobenius 自己同型による
捻りを除いて反対側の Koszul 解消に同定される。
従って表示した左右の Ext の式を得る。
双加群 Koszul 複体も各内部次数で有限な局所自由複体であり、
そのファイバーは Koszul 環の対角双加群の解消である。
分類論文の相対 Koszul 三次判定~\cite[相対 Koszul 三次判定]{QuadricIntrinsic2026} と同じ
ファイバーによる完全性の判定から、
各次数で局所分裂完全となり、任意の底変更と両立する。
各項は $\mathcal A_T^e$ 上の有限局所射影加群なので、
これは相対対角双加群の長さ四の射影解消である。

逆に (a), (b) を仮定する。
次数三、四の階数はファイバー上の Hilbert 級数から (1) となる。
全ファイバーが Koszul で次数三までが局所自由であるため、
分類論文の相対 Koszul 三次判定~\cite[相対 Koszul 三次判定]{QuadricIntrinsic2026} により
関係の交叉と二次双対の各次数は底変更と両立して局所自由となる。
ファイバー上の型と Frobenius 性から (2), (5) が従う。
正規性から (3)、正則性と商 $\mathcal B_T$ の局所自由性から
(4) が従う。従って有限条件が必要でもある。

最後に noetherian 性を示す。
$\mathcal A_T$ が左右 noetherian ならその商もそうである。
逆向きはアフィン局所に $g$ と最高次数直線束を自明化して考える。
正則正規性から $g$ による正規化自己同型が存在し、
$(g)$ 進付随次数環は $\mathcal B_T$ 上の skew polynomial 環である。
$\mathcal B_T$ が左右 noetherian ならこの付随次数環もそうである。
内部次数を固定すると $(g)$ 進フィルトレーションは有限だから、
斉次イデアルの初期生成元を有限個持ち上げることで
元の斉次イデアルも有限生成となる。
一般のイデアルには内部次数による最高次成分のイデアルを用いて
同じ持上げを適用できる。よって $\mathcal A_T$ も左右 noetherian である。
この議論は $\mathcal B_T$ の noetherian 性を用いており、
$\mathcal B$ 自身またはその幾何ファイバーだけの noetherian 性から
任意の $\mathcal B_T$ の noetherian 性を導いているのではない。
\end{proof}

\begin{remark}[相対被覆の存在として残る問題]
固定された平方根の相対軌道環 $\mathcal B$ に上の定理を適用できるとき、
被覆が fppf 局所に存在することは、有限表示の射
$E_{\mathcal B}\to S$ が fppf 局所に切断を持つことに等しい。
しかし、その射が滑らか・平坦・全射であるとは証明していない。
平方根そのものの相対存在と、得た被覆が相対 Serre 公式に適合することも
この定理の入力の外にある。
従ってこの有限表示への帰着だけを $\HHstack$ の代数性の証明としてはならない。
\end{remark}

\begin{example}[有限型射の fppf 像は代数的とは限らない]
\label{ex:fppf-image-nonalgebraic}
体 $k$ 上の有限型射
\[
 f:\mathbb A^2_{a,c}\longrightarrow\mathbb A^2_{a,b},
 \qquad (a,c)\longmapsto(a,ac)
\]
を考え、その fppf 層としての像を $I$ とする。
$I$ は $\mathbb A^2_{a,b}$ の部分層であるが、
代数空間では表現できない。
\end{example}
\begin{proof}
アフィン基底 $T=\operatorname{Spec}D$ 上で、
$(a,b)\in D^2$ が $I(T)$ に属することは $b\in aD$ と同値である。
実際、fppf 局所で $b=ac$ と書ければ、
$D/(a)$ における $b$ の像が忠実平坦な被覆上で零となる。
準連接層の切断の降下により既に零であり、$b\in aD$ となる。
逆向きは明らかである。

原点の双対数上の接ベクトルを
$a=u\varepsilon,b=v\varepsilon$ と書くと、
\[
 T_0 I\ \longrightarrow\ T_0\mathbb A^2=k^2
\]
の像は集合として
\[
 \{(u,v)\in k^2\mid u\ne0\}\ \cup\ \{(0,0)\}
\]
となる。これは線型部分空間ではない。
例えば $(1,0)$ と $(-1,1)$ は属するが、その和 $(0,1)$ は属さない。
もし $I$ が代数空間なら、単射 $I\to\mathbb A^2$ が誘導する接写像は
接ベクトル空間間の線型単射であり、その像は線型部分空間になる。
これは矛盾である。
\end{proof}

この例は、幾何点の構成可能な像や scheme-theoretic image を求めることと、
非被約基底上の fppf 局所存在条件を表すこととの相違を示す。
\emph{この例から $\HHstack$ 自身の非代数性は結論しない。}
本研究では、平方根と被覆を忘れた部分関手の代数性を得るために、
さらに局所切断、相対認識、または同値関係の表示を証明する必要がある。


% !TeX root = ../nc-quadric-surface-moduli.tex
% 2026-11-21 / v1.8
\section{反 Serre 偏極の認識と次数付き表示の変形}
\label{sec:antiserre-polarization}

全次数平坦な二次環から相対 Abel 圏を作れることと、
その環の次数移動が圏の内在的反 Serre 関手であることは異なる。
この相違は noncubic の二根対一致点で実際に現れる。
以下では全次数平坦な代数的な族を用い、
pointed Abel 圏を固定しても関係式空間の像が動くことを示す。

\begin{proposition}[pointed Abel 圏が一定となる非自明な二次環の族]
\label{prop:antiserre-polarization-family}
$k$ を標数零の代数閉体とし、
\[
 B=k[u,v]\langle c,d\rangle/(c^2-uv,d^2-uv),\qquad
 \iota(u)=u,\quad\iota(v)=v,\quad\iota(c)=-c,\quad\iota(d)=-d
\]
と置く。
$R=(B^{(2)})^{\iota}$ を
命題~\ref{prop:noncubic-antiserre-descent} の noncubic 反 Serre 環とする。
この表示では、$a\in B_{2m}$、$b\in B_{2n}$ の積は
\[
 a\mathbin{\circ} b=a\,\iota^m(b)
\]
である。
基底 $S=\operatorname{Spec}k[q,q^{-1}]$ 上で
\[
 \sigma_q(u)=qu,\qquad \sigma_q(v)=q^{-1}v,\qquad
 \sigma_q(c)=c,\qquad \sigma_q(d)=d
\]
と定める。
$\sigma_q$ は $\iota$ と可換な次数付き自己同型であり、
$R\otimes k[q,q^{-1}]$ の自己同型でもある。
その\emph{右} Zhang twist を $\mathcal R^{(q)}$ と書き、積を
\[
 a*_q b=\sigma_q^{\deg_R b}(a)\mathbin{\circ}b
 \tag{ZT}\label{eq:polarization-right-zhang}
\]
と定める。
この族は次を満たす。
\begin{enumerate}
\item 全次数有限自由で
$\operatorname{rank}\mathcal R^{(q)}_n=(2n+1)^2$ となり、
左右 noetherian な二次環の族である。
従って $S\to X_3$ および
$S\to[X_3/\operatorname{PGL}_9]$ を与える。
\item Zhang twist の同値により
\[
 (\operatorname{qgr}\mathcal R^{(q)},\pi\mathcal R^{(q)})
 \simeq
 (\operatorname{qgr}(R\otimes k[q,q^{-1}]),\pi(R\otimes k[q,q^{-1}]))
\]
である。この同値は基底上線型、pointed exact であり、
任意の底変更と両立する。
\item $q=1$ のファイバーは $R$ であるが、$q=2$ のファイバーは
$R$ に次数付き同型ではない。
しかも $q=2$ のファイバーは、いかなる $\HH$ 対象の
内在的反 Serre 環にも次数付き同型ではない。
\end{enumerate}
従って $X_3$ には、$\operatorname{qgr}$ と構造対象からは
$\HH$ 対象を与えながら、内在的反 Serre 環の表示ではない幾何点が存在する。
\end{proposition}
\begin{proof}
$uv$ は $\sigma_q$ で固定されるので、二つの二乗関係は保存される。
$\sigma_q$ と $\iota$ の可換性も表示から直接従う。
Zhang twist は次数別の加群を変えないため、平坦性と Hilbert 関数は
一定の族 $R\otimes k[q,q^{-1}]$ から従う。
$R=T(R_1)/(J)$ とすれば、二次関係は
\[
 J_q=(\sigma_q^{-1}\otimes\operatorname{id})(J)
\]
である。全次数の Zhang 同値から、この二次関係が全関係を生成する。
元の $B$ は中心 $k[u,v,h]$（$h=cd+dc$）上で有限である。
従って基底拡大後も左右 noetherian であり、
Veronese および Zhang twist の同値により $\mathcal R^{(q)}$ もそうである。
この議論は Noetherian 基底への底変更後にも適用できる。

次数付き加群の同値は Zhang の構成である
\cite[Theorem 3.1]{Zhang1996}。
係数環と左右の規約を明確にするため、ここでは具体式も記す。
次数 $n$ で $\phi_n=\sigma_q^{-n}$ とすると、
式~(\ref{eq:polarization-right-zhang}) の環は
\[
 a\diamond b=a\mathbin{\circ}\sigma_q^{-\deg_R a}(b)
\]
という通常の左次数による twist に $\phi$ で次数付き同型となる。
この後者の環に対して、次数付き右加群 $M$ の作用を
\[
 m\diamond a=m\,\sigma_q^{-\deg m}(a)
\]
と置けば、同じ次数付き加群上に新しい右作用を得る。
逆自己同型による同じ構成が逆関手となる。
従ってこれは係数環上で exact な同値であり、
自由加群を自由加群へ、torsion 加群を torsion 加群へ送る。
従って表示した pointed exact な $\operatorname{qgr}$ の同値を得る。
全構成は $\sigma_q$ と次数による有限の積から定義されるので、
底変更と両立する。

次数一の中心を比較して最後の主張を示す。
$R_1$ の基底を
\[
 (u^2,uv,v^2,uc,vc,ud,vd,cd,dc)
 \tag{P-basis}\label{eq:polarization-generator-basis}
\]
とする。
元の $R$ では $e=cd-dc$ が非零の中心一次元である。
実際、$e$ は $u,v$ と可換、$c,d$ と反可換であり、
$\iota(e)=e$ である。
従って上の基底の各元に対して、積 $\circ$ に入る $\iota$ の符号が
反可換性の符号を打ち消す。
$cd$ と $dc$ の線型独立性は PBW 基底から従う。

一方、上の九基底の $\sigma_q$-weight は
\[
 (2,0,-2,1,-1,1,-1,0,0)
\]
である。
基底元 $x$ の $\iota$-符号を $\eta_x\in\{1,-1\}$ とすると、
\[
 x*_q(uv)-(uv)*_q x
 =(q^{\operatorname{wt}(x)}-\eta_x)\,uvx.
\]
$q=2$ ではこの係数が零となる基底元は $uv,cd,dc$ だけである。
$B$ の PBW 基底から $uv$ の乗法は単射なので、
中心一次元があれば $z\in\langle uv,cd,dc\rangle$ である。
これらは $\sigma_q$ と $\iota$ で固定され、元の $B$ で $u^2$ と可換だから、
\[
 z*_q u^2-u^2*_q z=(1-q^2)u^2z.
\]
$u^2$ の乗法も単射なので、$q=2$ では $z=0$ となる。
従って
\[
 Z(R)_1\ne0,\qquad Z(\mathcal R^{(2)})_1=0
\]
であり、二つの環は次数付き同型ではない。

もし $\mathcal R^{(2)}$ が別の $\HH$ 対象
$(\mathcal C',\mathcal O')$ の反 Serre 環なら、
反 Serre 豊富性と再構成により
$(\mathcal C',\mathcal O')\simeq
(\operatorname{qgr}\mathcal R^{(2)},\pi\mathcal R^{(2)})$ である。
これは既に示した Zhang 同値によって元の pointed 圏と同値である。
内在的反 Serre 環は pointed exact 同値で不変なので
（定理~\cite[\texttt{thm:antiserre-invariant}]{QuadricIntrinsic2026}）、
$\mathcal R^{(2)}\simeq R$ となって矛盾する。
\end{proof}

\begin{proposition}[一定の Abel 変形へ移る非零接方向]
\label{prop:antiserre-polarization-tangent}
前命題の族を $q=1+\varepsilon$、$\varepsilon^2=0$ に制限する。
次数零の導分
\[
 \delta(u)=u,\quad\delta(v)=-v,\quad\delta(c)=\delta(d)=0
\]
を用いると、積の一次項は
\[
 \mu_1(a,b)=\deg_R(b)\,\delta(a)\mathbin{\circ}b.
\]
これは $\operatorname{HH}^2(R)_0$ の非零類を与える。
従って対応する $X_3$ の接方向は $\operatorname{PGL}_9$ 軌道の接方向ではない。
一方、対応する pointed Abel 圏の一次変形は自明である。
\end{proposition}
\begin{proof}
族全体が結合的なので $\mu_1$ は Hochschild $2$-cocycle である。
その非自明性は中心元上の交換子で検出できる。
\[
 U=u^4,\qquad V=v^4
\]
はともに $R$-次数二であり、元の $R$ の中心に属する。
実際、元の $B$ で中心、$\iota$ で固定され、
次数二の左因子に現れる $\iota^2$ も恒等写像である。
$\delta(U)=4U$、$\delta(V)=-4V$ なので、
変形した積で
\[
 U*V-V*U
 =\varepsilon\bigl(\mu_1(U,V)-\mu_1(V,U)\bigr)
 =16\varepsilon u^4v^4\ne0.
 \tag{P-comm}\label{eq:polarization-central-commutator}
\]
最後の非零性は中央の $u,v$ と交互語からなる PBW 基底による。

もし $\mu_1$ が Hochschild 境界なら、
$\operatorname{id}+\varepsilon h$ による一次の基底変更で
一定の積へ移せる。
しかし元の中心元 $U,V$ について、任意の一次基底変更が交換子へ与える補正は
\[
 [h(U),V]+[U,h(V)]=0
\]
である。
これは式~(\ref{eq:polarization-central-commutator}) に反する。
従って非零の Hochschild 類であり、特に次数一生成元の基底変更に由来する
$\operatorname{PGL}_9$ 軌道方向でもない。
pointed Abel 変形の自明性は前命題の Zhang 同値を双対数へ制限すれば従う。
\end{proof}

\begin{remark}[反 Serre 認識に対する意味]
上の一次変形では、構造対象付き Abel 圏が一定なので、
その内在的反 Serre 関手と反 Serre 軌道環も一定のものに同定される。
従って $\mathcal R^{(q)}$ の次数一移動を、
この内在的反 Serre 関手と同定することはできない。
全次数平坦性と $\operatorname{qgr}$ の再構成だけでは、
関係式空間から内在的 $\HHstack$ への全忠実性は得られない。
必要なのは、表示された次数移動が相対 Serre 双対から得られる
$S^{-1}[2]$ と一致するという追加の認識である。

中心元の交換子による上の証明は標数零で直接成立し、
接空間の数値比較を仮定しない。
独立な有限線型代数の確認では、既存の Grassmann 座標で
$\operatorname{PGL}_9$ 軌道行列の階数は $79$、
この方向を加えた階数は $80$ となる。
さらに六係数方向も加えた行列の階数は $81$、
そこへこの方向を加えると $82$ となる。
これらは
\texttt{computations/verify\_antiserre\_polarization.py}
による法 $101$ の厳密計算で確認した。
軌道行列の次数零導分 $\delta$ による非零核ベクトルを
整数上でも確認するので、最初の二つの階数 $79,80$ は
標数零でも正確な値である。
\end{remark}

% !TEX root = ../nc-quadric-moduli.tex
% 2026-11-22 / v1.11（相対モジュライ論文Bへ証明移管）
\section{有限 foundation 上の反 Serre 偏極の相対認識}
\label{sec:finite-polarization-recognition}

分類論文 \cite[反 Serre 偏極の節]{QuadricIntrinsic2026} の Zhang 族が示すように、平坦な pointed Abel 圏を構成することと、
表示された次数移動を内在的反 Serre 関手と同定することは別問題である。
本節では有限 foundation と次数移動の双加群表示が得られている場合に、
後者を非被約基底も含む有限個の方程式と開条件で判定する。
任意の $X_3$ 族からこの有限表示を構成することは、本節の仮定に含めない。

\begin{lemma}[最低次 Hom を線束にする scheme-theoretic 条件]
\label{lem:lowest-hom-line-locus}
$S$ を局所 noetherian スキーム、$K$ を $S$ 上 perfect 複体とする。
各幾何点で $H^i(K_{\bar s})=0$（$i<0$）、
$\dim H^0(K_{\bar s})\le1$ と仮定する。
このとき有限表示閉部分スキーム $j:Z\hookrightarrow S$ が存在して、
次を満たす。
\begin{enumerate}
\item $\mathcal L=H^0(\mathbf L j^*K)$ は線束であり、
その形成は $Z$ 上の任意の基底変更と可換する。
\item $f:T\to S$ が $Z$ を経由することと、全ての
$g:T'\to T$ に対し $H^0(\mathbf L(fg)^*K)$ が線束となることは同値である。
\end{enumerate}
\end{lemma}
\begin{proof}
perfect 複体の局所的な単位成分の消去により、
$K$ は局所的に有限自由加群の複体
$[K^0\xrightarrow{d^0}K^1\to\cdots\to K^b]$ で表せる。
特に負次数の項は不要である。
$\operatorname{rk}K^0=r$ とする。
仮定より、各幾何点で $d^0$ の階数は $r-1$ 以上である。
$r=0$ の領域には求める $Z$ の点がない。
$r\ge1$ の場合は、$d^0$ のある $(r-1)$ 次 minor が可逆となる
開集合で覆える（$r=1$ では空の minor を $1$ とする）。
その開集合上で行列の初等変形を行えば、
\[
 d^0=\begin{pmatrix}I_{r-1}&0\\0&v\end{pmatrix}
\]
と書ける。$v$ は一本の列であり、行数零の場合も許す。
この表示で $Z$ を $v$ の全成分が零となる閉部分スキームとする。
これは幾何点上の階数条件だけでなく、基底環での実際の消滅条件である。
その上では $\ker d^0$ は最後の一座標が張る直和因子であり、
任意の基底変更後も同じ線束を与える。

逆に任意の環への基底変更後、$H^0(K)$ が線束であるとする。
上の表示では、それは $\operatorname{Ann}(v)$ という基底環のイデアルである。
線束は基底環上忠実であり、$v$ の各成分はこのイデアルを零化するから、
各成分自体が零でなければならない。
従ってこの表示では、$H^0$ が線束であるという条件だけでも
上の閉部分スキームを経由する。普遍的特徴付けによって局所構成は貼り合う。
これは \cite[Tag 0BDH, Lemma 36.31.4]{StacksModuli} の
最低次数 $a=0$、階数 $r=1$ の場合の具体的な行列表示でもある。
\end{proof}

\begin{theorem}[scalar 自己同型を持つ双加群の比較部分スキーム]
\label{thm:bimodule-line-comparison}
$S$ を局所 noetherian スキーム、$\Lambda$ を有限局所自由
$\mathcal O_S$-代数とする。任意の $T\to S$ に対して
\[
 Z(\Lambda_T)=\mathcal O_T
\]
と仮定する。$P,Q$ を perfect な $\Lambda$-双加群とし、
それぞれ導来 tensor 積について可逆とする。
このとき $S$ 上の fppf 部分層
\[
 \mathscr E(P,Q)(T)=
 \begin{cases}
 \{*\},&P_T\simeq Q_T\text{ が }T\text{ の fppf 局所で成立する},\\
 \varnothing,&\text{それ以外}
 \end{cases}
\]
は有限表示の局所閉部分スキーム $S_{P\sim Q}\hookrightarrow S$ によって表される。
その上には標準的な線束
\[
 \mathcal L=H^0\mathbf R\mathcal Hom_{\Lambda^e}(P,Q)
\]
と評価同型 $P\otimes\mathcal L\simeq Q$ が存在する。
この線束と同型は、線束の同型を許して一意である。
全ての構成は任意の基底変更と可換する。
\end{theorem}
\begin{proof}
$\Lambda^e=\Lambda\otimes\Lambda^{\mathrm{op}}$ と置く。
$\Lambda^e$ は基底上有限局所自由であり、$P,Q$ はその perfect 加群なので、
\[
 K=\mathbf R\mathcal Hom_{\Lambda^e}(P,Q)
\]
は基底上 perfect で、任意の導来基底変更と可換する。
$P$ の可逆性により
\[
 \mathbf R\mathcal Hom_{\Lambda_T^e}(P_T,P_T)
 \simeq
 \mathbf R\mathcal Hom_{\Lambda_T^e}(\Lambda_T,\Lambda_T).
\]
右辺は負次数 cohomology を持たず、次数零は仮定より $\mathcal O_T$ である。
従って $\operatorname{Aut}(P_T)=\mathbb G_m(T)$ であり、
この scalar 記述は任意の基底変更後にも成立する。

perfect 複体の幾何ファイバーの cohomology 次元の上半連続性により、
\[
 U=\{s\in S\mid H^i(K_{\bar s})=0\ (i<0),\ 
                      \dim H^0(K_{\bar s})\le1\}
\]
は開集合である。比較同型が局所的に存在する任意の $T\to S$ は
$U$ を経由する。実際、その全幾何点で $K$ は $P$ の自己 Hom と同型となる。
上半連続性と局所的な非負次数表示は、\cite[Tag 0BDH, Lemma 36.31.1]{StacksModuli}
の有限自由複体の消去による。

補題~\ref{lem:lowest-hom-line-locus} を $K|_U$ に適用し、
閉部分スキーム $Z\hookrightarrow U$ と線束
$\mathcal L=H^0(K_Z)$ を得る。
負次数がないため、標準写像 $\mathcal L[0]\to K_Z$ があり、
tensor--Hom 随伴により評価射
\[
 e:P_Z\otimes\mathcal L\longrightarrow Q_Z
\]
を定める。その cone は基底上 perfect なので、それが零となる最大の
開部分 $W\subset Z$ が存在する。
この開条件は、有限自由表示の複体が完全になる条件、または幾何ファイバーの
全 Betti 数が零になる条件である。perfect 複体の Nakayama により、
任意の $T\to Z$ に対し $e_T$ が同型であることと、$T\to Z$ が $W$ を経由する
ことが同値となる。従って $W$ は $S$ の有限表示局所閉部分スキームである。

$T\to S$ 上 fppf 局所に $P_T\simeq Q_T$ とする。
この局所同型の下で $K_T$ の次数零 Hom は $\mathcal O_T$ となり、
任意の更なる基底変更でもそうである。線束性と評価同型は忠実平坦降下するので、
$T\to S$ はまず $Z$、次いで $W$ を経由する。
逆に $T\to W$ なら、線束 $\mathcal L_T$ を局所自明化することで
$P_T\simeq Q_T$ を得る。これで $W$ が所要の部分層を表すことが分かった。

任意の線束 $L$ と同型 $P_T\otimes L\simeq Q_T$ について、
相対 Hom を取ると $L\simeq H^0(K_T)$ となり、評価同型へ移る。
scalar な自己同型の記述から、この比較は一意である。
構成の特徴付けは任意の $T\to S$ を用いており、基底変更との整合性も従う。
\end{proof}

\begin{remark}[fppf 像の一般的障害との関係]
ここでは単に有限表示の比較パラメータ空間の像を取ったのではない。
自己 Hom の次数零が基底環そのものであるため、比較が存在する場合の
Hom は線束になり、その普遍性を有限行列の方程式で課すことができた。
同型の空間はこの部分スキーム上の $\mathbb G_m$-torsor となる。
有限型射の任意の fppf 像が代数的になるという主張ではない。
\end{remark}

\begin{corollary}[有限 foundation 上の反 Serre 比較]
\label{cor:foundation-antiserre-comparison-locus}
$\Lambda$ を分類論文 \cite{QuadricIntrinsic2026} で構成した相対 foundation 代数とし、
$P$ を表示された次数移動を表す可逆 perfect 双加群とする。
相対反 Serre 関手 $S^{-1}[2]$ の双加群を $Q$ とする。
次数移動が反 Serre 関手と基底線束 twist を除いて一致する条件は、
基底の有限表示局所閉部分スキームで表される。
\end{corollary}
\begin{proof}
foundation 代数の対角頂点成分は $\mathcal O_S$ である。
中心元は冪等元と可換なので頂点ごとの scalar であり、
頂点間の正階数の局所自由成分と可換することから全 scalar は等しい。
この議論は任意の基底変更後にも成立するので $Z(\Lambda_T)=\mathcal O_T$。
有限な相対 bar 解消による smooth proper 性と、仮定した双加群表示により
前定理を適用できる。
\end{proof}

\begin{remark}[次数の正規化と適用範囲]
反 Serre 軌道環の正規化で
$L_{\mathrm{deg}}=(\det R_1)^{-11}(\det R_2)^2$ とすると、
正規化した次数移動は $t\otimes L_{\mathrm{deg}}^{-1}$ を表す。
従って前定理で $P$ をこの次数移動、$Q=t$ とすれば、
比較線束は $L_{\mathrm{deg}}$ である。
大域的な $P\simeq Q$ だけを課してこの線束を落としてはならない。

本節は有限 foundation と次数移動の可逆 perfect 双加群表示を与えた後の
偏極認識を解決する。任意の $X_3$ 族に相対 foundation が存在すること、
heart の次数移動がその表示へ移ること、平方根・AS 被覆の局所存在が
同じ部分スキーム上で成立することは別に確認しなければならない。
従って本節だけから $X_{\mathrm{as}}$ 全体の表現可能性は結論しない。
\end{remark}

% !TeX root = ../nc-quadric-surface-moduli.tex
% 2026-11-22 / v1.31（論文Bへ証明を移管）
\section{非可換射影直線束の対角列に関する障害}
\label{sec:ruled-diagonal-obstruction}

以下は、幾何的 quintuple の行列から層双加群 $\mathcal E$ を作り、
$\operatorname{QGr}\Sym(\mathcal E)$ の素朴な対角列を cubic 偏極と同定する方針に対する反例である。
幾何的 quintuple 条件、加群の関係式、全次数の直接飽和をそれぞれ証明する。
既知の入力は非可換対称代数の Noetherian 性と標準二次解消であり、
出典を明記した。本例自体の新規性は未確認である。
この反例は cubic AS 正則代数自身の Noetherian 性を否定せず、
可換点の十分小さい開近傍に制限する別の議論も否定しない。

\subsection{楕円曲線上のデータと幾何的 quintuple}

$k$ を標数零の代数閉体とする。
滑らかな $(2,2)$ 曲線
\[
 i:C\hookrightarrow X\times Y,\qquad X=Y=\PP^1_k
\]
を取り、その二つの射影を $p:C\to X$, $q:C\to Y$ とする。
\[
 L_X=p^*\OO_X(1),\qquad L_Y=q^*\OO_Y(1),\qquad N=L_XL_Y^{-1}
\]
と置き、$N\not\simeq\OO_C$, $N^2\not\simeq\OO_C$ を仮定する。
このような曲線は存在する。例えば楕円曲線上の次数二線束 $L_X,L_Y$ を
$L_XL_Y^{-1}$ が二捩れでないように選べば、二つの二重被覆の積は
$\PP^1\times\PP^1$ への閉埋め込みを与える。
実際、異なる二点が両射影で同一視されるなら
その二点の和が $L_X,L_Y$ の双方に線形同値となる。
また両射影の微分が同じ点で消えるなら $L_X\simeq L_Y\simeq\OO_C(2z)$ となる。
従って積の射は単射かつ非分岐であり、閉埋め込みとなる。

\begin{lemma}\label{ruled-lem:resolution}
\[
 L=L_X^2L_Y^{-1},\qquad \mathcal E=i_*L
\]
と置くと、$p_*L\simeq\OO_X^2$, $q_*L\simeq\OO_Y^2$ であり、
ある二次元空間 $V_1,V_3$ に対して完全列
\begin{equation}\label{ruled-eq:matrix-resolution}
 0\longrightarrow\OO_{X\times Y}(-1,-1)\otimes V_3^\vee
 \xrightarrow{M}\OO_{X\times Y}\otimes V_1
 \longrightarrow\mathcal E\longrightarrow0
\end{equation}
が存在する。$\det M$ の零因子は $C$ である。
\end{lemma}

\begin{proof}
$L$ の次数は二であり、楕円曲線上で $h^0(L)=2$, $h^1(L)=0$ である。
$p_*L$ と $q_*L$ は階数二・次数零のベクトル束である。
さらに
\[
 H^0(X,p_*L(-1))=H^0(C,N)=0,
 \qquad H^0(Y,q_*L(-1))=H^0(C,N^2)=0.
\]
$\PP^1$ 上の分裂定理から両者は自明束となる。

次数二線束 $L$ は大域生成されるので、評価射
$V_1\otimes\OO_{X\times Y}\to i_*L$ は全射である。
ここで $V_1=H^0(C,L)$ とする。その核 $K$ は階数二の局所自由層である。
各 $p$-fiber または $q$-fiber への制限に Tor の項は生じない。
これは $i_*L$ がそれぞれの射影に関して平坦であるためである。
制限した評価射は大域切断上で同型であり、核は次数 $-2$、
$H^0=H^1=0$ の階数二束であるから、各 fiber 上で $\OO(-1)^2$ となる。
従って $K(1,0)$ は $Y$ からの引戻しであり、他方の fiber への制限から
$K\simeq\OO(-1,-1)^2$ を得る。
その係数空間を $V_3^\vee$ と取れば \eqref{ruled-eq:matrix-resolution} が得られる。
$L$ は滑らかな Cartier 因子 $C$ 上の線束なので、
この正方行列の行列式の因子は重複度一の $C$ である。
\end{proof}

$V_0=H^0(X,\OO_X(1))$, $V_2=H^0(Y,\OO_Y(1))$ と置く。
行列 \eqref{ruled-eq:matrix-resolution} は非零テンソル
\[
 w\in V_0\otimes V_1\otimes V_2\otimes V_3
\]
を定める。全ての $V_j$ は二次元である。
この記法では $x\in V_0^\vee$, $y\in V_2^\vee$ による評価は
\[
 M(x,y):V_3^\vee\longrightarrow V_1
\]
である。

\begin{proposition}\label{ruled-prop:geometric}
$(V_0,V_1,V_2,V_3,kw)$ は幾何的 quintuple である。
すなわち添字を $4$ を法として取り、任意の非零
$\lambda_j\in V_j^\vee$, $\lambda_{j+1}\in V_{j+1}^\vee$ に対して
$w$ の当該二因子による縮約は非零である。
\end{proposition}

\begin{proof}
$C$ は滑らかな既約 $(2,2)$ 曲線なので、両射影は有限である。
従って任意の $[x]\in X$ に対して $\det M(x,-)$ は恒等的には零でなく、
任意の $[y]\in Y$ に対しても $\det M(-,y)$ は恒等的には零でない。
隣接する四通りの縮約が零だと仮定したときの矛盾を記す。
\begin{itemize}
\item $(V_0,V_1)$：非零 $x,\beta$ による零縮約は
  $\beta M(x,y)=0$ を全ての $y$ について与える。
  よって $\det M(x,-)\equiv0$ となる。
\item $(V_1,V_2)$：非零 $\beta,y$ による零縮約は
  $\beta M(x,y)=0$ を全ての $x$ について与える。
  よって $\det M(-,y)\equiv0$ となる。
\item $(V_2,V_3)$：非零 $y,z$ による零縮約は
  $M(x,y)z=0$ を全ての $x$ について与える。
  よって $\det M(-,y)\equiv0$ となる。
\item $(V_3,V_0)$：非零 $z,x$ による零縮約は
  $M(x,y)z=0$ を全ての $y$ について与える。
  よって $\det M(x,-)\equiv0$ となる。
\end{itemize}
いずれも有限性に反する。
従って両 pushforward の自明性だけでなく、テンソル自体の全隣接縮約条件が成立する。
\end{proof}

\subsection{双対、正規化、および関係線束}

$X_{2r}=X$, $X_{2r+1}=Y$ とし、
\[
 \mathcal A=\Sym(\mathcal E)
\]
を標準的な非可換対称層 $\mathbb Z$-代数とする。
$\mathcal A_{i,i+1}=\mathcal E^{i*}$ は反復右双対であり、
二次関係は $X_i=X_{i+2}$ の対角上の $\OO_{X_i}$ である。
非可換対称代数の定義と twist については
\cite[\S5]{VanDenBergh2012P1} に従う。

$\mathcal E$ の右随伴双加群は、有限射の双対性から、$Y\times X$ 上で
\begin{equation}\label{ruled-eq:right-dual}
 \mathcal E^*=i^{\mathrm{op}}_*(L^{-1}L_Y^2)
\end{equation}
となる。実際、$q^!\OO_Y=\omega_C\otimes q^*\omega_Y^{-1}=L_Y^2$ である。
また
$\mathcal E^{**}=\OO_X(2)\otimes\mathcal E\otimes\OO_Y(-2)$
である \cite[Lemma 4.1.7]{VanDenBergh2012P1}。
従って $C$ 上の線束で表すと
\begin{equation}\label{ruled-eq:iterated-duals}
 \mathcal E^{2r*}=L N^{2r},\qquad
 \mathcal E^{(2r+1)*}=L^{-1}L_Y^2N^{-2r}.
\end{equation}
後者では support の二射影を逆順に読む。

各 $X_i$ 上の線束 $T_i=\OO_{X_i}(i)$ を用い、
\begin{equation}\label{ruled-eq:normalization}
 \mathcal A'_{ij}=T_i^{-1}\otimes\mathcal A_{ij}\otimes T_j
\end{equation}
と正規化する。ここで左・右の tensor は対応する射影からの引戻しを表す。
この twist は完全同値
\begin{equation}\label{ruled-eq:twist-equivalence}
 F:\Gr\mathcal A\xrightarrow{\sim}\Gr\mathcal A',\qquad
 F(P)_j=P_j\otimes T_j
\end{equation}
を与え、torsion 部分圏を保つ。

\begin{lemma}\label{ruled-lem:normalized}
正規化した隣接双加群は
\[
 \mathcal A'_{2r,2r+1}=i_*L_X^2,\qquad
 \mathcal A'_{2r+1,2r+2}=i^{\mathrm{op}}_*L_Y^2
\]
である。またその二次関係線束は
\[
 Q'_i=\OO_{X_i}(2)
\]
である。
\end{lemma}

\begin{proof}
偶数の場合、\eqref{ruled-eq:iterated-duals} と \eqref{ruled-eq:normalization} から
\[
 L_X^{-2r}(L N^{2r})L_Y^{2r+1}=L L_Y=L_X^2.
\]
奇数の場合も
\[
 L_Y^{-(2r+1)}(L^{-1}L_Y^2N^{-2r})L_X^{2r+2}
 =L^{-1}L_YL_X^2=L_Y^2.
\]
特に $i=1$ の twist は $\OO_Y(-1)\otimes\mathcal E^*\otimes\OO_X(2)$ であり、
単なる $\mathcal E^*\otimes\OO_X(1)$ ではない。
関係線束は元の $Q_i=\OO_{X_i}$ を twist した
$T_i^{-1}Q_iT_{i+2}=\OO_{X_i}(2)$ である。
\end{proof}

\subsection{対角列から見えない非零 coherent 対象}

有限二重被覆について
\begin{equation}\label{ruled-eq:double-cover-push}
 p_*\OO_C\simeq\OO_X\oplus\OO_X(-2),\qquad
 q_*\OO_C\simeq\OO_Y\oplus\OO_Y(-2)
\end{equation}
である。例えば $\OO(-2,-2)\to\OO\to i_*\OO_C$ を一方の射影で押し出すと
$0\to\OO\to p_*\OO_C\to\OO(-2)\to0$ を得る。
この拡大は $H^1(\PP^1,\OO(2))=0$ により分裂する。

\begin{proposition}\label{ruled-prop:invisible-module}
$\Gr\mathcal A'$ において
\[
 M_i=\begin{cases}
 \OO_{X_i}(-2)&i\ge0,\\
 0&i<0
 \end{cases}
\]
という次数成分を持つ noetherian 加群 $M$ が存在する。
その tails 対象 $\pi M$ は非零であり、全ての整数 $i$ に対して
\[
 (\omega\pi M)_i\simeq\OO_{X_i}(-2),\qquad
 \Hom_{\operatorname{QGr}\mathcal A'}(\pi e_i\mathcal A',\pi M)=0
\]
が成り立つ。
\end{proposition}

\begin{proof}
補題 \ref{ruled-lem:normalized} と \eqref{ruled-eq:double-cover-push} から、全ての $i\ge0$ で
\[
 M_i\otimes\mathcal A'_{i,i+1}
 \simeq\OO_{X_{i+1}}\oplus\OO_{X_{i+1}}(-2)
\]
となる。この直和の第二成分への全射を作用射
$\alpha_i:M_i\otimes\mathcal A'_{i,i+1}\to M_{i+1}$ とする。
各作用射は零でない定数倍を除いて一意である。

これらの射はまず tensor 層代数上の加群構造を定める。
二次関係の作用は、補題 \ref{ruled-lem:normalized} により
\[
 M_i\otimes Q'_i=\OO_{X_i}
 \longrightarrow M_{i+2}=\OO_{X_i}(-2)
\]
という層の射であるから零である。
$\mathcal A'$ の定義関係はこの二次関係により生成されるので、作用は
$\mathcal A'$ 上に降りる。
各 $\alpha_i$ は全射であり、双加群との tensor は右完全なので、
$M$ は次数零の $\OO(-2)$ により生成される。

$\mathcal E$ は左右局所自由・階数二である。
従って \cite[Theorem 1.2]{VanDenBergh2012P1} により $\mathcal A$ は noetherian であり、
twist 同値により $\mathcal A'$ も noetherian である。
$M$ は coherent 生成対象 $\OO(-2)\otimes e_0\mathcal A'$ の商なので noetherian である。
noetherian torsion 加群は上に有界であるが、$M_i\ne0$ は全ての $i\ge0$ で成立する。
従って $\pi M\ne0$ である。

隣接双加群を $E_i^\prime=\mathcal A'_{i,i+1}$ と略記する。
全ての $i\in\mathbb Z$ で $\widetilde M_i=\OO_{X_i}(-2)$ と置き、
\[
 \alpha_i:\widetilde M_i\otimes E_i^\prime
 \simeq\OO_{X_{i+1}}\oplus\OO_{X_{i+1}}(-2)
 \longrightarrow\OO_{X_{i+1}}(-2)
\]
を第二成分への非零定数倍射影とする。
関係の作用 $\widetilde M_i\otimes Q'_i=\OO\to\OO(-2)$ は零なので、
これは $\mathcal A'$ 加群である。$M=\widetilde M_{\ge0}$ であり、
$\widetilde M/M$ は上に有界だから torsion である。

次数 $i$ だけに $G$ を持つ加群を $S_i(G)$ と記す。
標準二次解消 \cite[Theorem 7.1.2]{VanDenBergh2012P1} を twist して、局所自由 $G$ に対し
\[
 0\to(G\otimes Q'_i)\otimes e_{i+2}\mathcal A'
 \to(G\otimes E_i^\prime)\otimes e_{i+1}\mathcal A'
 \to G\otimes e_i\mathcal A'\to S_i(G)\to0
\]
を得る。
各誘導関手は完全であり、次数評価との随伴および双加群の右随伴により、
$\mathbf R\Hom_{\Gr\mathcal A'}(S_i(G),\widetilde M)$ は
次の $X_i$ 上の複体 $\mathcal K_i^\bullet$ を $G^\vee$ と tensor したものの
導来大域切断で計算できる：
\[
 \mathcal K_i^\bullet=
 \left[
 \widetilde M_i\xrightarrow{\alpha_i^\flat}
 \widetilde M_{i+1}\otimes(E_i^\prime)^*
 \longrightarrow\widetilde M_{i+2}\otimes(Q'_i)^{-1}
 \right].
\]
ここでは誘導対象を射影対象と見なしていない。
誘導・次数評価・双加群随伴による導来 Hom と、上の有限解消の totalization を使っている。

関係の非退化性は
$E'_{i+1}\simeq(E_i^\prime)^*\otimes Q'_i$ を与える。
従ってこの複体の第二微分は
$\alpha_{i+1}\otimes(Q'_i)^{-1}$ であり、
\[
 \mathcal K_i^\bullet\simeq
 \left[
 \OO(-2)\longrightarrow\OO(-2)\oplus\OO(-4)
 \xrightarrow{(0,c_i)}\OO(-4)
 \right],\qquad c_i\in k^\times.
\]
第一微分は非零射 $\alpha_i$ の随伴であるから非零である。
$\Hom(\OO(-2),\OO(-4))=0$ なので、それは第一成分への非零定数倍包含である。
よって $\mathcal K_i^\bullet$ は分裂完全であり、
\[
 \mathbf R\Hom_{\Gr\mathcal A'}(S_i(G),\widetilde M)=0
\]
となる。

$\PP^1$ 上の任意の coherent $G$ は長さ一の有限局所自由解消を持つので、
この消滅は全ての coherent $G$ に成立する。
任意の noetherian torsion 加群は上と下に有界で、有限個の
$S_i(G)$ を successive quotient とする filtration を持つ。
従ってそのような $T$ に対して
\[
 \Hom(T,\widetilde M)=\Ext^1(T,\widetilde M)=0.
\]
$\Gr\mathcal A'$ は局所 Noetherian なので \cite[Theorem 1.2]{VanDenBergh2012P1}、
任意の torsion $T$ は noetherian 部分対象 $T_\lambda$ の filtered union である。
上の Hom 消滅から $\Hom(T,\widetilde M)=0$ である。
また $T$ の $\widetilde M$ による任意の拡大を各 $T_\lambda$ へ引き戻すと、
それぞれは一意に分裂する。これらの分裂は一意性によって整合するので、
filtered colimit を取れば元の拡大も分裂する。
従って $\Ext^1(T,\widetilde M)=0$ である。

これから $\widetilde M$ は飽和である。
実際、単位射 $\widetilde M\to\omega\pi\widetilde M$ の核と余核は torsion である。
Hom 消滅から核は零であり、Ext 一次消滅から余核による拡大は分裂する。
$\omega\pi\widetilde M$ は torsion-free なので、この余核も零である。
$\pi M=\pi\widetilde M$ により $\omega\pi M=\widetilde M$ を得る。
最後に随伴から
\[
 \Hom_{\operatorname{QGr}\mathcal A'}(\pi e_i\mathcal A',\pi M)
 \simeq H^0(X_i,(\omega\pi M)_i)
 =H^0(\PP^1,\OO(-2))=0
\]
となる。
\end{proof}

Nyman \cite[Theorem 2.5, Corollary 3.3]{Nyman2004} の一般理論からも、
Noetherian な層加群の高次数飽和が得られる。
ここではこの特定の加群について二次解消から直接計算し、
全整数での飽和と不可視性を示した。

\subsection{豊富性への帰結と適用範囲}

\begin{theorem}\label{ruled-thm:nonample}
$\qgr\mathcal A$ における列
\[
 K_n=\pi\bigl(\OO_{X_{-n}}(n)\otimes e_{-n}\mathcal A\bigr),\qquad n\in\mathbb Z,
\]
は ample sequence の生成条件、すなわち任意に大きい負添字からの生成条件を満たさず、
従って ample でない。偶数部分列も同じ生成条件を満たさない。
これは幾何的 quintuple を満たすテンソルから得られた反例である。
\end{theorem}

\begin{proof}
\eqref{ruled-eq:twist-equivalence} は $K_n$ を $\pi e_{-n}\mathcal A'$ に送る。
命題 \ref{ruled-prop:invisible-module} の非零対象 $\pi M$ に対し、
全ての $i$ で $\Hom(\pi e_i\mathcal A',\pi M)=0$ である。
従って全ての $i_j$ が十分大きいとき
$\bigoplus_j\pi e_{i_j}\mathcal A'\to\pi M$ という全射は存在しない。
これは \cite[\S8.1 の (A1)]{VanDenBergh2012P1} における ample sequence の生成条件を否定する。
偶数添字だけに制限しても同じである。
\end{proof}

\begin{remark}[正規化後の強豊富性との両立]
\cite[Lemma 8.1.1]{VanDenBergh2012P1} の仮定を正規化後の $\mathcal A'$ に対して確認できたとしても、得られる列は
\[
 \pi\bigl(\OO_{X_{-n}}(n)\otimes e_{-n}\mathcal A'\bigr)
\]
である。これを \eqref{ruled-eq:twist-equivalence} で戻すと
\[
 \pi\bigl(\OO_{X_{-n}}(2n)\otimes e_{-n}\mathcal A\bigr)
\]
となり、定理の $K_n$ とは異なる。
可換標準点 $\mathcal E=\OO_\Delta^2$ では、前者は
$\OO_{\PP^1\times\PP^1}(2n,n)$、$K_n$ は $\OO(n,n)$ に対応する。
従って正規化後に右大域生成性と高次コホモロジー消滅を証明しても、
元の対角偏極の豊富性は従わない。
\end{remark}

\paragraph{適用範囲.}
有限 foundation の Hom 代数が一致すること、あるいはその compact 導来圏が
Morita 同値であることから、指定された cubic 軌道列まで保つ
Noetherian heart の比較を結論してはならない。
本例はこの特定の対角列による全 geometric 領域での比較を妨げる。
別の列、別の heart、あるいは可換点の開近傍に限定した比較の可能性は残る。
また幾何的 quintuple に付随する cubic AS 正則 $\mathbb Z$-代数の
Noetherian 性そのものは既知の別の主張であり、ここでは否定していない。


% !TeX root = ../nc-quadric-moduli.tex
% 2026-11-21 / v1.33
\section{可換二次曲面を含む滑らかな Artin 開部分}
\label{sec:quadric-smooth-atlas}

\paragraph{範囲と依存関係.}
定理~\cite[\texttt{thm:quadric-finite-type-effectivity}]{QuadricIntrinsic2026} が与える実際の
$\mathfrak H_2$ 族と canonical compact 圏の比較から、
可換 quadric を含む有限型で滑らかな Artin 開部分を構成する。
最初に一般の効果的 foundation 族について、必要な入力を明記して表示定理を証明する。
$\mathfrak H_2$ 全体の代数性は結論しない。新規性は未確認である。

\subsection{入力の正確な形}

$k$ は標数零の代数閉体とする。
$\mathfrak H_2$ は局所有限型 $k$-スキームの big fppf site 上で定義された
pointed Abel 圏の族のスタックであり、既存の相対認識定理により
全忠実な射
\[
 \Phi:\mathfrak H_2\longrightarrow\mathfrak P
 :=[X_3/\operatorname{PGL}_9]
\]
を持つとする。$X_3$ は有限型分離 $k$ スキームであり、
$\mathfrak P$ の対角は affine かつ有限表示である。
したがって $\mathfrak H_2$ の Isom 層も同じ性質を持つ。
ここまでは既存の
\ref{prop:relative-h2-stack-recognition} と
\ref{prop:relative-h2-algebraicity-reduction} の結論である。

有限 foundation の標準表示空間を
\[
 S_0=\Hom_k(P\otimes W_1,W_2)
       \oplus\Hom_k(Q\otimes W_2,W_3)
       \simeq\mathbb A_k^{64}
\]
とする。各空間の次元は
$\dim P=\dim Q=\dim W_1=2$、$\dim W_2=4$、$\dim W_3=6$ である。
点 $(F,G)$ は四頂点の有向代数 $\Lambda$ を定め、
頂点 $1,2,3$ の部分代数では積
$\Lambda_{12}\otimes\Lambda_{23}\to\Lambda_{13}$ が同型である。
これらは既存の \cite[\texttt{prop:foundation-stack}]{QuadricIntrinsic2026} の表示である。
以下では $S\subset S_0$ を開部分とする。
一般の滑らかな foundation chart を用いる場合も、後述の基底の持上げを
その smooth chart の無限小持上げに置き換えれば同じ証明になる。

\begin{theorem}[効果的 foundation 族からの局所 atlas]
\label{thm:atlas-from-effective-foundation}
有限型スキーム $U$、\'etale 射 $a:U\to S$、および射
$q:U\to\mathfrak H_2$ が与えられているとする。
さらに $q$ の族の中に、元の site 内の任意の affine 基底変更と整合する対象の列
$(E_i)_{i\in\mathbb Z}$ があり、次の条件を満たすと仮定する。
\begin{enumerate}
 \item $E_0=\OO$ であり、各 $E_i$ は基底上平坦な noetherian 対象かつ
 canonical 導来圏の compact 対象である。
 \item 元の site 内の任意の affine 基底変更 $\operatorname{Spec} A\to U$ において
 $(E_{i,A})_{i\in\mathbb Z}$ は De Deken--Lowen の意味で geometric $(4,3)$-helix である。
 すなわち各四連続項は $D(\operatorname{Ind}\mathcal C_A)$ の compact generators である
 strong exceptional collection となり、前向きの導来 Hom は次数零の
 有限射影 $A$ 加群であり、helix の mutation 条件を満たす。
 \item $T_A=\bigoplus_{i=0}^3E_{i,A}$ に対して
 \[
  \End(T_A)\simeq a^*\Lambda\otimes_{\OO_U}A
 \]
 が、頂点と積を保ち基底変更と整合して与えられる。
 特に Hom の階数行列は
 \[
 \begin{pmatrix}
 1&2&4&6\\0&1&2&4\\0&0&1&2\\0&0&0&1
 \end{pmatrix}
 \]
 であり、$\Hom(E_1,E_2)\otimes\Hom(E_2,E_3)\to\Hom(E_1,E_3)$ は同型である。
\end{enumerate}
このとき $q$ はスキームで表現可能な smooth 射である。
その fppf 像 $\mathfrak V\subset\mathfrak H_2$ は開部分スタックであり、
有限型で滑らかな Artin スタックにより表現される。
具体的には
\[
 R=U\times_{\mathfrak H_2}U
   \simeq U\times_{\mathfrak P}U,
 \qquad
 \mathfrak V\simeq[U/R].
\]
$R$ は有限型スキームで、$R\to U$ の二つの射は smooth である。
$U$ に可換 quadric $Q$ を表す点があれば、$\mathfrak V$ は $Q$ の開近傍である。
\end{theorem}

この主張で $a$ が \emph{$U$ の全点で} \'etale であることは重要である。
可換点一つでの形式的変形関手の同型だけを仮定していない。

\subsection{指定された Artin 同型を延長する}

\begin{lemma}[foundation の持上げと指定同型]
\label{lem:atlas-artin-lift}
前定理の仮定の下で、$A'\twoheadrightarrow A$ を剰余体 $k$ を持つ
局所 Artin $k$ 代数の small extension とする。
$u_A:\operatorname{Spec} A\to U$、$c_{A'}\in\mathfrak H_2(A')$、
pointed exact 同型
\[
 \alpha_A:q(u_A)\xrightarrow{\sim}c_{A'}|_A
\]
が与えられれば、$u_A$ を延長する $u_{A'}:\operatorname{Spec} A'\to U$ と
$\alpha_A$ を延長する pointed exact 同型
\[
 \alpha_{A'}:q(u_{A'})\xrightarrow{\sim}c_{A'}
\]
が存在する。
\end{lemma}

\begin{proof}
\emph{(a) Abel 変形としての入力.}
$\mathcal G_{A'}=\operatorname{Ind}\mathcal C_{A'}$ と置く。
相対 $\mathfrak H_2$ 条件からこれは $A'$ 上 flat な Grothendieck 圏であり、
加群対象による基底変更により、その $A$ 上の還元は
$\mathcal G_A=\operatorname{Ind}\mathcal C_A$ である。
これは既存の \ref{lem:relative-h2-flat-center} と
相対 Abel 系の定義に含まれる。
したがって $\mathcal G_{A'}$ は $\mathcal G_A$ の flat Abel 変形である。
$A'\to A$ は coherent な可換環の冪零拡大であり、$A$ は $A'$ 上有限表示なので、
\cite[\S4.1]{DeDekenLowen2011} の基底条件を満たす。

\emph{(b) 対象と compact tilting 窓の持上げ.}
$\alpha_A$ により $E_{i,A}$ を $\mathcal C_A$ に移す。
これらは flat な対象であり、compact geometric helix を成す。
\cite[Theorem 5.13(1)--(4)]{DeDekenLowen2011} を基底環 $A$ と拡大 $A'\to A$ に適用する。
すると各 $E_{i,A}$ は、同型類一意に flat な対象 $E'_{i}$ へ持ち上がり、
指定した還元同型を選べる。持上げた四項窓は再び compact tilting generator であり、
前向き Hom は平坦で還元と可換する。
有限次元の剰余体上の Hom を持つ平坦 $A'$ 加群は有限自由なので、
\[
 \Lambda' =\End_{\mathcal G_{A'}}
       \left(\bigoplus_{i=0}^3 E'_i\right)
\]
は階数 $24$ の自由 $A'$ 代数となる。

この段階の対象の障害と選択を明記すると、$J=\ker(A'\to A)$ に対し
\[
 \mathbf R\Hom(E_{i,A},J\otimes_A^{\mathbf L}E_{i,A})
 \simeq J\otimes_A^{\mathbf L}\mathbf R\Hom(E_{i,A},E_{i,A})
 \simeq J[0]
\]
である。従って二次の障害も一次の持上げ差も消える。
同型そのものは一意ではない。
$E'_0$ は既存の標識 $\OO_{A'}$ と同型であり、
還元上の指定を延長する同型を選ぶ。
実際、自己同型群はそれぞれ $A'^\times$ と $A^\times$ であり、
冪零拡大に対して単元は持ち上がるので、任意の比較同型を所要の還元に補正できる。
この選択により以下の Morita 同値は pointed となる。

各 $E'_i$ が noetherian であることも確認できる。
$A'$ の極大イデアルによる有限フィルトレーションの商は、平坦性により
剰余体上の $E_i$ の有限直和である。
これらが noetherian なので $E'_i$ も noetherian である。

\emph{(c) foundation の座標の持上げ.}
頂点の恒等射で対角成分を $A'$ と同定する。
逆向き Hom は零であり、頂点 $1,2,3$ の積
\[
 \Hom(E'_1,E'_2)\otimes_{A'}\Hom(E'_2,E'_3)
       \longrightarrow\Hom(E'_1,E'_3)
\]
は還元上で同型なので Nakayama の補題により同型である。
五つの有限自由加群
\[
 \Hom(E'_0,E'_1),\quad\Hom(E'_0,E'_2),\quad\Hom(E'_0,E'_3),
 \quad\Hom(E'_1,E'_2),\quad\Hom(E'_2,E'_3)
\]
の基底を、$a\circ u_A$ の指定する基底を延長して選ぶ。
自由加群の基底は冪零拡大に対して持ち上がるので、ここに追加の障害はない。
残る二つの積の行列 $F',G'$ は $S_0(A')$ の点を与え、
$a\circ u_A$ を還元に持つ。
$S\subset S_0$ は開であり $\operatorname{Spec} A'$ の位相空間は一点なので、この点は $S(A')$ に入る。
得られた点を $s_{A'}$ と書く。
この構成は、代数の積を別途持ち上げる問題を解いているのではなく、
既に持ち上げた対象の実際の合成から積を取っている。
$\operatorname{HH}^3$ の消滅を新たな仮定に加える必要はない。

\emph{(d) \'Etale な根の chart への持上げ.}
$a:U\to S$ は \'etale なので
\[
 \begin{matrix}
 \operatorname{Spec} A&\longrightarrow&U\\
 \big\downarrow&&\big\downarrow a\\
 \operatorname{Spec} A'&\xrightarrow{s_{A'}}&S
 \end{matrix}
\]
には $u_A$ を延長する持上げ $u_{A'}$ が存在する。
これは全ての $U$ 上での \'etale 性を使う。
新しい閉点ごとに Serre 根の Hochschild 作用を再計算する必要はない。

\emph{(e) 指定された pointed exact 同型の持上げ.}
$q(u_{A'})$ と $c_{A'}$ の二つの compact tilting 窓の End 代数は、
上で選んだ基底により同じ $\Lambda'$ である。
標準 dg enhancement 上の導来 Morita によって、標識と四項窓を保つ同値
\[
 \beta_{A'}:
 D(\operatorname{Ind}\mathcal C_{q(u_{A'})})
 \xrightarrow{\sim}D(\operatorname{Ind}\mathcal C_{c_{A'}})
\]
を得る。その compact 部分の還元は、窓上の全ての射を含めて
$D(\alpha_A)$ と同型である。
実際、両者は同じ四つの projective 生成対象とその間の積に対応し、
選んだ基底の還元は $\alpha_A$ が誘導する比較と一致する。

この compact 同値は相対 Serre 双対を保つため、
$t=S^{-1}[2]$ と $\OO$ の全軌道 Hom、およびその積を保つ。
従って二つの相対反 Serre 環の同型、すなわち
\[
 \Phi(q(u_{A'}))\xrightarrow{\sim}\Phi(c_{A'})
 \quad\text{in }\mathfrak P(A')
\]
を与える。必要な次数線束の比較は既存の $\operatorname{PGL}_9$ による
正規化で吸収される。
$\Phi$ の全忠実性から、これは pointed exact 同型 $\alpha_{A'}$ に一意に持ち上がる。
還元が $\alpha_A$ になることも同じ全忠実性から従う。
ここでは単なる導来同値から heart の完全性を仮定していない。
二つの族が既に $\mathfrak H_2$ に属することと、相対反 Serre 再構成を用いている。
\end{proof}

\subsection{任意の有限型基底上の smoothness}

\begin{proof}[定理~\ref{thm:atlas-from-effective-foundation} の証明]
$T\to\mathfrak H_2$ を任意の局所有限型スキームからの射とする。
$\Phi$ の全忠実性により二ファイバー積の自然射は同値であり、
\[
 Z_T:=U\times_{\mathfrak H_2}T\simeq U\times_{\mathfrak P}T.
\]
右辺から $U\times_kT$ への射は $\mathfrak P$ の対角の基底変更であるから、
affine かつ有限表示である。従って $Z_T$ はスキームであり、
$Z_T\to T$ は局所有限表示である。
これで $q$ の表現可能性と有限表示性を、$\mathfrak H_2$ 自身の代数性を使わずに得る。

$T$ の有限型 affine 開部分ごとに議論すればよい。
このとき $Z_T$ も有限型 $k$ スキームなので、全ての閉点の剰余体は $k$ である。
$z\in Z_T(k)$ において small Artin extension $A'\twoheadrightarrow A$ と
可換図式
\[
 \begin{matrix}
 \operatorname{Spec} A&\longrightarrow&Z_T\\
 \big\downarrow&&\big\downarrow\\
 \operatorname{Spec} A'&\longrightarrow&T
 \end{matrix}
\]
を取る。上の点は、$u_A$、$T$ 側の族の還元、両者の指定された
pointed exact 同型を与える。
補題~\ref{lem:atlas-artin-lift} がこの図式の持上げを与える。
よって $Z_T\to T$ は各閉点で smooth である。
ここで用いているのは有限型 Noetherian 基底上の Artin 持上げ判定
\cite[Tag 02HY, Lemma 37.12.2]{StacksSmoothAtlas} である。
同補題が要求する $\beta:\operatorname{Spec} A'\to T$ の有限型性を仮定すると、
$A'$ は有限型 Artin $k$ 代数であり、その剰余体は $k$ である。
従って上の持上げは同補題の全ての試験図式を覆う。
非 smooth locus は閉であり、非空なら閉点を持つので、$Z_T\to T$ 全体が smooth となる。

以上が任意の $T$ に成り立つから、$q:U\to\mathfrak H_2$ は
表現可能な smooth 射である。
その fppf 像を $\mathfrak V$ とする。
各 $T\to\mathfrak H_2$ に対する $\mathfrak V\times_{\mathfrak H_2}T$ は、
smooth 射 $Z_T\to T$ の像に対応する開部分スキームである。
実際、この射はその開像上で smooth 全射であり、fppf 局所に section を持つ。
したがって $\mathfrak V\subset\mathfrak H_2$ は開部分スタックである。

$R=U\times_{\mathfrak H_2}U$ は上と同じ理由で有限型スキームであり、
二つの射 $R\rightrightarrows U$ は $q$ の基底変更なので smooth である。
よって $(U,R)$ は smooth groupoid であり、
\cite[Tags 04X0, 04TK]{StacksSmoothAtlas} により $[U/R]$ は Artin スタックとなる。
$q$ が $\mathfrak V$ を fppf に覆うことと Isom 層の一致から
$[U/R]\simeq\mathfrak V$ を得る。
$U$ は有限型であり、この smooth atlas により $\mathfrak V$ は有限型である。
さらに $U\to S$ は \'etale、$S$ は滑らかなので $U$ は $k$ 上 smooth であり、
従って $\mathfrak V$ も $k$ 上 smooth である。
\end{proof}

\begin{remark}[全 site への拡張の範囲]
上で構成する $[U/R]$ は通常の全スキーム上の Artin スタックである。
その局所有限型スキームの site への制限が $\mathfrak V$ を表す。
この構成のために $\mathfrak H_2$ 全体の任意基底への連続拡張や極限保存性を
先に証明する必要はない。
一方、$\mathfrak V$ の外側にある $\mathfrak H_2$ の対象については何も結論していない。
\end{remark}

\subsection{可換点近傍での仮定の確認}

\paragraph{\'Etale 性.}
既存の \cite[\texttt{prop:quadric-etale-morita-root}]{QuadricIntrinsic2026} は、平方根空間の写像が
可換点で \'etale であることを示した後、その \'etale locus に制限し、
可逆双加群の gerbe の \'etale 細分と比較線束の Zariski 自明化を行っている。
従ってそこで取る $U\to S$ は全点で \'etale である。
以後の $U^\circ\subset U$ への開制限もこの性質を保つ。

\paragraph{任意の閉ファイバーでの geometric helix.}
今回の cubic tails 比較は、$E_i=L_i$ と構造 row の対応を全整数 $i$ について与える。
すべての四項窓が compact tilting をなし、全ての $i\leq j$ に対し
\[
 \mathbf R\Hom(E_i,E_j)=\Hom(E_i,E_j)[0]
\]
となる。Serre と次数移動の比較は
\[
 E_{i+4}\simeq S^{-1}E_i[2]
\]
を与え、full exceptional 窓の Serre rotation は
\cite[Definition 5.12]{DeDekenLowen2011} の mutation 条件と一致する。
これらは $Q$ だけでなく、$U^\circ$ の全ての閉ファイバーで成立する。
元の site 内の任意の affine 基底変更でも同じ四項窓・Hom・完全列が存在するので、
補題~\ref{lem:atlas-artin-lift} は $A'\to A$ に直接適用できる。

\paragraph{平坦性と noetherian 性.}
各 row は次数成分ごとに基底上有限射影であり、
その像 $E_i$ への商関手は完全で tensor と可換する。
従って $E_i$ は基底上平坦である。
定理~\cite[\texttt{thm:quadric-finite-type-effectivity}]{QuadricIntrinsic2026} により
全有限型基底変更後の Noetherian 性が成立し、
これらは coherent heart の noetherian 対象である。
compact 性と unbounded 導来圏の生成は relative cubic tilting 定理による。
特異な基底において全ての coherent 対象が compact であるという主張は使わない。

\paragraph{豊富性と Z-generating 条件.}
全構造 Ext の公式と Noetherian heart の有限表示から、任意の coherent $M$ に対し
\[
 \Ext^p(E_i,M)=0\qquad(p>0,\ i\ll0)
\]
を得る。偶数部分 $E_{2n}=s^n\OO$ の AZ 豊富性は、列全体の
十分負の項の有限直和からの生成を与える。
Grothendieck 圏の任意の全射についても、$E_i$ からの一つの射の像は noetherian なので、
有限な noetherian 部分対象の全射へ帰着できる。
その核への上の Ext 一次消滅が tails-projectivity を与える。
付随する体上の cubic $\mathbb Z$ 代数は Noetherian であり、
$L_{\mathrm{tails}}=T_{\mathrm{tails}}$ である。
従って \cite[Corollary 3.16]{DeDekenLowen2011} の $\mathbb Z$-generating 条件を満たす。
前向き Hom の消滅は、有限次元 $k$ ベクトル空間 $V$ に対する
$\Ext^{1,2}(E_i,V\otimes_k E_j)=0$ $(i\leq j)$ も与える。
これで \cite[Theorems 4.5 and 5.15]{DeDekenLowen2011} による Abel 変形と
finite foundation 変形の比較も、全閉ファイバーで適用できる。

\paragraph{比較の対象.}
relation scheme $R=U^\circ\times_{\mathfrak P}U^\circ$ は
全ての pointed exact 同値を記録する。
foundation の選ばれた四対象を保つ線形群による商だけに置き換えていない。
とくに ruling の交換など、選ばれた窓を別の窓へ送る同値も
$\mathfrak P$ の Isom 層を通じて含まれる。

\begin{corollary}[可換 quadric の有限型 smooth 開近傍]
\label{cor:quadric-effective-smooth-neighborhood}
定理~\cite[\texttt{thm:quadric-finite-type-effectivity}]{QuadricIntrinsic2026} の
$q:U^\circ\to\mathfrak H_2$ に対し、$Q$ は $\mathfrak H_2$ の有限型で滑らかな Artin 開部分
$\mathfrak V_Q$ に属し、
\[
 U^\circ\longrightarrow\mathfrak V_Q,
 \qquad
 \mathfrak V_Q\simeq
 [U^\circ/(U^\circ\times_{\mathfrak P}U^\circ)]
\]
は smooth atlas とその groupoid 表示である。
\end{corollary}
\begin{proof}
定理~\cite[\texttt{thm:quadric-finite-type-effectivity}]{QuadricIntrinsic2026} が効果性を与え、
直前の各段落が定理~\ref{thm:atlas-from-effective-foundation} の仮定を確認する。
従って同定理を適用すればよい。
\end{proof}

\begin{remark}[代数性について残る範囲]
結論は $\mathfrak H_2$ 内部の開部分についてのものであり、
$\mathfrak P$ の開部分と同一視する主張ではない。
また全 $\mathfrak H_2$ の framed 実現部分層 $X_{\mathrm{as}}$ の表現可能性は、
$\mathfrak V_Q$ の外側を覆う同様の効果的 chart が得られるまで残る。
可換点だけの形式的比較、あるいは有限型の一族の存在だけから
この開性を主張してはならない。
本稿では任意の有限型 $T\to\mathfrak H_2$ に対する射 $Z_T\to T$ を実際に構成し、
その全閉点での Artin 持上げを確認したことが追加の根拠である。
\end{remark}


\begin{corollary}[可換点の局所次元]
\label{cor:quadric-open-stack-dimension}
上の $\mathfrak V_Q$ の可換点における接空間の同型類は九次元、
無限小自己同型は六次元であり、スタックの局所次元は三である。
必要なら $Q$ の連結開近傍に制限して、$\mathfrak V_Q$ を純次元三に取れる。
\end{corollary}
\begin{proof}
$\mathfrak V_Q\subset\mathfrak H_2$ が開であるため、
$Q$ を中心とする Artin 変形は両者で同じである。
系~\cite[\texttt{cor:quadric-formal-deformation-equivalence}]{QuadricIntrinsic2026} によれば、
これは pointed Abel 変形と同値であり、
接空間は $\operatorname{HH}^2(Q)$、無限小自己同型は
$\operatorname{HH}^1(Q)$ に対応する。その次元はそれぞれ $9,6$ である。
$\mathfrak V_Q$ は滑らかな Artin スタックなので、局所次元は $9-6=3$ となる。
滑らかなスタックの次元は局所定数なので、連結開近傍への制限で最後の主張を得る。
\end{proof}

% !TeX root = ../nc-quadric-surface-moduli.tex
% 2026-11-22 / v1.35（Keelerの刊行Corrigendumを明記）
\section{相対 cubic 環の普遍 Noetherian 性と有限型効果性}
\label{sec:relative-cubic-noetherian}

\paragraph{本稿の役割.}
全次数の局所自由性と標準 cubic 解消が既に構成された族に対し、
体上の Noetherian 性を各局所環へ持ち上げるだけでなく、
一つの Noetherian 基底上の全 projective row と column の Noetherian 性を導く。
周期性や Serre 根の選択は仮定せず、二周期の場合には通常の偶数環の左右 Noetherian 性も従う。
相対 point curve と正規族の構成を明記し、楕円部分の Noetherian 性には
Keeler の相対 Fujita 消滅を用いる。
linear 閉部分はテンソルの rank-one flattening により別に扱う。
最後に第\ref{sec:formal-noetherian-patching}節の形式的貼合せを適用する。
本稿の相対化および貼合せとの組合せの新規性は未確認である。

\subsection{仮定と結論}

$k$ を研究で固定した体、$R$ を可換 Noetherian $k$-代数、$S=\operatorname{Spec} R$ とする。
幾何ファイバーでは常に代数閉包へ移る。
正の $R$-線形 $\mathbb Z$-代数 $D$ について次を仮定する。
\begin{enumerate}
\item $D_{ii}=R$, $D_{ij}=0$ $(j<i)$ であり、全ての $D_{ij}$ は有限 projective $R$-加群である。
\item $V_i=D_{i,i+1}$ は階数二で、$D$ は $V_i$ により生成される。
  定義関係は階数二の直和因子
  $R_i\subset V_i\otimes V_{i+1}\otimes V_{i+2}$ により生成される。
\item 四次 syzygy は可逆加群 $W_i$ であり、普遍的に
  \[
  W_i\subset (R_i\otimes V_{i+3})\cap(V_i\otimes R_{i+1})
  \subset V_i\otimes\cdots\otimes V_{i+3}
  \]
  となる。縮約は同型
  \begin{equation}\label{noeth-eq:contraction}
  W_i\otimes V_{i+3}^{\vee}\simeq R_i,
  \qquad W_i\otimes V_i^{\vee}\simeq R_{i+1}
  \end{equation}
  を与える。
\item 全幾何ファイバーの cubic 代数は幾何的 quintuple に対応する AS 正則代数である。
  特に $W_i$ の任意の非零隣接二縮約は零でない。
\end{enumerate}
これらは既に得られた有限型 geometric cubic 開の相対 cubic 解消から供給する入力である。
\begin{theorem}[周期性を仮定しない普遍 Noetherian 性]
\label{thm:relative-geometric-cubic-noetherian}
上の仮定の下で、$D$ の全ての projective row と column は Noetherian である。
偶数添字に制限した $D^{\mathrm{ev}}_{ij}=D_{2i,2j}$ も同じ性質を持つ。
これらの結論は任意の可換 Noetherian $R$-代数 $R'$ への基底変更後にも成立する。
\end{theorem}

\begin{corollary}[二周期の偶数環]
\label{cor:relative-periodic-even-noetherian}
上の $D$ に積と整合する二周期同型 $\theta:D\xrightarrow{\sim}D(2)$ が指定されているとする。
\begin{equation}\label{noeth-eq:even-ring}
 B_n=D_{0,2n},\qquad
 a\star b=a\,\theta^m(b)\quad(a\in B_m, b\in B_n)
\end{equation}
と置くと、$B$ は左右 Noetherian である。
さらに任意の可換 Noetherian $R$-代数 $R'$ に対して
$B\otimes_RR'$ も左右 Noetherian である。
\end{corollary}

\begin{proof}
二周期同型は $\check B\simeq D^{\mathrm{ev}}$ を与える。
従って定理により $B$ の全ての斉次右イデアルと斉次左イデアルは有限生成である。
補題~\ref{lem:patch-graded-to-ungraded} により、次数を忘れても左右 Noetherian である。
任意の可換 Noetherian 基底変更後にも同じ議論が使える。
\end{proof}

以下ではまず nonlinear 部分で相対曲線と普遍正規族を構成し、
linear 閉部分と合わせて定理を証明する。
stalk ごとの Noetherian 性から開近傍上の Noetherian 性を推論する箇所はない。

\subsection{linear 閉部分と相対 point curve}

\subsubsection{閉部分の定義}

局所的に $W_0$ の生成元を $w$ とし、
\[
 w\in (V_0\otimes V_2)\otimes(V_1\otimes V_3)
\]
を $02\mid13$ flattening の $4\times4$ 行列として読む。
その $2\times2$ 小行列式が生成するイデアルを $I\subset R$ とする。
基底と $W_0$ の生成元の変更はこのイデアルを変えない。
\[
 Z=V(I),\qquad S_{\mathrm e}=S\setminus Z
\]
と置く。Van den Bergh の linear tensor の標準形
\cite[Lemma 5.2.4, (5.3)]{VdB2011} により、
$Z$ の幾何点はちょうど linear quintuple の点である。
実際、その標準形は
\[
 (x_0y_2-y_0x_2)(x_1y_3-y_1x_3)
\]
である。

これは reduced locus の定義だけではない。
$02\mid13$ の rank-one determinantal scheme は Segre scheme そのものであり、
$W_0$ が至る所非零であるため、$Z$ 上では可逆部分束
$A_0\subset V_0\otimes V_2$ と $A_1\subset V_1\otimes V_3$ が存在して
$W_0=A_0\otimes A_1$ となる。
局所的には一つの可逆な成分と小行列式の方程式を用いて
\begin{equation}\label{noeth-eq:linear-factorization}
 w=a_{02}\otimes b_{13}
\end{equation}
と因数分解でき、その二つの像線は一意なので重なりで貼り合わさる。
両双線形形式は非退化である。
もし一方がある幾何点で退化すれば、その零方向と隣接因子の任意の非零方向による
縮約が零となり、geometric 条件に反する。
従ってそれらの行列式は可逆である。
この議論は $Z$ が非被約の場合も適用できる。

\subsubsection{楕円部分の行列式}

以下しばらく $S_{\mathrm e}$ の任意の affine open を $S$ と書き直す。
$P_i=\PP(V_i^{\vee})$ とし、$P_i\times_S P_{i+1}$ 上で $W_i$ を最初の二因子により評価すると
\begin{equation}\label{noeth-eq:relation-matrix}
 \mu_i:W_i\otimes V_{i+3}^{\vee}\otimes\OO(-1,-1)
 \longrightarrow V_{i+2}\otimes\OO
\end{equation}
を得る。
geometric 条件により、これは全ての幾何点で零行列でない。
また nonlinear 条件と \cite[\S5.4, 特に Lemma 5.4.10]{VdB2011} により、
その行列式は各基底ファイバー上で恒等的に零ではない。
従って
\[
 C_i=V(\det\mu_i)\subset P_i\times_S P_{i+1}
\]
は相対有効 Cartier 因子であり、$S$ 上平坦な $(2,2)$ 曲線族である。
各幾何ファイバーは算術種数一の Gorenstein 曲線である。
ここで非被約性や可約性を除外しない。

\begin{lemma}\label{noeth-lem:scheme-shift}
$R_i$ の二つの multilinear relation により定義される
\[
 \Gamma_i\subset P_i\times_S P_{i+1}\times_S P_{i+2}
\]
の二つの隣接射影は同型
\[
 \Gamma_i\xrightarrow{\sim}C_i,
 \qquad \Gamma_i\xrightarrow{\sim}C_{i+1}
\]
を与える。従って canonical な同型 $\psi_i:C_i\xrightarrow{\sim}C_{i+1}$ がある。
これらは任意の基底変更と可換する。
\end{lemma}

\begin{proof}
最初の二座標を固定すると、$\Gamma_i$ の方程式は
$\lambda_{i+2}\mu_i=0$ である。
$\mu_i$ は至る所 rank $\ge1$ で、$C_i$ 上では行列式が零である。
ある成分が可逆である開集合で行基本変形・列基本変形を行えば、
$C_i$ 上の行列は $\operatorname{diag}(1,0)$ となる。
従ってその cokernel は線束であり、$\lambda_{i+2}$ は一意な projective 点を定める。
これは非被約基底上でも scheme-theoretic な同型 $\Gamma_i\simeq C_i$ を与える。

重なる二座標に対して同じ行列式を用いている点を具体的に確認する。
$\lambda_{i+1},\lambda_{i+2}$ を固定し、
\[
 M=W_i(-,\lambda_{i+1},\lambda_{i+2},-)
 \in V_i\otimes V_{i+3}
\]
を $2\times2$ 行列として読む。
$W_i\otimes V_{i+3}^{\vee}\simeq R_i$ によれば、$\Gamma_i$ の方程式は
$\lambda_iM=0$ である。
$W_i\otimes V_i^{\vee}\simeq R_{i+1}$ によれば、
$\Gamma_{i+1}$ の方程式は $M\lambda_{i+3}=0$ である。
従って両 triple のこの隣接 pair への像は、いずれも厳密に $\det M=0$ である。
別の $W_{i+1}$ による $R_{i+1}$ の表示を用いても、同じ relation 束の基底変更に過ぎず、
行列式のイデアルは変わらない。
geometric 条件により $M$ は零行列にならないので、
左 kernel と右 kernel はそれぞれ線束を与える。
これで第二の同型も scheme-theoretic に得られる。
全てが可逆成分による行列消去で記述されるため、任意の基底変更と整合する。
\end{proof}

\begin{remark}[二種類の $(2,2)$ 行列式を区別する]
ここで使う $C_i$ は\emph{隣接因子} $(i,i+1)$ 上の point curve である。
その各 $\PP^1$ への射影が有限であることは仮定せず、一般に成分を縮約し得る。
必要なのは隣接二射影を合わせた閉埋め込みと、$\OO(1,1)|_{C_i}$ の相対豊富性である。
これに対し、層双加群を作る際の\emph{反対因子} $(0,2)$ 上の行列式は
geometric 条件により両射影有限であるが、上の point curve と一様には同じでない。
特に linear 点では point scheme は曲面 $\PP^1\times\PP^1$ となり、
隣接行列式は恒等的に零となる。
従って本節の flat curve を linear 点までそのまま延長してはいけない。
\end{remark}

\subsection{相対 helix 座標代数と普遍正規族}

補題 \ref{noeth-lem:scheme-shift} により全 $C_i$ を $C=C_0$ と同定する。
その座標射を $\rho_i:C\to P_i$ とし、
\[
 \pi:C\to S,\qquad \mathcal L_i=\rho_i^*\OO_{P_i}(1),\qquad
 \mathcal L_{ij}=\mathcal L_i\otimes\cdots\otimes\mathcal L_{j-1}\quad(j>i)
\]
と置く。
幾何ファイバーではこれは Van den Bergh の admissible quadruple の elliptic helix そのものである
\cite[\S5.4, Lemma 5.5.1, Theorem 5.5.10]{VdB2011}。
従って各ファイバーの各既約成分 $E$ について
\begin{equation}\label{noeth-eq:parity-degrees}
 \deg(\mathcal L_{i+2}|_E)=\deg(\mathcal L_i|_E)
\end{equation}
である。
この数値等式は非被約ファイバーでもその reduced 成分上で読む。
本稿は相対線束としての helix 関係を fiberwise 同型から推論しない。
必要なのは \eqref{noeth-eq:parity-degrees} のみである。

\begin{lemma}\label{noeth-lem:cohomology-bc}
$j>i$ に対し
\[
 H_{ij}:=\pi_*\mathcal L_{ij}
\]
は階数 $2(j-i)$ の有限 locally free 層であり、
$R^{q}\pi_*\mathcal L_{ij}=0$ $(q>0)$ である。
この pushforward は任意の基底変更と可換する。
また $V_i\to\pi_*\mathcal L_i$ は同型で、$\mathcal L_i$ は相対大域生成される。
\end{lemma}

\begin{proof}
全幾何ファイバーで $\mathcal L_i$ は大域生成される次数二線束である。
積 $\mathcal L_{ij}$ も大域生成され、その次数は $2(j-i)>0$ である。
$(2,2)$ 曲線の tame 性と Riemann--Roch により
$H^1(\mathcal L_{ij,s})=0$, $h^0(\mathcal L_{ij,s})=2(j-i)$ である
\cite[Proposition 5.4.2, Lemma 5.5.3]{VdB2011}。
$C\to S$ は平坦・射影・有限表示で、$\mathcal L_{ij}$ は線束なので、
cohomology and base change により結論を得る。
同値に $R\pi_*\mathcal L_{ij}$ は perfect で、その全幾何ファイバーが次数零のみを持つため、
有限 locally free 層として表される。導来基底変更は通常の基底変更となる。
$V_i\to\pi_*\mathcal L_i$ は全ファイバーで同型であり、両辺が有限 locally free なので同型である。
評価射の全射性も幾何ファイバーで確認できる。
\end{proof}

$H_{ii}=R$, $H_{ij}=0$ $(j<i)$ とし、線束の積により $H$ を $\mathbb Z$-代数とする。
$\pi_*\OO_C=R$ は $(2,2)$ 因子の標準完全列から直接従う。
各 $R_i$ は $C$ 上で消えるので、canonical な準同型
\[
 D\longrightarrow H
\]
が得られる。

\begin{proposition}\label{noeth-prop:normal-family}
上の準同型は全成分で全射であり、核を $K$ とすると全 $K_{ij}$ は有限 projective である。
$K_{ij}=0$ $(j-i<4)$、$G_i:=K_{i,i+4}$ は可逆加群である。
さらに積は普遍同型
\begin{equation}\label{noeth-eq:normal-family}
 G_i\otimes_R D_{i+4,j}\xrightarrow{\sim}K_{ij}
 \xleftarrow{\sim}D_{i,j-4}\otimes_R G_{j-4}
 \qquad(j\ge i+4)
\end{equation}
を与える。従ってこれは左右 regular な正規族である。
構成、核、同型の全てが任意の基底変更と可換する。
\end{proposition}

\begin{proof}
幾何ファイバーでは \cite[Theorem 5.5.6(1)]{VdB2011} により $D\to H$ は全射である。
各成分の cokernel は有限 $R$-加群なので、全成分で全射である。
$D_{ij}$ と $H_{ij}$ は有限 projective であり、この全射は $R$-加群として分裂する。
従って $K_{ij}$ は有限 projective で、核は任意の基底変更と可換する。
低次数での消滅と $G_i$ の階数一性も同じ定理から従う。

式 \eqref{noeth-eq:normal-family} の射は $K$ が両側イデアルであることから定まる。
各幾何ファイバーで、左・右の核生成性と regular 性
\cite[Theorem 5.5.6(1),(3)]{VdB2011} により、両射は同型である。
両辺は有限 projective なので射自体が同型である。
この同型の tensor は任意の基底変更後も同型である。
\end{proof}

特に両側の表示を比較すると canonical な同型
\begin{equation}\label{noeth-eq:normal-commutation}
 G_i\otimes D_{i+4,j+4}\simeq D_{ij}\otimes G_j
\end{equation}
を得る。これが基底上の線束を保った正規交換式である。

\subsubsection{二周期性と通常の次数付き正規線}

この小節だけ、系~\ref{cor:relative-periodic-even-noetherian} の二周期同型 $\theta$ を仮定する。
相対 point scheme の構成は $D$ とその添字の同型に関して functorial である。
従って $\theta$ は $H$, $K$, $G_i$ にも二周期同型を誘導し、
$D\to H$ および \eqref{noeth-eq:normal-family} と整合する。
具体的には、$\theta$ が与える $C_2\to C_0$ と
$\psi_1\psi_0:C_0\to C_2$ の合成を $\sigma:C\to C$ とすれば、
\[
 \sigma^*\mathcal L_i\simeq\mathcal L_{i+2}
\]
となる。これは座標の tautological line bundle と $\theta$ の可換図式から得られる。

$\overline B_n=H_{0,2n}$ に \eqref{noeth-eq:even-ring} と同じ二周期積を入れると、
$\overline B$ は $\mathcal L_0\mathcal L_1$ と $\sigma$ による twisted homogeneous coordinate ring である。
また $G=G_0\subset B_2$ とすると
\begin{equation}\label{noeth-eq:graded-normal-line}
 \ker(B\to\overline B)=GB=BG,
 \qquad G\otimes B(-2)\simeq\ker(B\to\overline B)\simeq B(-2)\otimes G.
\end{equation}
これは $G_i$ の二周期同一視を \eqref{noeth-eq:normal-family} に適用したものである。
$G$ を自明化する開集合では、その生成元 $g\in B_2$ は左右 regular な正規元である。
全ての主張は任意の基底変更後も成立する。

\subsection{相対豊富性と座標代数の Noetherian 性}

\subsubsection{Keeler の消滅の適用}

用いる一次文献の正確な形は次である。
Keeler \cite[Theorem 1.5]{Keeler2003} は、可換 Noetherian 環上の射影スキーム、
豊富な線束 $H$、任意の coherent 層 $F$ に対し、ある $m_0$ が存在して
\begin{equation}\label{noeth-eq:keeler}
 H^q(F\otimes H^m\otimes N)=0
 \quad(q>0, m\ge m_0)
\end{equation}
が全ての relatively nef 線束 $N$ について成立することを与える。
arXiv v5 に付属する erratum と、2018年の正式な
Corrigendum \cite{Keeler2018Corrigendum} は原著の Lemma 3.3 を訂正し、
その定理への依存関係を再検証している。
従ってこの節で用いるのは訂正後にも有効な
Theorem 1.5 \cite[付属 Erratum, \S2]{Keeler2003} であり、
誤った旧 Lemma 3.3 の形は引用しない。

固定した $i$ に対し $H_i=\mathcal L_i\mathcal L_{i+1}$ は相対豊富である。
実際、これは隣接二座標による閉埋め込みの $\OO(1,1)$ である。
式 \eqref{noeth-eq:parity-degrees} から
\begin{equation}\label{noeth-eq:nef-prefix}
 \mathcal L_{i,i+2m}\otimes H_i^{-m}
\end{equation}
は各幾何ファイバーの各既約成分上で次数零であり、relatively numerically trivial、
従って relatively nef である。
奇数長の積ではさらに大域生成線束一つを tensor するので、残余因子はやはり nef である。
逆方向の積 $\mathcal L_{i-2m,i}$ についても同じ議論が働く。
従って \eqref{noeth-eq:keeler} により、前向き・後向き双方の積は ample sequence を成す。
これは全ての affine Noetherian 基底変更後にも適用できる。

\subsubsection{row ideal の有限生成を直接証明する}

\begin{lemma}\label{noeth-lem:row-noetherian}
$H$ は左右 Noetherian な $R$-線形 $\mathbb Z$-代数である。
偶数添字に制限した $H^{\mathrm{ev}}$ も左右 Noetherian である。
二周期同型が指定されている場合、その二周期環 $\overline B$ も左右 Noetherian である。
\end{lemma}

\begin{proof}
$\mathcal P_0=\OO_C$ とし、$\mathcal P_{j+1}=\mathcal P_j\mathcal L_j$ により全整数 $j$ の線束を定める。
すると $H_{ij}=H^0(C,\mathcal P_i^{-1}\mathcal P_j)$ $(j\ge i)$ である。
右 projective row $e_iH$ の任意の部分加群 $J$ を取り、
\[
 \mathcal I_j=\operatorname{im}\bigl(J_j\otimes_R\mathcal P_j^{-1}\longrightarrow\mathcal P_i^{-1}\bigr)
 \qquad(j\ge i)
\]
と置く。
$H_{ij}$ は有限 $R$-加群で $R$ は Noetherian なので、$J_j$ も有限である。
従って $\mathcal I_j$ は coherent である。
$\mathcal L_j$ は大域生成され、$J_jH_{j,j+1}\subset J_{j+1}$ だから
$\mathcal I_j\subset\mathcal I_{j+1}$ となる。
$\mathcal P_i^{-1}$ は Noetherian coherent 層なので、この列は $j\ge N$ で
一つの $\mathcal I$ に安定する。

$J_N$ の有限個の $R$-生成元を取れば全射
\[
 (\mathcal P_N^{-1})^{\oplus r}\twoheadrightarrow\mathcal I
\]
が得られる。その核 $F$ は coherent である。
前向きの ample sequence により $H^1(C,F\mathcal P_j)=0$ $(j\gg0)$ である。
従って $J_N$ の選んだ生成元による積は
\[
 H_{Nj}^{\oplus r}\twoheadrightarrow H^0(C,\mathcal I\mathcal P_j)
\]
を与える。
一方、像層の定義から $J_j\subset H^0(C,\mathcal I\mathcal P_j)$ であり、
積の像はイデアル性から $J_j$ に含まれる。
従って十分大きい $j$ では $J_j=J_NH_{Nj}$ である。
残る有限個の次数の有限 $R$-生成元を加えれば $J$ は有限生成となる。

左側は $He_i$ の部分加群 $J$ に対して
\[
 \mathcal I_j=\operatorname{im}(J_j\otimes_R\mathcal P_j\longrightarrow\mathcal P_i),\qquad j\le i
\]
を使う。$j\to-\infty$ に向かって像層が増大し、安定後は
後向きの ample sequence $\mathcal P_j^{-1}$ を用いて同じ議論を適用する。
従って全 projective row と column が Noetherian である。

偶数添字だけに制限すると、隣接線束は $\mathcal L_{2j}\mathcal L_{2j+1}$ に置き換わる。
これは大域生成され、前向き・後向きの豊富性も保たれるので、同じ証明がそのまま働く。
二周期同型が指定されていれば、$\Gr\overline B$ と $\Gr H^{\mathrm{ev}}$ は同値であるため、
$\overline B$ の全ての斉次右・左イデアルが有限生成となる。
最後に非負内部次数の最高次部分による通常の reduction により、
斉次イデアルの Noetherian 性から通常の左右 Noetherian 性が従う。
\end{proof}

\subsubsection{正規線からの持上げ}

\begin{lemma}\label{noeth-lem:normal-lift}
$A$ を非負次数付き $R$-代数とし、正次数にある可逆 $R$-部分加群 $G\subset A_d$ が
左右 regular な正規線を成すとする。
すなわち $GA=AG$ であり、積が
$G\otimes A(-d)\simeq GA\simeq A(-d)\otimes G$ を与えるとする。
$A/GA$ が左右 Noetherian なら $A$ も左右 Noetherian である。
\end{lemma}

\begin{proof}
$G$ を自明化する有限 affine open cover に制限すれば、
一つの左右 regular な正規元 $g$ を取れる。
$gA=Ag$ と両側 regular 性から、$ga=\alpha(a)g$ を満たす $\alpha(a)$ は一意に存在する。
逆向きの交換式も一意に解けるので $\alpha$ は次数を保つ自己同型である。
この自己同型を用いて
$\operatorname{gr}_{(g)}A\simeq(A/gA)[z;\alpha]$ であり、これは左右 Noetherian である。
任意の斉次イデアルの $(g)$-進初項イデアルから有限個の斉次生成元を lift する。
固定した内部次数 $n$ では $(g)^rA_n=0$ $(rd>n)$ なので、
初項を順に消す操作は有限回で止まり、元の斉次イデアルも有限生成となる。
さらに最高内部次数による reduction により通常のイデアルも有限生成となる。
有限 affine open cover の積は $R$ 上忠実平坦であるから、Noetherian 性は $A$ へ降下する。
\end{proof}

命題~\ref{noeth-prop:normal-family} の正則正規族と
補題~\ref{noeth-lem:row-noetherian} に、周期性を仮定しない
定理~\ref{thm:patch-normal-family} を $\delta=4$ として適用する。
これにより $S_{\mathrm e}$ 上の $D$ の全ての row と column は Noetherian である。
無限個の正規線 $G_i$ は同定や自明化を仮定せず、同定理の頂点線束 twist で扱う。
同じ構成と Keeler の定理は任意の可換 Noetherian 基底変更後に適用できるので、
楕円 open 上で普遍 Noetherian 性が得られる。
補題~\ref{noeth-lem:normal-lift} の通常の次数付き環としての形は、後で AS4 被覆に用いる。

\subsection{linear 閉部分}

ここでは再び元の $S$ に戻り、$Z=V(I)$ 上の代数を考える。
この小節では $R/I$ を $R$ と書き、全ての cubic データを $Z$ 上へ制限する。
式~\eqref{noeth-eq:linear-factorization} の二つの非退化線束を $A_0,A_1$ と書く。
$V_3^\vee$ による縮約から $R_0=A_0\otimes V_1$ を得る。
以下では二周期同型を用いず、普遍的な縮約同型~\eqref{noeth-eq:contraction} により
この表示を全整数へ伝播し、線束の twist を明示する。
添字順と異なる順に書いた tensor は、元の添字順への shuffle を介して読む。

\subsubsection{linear relation の両方向への伝播}

\begin{lemma}[一つの線型関係からの伝播]
\label{lem:linear-propagation-without-period}
ある $i$ において非退化な可逆部分束
$A_i\subset V_i\otimes_R V_{i+2}$ があり、
添字順への shuffle を介して
\[
 R_i=A_i\otimes_R V_{i+1}
\]
であるとする。このとき、一意な可逆部分束
$A_j\subset V_j\otimes_R V_{j+2}$ $(j\in\mathbb Z)$ が存在し、全て非退化で、
\begin{equation}\label{eq:linear-all-relations}
 R_j=A_j\otimes_R V_{j+1},\qquad
 W_j=A_j\otimes_R A_{j+1}
\end{equation}
となる。この構成は任意の基底変更と可換する。
ここで非退化とは、縮約
$A_j\otimes_R V_j^\vee\to V_{j+2}$ が同型であることをいう。
\end{lemma}

\begin{proof}
まず右向きに進める。
$W_i\subset R_i\otimes V_{i+3}$ と仮定の表示から、因子を並べ替えて
\[
 W_i\subset A_i\otimes(V_{i+1}\otimes V_{i+3})
\]
を得る。$A_i$ は可逆なので、これを tensor で外して
\[
 A_{i+1}:=A_i^{-1}\otimes W_i
       \longrightarrow V_{i+1}\otimes V_{i+3}
\]
と定める。
各幾何ファイバーで $W_i$ は非零なので、この射は各点で階数一である。
局所基底で少なくとも一つの成分が単元になることから、その像は可逆部分束である。
定義により $W_i=A_i\otimes A_{i+1}$ は部分束としての等式である。

$A_{i+1}$ がある幾何点で退化すると、ある非零の
$\lambda\in V_{i+1}^\vee$ が $A_{i+1}$ を零に縮約する。
任意の非零 $\mu\in V_i^\vee$ と合わせると $W_i$ の隣接二縮約も零になり、
geometric 条件に反する。
従って $A_{i+1}$ の縮約行列の行列式は至る所可逆であり、非退化である。
第二の同型 \eqref{noeth-eq:contraction} により
\[
 R_{i+1}
   =\operatorname{im}(W_i\otimes V_i^\vee)
   =A_{i+1}\otimes\operatorname{im}(A_i\otimes V_i^\vee)
   =A_{i+1}\otimes V_{i+2}.
\]
これで右向きの帰納法が進む。

左向きには $W_{i-1}\subset V_{i-1}\otimes R_i$ から
\[
 W_{i-1}\subset(V_{i-1}\otimes V_{i+1})\otimes A_i
\]
を得て、$A_{i-1}=W_{i-1}\otimes A_i^{-1}$ を定める。
同じ非退化性の議論と、$V_{i+2}^\vee$ による縮約を用いると
$R_{i-1}=A_{i-1}\otimes V_i$ を得る。
これで全整数に対する帰納法が完了する。
各段階は部分束の包含と可逆加群の tensor、普遍的な縮約同型だけで定まるので、
一意性と任意基底変更との整合性も従う。
\end{proof}

\subsubsection{係数線束を吸収する頂点 twist}

\begin{theorem}[周期を用いない線型標準形]
\label{thm:linear-normalization-without-period}
補題~\ref{lem:linear-propagation-without-period} の設定で、$D$ を $V_i$ により生成され、三次関係が $R_i$ である
正の向きの $R$ 線形 $\mathbb Z$ 代数とする。
$E=V_0$、$F=V_1$ と置く。
このとき頂点ごとの可逆加群 $N_i$ が存在し、
\[
 D'_{ij}=N_i^{-1}\otimes_R D_{ij}\otimes_R N_j
\]
は、偶数位置の矢空間が $E$、奇数位置の矢空間が $F$、
三次関係が同じ parity の二文字の交換関係である標準 linear $\mathbb Z$ 代数となる。
特に $n\geq0$ に対し
\begin{equation}\label{eq:linear-symmetric-components}
 D'_{i,i+n}\simeq
 \Sym^{\lceil n/2\rceil}E_{\bar i}
       \otimes_R\Sym^{\lfloor n/2\rfloor}E_{1-\bar i},
 \qquad E_0=E,\ E_1=F.
\end{equation}
全ての同定は積と整合する。
\end{theorem}

\begin{proof}
非退化線 $A_i\subset V_i\otimes V_{i+2}$ は同型
\[
 V_{i+2}\simeq V_i^\vee\otimes A_i
 \simeq V_i\otimes(\det V_i)^{-1}\otimes A_i
\]
を与える。二番目は階数二の標準的な外積による双対同定である。
この同定で $A_i$ の包含は、第一と第三の因子の標準交代テンソルに対応する。

$M_0=M_1=R$ とし、parity ごとに両方向へ再帰的に可逆加群 $M_i$ と同型
\[
 V_i\simeq E_{\bar i}\otimes M_i
\]
を選ぶ。具体的な再帰式は
\begin{equation}\label{eq:linear-arrow-lines}
 M_{i+2}=A_i\otimes(\det E_{\bar i})^{-1}\otimes M_i^{-1}.
\end{equation}
これにより関係線は
\[
 A_i\simeq\det(E_{\bar i})\otimes M_i\otimes M_{i+2}
\]
となり、その包含は交代テンソルの包含に一致する。
負の添字では \eqref{eq:linear-arrow-lines} を逆に解けばよい。
各 $M_i$ は有限回の tensor と双対で作られる可逆加群である。

次に $N_0=R$ として、全整数に対し
\begin{equation}\label{eq:linear-vertex-lines}
 N_{i+1}=N_i\otimes M_i^{-1}
\end{equation}
と定める。すると
\[
 N_i^{-1}\otimes V_i\otimes N_{i+1}\simeq E_{\bar i}.
\]
三次関係を同じ方法で移すと、係数線は
\[
 N_i^{-1}\otimes M_i\otimes M_{i+1}\otimes M_{i+2}\otimes N_{i+3}
 \simeq R
\]
となり、全て消える。
従って偶数 $i$ の関係は $\bigwedge^2E\otimes F$ を
$E\otimes F\otimes E$ へ shuffle して入れたもの、
奇数 $i$ の関係は $\bigwedge^2F\otimes E$ を
$F\otimes E\otimes F$ へ入れたものである。

この交換関係は path の偶数位置の文字同士と奇数位置の文字同士を独立に交換する。
同じ parity の隣り合う二文字は path の長さ三の部分に現れるため、
全ての並べ替えは与えられた三次関係から得られる。
従って tensor 代数の商は、二種類の対称積の tensor 積にちょうど一致し、
\eqref{eq:linear-symmetric-components} を得る。
この同定の下で path の連結は二重対称代数の積になる。
局所基底を使って確認する場合も、$E,F$ を自明化する有限個の開部分だけで足り、
全ての $A_i,M_i,N_i$ を自明化する必要はない。
\end{proof}

\subsubsection{各 row と column の普遍 Noetherian 性}

\begin{corollary}\label{cor:linear-universal-row-noetherian}
前定理で $R$ が Noetherian なら、$D$ の全ての row と column は Noetherian である。
同じ結論は任意の可換 Noetherian $R$ 代数への基底変更後にも成り立つ。
\end{corollary}

\begin{proof}
頂点線束 twist は加群圏の完全同値であり、各 row を基底の可逆加群による
tensor を除いて対応させる。従って $D'$ について証明すればよい。
可換 Segre 代数
\[
 T=\bigoplus_{n\geq0}\Sym^nE\otimes_R\Sym^nF
\]
は有限加群 $E\otimes F$ により生成される $R$ 代数であるから、Noetherian である。
また次数付き $T$ 加群
\[
 M_E=\bigoplus_{n\geq0}\Sym^{n+1}E\otimes\Sym^nF,
 \qquad
 M_F=\bigoplus_{n\geq0}\Sym^nE\otimes\Sym^{n+1}F
\]
は、それぞれ次数零の有限加群 $E$ と $F$ で生成されるので Noetherian である。

固定した $P'_i=e_iD'$ を終点の parity に分ける。
同じ parity の成分は $T$ のシフトであり、反対 parity の成分は
$M_E$ または $M_F$ のシフトである。
各 parity を保つ二段ずつの作用は、上で定義した $T$ の作用そのものである。
$P'_i$ の部分加群の任意の昇鎖は、両 parity に制限すると
二つの Noetherian $T$ 加群内の昇鎖となり、それぞれ停止する。
両方が停止した後は元の昇鎖も停止するので、$P'_i$ は Noetherian である。
固定した終点の column についても、始点の parity に分ければ
同じ三つの $T$ 加群が現れ、同じ議論が成り立つ。

基底変更後も、全ての line、非退化縮約、再帰式、
対称積と三次関係の同定は保たれる。
変更後の基底環が Noetherian であれば、その Segre 代数と二つの有限生成加群に
同じ証明を適用できる。
\end{proof}

\begin{remark}[二周期の場合の Zhang twist]
二周期同型 $\theta$ が指定される場合、$V_{2r}=E$, $V_{2r+1}=F$ と同定し、
$E,F$ と最初の二つの係数線束を有限開被覆上で自明化する。
非退化形式を $a=(\mathrm{id}\otimes P)\varepsilon_E$、
$b=(\mathrm{id}\otimes Q)\varepsilon_F$ と書く。$P,Q$ は可逆行列である。
添字 $2r,2r+1$ でそれぞれ $P^r,Q^r$ を使う隣接写像は、
\[
 (P^r\otimes P^{r+1})\varepsilon_E=\det(P)^r a
\]
と $Q$ に対する同式により、標準 linear 代数から元の $D$ への同型を与える。
この座標で $\theta$ は $P^{-1},Q^{-1}$ となるので、偶数環 $B_Z$ は
Segre 環に $(P^{-1},Q^{-1})$ が誘導する自己同型 $\tau$ による
$a*b=a\tau^{\deg a}(b)$ という Zhang twist として具体的に書ける。
主定理の row と column の証明は、この二周期性を用いない。
\end{remark}

\subsection{形式貼合せと普遍性}

ここから再び元の基底環 $R$ とそのイデアル $I$ を用いる。
第\ref{sec:formal-noetherian-patching}節の $\mathbb Z$ 代数に対する定理~\ref{thm:patch-z-algebra}を用いる。

\begin{proof}[定理 \ref{thm:relative-geometric-cubic-noetherian} の証明]
$I=(f_1,\ldots,f_r)$ と有限個の実際の生成元を取る。
各 $D(f_a)$ は $S_{\mathrm e}$ に含まれるので、相対曲線商と正規族による結果から
$D[f_a^{-1}]$ の全ての row と column は Noetherian である。
系~\ref{cor:linear-universal-row-noetherian} により $D/ID$ も同じ性質を持つ。
全 $D_{ij}$ は有限 projective なので、定理~\ref{thm:patch-z-algebra} を適用して
$D$ の全 row と column の Noetherian 性を得る。

偶数添字への制限もこの性質を保つ。
実際、$e_iD^{\mathrm{ev}}$ の部分加群 $J$ から $e_{2i}D$ 内に生成される部分加群を
$\widetilde J$ とすると、偶数終点では $(\widetilde J)_{2j}=J_j$ である。
従って $e_iD^{\mathrm{ev}}$ の部分加群の昇鎖は、$e_{2i}D$ 内へ拡大して安定させた後、
偶数終点へ制限すれば安定する。column にも同じ議論が使える。

任意の可換 Noetherian $R$-代数 $R'$ に対しても、
仮定の cubic データと rank-one ideal の定義は基底変更で保たれる。
楕円 open の相対曲線・cohomology base change・正規族は任意の基底変更と可換し、
Keeler の定理は $R'$ 自身に適用できる。
linear 閉部分の実際の tensor factorization、両方向への伝播、頂点線束 twist も保たれる。
従って $IR'$ と $D\otimes_RR'$ に同じ議論を適用できる。
\end{proof}

\begin{corollary}[AS4 被覆への帰結]
系~\ref{cor:relative-periodic-even-noetherian} の $B$ が、同じ基底上の次数付き代数 $A$ の regular normal な二次線による商であり、
その regular 性と商の同定が任意の基底変更と可換するとする。
このとき $A$ も任意の可換 Noetherian 基底変更後に左右 Noetherian である。
\end{corollary}

\begin{proof}
各基底変更後に系~\ref{cor:relative-periodic-even-noetherian} と補題~\ref{noeth-lem:normal-lift} を適用する。
\end{proof}

\paragraph{ここで証明していない事項.}
この結果は Noetherian 性に関するものであり、任意の導来 Morita 同値が
指定された相対 heart や偏極を保つことまでは主張しない。
AS4 被覆の全次数平坦性、正則正規元、相対 Gorenstein 解消は前提となる別の構成結果である。


\subsection{可換二次曲面の有限型近傍への適用}

\begin{theorem}[可換二次曲面の有限型相対実現]
\label{thm:quadric-finite-type-effectivity}
命題~\ref{prop:quadric-zariski-homological-as-cover} と
命題~\ref{prop:quadric-canonical-cubic-tails} の共通の有限型 affine 開
$U^\circ$ を取る。この開は可換二次曲面 $Q$ の点 $u$ を含む。
このとき以下が成立する。
\begin{enumerate}
\item 任意の可換 Noetherian affine 基底変更
$T=\operatorname{Spec}R\to U^\circ$ に対し、軌道環 $B_T$ と同じ基底上の被覆 $A_T$ は
左右 Noetherian である。
\item $T$ が局所有限型 $k$-site に属する場合、
$(\qgr B_T,\pi B_T)$ とその全有限型基底変更は
定義~\ref{def:relative-h2} の全公理を満たす。
これらは基底変更・fppf 降下と整合し、射
\[
 f:U^\circ\longrightarrow\mathfrak H_2
\]
を定める。$u$ の像は $(\coh Q,\OO_Q)$ である。
\item この族の canonical compact 圏は $\operatorname{Perf}\Lambda_T$ であり、
次数移動は既存の Morita 根 $\mathsf s_T$ に対応する。
その全軌道環は $B_T$、同じ根の AS4 被覆は $A_T$ である。
\end{enumerate}
元の site に属さない Noetherian $R$ に対しては、
有限型 $R$-test 上で同じ全公理を満たす系という意味で (2) を読む。
本定理だけから $\mathfrak H_2$ 全体の代数性は結論しない。
\end{theorem}

\begin{proof}
右加群の積の規約を先に固定する。
命題~\ref{prop:quadric-canonical-cubic-tails} のように
\[
 C_{ij}=D_{-j,-i}
       =\Hom(L_{-j},L_{-i})
\]
を用いる。$C$ の二周期同型は $\mathsf s^{-1}$ によるものである。
$C$ は普遍的な三次表示・四次重なり・幾何的 quintuple を持ち、
全成分が有限 locally free である。
定理~\ref{thm:relative-geometric-cubic-noetherian} と系~\ref{cor:relative-periodic-even-noetherian} を $C_T$ に適用する。
その偶数環と既存の右軌道環の同型は
\[
 C_{0,2n}=\Hom(L_{-2n},L_0)
 \xrightarrow{\mathsf s^n}
 \Hom(L_0,\mathsf s^nL_0)=B_n
\]
である。$a\in C_{0,2m}$, $b\in C_{0,2n}$ に対する積
$a\circ\mathsf s^{-m}(b)$ は、この写像により
$\mathsf s^n(\mathsf s^m a)\circ\mathsf s^n b$ に写る。
これは本文の右軌道積 $x*y=\mathsf s^{\deg y}(x)\circ y$ そのものである。
従って $B_T$ は左右 Noetherian である。

$A_T$ は普遍的に regular normal な二次線を持ち、その商が $B_T$ である。
補題~\ref{noeth-lem:normal-lift} により $A_T$ も左右 Noetherian である。
相対 Koszul・Frobenius 解消は既存の homological 被覆構成から得られるので、
これで定義上の相対 AS4 被覆となる。

定理~\ref{thm:quadric-noetherian-relative-promotion} を適用すると、
coherent Hom の有限性と flat base change、module-object 基底変更、fppf 降下、
構造対象の平坦性、canonical compact 圏、Serre 関係、AZ 豊富性、
全軌道環と同じ根の被覆が、同一の有限型開全体とその有限型基底変更で成立する。
全ての構成は重なりで canonical に整合するので、相対対象と上の射 $f$ が得られる。
体上の構造対象の rank one torsion-free 性と full bounded Serre 双対は、
同じ昇格定理における体上分類の適用により確認されている。
最後の主張は命題~\ref{prop:quadric-canonical-cubic-tails} による。
\end{proof}

% !TeX root = ../nc-quadric-surface-moduli.tex
% 2026-11-21 / v1.34（論文Bへ証明移管）
\section{幾何的二周期 cubic 族の普遍 canonical hull}
\label{sec:universal-cubic-hull}
\paragraph{主張の範囲.}
任意の幾何的二周期 cubic 族について、全幾何ファイバーの canonical hull を
有限 projective 加群間の射で相対化する。
可換二次曲面の特殊性、Artin 厚み、Krull 交叉定理を用いない。
体上の存在定理は Van den Bergh の Theorem~6.1 に依存し、
本節ではそこで存在する hull の六次元関係空間を、逆 coevaluation の核として
正確に同定する。相対化の新規性は未確認である。
二周期の存在と、その二乗が反 Serre 関手であることは別条件である。

\subsection{入力と定理}

$k$ を標数零の体、$R$ を可換 Noetherian $k$-代数とする。
$D$ を正の $R$-線形 $\mathbb Z$-代数とし、次を仮定する。
\begin{enumerate}
\item $D_{ii}=R$、$D_{ij}=0$ $(j<i)$、全ての $D_{ij}$ は有限 projective である。
  $V_i=D_{i,i+1}$ の階数は二であり、$D$ は三次関係の直和因子
  $R_i\subset V_i\otimes V_{i+1}\otimes V_{i+2}$、$\rank R_i=2$ により提示される。
\item 四次の交わり
  \[
  W_i=(R_i\otimes V_{i+3})\cap(V_i\otimes R_{i+1})
  \]
  は可逆 $R$-加群である。交わりは任意の基底変更と可換し、縮約により
  \[
  W_i\otimes V_{i+3}^{\vee}\xrightarrow{\sim}R_i,
  \qquad W_i\otimes V_i^{\vee}\xrightarrow{\sim}R_{i+1}
  \tag{UH-1}\label{hull-eq:coevaluation}
  \]
  が得られる。
\item 全幾何ファイバーは geometric quintuple に対応する標準 cubic AS 正則
  $\mathbb Z$-代数である。標準 cubic 解消は各内部次数で有限 projective の
  分裂完全複体となり、任意の基底変更に整合する。
\item 積と整合する実際の二周期同型 $\beta:D\xrightarrow{\sim}D(2)$ がある。
\end{enumerate}
これらは補題~\ref{lem:relative-geometric-cubic-flatness} の相対 cubic 解消から得られる入力である。
最後の条件は、幾何ファイバーごとの二周期同型の存在より強い。
本節の積の規約は
\[
 B_n=D_{0,2n},\qquad a\star b=a\,\beta^m(b)
 \quad(a\in B_m,\ b\in B_n)
 \tag{UH-2}\label{hull-eq:product}
\]
である。右軌道積 $a*b=s^n(a)\circ b$ と比較する際は、
式 (\ref{eq:cubic-tails-right-convention}) の
$C_{ij}=D_{-j,-i}$、$\beta_C=s^{-1}$ を用いて規約を合わせる。
偶数環と右軌道環の積の一致は定理~\ref{thm:quadric-finite-type-effectivity} の証明と同じである。

$U_i=V_i\otimes V_{i+1}$ および
\[
 S_i=R_i\otimes V_{i+3}+V_i\otimes R_{i+1}
       \subset U_i\otimes U_{i+2}
 \tag{UH-3}\label{hull-eq:relations}
\]
と置く。$S_i$ は階数七の直和因子である。
$U=U_0$ とし、$U\otimes U\xrightarrow{1\otimes\beta}U_0\otimes U_2$
に沿って $S_0$ を $U^{\otimes2}$ の部分加群とも考える。
この規約で $B=T_R(U)/(S_0)$ である。

\begin{theorem}[一般二周期 cubic 族の canonical AS4 被覆]
\label{thm:universal-cubic-hull}
上の仮定の下で、基底変更に自然な全射
\[
 \Phi_i:S_i\longrightarrow W_i,\qquad
 \Phi_i|_{W_i}=2\,\mathrm{id}_{W_i}
 \tag{UH-4}\label{hull-eq:Phi}
\]
があり、$H_i=\ker\Phi_i$ は階数六の直和因子である。
$\beta(H_i)=H_{i+2}$ であり、
\[
 A=T_R(U)/(H_0)\longrightarrow B
 \tag{UH-5}\label{hull-eq:hull}
\]
は次を満たす。
\begin{enumerate}
\item 各 $A_n$ は階数 $\binom{n+3}{3}$ の有限 projective 加群であり、
  任意の基底変更後にも同じ提示と次数別比較が成立する。
\item $P=A^!$ は階数列 $(1,4,6,4,1)$ の有限 projective な二次代数で、
  $P_n=0$ $(n\ge5)$ である。$P$ は相対 Koszul Frobenius であり、
  $A$ の左右および双加群 Koszul 解消は長さ四で、任意の基底変更後にも完全である。
\item $G=S_0/H_0\cong W_0$ は $A_2$ の直和因子である。
  $J=\ker(A\to B)$ に対し、全次数で普遍的な同型
  \[
   G\otimes_R A_{n-2}\xrightarrow{\sim}J_n
     \xleftarrow{\sim}A_{n-2}\otimes_R G
   \tag{UH-6}\label{hull-eq:regular-normal}
  \]
  がある。従って $A/(G)=B$、$GA=AG$ であり、$G$ の任意の局所生成元は
  左右正則正規な二次の元である。
\item 任意の可換 Noetherian $R$-代数 $R'$ に対して、$A_{R'}$ は左右 Noetherian である。
\end{enumerate}
各幾何ファイバーの被覆は Van den Bergh の canonical hull を
指定された二周期に沿って graded 化したものに一致する。
ここで Serre 平方根に関する仮定は用いない。
\end{theorem}

Noetherian 性には定理~\ref{thm:relative-geometric-cubic-noetherian} を用いる。
それ以外の結論は本節の有限 projective と tensor の議論による。
相対 AS4 という語は基底の解消を指し、特異な $R$ 上で環 $A$ 自身の
絶対大域次元が四であると主張しない。

\subsection{逆 coevaluation と次数三の正規性}

\subsubsection{評価の規約と線束値の写像}

\eqref{hull-eq:coevaluation} の逆を使うと
\[
 R_i\otimes V_{i+3}\cong
 W_i\otimes V_{i+3}^{\vee}\otimes V_{i+3}\longrightarrow W_i
\]
および
\[
 V_i\otimes R_{i+1}\cong
 V_i\otimes W_i\otimes V_i^{\vee}\longrightarrow W_i
\]
という二つの評価写像が得られる。
両者の $W_i$ への制限は階数二の恒等写像の trace、すなわち $2\,\mathrm{id}$ である。
従って交わりで一致し、\eqref{hull-eq:Phi} を定める。
この全射は $W_i\hookrightarrow S_i$ の半分を切断に持ち、
$S_i=H_i\oplus W_i$ が成り立つ。
構成は有限 projective の評価と普遍的な交わりのみを用いるため、
任意の基底変更および $D$ の同型に自然である。

局所生成元 $w_i$ を選び $\Phi_i(x)=\phi_i(x)w_i$ と書くと、
\[
 w_i=\sum_{a=1}^2r_a\otimes v_a
 \quad\Longrightarrow\quad
 \phi_i(r_a\otimes v_b)=\delta_{ab}.
 \tag{UH-7}\label{hull-eq:inverse-coeval}
\]
他方の分解でも同じ規約を用いる。
特に $\phi_i(w_i)=2$ であり、$w_i$ を $u w_i$ に変えると
$\phi_i$ は $u^{-1}\phi_i$ になる。
従って本来の大域対象は単一の scalar-valued functional ではなく、
\eqref{hull-eq:Phi} とその双対線 $W_i^{\vee}\subset S_i^{\vee}$ である。

\subsubsection{四周期と二周期の区別}

$W_0,W_1$ を局所的に自明化し、$w_{i+2}=\beta(w_i)$ と選ぶ。
$w_i\in V_i\otimes R_{i+1}$ と
$w_{i+1}\in R_{i+1}\otimes V_{i+4}$ の二つの非退化 tensor の比較により、
一意な同型 $\theta_i:V_i\xrightarrow{\sim}V_{i+4}$ が定まる。
具体的には $w_i=\sum_a v_a\otimes r_a$ と書いたとき
$w_{i+1}=\sum_a r_a\otimes\theta_i(v_a)$ とする。
四回の巡回から $\theta_i\otimes\theta_{i+1}\otimes\theta_{i+2}$ は
$R_i$ を $R_{i+4}$ に移す。
従って $\theta:D\to D(4)$ は代数同型である。
自然性から
\[
 \theta_{i+2}\beta_i=\beta_{i+4}\theta_i
 \tag{UH-8}\label{hull-eq:commute}
\]
であり、選んだ生成元の下で $\theta$ と $\beta$ は可換する。
この $\theta$ は \cite[\S5.2, (5.2)]{VdB2011} の canonical 四周期である。
一般に $\theta=\beta^2$ は成立しない。

\begin{lemma}[偶数環の相対 Koszul 性]
\label{hull-lem:even-koszul}
$B$ は相対 Koszul であり、$E=B^!$ の全次数は有限 projective で、
任意の基底変更と両立する。階数列は
\[
 (\rank E_n)_{n\ge0}=(1,4,7,8,8,\ldots).
 \tag{UH-9}\label{hull-eq:E-ranks}
\]
\end{lemma}
\begin{proof}
体上で、標準 cubic 解消を偶数頂点と奇数頂点に分ける。
奇数頂点の加群には、$V$ と $R$ を交互に用いる偶数環上の線型解消が得られる。
偶数部分の二つの短完全列に Tor を適用すると、偶数環の単純加群の
Tor は対角次数に集中する。これは \cite[Proposition 2.1]{BauwensVanDenBergh1998}
の証明に各 $V_i,R_i,W_i$ の添字を付けたものであり、実際の一周期を必要としない。
従って全幾何ファイバーの $B$ は Koszul である。

$B_1,B_2,B_3$ は仮定により有限 projective である。
\cite[Chapter 6, Theorem 7.3, Propositions 7.1--7.2]{PolishchukPositselski2005}
を各局所環で適用すれば、関係部分加群の全交わりと商が直和因子となる
（補題~\ref{lem:relative-koszul-degree-three}）。
これにより $E$ の次数別 projectivity、任意の基底変更との整合、
左右の普遍 Koszul 解消が得られる。
各ファイバーでは $\dim B_n=(n+1)^2$ だから
$h_E(t)=h_B(-t)^{-1}=(1+t)^3/(1-t)$ となり、\eqref{hull-eq:E-ranks} が従う。
\end{proof}

\begin{lemma}[普遍的な次数三分解]
\label{hull-lem:K3}
$E_3^{\vee}$ の対応する添字付き成分は
\begin{align*}
 K_{3,i}
 &= (S_i\otimes V_{i+4}\otimes V_{i+5})
       \cap(V_i\otimes V_{i+1}\otimes S_{i+2})\\
 &= (R_i\otimes R_{i+3})
       \oplus(V_i\otimes W_{i+1}\otimes V_{i+5}).
 \tag{UH-10}\label{hull-eq:K3}
\end{align*}
この同定は任意の基底変更と両立する。
\end{lemma}
\begin{proof}
右辺の二つの部分加群は $W_{i+1}$ の二つの包含により左辺に入る。
左辺は補題~\ref{hull-lem:even-koszul} により階数八の projective 加群であり、
右辺は階数四の projective 加群二つの直和からの射である。
従って各幾何ファイバーで直和であることを示せば十分である。

体上で二つの像の交わりを取ると、
$R_iV_{i+3}\cap V_iR_{i+1}=W_i$ により、それは
$W_iV_{i+4}V_{i+5}$ に入る。
一方、右 cubic 解消の次数六の部分は
\[
 W_i\otimes V_{i+4}\otimes V_{i+5}
 \hookrightarrow R_i\otimes D_{i+3,i+6}
\]
を単射にする。
右辺は $R_iV_{i+3}V_{i+4}V_{i+5}$ を $R_iR_{i+3}$ で割ったものだから、
$W_iV_{i+4}V_{i+5}\cap R_iR_{i+3}=0$ である。
従って交わりは零となる。次元八より同型である。
この議論は \cite[(2.7), Lemmas 2.3--2.4]{BauwensVanDenBergh1998} の添字付きの証明でもある。
有限 projective の間の同型判定を各ファイバーで行えば主張を得る。
\end{proof}

\begin{lemma}[functional 線の普遍正規性]
\label{hull-lem:normal}
\eqref{hull-eq:Phi} の双対は可逆加群 $N=W_0^{\vee}\subset E_2$ を与え、
$NE=EN$ が任意の基底変更後に成立する。
\end{lemma}
\begin{proof}
局所的な $w_i$ の選択の下で
\[
 c_L=\phi_i\otimes1:K_{3,i}\longrightarrow U_{i+4},\qquad
 c_R=1\otimes\phi_{i+2}:K_{3,i}\longrightarrow U_i
\]
と置く。$\Theta_i=\theta_i\otimes\theta_{i+1}$ とすると
\[
 c_L=\Theta_i c_R
 \tag{UH-11}\label{hull-eq:normal-contraction}
\]
である。これはファイバーでのみ成立する式ではなく、次の tensor 計算による恒等式である。

$\theta$ により四つ先の空間を同一視する。
$V_iW_{i+1}V_{i+5}$ 上では二つの縮約はいずれも
coevaluation と evaluation の合成なので同じ恒等射になる。
第一成分では、四つの空間の基底を $u_a,v_b,x_c,y_d$ とし、
$w_i=\sum t_{abcd}u_a v_b x_c y_d$ と書く。
\[
 r_d=\sum_{a,b,c}t_{abcd}u_av_bx_c\in R_i,\qquad
 t_c=\sum_{d,a,b}t_{abcd}y_du_av_b\in R_{i+3}
\]
はそれぞれ基底である。両縮約は $r_d\otimes t_c$ を
$\sum_{a,b}t_{abcd}u_av_b$ へ送るため、\eqref{hull-eq:normal-contraction} が成立する。

\eqref{hull-eq:product} に沿って graded 化する。
$T=\beta^{-2}\theta$ が $U$ に誘導する可逆写像を同じ記号で表すと、
\eqref{hull-eq:normal-contraction} の双対は
\[
 \phi\xi=(\xi\circ T)\phi\qquad(\xi\in E_1)
 \tag{UH-12}\label{hull-eq:graded-normality}
\]
となる。$T$ は代数自己同型から来るので、その双対は $E$ の自己同型を定める。
一次生成性から $\phi E=E\phi$ を得る。
生成元を替えると $\phi$ が単元倍されるだけなので、これは大域的な線 $N$ の正規性を意味する。
計算は基底上の恒等式であり、任意の基底変更後も同じである。
\end{proof}

\subsection{全幾何ファイバーでの canonical hull との同定}

本節のみ基礎環を代数閉体 $K$ とする。
\cite[Theorem 6.1]{VdB2011} により、canonical 四周期 $\theta$ の逆の
正規化自己同型の類を持つ、偶数 $\mathbb Z$-代数 $D^{\mathrm{ev}}$ の hull が存在する。
以下では、そこに現れる関係空間が実際に $H_i=\ker\phi_i$ であることを証明する。
これは、既知の hull の存在だけから候補の同定を省略しないための一段である。

\begin{lemma}[正規化式が決める評価の一意性]
\label{hull-lem:fibre-identification}
canonical hull の六次元関係空間を $H'_i\subset S_i$ とする。
このとき $H'_i=H_i$ である。
\end{lemma}
\begin{proof}
$H'_i=\ker\psi_i$ となる非零写像 $\psi_i:S_i\to K$ を選ぶ。
$g_i\in S_i/H'_i$ を $\psi_i(g_i)=1$ で正規化する。
この記法では、$S_i$ の元の hull における像は $\psi_i$ 倍の $g_i$ である。
canonical normalizer の意味から、ある $\lambda_i\in K^\times$ に対し
\[
 g_i u=\lambda_i\Theta_i^{-1}(u)g_{i+2}
 \qquad(u\in U_{i+4})
 \tag{UH-13}\label{hull-eq:hull-normalizer}
\]
である。scalar 自己同型の不定性は $\lambda_i$ に吸収した。

$z\in K_{3,i}$ を二通りに hull の次数三へ送れば
\[
 g_i(\psi_i\otimes1)(z)
   =(1\otimes\psi_{i+2})(z)g_{i+2}
\]
となる。$g_{i+2}$ の右正則性と \eqref{hull-eq:hull-normalizer} より
\[
 (\psi_i\otimes1)(z)
   =\lambda_i^{-1}\Theta_i(1\otimes\psi_{i+2})(z).
 \tag{UH-14}\label{hull-eq:psi-normalizer}
\]

これを第二直和成分 $V_iW_{i+1}V_{i+5}$ に制限する。
coevaluation によって、左辺は $F_i\otimes\mathrm{id}$、右辺は
$\lambda_i^{-1}\mathrm{id}\otimes G_i$ という形に書ける。
ここでは $\theta_i,\theta_{i+1}$ により始点と終点を同一視した。
二つの非零有限次元空間の endomorphism が
$F\otimes\mathrm{id}=\mathrm{id}\otimes G$ を満たせば、
行列単位との比較により $F,G$ は同じ scalar の恒等写像である。
従って
\[
 \psi_i|_{V_iR_{i+1}}=a_i\phi_i|_{V_iR_{i+1}},\qquad
 \psi_{i+2}|_{R_{i+2}V_{i+5}}
       =b_i\phi_{i+2}|_{R_{i+2}V_{i+5}}
\]
となる。$i$ を二つ戻した式も用いれば、$\psi_i$ の二つの制限は
それぞれ逆 coevaluation の scalar 倍である。
交わり $W_i$ 上で両者は同じ $\psi_i(w_i)$ を与え、
$\phi_i(w_i)=2\ne0$ なので二つの scalar は一致する。
よって $\psi_i=a_i\phi_i$ である。
$\psi_i\ne0$ から $a_i\ne0$、従って $H'_i=H_i$ を得る。
\end{proof}

\begin{remark}[Bauwens--Van den Bergh の式との直接照合]
\cite[Proposition 2.7, (2.11), Remark 3.3]{BauwensVanDenBergh1998} は graded cubic の場合、
\[
 w=x^tf=(Qf)^tx,\qquad
 \psi(f_i x_j)=\eta(\Lambda^{-1})_{ji},\qquad
 \psi(x_i f_j)=\eta(\Lambda Q^{-t})_{ij}
\]
と書く。canonical choice は $\Lambda=Q^t$ であるから
\[
 \psi(x_i f_j)=\eta\delta_{ij},\qquad
 \psi(f_i x_j)=\eta(Q^{-1})_{ij}.
 \tag{UH-15}\label{hull-eq:BV-convention}
\]
これは $w$ の二つの分解に対する \eqref{hull-eq:inverse-coeval} と正確に一致する。
原典では偶数環の元の次数を元の cubic 環の次数で記している箇所があるため、
そこに現れる次数四の functional は本節の $E_2$ である。
一般 $\mathbb Z$-代数での同定は補題~\ref{hull-lem:fibre-identification} が直接与える。
\end{remark}

\begin{corollary}[体上の双対と正則性]
\label{hull-cor:fibre-dual}
任意の実際の二周期 $\beta$ に対して、$T(U)/(H_0)$ は canonical hull の
graded 化である。$\phi\in E_2$ は左右正則正規で、
$E/(\phi)$ は階数列 $(1,4,6,4,1)$ の Koszul Frobenius 代数である。
\end{corollary}
\begin{proof}
$H_i$ の構成の自然性により $\beta(H_i)=H_{i+2}$ であるから、
二周期は hull の提示そのものに持ち上がる。
これは \cite[Lemmas 6.4--6.5, Proposition 6.6]{VdB2011} の構成と一致する。
canonical hull は標準四次元 AS 正則なので Koszul であり、
その二次双対 $P$ は Frobenius、Hilbert 級数 $(1+t)^4$ を持つ。
二次提示から $P=E/(\phi)$ である。
補題~\ref{hull-lem:normal} より $\phi$ は正規である。
\[
 h_E(t)=\frac{(1+t)^4}{1-t^2},\qquad h_P(t)=(1+t)^4
\]
なので、各次数の商と積の次元を比べると
$E_{n-2}\xrightarrow{\phi\cdot}E_n$ および
$E_{n-2}\xrightarrow{\cdot\phi}E_n$ は単射となる。
これで正則性が従う。原典でも \cite[Propositions 2.7, 3.2]{BauwensVanDenBergh1998} が同じ双対関係を示す。
\end{proof}

\subsection{非被約基底上での相対化}

再び一般の Noetherian $R$ に戻る。
定理~\ref{thm:universal-cubic-hull} の証明を完結する。

\subsubsection{まず双対の商を平坦にする}

補題~\ref{hull-lem:even-koszul} より $E_n$ は有限 projective である。
補題~\ref{hull-lem:normal} の正規線 $N=W_0^{\vee}$ による左右乗法
\[
 N\otimes E_{n-2}\longrightarrow E_n,
 \qquad E_{n-2}\otimes N\longrightarrow E_n
 \tag{UH-16}\label{hull-eq:E-inject}
\]
は、系~\ref{hull-cor:fibre-dual} により全幾何ファイバーで単射である。
有限 projective 加群間の写像のファイバー単射性から、
これらは局所直和因子への単射になる。
従って
\[
 P=E/(N)
 \tag{UH-17}\label{hull-eq:P}
\]
の全次数は有限 projective であり、商は任意の基底変更と可換する。
\eqref{hull-eq:E-ranks} の差から階数列は $(1,4,6,4,1,0,\ldots)$ となる。
ここでは商のファイバー次元が一定であることだけから平坦性を推論していない。
既に projective である加群間の写像 \eqref{hull-eq:E-inject} を用いている。

$P$ は二次代数であり、$P_1,P_2,P_3$ は projective、全幾何ファイバーは Koszul である。
従って再び \cite[Chapter 6, Theorem 7.3]{PolishchukPositselski2005} を \emph{$P$ に}適用する。
全ての関係交わりが直和因子となるため、二次双対
\[
 P^!=T_R(U)/(H_0)=A
\]
の全次数も有限 projective で、任意の基底変更に整合する。
全ファイバーでの canonical hull の Hilbert 級数により
$\rank A_n=\binom{n+3}{3}$ となる。
この順序により、$A_3$ の局所自由性を証明前に仮定する循環を避けている。

\subsubsection{Frobenius、AS 解消、二次切断}

$L=P_4$ は可逆加群である。
各積の adjoint
\[
 P_i\longrightarrow\operatorname{Hom}_R(P_{4-i},L)
 \qquad(0\le i\le4)
 \tag{UH-18}\label{hull-eq:Frobenius}
\]
は有限 projective 加群間の射で、全幾何ファイバーで同型である。
従って同型であり、任意の基底変更後もそうである。
これが相対 Frobenius 性である。
Koszul 複体の各内部次数は有限な projective 複体で、
そのファイバーは体上の Koszul 解消だから分裂完全となる。
$P_i=0$ $(i>4)$ より、左右および対角双加群の解消は長さ四である。
Frobenius pairing による双対化から
\[
 \underline{\operatorname{Ext}}_A^j(R,A)
 \cong
 \begin{cases}
 L\otimes_RR(4),&j=4,\\
 0,&j\ne4
 \end{cases}
 \tag{UH-19}\label{hull-eq:AS}
\]
を左右に得る。ここで表示は対応する側の加群としてのものとし、
双加群同定には通常の Frobenius 自己同型を含める。
解消は任意の基底変更後にも完全なので、これらは相対 AS4 条件である。

$A\to B$ は提示から全射である。
$A_n,B_n$ が有限 projective だから $J_n=\ker(A_n\to B_n)$ も projective であり、
核は任意の基底変更と可換する。
$G=S_0/H_0\cong W_0$ は $J_2$ と同一視される。
左右乗法 \eqref{hull-eq:regular-normal} は有限 projective 加群間の射で、
各幾何ファイバーでは canonical hull の正則二次切断の同型である。
従って全次数で同型であり、任意の基底変更後も同型である。
これにより正則性、正規性、商 $A/(G)=B$ を同時に得る。
特に局所生成元 $g$ について $gA=Ag$ と両側正則性があるので、
$ga=\nu(a)g$ によって定まる写像 $\nu$ は全単射の代数自己同型である。
規約 \eqref{hull-eq:graded-normality} では、$\nu|_U=T^{-1}$ となる。
実際、$\phi(c)=1$ となる $c\in S_0$ で $g$ を表し、
$c\otimes u-\nu(u)\otimes c=h_L+h_R$、
$h_L\in H_0\otimes U$、$h_R\in U\otimes H_0$ と書く。
すると $z=c\otimes u-h_L=\nu(u)\otimes c+h_R\in K_3$ であり、
二つの縮約はそれぞれ $u$ と $\nu(u)$ になる。
\eqref{hull-eq:normal-contraction} を graded 化した等式から $u=T\nu(u)$ を得る。
この確認も幾何ファイバーへの還元を用いない。

\subsubsection{有限判定と Noetherian 性}

定理~\ref{thm:relative-finite-as-cover} の五条件は、本構成では全て基底上で成立する。
具体的には
\[
 \rank A_3=20,\quad\rank A_4=35,\qquad
 \rank P_3=4,\quad\rank P_4=1,\quad P_5=0,
\]
左右の $G$ 乗法の像の一致、次数一・二に対する局所直和因子への単射、
および \eqref{hull-eq:Frobenius} である。
従って canonical な階数六部分束 $H_0\subset S_0$ は、
被覆の有限判定部分スキームへの実際の切断を与える。
全ファイバーで有限条件が成り立つという主張のみで代用していない。

最後に任意の Noetherian $R'$ に移る。
系~\ref{cor:relative-periodic-even-noetherian} より $B_{R'}$ は左右 Noetherian である。
$G_{R'}$ を有限アフィン開被覆上で自明化する。
補題~\ref{noeth-lem:normal-lift} の議論を適用する。
正則正規な正次数元 $g$ に対して、$\operatorname{gr}_{(g)}A_{R'}$ は
$B_{R'}$ 上の skew polynomial 環となるため左右 Noetherian である。
固定内部次数では $(g)$ 進フィルトレーションは有限なので、
初期斉次生成元を有限個持ち上げる還元は有限回で止まる。
これにより斉次イデアルは有限生成、最高内部次数による還元で一般のイデアルも有限生成となる。
基底の有限開被覆から忠実平坦に降下すれば $A_{R'}$ は左右 Noetherian である。
これで定理~\ref{thm:universal-cubic-hull} の全ての主張を得た。

\subsection{Serre 平方根の条件は独立である}

\begin{remark}[二周期と反 Serre 関手]
任意の実際の二周期 $\beta$ が上の AS4 被覆を持つことと、
そこから来る exact 自己同値 $s$ が
\[
 s^2\simeq S^{-1}[2]
 \tag{UH-20}\label{hull-eq:Serre-root}
\]
を満たすことは別である。
後者は本定理の仮定にも結論にも含まれない。
有限 foundation $\Lambda$ 上で $s$ が可逆 perfect 双加群 $P_s$ により表され、
$\mathcal S_\Lambda=\mathbf R\operatorname{Hom}_R(\Lambda,R)$ が Serre 双加群である場合、
正確な追加条件は
\[
 P_s\mathbin{\otimes^{\mathbf L}}_\Lambda P_s
    \simeq\mathcal S_\Lambda^{-1}[2]
 \quad\text{in }\operatorname{Perf}(\Lambda^e).
 \tag{UH-21}\label{hull-eq:Morita-Serre}
\]
右辺の逆は可逆双加群の Morita 逆である。
根・根同型を fppf 局所の存在条件として扱う定義では、
基底線束 $\ell$ を許した
$P_s\mathbin{\otimes^{\mathbf L}} P_s\simeq\mathcal S_\Lambda^{-1}[2]\otimes_R\ell$ を
局所自明化して同じ条件を表せる。
この比較は定理~\ref{thm:bimodule-line-comparison} と
系~\ref{cor:foundation-antiserre-comparison-locus} で認識する対象である。
単に対象ごとに $s(F_0)\cong F_2$, $s(F_1)\cong F_3$ であることだけでは
\eqref{hull-eq:Morita-Serre} は従わない。
\end{remark}

\begin{example}[period two は自動的な Serre 根ではない]
$Q=\mathbb P^1\times\mathbb P^1$、$\sigma\in\operatorname{Aut}(Q)$ を
$\sigma^2\ne1$ かつ両 ruling を保つ自己同型とする。
例えば標数零で
$\sigma([x_0:x_1],[y_0:y_1])=([2x_0:x_1],[y_0:y_1])$ とできる。
$s=(-\otimes\OO_Q(1,1))\circ\sigma^*$ は標準 cubic helix の各対象を二項先へ送り、
実際の二周期を与える。
$s^{-n}\OO_Q\cong\OO_Q(-n,-n)$ なので $(\OO_Q,s)$ は AZ ample でもある。
従ってその偶数軌道環には本定理の canonical AS4 被覆がある。
しかし
\[
 s^2=(-\otimes\OO_Q(2,2))\circ\sigma^{2*},\qquad
 S^{-1}[2]=(-\otimes\OO_Q(2,2))
\]
であり、動かされる点の skyscraper sheaf を比べれば同型でない。
従って同じ指定二周期を H2 の witness とするには、
\eqref{hull-eq:Serre-root} を別途確かめる必要がある。
この例は underlying $(\coh Q,\OO_Q)$ が H2 に属することを否定しない。
\end{example}

\paragraph{一般 cubic 点への適用で残る入力.}
本節により、実際の geometric cubic 二周期族が構成された後の AS4 被覆は
Q 固有の障害ではなくなる。
一般の内在的 H2 族に適用するには、構造対象を含む適合する四項 foundation、
その上で代数化された実際の Morita 根、
および \eqref{hull-eq:Morita-Serre} と全軌道環の比較を供給する必要がある。
これらの存在範囲および根を忘れた相対実現部分層の代数性は、本節では主張しない。


% !TEX root = ../nc-quadric-surface-moduli.tex
% 2026-11-21 / v1.35
\section{一般の幾何的 cubic 族の相対内在的実現}
\label{sec:relative-cubic-families}

本節では、可換二次曲面を中心とするという制限を外し、
普遍 Noetherian 性と canonical hull を相対 tails の認識に結び付ける。
Serre 平方関係は、まず指定された周期の独立な入力として明記する。
次節の Picard 構成が非線型の普遍族上でこの入力を供給する。

\subsection{周期を選ぶ前の canonical compact 圏}

$R$ を可換 Noetherian $k$-代数とし、
$C$ を定理~\ref{thm:relative-geometric-cubic-noetherian} の仮定を満たす
正の向きの $R$-線形 $\mathbb Z$-代数とする。
右表現可能加群を $P_i=e_iC$、$\mathcal O_i=\pi P_i$ と書く。
三次関係を $R_i$、四次 overlap を $W_i$ とする。

\begin{proposition}[一般 geometric cubic 族の canonical tails]
\label{prop:general-cubic-canonical-tails}
任意の可換 Noetherian $R$-代数 $R'$ 上で
$\operatorname{QGr}(C_{R'})$ は locally Noetherian であり、
任意の連続四項の直和はその非有界導来圏の compact tilting generator である。
特に
\[
 T=\bigoplus_{i=0}^3\mathcal O_{-i},\qquad
 \Gamma=\operatorname{End}(T)
\]
と置くと、$\Gamma$ は階数 $24$ の有限 projective $R$-代数であり、
\[
 D(\operatorname{QGr}C)\simeq D(\operatorname{Mod}\Gamma),\qquad
 D(\operatorname{QGr}C)^c\simeq\operatorname{Perf}\Gamma.
 \tag{GC-Morita}\label{eq:general-cubic-canonical-morita}
\]
この比較は実際の module-object 引戻しの導来関手と整合する。
compact 圏は相対 smooth proper で、相対 Serre 関手とその任意基底変更を持つ。
この命題では二周期の存在を仮定しない。
\end{proposition}

\begin{proof}
全成分の projectivity と四次 overlap に対し、
補題~\ref{lem:relative-geometric-cubic-flatness} の標準 cubic 解消を用いる。
その各内部次数は $R$ 上分裂完全である。
双対複体にも同じ fiber 判定を適用すると
\[
 \operatorname{Ext}^q_{\operatorname{Gr}C}(S_i,P_j)=0
 \quad\text{unless }(q,j)=(3,i+4),\qquad
 \operatorname{Ext}^3(S_i,P_{i+4})=W_i^\vee
\]
を得る。実際、この双対複体の各項は有限 projective である。
消滅する fiber 複体は局所分裂完全であり、
例外の $(q,j)=(3,i+4)$ では最終項が $W_i^\vee$、
直前の項は正の向きにより零である。
これらの議論は任意の基底変更後にも成立する。
従って定理~\ref{thm:relative-cubic-tails-tilting} の全入力が揃い、
式~\eqref{eq:general-cubic-canonical-morita} を得る。

定理~\ref{thm:relative-geometric-cubic-noetherian} により
$C_{R'}$ の各 row は Noetherian である。
したがって $\operatorname{QGr}C_{R'}$ は locally Noetherian で、
$\mathcal O_i$ は Noetherian 生成系をなす。
その有限 cones と直和因子である compact 対象は
$D^b(\operatorname{qgr}C_{R'})$ に属する。

四項窓の前向き Hom の階数は $1,2,4,6$ である。
対角を $E=R^4$ とし、厳密に上三角な部分を $J$ とすると、
$\Gamma=E\oplus J$、各 $e_iJe_j$ は有限 projective で、
$J^{\otimes_E4}=0$ である。
$E$ 上の normalized bar 解消は長さ三であり、
各項 $\Gamma\otimes_EJ^{\otimes_Er}\otimes_E\Gamma$ は
有限生成 projective な $\Gamma$ 双加群となる。
従って $\Gamma$ は相対 smooth、$R$ 上有限 projective なので proper である。
これは補題~\ref{lem:foundation-relative-smooth-proper} と同じ直接の議論である。
相対 Serre 双加群は
$\mathcal S_\Gamma=\mathbf R\operatorname{Hom}_R(\Gamma,R)$
であり、任意の基底変更と整合する。

最後に各 $\mathcal O_i$ は、全成分が基底上平坦な $P_i$ の像なので、
module-object の意味で基底上平坦である。
通常の引戻しがこの tilting summand 上で導来引戻しを計算する。
summand と全ての射を基底変更すると、変更後の $T$ と $\Gamma$ を得る。
よって Morita 表示における実際の導来引戻しは
$-\otimes_\Gamma^{\mathbf L}\Gamma_{R'}$ であり、
compact 圏と Serre 双加群の比較もこの引戻しに対して成立する。
\end{proof}

\subsection{実際の二周期からの全相対条件}

$C$ に二周期同型 $\beta:C\to C(2)$ を指定し、
\[
 B_n=C_{0,2n},\qquad a\star b=a\,\beta^m(b)
 \quad(a\in B_m,\ b\in B_n)
 \tag{GC-prod}\label{eq:general-cubic-period-product}
\]
と定める。$\check B\simeq C^{\mathrm{ev}}$ の下で、
$B$ の次数移動 $(1)$ は $C$ の right tails 上で
$\mathcal O_i\mapsto\mathcal O_{i-2}$ という完全自己同値 $s$ を与える。
矢も $\beta^{-1}$ で移す。単に対象の置換を指定しているのではない。
式~\eqref{eq:general-cubic-canonical-morita} によって
$s$ を表す可逆 perfect $\Gamma$ 双加群を $P_s$ と書く。

\begin{theorem}[Serre 適合二周期族の相対効果性]
\label{thm:geometric-cubic-relative-effectivity}
上の $C,\beta$ に対し、基底の fppf 局所で
\[
 P_s\otimes_\Gamma^{\mathbf L}P_s
       \simeq\mathcal S_\Gamma^{-1}[2]
 \tag{GC-root}\label{eq:general-cubic-serre-input}
\]
が与えられ、任意の基底変更ではこの同型を引き戻すとする。
このとき有限型 $R$-基底上の系
\[
 \mathcal C_{R'}=\operatorname{qgr}(B\otimes_RR'),\qquad
 \mathcal O_{R'}=\pi(B\otimes_RR')
 \tag{GC-family}\label{eq:general-cubic-relative-family}
\]
は定義~\ref{def:relative-abelian-system}、
\ref{def:relative-h2} の全条件を満たす。
特に $\operatorname{Spec}R$ が元の局所有限型 $k$-site に属すれば、
実際の $\mathfrak H_2(R)$ の対象を与える。
平方根、平方同型、二周期、cubic helix および AS4 被覆は、
構成に用いる witness であり、出力の標識には加えない。

各 affine Noetherian 基底変更上で全相対 $\chi_q$ も成立する。
$R$ が元の site に属さない場合は、
有限型 $R$-test 上で同じ全公理が成立するという意味に限る。
\end{theorem}

\begin{proof}
\emph{Noetherian heart と基底変更。}
系~\ref{cor:relative-periodic-even-noetherian} により、
全ての Noetherian $R'$ について $B'=B\otimes_RR'$ は左右 Noetherian である。
従って $\operatorname{QGr}B'$ は locally Noetherian、
その Noetherian 部分圏は $\operatorname{qgr}B'$ であり、
この部分圏の Ind 圏が $\operatorname{QGr}B'$ となる。
偶数頂点への制限は
$\operatorname{QGr}C_{R'}\simeq\operatorname{QGr}B'$ を与える。
この同値は表現可能加群と全ての射を保ち、任意の基底変更と整合する。

$B'$ は有限 projective な一次成分で生成される。
補題~\ref{lem:qgr-module-object-basechange} により、
module-object 基底変更と tensor の結合則による coherent 比較が得られる。
flat base change は完全、faithfully flat base change はさらに保守的である。
したがって系~\ref{cor:quadric-local-tails-fppf-descent} の
comonadic な降下の証明がそのまま適用でき、
対象・射・完全列の有効 fppf descent を得る。
この議論に可換 quadric の特殊ファイバーは用いない。
補題~\ref{lem:relative-cubic-coherent-cohomology} により、
全ての coherent Hom は有限で、flat base change と可換する。

\emph{canonical compact 圏と相対 Serre。}
前命題と偶数制限により、canonical compact 圏は
$\operatorname{Perf}\Gamma_{R'}$ である。
これは相対 smooth proper で、実際の module-object 導来引戻し、
構造対象、および Serre 双加群と整合する。
構造対象は基底上平坦で、
\[
 \mathbf R\operatorname{Hom}(\mathcal O_{R'},\mathcal O_{R'})
       \simeq R'[0]
\]
である。従って (RQ1) と相対 Serre 条件が成立する。
式~\eqref{eq:general-cubic-serre-input} の引戻しにより
$s^2\simeq S^{-1}[2]$ が同じ canonical compact 圏上で成り立つ。

\emph{AZ 豊富性と全軌道環。}
有限生成次数付き $B'$ 加群の十分高い truncation は有限生成であり、
任意に大きい次数以上の有限個の斉次生成元を取れる。
その像を $\operatorname{qgr}B'$ へ移すと AZ (A1) を得る。
任意の全射の kernel は Noetherian なので、
補題~\ref{lem:relative-cubic-coherent-cohomology} の一様な
$\operatorname{Ext}^1$ 消滅を適用して AZ (A2) を得る。
また式~\eqref{eq:cubic-tails-line-cohomology} より
\[
 \operatorname{Hom}(\mathcal O_0,s^n\mathcal O_0)
       =C_{-2n,0}\xrightarrow[\sim]{\beta^n}C_{0,2n}=B'_n .
\]
ここでは非負次数の全てに対し同型があり、低次数の切捨てはない。
右軌道積 $x*y=s^n(x)\circ y$ をこの同型で送ると、
$x\star y=x\,\beta^m(y)$（$x\in B'_m,\ y\in B'_n$）となる。
従って全軌道環と式~\eqref{eq:general-cubic-period-product} は、
全次数・積・任意の基底変更に関して一致する。

\emph{同じ根に対する相対 AS4 被覆。}
定理~\ref{thm:universal-cubic-hull} を $C,\beta$ に適用する。
得られる $A\to B$ は全次数 projective な canonical hull であり、
可逆部分束 $G\subset A_2$ による左右正則正規切断である。
$G$ を有限 affine 開被覆で自明化すれば、実際の二次の元 $g$ を取れる。
各 $A_{R'}$ は左右 Noetherian、相対 Koszul Frobenius であり、
全ての解消と二次切断は任意の基底変更後も完全である。
これが同じ次数移動 $s$ に対する (RQ4) を与える。

\emph{幾何ファイバーと全相対 $\chi$。}
各幾何ファイバーの $C$ は geometric cubic AS 正則
$\mathbb Z$-代数である。
定理~\cite[\texttt{thm:all-admissible-inverse}]{QuadricIntrinsic2026} と偶数制限により、
その構造対象を持つ tails は体上の $\mathcal H_2$ に属する。
特に Gabriel--Krull rank one torsion-free 条件と、
fiber の全 $D^b$ に関する Serre 条件が成立する。
その fiber は上記の module-object 基底変更によるものと一致するので、
(RQ2) も満たされる。

最後に全 $\chi_q$ については
補題~\ref{lem:relative-cubic-global-chi} の証明を適用する。
そこで用いる入力は、同じ affine 基底での Noetherian 性、
構造対象間の普遍 cohomology 公式、coherent 対象の一様 Serre 消滅、
$\operatorname{cd}\omega\le2$ および有限 projective な各 $B'_n$ である。
これらは本節で全て確認されている。
有限幅の局所 cohomology spectral sequence によるその証明は、
$Q$ を中心とする族という条件を用いない。
反対側にも同じ議論を適用して左右の全相対 $\chi_q$ を得る。
\end{proof}

\begin{remark}[Hom の向きで表示する場合]
\label{rem:general-cubic-right-orientation}
framed 五つ組を $D_{ij}=\operatorname{Hom}(E_i,E_j)$ として
回復する場合には、右加群のために
\[
 C_{ij}=D_{-j,-i},\qquad \beta_C=\beta_D^{-1}
\]
を用いる。このとき $E_i$ は $\pi e_{-i}C$ に対応し、
四項 foundation の頂点・全 Hom・積が一致する。
canonical 四周期も反転し、$\beta_D^2=\theta_D$ が
頂点スカラーを除いて成立すれば、$C$ 側でも対応する平方関係が成立する。
相対 Serre の比較は $C$ の canonical compact 圏で行う。
反対の積を暗黙に同じ環と同一視しない。
\end{remark}

% !TEX root = ../nc-quadric-moduli.tex
% 2026-11-21 / v1.35
\section{相対 Picard 作用による cubic Serre 平方根の構成}
\label{sec:relative-picard-roots}

\paragraph{主張の範囲と依存.}
本稿は標数零における\emph{非線型} geometric cubic 五つ組の普遍族上に、
étale 局所に実際の二周期を構成する。
その平方は四次 tensor の巡回から定まる canonical 四周期に
頂点スカラーを除いて一致する。
全ての恒等式を被約な普遍基底上で証明してから引き戻すので、
非被約な Artin 基底も結論に含む。
原典の既知入力は、Artin--Tate--Van den Bergh の
\emph{相対} Picard 作用と、Van den Bergh の geometric 五つ組の構成である。
相対 tails と有限 foundation の比較、および Serre の双加群比較には、
定理~\ref{thm:relative-cubic-tails-tilting} と定理~\ref{thm:bimodule-line-comparison} を用いる。
本稿は任意の二周期が Serre 平方根であるとは主張しない。
相対化の新規性は未確認である。

\subsection{一次資料の正確な内容}

Van den Bergh \cite[Proposition 5.6.1]{VdB2011} は、
代数閉体 $k$、$\operatorname{char}k\ne2$ 上の全ての quadric
$\mathbb Z$ 代数 $D$ に対して $D\simeq D(2)$ を示す。
全 $\mathbb Z$ 代数の二周期性が主張されている。
その後の Corollary 5.6.2 が偶数部分の通常の次数付き環による表示を与える。

非線型の場合、楕円 helix $(C,(\mathcal L_i)_i)$ に対し
\[
 N=\mathcal L_2\mathcal L_0^{-1}\in\Pic^0(C),\qquad
 \mathcal L_{i+2}\simeq\mathcal L_iN
\]
である。$a\in\Pic^0(C)$、$a^{-2}=N$ を選び、
\cite[Theorem 4.2.3]{VdB2011} の作用
\[
 \eta:\Pic^0(C)\longrightarrow\Aut(C),\qquad
 \eta(a)^*L=L\otimes a^{-\deg L}
 \tag{PR-1}\label{eq:field-eta}
\]
を用いる。$\sigma=\eta(a)$ は各次数二の束 $\mathcal L_i$ を $\mathcal L_{i+2}$ へ移す。
原典の Proposition 6.6 は二周期を標準 hull へ持ち上げるが、
任意の二周期 $\beta$ に対する $\beta^2=\theta$ を主張してはいない。

\begin{proposition}[ATV の相対 Picard 作用]\label{prop:ATV-relative-action}
$C\to S$ を $\PP^2_S$ の三次因子、または
$\PP^1_S\times_S\PP^1_S$ の bidegree $(2,2)$ の因子の平坦族とする。
$P=\Pic^0_{C/S}$ を、各幾何ファイバーの各被約既約成分で次数零の
直線束の類を表す相対 Picard 部分群とする。
$P$ の $C$ への作用があり、その構成は任意の底変更および族の同型に自然である。
作用の符号を \eqref{eq:field-eta} と合わせて $\eta$ と書く。
\end{proposition}

\begin{proof}[相対作用の構成]
\cite[p.~355, Proposition 5.5 and Corollary 5.7]{ArtinTateVanDenBergh1991} は初めから
任意の基底 $S$ 上の因子の族 $C/S$ を扱う。
相対多重次数零の直線束 $a$ と section $p:S\to C$ に対し、
$a\otimes I_p$ が局所的にある一意な section $p'$ の ideal と同型になることを示す。
具体的には
\[
 \pi_*\mathcal Hom(a\otimes I_p,\OO_C)
\]
が階数一で、その局所生成射が単射になり、余核が section の構造層となる。
一意性によりこの操作が貼り合わさり、任意の底変更と両立する群作用となる。
原典の (5.8) の $\eta_a$ はこの作用に $a^{-1}$ を入れた符号である。
したがって smooth elliptic fiber 上では usual translation となり、
Lemma 5.10(iii) が \eqref{eq:field-eta} を与える。
ideal と tensor による一意な構成なので、曲線族の同型への自然性も成り立つ。

Corollary 5.7 は相対的である一方、Lemma 5.10(iii) の記述は体上である。
以下では、この違いを埋める普遍基底上の議論を明記する。
\end{proof}

\subsection{被約な普遍基底で証明してから底変更する}

以下、$k$ は代数閉体、$\operatorname{char}k=0$ とする。
四つの二次元空間 $V_0,\ldots,V_3$ を固定し、
\[
 \mathscr G\subset\PP(V_0\otimes V_1\otimes V_2\otimes V_3)
\]
を geometric 五つ組の開部分とする。
$02\mid13$ flattening の階数一の閉部分を除き、
\[
 S=\mathscr G\setminus\mathscr G_{\mathrm{lin}}
\]
と置く。$S$ は滑らかで被約、既約な有限型 $k$ スキームである。
以下の構成は $S$ の affine open 上で行って貼り合わせてもよい。

普遍直線 $W\subset(V_0\otimes\cdots\otimes V_3)\otimes\OO_S$ を巡回し、
$V_{i+4}=V_i$ として $W_i$ と三次関係 $R_i$ を定める。
これにより普遍 cubic $\mathbb Z$ 代数 $D$ を得る。
$W_i$ は全て同じ基底線 $W$ の巡回像であり、この共通の同定の下で
canonical 四周期
\[
 \theta:D\longrightarrow D(4)
 \tag{PR-2}\label{eq:theta}
\]
は各 $V_i=V_{i+4}$ 上の恒等である。
この記述に $W$ の大域的な生成元は必要ない。
局所生成元 $w_i$ を使えば
\[
 w_i=\sum_a v_a\otimes r_a,
 \qquad w_{i+1}=\sum_a r_a\otimes\theta_i(v_a)
\]
という \cite[\S5.2]{VdB2011} の定義そのものである。

\begin{lemma}[普遍族の次数別平坦性]\label{lem:universal-flatness}
各 $D_{ij}$ は有限局所自由で、標準 cubic 解消は各内部次数で分裂完全である。
その双対複体も AS 双対の所定の次数以外で完全で、残る cohomology は可逆層である。
従ってこれらの提示・解消・双対・関係の交わりは任意の底変更と両立する。
\end{lemma}

\begin{proof}
固定した $i,j$ に対し、$D_{ij}$ は有限個の tensor 積と三次関係からなる
有限自由加群間の行列の余核として提示される。
全ての幾何ファイバーは cubic AS 正則で、$\dim D_{ij}$ は一定である。
従って提示行列のファイバー階数は一定である。
適切な階数の小行列式が単元となる近傍で行列を掃き出すと、
余分な block の全成分は全ての点で零となる。
$S$ は被約なので、その block 自体が零である。
従って余核は局所自由で、提示は局所的に標準形へ分裂する。

標準 cubic 複体の固定内部次数は有限局所自由複体となり、
全幾何ファイバーで完全である。
そのような有限複体は局所的に分裂完全であるから、任意底変更後も完全である。
この解消を使って定義される kernels、images および四次の overlap も
直和因子となり、任意底変更と両立する。
双対複体の固定内部次数も有限局所自由複体である。
各幾何ファイバーの AS 双対により、その各 cohomology の次元は一定であり、
これと各項の階数から各微分の階数も一定となる。
被約性を用いた同じ行列の標準形により、像・核・cohomology は局所直和因子となり、
その形成は任意底変更と両立する。
特に右単純対象 $S_i$ と $P_j=e_jD$ について
\[
 \operatorname{Ext}^q_{\operatorname{Gr}D}(S_i,P_j)=0
 \quad\text{unless }q=3,\ j=i+4,
 \qquad
 \operatorname{Ext}^3(S_i,P_{i+4})\simeq W_i^\vee.
\]
最後の同定は解消の最高項への Hom から得られる。
各固定次数ごとの議論で十分であり、無限個の次数を同時に自明化する
一つの開被覆は必要ない。
\end{proof}

普遍 point curve を $\pi:C\to S$ とする。
隣接二因子の評価による $2\times2$ 行列の determinant が
$\PP(V_0^\vee)\times\PP(V_1^\vee)$ 内に $C$ を定める。
geometric 条件は、determinant が零である点でもこの行列の階数が一であることを保証する。
三点 correspondence の隣接二点への射影は、単元小行列式による掃き出しから
スキームとして同型になる。
従って全 $\mathcal L_i$ を同じ $C$ 上の直線束として定義できる。
これは第~\ref{sec:relative-cubic-noetherian}節の相対 point-curve 構成である。

$S$ は線型 locus を除いているので、determinant は各幾何ファイバーで非零である。
したがって $C/S$ は平坦な relative Cartier divisor である。
その標準 exact sequence
\[
 0\longrightarrow\OO_{\PP^1_S\times\PP^1_S}(-2,-2)\otimes\pi^*\ell
 \longrightarrow\OO_{\PP^1_S\times\PP^1_S}
 \longrightarrow\OO_C\longrightarrow0
 \tag{PR-3}\label{eq:curve-sequence}
\]
（$\ell$ は基底の可逆層）から、$\pi_*\OO_C=\OO_S$ が普遍的に成り立ち、
$R^1\pi_*\OO_C$ は可逆層、$R^{>1}\pi_*\OO_C=0$ となる。

\begin{lemma}[被約基底での相対 Picard の一致判定]\label{lem:reduced-equalizer}
$X$ を被約スキーム、$Y\to X$ を相対 diagonal が immersion である代数空間とする。
二つの section $f,g:X\to Y$ が全幾何点上で一致すれば、$f=g$ である。
この判定は $Y=\Pic_{C/X}$ に適用できる。
\end{lemma}

\begin{proof}
equalizer $E\to X$ は diagonal の底変更なので immersion である。
仮定よりその像は $|X|$ 全体である。
従って局所的には $X$ の全ての点で零となる ideal による閉 immersion であり、
$X$ の被約性からその ideal は零である。
相対 Picard の diagonal は quasi-compact immersion である
\cite[Lemma 108.9.2]{StacksRelativePicard}。
相対 Picard 全体の separatedness は必要ない。
\end{proof}

\begin{lemma}[普遍 helix の相対的な差]\label{lem:relative-helix-difference}
$N=[\mathcal L_2\mathcal L_0^{-1}]$ は $P=\Pic^0_{C/S}$ の section であり、相対 Picard 群で
\[
 [\mathcal L_{i+2}\mathcal L_i^{-1}]=N
 \tag{PR-4}\label{eq:helix-difference}
\]
が全ての $i$ について成り立つ。
\end{lemma}

\begin{proof}
各幾何ファイバーの admissibility により、$\mathcal L_2$ と $\mathcal L_0$ の多重次数は等しい。
従って $N$ は $P$ に属する。
体上の elliptic helix は
$\mathcal L_i\mathcal L_{i+1}^{-1}\mathcal L_{i+2}^{-1}\mathcal L_{i+3}\simeq\OO_C$
を満たすので、隣接する二つの差
$[\mathcal L_{i+2}\mathcal L_i^{-1}]$ と $[\mathcal L_{i+3}\mathcal L_{i+1}^{-1}]$ は等しい。
全ての差は $N$ に等しく、被約な $S$ に補題~\ref{lem:reduced-equalizer} を適用して
\eqref{eq:helix-difference} を相対的に得る。
有限個の $i$ についてこの議論を行えば、以下の二周期の構成には十分である。
\end{proof}

\subsection{平方根の étale パラメータ空間}

\begin{lemma}[相対 Picard の二等分]\label{lem:picard-halves}
$P=\Pic^0_{C/S}$ は滑らかな相対次元一の可換群代数空間であり、
\[
 [-2]:P\longrightarrow P
\]
は étale かつ全射である。
従って
\[
 T=S\mathbin{\times}_{N,P,[-2]}P
 \tag{PR-5}\label{eq:T}
\]
は $S$ 上 étale 全射であり、普遍 section $a\in P(T)$ は
$a^{-2}=N$、すなわち $a^2=N^{-1}$ を満たす。
\end{lemma}

\begin{proof}
\eqref{eq:curve-sequence} と曲線の相対 Picard の一般定理により
$\Pic_{C/S}$ は代数空間で局所有限表示であり、相対次元一の曲線であるため滑らかである
\cite[Lemma 108.9.5]{StacksRelativePicard}。
相対 Picard は quasi-separated でもあるため、
\cite[Appendix A, Lemma A.1]{HolmesMolchoOrecchiaPoiret2023} を適用できる。
その fiberwise identity component $P$ は一意な開部分群である。
開部分と幾何ファイバーの恒等成分による特徴付けの一意性から、任意の底変更と両立する。
各ファイバーの接空間は $H^1(C_s,\OO_{C_s})$ なので相対次元は一である。

ここで多重次数零と identity component の一致は非被約ファイバーも含む。
体上で nilradical を $I\subset\OO_C$ とすると、標数零の有限な対数級数により
$1+I\simeq I$ が成り立つ。
$H^2(C,I)=0$ なので $\Pic C\to\Pic C_{\mathrm{red}}$ は全射であり、
その核は $H^1(C,I)$ から来る連結な加法群である。
被約曲線では正規化と conductor の単元の完全列により、
多重次数零の部分は、正規化の各成分の Jacobian の積の
連結な torus・unipotent 群による拡大である。
従って多重次数零の部分はちょうど連結成分 $\Pic^0(C)$ となる。
相対的には滑らかな群代数空間の fiberwise identity component の開部分群を取っている。

二倍写像の相対微分は群の単位元で $2\,\mathrm{id}$、他の点ではその平行移動である。
$2$ は可逆なので、滑らかな相対一次元空間の間の $[-2]$ は étale となる。
各幾何ファイバーでその像は開部分群である。
ファイバー $P_s$ は連結なので、この像は全体である。
従って $[-2]$ は全射である。
\eqref{eq:T} の主張は底変更から従う。
\end{proof}

$T$ は $k$ 上滑らかで被約である。
必要なら $T$ の étale scheme chart を取り、以下では $T$ 自体をスキームとして扱う。
命題~\ref{prop:ATV-relative-action} により、
普遍 Picard class $a$ から実際の曲線の自己同型
\[
 \sigma=\eta(a):C_T\longrightarrow C_T
 \tag{PR-6}\label{eq:sigma}
\]
が得られる。
ここでは $a$ を表す大域的な直線束を選ぶ必要はない。
群作用そのものが相対 Picard 空間上に定義されている。

\begin{lemma}[線束の比較と二周期]\label{lem:relative-beta}
$T$ の Zariski 開被覆上で同型
\[
 \varepsilon_i:\sigma^*\mathcal L_i\xrightarrow{\sim}\mathcal L_{i+2}
 \qquad(i=0,1,2,3)
 \tag{PR-7}\label{eq:epsilons}
\]
を選べ、それらが実際の代数同型
\[
 \beta:D_T\xrightarrow{\sim}D_T(2)
 \tag{PR-8}\label{eq:beta}
\]
を与える。
\end{lemma}

\begin{proof}
全幾何点で \eqref{eq:field-eta} と $\deg\mathcal L_i=2$ により
\[
 [\sigma^*\mathcal L_i]=[\mathcal L_i]a^{-2}=[\mathcal L_i]N=[\mathcal L_{i+2}]
\]
である。
被約な $T$ 上で補題~\ref{lem:reduced-equalizer} を適用すると、これは
相対 Picard 空間の section としての等式である。

従って $\sigma^*\mathcal L_i\otimes\mathcal L_{i+2}^{-1}$ は基底の直線束の引戻しである。
詳しくは、相対 Picard class が零なので fppf 局所に自明となる。
$\pi_*\OO_{C_T}=\OO_T$ が普遍的に成り立つことから、その自明化の変換関数は
基底から来るため、直線束そのものが基底へ降下する。
四つの比較直線束を Zariski 局所に自明化して\eqref{eq:epsilons} を得る。

$\pi_*\mathcal L_i$ は $V_i\otimes\OO_T$ と同一視される。
これは全ファイバーで $H^1(\mathcal L_i)=0$、$h^0(\mathcal L_i)=2$ であり、
普遍 evaluation との比較がファイバー上同型になることから従う。
\eqref{eq:epsilons} による切断の引戻しは可逆写像
$\beta_i:V_i\to V_{i+2}$ を与える。
三次関係は
\[
 R_i=\ker\left(V_i\otimes V_{i+1}\otimes V_{i+2}
   \longrightarrow\pi_*(\mathcal L_i\mathcal L_{i+1}\mathcal L_{i+2})\right)
\]
であるから、$\beta_0,\ldots,\beta_3$ は $R_0,R_1$ を $R_2,R_3$ へ移す。
従ってその tensor 積は $W_0$ を $W_2$ へ移す。
四周期により $\beta_{i+4}=\beta_i$ と延長する。
巡回を $c$ と書けば、対応する tensor 積の射について
\[
 c(\beta_0\otimes\beta_1\otimes\beta_2\otimes\beta_3)
  =(\beta_1\otimes\beta_2\otimes\beta_3\otimes\beta_0)c
\]
が成り立つ。
従って $W_0$ を $W_2$ へ移すことから、巡回された全ての $W_i$、
従って全ての三次関係を保つ。
逆写像も同じ性質を持ち、\eqref{eq:beta} を得る。
最初の隣接 pair に対する曲線の作用は、構成から $\sigma$ の作用である。
三点 correspondence の二つの隣接 pair への射影はスキーム同型であり、
代数同型による作用と可換する。
この自然性により、全ての添字で同じ $\sigma$ が誘導される。
従って四周期に延長した $\beta$ の平方も、切断上の $\sigma^{2*}$ と projectively 一致する。
\end{proof}

\subsection{平方と canonical 四周期の比較}

\begin{theorem}[相対的な適合 Serre 根]\label{thm:relative-compatible-root}
上の $T$ をさらに Zariski 局所化した基底上で、
\eqref{eq:beta} は頂点対角型の自己同型 $\delta$ に対し
\[
 \beta^2=\delta\theta
 \tag{PR-9}\label{eq:beta-square}
\]
を満たす。
従って tails 上では $(2)_\beta^2\simeq(4)_\theta$ である。
この同型は基底上に構成されており、任意の底変更後にも成り立つ。
特に非被約基底へ引き戻しても、相対 Serre 関手を持つ有限 foundation の比較の下で
\[
 s=(2)_\beta,\qquad s^2\simeq S^{-1}[2]
 \tag{PR-10}\label{eq:Serre-square}
\]
を得る。
\end{theorem}

\begin{proof}
\emph{(a) 四周期と Picard の符号.}
$\eta$ は可換群の作用なので、$T$ 上で厳密に
\[
 \sigma^2=\eta(a)^2=\eta(a^2)=\eta(N^{-1})
 \tag{PR-11}\label{eq:sigma-square}
\]
である。
ただしこの式だけで canonical $\theta$ との同定を済ませない。
その同定には次の tensor の比較を使う。

\emph{(b) smooth elliptic locus での tensor 比較.}
$S_{\mathrm{sm}}\subset S$ を $C/S$ が滑らかな曲線の locus とする。
これは非空開部分で、既約な $S$ において稠密である。
その上の各幾何点では、補題~\cite[\texttt{lem:elliptic-period-square}]{QuadricIntrinsic2026} の次の計算を使う。
原点を選び、$\mathcal L_i$ の次数二の因子和を $b_i$ と書くと
$b_{i+2}-b_i$ は一定である。
$\sigma$ は $2v=b_i-b_{i+2}$ を満たす平行移動 $p\mapsto p+v$ である。
$\mathcal L_{i+2}$ の二重被覆の involution を
$\iota_{i+2}(p)=b_{i+2}-p$ とする。

$w_i$ を一般点 $p$ の $\mathcal L_{i+2}$ の評価で縮約し、
$p'=\iota_{i+2}(p)$ を使って二つの rank-one summand に分ける。
隣接二因子の縮約の非零性により、その両係数は非零である。
$w_{i+2}$ が $w_i$ の二回巡回であることから、
$\theta_i$ による $\mathcal L_{i+4}$ の $p$ における評価の引戻しは
$\mathcal L_i$ の $p'$ における評価である。
後者は $\mathcal L_i$ の involution を使うと
\[
 \iota_i(p')=p+b_i-b_{i+2}=p+2v
\]
における評価と同じである。
従って $\theta_i$ と $\beta_{i+2}\beta_i$ は $V_i$ 上で projectively 等しい。
この比較は切断の引戻しという \eqref{eq:epsilons} の向きと一致する。

\emph{(c) 普遍基底上での延長.}
$T\to S$ は étale なので、$T_{\mathrm{sm}}$ は $T$ の各既約成分で稠密である。
実際、$T$ は滑らかであり、その各成分の像は $S$ の非空開部分を含む。
四つの projective linear map
\[
 [\beta_{i+2}\beta_i],\ [\theta_i]
 \quad(i\bmod4)
\]
は $\PGL(V_i,V_{i+4})$ という separated な空間への射である。
(b) により $T_{\mathrm{sm}}$ の全幾何点で一致し、
$T$ は被約なので、全 $T$ で一致する。
ここで被約なのは普遍根の基底 $T$ である。
任意の非被約な底変更先で、幾何ファイバーの一致だけから結論しているのではない。

\emph{(d) 頂点スカラーの明示.}
したがって一意な単元 $c_i\in\OO_T^\times$ があり、
\[
 \beta_{i+2}\beta_i=c_i\theta_i,
 \qquad c_{i+4}=c_i
\]
となる。
$u_0=1$、$u_{i+1}=c_i^{-1}u_i$ と全整数について再帰的に定める。
各 $u_i$ は有限積からなる単元である。
すると
\[
 \delta|_{D_{ij}}=\frac{u_i}{u_j}\,\mathrm{id}_{D_{ij}}
\]
は代数自己同型で、一次生成性により \eqref{eq:beta-square} が成り立つ。
右加群の $i$ 成分を $u_i$ 倍する射は、$\delta$ による引戻しと恒等関手の自然同型を与える。
従って tails 上の平方同型を得る。

\emph{(e) 基底上の Serre 双加群比較.}
定理~\ref{thm:relative-cubic-tails-tilting} により、連続四項の直和は
相対 compact tilting generator となる。
その有限 foundation 代数を $\Lambda$ とし、
$s^2$ と相対反 Serre 関手 $S^{-1}[2]$ を表す
可逆 perfect $\Lambda$-双加群を $P,Q$ とする。
$P$ は実際の次数移動 $s=(2)_\beta$ を標準 dg Morita 同値で移したものである。
$Q$ は $\Lambda$ の相対 Serre 双加群の逆と $[2]$ から得られる。
有限 foundation の相対 smooth proper 性
（補題~\ref{lem:foundation-relative-smooth-proper}）により、
これらは基底変更と整合する可逆 perfect 双加群である。

全ての幾何点で (a)--(d) が成り立つので、
\emph{ここで構成した同じ $\beta$} に
命題~\cite[\texttt{prop:admissible-period-h2}]{QuadricIntrinsic2026} を適用できる。
従って $P_{\bar t}\simeq Q_{\bar t}$ である。
これは各ファイバーに何らかの根が存在するという主張より強く、
普遍次数移動 $s$ の平方そのものの比較である。

任意の底変更後にも $Z(\Lambda)=\mathcal O_T$ である。
実際、中心元は四つの頂点で scalar であり、
隣接する忠実な有限局所自由 Hom 加群との可換性により全 scalar が等しい。
従って定理~\ref{thm:bimodule-line-comparison} を $P,Q$ に適用できる。
比較 locus $T_{P\sim Q}\hookrightarrow T$ は有限表示局所閉 immersion で、
全ての幾何点を含む。
$T$ は被約なので、補題~\ref{lem:reduced-equalizer} の proof と同じ ideal の議論により
$T_{P\sim Q}=T$ となる。
この locus の構成は、最低次 Hom の行列を基底環で零にする条件を含むので、
非被約基底に対しても有効な scheme-theoretic 比較である。

したがって比較線束
\[
 \ell=H^0\mathbf R\mathcal Hom_{\Lambda^e}(P,Q)
\]
と評価同型 $P\otimes\ell\simeq Q$ を得る。
さらに Zariski 局所に $\ell$ を自明化すれば
\eqref{eq:Serre-square} が成り立つ。
全ての射は普遍根の基底上に構成されているため、
任意の非被約基底への引戻し後も同じ同型が成り立つ。
\end{proof}



\begin{remark}[比較線束の局所自明化]
相対 Picard の等式からは、まず二つの直線束の差が基底から来ることが分かる。
\eqref{eq:epsilons} の前に四つの比較直線束を自明化したことが必要である。
その選択を変えると $\beta_i$ は基底単元倍されるが、
projective な平方比較は変わらず、(d) の頂点スカラーがその差を吸収する。
さらに (e) では Serre 双加群の比較線束 $\ell$ を自明化している。
従って幾何ファイバー上の平方根の存在だけから、
基底線束を省略した大域的な平方同型を主張してはいない。
必要なのは有限個の基底直線束の局所自明化である。
\end{remark}

\begin{remark}[四項 foundation と右加群の向き]
Abelian helix から $D_{ij}=\Hom(E_i,E_j)$ を回復する場合は、
右 tails の構成に添字を反転した反対代数
\[
 C_{ij}=D_{-j,-i},\qquad \beta_C=\beta_D^{-1}
\]
を用いる。
すると $E_i$ は $\pi e_{-i}C$ に対応し、四項窓の自己準同型代数は
元の foundation の積と一致する。
$\beta_D^2=\delta\theta_D$ を反転すれば、$C$ 上でも
$\beta_C^2=\theta_C$ が頂点スカラーを除いて成り立つ。
ここで $\theta_C$ は $\theta_D^{-1}$ に対応する canonical 四周期である。
(e) の双加群 $P,Q$ とその比較 locus は、この $C$ の右 tails 上で直接取れる。
従って Serre の比較と右軌道積の規約は同じ向きで成立する。
\end{remark}

\subsection{Artin 持上げと atlas 構成への入力}

\begin{corollary}[有限型の étale 根 chart]\label{cor:finite-type-root-chart}
$S$ を覆う有限個の affine scheme $T_\alpha$ と étale 射 $T_\alpha\to S$ を選び、
その各々で上の実際の二周期および Serre 平方同型を構成できる。
従って $\coprod_\alpha T_\alpha$ は有限型の滑らかな $k$ スキームであり、
全ての非線型 cubic geometric 五つ組の点を覆う。
\end{corollary}

\begin{proof}
$T\to S$ は étale 全射であり、その scheme chart の affine open の像は
$S$ を開被覆する。$S$ は準コンパクトなので有限個を選べる。
各 affine chart は $S$ 上有限型に取れる。
四つの Picard 比較直線束と Serre 比較線束 $\ell$ を
有限個の affine open で自明化し、必要な被覆を細分する。
\end{proof}

\begin{corollary}[指定した Picard 根の全 Artin 持上げ]\label{cor:all-artin-root}
$x\in S(k)$ と $a_x\in\Pic^0(C_x)$、$a_x^{-2}=N_x$ を固定する。
\eqref{eq:T} の中の $a_x$ に対応する点を含む affine étale scheme chart を選び、
その点を $t$ とする。必要なら前の有限被覆にこの chart を一枚追加する。
任意の局所 Artin $k$-代数 $A$ と、その閉点が $x$ となる framed geometric
五つ組 $\operatorname{Spec} A\to S$ に対し、$t$ を指定した $\operatorname{Spec} A\to T$ の持上げは一意に存在する。
対応する二周期は Serre 適合であり、任意の Artin 商・小拡大に整合する。
\end{corollary}

\begin{proof}
étale 射は nilpotent closed immersion に対して形式的 étale である。
閉点の持上げ $t$ から $\operatorname{Spec} A\to T$ が一意に得られる。
さらに選んだ Zariski 近傍は $t$ を含むので、$\operatorname{Spec} A$ 全体がその近傍へ入る。
そこで構成した \eqref{eq:beta-square} と \eqref{eq:Serre-square} を引き戻す。
\end{proof}

\paragraph{Abelian 変形と atlas への適用.}
相対根の構成自体は、非線型 locus 全体で特異・可約・非被約な曲線ファイバーを含む。
平坦 Abelian Artin 変形からの適用では、
補題~\cite[\texttt{lem:cubic-artin-overlap-recovery}]{QuadricIntrinsic2026} と命題~\cite[\texttt{prop:framed-cubic-artin-lift}]{QuadricIntrinsic2026} が
DDL による helix の持上げから、可逆部分束 $W_0$ とその overlap・縮約を含む
指定基底付き五つ組を回復する。
補題~\ref{lem:artin-two-heart-comparison} は、同じ四項 foundation を持つ二つの圏の
標準 dg Morita 同値が、指定された exact な還元を保って heart-exact となることを示す。
従って系~\ref{cor:all-artin-root} の根を元の Abelian 変形へ移せる。
この指定同型の持上げと有限型効果性を合わせた smooth atlas は
定理~\ref{thm:quintuple-chart-artin-lifting} と定理~\ref{thm:cubic-smooth-open} で構成する。
線型 cubic 点は全て pointed $(\operatorname{coh}(\PP^1\times\PP^1),\OO)$ に対応し、
既存の可換二次曲面の étale Morita-root chart を用いればよい。


% !TeX root = ../nc-quadric-surface-moduli.tex
% 2026-11-22 / v1.1（指定Artin同型のsmoothness判定を共有補題に統合）
\section{一般 cubic helix の持上げと cubic locus の滑らかな開性}
\label{sec:cubic-smooth-atlas}

\paragraph{主張の範囲と依存関係.}
本節では、一般 geometric cubic 点について、De Deken--Lowen の helix 持上げから
四つの隣接 Hom 空間と四次 syzygy の直線を回復する。
可換二次曲面上の Künneth 計算は用いず、指定された pointed 同型と Hom の基底を保って
五つ組の有限パラメータ空間へ持ち上げる。
相対 Picard 根と定理~\ref{thm:geometric-cubic-relative-effectivity} による実際の族を入力とし、
幾何点がちょうど cubic 対象である有限型で滑らかな Artin 開部分を得る。
noncubic 境界の相対認識と $\mathfrak H_2$ 全体の代数性は別の問題である。
本節の相対回復と atlas 構成の新規性は未確認である。

\subsection{五つ組と四項 foundation}

$k$ は標数零の代数閉体とする。
四つの二次元 $k$ 空間 $U_0,U_1,U_2,U_3$ を固定し、
\[
 S\subset\mathscr G\subset
 \mathbb P(U_0\otimes_kU_1\otimes_kU_2\otimes_kU_3)
\]
を geometric 五つ組の開部分 $\mathscr G$ の任意の開部分とする。
ここで $\mathbb P$ は直線をパラメータ化する規約であり、
$S$ 上の tautological 直線を $W$ と書く。
非線型 locus に制限する場合も、以下の持上げ補題はそのまま適用できる。

$S$ 上で縮約により
\[
 R_0=\operatorname{im}
 \bigl(W\otimes U_3^\vee\longrightarrow U_0\otimes U_1\otimes U_2\bigr)
\]
を定める。geometric 条件によりこれは階数二の部分束である。
五つ組の四項 foundation 代数 $\Lambda(W)$ は、頂点 $0,1,2,3$、
隣接矢空間 $U_0,U_1,U_2$ を持ち、三次関係が $R_0$ である代数である。
従ってその前向き成分は
\[
 \Lambda(W)_{ii}=\OO_S,\quad
 \Lambda(W)_{i,i+1}=U_i,\quad
 \Lambda(W)_{i,i+2}=U_i\otimes U_{i+1},\quad
 \Lambda(W)_{03}=(U_0\otimes U_1\otimes U_2)/R_0.
\]
全階数は $24$ である。積は path の連結であり、圏の Hom と比較するときは合成による積を用いる。
同じ右加群規約を両辺で固定する。

普遍 cubic 代数 $D(w)$ を前向きの $\Hom(E_i,E_j)$ の順で読む場合、
右 tails の入力は、添字を反転した反対代数
\begin{equation}\label{eq:quintuple-right-tails-convention}
 C_{ij}=D(w)_{-j,-i},\qquad E_i=\pi e_{-i}C
\end{equation}
である。定理~\ref{thm:relative-cubic-tails-tilting} の全構造 Hom の計算により
\[
 \Hom(E_i,E_j)=C_{-j,-i}=D(w)_{ij}\qquad(i\le j)
\]
となり、上の foundation と一致する。
実際の二周期 $\beta_D:D\to D(2)$ に対しては
$\beta_C=\beta_D^{-1}:C\to C(2)$ を用いる。
この変換は反転した canonical 四周期との平方比較を保つ。
相対族への昇格では、定理~\ref{thm:geometric-cubic-relative-effectivity} をこの $C$ に適用する。
偶数環と右軌道環の比較は定理~\ref{thm:quadric-finite-type-effectivity} の証明と同じで、
右軌道積は $a*b=s^{\deg b}(a)\circ b$ である。
$D$ の通常の右 tails と $C$ の右 tails を同一視することはしない。

\subsection{五つの対象からの有限次数回復}

\begin{lemma}[四次 overlap の普遍的回復]
\label{lem:cubic-artin-overlap-recovery}
$B$ を剰余体 $k$ の可換局所 Artin $k$ 代数とし、
$\mathcal D$ を $B$ 線形 Grothendieck 圏とする。
$E_0,\ldots,E_4$ は compact な $B$-flat 対象であり、
その導来 Hom は前向きに次数零へ集中すると仮定する。
また $(E_0,E_1,E_2,E_3)$ と $(E_1,E_2,E_3,E_4)$ は strong exceptional collections とする。
\[
 D_{ij}=\Hom_{\mathcal D}(E_i,E_j)\quad(0\le i\le j\le4)
\]
は有限自由で、剰余体上への基底変更と可換し、
その剰余体上の積は標準 geometric cubic 代数の対応する窓と一致するとする。
このとき、次の構成が成立する。
\begin{enumerate}
\item $V_i=D_{i,i+1}$ は階数二である。二重積は同型、三重積と四重積は全射であり、
\[
 \operatorname{rk}D_{i,i+r}=2,4,6,9\qquad(r=1,2,3,4)
\]
である（添字が $0,\ldots,4$ に入る範囲で読む）。
\item 三重積の核
\[
 R_i=\ker\bigl(V_i\otimes V_{i+1}\otimes V_{i+2}\longrightarrow D_{i,i+3}\bigr),
 \qquad i=0,1,
\]
は階数二の直和因子である。
\item $T_4=V_0\otimes V_1\otimes V_2\otimes V_3$ と置く。
\[
 W_0=(R_0\otimes V_3)\cap(V_0\otimes R_1)\subset T_4
\]
は階数一の直和因子であり、積による列
\begin{equation}\label{eq:artin-overlap-complex}
0\longrightarrow W_0\longrightarrow
 (R_0\otimes V_3)\oplus(V_0\otimes R_1)
 \longrightarrow T_4\longrightarrow D_{04}\longrightarrow0
\end{equation}
は分裂完全な有限自由複体である。
最初の射は $w\mapsto(w,-w)$、中央の射は二つの包含の和である。
\item 縮約は同型
\begin{equation}\label{eq:artin-contraction-isos}
 W_0\otimes V_3^\vee\xrightarrow{\sim}R_0,
 \qquad
 W_0\otimes V_0^\vee\xrightarrow{\sim}R_1
\end{equation}
を与える。従って $(V_0,V_1,V_2,V_3,W_0)$ は $B$ 上の geometric 五つ組である。
\end{enumerate}
これらの有限加群、核、交わり、縮約の構成は、任意の可換 $B$ 代数への
基底変更と可換する。
圏の Hom の基底変更と比較するときは、与えられた相対 compact 圏の基底変更同型を用いる。
\end{lemma}

\begin{proof}
特殊ファイバーの標準 cubic 代数には二次関係がなく、三次関係は階数二であり、
最初の五つの Hilbert 数は $1,2,4,6,9$ である。
標準 cubic 解消から、四次の関係は二つの三次関係の積の和で、
その交わりは一次元である。
これが本補題で用いる体上の入力であり、
補題~\ref{lem:relative-geometric-cubic-flatness} で用いた標準 cubic 解消の最初の四次に当たる
\cite[arXiv v3, Theorems 5.5.7--5.5.10]{VdB2011}。

二重積は同階数の有限自由加群間の射で、特殊ファイバーで同型なので同型である。
三重積と四重積の cokernel は有限加群であり、特殊ファイバーで零である。
Nakayama の補題により各積は全射である。
その標的は自由なので全射は分裂し、三次の核 $R_i$ は階数二の自由直和因子となる。
四次積の核を
\[
 K_0=\ker(T_4\longrightarrow D_{04})
\]
とすると、これも階数七の自由直和因子である。
ここまでの核の形成は分裂完全性により任意の基底変更と可換する。

合成の結合律により、二つの部分束 $R_0\otimes V_3$ と $V_0\otimes R_1$ は
$K_0$ に入る。従って自然な射
\[
 q:(R_0\otimes V_3)\oplus(V_0\otimes R_1)\longrightarrow K_0
\]
がある。特殊ファイバーの三次提示により $q\otimes_Bk$ は全射である。
再び Nakayama により $q$ は全射となり、標的の自由性により分裂する。
従って $\ker q$ は階数一の自由加群である。

二つの部分束の $T_4$ への包含は単射なので、
$\ker q$ を第一成分へ射影する射は単射で、その像は二つの部分束の交わりにちょうど等しい。
これにより $\ker q\simeq W_0$ を得る。
この射は剰余体上で非零の直線を与えるので、局所基底で少なくとも一つの係数は単元である。
従って $W_0\subset T_4$ 自体も直和因子である。
列~\eqref{eq:artin-overlap-complex} は、いま得た二つの分裂短完全列をつないだものとなる。

$W_0\subset R_0\otimes V_3$ と $W_0\subset V_0\otimes R_1$ を縮約すると
\eqref{eq:artin-contraction-isos} を得る。
両辺は階数二の自由加群であり、特殊ファイバーで標準 cubic 解消の非退化 pairing になる。
従って行列式は単元で、どちらも同型である。
基底を選んだ四つの $V_i$ の下で $W_0$ は射影空間の $B$ 点を定める。
その閉点は geometric open に属し、$\operatorname{Spec} B$ の位相空間は一点なので、$B$ 点全体が同じ開に入る。
これにより geometric 条件が成り立つ。

最後に、全ての核を分裂短完全列によって構成したので、任意の基底変更後も同じ核となる。
特に基底変更後の $\ker q$ も二つの部分束の交わりに同定される。
縮約同型と三次提示の同定も tensor により保たれる。
\end{proof}

\begin{corollary}[四項 foundation の完全な同定]
\label{cor:quintuple-four-window}
前補題で四つの基底同型 $U_i\otimes B\xrightarrow{\sim}V_i$ を選び、
それによる五つ組の $B$ 点を $w_B$ とする。
このとき合成は、頂点と全ての Hom の作用を保つ同型
\[
 \Lambda(w_B)\xrightarrow{\sim}
 \End_{\mathcal D}\left(\bigoplus_{i=0}^3E_i\right)
\]
を与える。
\end{corollary}

\begin{proof}
長さ一の成分は選んだ基底同型である。
長さ二の成分は補題の二重積の同型であり、長さ三の成分は三重積の全射の核が
$R_0=W_0\otimes V_3^\vee$ であることから得られる。
従って全ての頂点成分で同型であり、積は実際の合成から定めたため整合する。
四項 foundation の比較には、五つ組から再構成した無限 $\mathbb Z$ 代数と
全ての $\Hom(E_i,E_j)$ を先に同定する必要はない。
\end{proof}

\subsection{指定された還元同型と基底を持つ持上げ}

\begin{proposition}[一般 cubic 点での framed Artin 持上げ]
\label{prop:framed-cubic-artin-lift}
$A'\twoheadrightarrow A$ を剰余体 $k$ を持つ局所 Artin $k$ 代数の冪零拡大とする。
$\mathcal G_A$ は flat な $A$ 線形 Grothendieck 圏であり、
$A$-flat な対象の geometric $(4,3)$-helix $(E_{i,A})_{i\in\mathbb Z}$ を持つとする。
全ての前向き Hom は有限自由で基底変更と可換し、閉ファイバーは標準 geometric cubic helix であるとする。
特に各四項窓は $D(\mathcal G_A)$ の compact tilting generator である。

$\mathcal G_{A'}$ を $\mathcal G_A$ の flat Grothendieck 変形とする。
同時に指定された還元同型を持つ flat 構造対象 $\OO_{A'}$ があり、
$E_{0,A}=\OO_A$ と同定されているとする。
四つの基底同型
\[
 \varphi_{i,A}:U_i\otimes A\xrightarrow{\sim}
 \Hom(E_{i,A},E_{i+1,A}),\qquad i=0,1,2,3,
\]
を指定し、回復される五つ組の点を $w_A\in S(A)$ とする。
このとき次が存在する。
\begin{enumerate}
\item flat な geometric helix $(E'_{i})_{i\in\mathbb Z}$ と指定還元同型
  $A\otimes_{A'}E'_i\xrightarrow{\sim}E_{i,A}$。
  $E'_0=\OO_{A'}$ とでき、その還元同型も指定したものに取れる。
\item $\varphi_{i,A}$ を延長する基底同型
\[
 \varphi_{i,A'}:U_i\otimes A'\xrightarrow{\sim}\Hom(E'_i,E'_{i+1})
\]
と、これらから回復される $w_{A'}\in S(A')$。
その $A$ への還元は\emph{指定された点} $w_A$ に等しい。
\item 頂点と全ての積を保つ同型
\[
 \Lambda(w_{A'})\xrightarrow{\sim}\End\left(\bigoplus_{i=0}^3E'_i\right)
\]
であって、その還元は $\varphi_{i,A}$ と元の合成から定まる同型に一致するもの。
\end{enumerate}
$W_{0,A}$ の生成元も指定されていれば、その生成元を $W_{0,A'}$ の生成元へ持ち上げられる。
射影空間 $S$ の点の持上げはこの追加の生成元の選択を必要としない。
\end{proposition}

\begin{proof}
De Deken--Lowen の基底記法で $(R,k)=(A',A)$ と置く。
両環は coherent で、$A$ は $A'$ 上有限表示、核は冪零だから、
\cite[\S4.1]{DeDekenLowen2011} の仮定を満たす。
各 $E_{i,A}$ は flat な heart の対象なので有限 flat 次元を持つ。
従って \cite[Theorem 5.13(1)--(4)]{DeDekenLowen2011} を\emph{基底 $A$ に直接}適用できる。
この定理は flat 対象としての持上げ、compact generation、strong 性、
geometric 性、および前向き Hom の平坦性と還元同型を与える。

平坦な $A'$ 加群で剰余体上の次元が有限なら、加群は有限自由である。
実際、剰余体上の有限基底を持ち上げると、余核は $M=\mathfrak mM$ を満たし、
$\mathfrak m$ の冪零性から零となる。有限平坦加群は局所環上自由である。
よって全ての前向き Hom は元と同じ階数の有限自由加群であり、
$E'_0,\ldots,E'_4$ は補題~\ref{lem:cubic-artin-overlap-recovery} の仮定を満たす。

対象の持上げと指定還元を区別する。
平方零の一段では、各 $E_{i,A}$ について、核 $J=\ker(A'\to A)$ を用いた障害複体は
\[
 \mathbf R\Hom(E_{i,A},J\otimes_A^{\mathbf L}E_{i,A})
 \simeq J\otimes_A^{\mathbf L}\mathbf R\Hom(E_{i,A},E_{i,A})
 \simeq J[0]
\]
である（一般の冪零拡大では平方零段階へ分ける）。
従って対象の障害と持上げ差に当たる $\Ext^2,\Ext^1$ は零である。
持上げの\emph{同型類}は一意だが、同型そのものの一意性は用いない。
$\OO_{A'}$ は同じ例外対象の flat な持上げだから、$E'_0$ と同型である。
$\Aut(E'_0)=A'^\times\to A^\times=\Aut(E_{0,A})$ は全射なので、
比較同型を単元倍で調整し、所要の還元同型を持つ $E'_0=\OO_{A'}$ を選べる。

基底付き自由加群の持上げでは
$\mathrm{GL}_2(A')\to\mathrm{GL}_2(A)$ が全射である。
従って四つの $\varphi_{i,A}$ をそれぞれ $\varphi_{i,A'}$ へ延長できる。
補題~\ref{lem:cubic-artin-overlap-recovery} によって得た直線 $W'_{0}$ を
これらの基底で固定された $U_0\otimes\cdots\otimes U_3$ の部分束へ戻す。
その還元が $W_{0,A}$ となることは、核と交わりの基底変更の整合性による。
従って $w_{A'}$ は $w_A$ の実際の持上げである。
$S$ は開部分で、$\operatorname{Spec} A'$ と $\operatorname{Spec} A$ の位相空間は同じなので、
$w_{A'}$ は $S(A')$ に入る。
直線の生成元の持上げも単元の持上げにより可能である。
最後の End 代数の同定は系~\ref{cor:quintuple-four-window} による。

なお $E'_i$ の noetherian 性が必要な場合は、
$A'$ の極大イデアルによる有限フィルトレーションを用いる。
flat 性により各商は閉ファイバーの $E_i$ の有限直和であるから、
閉ファイバーの $E_i$ が noetherian なら $E'_i$ も noetherian である。
\end{proof}

\subsection{pointed exact 同型の還元を保つ比較}

\begin{lemma}[Artin 基底上の二つの heart の比較]
\label{lem:artin-two-heart-comparison}
$B$ は局所 Artin 環とし、$\mathcal G_1,\mathcal G_2$ を flat な $B$ 線形
Grothendieck 圏とする。それぞれ flat compact tilting generator $T_1,T_2$ を持ち、
その End 代数は $B$ 上有限自由で、Artin 商への基底変更と可換するとする。
各 Artin 商でもこれらの還元は compact tilting generators と仮定する。
頂点を保つ End 代数の同型から標準 dg Morita 同値
\[
 F:D(\mathcal G_1)\xrightarrow{\sim}D(\mathcal G_2)
\]
を定める。
その閉ファイバー $F_k$ が指定された exact 同値の導来化なら、$F$ も heart-exact である。
さらに $B\twoheadrightarrow A$ 上の還元 Morita 同値が
指定された exact 同値 $\alpha_A$ の導来化に一致すれば、
$F$ の heart への制限は $\alpha_A$ を延長する。
指定された構造対象とその比較同型を保つ場合は pointed 同値となる。
\end{lemma}

\begin{proof}
$\Lambda_r=\End(T_r)$ とし、両代数を指定同型で $\Lambda_B$ と同定する。
閉ファイバーの包含 $i_r:\mathcal G_{r,k}\to\mathcal G_r$ は
$\Lambda_B\to\Lambda_k$ に沿う restriction of scalars に対応する。
実際、$T_r$ の flat 性と導来随伴により
\[
 \mathbf R\Hom_{\mathcal G_r}(T_r,i_rM)
 \simeq\mathbf R\Hom_{\mathcal G_{r,k}}(T_{r,k},M)
\]
であり、作用は $\Lambda_k$ を通る。
標準 dg Morita の $F$ はこの同定を通じて恒等関手であるから、
$Fi_1\simeq i_2F_k$ となる。
同じ議論は $B\to A$ にも成立する。

$B$ の極大イデアルを $\mathfrak m$ とする。
任意の heart の対象 $M$ は有限フィルトレーション
\[
 M\supset\mathfrak mM\supset\cdots\supset\mathfrak m^NM=0
\]
を持ち、その商は閉ファイバーの heart の対象である。
$F_k$ の exact 性と上の可換同型により、各商の $F$ による像は heart に入る。
heart は拡大に閉じるから $F(M)$ も heart に入り、逆関手にも同じ議論が使える。
従って $F$ は完全同値を誘導する。
$A$ への還元の一致は、上の restriction of scalars との整合性と
$F_A\simeq D(\alpha_A)$ から従う。
これは一般の導来同値を無条件に heart-exact とする主張ではなく、
Artin の有限冪零フィルトレーションと指定された exact な還元を用いた結論である。
\end{proof}

\begin{theorem}[五つ組への étale chart の指定 Artin 持上げ]
\label{thm:quintuple-chart-artin-lifting}
$S\subset\mathscr G$ を上の開部分とし、$a:U\to S$ を有限型の étale 射とする。
実際の相対族 $q:U\to\mathfrak H_2$ が構成されており、任意の affine 基底変更に整合する
flat compact geometric cubic helix $(E_i)$、$E_0=\OO$ を持つとする。
さらに $a$ の普遍五つ組とこの helix の五つ組が同定され、
その四項窓の End 代数は頂点と積を保って $a^*\Lambda(W)$ と同定されるとする。
全ての幾何点で helix は標準 geometric cubic helix と仮定する。

$A'\twoheadrightarrow A$ を局所 Artin $k$ 代数の small extension とし、
$u_A\in U(A)$、$c_{A'}\in\mathfrak H_2(A')$、および\emph{指定された} pointed exact 同型
\[
 \alpha_A:q(u_A)\xrightarrow{\sim}c_{A'}|_A
\]
を取る。
このとき $u_A$ を延長する $u_{A'}\in U(A')$ と、$\alpha_A$ を延長する pointed exact 同型
\[
 \alpha_{A'}:q(u_{A'})\xrightarrow{\sim}c_{A'}
\]
が存在する。いずれの段階にも Hochschild 一次または二次上の根の作用の仮定は不要である。
\end{theorem}

\begin{proof}
$\alpha_A$ により普遍 helix とその Hom の基底を $c_{A'}|_A$ へ移す。
相対族の module-object 基底変更の公理から
$\operatorname{Ind}(c_{A'})$ は $\operatorname{Ind}(c_{A'}|_A)$ の flat Grothendieck 変形である。
従って命題~\ref{prop:framed-cubic-artin-lift} を適用でき、
$a(u_A)$ を延長する framed 五つ組 $w_{A'}\in S(A')$ を得る。
同時に、持ち上げた対象 $E'_0=\OO_{c_{A'}}$, $E'_1,E'_2,E'_3$ の End 代数は
$\Lambda(w_{A'})$ と同定される。

$a$ の formal étale 性を用いると、可換図式
\[
 \begin{matrix}
 \operatorname{Spec} A&\xrightarrow{u_A}&U\\
 \big\downarrow&&\big\downarrow a\\
 \operatorname{Spec} A'&\xrightarrow{w_{A'}}&S
 \end{matrix}
\]
の持上げ $u_{A'}$ が存在する。
この段階では新たな根や二周期を変形圏の内部で解く必要はない。
それらは、既に構成された $U$ 上の相対族を引き戻すことで供給される。

$q(u_{A'})$ の普遍四項窓と $(E'_0,E'_1,E'_2,E'_3)$ は、
同じ $\Lambda(w_{A'})$ を End 代数に持つ。
標準 dg Morita により、構造対象を保つ導来同値 $F_{A'}$ を得る。
その還元が\emph{指定された} $D(\alpha_A)$ と同型になることを確認する。
隣接 Hom に対する作用は、$\varphi_{i,A'}$ の還元を指定された $\varphi_{i,A}$ に取ったため一致する。
四項窓内の他の全 Hom は二重積の同型または三重積の全射により生成される。
従って対角、全ての頂点成分、および積の作用まで一致する。
標準 dg Morita 表示では、compact generator とその End 代数への作用が一致する二つの関手は
同じ双加群によって表され、自然同型である。
これにより $F_A\simeq D(\alpha_A)$ を得る。
構造対象上の比較は最初に指定した pointing へ合わせてある。

補題~\ref{lem:artin-two-heart-comparison} によって $F_{A'}$ は heart-exact である。
その制限が必要な $\alpha_{A'}$ となる。
Grothendieck 圏の $A'$ 線形完全同値は任意の代数の加群対象の圏にも完全同値を誘導する。
相対 $\mathfrak H_2$ 系はこの module-object 基底変更を満たすから、これは一つの heart だけでなく
全ての基底変更を含む相対系の同型である。
完全同値は noetherian 対象を保ち、標識も保つ。
従って相対 $\mathfrak H_2$ における指定された同型の持上げを得る。
\end{proof}

\subsection{smoothness の帰結と限界}

\begin{corollary}[全 closed cubic 点での smooth chart]
\label{cor:quintuple-effective-smooth-chart}
前定理の入力に加え、既存の相対認識から全忠実な射
\[
 \Phi:\mathfrak H_2\longrightarrow\mathfrak P=[X_3/\operatorname{PGL}_9]
\]
があり、$\mathfrak P$ の対角は affine かつ有限表示であるとする。
このとき $q:U\to\mathfrak H_2$ はスキームで表現可能な smooth 射である。
その fppf 像は $\mathfrak H_2$ の開部分スタックで、有限型で滑らかな Artin スタックである。
\end{corollary}

\begin{proof}
前定理は、$u_A$ の持上げだけでなく、
指定された pointed exact 同型 $\alpha_A$ の持上げも与える。
$\Phi$ の全忠実性と affine 有限表示対角を使うと、
補題~\ref{lem:pointed-artin-lifting-smooth-bridge} の仮定が成立する。
従って $q$ は表現可能な smooth 射で、その像は
$R=U\times_{\mathfrak P}U$ により表示された開部分スタックである。
$S$ は射影空間の開部分で $k$ 上 smooth、
$U\to S$ は étale なので、開像も smooth である。
\end{proof}

\subsection{全 cubic locus の有限型 smooth 開部分}

\begin{theorem}[cubic 対象全体の smooth Artin 開部分]
\label{thm:cubic-smooth-open}
$\mathfrak H_2$ には有限型で滑らかな Artin 開部分
\[
 \mathfrak V_{\mathrm{cub}}\subset\mathfrak H_2
\]
が存在し、その幾何点はちょうど pointed cubic 表示を持つ対象である。
より正確には、任意の局所有限型基底 $T$ 上の相対 $\mathfrak H_2$ 対象が
$\mathfrak V_{\mathrm{cub}}(T)$ に属するための必要十分条件は、
全ての幾何ファイバーが pointed cubic 表示を持つことである。

有限型の滑らかなスキーム $U_{\mathrm{cub}}$ と representable smooth 全射
\[
 q_{\mathrm{cub}}:U_{\mathrm{cub}}\longrightarrow\mathfrak V_{\mathrm{cub}}
\]
を取れる。その同値関係は全ての pointed exact 同値を用いる
\[
 R_{\mathrm{cub}}=U_{\mathrm{cub}}\times_{\mathfrak H_2}U_{\mathrm{cub}}
 \simeq U_{\mathrm{cub}}\times_{\mathfrak P}U_{\mathrm{cub}},\qquad
 \mathfrak V_{\mathrm{cub}}\simeq[U_{\mathrm{cub}}/R_{\mathrm{cub}}]
\]
であり、選んだ五つ組や根を保つ同値だけに限定しない。
\end{theorem}

\begin{proof}
\emph{非線型 cubic 点を覆う実際の族.}
$S_{\mathrm e}=\mathscr G\setminus\mathscr G_{\mathrm{lin}}$ とする。
系~\ref{cor:finite-type-root-chart} により、$S_{\mathrm e}$ を覆う有限個の
有限型 affine étale chart $T_\alpha\to S_{\mathrm e}$ を選べる。
各 chart には補題~\ref{lem:relative-beta} の実際の二周期と、
定理~\ref{thm:relative-compatible-root} の相対 Serre 平方同型がある。
式~\eqref{eq:quintuple-right-tails-convention} により反転した反対代数 $C$ と
$\beta_C=\beta_D^{-1}$ を使い、
定理~\ref{thm:geometric-cubic-relative-effectivity} を適用する。
これで実際の射
\[
 q_\alpha:T_\alpha\longrightarrow\mathfrak H_2
\]
が得られる。
普遍 cubic 解消、全 Hom の計算、canonical compact 比較により、
全ての四連続窓は flat compact 対象による full strong collection である。
二周期は $s(E_i)=E_{i+2}$ を与え、Serre 平方同型から $E_{i+4}\simeq S^{-1}E_i[2]$ を得る。
従って Serre と mutation の標準公式により geometric $(4,3)$-helix となる。
前向き Hom は全て有限 locally free で任意の基底変更と整合するから、
この族は定理~\ref{thm:quintuple-chart-artin-lifting} の全仮定を満たす。
従って各 $q_\alpha$ は系~\ref{cor:quintuple-effective-smooth-chart} により
representable smooth である。

\emph{線型点と全 cubic 幾何点の被覆.}
線型 geometric cubic 代数の pointed 圏は
$(\operatorname{coh}(\mathbb P^1\times\mathbb P^1),\OO)$ である。
この点には系~\ref{cor:quadric-effective-smooth-neighborhood} の有限型 smooth chart
$q_Q:U_Q\to\mathfrak H_2$ があり、$U_Q$ は滑らかである。
その構成では命題~\ref{prop:quadric-geometric-quintuple-open} の開部分へ制限しているので、
$U_Q$ の全ての幾何ファイバーも cubic である。

任意の非線型 pointed cubic 表示 $C_0$ に対し、
$D_0=C_0^{\mathrm{op,ref}}$ は geometric 五つ組から得られる cubic 代数である。
その五つ組の基底を選べば $S_{\mathrm e}$ の幾何点を得る。
$T_\alpha$ の被覆からその上の幾何点を取ると、構成された右 tails の代数は
$D_0^{\mathrm{op,ref}}=C_0$ に戻る。
従って構造対象も保った元の pointed 圏が $q_\alpha$ の像に入る。
これは定理~\cite[\texttt{thm:complete-h2-classification}]{QuadricIntrinsic2026} の cubic 側の回復を使うが、
noncubic 側も cubic 表示を持つとは仮定しない。

有限和
\[
 U_{\mathrm{cub}}=U_Q\amalg\coprod_\alpha T_\alpha
\]
とその射 $q_{\mathrm{cub}}$ を取る。
これは有限型で滑らかなスキームからの representable smooth 射であり、
その像の幾何点はちょうど cubic 対象である。
その fppf 像を $\mathfrak V_{\mathrm{cub}}$ と定める。
系~\ref{cor:quintuple-effective-smooth-chart} の像と groupoid の議論により、
これは有限型で滑らかな Artin 開部分であり、表示された quotient と一致する。

\emph{任意の基底上での内在的な記述.}
$T\to\mathfrak H_2$ を局所有限型の test とする。
その $\mathfrak V_{\mathrm{cub}}$ への逆像は、smooth 射
\[
 U_{\mathrm{cub}}\times_{\mathfrak H_2}T\longrightarrow T
\]
の開像である。
全ての幾何ファイバーが cubic なら、この射は幾何点上で全射である。
従って smooth 全射であり、fppf 局所に section を持つ。
これにより $T$ 上の相対対象は $\mathfrak V_{\mathrm{cub}}$ に属する。
逆向きは $U_{\mathrm{cub}}$ の全幾何ファイバーが cubic であることから従う。
\end{proof}

\begin{corollary}[全 cubic 点の形式変形と次元]
\label{cor:cubic-formal-h2-abelian}
$x=(\mathcal C,\OO)$ を $\mathfrak V_{\mathrm{cub}}(k)$ の点とする。
剰余体 $k$ の局所 Artin 基底上で、閉ファイバーの同定を持つ相対 $\mathfrak H_2$ 変形から
flat pointed Abel 変形への忘却は groupoid の同値である。
従ってこの形式変形問題は無障害であり、四項 foundation 代数を $\Lambda$ とすると、
その接空間と無限小自己同型空間はそれぞれ
\[
 \operatorname{HH}^2(\Lambda),\qquad\operatorname{HH}^1(\Lambda)
\]
である。さらに $\mathfrak V_{\mathrm{cub}}$ は純次元三である。
ここでは根・根の平方同型・AS4 被覆は標識として固定しない。
\end{corollary}

\begin{proof}
まず忘却の全忠実性を確認する。
相対系の affine test $\operatorname{Spec} B\to\operatorname{Spec} A$ での Grothendieck 圏は、
加群対象による基底変更の公理により、$A$ 上の圏における $B$ 加群対象の圏と同定される。
従って $A$ 線形の pointed 完全同値 $F_A$ は、任意の有限型 $A$ 代数 $B$ に対し、
$B$ の作用を $F_A$ で移すことで完全同値 $F_B$ に自然に延長する。
これらは代数準同型に沿う基底変更と可換し、自然同型についても同じ構成が働く。
完全同値は noetherian 対象を保ち、$B\otimes_A\OO_A$ を対応する構造対象へ送る。
非 affine test では相対系の fppf 降下により貼り合わさる。
よって一つの Artin heart における同値は全有限型 test 上の相対系の同値を定め、
相対系の射も $A$ 上の射から決まる。これが必要な全忠実性である。

非線型 cubic 点で本質的全射性を示す。
前定理で $x$ を表す点 $u\in T_\alpha(k)$ と、その pointed 比較を一つ選ぶ。
任意の flat pointed Abel 変形 $(\mathcal G_A,\OO_A)$ に、
命題~\ref{prop:framed-cubic-artin-lift} を $A\twoheadrightarrow k$ として適用する。
この命題の証明は、変形先が既に相対 $\mathfrak H_2$ 対象であることを必要としない。
従って閉ファイバーの指定基底と pointing を延長する $w_A\in S_{\mathrm e}(A)$ を得る。
$T_\alpha\to S_{\mathrm e}$ の étale 性により、$u$ を延長する $u_A\in T_\alpha(A)$ が存在する。
$q_\alpha(u_A)$ と $\mathcal G_A$ の四項窓は同じ End 代数を持つので、
補題~\ref{lem:artin-two-heart-comparison} による pointed 完全同値を得る。
従ってこの Abel 変形は実際の相対 $\mathfrak H_2$ 族 $q_\alpha(u_A)$ の忘却像に入る。
上の加群対象による自然延長により、比較は全有限型 test を含む相対系の比較となる。
線型点は可換二次曲面であり、既存の
定理~\ref{thm:quadric-all-artin-h2} と
系~\ref{cor:quadric-formal-deformation-equivalence} を用いる。

De Deken--Lowen の変形比較
\cite[Theorem 5.15]{DeDekenLowen2011} と
補題~\ref{lem:foundation-artin-deformation} により、pointed Abel 変形の接・無限小自己同型・障害は
それぞれ foundation の Hochschild 次数二・一・三で与えられる。
構造対象は例外的なので、その flat な持上げの同型類に追加の変形方向や障害はない。
また $\operatorname{HH}^0(\Lambda)=k\to\End(\OO)=k$ は恒等を送る同型であり、
pointing の scalar 選択は関手の scalar 自然同型で吸収される。
従って pointing による無限小自己同型空間の追加もない。
命題~\ref{prop:foundation-hochschild} の相対正規化 bar 複体の計算は、全ての cubic foundation について
\[
 \operatorname{HH}^n(\Lambda)=0\quad(n\ge3),\qquad
 \dim\operatorname{HH}^2(\Lambda)-\dim\operatorname{HH}^1(\Lambda)=3
\]
を与える。従って形式変形問題は無障害である。
$\mathfrak V_{\mathrm{cub}}$ は既に smooth Artin スタックなので、各幾何点での次元は
接空間の次元から無限小自己同型空間の次元を引いたものに等しく、三となる。
以上から純次元三を得る。
\end{proof}

\begin{remark}[noncubic 境界との区別]
$\mathfrak V_{\mathrm{cub}}$ の幾何点は cubic 対象だけであり、
定理~\cite[\texttt{thm:complete-h2-classification}]{QuadricIntrinsic2026} の noncubic 対象は含まない。
従ってこの定理は $\mathfrak H_2$ 全体の代数性や、noncubic 境界に沿う相対認識を結論しない。
また $\mathfrak V_{\mathrm{cub}}$ を $[X_3/\operatorname{PGL}_9]$ の開部分と同一視する主張でもない。
結論に必要なのは、相対 Picard 作用と二周期の実際の構成、同じ根の相対 AS4 被覆、
全 Noetherian 性、canonical tails 比較、そして本節の指定 Artin 同型の持上げを合わせたことである。
\end{remark}

% !TEX root = ../nc-quadric-moduli.tex
% 2026-11-22 / v1.4（不変式の全証明は相対モジュライ別稿へ移管）
\section{中央二重平方族の商スタックと慣性}
\label{sec:noncubic-stack}

$W=k u\oplus k v$ とし、非退化二次形式の開集合を
$Q^{\mathrm{nd}}=\operatorname{Sym}^2W\setminus V(\operatorname{disc})$ と書く。
\[
 P=Q^{\mathrm{nd}}\times Q^{\mathrm{nd}},\qquad
 \widetilde G=(\operatorname{GL}(W)\times\mathbb G_m^2)\rtimes S_2
\]
と置く。$S_2$ は二つの形式と二つのスカラーを交換し、恒等成分の作用は
\[
 (A,a_1,a_2)\cdot(f_1,f_2)
 =\bigl(a_1^{-2}\operatorname{Sym}^2(A)f_1,
        a_2^{-2}\operatorname{Sym}^2(A)f_2\bigr)
\]
である。これは環
\[
 B_f=k[u,v]\langle c,d\rangle/(c^2-f_1,d^2-f_2)
\]
の $W$ の基底変更、$c,d$ のスカラー変更と交換に対応する。
中央部分群
\[
 D=\{(\lambda I_W,\lambda,\lambda)\mid\lambda\in\mathbb G_m\}
\]
は $P$ に自明に作用する。$G=\widetilde G/D$ とする。
$G$ は $W\oplus kc\oplus kd$ の分解を保ち、後二因子を交換してよい
$\operatorname{PGL}_4$ の閉部分群であり、滑らかな五次元アフィン代数群である。
$G$ にも $c,d$ の符号変更から来る $\mu_2^2$ が全点の慣性として残る。
これは真の pointed 自己同値なので除去しない。
$D$ を除く根拠は、作用の非効果性だけでなく、以下で確認する tails 上の自然同型である。

\begin{proposition}[表示から構成するスタック]
\label{prop:noncubic-presentation-stack}
\[
 \mathfrak N^{\mathrm{pres}}=[P/G]
\]
は滑らかな有限型 Artin スタックであり、次元は一である。
$P$ 上の普遍次数付き環は全ての次数で有限自由である。
その幾何点の同型類は全 noncubic $\mathcal H_2$ 対象の pointed exact 同値類と一致する。
さらに幾何点での慣性は、pointed exact 自己同値の自然同型類の群と一致する。
ここで pointed 射は $(F,\phi:F\mathcal O\simeq\mathcal O')$、
二射は $\phi'\eta_{\mathcal O}=\phi$ を満たす自然同型とする。
\end{proposition}

\begin{proof}
$P$ は六次元の滑らかな開集合であり、$G$ は五次元の滑らかな代数群なので、
商スタックの一般論により最初の主張を得る
\cite[Tag 0DLS]{StacksModuli}。

普遍環は係数環上で $u^a v^b w$ を基底とする。
ここで $w$ は $c,d$ の交互語であり、空語も許す。
$c^2\mapsto f_1$, $d^2\mapsto f_2$ の書換えは長さを下げ、
$c^3,d^3$ の重なりは係数の中心性によって解消する。
従って基底は任意の底変更と整合し、各次数は有限自由である。

幾何点の同型類の対応には定理~\cite[\texttt{thm:noncubic-classification}]{QuadricIntrinsic2026} を用いる。
以下では既存の命題~\cite[\texttt{prop:first-orbit-object}]{QuadricIntrinsic2026}、
系~\cite[\texttt{cor:first-orbit-rank-invariance}]{QuadricIntrinsic2026} と
定理~\cite[\texttt{thm:ore-anticommuting-central}]{QuadricIntrinsic2026} を使って慣性を確認する。
最初の命題の内在的認識によって、任意の pointed exact 自己同値は
適合軌道対象の全列を同型を除いて保つ。
中央二重平方環 $B_f$ の通常のシフトを適合根であると主張する必要はない。
実際、定理~\cite[\texttt{thm:ore-anticommuting-central}]{QuadricIntrinsic2026} で構成した適合 Zhang twist は、
その添字付き環を $B_f$ の添字付き環と同定し、全自由シフトを対応させる。

\paragraph{曲線を点ごとに固定する変換.}
二つの根対が異なる場合には、一根共有の場合を含め、
$f_1,f_2,f_1f_2$ のいずれも $k(u,v)$ の平方ではない。
これは各根に対応する一次素因子の付値の偶奇を比較すれば従う。
従って $k(u,v)(\sqrt{f_1},\sqrt{f_2})$ は四次の体拡大である。
可換環
\[
 k[u,v,c,d]/(c^2-f_1,d^2-f_2)
\]
は $k[u,v]$ 上の基底 $1,c,d,cd$ を持つので、この体に単射し、整域である。
その斉次イデアルに一次式はないため、対応する整曲線 $E_f$ は $\mathbb P^3$ を張る。
射影変換を表す可逆行列の固定点集合は、その固有空間の射影化の有限和である。
$E_f$ が点ごとに固定されれば、その整性により一つの固有空間に全体が含まれる。
従って行列はスカラーであり、射影変換は恒等である。

根対が一致する場合には、座標と定数倍を選んで $f_1=f_2=uv$ とできる。
$E_f$ は平面 $d=c$ と $d=-c$ にある二つの滑らかな円錐曲線
\[
 C_+=\{d=c,\ c^2=uv\},\qquad
 C_-=\{d=-c,\ c^2=uv\}
\]
の和である。それぞれがその平面を張るため、点ごとに固定する行列の
二つの三次元線型部分空間への制限はスカラーとなる。
二つの部分空間は二次元空間 $\langle u,v\rangle$ で交わるので、
これらのスカラーは等しい。両部分空間の和は四次元全体であり、
この場合にも行列はスカラーである。

\paragraph{自己同値群への帰着.}
軌道対象を $L_i$ と書き、$F(L_i)\simeq L_i$ を選ぶ。
$i=0$ では与えられた標識同型を使う。
これらは添字付き Hom 環の自己同型を与える。
定理~\cite[\texttt{thm:ore-anticommuting-central}]{QuadricIntrinsic2026} の双線形化計算によると、
point correspondence の一次元部分は $E_f$ 上の対角であり、
残りは互いに移る二つの孤立点である。
従って隣接する二つの線型写像 $g_i,g_{i+1}$ は、射影化すると $E_f$ 上で一致する。
前段より $g_i=b_i g$ と書ける。
$L_i$ の同型を $t_i$ 倍すると隣接成分の作用は $t_{i+1}/t_i$ 倍されるため、
$t_0=1$ として $t_{i+1}/t_i=b_i^{-1}$ を全整数について逐次選べば、
全ての隣接成分の作用を同じ $g$ にできる。
一次生成性より全成分の作用が $g$ の誘導する作用と一致する。
Artin--Zhang 再構成をこの豊富な列に適用すると、これは $F$ と次数付き自己同型から
誘導される関手の pointed 自然同型を与える。
すなわち同型類の存在だけでなく、各自己同値そのものをこの方法で表せる。

次数付き自己同型の核については補題~\cite[\texttt{lem:scalar-tails}]{QuadricIntrinsic2026} の証明をそのまま使える。
その証明が使うのは $\operatorname{End}(L_i)=k$、一次生成性、
および非零の次数一元の左乗法が tails で零射でないことだけである。
これらは $B_f$ でも例外性、上の Zhang 同値と domain 性から成立する。
具体的には $(-)_\sigma\simeq\mathrm{Id}$ の $L_i$ 成分を、
基準同型 $b\mapsto\sigma^{-1}(b)$ の $c_i$ 倍と書くと、
$a\in B_1$ に対する自然性が
\[
 c_{i+1}\sigma^{-1}(a)=c_i a
\]
を与える。よって $\sigma$ は次数スカラーである。
逆向きは第 $n$ 次成分を $\lambda^{-n}$ 倍する自然同型で与えられ、
$i=0$ の基準同型と整合する。
異なる標識同型はスカラー自然同型で互いに同型となるため、余分な
$\mathbb G_m$ は残らない。
また恒等関手の自然自己同型は豊富な列上の定数スカラーで決まり、
$\mathcal O$ 上で恒等という条件はその値を一にする。
従って、標識と整合する二射を取った pointed groupoid の自己同型群は
\[
 \operatorname{Aut}_{\mathrm{gr}}(B_f)/D,
\]
二自己同型群は自明である。

\end{proof}

\begin{proposition}[根配置スタックへの射]
\label{prop:noncubic-sign-gerbe}
$\mathcal D=\mathbb P(\operatorname{Sym}^2W)\setminus V(\operatorname{disc})$ とし、
\[
 \mathfrak R=[(\mathcal D\times\mathcal D)/
       (\operatorname{PGL}(W)\times S_2)]
\]
と置く。自然な射 $\mathfrak N^{\mathrm{pres}}\to\mathfrak R$ は gerbe であり、
二つの形式を順序付けた局所表示上では相対慣性が $\mu_2^2$ である。
$S_2$ はこの二因子を交換する。
従って根配置の商スタック自体は pointed 慣性を保持しない。
\end{proposition}

\begin{proof}
$G\to\operatorname{PGL}(W)\times S_2$ の核は二次元トーラスであり、
形式の二つのスカラーに対して逆二乗により作用する。
その作用の核は $\mu_2^2$ で、$c,d$ の独立の符号変更を表す。
$P\to\mathcal D^2$ は二つの形式の代表を選ぶトーラス torsor である。
任意の根配置の同型は、射影変換の線型持上げと単元の平方根を選べば
局所的に $G$ の元へ持ち上がる。
標数零ではこれらの選択は fppf、実際には適当な\'etale 局所化で可能である。
従って対象と同型の両方が局所的に持ち上がり、gerbe の条件を満たす
\cite[Tag 06PD]{StacksModuli}。
\end{proof}

\begin{example}[一般慣性と Deligne--Mumford 性の限界]
四点が相異なる一般の場合、二点集合の非順序対を保つ射影変換群は
Klein 四群である。従って pointed 慣性の位数は十六である。
この一般性は、四点の交比の六つの値が相異なることから得られ、
二対を保つ順列は Klein 四群となる。
より具体的に $f_1=uv$, $f_2=(u-v)(u-\lambda v)$ とし、$r^2=\lambda$ を選ぶ。
行列は列 $(u,v,c,d)^{\mathsf t}$ への置換を表し、
例えば以下では $R(u)=\lambda v$, $R(v)=u$ である。
二つの自己同型 $R,T$ を
\[
 R|_W=\begin{pmatrix}0&\lambda\\1&0\end{pmatrix},\quad
 R(c)=rc,\quad R(d)=rd,
\]
\[
 T|_W=\begin{pmatrix}\lambda&-\lambda\\1&-\lambda\end{pmatrix},\quad
 T(c)=rd,\quad T(d)=r(\lambda-1)c
\]
で定める。置換行列の二乗はそれぞれ $\lambda I_4$ と $\lambda(\lambda-1)I_4$ であり、
次数スカラーを除くと $R^2=T^2=1$ である。
$c,d$ の符号変更を $e_1,e_2$ とすれば、
$R$ は両方と可換、$T$ は両者を交換し、$[R,T]=e_1e_2$ となる。
$\langle e_1,T\rangle$ は位数八の二面体群 $D_8$ で、$Re_1$ は独立の中央対合なので、
一般慣性は $D_8\times C_2$ である。
ここで $D_8$ は位数八の二面体群を表す。
係数恒等式とこれらの射影群関係は
\texttt{computations/verify\_stack\_symmetries.py} で厳密に検算できる。

一根共有の代表 $f_1=uv$, $f_2=u(u-v)$ では、
$u\mapsto u$, $v\mapsto u-v$, $c\leftrightarrow d$ と独立の符号変更により
慣性は $D_8$ となる。
一方、$f_1=f_2=uv$ では
\[
 u\mapsto t u,\quad v\mapsto t^{-1}v,\quad c\mapsto c,\quad d\mapsto d
 \qquad(t\in\mathbb G_m)
\]
が非自明な一次元部分群を与える。
従って $\mathfrak N^{\mathrm{pres}}$ 全体は Deligne--Mumford ではない。
二つの根対が異なる開部分では慣性は有限なので Deligne--Mumford となる。
\end{example}

\subsection{良モジュライ空間と粗モジュライ空間の障害}
\label{subsec:noncubic-good-moduli}

$f=a u^2+buv+c v^2$ と $g=a'u^2+b'uv+c'v^2$ に対し
\[
 \Delta(f)=b^2-4ac,\qquad
 D(f,g)=bb'-2ac'-2a'c,
\]
\[
 \chi(f,g)=\frac{D(f,g)^2}{\Delta(f)\Delta(g)}
\]
と置く。$D$ は判別式の極化であり、線型変換 $A$ の下で
$\Delta$ とともに $\det(A)^2$ 倍される。
$\chi$ は両形式の定数倍と交換の下でも不変なので、$P$ 上の正則な $G$-不変関数である。

\begin{proposition}
\label{prop:noncubic-good-moduli}
$k$ を標数零の代数閉体とする。このとき
\[
 k[P]^G=k[\chi].
\]
従って $\chi:\mathfrak N^{\mathrm{pres}}\to\mathbb A^1_\chi$ は
Alper の意味の good moduli space である。
一方、$\mathfrak N^{\mathrm{pres}}$ 全体は、幾何点の同型類を一対一に分類する
粗モジュライ代数空間を持たない。
\end{proposition}

\begin{proof}
不変式環 $k[P]^G=k[\chi]$ の完全な計算と
$\chi:\mathfrak N^{\mathrm{pres}}\to\mathbb A^1$ の
good moduli 性は既に本稿の命題~\ref{prop:moduli-noncubic-good-moduli} により証明した。
四根が相異なる通常形において
\[
 \chi\bigl(uv,(u-v)(u-\lambda v)\bigr)
 =\frac{(1+\lambda)^2}{(1-\lambda)^2}.
\]
一根共有の類と二根対が一致する類ではいずれも $\chi=1$ である。
これらは異なる幾何点であるが good moduli space 上では同じ点になる。
\cite[Theorem 6.6]{Alper2009} の規約内では、good moduli space の普遍性から、
代数空間への射はこの二点を分離できない。
任意の代数空間を許す場合も、例~\ref{ex:noncubic-specialization} の
穿孔直線上で fppf 局所的に定値な族と補題~\ref{lem:constant-punctured-line-space} により
同じ非分離が従う。従って幾何点の同型類に対して全単射となる
粗モジュライ代数空間は存在しない。
この議論を全内在的モジュライへ適用した形は、系~\ref{cor:global-h2-no-coarse-space} に置く。
\end{proof}

両者とその間の特殊化を保つ対象が $\mathfrak N^{\mathrm{pres}}$ であり、
$\mathbb A^1_\chi$ はその good moduli space である。
本節の不変式計算は表示スタック上で行ったものである。
定理~\ref{thm:noncubic-intrinsic-presentation-equivalence} により、元の有限型 site 上では
この表示商は内在的被約 noncubic 境界に同値となる。
その良モジュライ射の特殊 fiber と全慣性は第~\ref{sec:noncubic-boundary-comparison}節で記述する。

\begin{example}[共有根から根対一致への特殊化]
\label{ex:noncubic-specialization}
$\mathbb A^1_t$ 上で $f_1=uv$, $f_2=u(v-tu)$ とする。
両形式の判別式は常に一である。$t\ne0$ の全ての点は一根共有の軌道に属し、
$t=0$ では二つの根対が一致する。命題~\ref{prop:noncubic-presentation-stack} の基底により
対応する次数付き環は各次数で自由であり、これは $P$ 上の実際の平坦族である。
従って一根共有の軌道の閉包は、根対一致の軌道を含む。
また $\operatorname{Sym}^2\mathcal D$ を粗なスキームとして先に取ると、
一致した二対の $S_2$ 慣性も失うため、二対の交換もスタックとして割る必要がある。
\end{example}

\subsection{表示族における対合と正則切断の局所構成}
\label{subsec:noncubic-relative-root}

\begin{lemma}
\label{lem:noncubic-local-exchange}
標数零の体 $k$ 上のスキーム $S$ において、非退化な二元二次形式
$f_1,f_2$ の族が与えられたとする。
$S$ の \'etale 局所化の後、二形式を交換する線型対合 $L$ と単元 $\eta$ を選び、
\[
 L^2=1,\qquad Lf_1=\eta f_2,\qquad Lf_2=\eta^{-1}f_1
\]
とできる。さらに単元 $\gamma$、$\gamma^2=\eta$ を選べば
\[
 \sigma|_W=L,\qquad \sigma(c)=\gamma d,\qquad
 \sigma(d)=-\gamma^{-1}c
\]
は中央二重平方環の族の次数付き自己同型であり、
$\sigma^2$ は $W$ 上で恒等、$c,d$ 上で符号反転となる。
この構成は二つの根対が一致するファイバーの近傍でも有効である。
\end{lemma}

\begin{proof}
判別式が可逆なので、各形式の二つの根は \'etale 局所に順序付けられる。
第一根対を $\{0,\infty\}$ とする座標を選ぶ。
第二根対の一方を有限座標 $a$、他方を逆座標 $b$ で表示するように
根の順序と開集合を選べば、第二根対は $\{a,1/b\}$ と書ける。
ここで $b=0$ は第二根が無限遠にあることを表す。
二根が相異なる条件は
\[
 \theta=1-ab\in\mathcal O_S^\times
\]
である。従って、適当な単元 $\rho_1,\rho_2$ により
\[
 f_1=\rho_1uv,\qquad
 f_2=\rho_2 q,\qquad q=(u-av)(bu-v)
\]
と表示できる。この段階では形式の代表を平方根で正規化していない。
置換行列
\[
 A=\begin{pmatrix}1&-a\\ b&-1\end{pmatrix}
\]
について直接計算すると
\[
 A^2=\theta I_2,\qquad A(uv)=q,\qquad A(q)=\theta^2uv.
\]
$\delta^2=\theta$ を満たす単元 $\delta$ を \'etale 局所に選び、
\[
 L=\delta^{-1}A,\qquad \eta=\frac{\rho_1}{\rho_2\theta}
\]
と置けば所要の三つの恒等式を得る。
さらに $\gamma^2=\eta$ を選ぶ。
$\theta,\eta$ は単元であり、標数零なので、これらの平方根を付加する射は
有限 \'etale である。

$r_c=c^2-f_1$, $r_d=d^2-f_2$ と書くと、定義した $\sigma$ は
\[
 \sigma(r_c)=\eta r_d,\qquad
 \sigma(r_d)=\eta^{-1}r_c
\]
を満たすため商環へ降りる。$\sigma^2$ の式は定義から従う。
根対が一致する点では $a=b=0$ と選べ、このとき $\theta=1$、
$A=\operatorname{diag}(1,-1)$ である。従って一致境界で分母は消えない。
\end{proof}

\begin{proposition}[局所的な正則切断の族]
\label{prop:noncubic-relative-regular-cut}
前補題の選択を行ったアフィン底 $\operatorname{Spec}R$ 上で、
$u,v$ を中央とし、一つの関係 $r_c-\eta r_d=0$ だけを課した環を $A'$ と置く。
\[
 g'=r_c+\eta r_d
\]
は中央正則元であり、$A'/(g')\simeq B_f$ である。
$A'$、$B_f$ の各次数は有限自由 $R$-加群で、これらの記述と正則切断の完全列は
任意の底変更と整合する。
$\sigma$ は $A'$ に延びて $g'$ を固定するので、
Zhang twist 後にも
\[
 B_f^\sigma\simeq (A')^\sigma/(g')
\]
という平坦な正則正規切断の族を得る。
その各幾何ファイバーは定理~\cite[\texttt{thm:ore-anticommuting-central}]{QuadricIntrinsic2026} の適合表示となる。
\end{proposition}

\begin{proof}
$x=c+\gamma d$, $y=c-\gamma d$ と置く。$2\gamma$ は単元なので可逆な基底変更であり、
\[
 xy+yx=2H(u,v),\qquad H=f_1-\eta f_2
\]
を得る。従って $A'$ は $R[u,v,x]$ の Ore 拡大であり、
$\tau(u)=u$, $\tau(v)=v$, $\tau(x)=-x$、
$\partial(u)=\partial(v)=0$, $\partial(x)=2H$ を用いて
\[
 A'=R[u,v,x][y;\tau,\partial]
\]
と表示される。この定義から $u^a v^b x^i y^j$ が $R$ 上の基底となる。
また
\[
 g'=\frac{x^2+y^2}{2}-(f_1+\eta f_2)
\]
である。$x^2,y^2$ は中央であるから $g'$ は中央であり、
$y$ に関する最高次数の係数が可逆元 $1/2$ であるため、
$g'$ の乗法は $R$ に零因子がある場合にも単射である。
$r_c-\eta r_d$ と $r_c+\eta r_d$ を同時に課すと、
$2\eta$ の可逆性によって $r_c=r_d=0$ に等しくなる。
商 $B_f$ の $R$-基底は $u^a v^b$ と $c,d$ の交互語との積である。
従って各次数の有限自由性と任意の底変更に対する正則性が従う。

前補題の関係から
$\sigma(r_c-\eta r_d)=-(r_c-\eta r_d)$、$\sigma(g')=g'$ であり、
$\sigma$ は $A'$ 上でも定義される。
Zhang twist は底上の次数付き加群を変えない。
$g'$ は twist 後には正規元で、掛け算は元の $g'$ の乗法と
自己同型 $\sigma^2$ との合成であるため正則性も保つ。
各幾何ファイバーでの AS 正則性、型 $(1,4,6,4,1)$ と適合根の性質は
定理~\cite[\texttt{thm:ore-anticommuting-central}]{QuadricIntrinsic2026} の Ore 表示と Nakayama 計算に帰着する。
\end{proof}

\begin{remark}[相対主張の範囲]
前二結果は、明示された二形式の族から、底変更と整合する次数付き環、
正則切断、およびその幾何ファイバー上の適合根を局所的に構成する。
左右 noetherian 性は noetherian 基底では後続の
命題~\ref{prop:noncubic-relative-noetherian} で証明するが、
相対 Serre 関手は後の Clifford trace により計算し、
canonical compact 圏を含む全相対条件は
第\cite[\texttt{sec:relative-noncubic-families}]{QuadricIntrinsic2026}節で証明する。また、任意の内在的 $\mathcal H_2$ 族がこの正規形へ
\'etale 局所に還元できるという認識定理も含めていない。
局所的な平方根の存在と、内在的な族の相対分類を区別する必要がある。
\end{remark}

\begin{remark}[相対版として未証明の部分]
上で無条件に構成したのは、普遍中央二重平方環とその射影化した座標変換からなる
表示スタックである。普遍環は $P$ 上、あるいは $[P/\widetilde G]$ 上にある。
$D$ は次数 $n$ に重み $n$ で作用するため、同じ環束が剛体化後へ
そのまま降下するとは主張しない。幾何点の同型類と慣性は正しいが、任意の基底上の
$\mathcal O$-marked 圏の族との同値を主張するには、相対再構成と相対正規形が必要である。
特に「全ての幾何ファイバーが noncubic」という条件だけでは、
階数四を検出する小行列式の冪零成分が消えるとは限らない。
例えば双対数上の行列は全幾何ファイバーで階数三でも、四次小行列式が
非零の冪零元となり得る。階数三の閉部分スタックを定めるなら、
その小行列式をスキーム論的に零とする条件が必要である。
このため、点の分類をそのまま全 family の分類へ読み替えない。
\end{remark}



\subsection{平方根の選択に依存しない反 Serre 環の相対表示}
\label{subsec:noncubic-antiserre-descent}

ここでは表示スタックの上に、適合平方根を族として選ばずに
反 Serre 切断環を復元する。これは単なる幾何点での同値類の記述より強いが、
内在的な相対 Abel 圏の再構成とは区別する。
$\widetilde G$、$D$、$G$、$P$ は本節の冒頭の記号とする。
$P$ 上の普遍二重平方環を $\mathcal B$ とし、
\[
 \iota(u)=u,\quad \iota(v)=v,\quad
 \iota(c)=-c,\quad\iota(d)=-d
\]
と置く。これは全ての基底変更および二形式の交換と可換する
次数付き代数自己同型である。
$W$ は $\widetilde G$ の標準二次元表現であり、
$\mathcal D=\det W$ を $P$ 上の同変線束とする。

\begin{proposition}[noncubic 反 Serre 環の相対降下]
\label{prop:noncubic-antiserre-descent}
$P$ 上で
\[
 \mathcal R_n=\mathcal B_{2n}\otimes\mathcal D^{-n}
 \quad(n\ge0),\qquad
 (a\otimes\ell^{-n})\star(b\otimes\ell^{-m})
 =a\,\iota^n(b)\otimes\ell^{-(n+m)}
 \tag{NC-AS}\label{eq:noncubic-antiserre-product}
\]
と定める。線束因子の積は自然なテンソル積と解釈する。
すると以下が成立する。
\begin{enumerate}
\item $\mathcal R=\bigoplus_{n\ge0}\mathcal R_n$ は結合的な連結
  次数付き代数であり、任意の底変更と両立する。
  全次数で有限局所自由で
  \[
    \operatorname{rank}\mathcal R_n=(2n+1)^2.
  \]
\item $\mathcal R$ の $\widetilde G$-線型化では $D$ が自明に作用する。
  従って $\mathcal R$ は rigidification
  $\mathfrak N^{\mathrm{pres}}=[P/G]$ 上の
  通常の準連接次数付き代数として降下する。
\item 任意の幾何点において、$\mathcal D$ の線維を自明化した
  $\mathcal R$ は、対応する $\HH$ 対象の内在的反 Serre 環
  \[
    \bigoplus_{n\ge0}
       \operatorname{Hom}(\mathcal O,(S^{-1}[2])^n\mathcal O)
  \]
  と次数付き代数として同型である。
  この主張では適合平方根そのものを追加標識として選ばない。
\end{enumerate}
\end{proposition}
\begin{proof}
$\mathcal B$ の基底は $u^av^bw$（$w$ は $c,d$ の交互語）であり、
その Hilbert 級数は
\[
 H_{\mathcal B}(z)
 =\frac1{(1-z)^2}\left(1+\frac{2z}{1-z}\right)
 =\frac{1+z}{(1-z)^3}.
\]
従って第 $j$ 次成分の階数は $(j+1)^2$ である。
$\iota$ は代数自己同型なので $(\mathrm{NC\text{-}AS})$ の
積の結合律は
\[
 (a\iota^n(b))\iota^{n+m}(c)
 =a\iota^n(b\iota^m(c))
\]
から従う。すべては普遍的な基底と積により定義され、
非被約基底を含む任意の底変更と両立する。

$\widetilde G$ は $\mathcal B_j$ に次数付きに作用し、$\iota$ と可換する。
$D$ の元 $(\lambda I_W,\lambda,\lambda)$ は
$\mathcal B_{2n}$ に重み $2n$、$\det W$ に重み $2$ で作用する。
従って $\mathcal B_{2n}\otimes(\det W)^{-n}$ 上の
$D$-作用は自明であり、積も $\widetilde G$-同変である。
$D$ が $P$ 上でも自明に作用するので、
$\widetilde G$-同変準連接層のうち $D$ が自明に作用するものは
$G=\widetilde G/D$-同変層へ降下する。

最後に、補題~\ref{lem:noncubic-local-exchange} により
\'etale 局所では適合する次数付き自己同型
$\sigma$ を選べて、$\sigma^2=\iota$ である。
通常の Zhang twist の積を
$a*_\sigma b=a\sigma^{\deg a}(b)$ と書くと、
偶数次数成分上で
\[
 a*_\sigma b=a\,\sigma^{2n}(b)
              =a\,\iota^n(b)
       \qquad(a\in\mathcal B_{2n})
\]
となる。従って
\[
 (\mathcal B^\sigma)^{(2)}
 \simeq\bigl(\mathcal B^{(2)},\star\bigr)
\]
は $\sigma$ の選択によらない。
幾何点では $\mathcal B^\sigma$ が適合する
$(\mathcal O,s)$ の全切断環であり、
$s^2\simeq S^{-1}[2]$ なのでその偶数部分は
内在的反 Serre 環そのものである。
$\mathcal D$ の線維の自明化は各次数の共通次数スカラーを変えるだけで、
幾何点の代数同型類には影響しない。
\end{proof}

\begin{proposition}[九生成元の相対二次関係と商スタックへの射]
\label{prop:noncubic-antiserre-grassmann}
上の $\mathfrak N^{\mathrm{pres}}$ 上の代数 $\mathcal R$ は
次数一で生成される。乗法
\[
 \mathcal R_1\otimes\mathcal R_1
    \longrightarrow\mathcal R_2
\]
は任意の底変更後も全射であり、その核 $\mathcal J$ は
階数 $56$ の部分束である。
従って標準九次元空間 $U$ を固定すれば、自然な射
\[
 \Psi_{\mathrm{nc}}:
 \mathfrak N^{\mathrm{pres}}
 \longrightarrow
 \left[\operatorname{Gr}(56,U\otimes U)/\operatorname{PGL}(U)\right]
 \tag{NC-Gr}\label{eq:noncubic-antiserre-grassmann}
\]
を得る。幾何点上ではこの射は反 Serre 環の
二次関係空間を与えるので、全体の
反 Serre 商スタック候補の noncubic chart を具体化する。
\end{proposition}
\begin{proof}
$\mathcal B$ は次数一生成である。
任意の長さ $2n$ の積を隣接する二文字ずつに括り、
偶数部分の $n$ 個の元 $a_1,\ldots,a_n$ と書く。
$\iota$ は可逆なので
\[
 a_1\star\iota^{-1}(a_2)\star
       \iota^{-2}(a_3)\star\cdots\star
       \iota^{-(n-1)}(a_n)
 =a_1a_2\cdots a_n .
\]
よって $\mathcal R$ は次数一生成である。
特に $\mathcal R_1^{\otimes2}\twoheadrightarrow\mathcal R_2$ は
全射であり、両辺の階数は $81,25$ だから
$\mathcal J$ は階数 $56$ の局所直和因子である。
この結論は各次数の局所自由性と積の任意底変更との整合性から、
全ての基底上で成立する。

$\mathcal R_1$ の射影枠束は $\operatorname{PGL}(U)$-torsor である。
$\mathcal J\subset\mathcal R_1^{\otimes2}$ はその枠に関して
同変な Grassmannian の点を与えるので
(\ref{eq:noncubic-antiserre-grassmann}) を得る。
各幾何点での識別は直前の命題による。
この射の存在は、その像が内在的な相対 $\HH$ locus
として開・閉・有限表示であることまでは主張しない。
\end{proof}

\begin{corollary}[noncubic の全反 Serre 環の相対二次提示]
\label{cor:noncubic-relative-full-quadratic-presentation}
命題~\ref{prop:noncubic-antiserre-grassmann} の
九生成元からの自然な射は、全次数で同型
\[
 T_{\mathcal O}(\mathcal R_1)/(\mathcal J)
       \xrightarrow{\ \sim\ }\mathcal R,
 \qquad
 \mathcal J=\ker(\mathcal R_1^{\otimes2}\to\mathcal R_2)
 \tag{NC-quad}\label{eq:noncubic-relative-quadratic-presentation}
\]
である。この同型は任意の基底変更と両立する。
従って射 $\Psi_{\mathrm{nc}}$ は、
二次関係だけで表される反 Serre 環の
\emph{全次数}の平坦な族を与えている。
\end{corollary}
\begin{proof}
$\mathcal B$ は次数一生成かつ二次関係で定義されている。
$\mathcal R$ はその偶数 Veronese を
次数自己同型 $\iota$ で twist し、
線束 $\det W$ の次数付き冪で捩ったものなので、
補題~\cite[\texttt{lem:relative-quadratic-veronese}]{QuadricIntrinsic2026} が適用できる。
線束は局所的に自明化でき、二次表示の同型は
自明化の変更と整合するため、表示スタック全体へ降下する。
$\mathcal B_4$ は任意基底上で自由であり、
$\mathcal J$ の底変更も成立する。
\end{proof}

\begin{proposition}[noncubic Grassmann 射の明示的な affine cell]
\label{prop:noncubic-grassmann-explicit-cell}
$W=ku\oplus kv$ の基底を固定し、
$\det W$ を $u\wedge v$ で自明化する。
二形式 $(f_1,f_2)$ の係数に関する六次元アフィン空間
\[
 \mathbb A^6=\operatorname{Sym}^2W\oplus\operatorname{Sym}^2W
\]
全体まで $\mathcal R$ とその二次関係を延長する。
次数一の九生成元を
\[
 \begin{gathered}
 x_0=u^2,\quad x_1=uv,\quad x_2=v^2,\\
 p_0=uc,\quad p_1=vc,\quad q_0=ud,\quad q_1=vd,\quad z=cd,\quad z'=dc
 \end{gathered}
 \tag{NC-9}\label{eq:noncubic-nine-generators}
\]
と置く（右辺はすべて $\mathcal B_2$ の元で、
$\mathcal R$ の積には $\star$ を用いる）。
以下の二十五個の積
\[
\begin{gathered}
 x_0\star x_0,\ x_0\star x_1,\ x_0\star x_2,\
 x_1\star x_2,\ x_2\star x_2,\\
 x_0\star p_0,\ x_0\star p_1,\
 x_1\star p_1,\ x_2\star p_1,\quad
 x_0\star q_0,\ x_0\star q_1,\
 x_1\star q_1,\ x_2\star q_1,\\
 x_i\star z,\ x_i\star z'\quad(i=0,1,2),\\
 p_0\star z',\ p_1\star z',\
 q_0\star z,\ q_1\star z,\quad
 z\star z,\ z'\star z'
\end{gathered}
\tag{NC-25}\label{eq:noncubic-25-pivots}
\]
は、どの $(f_1,f_2)$ 上でも $\mathcal R_2$ の基底となる。
従って積写像 $\mathcal R_1^{\otimes2}\to\mathcal R_2$ の
固定した $25\times25$ minor は $\pm1$ である。

この minor に対応する
$\operatorname{Gr}(56,81)$ の affine cell を
$\mathbb A^{25\cdot56}$ と同定すれば、
\[
 \Phi:\mathbb A^6\longrightarrow\mathbb A^{25\cdot56}
 \tag{NC-cell}\label{eq:noncubic-cell-map}
\]
は全係数環にわたり多項式であり、
実際には\emph{閉埋込み}である。
従って非退化二形式の開集合 $P\subset\mathbb A^6$ は、
固定された同じ cell への局所閉埋込みとして実現される。
ここで述べる埋込みは\emph{基底を固定した表示空間}の埋込みであり、
その群商が内在的な $\HHstack$ へ閉埋込みとなるとは
主張しない。
\end{proposition}
\begin{proof}
$\mathcal B$ の $u^av^bw$（$w$ は $c,d$ の交互語）
という単項式基底は、二形式が退化しても成立する。
次数四の二十五個の基底を、交互語の長さ
$0,1,2,3,4$ ごとに並べると、それぞれ
\[
 5,\quad 8,\quad 6,\quad4,\quad2
\]
個である。
式~(\ref{eq:noncubic-25-pivots}) の積は、
この順にちょうど全基底単項式を列挙する。
右因子が $p_i,q_i$ の場合だけ
$\iota$ によって負符号が付き、それ以外は正符号となる。
どの積にも $c^2$ や $d^2$ の簡約は生じず、
したがってこの $25$ 列の行列は、
上記単項式基底で符号付き置換行列である。
その行列式は $\pm1$ で、二形式の係数に依存しない。

この基底を使えば、残り $56$ 個の二語積を
$c^2=f_1,d^2=f_2$ によって簡約した係数を
$\mathbb A^{25\cdot56}$ の座標として取れる。
簡約はモニックなので、各座標は二形式の六係数の
多項式である。従って $\Phi$ は多項式射である。

閉埋込み性には、六係数をこの座標から回復すればよい。
\[
 f_1=a u^2+buv+c v^2,\qquad
 f_2=a' u^2+b'uv+c'v^2
\]
と置くと
\[
 p_0\star p_0=-u^2f_1
   =-a u^4-bu^3v-cu^2v^2,\qquad
 q_0\star q_0=-u^2f_2
   =-a'u^4-b'u^3v-c'u^2v^2.
\]
右辺の三基底単項式は
$x_0\star x_0,x_0\star x_1,x_0\star x_2$
に対応する固定 pivot である。
従って $\Phi$ の行列座標のうち六つが
$-a,-b,-c,-a',-b',-c'$ となる。
この六座標の射影
$\mathbb A^{25\cdot56}\to\mathbb A^6$
と $\Phi$ の合成は、符号反転を除き恒等である。
従って $\Phi$ はアフィン空間間の多項式写像のグラフとして
閉埋込みになる。$P$ は判別式の非消滅開であるため、
その制限は局所閉埋込みである。
\end{proof}

\subsection{noncubic の Koszul 性と一次変形の全次数平坦性}
\label{subsec:noncubic-koszul-infinitesimal}

ここでは三次 Hilbert 条件に対する接計算を、
全次数の一次変形と Hochschild コホモロジーへ結び付ける。
この結び付きは Koszul 変形原理の一次版であり、
高次の変形を自動的に積分するという意味ではない。

\begin{proposition}[noncubic 反 Serre 環の Koszul 性]
\label{prop:noncubic-antiserre-koszul}
標数零の代数閉体 $k$ 上で、任意の二元二次形式
$f_1,f_2\in k[u,v]_2$ に対し
\[
 B=k[u,v]\langle c,d\rangle/(c^2-f_1,d^2-f_2),
 \qquad
 R=\bigl(B^{(2)}\bigr)^\iota,\quad
 \iota|_{k[u,v]}=1,\quad \iota(c)=-c,\ \iota(d)=-d
\]
と置く。ここで $\iota$ は偶数 Veronese に作用する
次数を保つ自己同型とし、最後の上添字は Zhang twist を表す。
このとき $B$ および $R$ は Koszul 代数である。
特に本稿の全 noncubic $\HH$ 対象の反 Serre 環は Koszul である。
\end{proposition}
\begin{proof}
$u,v,c,d$ の自由結合代数において
\[
 vu\mapsto uv,\quad
 cu\mapsto uc,\quad cv\mapsto vc,\quad
 du\mapsto ud,\quad dv\mapsto vd,\quad
 c^2\mapsto f_1,\quad d^2\mapsto f_2
\]
という二次書換え系を取る。
次数辞書式順序 $d>c>v>u$ で左辺は先頭項となる。
重複のうち $c^3,d^3$ については、
$f_1,f_2$ が中央の $u,v$ の多項式なので
左右の簡約は同じ結果になる。
その他の長さ三の重複は $u,v$ を左へ移す規則と
二乗規則の交換、または $u,v$ の可換性の整合であり、
いずれも同じ正規形へ簡約される。
従って Diamond Lemma により
$u^av^bw$（$w$ は $c,d$ の交互語）が基底となり、
表示は二次 Gröbner (PBW) 基底を持つ。
二次 PBW 代数の Koszul 性は
Priddy の定理
\cite[Chapter 4, Theorem 3.1]{PolishchukPositselski2005}
によるから $B$ は Koszul である。

Koszul 代数の Veronese 部分環も Koszul であることは
\cite[Chapter 3, Proposition 2.2(i)]{PolishchukPositselski2005}
により、$B^{(2)}$ は Koszul となる。
Zhang twist は次数付き加群圏と線型自由解消を移すので
Koszul 性を保ち、$R$ も Koszul である。
noncubic の反 Serre 環との同一視は
命題~\ref{prop:noncubic-antiserre-descent} による。
\end{proof}

\begin{lemma}[Koszul 二次環の一次変形における三次判定]
\label{lem:koszul-first-order-degree-three}
$R_0=T_k(V)/(J)$ を局所有限な連結 Koszul 二次代数、
$A=k[\varepsilon]/(\varepsilon^2)$ とする。
$J_A\subset V_A^{\otimes2}$ を $J$ に特殊化する
階数 $\dim J$ の直和因子とし、
$R_A=T_A(V_A)/(J_A)$ とする。
このとき以下は同値である：
\begin{enumerate}
\item $(R_A)_3$ が $A$ 上平坦である。
\item 全ての $n\ge0$ で $(R_A)_n$ が $A$ 上平坦である。
\end{enumerate}
従って三次 Hilbert 関数の determinantal locus の
$R_0$ における接空間は、
\emph{全次数平坦な一次の次数付き二次代数変形}を
正確に記述する。
これは双対数上の結論であり、二次・三次以降の
Artin 小拡大への持上げは主張しない。
\end{lemma}
\begin{proof}
$E=V^{\otimes2}/J$ とし、$J_A$ を
$j-\varepsilon\widetilde\phi(j)$ の張るグラフとして表す。
ここで $\phi\in\operatorname{Hom}_k(J,E)$ である。
二次関係の三次における重なりは
\[
 K_3=(V\otimes J)\cap(J\otimes V)\subset V^{\otimes3}
\]
で測られる。
$R_0$ の Koszul 双加群複体の内部次数零の
Hochschild コチェイン複体の最初の項は
\[
 C^1=\operatorname{Hom}(V,R_{0,1}),\quad
 C^2=\operatorname{Hom}(J,R_{0,2}),\quad
 C^3=\operatorname{Hom}(K_3,R_{0,3}).
 \tag{Kos-C}\label{eq:noncubic-koszul-hoch-cochain}
\]
三次関係を $A$ 上で持ち上げたときの整合性は、
$\phi$ の微分
$d^2\phi\in\operatorname{Hom}(K_3,R_{0,3})$
が零になる条件である。
実際、一つの三次重なりを
$V\otimes J$ と $J\otimes V$ の二通りに書き、
関係式 $j-\varepsilon\widetilde\phi(j)$ を代入すると、
両者の差の $\varepsilon$ 次係数が
$d^2\phi$ となる。
よって $(R_A)_3$ の平坦性は $d^2\phi=0$ と同値である。

Koszul 性から上の双加群複体は対角双加群の
projective resolution であり、bar 解消への比較は
Hochschild コチェイン複体の準同型を与える。
従って $d^2\phi=0$ なら、
$\phi$ は通常の次数零 Hochschild $2$-cocycle
$\mu_1:R_0\otimes R_0\to R_0$ の制限として表せる。
比較の変更によって生じる $d^1$-境界は
次数一生成元の $1+\varepsilon\theta$ という
基底変更で補正できる。
ゆえに $\phi$ に対応する二次関係を持つよう
$\mu_1$ と生成元の lift を選べる。

このとき $R_0\otimes_k A$ 上の積
\[
 a*b=ab+\varepsilon\mu_1(a,b)
\]
は $\varepsilon^2=0$ と cocycle 条件により結合的で、
全次数自由な次数付き $A$-代数 $R_{\mu_1}$ を与える。
上述の生成元の選択から、
$R_A\twoheadrightarrow R_{\mu_1}$ という
次数付き $A$-代数の全射が存在する。
各次数で両者の剰余体への制限は同じ $R_{0,n}$ であり、
$R_{\mu_1,n}$ は自由なので、
局所 Artin 環上の Nakayama 補題によって
この全射は各次数で同型である。
従って $(R_A)_n$ は全次数で自由となる。
逆向きは自明である。
\end{proof}

\begin{remark}[既知の Koszul 変形原理との関係]
形式的な Koszul 変形について、三次までの平坦性が
全次数の平坦性を含意することは Drinfeld の
Koszul deformation principle として既知である
\cite[Chapter 6, Theorem 2.1 と同章導入]{PolishchukPositselski2005}。
本補題は Koszul 双加群解消による一次定式化の直接証明である。
任意の非被約 Noetherian 基底についての強化は、
同書 Chapter 6, Theorem 7.3 を用いる
補題~\ref{lem:relative-koszul-degree-three} で与える。
\end{remark}

\subsection{三次 Hilbert 条件の接空間と非cubic表示の法方向}
\label{subsec:noncubic-tangent-degree-three}

前節の固定 Grassmann cell の局所座標を用いて、
相対反 Serre 環の三次成分が所望階数
$49=(2\cdot3+1)^2$ を保つ一次変形を調べる。
これだけで相対 $\HHstack$ を構成できるとは主張しない。

$V=k^9$、$J\subset V^{\otimes2}$ を階数 $56$ の関係部分空間、
$E=V^{\otimes2}/J$ とし、自然な写像
\[
 L_J:V\otimes J\longrightarrow E\otimes V,\qquad
 v\otimes r\longmapsto(\pi\otimes1)(v\otimes r)
 \tag{H3-L}\label{eq:noncubic-degree-three-relation-map}
\]
を考える。$R=T(V)/(J)$ とすると
$R_3=\operatorname{coker}L_J$ である。
$E\otimes V$ は次元 $25\cdot9=225$、
$V\otimes J$ は次元 $9\cdot56=504$ なので、
$\dim R_3=49$ は $\operatorname{rank}L_J=176$ と同値になる。
したがって Grassmannian 上に
$\operatorname{rank}L_J\le176$ で定まる
閉 determinantal locus $\mathcal D_{3,\ge49}$ がある。
階数がちょうど $176$ の点で
その接空間は、一般の determinantal 接空間の公式から
\[
 T_J\mathcal D_{3,\ge49}
 =\left\{
 \phi\in\operatorname{Hom}(J,E):
 \overline{\delta L_\phi}:
   \ker L_J\to\operatorname{coker}L_J
 \text{ が零}
 \right\}.
 \tag{H3-T}\label{eq:noncubic-third-degree-tangent}
\]

\begin{proposition}[noncubic の接方向の分離と三次条件の正確な接次元]
\label{prop:noncubic-third-degree-tangent}
$k$ を標数零の代数閉体とし、
命題~\ref{prop:noncubic-grassmann-explicit-cell} の
六係数から得られる関係空間を $J(f_1,f_2)$ とする。
次の三つの点を考える：
\[
\begin{array}{c|c|c}
 &f_1&f_2\\ \hline
 x_{\mathrm{dis}}&uv&u^2+v^2\\
 x_{\mathrm{shr}}&uv&u(v-u)\\
 x_{\mathrm{eq}}&uv&uv .
\end{array}
\tag{H3-types}\label{eq:noncubic-tangent-examples}
\]
それぞれで、$\mathrm{PGL}_9$ 軌道の接空間と
六係数の像の接空間の和は、正確に次元 $81$ である。
個々の軌道の接次元は順に $80,80,79$、
六係数の接像はそれぞれ六次元であり、
交わりの次元は順に $5,5,4$ である。

また上の三点で
$\mathcal D_{3,\ge49}$ の接空間はそれぞれ
\[
 \dim T_{x_{\mathrm{dis}}}\mathcal D_{3,\ge49}=83,
 \quad
 \dim T_{x_{\mathrm{shr}}}\mathcal D_{3,\ge49}=83,
 \quad
 \dim T_{x_{\mathrm{eq}}}\mathcal D_{3,\ge49}=84
 \tag{H3-bound}\label{eq:noncubic-tangent-bounds}
\]
と\emph{正確に}決まる。
従って、三次 Hilbert 条件を保つ一次方向を
noncubic 表示と基底変更から来る方向で割った空間の次元は、
この三点でそれぞれ\emph{正確に} $2,2,3$ である。
二形式の非交差な根対を持つ一般点についても
上界 $83$ は成立する。

最後の等号は三次 Hilbert determinantal locus に関するものである。
三次 Hilbert 条件は $\HH$ 公理の一部分にすぎず、
四次以上の平坦性・相対 heart・Serre・
AZ/AS 条件を含めない。
\end{proposition}
\begin{proof}
六係数の接像が六次元であることは、
命題~\ref{prop:noncubic-grassmann-explicit-cell} の
閉埋込み性による。
$\mathrm{PGL}_9$ 軌道の次元は pointed 自己同値群の
幾何点慣性の計算で求められる。
異なる根対（交わりが空または一根）の安定化群は有限、
一致した根対では一次元である。
従って軌道次元は $80,80,79$ となる。

相対二形式表示の群 $G$ は次元五で、
$\Phi:P\to\operatorname{Gr}(56,81)$ は $G$-同変である。
$G$ の根配置点の軌道次元は、
異なる根対の場合は五、一致する場合は四である。
これらの軌道の接方向は、$\mathrm{PGL}_9$ 軌道と
六係数像の接空間の交わりに入る。
したがって両者の和の次元は
それぞれ高々 $80+6-5=81$ と
$79+6-4=81$ である。

一方、固定25個の単位 pivot を用い、
$T_J\operatorname{Gr}(56,81)\simeq\operatorname{Hom}(J,E)$
の1400個の座標で、$\mathfrak{pgl}_9$ の
80個の行列単位と六係数の微分を計算すると、
係数を法 $101$ に還元した正確な行列の階数は、
三点とも $81$ である。
この検算は
\texttt{computations/verify\_noncubic\_degree3\_tangent.py}
に保存した。
整数係数行列の標数零における階数は
法 $101$ の階数以上なので、前段の上界と併せて
標数零の和の接空間は正確に次元 $81$ と分かる。

次に式~(\ref{eq:noncubic-third-degree-tangent}) の線形化を行う。
三点とも $L_J$ は階数 $176$ であり、
$\dim\ker L_J=328$、
$\dim\operatorname{coker}L_J=49$ である。
二次関係 $J$ を $25\times56$ の Grassmann affine 座標で動かしたとき、
$\overline{\delta L_\phi}$ は、
$328\cdot49$ 個の線形方程式として表示できる。
これらを1400変数上の行列として法 $101$ で
\emph{全行を}正確な行基本変形で処理すると、
三点での階数はそれぞれ
\[
 1317,\qquad 1317,\qquad 1316 .
\]
この計算ではランダムなサンプリングを行わず、
$\ker L_J$ の328列と余核の49座標から
全 $16072$ 行を構成し、
既知の86個の群・係数方向が全て線形条件を
満たすことも確認している。
同じスクリプトに全有限体計算を収めた。
整数係数の行列で行っているため、
標数零での線形条件の階数はそれぞれ
この値以上である。これが接次元の上界を与える。

さらに、法101の行基本変形から得た
$1400$ 成分の核ベクトル
$83,83,84$ 本を有理数へ復元する。
非零成分は分子の絶対値高々 $5$、
分母高々 $2$ の有理数で再構成できる。
各列を分母の最小公倍数で払って
\emph{整数ベクトル}にした後、
$\ker L_J$ の基底も整数係数へ持ち上げて
検算を行う。三次関係の線形化に対応する
\emph{整数係数}の全 $16072\times1400$ 行列を
独立に構成し直し、すべての整数核ベクトルについて
通常の整数行列積が\emph{正確に零}になることを確認した。
これらの列は元の自由変数位置に非零の
対角小行列を持つため、有理数上線型独立である。
この整数検証を上の計算スクリプトに収めた。
従って標数零での核次元はそれぞれ少なくとも
$83,83,84$ でもある。
前段の上界と併せて等号を得る。

非交差根対の一般点では、上記の基準点で
階数 $1317$ の minor が非零であることから、
階数の下半連続性により同じ上界が
非空 Zariski 開上で成立する。
\end{proof}

\begin{remark}[接空間の正確な解釈]
この命題は、$\mathrm{PGL}_9$ 軌道と
noncubic 六係数族の和の接空間だけでなく、
三次 Hilbert determinantal locus の
三基準点の接次元も\emph{正確に}決める。
ただし法101での階数の一致だけでは
標数零での等号は証明できないため、
有理核の再構成と整数行列積による零性の検証を必要とした。
これを、四次以降のHilbert条件や内在的
$\HHstack$ の接空間次元と同一視しない。
非交差根対の\emph{一般点}については、
少なくとも次元83の上界が従うが、
全非交差点で同じ次元だとはここでは主張しない。
\end{remark}

\begin{theorem}[noncubic 反 Serre 環の低次 Hochschild コホモロジー]
\label{thm:noncubic-hh01-hh2}
標数零の代数閉体上、式~(\ref{eq:noncubic-tangent-examples})
の三つの noncubic 反 Serre 環を
$R_{\mathrm{dis}},R_{\mathrm{shr}},R_{\mathrm{eq}}$
とする。$HH^i(R)_0$ は内部次数零の
Hochschild コホモロジーを表す。
このとき
\[
\begin{array}{c|ccc}
 &R_{\mathrm{dis}}&R_{\mathrm{shr}}&R_{\mathrm{eq}}\\ \hline
 \dim HH^1(R)_0 &1&1&2\\
 \dim HH^2(R)_0 &3&3&5
\end{array}
\tag{NC-HH}\label{eq:noncubic-hh-table}
\]
である。

前二点の $HH^1_0$ の一次元は
次数スカラーに対応する Euler 導分である。
二根対一致点では、これに加えて
連続自己同型の一次元が現れる。
$HH^2_0$ は全次数で平坦な
\emph{次数付き結合代数の一次変形}を
基底変更で割った接空間を記述する。
これを、そのまま内在的な
$\HHstack$ の相対 Abel 変形の接空間と
同一視してはならない。
\end{theorem}
\begin{proof}
命題~\ref{prop:noncubic-antiserre-koszul} により
三つの $R$ は Koszul である。
従って式~(\ref{eq:noncubic-koszul-hoch-cochain}) の
Koszul 双加群解消が内部次数零の
Hochschild コホモロジーを計算する。

ここで $\dim V=9$、
$\dim J=56$、$\dim R_2=25$、$\dim R_3=49$ である。
三次関係の自然な写像
$V\otimes J\to (V^{\otimes2}/J)\otimes V$
の階数は $176$ なので
\[
 \dim K_3
  =2(9\cdot56)-(9^3-49)=328.
\]
ゆえに
\[
 \dim C^1=81,\quad
 \dim C^2=56\cdot25=1400,\quad
 \dim C^3=328\cdot49=16072.
\]
前述の三次 Hilbert 条件の微分
$d^2:C^2\to C^3$ の階数は
命題~\ref{prop:noncubic-third-degree-tangent}
の整数核検証によって
\[
 1317,\quad1317,\quad1316
\]
と標数零で正確に決まっている。
一方 $d^1:C^1\to C^2$ は生成元の
基底変更による微分であり、
その階数は同じ命題の $\mathrm{PGL}_9$ 軌道次元
$80,80,79$ と一致する。
従って
\[
 \dim HH^2(R)_0
  =\dim\ker d^2-\operatorname{rank}d^1
  =(83-80,\ 83-80,\ 84-79)=(3,3,5).
\]

また $C^0=R_0=k$ から $C^1$ への
微分は内微分、すなわちスカラーとの交換なので零である。
よって
\[
 \dim HH^1(R)_0
 =\dim\ker d^1=(81-80,\ 81-80,\ 81-79)=(1,1,2).
\]
生成元 $V$ に恒等写像を作用させる
Euler 導分は、その次数 $n$ 成分で $n$ 倍として作用する。
従って $HH^1_0$ の一般点の一次元を張る。
二根対一致点の追加方向は
命題~\ref{prop:noncubic-presentation-stack} に記述した
正次元の pointed 慣性に対応する。
最後の一次変形の解釈は
補題~\ref{lem:koszul-first-order-degree-three} による。
\end{proof}

\begin{remark}[形式的な持上げと次数付き代数の局所模型]
今回得た $HH^2_0$ の $3,3,5$ 次元は
真の\emph{一次}平坦次数付き代数変形である。
それらが $k[\varepsilon]/(\varepsilon^3)$ や
$k[[t]]$ へ持ち上がるとは限らない。
後続の障害は $HH^3(R)_0$ 中にあり、
一次 cocycle の Gerstenhaber 平方の類を調べる必要がある。
また、次数付き環の変形が
相対 $\HH$ の heart・平方根・AS 正則被覆を持つかは
なお独立の認識問題である。
したがって今回の定理は内在的 $\HHstack$
の smoothness や成分数を証明するものではない。
\end{remark}

\begin{proposition}[noncubic 横断方向の三次平坦性の二次持上げ]
\label{prop:noncubic-transverse-second-order}
式~(\ref{eq:noncubic-tangent-examples}) の三つの点を取る。
各点において
\[
 T_x\mathcal D_{3,\ge49}
 \big/
 \bigl(T_x(\operatorname{PGL}_9\cdot x)
     +d\Phi(T_xP)\bigr)
\]
の次元は $r_x=2,2,3$ である。
それぞれで、この商へ同型に写る
$\mathbb Q$-上定義された $r_x$ 次元部分空間 $N_x$
を選べて、任意の $\phi\in N_x$ に対し
$k[t]/(t^3)$ 上の関係部分束
\[
 J_t=J+t\phi+t^2\psi
 \quad(\bmod\,t^3)
\]
で $R_t=T_{k[t]/t^3}(V\otimes k[t]/t^3)/(J_t)$
の\emph{三次成分}が自由階数 $49$ となるものが存在する。
言い換えれば、三次 Hilbert determinantal locus の
二次障害写像は選んだ $N_x$ 上で恒等的に零である。

ここで直接構成する持上げは三次成分に関する二次 jet である。
系~\ref{cor:noncubic-second-order-all-degrees} により、
これらは全次数で平坦な代数の変形となる。
任意の内在的 $\HHstack$ 族への持上げ、
および無限次の形式変形の存在は主張しない。
\end{proposition}
\begin{proof}
命題~\ref{prop:noncubic-third-degree-tangent} の
整数行列で、$\ker d^2$ と
群作用・六係数方向の和を計算する。
法 $101$ の行簡約で商の基底候補を選び、
分母を払って整数係数の
$\phi_1,\ldots,\phi_{r_x}\in\ker d^2$
を得る。実際に整数行列との積が零であることを
直接検証し、商での独立性は法 $101$ の
線型独立性で確認できる。

$L_J$ の階数 $176$ を与える小行列を固定し、
その行集合を $I$、列集合を $C$ と書く。
$A=L_0|_{I,C}$ は整数係数行列であり、
今回の三点では $\det A=\pm1$ と
整数係数逆行列の等式を確認した。
$K_0$ を $\ker L_0$ の列基底、
$Q_3$ を $\operatorname{coker}L_0$ への
PBW 基底による射影とする。
一次変形 $\phi$ に対して
\[
 L_t=L_0+tL_1(\phi)
      +t^2\bigl(L_1(\psi)+L_2(\phi)\bigr)
      \pmod{t^3}
\]
と書く。$L_2(\phi)$ は二次関係の
Grassmann affine 座標の積から来る
$\phi$ に関して二次の項である。

一次整合 $Q_3L_1(\phi)K_0=0$ が成り立つので、
kernel 基底の一次補正を
\[
 K_1=-i_C\,A^{-1}(L_1(\phi)K_0)|_I
\]
と選べる。従って二次持上げ条件は
\[
 d^2\psi+
 \mathcal Q_x(\phi)=0,\qquad
 \mathcal Q_x(\phi)
  =Q_3\left(
      L_2(\phi)K_0
      -L_1(\phi)i_CA^{-1}(L_1(\phi)K_0)|_I
      \right)
 \tag{NC-Ob2}\label{eq:noncubic-degree-three-second-obstruction}
\]
という $\operatorname{Hom}(\ker L_0,\operatorname{coker}L_0)$
の線形方程式に帰着する。
ここで $d^2$ は式~(\ref{eq:noncubic-koszul-hoch-cochain})
の微分である。
この表示は小行列 $A$ に関する Schur complement の
二次 Taylor 係数に等しい。

各基底ベクトル $\phi_i$ と、
異なる二基底ベクトルの和 $\phi_i+\phi_j$ に対し、
法 $101$ の線形方程式を解いて
$\psi_i,\psi_{ij}$ の候補を求めた。
さらに、$\phi$・$L_0$・$K_0$・$Q_3$・
$d^2$ を整数係数で再構成し、
必要に応じて $101^2$ を法とする有理数復元を行い、
\[
 d^2\psi_i=-\mathcal Q_x(\phi_i),\qquad
 d^2\psi_{ij}=-\mathcal Q_x(\phi_i+\phi_j)
\]
が\emph{厳密な有理数の等式}となることを検証した。
この計算は
\path{computations/verify_noncubic_second_order.py}
で再現できるようにした。
非交差・一根共有点では三組ずつ、
二根対一致点では六組の補正を検証し、
一致点では一部に分母 $2$ の係数を用いた。
有限体で障害が零というだけで
標数零の零性を主張しているのではない。

$\mathcal Q_x(\phi)$ は $\phi$ に関して二次式なので、
各基底方向とその二つずつの和で
障害の類が零であることは、
すべての
$\phi=\sum_i a_i\phi_i\in N_x$
に対して障害の類が零であることと同値である。
例えば二次補正として
\[
 \psi(\phi)=
 \sum_i a_i^2\psi_i
 +\sum_{i<j}a_ia_j(\psi_{ij}-\psi_i-\psi_j)
\]
を選べば式~(\ref{eq:noncubic-degree-three-second-obstruction})
が成立する。従って任意の横断方向で
三次 Hilbert 条件を二次まで保つ族を得る。
\end{proof}

\begin{remark}[二次持上げと実際のモジュライとの距離]
命題~\ref{prop:noncubic-transverse-second-order} は、
三次条件の一次法方向が二次障害によって
直ちに消えるわけではないことを示す。
今回の証明から $\operatorname{HH}^3(R)_0$ 全体の
消滅や、相対 Abel 圏としての変形の unobstructedness
は従わない。
次の課題は三次・それ以上の障害、および
foundation から得られる実際の相対反 Serre
族の横断接像との比較である。
\end{remark}

\begin{lemma}[二重平方環の中心 Clifford 表示]
\label{lem:noncubic-relative-clifford-center}
任意の可換環 $A$ 上の二元二次形式 $f_1,f_2$ に対し、
\[
 B=A[u,v]\langle c,d\rangle/(c^2-f_1,\ d^2-f_2),
 \qquad h=cd+dc
\]
とする。$h$ は中心元であり、
\[
 C=A[u,v,h]\subset B,\qquad
 B=C\oplus Cc\oplus Cd\oplus C(cd)
 \tag{NC-Cl}\label{eq:noncubic-relative-clifford}
\]
は次数付き $C$-加群として自由な直和である
（$\deg u=\deg v=\deg c=\deg d=1,\ \deg h=2$）。
特に $B$ は二元二次形式
$c^2=f_1,\ d^2=f_2,\ cd+dc=h$
の Clifford 代数として、中心多項式環上の階数四の代数である。
この表示は任意の底変更と可換する。
\end{lemma}
\begin{proof}
$u,v$ は中心なので、$c^2=f_1,d^2=f_2$ から
$[c,h]=[d,h]=0$ が直接従う。
関係
\[
 c^2=f_1,\qquad d^2=f_2,\qquad dc=h-cd
\]
を使うと、任意の語は $C$ 上 $1,c,d,cd$ の
線型結合に整理できる。従って右辺から $B$ への自然射は全射である。
一方、交互語基底によって $B$ の各次数は $A$-自由であり、
Hilbert 級数は $(1+z)/(1-z)^3$ である。
右辺も $A$-自由で、その Hilbert 級数は
\[
 \frac{1+2z+z^2}{(1-z)^2(1-z^2)}
 =\frac{1+z}{(1-z)^3}
\]
と一致する。
各次数では同じ階数の有限自由 $A$-加群の間の全射なので
同型である。従って $C$ も本当に $B$ の多項式部分環であり、
全表示は任意の底変更に両立する。
\end{proof}

\begin{proposition}[Clifford trace と相対 Serre 双加群]
\label{prop:noncubic-relative-serre-on-perfect}
$A$ を標数零の体上の可換 noetherian 環とし、
$\mathcal B=B$ を補題~\ref{lem:noncubic-relative-clifford-center}
の環とする。$C=A[u,v,h]$ に
$\deg(u,v,h)=(1,1,2)$ を入れる。
次数付き $C$-線型写像
\[
 \operatorname{tr}:B\longrightarrow C(-2)
\]
を、基底 $(1,c,d,cd)$ における $cd$ の係数を取る写像とする。
すると積による pairing は完全で、
\[
 \operatorname{tr}(ab)=\operatorname{tr}(b\iota(a)),
 \qquad
 \operatorname{Hom}_C(B,C(-4))
       \simeq B_{\iota}(-2)
 \tag{NC-Serre}\label{eq:noncubic-relative-serre-bimodule}
\]
という次数付き $B$-双加群の同型を与える。
ここで $B_\iota$ は右作用を $\iota$ で twist した双加群である。

重み付き射影スタック
$\mathcal X=\mathbb P_A(1,1,2)$ 上で、
$B$ から作る局所自由な階数四の代数層を $\widetilde B$ とする。
$\operatorname{Perf}(\mathcal X,\widetilde B)$ は相対 Serre 双対を持ち、
その関手は
\[
 S_A\simeq(-)_\iota(-2)[2],
 \qquad S_A^{-1}[2]\simeq(-)_\iota(2)
 \tag{NC-S}\label{eq:noncubic-relative-serre-functor}
\]
である。これらの pairing と双加群の構成は任意の
noetherian 底変更に両立する。
さらに補題~\ref{lem:noncubic-local-exchange} で得た
$\sigma^2=\iota$ により、適合平方根の候補は
\'etale 局所に exact 関手 $(-)_\sigma(1)$ として得られる。
\end{proposition}
\begin{proof}
$\operatorname{tr}$ の積 pairing は、$C$-基底
$1,c,d,cd$ 上で行列式が単元となる。
実際、$c^2=f_1,d^2=f_2,dc=h-cd$ により
$1$ は $cd$ と、$c$ は $d$ と、$d$ は $c$ と、
$cd$ は $1$ とそれぞれ係数が $\pm1$ の主項で対になる。
この pairing は $\iota$-cyclic である。
生成元 $u,v,h$ は中央で $\iota$ に固定され、
$c,d$ については
$\operatorname{tr}(cd)=1$、
$\operatorname{tr}(dc)=-1$ から
$\operatorname{tr}(ab)=\operatorname{tr}(b\iota(a))$
を確認できる。生成元での確認から全元に拡張する。
従って写像 $a\mapsto[b\mapsto\operatorname{tr}(ba)]$ は
$C$ 上の Frobenius 双対を与え、
右作用は $\iota$ で twist される。
重み $2$ の trace から
$\operatorname{Hom}_C(B,C)\simeq B_\iota(2)$、
$C$ の canonical shift $-4$ を加えると
(\ref{eq:noncubic-relative-serre-bimodule}) を得る。

$\mathcal X$ は基底上滑らかな相対次元二の
proper な重み付き射影 Deligne--Mumford スタックで、
$\omega_{\mathcal X/A}\simeq\mathcal O_{\mathcal X}(-4)$ である。
上の Clifford 表示により $\widetilde B$ は
$\mathcal O_{\mathcal X}$ 上有限局所自由である。
従って $\widetilde B$ 上 perfect な二対象 $P,Q$ に対し、
$\mathcal X/A$ の相対 Grothendieck--Serre 双対性と
有限局所自由代数の tensor--Hom 随伴から
\[
 \mathbf R\operatorname{Hom}_A
 \bigl(\mathbf R\operatorname{Hom}_{\widetilde B}(P,Q),A\bigr)
 \simeq
 \mathbf R\operatorname{Hom}_{\widetilde B}
 \bigl(Q,P\otimes_{\widetilde B}^{\mathbf L}
       \operatorname{Hom}_{\mathcal O_{\mathcal X}}
         (\widetilde B,\omega_{\mathcal X/A})[2]\bigr).
\]
ここで左辺・右辺の $\mathbf R\operatorname{Hom}$ は
$\mathcal X$ 上の大域導来 Hom を含む。
双加群の同型により右辺は
$\mathbf R\operatorname{Hom}_{\widetilde B}
 (Q,P_\iota(-2)[2])$ となる。
$\iota$ は次数付き自己同型なので twist と次数移動は
perfect 対象の自己同値であり、これが相対 Serre 関手である。
相対 dualizing 線束、Clifford trace、局所自由双加群は
底変更するので、双対性も noetherian 底変更と可換する。
最後に $\sigma^2=\iota$ と shift との可換性から
$(-)_\sigma(1)$ の二乗が
$(-)_\iota(2)=S_A^{-1}[2]$ となる。
この結論は perfect 圏の Serre 作用を計算するものであり、
全ての coherent $\widetilde B$-加群が perfect であることを
別の証明なしに主張するものではない。
\end{proof}

\begin{corollary}[noncubic 族の相対反 Serre 軌道の完全な計算]
\label{cor:noncubic-relative-orbit-hom}
前命題と同じ仮定の下、$\mathcal X=\mathbb P_A(1,1,2)$、
$\mathcal O_B=\widetilde B$ と置く。
perfect 圏上の相対反 Serre 作用
$\mathsf t=S_A^{-1}[2]$ に対し、全ての $n\ge0$ で
\[
 \mathbf R\operatorname{Hom}_{\mathcal X,\widetilde B}
    (\mathcal O_B,\mathsf t^n\mathcal O_B)
       \simeq B_{2n}[0]
 \tag{NC-orbit}\label{eq:noncubic-relative-orbit-hom}
\]
が成立し、任意の noetherian 底変更と可換する。
積は (\ref{eq:noncubic-antiserre-product}) の
$\iota$-twist された積（行列式線束の twist を除く）である。
従って命題~\ref{prop:noncubic-antiserre-descent}
の相対環は、幾何点で一致するだけでなく、
明示された noncubic 表示族の perfect 圏の
\emph{実際の相対反 Serre 軌道環} を降下したものである。
\end{corollary}
\begin{proof}
式~(\ref{eq:noncubic-relative-serre-functor}) より
$\mathsf t^n\mathcal O_B$ の下にある
$\mathcal O_{\mathcal X}$-加群は
$\widetilde B(2n)$ である。
$B$ の中心 Clifford 表示により
\[
 \widetilde B(2n)\simeq
 \mathcal O_{\mathcal X}(2n)
 \oplus\mathcal O_{\mathcal X}(2n-1)^{\oplus2}
 \oplus\mathcal O_{\mathcal X}(2n-2).
\]
重み付き多項式環 $C=A[u,v,h]$ の相対 \'Cech 複体は、
$\mathcal O_{\mathcal X}(j)$ の $H^1$ が任意の $j$ で零、
$H^2$ が $j\ge-3$ で零であることを示す。
$n\ge0$ なら上の最小の twist は $-2$ である。
従って高次大域コホモロジーは全て零であり、
$H^0$ はそれぞれ $C_{2n}$、$C_{2n-1}$、$C_{2n-2}$
（負の次数は零）である。
それらの直和は Clifford 表示の $B_{2n}$ となる。
各重みについて単項式で記述される weighted \v{C}ech 複体を
直接計算すると、消滅する項は底環上分裂可能な複体の一部であり、
残る $H^0$ は上記の有限自由加群となる。
従って任意の noetherian 底変更と整合する。
積は $\mathsf t=(-)_\iota(2)$ の functorial な合成から
$a\star b=a\iota^{\deg_{\mathcal R}a}(b)$ と計算できる。
\end{proof}

\begin{proposition}[noetherian 基底上の noncubic 表示の有限性]
\label{prop:noncubic-relative-noetherian}
任意の可換 noetherian 基底に引き戻した普遍環 $\mathcal B$、
および上記の反 Serre 環 $\mathcal R$ は左右 noetherian である。
より精密には $\mathcal B$ は中央多項式環上階数四の有限自由代数、
$\mathcal R$ は同じ中央多項式環の四次 Veronese 上
有限生成加群である。
これらの主張は相対 Serre 双対性や相対 heart の認識を仮定しない。
\end{proposition}
\begin{proof}
noetherian なアフィン基底 $\operatorname{Spec}A$ 上で調べる。
前補題により
$C=A[u,v,h]$ は $\mathcal B$ の中央に入り、
$\mathcal B$ は $C$ 上階数四の自由加群である。
$C$ は可換 noetherian なので、$\mathcal B$ は左右 noetherian である。

さらに $\mathcal R$ を、偶数 Veronese
$\mathcal B^{(2)}$ の $\iota$ による Zhang twist とみなす。
$C^{(4)}=\bigoplus_{m\ge0}C_{4m}$ は有限生成の
可換 noetherian 環である。
$\iota$ は $C$ 上恒等であり、$C_{4m}$ の元は
$\mathcal R$ の次数 $2m$ に入る。
この次数では $\iota^{2m}=1$ だから
(\ref{eq:noncubic-antiserre-product}) の
左・右の積は通常の積と一致する。
従って $C^{(4)}$ は $\mathcal R$ の中心に入る。
$C$ は $C^{(4)}$ 上有限生成加群であり
（任意の次数四以上の単項式は次数四の単項式で割れる）、
$\mathcal B$ は $C$ 上階数四の自由加群である。
従って偶数次数部分 $\mathcal R$ も
$C^{(4)}$ 上有限生成加群となる。
この有限性から $\mathcal R$ も左右 noetherian である。
線束 $\mathcal D$ の twist はアフィン基底上で局所的に
自明化できるため、結論は表示スタックの任意の
noetherian 基底への引戻しでも成立する。
\end{proof}

\begin{proposition}[平方根を大域選択しない noncubic 相対 heart]
\label{prop:noncubic-relative-heart-veronese}
$\mathfrak N^{\mathrm{pres}}=[P/G]$ 上の
命題~\ref{prop:noncubic-antiserre-descent} の
次数付き環 $\mathcal R$ は、任意の可換 noetherian 基底に引き戻すと
左右 noetherian、次数一生成であり、各次数は有限局所自由となる。
従って
\[
  (\mathcal A_T,\mathcal O_T)
  =\bigl(\operatorname{qgr}\mathcal R_T,\,
          \pi\mathcal R_T\bigr)
  \tag{NC-heart}\label{eq:noncubic-relative-heart}
\]
という構造対象付き noetherian Abel 圏を、基底上の
表示データから直接構成できる。
これは $\sigma$ や ambient の大域的な選択を必要としない。

$T$ の \'etale 局所化後、
命題~\ref{prop:noncubic-relative-regular-cut} の
適合 Zhang twist $\mathcal B^\sigma$ を取ると、
構造対象を保つ exact 同値
\[
 \operatorname{qgr}\mathcal B^\sigma
       \simeq\operatorname{qgr}\mathcal R_T
\]
があり、左辺の二次数 shift と右辺の一次数 shift は対応する。
全幾何ファイバーは当初の noncubic $\HH$ 対象である。
さらに perfect 対象の相対反 Serre 関手は右辺の
次数 shift $(1)$ に対応する。
\end{proposition}
\begin{proof}
次数別局所自由性、次数一生成、noetherian 性はそれぞれ
命題~\ref{prop:noncubic-antiserre-descent}、
命題~\ref{prop:noncubic-antiserre-grassmann}、
命題~\ref{prop:noncubic-relative-noetherian} による。
従って $\operatorname{qgr}\mathcal R_T$ は
有限生成次数付き右加群を有限次数の torsion 加群で割った
noetherian Abel 圏であり、構造対象 $\pi\mathcal R_T$ を持つ。

\'etale 局所に $\det W$ を自明化し、
補題~\ref{lem:noncubic-local-exchange} の $\sigma$ を取れば
$\mathcal R_T=(\mathcal B^\sigma)^{(2)}$ である。
$\mathcal B^\sigma$ は noetherian で次数一生成なので、
次数二 Veronese の quotient-category 同値を適用できる。
体上での通常の主張は Verevkin~\cite[Theorem 4.4]{Verevkin1992}
（定理の引用は後続文献の照合による）に対応する。
今回必要な noetherian 底上でも、
次数一生成性から高次数の奇数成分が直前の偶数成分と次数一の積で
生成されること、$\mathcal B^\sigma$ が偶数 Veronese 上
有限生成加群であることを用いれば、有限生成加群の
Veronese 関手が torsion を保ち、商圏で忠実充満かつ本質的全射となる
同じ証明が成立する。
具体的には次数付き加群 $M$ を
$\bigoplus_{n}M_{2n}$ に移し、逆方向は偶数部分環から元の環への
誘導加群を取る。両方向の自然比較の核と余核は高次数で零となり、
$qgr$ の同値を与える。正則右加群と
二次数 shift をそれぞれ正則 Veronese 加群と
一次数 shift に送る。
ここで $\mathcal R_T$ は偶数部分そのものなので、
全ての比較はもとの環と底変更の形成に両立する。
perfect 圏の Serre の表示は
命題~\ref{prop:noncubic-relative-serre-on-perfect} により
$S^{-1}[2]\simeq(-)_\iota(2)$ であり、
Zhang twist および Veronese の同値によって $(1)$ に移る。
各幾何ファイバーの適合性は既存の逆構成定理による。
\end{proof}

\begin{remark}[明示表示族と横断認識の区別]
降下した反 Serre 環と $\qgr$ heart は、
定理~\cite[\texttt{thm:noncubic-relative-effectivity}]{QuadricIntrinsic2026} により、
全ての相対 $\mathfrak H_2$ 条件を備えた実際の表示族を与える。
同定には weighted Clifford 表示の canonical compact 圏・相対 smoothness、
同じ根の Serre 双対・AZ・全軌道、および Ore/PBW 被覆の全相対条件を用いる。
ただし任意の内在的 noncubic 族のこの表示への局所還元、
表示射の相対全忠実性、noncubic 特殊点を持つ一般有限型の横断族の認識は未証明である。
Artin 基底上の横断変形は系~\ref{cor:noncubic-all-artin-h2} により認識できる。
全幾何ファイバーが noncubic でも、非被約基底上の四次 minor が
スキーム論的に零であるとは限らない。
従って次元一の表示族から全境界近傍の代数性を結論しない。
\end{remark}

% !TEX root = ../nc-quadric-moduli.tex
% 2026-11-21 / v1.39
\section{noncubic 境界の適合根と Hochschild 作用}
\label{sec:noncubic-hh-roots}


\paragraph{結論と範囲.}
中央二重平方表示の二つの根対を交換する\emph{trace 零の}線型対合を選ぶと、
対応する適合 Serre 根は有限 foundation の $\operatorname{HH}^1,\operatorname{HH}^2$ に恒等に作用する。
ただし根対が一致するときは対角の交換対合を選ぶ。
trace 零だけでは十分でないことと、選択条件の具体的な行列式を後で示す。
この計算は noncubic 表示の一次元の変形だけを扱うものではない。
foundation の全変形接空間、すなわち根対が異なる場合の三次元、
一致する場合の四次元を扱う。
従って既存の Artin 根持上げ補題を適用すると、指定した良い根は
全ての Artin Abel 変形へ持ち上がり、その Morita 類は一意である。
さらに有限 foundation の各 noncubic 点には、
この根と実際の Serre 平方同型を持つ有限型 étale 近傍が存在する。
同じ根に対する AS4 被覆の横断変形への持上げと、有限型基底上の全相対認識は
本節の結論に含めない。
本節の計算の新規性は未確認である。

\subsection{四項 foundation の二項移動}

$k$ を標数零の代数閉体とする。
full geometric helix の四項を $E_0,E_1,E_2,E_3$ とし、
\[
A=\Hom(E_0,E_1),\quad B=\Hom(E_1,E_2),\quad
C=\Hom(E_2,E_3),\quad D=\Hom(E_0,E_2)
\]
と置く。$A,B,C$ は二次元、$D$ は四次元である。
$W=\Hom(E_0,E_3)$ は六次元であり、
\[
F:A\otimes B\longrightarrow D,\qquad
G:C\otimes D\longrightarrow W
\]
が foundation の積を与える。
因子の並べ替えを省略して書く。
$G$ は全射とし、$K=\ker G$ は二次元とする。
包含 $K\subset C\otimes D$ を curry した写像を
\[
j:D^{\vee}\longrightarrow C\otimes K^{\vee}
\]
と書く。さらに
\begin{equation}\label{eq:nchh-delta}
\delta=(-F^{\vee},j):D^{\vee}\longrightarrow
H=(A^{\vee}\otimes B^{\vee})\oplus(C\otimes K^{\vee})
\end{equation}
は単射とする。これらは、以下で使う noncubic 点の近傍で開条件である。
$\mathscr L=\operatorname{im}\delta\in\Gr(4,H)$ と置く。

\begin{lemma}[二項移動の有限線型代数]\label{lem:nchh-gale}
$E_{i+4}=S^{-1}E_i[2]$ により延長し、四項窓を
$(E_2,E_3,E_4,E_5)$ へ移す。
対応するデータは
\begin{align}
A'&=C,& B'&=K^{\vee},& C'&=A,\label{eq:nchh-new-small}\\
D'&=H/\mathscr L,&
F'&:C\otimes K^{\vee}\longrightarrow D',\label{eq:nchh-new-D}\\
K'&=\operatorname{im}\bigl(B^{\vee}\longrightarrow A\otimes D'\bigr),\label{eq:nchh-new-K}
\end{align}
である。\eqref{eq:nchh-new-K} は $A$ の coevaluation と第一 summand の商写像の合成である。
$K'\simeq B^{\vee}$ となる開近傍で、この新しい四平面は
\begin{equation}\label{eq:nchh-gale}
\mathscr L'=\{(-Y,X):(X,Y)\in\mathscr L^{\perp}\}
 \subset(C^{\vee}\otimes K)\oplus(A\otimes B)
\end{equation}
である。ここで $\perp$ は $H$ と $H^{\vee}$ の標準評価による annihilator を表す。
この記述は任意の局所 Artin 変形にも同じ式で適用できる。
\end{lemma}

\begin{proof}
符号も含めて相対 bar 複体で確認する。
積を $\Gamma_{ij}=\Hom(E_i,E_j)$ とする上三角代数として書き、
左 projective $P_i=\Gamma e_i$ を使えば $E_i$ に対応する。
$I_i=\Hom_k(e_i\Gamma,k)=S(P_i)$ と置くと
\[
\Hom(E_j,E_4)=\Ext^2_{\Gamma}(I_j,P_0).
\]
$I_3$ は頂点三の単純加群であり、bar 複体の可消去な成分を除くと
\[
\Ext^2(I_3,P_0)=\operatorname{coker}(G^{\vee})=K^{\vee}
\]
である。

$I_2$ は頂点二で $k$、頂点三で $C^{\vee}$ を持ち、
$C$ の作用は評価である。
bar cochain の次数一は
$D^{\vee}\oplus\Hom(W,C)$、次数二は
\[
(A B)^{\vee}\oplus\Hom(A B C,C)\oplus\Hom(D C,C),
\]
次数三は $\Hom(A B C,C)$ である。
次数二の元を $(r,g,h)$ とすると、次の微分は
\[
(r,g,h)\longmapsto g-h(F\otimes1_C)-r\otimes1_C.
\]
従って kernel では中央成分 $g$ を消去できる。
次数一の $(a,b)$ の像は、残る二成分で
\[
(-aF,-bG+a\otimes1_C).
\]
$bG$ を除くことは $h$ を $K$ へ制限することに等しく、
残る $a$ の像はちょうど $(-F^{\vee}a,j(a))$ である。
従って \eqref{eq:nchh-new-D} を得る。
前合成 $E_2\to E_3$ は第二 summand の包含なので、$F'$ も表示した写像となる。

次に $0\to A\otimes P_0\to P_1\to S_1\to0$ を用いる。
頂点一、二、三の部分代数は道代数なので
$\Ext^1(I_2,S_1)=B^{\vee}$、$\Ext^2(I_2,S_1)=0$ である。
長完全列における境界写像は
\[
B^{\vee}\longrightarrow A\otimes\Ext^2(I_2,P_0)=A\otimes D'
\]
であり、bar 複体で代表を取れば
$b^{\vee}\mapsto\sum_a a\otimes[a^{\vee}\otimes b^{\vee}]$ と計算される。
これが \eqref{eq:nchh-new-K} である。
$G'$ は $A\otimes D'$ をその像で割った六次元商である。
また $\Hom(E_4,E_5)=A$ は Serre 移動による同定を使う。

$D'^{\vee}$ は $\mathscr L^{\perp}\subset
(A\otimes B)\oplus(C^{\vee}\otimes K)$ である。
$F'^{\vee}$ は第二成分 $Y$ への射影、$j'$ は第一成分 $X$ への射影だから、
$(-F'^{\vee},j')$ の像は \eqref{eq:nchh-gale} となる。
全て有限 projective な複体と直和因子の消去であるから、
閉ファイバーで必要な階数が成立する任意の Artin 変形でも同じ計算が成立する。
\end{proof}

\begin{remark}[全変形空間を使うこと]
$G$ の商の基底を忘れ、代わりに $K$ の基底を選ぶと、
foundation のデータは $(F,j)$ である。
$\delta$ が単射である開部分で $D$ の基底変更を除くことは、
四平面 $\mathscr L\subset H$ を取ることに等しい。
従って全 foundation 変形の局所模型は
$\Gr(4,8)$ の開部分を
\[
\widetilde G=\GL(A)\times\GL(B)\times\GL(C)\times\GL(K)
\]
で割ったものとなる。
第一 summand への作用は $g_A^{-t}\otimes g_B^{-t}$、
第二 summand への作用は $g_C\otimes g_K^{-t}$ である。
自明に作用する三次元中心を除いた群を $G_{\mathrm{eff}}$ と書く。
その次元は十三で、Grassmannian の次元は十六である。
この模型には noncubic 表示を保たない全ての横断方向が含まれている。
\end{remark}

\subsection{中央二重平方と良い根の選択}

\begin{lemma}[trace 零の交換対合]\label{lem:nchh-exchange}
二つの非退化二元二次形式 $f_1,f_2$ に対し、座標とそれぞれの平方根変数の
定数倍を選び、
\begin{equation}\label{eq:nchh-normal-forms}
f_1=uv,\qquad f_2=\ell_1\ell_2,
\qquad
\begin{pmatrix}\ell_1\\\ell_2\end{pmatrix}
=L_0\begin{pmatrix}u\\v\end{pmatrix},\qquad
L_0=\begin{pmatrix}\alpha&\beta\\\gamma&-\alpha\end{pmatrix},\quad
\alpha^2+\beta\gamma=1
\end{equation}
とできる。二根対が一致する場合にも $L_0=\operatorname{diag}(1,-1)$ と選べる。
\end{lemma}
\begin{proof}
既存の 補題~\ref{lem:noncubic-local-exchange} の行列
$\left(\begin{smallmatrix}1&-a\\b&-1\end{smallmatrix}\right)$ を使う。
その平方は単元 $(1-ab)I$ であり、平方根で割れば trace 零の対合となる。
$L_0f_1=\eta f_2$ の単元 $\eta$ を $d$ の定数倍に吸収する。
この構成は二根対一致の場合にも $a=b=0$ で成立する。
$f_1$ の非退化性により $f_1=uv$ と正規化できる。
\end{proof}

以下、
\[
\mathcal B=k[u,v]\langle c,d\rangle/(c^2-uv,d^2-\ell_1\ell_2)
\]
と置く。$u,v$ は中央である。
\begin{equation}\label{eq:nchh-sigma}
\sigma(u,v)=L_0(u,v),\qquad\sigma(c)=d,\qquad\sigma(d)=-c
\end{equation}
は次数付き自己同型で、$\sigma^2=\iota$、ただし
$\iota(u)=u,\iota(v)=v,\iota(c)=-c,\iota(d)=-d$ である。
命題~\ref{prop:noncubic-relative-serre-on-perfect} の双加群計算から
$S=(-)_\iota(-2)[2]$ であり、
\begin{equation}\label{eq:nchh-good-root}
s=(-)_\sigma(1),\qquad s^2=S^{-1}[2]
\end{equation}
は適合根である。同じ根の AS4 表示は定理~\cite[\texttt{thm:ore-anticommuting-central}]{QuadricIntrinsic2026} の Ore/Zhang 構成が与える。

\begin{proposition}[根の作用と交代 Gale 双対]\label{prop:nchh-root-gale}
\eqref{eq:nchh-good-root} の Hochschild 作用は、前節の Grassmann 模型で
\begin{equation}\label{eq:nchh-involution}
\tau(\mathscr L)=J\mathscr L^{\perp},\qquad
J(X,Y)=(D_0Y,-D_0X),\qquad D_0=\operatorname{diag}(1,-1)
\end{equation}
の微分により与えられる。$X,Y$ を二次正方行列として書いている。
$J:H^{\vee}\to H$ は交代同型であり、中央二重平方の点 $\mathscr L$ は
$\tau$ の固定点、すなわち対応する交代形式の Lagrangian 四平面である。
\end{proposition}

\begin{proof}
行列分解を
\[
\Phi_c=\begin{pmatrix}c&u\\-v&-c\end{pmatrix},\qquad
\Phi_d=\begin{pmatrix}d&\ell_1\\-\ell_2&-d\end{pmatrix}
\]
とする。$M_c=\operatorname{coker}\Phi_c$、$M_d=\operatorname{coker}\Phi_d$ とし、
normal foundation を
\[
(\OO,\pi M_c,\OO(1),\pi M_d(1))
\]
と取る。
$A,C$ はそれぞれ $M_c,M_d$ の二つの次数零生成元の空間、
$B$ は $\Phi_c$ の二つの行、$K$ は $\Phi_d$ の二つの関係列である。
従って \eqref{eq:nchh-delta} の係数行列は
\begin{equation}\label{eq:nchh-plane}
\lambda\longmapsto
\bigl(-\Phi_c^{t}(\lambda),\Phi_d(\lambda)\bigr)
\quad\bigl(\lambda\in\langle u,v,c,d\rangle^{\vee}\bigr).
\end{equation}

二項移動の $D'=H/\mathscr L$ を $\langle u,v,c,d\rangle$ と同定する写像は、
二つの summand 上で次の係数行列を持つ。
\begin{equation}\label{eq:nchh-quotient-identification}
h_0=-D_0\Phi_c=
\begin{pmatrix}-c&-u\\-v&-c\end{pmatrix},\qquad
h_1=\Phi_d^tD_0=
\begin{pmatrix}d&\ell_2\\\ell_1&d\end{pmatrix}.
\end{equation}
実際、$h_0F^{\vee}=h_1j$ であり、両辺は
$u^{\vee}\mapsto-v$, $v^{\vee}\mapsto u$, $c^{\vee},d^{\vee}\mapsto0$
である。ここで $\alpha^2+\beta\gamma=1$ を使う。
$h_0,h_1$ の像は $u,v,c,d$ を張るため、商からの写像は同型である。
$F'$ は $h_1$、$K'$ の二つの関係列は $h_0=-D_0\Phi_c$ に対応する。
これらは移動した実際の Hom と関係の行列である。

右加群の規約も確認する。
$(M_c)_\sigma$ の提示は
$\sigma^{-1}\Phi_c=-D_0\Phi_dD_0$ であり、
生成元を $D_0$ 倍すると $M_d$ へ移る。
一方 $\sigma^{-1}\Phi_d=\Phi_c$ なので $(M_d)_\sigma\simeq M_c$ である。
また canonical 四移動 $t=(-)_\iota(2)$ では
$(M_c)_\iota\simeq M_c$ の生成元比較が $D_0$ となる。
\eqref{eq:nchh-quotient-identification} の $h_1$ は
$K^{\vee}$ を $D_0$ 倍した $\Phi_d$ の行へ対応させ、
$K'$ は $-D_0\Phi_c$ の関係列である。
従って $s$ による四つの二次元空間の比較は、順に
\begin{equation}\label{eq:nchh-four-comparisons}
(D_0,-I_2,D_0,-I_2)
\end{equation}
となる。
この比較を \eqref{eq:nchh-gale} に入れると、第一 summand は $D_0Y$、
第二 summand は $-D_0X$ となり、\eqref{eq:nchh-involution} を得る。

直接の固定点確認もできる。
$\mathscr L^{\perp}$ の元は
\[
X=\begin{pmatrix}a&-\beta r-\alpha s\\\alpha r-\gamma s&a\end{pmatrix},\qquad
Y=\begin{pmatrix}b&r\\s&b\end{pmatrix}
\]
の形である。$J(X,Y)$ は \eqref{eq:nchh-plane} における
$(\lambda_u,\lambda_v,\lambda_c,\lambda_d)=(s,r,-b,-a)$ の像である。
従って $\tau(\mathscr L)=\mathscr L$。
行ごとの座標で $D_0$ の左乗法は対称行列
$D_L=\operatorname{diag}(1,1,-1,-1)$ であり、
\[
J=\begin{pmatrix}0&D_L\\-D_L&0\end{pmatrix},\qquad J^t=-J.
\]
以上は foundation の全 Artin 変形を移動してから、閉ファイバーを指定した
$s$ で戻す構成である。従ってその接作用は、Morita 不変な
Hochschild 作用そのものである。
\end{proof}

\subsection{全 Hochschild 接空間への作用}

\begin{theorem}\label{thm:noncubic-hh-trivial-root}
任意の noncubic $\mathcal H_2$ 対象の有限 foundation を $\Gamma$ とする。
補題~\ref{lem:nchh-exchange} の trace 零の交換対合を用い、
根対が一致する場合には $L_0=D_0$ と選ぶ。
この選択による適合根 $s$ は
\[
s_*|_{\operatorname{HH}^1(\Gamma)}=1,\qquad s_*|_{\operatorname{HH}^2(\Gamma)}=1
\]
を満たす。
二つの根対が異なる場合は $\dim\operatorname{HH}^1=0$, $\dim\operatorname{HH}^2=3$、
一致する場合はそれぞれ $1,4$ である。
\end{theorem}

\begin{proof}
二つの接複体を同変に比較する。
Grassmannian の接空間 $T$ は十六次元である。
\eqref{eq:nchh-involution} は Lagrangian 点における交代 annihilator involution なので
\[
T^+\simeq\operatorname{Sym}^2\mathscr L^{\vee},\qquad
T^-\simeq\bigwedge^2\mathscr L^{\vee},\qquad
(\dim T^+,\dim T^-)=(10,6).
\]
実際、Lagrangian complement を選んで近傍を行列 $Z$ の graph と書けば、
微分は $Z\mapsto Z^t$ である。

次に基底変更群を調べる。
$\operatorname{Lie}\widetilde G=\mathfrak{gl}(A)\oplus\mathfrak{gl}(B)
\oplus\mathfrak{gl}(C)\oplus\mathfrak{gl}(K)$ 上の作用は
\begin{equation}\label{eq:nchh-lie-action}
(a,b,c,k)\longmapsto
(D_0cD_0,-k^t,D_0aD_0,-b^t).
\end{equation}
二つの対を交換するので、その正負固有空間の次元はともに八である。
自明な中心は、scalar 四つ組について
\[
-a-b=c-k
\]
という一つの条件を満たす三次元空間である。
その正固有空間は $a=c$, $b=-k$, $a+b=0$ により一次元、
負固有空間は $a=-c$, $b=k$ により二次元である。
従って有効群の Lie 代数 $\mathfrak g$ について
\begin{equation}\label{eq:nchh-lie-dim}
(\dim\mathfrak g^+,\dim\mathfrak g^-)=(7,6).
\end{equation}

foundation の全変形を表すことから、同変な完全列
\begin{equation}\label{eq:nchh-tangent-sequence}
0\longrightarrow\operatorname{HH}^1(\Gamma)\longrightarrow\mathfrak g
 \longrightarrow T\longrightarrow\operatorname{HH}^2(\Gamma)\longrightarrow0
\end{equation}
がある。これは有向代数の基底変更を除いた結合代数変形の接複体である。
前節の Grassmann 模型は $G$ の全射性と $\delta$ の単射性という開条件だけで得られ、
この完全列に横断方向を除く追加条件は入っていない。
ここで $\OO$ は exceptional であり、
$\operatorname{HH}^0(\Gamma)=Z(\Gamma)\to\operatorname{End}(\OO)$ は $k\to k$ の同型である。
従って pointed 化による追加の一次核はなく、
定理~\ref{thm:artin-foundation-lift} の同定の下で
pointed 慣性の Lie 代数は $\operatorname{HH}^1(\Gamma)$ である。

根対が異なる場合、命題~\ref{prop:noncubic-presentation-stack} の慣性計算により
pointed 慣性は有限である。標数零なので $\operatorname{HH}^1(\Gamma)=0$ である。
\eqref{eq:nchh-tangent-sequence} の負部分は六次元空間から六次元空間への単射であり同型、
従って $\operatorname{HH}^2(\Gamma)^-=0$ となる。正部分は $10-7=3$ 次元である。

根対が一致する場合は $f_1=uv$, $f_2=-uv$, $L_0=D_0$ と取る。
pointed 慣性の連結成分は
$u\mapsto zu$, $v\mapsto z^{-1}v$, $c,d\mapsto c,d$ のトーラスである。
これは $\sigma$ と可換であり、次数移動とも可換なので $s$ の共役作用は恒等である。
従って一次元の $\operatorname{HH}^1(\Gamma)$ は正固有空間に入る。
線型模型でも、$T_z=\operatorname{diag}(z,1)$ に対する安定化群の表示は
\[
(g_A,g_B,g_C,g_K)=(T_z,T_z^{-1},T_z,T_z),
\]
であり、\eqref{eq:nchh-lie-action} で固定される。
負部分は再び同型、正部分の余核は $10-(7-1)=4$ 次元である。
これで両次数への恒等作用が従う。
\end{proof}

\begin{corollary}[全 Artin 変形への根の持上げ]\label{cor:noncubic-root-all-artin}
前定理の根 $s$ は、閉ファイバーを固定した任意の局所 Artin Abel 変形へ、
相対 Serre 平方根として持ち上がる。
その Morita 自己同値の同型類は一意である。
この一意性は平方同型自体の一意な選択を意味しない。
\end{corollary}
\begin{proof}
既存の 定理~\ref{thm:artin-foundation-lift} により、全 Abel 変形を有限自由相対 smooth な
foundation 変形で計算できる。
系~\ref{cor:hhtrivial-root-artin} と定理~\ref{thm:noncubic-hh-trivial-root} の $\operatorname{HH}^2$ 作用から
根の存在を得る。
命題~\ref{prop:artin-serre-root-uniqueness} と $\operatorname{HH}^1$ 作用から同型類の一意性を得る。
補題~\ref{lem:artin-morita-heart-exact} により、持ち上がった Morita 関手は
元の Abel 変形の heart 上で完全である。
本証明は AS4 被覆の横断変形への持上げを与えていない。
\end{proof}

\subsection{全 foundation 変形の Lagrangian 正規形}
\label{subsec:noncubic-lagrangian-slice}

前の接計算は、有限型の具体的な局所表示に強められる。
この表示は foundation の全横断方向を含むが、
それだけで相対 $\mathcal H_2$ の認識を与えるものではない。

四つの二次元空間をそれぞれ $k^2$ と同定し、
第\ref{sec:noncubic-hh-roots}節の $H$ と $J$ を使う。
$\omega$ を行列 $J^{-1}$ の非退化交代形式とする。
$X=\Gr(4,H)$、$L=\operatorname{LG}(4,H,\omega)$ と置けば、
$L$ は $\tau(\mathscr L)=J\mathscr L^\perp$ の固定部分スキームである。
基底変更群の作用は
\[
 \rho(g_A,g_B,g_C,g_K)
 =\operatorname{diag}(g_A^{-t}\otimes g_B^{-t},
                      g_C\otimes g_K^{-t})
\]
である。$X$ 上で自明に作用する中心を
\[
 C_0=\{(aI,bI,cI,dI):a^{-1}b^{-1}=cd^{-1}\}
       \simeq\mathbf G_m^3,
 \qquad G=\widetilde G/C_0
\]
と書く。

\begin{proposition}[Lagrangian な有限型 \texorpdfstring{\'etale}{etale} 表示]
\label{prop:noncubic-lagrangian-etale-slice}
群対合
\begin{equation}\label{eq:noncubic-group-involution}
 \widetilde\theta(g_A,g_B,g_C,g_K)
 =(D_0g_CD_0,g_K^{-t},D_0g_AD_0,g_B^{-t})
\end{equation}
は $G$ の対合 $\theta$ を誘導し、
$\tau(g\mathscr L)=\theta(g)\tau(\mathscr L)$ を満たす。
その固定群の単位成分を $G^+=(G^\theta)^\circ$ と書くと、
\[
 G^+\simeq
 (\GL_2\times\GL_2)/(\mathbf G_m\times\mu_2),
 \qquad \dim G^+=7.
\]
ここで中央部分群の埋込みは
$(z,\epsilon)\mapsto(zI,\epsilon z^{-1}I)$ である。

$G$ の Lie 代数の負固有空間を $\mathfrak g^-$ とする。
自然な軌道微分から得る束写像
\begin{equation}\label{eq:noncubic-normal-orbit-map}
 q:\mathfrak g^-\otimes\mathcal O_L
   \longrightarrow N_{L/X}
   \simeq\bigwedge^2\mathscr U^\vee
\end{equation}
を考える。$\mathscr U$ は $L$ 上の tautological 四平面束である。
その行列式の非消滅開を $L^{\mathrm{tr}}$ とすると、
\[
 [L^{\mathrm{tr}}/G^+]\longrightarrow[X/G]
 \tag{LG-slice}\label{eq:noncubic-lagrangian-stack-map}
\]
は表現可能な \texorpdfstring{\'etale}{etale} 射である。
双方のスタックの次元は三である。
$\det q$ は、定数一次元因子を除けば
Pl\"ucker 線束の三乗 $\mathcal O_L(3)$ の切断である。

中央二重平方の trace 零正規形
\eqref{eq:nchh-plane} が $L^{\mathrm{tr}}$ に属する必要十分条件は
$\alpha\ne0$ である。
根対が異なればこの条件は自動的に成立し、
根対が一致するときは補題~\ref{lem:nchh-exchange} の通り
$L_0=D_0$ と選べば成立する。
従ってその近傍の全 foundation 変形は、
\texorpdfstring{\'etale}{etale} 局所に Lagrangian 四平面で表示できる。
同じ結論は、補題~\ref{lem:nchh-gale} の階数条件を課した
foundation の開部分へ制限したときにも成立する。
\end{proposition}

\begin{proof}
行列を掛けると
\[
 \rho(\widetilde\theta(g))
       =J\rho(g)^{-t}J^{-1}
\]
を得る。この式は $\widetilde\theta^2=1$、
$\widetilde\theta(C_0)=C_0$ と上の同変性を示す。

$\widetilde G$ の固定群は
\[
 (a,b)\longmapsto(a,b,D_0aD_0,b^{-t})
 \qquad(a,b\in\GL_2)
\]
の像である。その $C_0$ との共通部分は
$a=zI$、$b=\epsilon z^{-1}I$、$\epsilon^2=1$ で尽くされる。
従って $G$ 内の像は上で表示した滑らかな連結七次元群である。
標数零の群の固定部分スキームは滑らかであり、その Lie 代数は
$\mathfrak g^+$ である。
\eqref{eq:nchh-lie-dim} によりこの次元は七なので、
上の像は $(G^\theta)^\circ$ と一致する。

$\tau$ の $L$ に沿う微分には、
接束 $T_L\simeq\operatorname{Sym}^2\mathscr U^\vee$ と
法束 $N_{L/X}\simeq\bigwedge^2\mathscr U^\vee$ がそれぞれ
正・負固有束として現れる。
同変性により軌道微分は正負分解を保つため、
\eqref{eq:noncubic-normal-orbit-map} が得られる。
両辺の階数は六であり、その可逆性は $G^+$-不変な開条件である。
$\det(\bigwedge^2\mathscr U^\vee)
=(\det\mathscr U^\vee)^3$ だから、行列式の線束も主張の通りである。

閉部分群 $G^+\subset G$ による associated bundle
$G\times^{G^+}L^{\mathrm{tr}}$ と軌道写像
\[
 m:G\times^{G^+}L^{\mathrm{tr}}\longrightarrow X,
 \qquad[g,\mathscr L]\longmapsto g\mathscr L
\]
を考える。これは滑らかな十六次元空間の間の射である。
$[1,\mathscr L]$ における接写像は
\[
 (\mathfrak g\oplus T_{\mathscr L}L)\big/
 \{(a,-a\mathscr L):a\in\mathfrak g^+\}
 \longrightarrow T_{\mathscr L}X,
 \qquad[(a,v)]\longmapsto a\mathscr L+v.
\]
正部分では恒等、負部分では $q_{\mathscr L}$ となる。
従って $L^{\mathrm{tr}}$ 上では同型である。
$G$-同変性により全ての点で同型なので、$m$ は \texorpdfstring{\'etale}{etale} である。
この射は \eqref{eq:noncubic-lagrangian-stack-map} を atlas $X\to[X/G]$
で基底変更したものだから、同スタック射も表現可能かつ
\texorpdfstring{\'etale}{etale} である。
次元は $10-7=16-13=3$ である。

最後の主張は一般係数で直接確認できる。
$\mathfrak g^-$ の元を
$(a,b,-D_0aD_0,b^t)$、$a,b\in\mathfrak{sl}_2$ で代表し、
各 $\mathfrak{sl}_2$ の基底を
$\operatorname{diag}(1,-1),E_{12},E_{21}$ の順に取る。
\eqref{eq:nchh-plane} の係数行列を $M$ とすると、
法方向の交代形式は $M^tJ^{-1}\rho_*(a,b,-D_0aD_0,b^t)M$ である。
$u^\vee,v^\vee,c^\vee,d^\vee$ の順に対する六つの上三角成分を並べれば、
$q_{\mathscr L}$ の行列は
\[
 Q_{\alpha,\beta,\gamma}=
 \begin{pmatrix}
 -2\alpha&0&0&2\alpha&0&0\\
 0&\gamma&-\alpha&0&\alpha&-\gamma\\
 0&0&-1&0&-1&0\\
 0&-\alpha&-\beta&0&\beta&\alpha\\
 0&1&0&0&0&1\\
 -2&0&0&-2&0&0
 \end{pmatrix},
 \qquad
 \det Q_{\alpha,\beta,\gamma}
       =-32\alpha(\alpha^2+\beta\gamma)=-32\alpha.
\]
一方 $\alpha=0$ なら $\beta\gamma=1$ なので
$f_2=(\beta v)(\gamma u)=uv=f_1$ である。
従って根対非一致の全場合には $\alpha\ne0$ であり、
一致の場合にも対角選択 $\alpha=1,\beta=\gamma=0$ が使える。
foundation を再構成するための追加条件は有限写像の階数条件で開であり、
それらを満たす点の近傍に上の射を制限すればよい。
\end{proof}

\begin{remark}[この表示が与えるもの]
\label{rem:noncubic-lagrangian-scope}
上の射の像は foundation スタックの開部分であり、
全 noncubic 点を含む。その smooth 性は有限代数表示についての主張である。
実際の Morita 平方根との比較には
系~\ref{cor:noncubic-etale-morita-root} を使い、
必要に応じて二つの \texorpdfstring{\'etale}{etale} 近傍の積を取る。
全正軌道 Hom、相対 heart、同じ根の AS4 被覆の存在を、
この Lagrangian 正規形だけから結論してはならない。
特に \eqref{eq:noncubic-lagrangian-stack-map} の標的を
内在的 $\mathfrak H_2$ に置き換えたとは主張していない。
\end{remark}

\begin{remark}[trace 零だけでは良い根を特徴付けない]
\label{rem:noncubic-offdiagonal-root}
根対一致の場合、$L_0=\left(\begin{smallmatrix}0&1\\1&0\end{smallmatrix}\right)$
も trace 零の交換対合であるが、上の \texorpdfstring{\'etale}{etale} 開には属さない。
対応する根は連結慣性
$u\mapsto zu,v\mapsto z^{-1}v$ を逆元へ移すため、
$\operatorname{HH}^1$ に $-1$ で作用する。
\eqref{eq:nchh-tangent-sequence} の正負分解より、この場合
\[
 (\dim\operatorname{HH}^2_+,\dim\operatorname{HH}^2_-)=(3,1).
\]
従ってこの根は全ての一次 Abel 変形に延長するわけではない。
本稿で全 Artin 延長に使う良い根は、根対一致時には
$L_0=D_0$ を選ぶという条件を含む。
\end{remark}

\subsection{Hochschild 二次に自明な根の étale 代数化}
\label{subsec:hh2-etale-morita-root}

次の命題は noncubic 点に限らない。
根の存在する空間から étale な切片を取ることにより、
Hochschild 一次への作用を仮定せずに根を代数化する。

\begin{proposition}[Hochschild 二次に自明な根の étale 代数化]
\label{prop:hh2-etale-morita-root}
$S$ を標数零の代数閉体 $k$ 上の局所有限型 scheme とし、
$\Lambda$ を $S$ 上の有限局所自由代数とする。
次を仮定する。
\begin{enumerate}
\item $\operatorname{Perf}(\Lambda)$ は相対 smooth proper であり、
任意の基底変更後も
$\operatorname{HH}^{<0}(\Lambda)=0$、
$\operatorname{HH}^0(\Lambda)=\OO_S$ である。
\item 相対 inverse Serre $t=S_\Lambda^{-1}[2]$ は
可逆 perfect 双加群で表され、任意の導来基底変更と可換する。
\item $x\in S(k)$ と
$p_0\in\operatorname{DPic}_k(\Lambda_x)$ が与えられ、
$p_0^2=t_x$、
$(p_0)_*|_{\operatorname{HH}^2(\Lambda_x)}=1$ である。
\end{enumerate}
$\mathcal G\to S$ を Morita 自己同値の étale 局所同型類の代数空間、
$\tau=[t]$ とし、
\[
 \mathcal Z=
 \mathcal G\times_{\mathrm{sq},\mathcal G,\tau}S
\]
を平方根空間とする。
このとき $\mathcal Z\to S$ は $p_0$ で smooth である。
さらに有限型 affine な étale 近傍 $(U,u)\to(S,x)$ 上に、
可逆 perfect $\Lambda_U$-双加群 $P$ と実際の同型
\[
 P\otimes_{\Lambda_U}^{\mathbf L}P\simeq t_U
 \tag{N-root}\label{eq:hh2-etale-actual-root}
\]
が存在し、$[P_u]=p_0$ となる。
追加で $(p_0)_*|_{\operatorname{HH}^1(\Lambda_x)}=1$ なら、
$\mathcal Z\to S$ 自体が $p_0$ で étale であり、
同じ閉根を持つ Artin 持上げの Morita 同型類は一意である。
\end{proposition}

\begin{proof}
命題~\ref{prop:quadric-etale-morita-root} の代数空間の構成は、
可換 quadric という仮定を使わず、仮定 (1), (2) に適用できる。
詳しくは \cite[著者公開版 April 2007, Corollary 3.24,
Lemma 3.25]{TV2007} により $\mathcal G$ は $S$ 上局所有限表示の代数空間である。
同論文 Corollary 3.22 により、実際の可逆 perfect 双加群のスタックから
$\mathcal G$ への射は étale 局所に中立な $\mathbf G_m$-gerbe である。
$S$ は局所 Noetherian なので $\mathrm{sq}$ も局所有限表示であり、
$\mathcal Z$ は $S$ 上局所有限表示となる。
ここでは基底の線束による tensor が消える étale 層化を用いる。

$p_0$ を閉点とする任意の局所 Artin 変形
$\operatorname{Spec} A\to\mathcal Z$ と、これに整合する小拡大
$A'\twoheadrightarrow A$ 上の $\operatorname{Spec} A'\to S$ を取る。
剰余体は $k$ とする。
$A,A'$ は strictly henselian で、上の gerbe は中立であり、
線束も自明だから、指定された点は実際の
$p_A\in\operatorname{DPic}_A(\Lambda_A)$、
$p_A^2=t_A$ で表せる。
補題~\ref{lem:root-obstruction-composition} により、
$p_A$ の持上げ障害
$b\in\operatorname{HH}^2(\Lambda_x)\otimes_k\ker(A'\to A)$ は
\[
 b+(p_0)_*b=0
\]
を満たす。
仮定 (3) と標数零により $b=0$ である。
さらに命題~\ref{prop:artin-serre-square-correction} は、
指定された $p_A$ を保ったまま平方関係を補正する。
従ってこの小拡大に対し、指定された点を保つ
$\operatorname{Spec} A'\to\mathcal Z$ が存在する。

$(V,v)\to(\mathcal Z,p_0)$ を étale scheme chart とする。
形式的 étale 性により上の持上げは、指定された
$\operatorname{Spec} A\to V$ も保って $V$ へ持ち上がる。
閉点での小拡大判定
\cite[Tag 02HX, Lemma 37.12.1]{StacksSmoothAtlas}
により $V\to S$ は $v$ で smooth であり、
従って $\mathcal Z\to S$ は $p_0$ で smooth である。

$(V,v)$ を smooth locus に縮める。
$v$ は fiber の閉点であり、$\kappa(v)=\kappa(x)=k$ だから、
\cite[Tag 057G, Lemma 37.38.5]{StacksSmoothAtlas}
により、$v$ を含む immersion $U_0\to V$ で
$U_0\to S$ が étale となるものが存在する。
これは指定された根を通る étale な切片であり、
根の一意性を要求しない。
$U_0$ を通じて上の gerbe を étale 局所に中立化し、
実際の可逆双加群 $P$ を取る。
定理~\ref{thm:bimodule-line-comparison} は比較線束
\[
 L=H^0\mathbf R\mathcal Hom_{\Lambda_U^e}
       (P\otimes_{\Lambda_U}^{\mathbf L}P,t_U)
\]
と評価同型
$(P\otimes_{\Lambda_U}^{\mathbf L}P)\otimes L\simeq t_U$
を与える。
$u$ の Zariski 近傍で $L$ を自明化すれば
\eqref{eq:hh2-etale-actual-root} を得る。
各細分は étale であり、$u$ を含む affine 開へ縮めて
$U$ を有限型にできる。
その像は $x$ を含む開部分なので、必要なら $S$ をこの像へ縮める。

追加の Hochschild 一次の仮定の下では、
命題~\ref{prop:artin-serre-root-uniqueness} により
指定された閉根を持つ Artin 持上げの同型類は一意である。
従って smooth な $\mathcal Z\to S$ の $p_0$ での相対接空間は零であり、
その相対次元は零である。よってこの射自体が $p_0$ で étale となる。
平方同型そのものを一意に選ぶという意味ではない。
\end{proof}

\begin{remark}[Hochschild 一次の役割]
\label{rem:hh2-etale-root-tangent}
同じ閉根を持つ根の一次的な差は
$a\in\operatorname{HH}^1(\Lambda_x)$ で
$a+(p_0)_*a=0$ を満たすものである。
従って平方根空間の相対接空間は
$\ker(1+(p_0)_*)$ である。
Hochschild 一次への恒等作用はこの相対接空間を消し、
根の空間自体を étale にするために使う。
étale 近傍上で一つの根を選ぶ存在定理には不要である。
\end{remark}

\begin{corollary}[全 noncubic 点の有限型 étale Morita 根]
\label{cor:noncubic-etale-morita-root}
$S$ を有限 foundation スタックの有限型 affine chart とし、
$x\in S(k)$ を任意の noncubic $\mathcal H_2$ 対象の
normal foundation に対応する点とする。
定理~\ref{thm:noncubic-hh-trivial-root} の良い根 $s_0$ を選ぶ。
このとき $x$ の有限型 affine étale 近傍 $(U,u)\to(S,x)$ 上に、
$s_0$ を閉 fiber とする可逆 perfect 双加群 $P$ と
実際の平方同型 \eqref{eq:hh2-etale-actual-root} が存在する。
平方根空間自体もこの点で étale である。

四つの普遍 projective 対象 $F_0,\ldots,F_3$ を、
閉 fiber で
$(\OO,\pi M_c,\OO(1),\pi M_d(1))$ に対応させる。
近傍をさらに縮めれば、$P$ による自己同値 $\mathsf s$ は
\[
 \mathsf s F_0\simeq F_2,\qquad
 \mathsf s F_1\simeq F_3
 \tag{N-window}\label{eq:noncubic-etale-root-window}
\]
を満たす。これらの同型は任意の導来基底変更に保たれる。
\end{corollary}

\begin{proof}
補題~\ref{lem:foundation-relative-smooth-proper} と
補題~\ref{lem:foundation-relative-antiserre-perfect} が前命題の
相対 smooth proper 性と Serre の基底変更を与える。
任意の基底変更後に中心が基底環であり負の Hochschild 次数が零であることは、
命題~\ref{prop:quadric-etale-morita-root} の有向代数の中心の計算による。
定理~\ref{thm:noncubic-hh-trivial-root} は両 Hochschild 次数への恒等作用を与えるので、
前命題を適用できる。

最後の対象比較は有限個の perfect 複体の開性による。
$a=0,1$ に対し
$\mathbf R\Hom(\mathsf sF_a,F_{a+2})$ は底上 perfect で、
命題~\ref{prop:nchh-root-gale} の実際の根の比較により
閉 fiber は $k[0]$ である。
従って近傍で次数零の線束となる。
閉 fiber の同型に対応する切断を選び、
その射の cone に perfect Nakayama を適用すると
\eqref{eq:noncubic-etale-root-window} が得られる。
この同型は derived tensor 後も同型のままである。
\end{proof}

\begin{remark}[有限型近傍でまだ結論していないこと]
\label{rem:noncubic-etale-root-scope}
上の $U$ は有限代数と perfect Morita 根の実際の有限型族である。
この構成だけから相対 Abel heart、
全正軌道 Hom の同じ開近傍での次数零集中、
全軌道環の平坦性・Noetherian 性・AZ 豊富性、
同じ根に対する相対 AS4 被覆は従わない。
従って $U\to\mathfrak H_2$ や
noncubic 点を含む smooth Artin 開近傍を、ここでは構成したとは主張しない。
明示的な noncubic 二形式表示族の効果性と、
cubic を横断するこの étale 近傍の効果性も区別する。
\end{remark}

\subsection{根の有限性、計算の検証と残る問題}

根対が異なる場合、命題~\cite[\texttt{prop:first-orbit-object}]{QuadricIntrinsic2026} により任意の適合根の
$\OO$ 上の像は同じ内在的対象である。
従って二つの根の差は pointed 自己同値を与え、全根の同型類は
有限な pointed 自己同値群の torsor の部分集合に入る。
一般点でのこの群は既存原稿では $D_8\times C_2$（位数十六）であり、
二つの平方根変数の符号群 $\mu_2^2$ だけに限定してはいけない。
本節は全ての根の作用を恒等と主張せず、上記の trace 零の選択による
一つの根を特定して証明した。
他の根に対する作用は、それとの差である pointed 自己同値の作用との合成である。

補足の厳密有理計算
\texttt{computations/verify\_noncubic\_gale\_root.py} は、非交差、一根共有、一致の各配置で
$\mathscr L$、$J$、Grassmann tangent の十六次元行列、
基底変更軌道の全十六生成元を構成する。
浮動小数点や乱数を使わず、
\[
\begin{array}{c|ccc}
 &\text{非交差}&\text{一根共有}&\text{一致}\\\hline
\dim(\text{orbit}^+,\text{orbit}^-)&(7,6)&(7,6)&(6,6)\\
\dim(\operatorname{HH}^2_+,\operatorname{HH}^2_-)&(3,0)&(3,0)&(4,0)
\end{array}
\]
を確認する。一般の配置への証明はこの三例の計算から外挿せず、
定理~\ref{thm:noncubic-hh-trivial-root} の同変接複体と安定化群の議論による。

\paragraph{同じ根の AS4 被覆は別の問題である.}
持ち上げた同じ根の軌道環の $B_1$ は四次元の $D$ であり、
$B_2=\Hom(E_0,E_4)$ と積は同じ有限 foundation の bar 複体から計算できる。
根を適合するように twist する前の中央二重平方環では、二次双対に正規な二次元の
pencil が存在するが、固定した適合根の $\mathcal B^\sigma$ でも
同じ pencil が残るとは限らない。ambient の関係線の $\sigma$ 不変性が必要である。
非交差・一根共有・一致の三つの有理配置では、三次平坦性を保つ
関係空間の接次元 $18,18,19$ の全方向に対し、選んだ projective 正規線と
その normalizer が一次に一意に持ち上がることを
\texttt{computations/verify\_noncubic\_cover\_tangent.py} の有限線型代数で確認した。
これはこの特定の被覆の一次計算であり、それ自体から高次存在は従わない。
次節では別の中央 AS4 被覆を全配置で構成し、
系~\ref{cor:noncubic-all-artin-h2} で同じ良い根に対する全 Artin 延長を証明する。
有限型近傍の一様な効果性はなお別問題である。

% 2026-11-21 / v1.40
% !TEX root = ../nc-quadric-surface-moduli.tex
\section{noncubic の良い根に対する中央 AS4 被覆}
\label{sec:noncubic-central-hull}

\paragraph{範囲.}
以下の構成は、中央二重平方環を与える従来の Ore 被覆とは異なる被覆である。
同じ trace 零交換対合で選んだ Serre 根の軌道環について、切断元自体が中央となる
AS4 被覆を与える。既存の Ore ambient の中で切断元を中央化する主張ではない。
新規性は未確認である。

基礎体を標数零の代数閉体とし、\(a,b\in k\)、
\(\theta=1-ab\ne0\)、\(\delta^2=\theta\) と置く。
二つの非退化二元二次形式は座標・定数倍を選べば
\[
 f_1=uv,\qquad f_2=(u-av)(bu-v)
       =bu^2-(1+ab)uv+av^2
\]
と書ける。二根対の一致も \(a=b=0\) に含まれる。
\[
 B_0=k[u,v]\langle c,d\rangle/(c^2-f_1,d^2-f_2)
\]
で \(u,v\) は中央とする。良い交換自己同型は
\[
 \sigma(u)=(u-av)/\delta,\quad
 \sigma(v)=(bu-v)/\delta,\quad
 \sigma(c)=d/\delta,\quad \sigma(d)=-\delta c
\]
であり、\(\sigma^2=\iota\)、\(\iota(u,v,c,d)=(u,v,-c,-d)\) である。
以下では \(x*y=x\sigma^{\deg x}(y)\) の Zhang 規約を使う。

\begin{theorem}\label{thm:noncubic-central-hull}
次の六関係環を定める。
\[
 A_0=T(u,v,c,d)\Big/
 \left(
 \begin{gathered}
 [u,c],[u,d],[v,c],[v,d],\quad c^2-\frac{uv+vu}{2},\\
 d^2-bu^2+\frac{1+ab}{2}(uv+vu)-av^2
 \end{gathered}
 \right).
\]
\(A_0\) は左右 Noetherian な quadratic Koszul AS 正則環で、型は
\((1,4,6,4,1)\) である。
\(g_0=uv-vu\) は左右正則正規で、
\[
 A_0/(g_0)=B_0,\qquad
 g_0u=-ug_0,\quad g_0v=-vg_0,\quad
 g_0c=cg_0,\quad g_0d=dg_0.
\]
\(\sigma\) は \(A_0\) に延長し、\(\sigma(g_0)=-g_0\) となる。
従って \(A=A_0^\sigma\) における \(g=g_0\) は中央正則二次元であり、
\[
 B_0^\sigma=A/(g).
\]
特に、指定した良い Serre 根の全軌道環に対する中央 AS4 被覆が存在する。
さらに \(E=(B_0^\sigma)^!\) の二次中心は一次元であり、
この被覆の双対線がその全体である。
\end{theorem}

\begin{proof}
まず untwisted な双対 \(E_0=B_0^!\) を計算する。
\(U,V,C,D\) を \(u,v,c,d\) の双対とし、\(p=C^2,q=D^2\) と置く。
二次関係の直交補空間から
\begin{gather*}
 CD=DC=0,\qquad
 UC=-CU,\ VC=-CV,\ UD=-DU,\ VD=-DV,\\
 U^2=-bq,\qquad V^2=-aq,\qquad
 UV+VU=-p+(1+ab)q
\end{gather*}
を得る。\(p,q\) は中央、\(pq=0\) である。
\(E_0\) は Koszul で、Hilbert 級数は
\(h_{E_0}(t)=(1+t)^4/(1-t^2)\) である。
既存の Ore 被覆の双対により \(\zeta=p+q\) は正則である
（定理~\cite[\texttt{thm:ore-anticommuting-central}]{QuadricIntrinsic2026} の双対表示）。
具体的には $E_0/(p+q)$ は同定理の Ore ambient の二次双対で、
Hilbert 級数は $(1+t)^4$ である。$p+q$ は中央だから、その乗法の完全列と
$h_{E_0}(t)=(1+t)^4/(1-t^2)$ の係数比較でも左右の核が零と分かる。

\(\chi=UV-VU\) と置くと、Clifford の二生成元計算から
\[
 \chi U=-U\chi,\quad \chi V=-V\chi,\quad
 \chi C=C\chi,\quad \chi D=D\chi,\qquad
 \chi^2=p^2+\theta^2q^2.
\]
\(\chi^2(p^2+\theta^{-2}q^2)=(p+q)^4=\zeta^4\) だから、
\(\zeta\) の左右正則性により \(\chi\) も左右正則である。
従って \(P_0=E_0/(\chi)\) は正則正規二次商であり、Koszul である
（命題~\ref{prop:all-h2-antiserre-koszul} の環変換による証明を適用する）。
その Hilbert 級数は \((1+t)^4\) である。

\(P_0\) の Frobenius 性を明示的に確認する。
以下、\(P_0\) 内の同じ文字を用い、最高次数の \(q^2\) の係数を trace とする。
\(P_0\) では \(UV=VU\)、\(p^2=-\theta^2q^2\) である。
関係から次数四の全ての語は \(q^2\) の倍数へ還元される。
交差項の消滅を明記すると、
\[
 pU=(1+ab)qU+2bqV,\qquad pV=2aqU+(1+ab)qV
\]
を $D$ 倍した二式の係数行列は行列式 $\theta^2$ を持つ。
$pD=0$ なので $qUD=qVD=0$ となる。$qC=0$ と同じ二式から
$pUC=pVC=0$ も従い、残る語は $p^2,q^2$ の倍数に還元される。
Hilbert 級数によって \(\dim(P_0)_4=1\) なので \(q^2\ne0\) である。
次数三の基底を
\[
 Uq,\ Vq,\ C^3,\ D^3
\]
と取れる。実際、これらが全てを張ることは上の関係から従い、次元は四である。
次数一の基底 \((U,V,C,D)\) との積の行列は
\[
 \begin{pmatrix}
 -b&(1+ab)/2&0&0\\
 (1+ab)/2&-a&0&0\\
 0&0&-\theta^2&0\\
 0&0&0&1
 \end{pmatrix}.
\]
左上 block の行列式は \(-\theta^2/4\) なので非退化である。
逆順の積も、交差反交換の符号を除いて同じ非退化行列を与える。
次数二には基底 \((p,q,UC,VC,UD,VD)\) を取り、その積の行列は次の
三 block の直和である。
\[
 \operatorname{diag}(-\theta^2,1),\qquad
 \begin{pmatrix}0&-\theta^2/2\\-\theta^2/2&0\end{pmatrix},\qquad
 \begin{pmatrix}b&-(1+ab)/2\\-(1+ab)/2&a\end{pmatrix}.
\]
いずれも非退化である。これで \(P_0\) は Koszul Frobenius、最高次数四となる。

直交補空間を再び取ると \(P_0^!\) は定理の六関係環 \(A_0\) である。
Koszul 双対性と Frobenius pairing により \(A_0\) は型
\((1,4,6,4,1)\) の AS 正則環で、\(h_{A_0}(t)=(1-t)^{-4}\) となる。
\(g_0\) の正規性は表示からも直接見える。
\(c^2=(uv+vu)/2\) は \(u,v\) と可換であるから、
\([u,uv+vu]=ug_0+g_0u=0\)、
\([v,uv+vu]=-vg_0-g_0v=0\) である。
交差交換関係から \(g_0\) は \(c,d\) と可換である。
商は明らかに \(B_0\) であり、
\(h_{B_0}(t)=(1-t^2)h_{A_0}(t)\) を各次数の正規元商の完全列へ適用すれば、
\(g_0\) の左右乗法の核は零である。
\(B_0\) の左右 Noetherian 性と正則正規元による標準的な有限次数フィルターの
持上げから、\(A_0\) も左右 Noetherian となる。

最後に twist を確認する。
対称二次形式 \((uv+vu)/2\) と
\(bu^2-(1+ab)(uv+vu)/2+av^2\) は、\(\sigma\) により係数
\(\theta^{-1},\theta\) を伴って交換される。
従って六関係は保存され、\(\sigma\) は \(A_0\) の自己同型である。
\(\det(\sigma|_{\langle u,v\rangle})=-1\) なので \(\sigma(g_0)=-g_0\) である。
\(g_0\) の normalizer は \(-\iota\) だから、一次元 \(x\) について
\[
 g_0*x=g_0\sigma^2(x)=g_0\iota(x)=-xg_0
             =x\sigma(g_0)=x*g_0.
\]
これで \(g\) は中央である。twist は Koszul 性、AS 型、Noetherian 性、正則性を保つ。

中心の一意性を示す。
\(\alpha=\sigma^{-t}\) を双対空間上の自己同型とすると、二次関係の直交補空間より
\[
 (B_0^\sigma)^!\cong E_0^\alpha,\qquad
 \alpha^2=\iota,\quad\alpha(\chi)=-\chi.
\]
従って \(\chi\) の twisted な像は中央であり、この線の商は \(P_0^\alpha\) である。
\((E_0)_2\) の基底は
\(p,q,\chi,UC,VC,UD,VD\) である。
\(\alpha^4=1\) なので、twisted center の元を \(\alpha\)-固有元
\(\alpha(z)=\lambda z\) に分解できる。
中心条件は
\[
 zx=\lambda xz\quad(x=U,V),\qquad
 zx=-\lambda xz\quad(x=C,D)
\]
となる。
\(\lambda=\pm i\) なら \(z\) は四つの交差項の線形結合である。
\(C,D\) との式と \(CD=0\) より各係数は零となる。
この線形独立性は \(\zeta\) で局所化し、直交する \(p/\zeta,q/\zeta\) の二成分で
\(C,D\) をそれぞれ可逆にして確認できる。
各成分の二元 Clifford 表示では \(U,V\) は独立であり、
非退化二次形式により \(\chi/p\)、\(\chi/q\) は scalar ではない。
\(\lambda=1\) なら \(z\in\langle p,q,\chi\rangle\) で、\(C,D\) との反交換条件は
上の二成分で \(z=0\) を与える。
\(\lambda=-1\) なら \(U,V\) との反交換条件で \(p,q\) の係数が零となり、
\(z\in k\chi\) となる。
\(\zeta\) の正則性により局所化前にも同じ結論となる。
よって二次中心はちょうどこの一本の線である。
\end{proof}

\begin{lemma}[良い noncubic 根では二次反中心が消える]
\label{lem:noncubic-no-quadratic-anticenter}
定理~\ref{thm:noncubic-central-hull} の \(E=(B_0^\sigma)^!\) について
\[
 Z(E)_2=k\chi,\qquad
 Z^-(E)_2:=\{z\in E_2\mid zx=-xz\ (x\in E_1)\}=0
\]
である。最初の式の \(\chi\) は同定 \(E=E_0^\alpha\) による像を表す。
\end{lemma}
\begin{proof}
中心の式は前定理による。反中心も \(\alpha\) により保たれる。
\(\alpha(z)=\lambda z\) とすれば、反中心の条件は untwisted な積で
\[
 zx=-\lambda xz\quad(x=U,V),\qquad
 zx=\lambda xz\quad(x=C,D)
\]
となる。\(\alpha^4=1\) だから \(\lambda\in\{1,-1,i,-i\}\) ごとに調べればよい。
\(\lambda=\pm i\) の場合は \(z\) が交差四項の線形結合となり、
\(C,D\) との式、\(CD=0\)、\(\lambda\ne-1\) により全ての係数が零になる。
\(\lambda=1\) の場合は \(z\in\langle p,q,\chi\rangle\) であり、
\(U,V\) との反交換条件が \(p,q\) の係数を消す。
残る \(\chi\) は \(\alpha\)-固有値 \(-1\) を持つため、この場合にも \(z=0\) である。
\(\lambda=-1\) では \(U,V\) との可換条件が \(\chi\) の係数を消し、
\(C,D\) との反交換条件が \(p,q\) の係数を消す。
これらの係数比較は \(\zeta=p+q\) で局所化した二つの Clifford 成分で行える。
\(\chi\) は正則であり、各成分の \(U,V\) は非零であることに注意する。
\(\zeta\) も正則なので局所化前の反中心は零である。
\end{proof}


\begin{remark}[相対化と横断族への適用で分ける入力]
明示 noncubic 表示族上では、\(2\theta\) が単元となる任意の可換 Noetherian 基底に
上の構成を相対化できる。双対の各次数の projectivity は既存の相対 Koszul 補題による。
\(\chi\) 乗法は各幾何 fiber で単射であるから局所直和因子への単射となり、商も
任意の基底変更と両立する有限 projective 加群となる。
上の pairing の行列式は単元なので、相対 Frobenius 条件も成立する。
商に PP の相対 Koszul 定理を適用することで双対 ambient の全次数平坦性を得る。
正規元商からの Noetherian 持上げは任意の Noetherian 基底変更後にも適用できる。

横断族では、この計算だけから全軌道の平坦性や canonical heart は従わない。
ただし、被約有限型基底 \(S\) 上にその良い根の平坦な Koszul 軌道環 \(B\) があり、
\(E=B^!\) の二次・三次が階数七・八の束であると仮定する。
稠密な cubic 開部分上で同じ根の canonical hull が中央であるなら、
commutator 写像 \(E_2\to\mathcal Hom(E_1,E_3)\) の階数は稠密開で高々六である。
被約性により七次 minors は全体で零であり、上の定理は noncubic 点で階数六を与える。
従ってその点の近傍で核は任意の基底変更と両立する可逆部分束となる。
これが延ばすべき中央線であり、従来の Ore 被覆の正規線ではない。
局所環上では、次節の定理~\ref{thm:noncubic-local-central-hull} がこの入力を確認する。
全軌道平坦性、同じ根と cubic canonical 周期との比較、有限 AS4 条件を満たす
単一近傍の構成、Noetherian heart の認識は、この線の延長と別に記載する必要がある。
\end{remark}

% !TEX root = ../nc-quadric-moduli.tex
% 2026-11-21 / v1.40
\section{五頂点窓と noncubic 変形の全 Artin 認識}
\label{sec:noncubic-local-tails}

\paragraph{局所環で成立することと、一様な開で残ること.}
本節では系~\ref{cor:noncubic-etale-morita-root} の有限型近傍の
\emph{指定された noncubic 点における局所環}を取る。
その同じ局所環上で、全ての前向き Hom、全軌道環、canonical
Grothendieck 導来圏と Noetherian heart を構成する。
極小解消の Betti 数を族の中で固定せず、五頂点の非極小 bar 解消を使う。
さらに二次双対の中心・反中心を使って同じ根の中央 AS4 被覆を局所環上に構成し、
全 noncubic 点で任意の Artin 変形を相対 $\mathfrak H_2$ として認識する。
有限型の一つの開近傍上での全次数一様性と全幾何 fiber の $\HH$ 所属は別問題であり、
本節だけから noncubic 点を含む $\mathfrak H_2$ の開 chart は得られない。

\subsection{極小 Betti 数に依存しない bar 解消}

以下で $C$ は正の向きの $R$-線形 $\mathbb Z$ 代数、
$P_i=e_iC$、$S_i$ は頂点 $i$ に $R$ を置く加群とする。
積は $C_{ij}\otimes C_{jk}\to C_{ik}$ と書く。

\begin{lemma}[局所 perfect 複体の次数零集中]
\label{lem:local-perfect-fibre-concentration}
$(R,\mathfrak m,k)$ を可換局所環、$P\in\operatorname{Perf}(R)$ とする。
$P\otimes_R^{\mathbf L}k$ が次数零の $r$ 次元ベクトル空間に
準同型ならば $P\simeq R^{\oplus r}[0]$ である。
同型の選択は非標準的でよいが、任意の $R$-代数 $R'$ に対して
$P\otimes_R^{\mathbf L}R'\simeq (R')^{\oplus r}[0]$ となる。
\end{lemma}
\begin{proof}
有界な有限自由複体で $P$ を表す。
微分行列に単元成分があるたびに基底変換して
$R\xrightarrow{\mathrm{id}}R$ の可縮直和因子を取り除く。
有限回で全微分の成分が $\mathfrak m$ に属する複体が得られる。
その剰余体への還元の微分はすべて零なので、仮定により
非零項は次数零の階数 $r$ の自由加群のみである。
従って元の複体は $R^{\oplus r}[0]$ と準同型である。
有限自由複体の可縮因子の消去は任意の基底変更後も保たれる。
\end{proof}

\begin{lemma}[五頂点 bar 解消の持上げ]
\label{lem:five-window-bar-lifting}
次のいずれかを仮定する。
\begin{enumerate}
\item $R$ は可換局所環であり、全 $C_{ij}$ は有限自由である。
  以下の fiber 条件を剰余体上で仮定する。
\item $R$ は任意の可換環であり、全 $C_{ij}$ は有限 projective である。
  以下の fiber 条件を全ての幾何点上で仮定する。
\end{enumerate}
fiber 条件は、各 $S_i$ が長さ三の有限生成 projective 解消を持ち、
その項に現れる表現可能加群が $P_i,\ldots,P_{i+4}$ に限られ、
最後の項が $P_{i+4}$ 一個であり、さらに
\[
 \operatorname{Ext}_{C_{\bar x}}^q(S_i,P_j)=0
 \quad\text{unless }(q,j)=(3,i+4),\qquad
 \operatorname{Ext}_{C_{\bar x}}^3(S_i,P_{i+4})\simeq\kappa(\bar x)
 \tag{BW-fiber}\label{eq:five-window-fiber-dual}
\]
となることである。最後の項以外の極小 Betti 数は一定と仮定しない。

このとき $S_i$ は $C$ 上 perfect で、projective dimension は高々三である。
解消と AS 双対条件は任意の $R$-基底変更後にも成立し、
\[
 \mathbf R\operatorname{Hom}_{C}(S_i,P_j)\simeq
 \begin{cases}\mathcal L_i[-3],&j=i+4,\\0,&j\ne i+4\end{cases}
 \tag{BW-dual}\label{eq:five-window-relative-dual}
\]
である。ここで $\mathcal L_i$ は可逆 $R$-加群である。
\end{lemma}

\begin{proof}
まず指定された体上の解消を、頂点を保って極小化する。
同一頂点間の可逆な次数零成分を消去するだけなので、支持は依然として
$i,\ldots,i+4$ に含まれる。
$j\le i+4$ に対して
$\operatorname{Hom}(P_j,P_{i+4})$ は $j=i+4$ の場合にだけ非零である。
さらに極小微分の同一頂点成分は零である。
従って極小解消中の $P_{i+4}$ の次数 $p$ における重複度は
$\dim\operatorname{Ext}^p(S_i,P_{i+4})$ に等しい。
\eqref{eq:five-window-fiber-dual} は、この重複度が $p=3$ で一、
それ以外で零であることを保証する。
特に低い解消次数にも $P_{i+4}$ が現れ得る、という可能性はここで排除される。

$W_i$ を頂点 $i,\ldots,i+4$ の有限部分代数とし、
$E_i=\bigoplus_{a=i}^{i+4}Re_a$、$J_i=(W_i)_{>0}$ と置く。
$S_i$ の $E_i$ に相対的な normalized bar 解消を $C$ へ誘導する。
得られる増大複体の第 $p$ 項は
\[
 \mathcal B_{i,p}=
 \bigoplus_{i=i_0<i_1<\cdots<i_p\le i+4}
 C_{i_0i_1}\otimes_R\cdots\otimes_R C_{i_{p-1}i_p}
       \otimes_RP_{i_p},\qquad 0\le p\le4,
 \tag{BW-bar}\label{eq:five-window-bar}
\]
である。$p=0$ は $P_i$ とする。
微分は隣接因子の積および最後の $P_{i_p}$ への作用による交代和である。
全ての係数は有限 projective であり、各頂点成分は有界な有限 projective
$R$-複体となる。

指定された fiber 上で、仮定した長さ三の解消を $W_i$ へ制限する。
現れる全 projective の添字が窓に含まれるため、これは $W_i$ 上でも
projective 解消である。その bar 解消との homotopy 同値を $C$ へ誘導すると、
誘導された bar 複体は元の $C$ 上の解消と homotopy 同値になる。
従って \eqref{eq:five-window-bar} の fiber は全頂点成分で完全である。

有限 projective 複体の fiber 判定により、増大 bar 複体の各頂点成分は
分裂完全である。局所環の場合には補題~\ref{lem:local-perfect-fibre-concentration}を
各頂点の増大複体に適用する。
任意の可換基底の場合には、全幾何 fiber の完全性を
各素イデアルにおける剰余体への忠実平坦な拡大から確認して
局所環で同じ操作を行う。
これは perfect 複体の所定次数・所定階数への集中の開性
\cite[Tag 0BDK, Lemma 36.31.3]{StacksPerfectFibers}
の初等的な場合である。
従って bar 複体は全 $C$ 上で完全であり、任意基底変更後にも完全である。
無限個の頂点に対する開近傍を交叉していない。
局所環の場合には常に同じ局所環を使い、一般基底の場合には既に全 fiber に
仮定を置いている。

長さ四を三へ短縮できることを確認する。
第4項は $P_{i+4}$ の有限 projective 係数倍だけからなる。
微分の第3項にある $P_{i+4}$ 係数への成分は、fiber 上で
$\operatorname{Tor}^{W_i}_4(S_i,S_{i+4}^{\ell})=0$ により単射である。
ここで $S_{i+4}^{\ell}$ は左単純加群である。
従ってこの係数写像は局所直和因子への単射であり、
第4項を projective な可縮直和成分として消去できる。
残る長さ三の複体は有限生成 projective からなる。
局所的な短縮から大域的にも projective dimension が高々三であることが従う。
実際、元の長さ四の有限 projective 解消の第3 syzygy は有限表示で、
中心に沿う局所化で projective となるので projective である。

最後に $\operatorname{Hom}_C(\mathcal B_{i,\bullet},P_j)$ を取る。
各項は有限 projective $R$-加群であり、その指定 fiber は
\eqref{eq:five-window-fiber-dual} によって零または階数一の単一 cohomology を持つ。
同じ有限複体の消去により \eqref{eq:five-window-relative-dual} を得る。
この複体の準同型と line の同定は任意 tensor に保たれる。
\end{proof}

\begin{remark}[二次曲面の固定階数の五頂点複体]
\label{rem:five-window-fixed-ranks}
$\operatorname{rank} C_{i,i+r}=1,2,4,6,9$（$r=0,1,2,3,4$）の場合、
\eqref{eq:five-window-bar} は局所的に
\begin{align*}
0\longrightarrow P_{i+4}^{16}&\longrightarrow
 P_{i+3}^{8}\oplus P_{i+4}^{48}
 \longrightarrow P_{i+2}^{4}\oplus P_{i+3}^{16}\oplus P_{i+4}^{40}\\
 &\longrightarrow P_{i+1}^{2}\oplus P_{i+2}^{4}
       \oplus P_{i+3}^{6}\oplus P_{i+4}^{9}
 \longrightarrow P_i\longrightarrow S_i\longrightarrow0
\end{align*}
となる。最終添字 $i+4$ の係数複体の階数は、homological degree
$1,2,3,4$ の順に $9,40,48,16$ である。
極小形が cubic 型から noncubic 型へ変わっても、この bar 複体の階数は変わらない。
本節の方法が要求するのはこの固定階数の複体の完全性である。
\end{remark}

\begin{lemma}[五頂点窓の torsion 核と compact 性]
\label{lem:five-window-compact-localization}
補題~\ref{lem:five-window-bar-lifting} の仮定の下で、
$\operatorname{Gr}C$ の次数 torsion 加群の Serre 部分圏を
$\operatorname{Tors}C$ とすると、
\[
 D(\operatorname{QGr}C)\simeq
 D(\operatorname{Gr}C)/\operatorname{Loc}\{S_i\mid i\in\mathbb Z\}.
 \tag{BW-loc}\label{eq:five-window-compact-localization}
\]
各 $\mathcal O_i=\pi P_i$ は $D(\operatorname{QGr}C)$ の compact 対象であり、
それら全体は同圏を localizing subcategory として生成する。
この結論に $C$ の Noetherian 性は不要である。
\end{lemma}
\begin{proof}
正の向きの $\mathbb Z$-代数の次数 torsion 加群は、
斉次元が生成する有限頂点支持部分加群の filtered colimit である。
有限頂点支持加群は最大の頂点から順に濾過でき、
逐次商は $S_i\otimes_R M_i$ と書ける。
任意の $R$-加群 $M_i$ の自由解消を使うと、
$S_i\otimes_RM_i$ は $\operatorname{Loc}\{S_i\}$ に属する。
$S_i$ は頂点 $i$ に $R$ を置く $R$-平坦加群なので、
ここで通常 tensor と導来 tensor は一致する。
filtered colimit の完全性から全 torsion 加群もこの localizing subcategory に属する。
cohomology が torsion である任意の導来複体は、
有限範囲の標準切断と、それらの homotopy colimit により回復できるため、
同じ localizing subcategory に属する。

Gabriel 商 $\pi:\operatorname{Gr}C\to\operatorname{QGr}C$ は完全で、
その右随伴 $\omega$ は fully faithful である。
$\omega$ は単射対象を単射対象へ送るため、
導来 counit $\pi\mathbf R\omega\to\mathrm{id}$ は同型となる。
従って導来局所化の核は torsion cohomology を持つ複体全体であり、
上記の Verdier 商を得る。

$D(\operatorname{Gr}C)$ は projective representables $P_i$ を
compact generators とする。
前補題により $S_i$ は有限 projective 解消を持ち compact である。
Neeman \cite[Theorem 2.1]{Neeman1992} の compact localization 定理により
商でも $\pi P_i$ は compact となる。
$\pi$ は直和を保ち、$P_i$ が元の導来圏を生成するので、
これらの像は商も生成する。
\end{proof}

\begin{proposition}[五頂点解消による canonical tails]
\label{prop:five-window-canonical-tails}
前補題の仮定に加え、$C$ が距離高々二の成分で生成されるとする。
$\mathcal O_i=\pi P_i\in\operatorname{QGr}C$ と置くと、
\[
 \operatorname{Ext}^q(\mathcal O_i,\mathcal O_j)\simeq
 \begin{cases}
 C_{ji},&q=0,\\
 \operatorname{Hom}_R(C_{i,j-4},\mathcal L_{j-4}),&q=2,\ j\ge i+4,\\
 0,&\text{otherwise}
 \end{cases}
 \tag{BW-coh}\label{eq:five-window-tails-cohomology}
\]
である。$\mathbf R\omega$ の cohomological dimension は高々二であり、
任意の連続四項の直和 $T_a=\bigoplus_{i=a}^{a+3}\mathcal O_i$ は
$D(\operatorname{QGr}C)$ の compact tilting generator となる。
従って
\[
 D(\operatorname{QGr}C)\simeq D(\operatorname{Mod}\operatorname{End}(T_a)),
 \qquad D(\operatorname{QGr}C)^c\simeq\operatorname{Perf}\operatorname{End}(T_a).
 \tag{BW-Morita}\label{eq:five-window-tails-morita}
\]
全構成は任意の導来基底変更と整合する。
Noetherian 性はこの命題の仮定にも結論にも含めない。
\end{proposition}

\begin{proof}
定理~\ref{thm:relative-cubic-tails-tilting} の局所化の証明が、
極小解消の形によらず適用できることを確認する。
有限 truncation $P_i/(P_i)_{\ge n}$ は
$C_{il}\otimes_RS_l$ を逐次商に持つ有限 filtration により
projective dimension 高々三の perfect 対象である。
従って tail $(P_i)_{\ge n}$ は perfect で projective dimension 高々二となる。
距離二までの有限生成性から全 tail は有限生成で、次数 torsion は localizing
subcategory をなす。

torsion radical の各頂点での導来関手は、有限 truncation からの Ext の
filtered colimit として計算される。
上に有界な対象への十分高い tail からの全 Ext は零であり、
一般の torsion 対象はこのような有界部分対象の filtered colimit である。
有限 truncation の perfectness を使えば正次数の導来 torsion は
torsion 対象上で消える。従って section functor も tail からの Ext の
filtered colimit で計算でき、$\operatorname{cd}\tau\le3$、
$\operatorname{cd}\omega\le2$ を得る。
各導来関手は直和と可換する。有限 cohomological dimension により、
unbounded $\mathbf R\omega$ も直和と可換し、
\[
 \mathbf R\tau X\longrightarrow X\longrightarrow\mathbf R\omega\pi X\longrightarrow
\]
が成立する。全 $\mathcal O_i$ の compact 性と生成性は、
この直和保存と独立に補題~\ref{lem:five-window-compact-localization} からも従う。
\eqref{eq:five-window-relative-dual} と有限 truncation の filtration を使うと
\eqref{eq:five-window-tails-cohomology} を得る。
また同じ AS 双対計算により $\mathbf R\tau P_j=N_j[-3]$ であり、
\[
 (N_j)_i=\begin{cases}
 \operatorname{Hom}_R(C_{i,j-4},\mathcal L_{j-4}),&i\le j-4,\\
 0,&i>j-4.
 \end{cases}
\]
したがって $N_j$ は上に有界な次数 torsion 加群である。
torsion 核全体の記述と compact 性は
補題~\ref{lem:five-window-compact-localization} で独立に証明した。
ここでは無界複体を有界切断の無条件な逆極限で置き換えていない。

四項による生成性のため、bar 複体を最終添字で filter する。
最終添字が $i+4$ の係数複体を $V_i^\bullet$ とする。
これは $W_i$ の bar 解消に $-\otimes_{W_i}S_{i+4}^{\ell}$ を適用した複体である。
前補題の証明で確認した極小解消の終頂点条件から、
指定 fiber は階数一の $\operatorname{Tor}_3$ だけを持つ。
相対的にも、支持が $i,\ldots,i+4$ に限られることから
\[
 \operatorname{Hom}_R(V_i^\bullet,R)
   =\operatorname{Hom}_C(\mathcal B_{i,\bullet},P_{i+4})
   \simeq\mathcal L_i[-3].
\]
有限 projective 複体の双対を再び取れば
$V_i^\bullet\simeq\mathcal L_i^{\vee}[3]$ となる。
ここでは bar 第 $p$ 項を cohomological degree $-p$ に置く。
最終添字 $\le i+3$ の部分複体は $P_i,\ldots,P_{i+3}$ の thick 閉包に属する。
bar 複体全体の tails 像は零だから、上の filtration の三角により
$\mathcal O_{i+4}$ も $\mathcal O_i,\ldots,\mathcal O_{i+3}$ の thick 閉包に属する。
逆方向には $P_i$ を bar 複体から分離すれば
$\mathcal O_i$ が次の四項で生成される。
従って任意の四項窓は全 $\mathcal O_i$ を生成する。
四項間の higher Ext は \eqref{eq:five-window-tails-cohomology} で零なので
導来 Morita 定理~\cite[\S\S4.3, 8.2]{Keller1994} が適用できる。

全解消は任意基底変更で完全である。導来拡大は $P_i,S_i$ を変更後の同じ対象へ送り、
$\operatorname{loc}\{S_i\}$ による局所化へ降りる。
四項窓の自己準同型環も通常の tensor で引き戻されるため、
\eqref{eq:five-window-tails-morita} は実際の導来基底変更と整合する。
\end{proof}

\subsection{良い noncubic 根の局所環上での全軌道と heart}

\begin{theorem}[noncubic 局所環における全次数 tails 認識]
\label{thm:noncubic-local-tails-effectivity}
系~\ref{cor:noncubic-etale-morita-root} の $(U,u)$ と Morita 根 $\mathsf s$ を取り、
$R=\mathcal O_{U,u}$ とする。
より一般に、閉点を $u$ へ送る Noetherian 局所 $\mathcal O_{U,u}$-代数を $R$ としてもよい。
$\Lambda=\Lambda_U\otimes R$ とし、
\[
 E_{2a}=\mathsf s^aF_0,\quad E_{2a+1}=\mathsf s^aF_1,
 \quad C_{ij}=\operatorname{Hom}_{\Lambda}(E_{-j},E_{-i})\quad(i\le j)
\]
と置く。$C_{ij}=0$（$i>j$）とし、積は射の合成とする。
次が成立する。
\begin{enumerate}
\item 全 $i\le j$ で $\mathbf R\operatorname{Hom}(E_i,E_j)$ は有限自由 $R$-加群の
  次数零成分に集中する。階数は閉 fiber の normal helix と等しく、
  偶数差 $2n$ で $(n+1)^2$、奇数差 $2n+1$ で $(n+1)(n+2)$ である。
  この同定と積は任意の導来基底変更と整合する。
\item 全軌道環
  $B=\bigoplus_{n\ge0}\operatorname{Hom}(F_0,\mathsf s^nF_0)$ は
  一次生成・二次・相対 Koszul で、全 $B_n$ は階数 $(n+1)^2$ の自由加群である。
  $B$ は左右 Noetherian であり、任意の有限型可換 $R$-代数への変更後も同様である。
\item 任意の可換 $R$-代数 $R'$ について canonical な比較
  \[
  D(\operatorname{QGr}B_{R'})\simeq D(\operatorname{Mod}\Lambda_{R'}),
  \qquad D(\operatorname{QGr}B_{R'})^c\simeq\operatorname{Perf}\Lambda_{R'}
  \]
  があり、$\pi B_{R'}$ を $F_0$ に、次数移動 $(1)$ を $\mathsf s$ に送る。
  特に compact 圏の canonical Serre 関手について $(1)^2\simeq S^{-1}[2]$ である。
\item 有限型可換 $R$-代数 $R'$ 上で $\mathcal C_{R'}=\operatorname{qgr}B_{R'}$ は
  Noetherian な相対 Abel 圏の系をなし、module-object 基底変更と fppf descent を持つ。
  coherent 対象間の Hom は有限で平坦基底変更と可換する。
  $\pi B_{R'}$ は平坦かつ exceptional で、$(\pi B_{R'},(1))$ は AZ ample である。
  左右の全 $\chi_q$ が成立し、canonical compact 圏は $\operatorname{Perf}\Lambda_{R'}$ である。
\end{enumerate}
この定理だけには同じ根の AS4 被覆を含めない。標準局所基底では、
後の定理~\ref{thm:noncubic-local-central-hull} がその被覆を与える。
また $R$ は一般に元の局所有限型 $k$-site の対象ではないため、
以上をそのまま $U$ の開近傍上の $\mathfrak H_2$ 族と呼ばない。
\end{theorem}

\begin{proof}
$\Lambda$ は相対 smooth proper であり、$E_i$ は perfect なので、
任意の $i,j$ に対する derived Hom は perfect $R$-複体である。
閉 fiber の列は定理~\cite[\texttt{thm:normal-helix}]{QuadricIntrinsic2026} の geometric helix である。
したがって前向き Hom の閉 fiber は指定階数の次数零自由加群である。
補題~\ref{lem:local-perfect-fibre-concentration}を
各組 $(i,j)$ の derived Hom に適用すれば、
同じ局所環 $R$ 上で全ての前向き Hom は次数零の自由加群となる。
自由基底の選択は非標準的でも、derived Hom の基底変更写像は canonical であり、
任意の $R$-代数への変更後も次数零への集中を保つ。
無限個の $(i,j)$ への個別適用は無限個の有限型開を交叉することではない。
全有限型開への延長は後の有限型効果性定理の別の内容である。

同じ Nakayama の議論により $C$ は距離高々二の成分で生成され、
奇数始点では次の偶数頂点を経由する射だけで正距離成分が生成される。
これは対応する全射の余核が各終点について有限生成で、閉 fiber で零だからである。
系~\cite[\texttt{cor:normal-cubic-defect}]{QuadricIntrinsic2026} の特殊 fiber の解消は
\[
 0\to P_{i+4}\to P_{i+3}^{2}\oplus P_{i+2}^{d_i}
 \to P_{i+1}^{2}\oplus P_{i+2}^{d_i}\to P_i\to S_i\to0,
 \quad d_i=\begin{cases}1&i\text{ even},\\0&i\text{ odd}\end{cases}
\]
であり、その AS 双対条件も成立する。
従って補題~\ref{lem:five-window-bar-lifting} と
命題~\ref{prop:five-window-canonical-tails} が適用できる。
相対的に $P_{i+2}$ の追加項を極小項として延ばす必要はない。
その二項が generic fiber で可縮になることも非極小 bar 複体が許している。

$B$ の各次数は (1) により有限自由である。
積 $B_1\otimes B_n\to B_{n+1}$ は特殊 fiber で全射なので全射となる。
$K=\ker(B_1^{\otimes2}\to B_2)$ は階数七の直和因子である。
$T_R(B_1)/(K)\to B$ の各次数は有限生成加群から自由加群への全射で、
その kernel は特殊 fiber で零である。
分裂と Nakayama により kernel は零、従って $B$ は二次である。
特殊 fiber の Koszul 性は命題~\ref{prop:all-h2-antiserre-koszul} による。
最初の三次数は自由なので、補題~\ref{lem:relative-koszul-degree-three} の証明にある
局所 Koszul 変形原理を $R$ へ直接適用し、全相対 Koszul 解消と任意底変更整合性を得る。
特殊 fiber は左右 Noetherian であるから、
補題~\ref{lem:local-commutative-fibre-noetherian} が $B$ の左右 Noetherian 性を与える。
有限型基底変更には補題~\ref{lem:local-finite-type-strong-noetherian} を使う。

偶数頂点 restriction は $\operatorname{QGr}C\simeq\operatorname{QGr}B$ を与える。
実際、十分長い道は必ず偶数頂点を通り、奇数始点の射は直後の偶数頂点を経由する。
restriction の左随伴 induction の unit は同型、counit の kernel と cokernel は
偶数成分が零なので torsion である。induction も torsion を torsion へ送る。
これらの事実は全基底変更後に保たれる。
四項窓 $i=-3,\ldots,0$ の End 代数は、合成も含めて元の $\Lambda$ である。

全軌道対象と次数移動の比較を明示する。
$P_i\mapsto E_{-i}$ とする cocontinuous derived functor を取る。
前向き Hom の次数零集中により、有向 dg 部分圏の truncation と $H^0$ の zigzag を通じて
この functor を構成できる。
各 $S_i$ の五頂点 bar 解消の像は有限個の perfect $\Lambda$-対象からなる複体である。
特殊 fiber 上では normal helix の tails 比較の下で単純 torsion 加群の像なので零となる。
よって perfect Nakayama により $R$ 上でも零である。
導来 torsion kernel が $\operatorname{loc}\{S_i\}$ であるため、functor は tails へ降りる。
四項 generator 上で元の $\Lambda$ の同定を与えるから、これは上の Morita 同値である。
この議論で同値にした heart を先に仮定していない。
全 $P_i$ の像と二周期移動が指定された $E_{-i}$ と $\mathsf s$ になるので、
$(1)\leftrightarrow\mathsf s$ および実際の平方同型も得る。
有限複体の消滅と全成分の自由性は任意基底変更に保たれるため、(3) が従う。

有限型 $R'$ 上で $B_{R'}$ は左右 Noetherian である。
各 $\pi P_i$ の偶数 restriction は有限生成 $B_{R'}$-加群の像なので
Noetherian な生成系をなす。
式~\eqref{eq:five-window-tails-cohomology} と
補題~\ref{lem:relative-cubic-coherent-cohomology} の証明は
Noetherian な生成系とこの cohomology 公式だけを使う。
従ってそのまま一般 coherent Hom の有限性・平坦基底変更・一様消滅に適用できる。
偶数 truncation の生成性と合わせて AZ 豊富性を得る。
また補題~\ref{lem:relative-cubic-global-chi} の証明も同じ torsion 幅と消滅を用いるため適用できる。
左側には反対代数と添字反転に対する同じ bar 議論を使う。
module-object 基底変更は補題~\ref{lem:qgr-module-object-basechange}、
fppf descent は完全かつ保守的な平坦引戻しによる既存の証明を用いる。
構造対象の平坦性は $B_n$ の自由性、exceptional 性は四項窓の Hom 公式による。
tilting summand は基底上平坦なので、非平坦変更にも実際の導来引戻しが
Morita 表示の $-\otimes_{\Lambda_{R'}}^{\mathbf L}\Lambda_{R''}$ と一致する。
これで (4) を得る。
\end{proof}

\begin{remark}[局所定理と有限型一様化の役割]
有限型近傍上の各軌道 Hom は perfect であり、固定した有限個については開性を使える。
本定理は、その無限個の開条件の共通部分が開であるとは述べない。
AS4 被覆の延長が得られても、有限型近傍で使う全軌道の一様平坦性・表示比較と
全幾何 fiber の rank-one torsion-free 条件は別途確認する。
一方、同じ局所環上では全次数の bar 解消と Noetherian heart を既に得たので、
以後の局所 AS4 構成で canonical compact 圏を独立の未証明入力に戻す必要はない。
有限型一様化は定理~\ref{thm:noncubic-finite-type-effectivity} で、
全幾何 fiber の有限 foundation 認識を新しい入力として証明する。
\end{remark}


\subsection{中央被覆の局所延長と全 Artin 認識}

\begin{lemma}[cubic 根の canonical hull の scalar normalizer]
\label{lem:cubic-compatible-scalar-normalizer}
体 \(K\) 上の geometric cubic \(\mathbb Z\)-代数に実際の二周期 \(\beta\) があり、
対応する tails 自己同値 \(s\) が \(s^2\simeq S^{-1}[2]\) を満たすとする。
その偶数環 \(B\) の canonical hull の双対線 \(K\phi\subset B^!_2\) は、
ある \(\lambda\in\{1,-1\}\) に対して
\[
 \phi x=\lambda x\phi\qquad(x\in B^!_1)
\]
を満たす。ここでは最初から \(\lambda=1\) とは主張しない。
\end{lemma}
\begin{proof}
canonical 四周期を \(\theta\) とする。
補題~\cite[\texttt{lem:cubic-serre-period}]{QuadricIntrinsic2026} により、これは tails 上で反 Serre 関手を表す。
定理~\ref{thm:universal-cubic-hull} の正規化式より、
\(\phi x=(x\circ T)\phi\)、\(T=\beta^{-2}\theta|_{B_1}\) である。
\(\beta^2\) と \(\theta\) は tails で同じ反 Serre 関手を表す。
従ってその比は tails に恒等関手を誘導する。
geometric helix の各対象の自己同型は scalar であり、自然同型をそれらの
対象と全ての前向き Hom に制限すれば、この比の \(\mathbb Z\)-代数上の作用は
頂点 scalar によるものとなる。特に \(B_1=D_{0,2}\) 上の \(T\) は
一つの非零 scalar \(\lambda\) である。
\(\phi\) は正則であり、次数二の元であるため、同じ交換式を二度使って
\(\phi^2=\lambda^2\phi^2\) を得る。従って \(\lambda^2=1\) である。
必要なら代数閉体へ拡大して上の比較を行い、scalar 性と式は忠実平坦に降下する。
\end{proof}

\begin{theorem}[局所基底上の noncubic 横断中央被覆]
\label{thm:noncubic-local-central-hull}
\((R,\mathfrak m,k)\) を標数零の体上の Noetherian 正則局所環、\(K=\operatorname{Frac}R\)
とする。二次環 \(B=T_R(V)/(J)\) について次を仮定する。
\begin{enumerate}
\item \(V,J\) は階数四、七の自由加群で、\(J\subset V^{\otimes2}\) は直和因子である。
全ての \(B_n\) は階数 \((n+1)^2\) の有限自由加群であり、\(B\) は相対 Koszul である。
特に \(E=B^!\) の全次数は有限自由、任意の底変更と可換し、その階数は
\(1,4,7,8,8,\ldots\) である。
\item 閉 fiber \(B_k\) は定理~\ref{thm:noncubic-central-hull} の良い根の
noncubic 軌道環である。
\item generic fiber \(B_K\) は geometric cubic helix を二項移す実際の根の偶数環であり、
その根の平方は反 Serre 関手に同型である。
\end{enumerate}
このとき \(B\) は相対 quadratic Koszul AS4 環 \(A\) と中央正則二次元 \(g\) により
\[
 B\cong A/(g)
\]
と表示される。全次数の有限自由性、左右および双加群 Koszul 解消、中央正則切断は
任意の基底変更と両立する。\(B\) が左右 Noetherian なら \(A\) もそうである。
この定理は有限型開近傍上の一様な Noetherian 性をそれ自体では主張しない。
\end{theorem}

\begin{proof}
二つの有限自由加群間の射
\[
 c_\pm:E_2\longrightarrow\operatorname{Hom}_R(E_1,E_3),\qquad
 c_\pm(z)(x)=zx\pm xz
\]
を取る。閉 fiber では \(c_+\) が単射であるため、階数七の minor が単元となる。
よって \(c_+\) は \(R\) 上で split 単射であり、generic fiber にも反中心はない。
前補題より generic fiber の canonical 正規線の scalar は \(\lambda=1\) となり、
従って \(c_-\otimes K\) の核は非零である。

閉 fiber では \(\dim\ker(c_-\otimes k)=1\) なので、階数六の minor は単元である。
generic fiber での核の非零性から七次 minors は \(K\) 上で全て零であり、
\(R\hookrightarrow K\) より \(R\) 上でも零である。
単元 minor で掃き出すと \(c_-\) は階数六の直和因子への射となる。
従って \(N=\ker c_-\) は階数一の直和因子であり、任意の底変更と両立する。
生成元を \(\phi\) とする。\(\phi_k\) は前定理の \(\chi\) の非零倍、
\(\phi_K\) は cubic canonical 正規線の非零倍である。

\(P=E/(\phi)\) と置く。
各 \(E_n\) は有限自由であり、閉 fiber の \(\phi_k=\chi\) は前定理により全次数で正則である。
従って各乗法
\[
 E_{n-2}\xrightarrow{\phi\cdot}E_n
\]
は閉 fiber で単射であり、最大 minor の一つが \(R\) の単元となる。
よってこれは split 単射であり、その余核は有限自由となる。
\(\phi\) の中央性により余核は正確に \(P_n\) である。
これで \(\phi\) の全次数の左右正則性、\(P_n\) の自由性、任意の基底変更との整合を同時に得る。
階数の差を取ると \(P\) の階数列は \((1,4,6,4,1,0,\ldots)\) である。
ここでは未知の商の fiber 次元一定から平坦性を推論していない。
既に自由と分かっている \(E_{n-2},E_n\) 間の写像の特殊 fiber の単射性を使っている。

閉 fiber \(P_k\) は前定理により Koszul であり、\(P_1,P_2,P_3\) は自由である。
Polishchuk--Positselski の局所環上の定理
\cite[Chapter 6, Theorem 7.3]{PolishchukPositselski2005} により、
\(P\) は相対 Koszul で、双対 \(A=P^!\) の全次数は任意の基底変更と両立する有限自由加群となる。
Frobenius pairing は閉 fiber で非退化なので、その行列式は \(R\) の単元であり、
\(P\) は相対 Koszul Frobenius、階数列 \((1,4,6,4,1)\) となる。
従って \(A\) は所要の相対 AS4 型を持つ。

切断の中央性を tensor として確認する。
\(H\subset J\) を \(\phi\) の核とすれば \(A=T(V)/(H)\)、\(J/H\) の生成元を \(g\) とできる。
\(K_3=(J\otimes V)\cap(V\otimes J)=E_3^\vee\) と置く。
\(\phi\) の中央性は二つの縮約
\(c_L=\phi\otimes1,c_R=1\otimes\phi:K_3\to V\) の一致を意味する。
\(\phi E_1\subset E_3\) が直和因子への単射なので、この共通写像は全射である。
\(\phi(r)=1\) となる \(r\in J\) を選ぶと、任意の \(v\in V\) に対して
\[
 z=r\otimes v+h_L=v\otimes r+h_R\in K_3,
 \qquad h_L\in H\otimes V,\quad h_R\in V\otimes H
\]
を取れる。\(A_3\) 内で \(gv=vg\) となり、一次生成性から \(g\) は中央である。
商の提示から \(A/(g)=B\) である。
両者の全次数は自由で、\(h_B=(1-t^2)h_A\) だから、
各次数の \(g\) 乗法は等階数自由加群間の全射として核零となる。
その像と商も自由なので任意の底変更後に正則性を保つ。
最後に \(B\) が Noetherian なら、中央正則元の有限内部次数フィルターを使って
\(A\) へ左右 Noetherian 性を持ち上げられる。
\end{proof}

\begin{remark}[Artin 結論と有限型一様性]
この定理の入力を有限 foundation の局所環上で確認できれば、
任意の Artin 厚みへの引戻しは、同じ良い根の AS4 被覆を与える。
しかし局所環上の全次数平坦性・Noetherian 性の証明だけから、
単一の有限型開近傍上で全次数の認識が成立すると推論しない。
有限型近傍の全 orbit、canonical tails と普遍 Noetherian 性はそれぞれ別の入力である。
\end{remark}

\begin{corollary}[全 noncubic 点の Artin 認識]
\label{cor:noncubic-all-artin-h2}
標数零の代数閉体上の任意の noncubic \(\mathcal H_2\) 対象
\((\mathcal C,\mathcal O)\) を取る。
その任意の flat pointed Abel Artin 変形は、指定した良い根の持上げの全軌道環に対する
中央正則二次 AS4 被覆を持ち、全有限型基底変更を備えた相対 \(\mathfrak H_2\) 族へ昇格する。
従って、局所 Artin 基底上で相対 \(\mathfrak H_2\) 変形から pointed Abel 変形への忘却は
変形 groupoid の同値である。
根対が異なる点では接次元三・無限小慣性零、根対が一致する点では接次元四・無限小慣性一次元であり、
いずれも変形は無障害である。
\end{corollary}
\begin{proof}
有限 foundation の滑らかな既約 chart \(S\) とその noncubic 点 \(x\) を取り、
系~\ref{cor:noncubic-etale-morita-root} の有限型 étale 近傍 \((U,u)\to(S,x)\) で
良い根を実際の perfect Morita 双加群として構成する。
\(u\) を含む既約成分へ制限し、\(R=\mathcal O_{U,u}\) と置く。
これは正則局所整域である。
定理~\ref{thm:noncubic-local-tails-effectivity} が全前向き Hom、全軌道環、
相対 Koszul 性、Noetherian 性、同じ根の canonical tails 比較を与える。
foundation の cubic 開部分は稠密であるから、この \(R\) の generic fiber は cubic である。
local tails 比較における \(E_{i+4}=S^{-1}E_i[2]\) と最初の四項窓の一致により、
generic fiber の列はその標準 cubic helix と同じである。
定義した列 \(E_{2a}=\mathsf s^aF_0\)、\(E_{2a+1}=\mathsf s^aF_1\) に対し
同じ根が実際に二項移動を与え、その二乗は反 Serre 関手であるため、
定理~\ref{thm:noncubic-local-central-hull} の全入力が成立する。
従って \(R\) 上に中央 AS4 被覆を得る。

任意の局所 Artin \(k\)-代数 \(T\) 上の pointed Abel 変形を取る。
定理~\ref{thm:artin-foundation-lift} によりその foundation を持ち上げ、
閉 fiber の指定同型に沿って各有限自由成分の基底を選ぶと、
閉点を \(x\) へ送る \(\operatorname{Spec}T\to S\) を得る。
\(U\to S\) の étale 性により、指定された閉点 \(u\) を保つ一意な lift
\(\operatorname{Spec}T\to U\) が存在し、その射は \(R\) を経由する。
\(R\) 上の全 orbit と被覆を \(T\) へ引き戻す。
根の Morita 類の一意性は系~\ref{cor:noncubic-root-all-artin} による。
Abel heart 上の exactness と AZ 豊富性は
補題~\ref{lem:artin-morita-heart-exact}、補題~\ref{lem:artin-ample-lift} による。
従って得た環は、元の Abel 変形における同じ根の全軌道環である。
命題~\ref{prop:artin-witness-relative-promotion} を適用すれば、
任意の有限型 \(T\)-代数を含む相対 \(\mathfrak H_2\) 族への昇格を得る。

全忠実性は module-object 基底変更系の一意性から従う。
根と AS4 被覆は存在条件なので、同値に追加の保存条件を課さない。
接と無限小慣性はそれぞれ \(\operatorname{HH}^2\)、\(\operatorname{HH}^1\) であり、
その次元は定理~\ref{thm:noncubic-hh-trivial-root} による。
\(\operatorname{HH}^{\ge3}=0\) は有限 foundation の計算
（命題~\ref{prop:foundation-hochschild}）によるため無障害である。
\end{proof}

\begin{remark}[形式的完了と代数性の区別]
全 cubic 点における既存の形式認識と前系を合わせると、全 \(\mathcal H_2\) 幾何点で
\(\mathfrak H_2\) と pointed Abel 変形の Artin groupoid が一致する。
この結論から内在的スタック全体の有限型代数性はまだ従わない。
局所環上で同時に成り立つ全軌道・Noetherian 条件を、単一の有限型 étale 近傍へ
一様に広げる段階と、その相対認識を別に閉じる必要がある。
\end{remark}

\subsection{中央二次元の全次数正則性の有限型一様化}

\begin{lemma}[有限な二次商から中央正則線を一様化する]
\label{lem:spread-central-regular-line}
$R$ を可換 Noetherian 環、$V$ を有限 $R$-加群とし、
$J\subset V^{\otimes2}$ とし、$E=T_R(V)/(J)$ を次数一生成の二次 $R$-代数とする。
特に $E_{<0}=0$、$E_0=R$、$E_1=V$ は有限生成である。
$\phi\in E_2$ は中央であると仮定し、$P=E/(\phi)$ と置く。
$P_0,\ldots,P_4$ は有限 projective $R$-加群、$P_5=0$ とする。
さらに、ある $\mathfrak p\in\operatorname{Spec}R$ において
$\phi$ の $E\otimes_RR_{\mathfrak p}$ 上の乗法が単射であると仮定する。
これは剰余体上だけの正則性ではなく、局所環上の正則性の仮定である。

このとき $f\in R\setminus\mathfrak p$ が存在し、単一の開 $D(f)$ 上で
\begin{enumerate}
\item $\phi$ は全次数で左右正則である。
\item 全ての $(E_f)_n$ は有限 projective $R_f$-加群であり、
  \[
  0\longrightarrow(E_f)_{n-2}\xrightarrow{\phi\cdot}(E_f)_n
       \longrightarrow(P_f)_n\longrightarrow0
  \tag{SC-exact}\label{eq:spread-central-regular-sequence}
  \]
  は $R_f$ 上分裂完全である。負次数は零とする。
\item 任意の可換 $R_f$-代数への変更後にも、中央正則性、
  上の全次数の完全列、商 $E_f/(\phi)=P_f$ が保たれる。
\end{enumerate}
さらに $\deg t=2$、$t$ の作用を $\phi$ とすれば、非標準的な次数付き同型
\[
 E_f\simeq\bigoplus_{i=0}^4(P_i)_f\otimes_{R_f}R_f[t](-i)
 \tag{SC-module}\label{eq:spread-central-projective-module}
\]
がある。従って $E_f$ は $R_f[t]$ 上有限 projective でもある。
各 $P_i$ が階数 $1,4,6,4,1$ を持つ場合、$E_n$ の階数は
$1,4,7,8,8,\ldots$ となる。
この補題は $P$ または $E$ の Koszul 性を結論しない。
したがって二次双対 $P^!$ の AS 正則性や全次数平坦性も、この補題だけからは従わない。
\end{lemma}

\begin{proof}
$P$ は次数一生成なので、$P_5=0$ から $P_n=0$（$n\ge5$）である。
各 $P_i$（$0\le i\le4$）の有限生成元を $E_i$ へ持ち上げる。
$t$ を次数二の中央変数とし、$t$ の作用を $\phi$ の乗法で定めると、
これらの有限個の持上げは $E$ を $R[t]$-加群として生成する。
実際、任意の斉次元をまず $P_n$ への像で消し、残りを
$\phi E_{n-2}$ と書く。この操作を内部次数について繰り返せば有限回で止まる。
この段階では $\phi$ の正則性を使っていない。

$R[t]$ は Noetherian で $E$ はその有限加群だから、
\[
 K=\ker\bigl(\phi\cdot:E(-2)\longrightarrow E\bigr)
\]
も有限 $R[t]$-加群である。
局所化は完全なので、仮定により $K\otimes_RR_{\mathfrak p}=0$ である。
$K$ の有限個の $R[t]$-生成元を $z_1,\ldots,z_m$ とする。
各 $z_a$ は $\mathfrak p$ の外のある $f_a\in R$ によって消えるため、
$f=\prod_af_a$ とすれば $K_f=0$ である。
$K=0$ の場合は $f=1$ とする。
この一つの分母が全内部次数の kernel を同時に消す。
中央性により左右の正則性は同じである。

中央性と $P=E/(\phi)$ から
\eqref{eq:spread-central-regular-sequence} が完全となる。
全 $P_n$ は projective で、$n\ge5$ では零だから、これらの列は
$R_f$-加群として分裂する。
$E_{<0}=0$ から二つの parity で帰納すれば、全 $E_n$ は有限 projective となる。
各 $i=0,\ldots,4$ で $E_i\to P_i$ の $R_f$-線形切断を選び、
これを $R_f[t]$-線形に延ばすと \eqref{eq:spread-central-projective-module} の右辺から左辺への射を得る。
両辺の次数 $n$ の成分は、次数 $n-2$ の成分と $P_n$ による
\eqref{eq:spread-central-regular-sequence} 型の短完全列を持つ。
従って次数に関する帰納法でこの射は同型となる。
切断の選択によるため同型は非標準的である。
分裂完全列は任意の tensor 後も完全なので、非平坦な基底変更についても (3) が従う。
階数の主張は
\[
 \sum_{n\ge0}\operatorname{rank}(E_f)_n\,z^n
 =\frac{1+4z+6z^2+4z^3+z^4}{1-z^2}
\]
から得られる。
\end{proof}

\begin{remark}[中央正則性と Koszul 性の入力]
この補題の有限性は、中央二次元による商が有限個の次数にしか台を持たないことから来る。
したがって $E_n$ ごとに別々の開近傍を選ぶ必要がない。
有限型近傍上で $P_0,\ldots,P_5$ の指定階数と中央性を確立できれば、
局所環で得た $\phi$ の正則性を、全次数について同じ開へ延ばせる。
一方、指定点での $P$ の Koszul 性をその開全体へ延ばすことは別である。
既存の局所 Koszul 変形原理を用いるには、適用先の各局所環の剰余体上の
Koszul 性、またはそれを与える一様な伝播補題がさらに必要である。
第\ref{sec:noncubic-finite-type}節では、全 fiber の有限認識からこの入力を供給し、
補題~\ref{lem:spread-local-central-hull} で同じ根の AS4 被覆まで延長する。
\end{remark}

% !TEX root = ../nc-quadric-moduli.tex
% 2026-11-21 / v1.36
\section{明示 noncubic 表示族の相対内在的実現}
\label{sec:relative-noncubic-families}

第~\cite[\texttt{sec:noncubic-stack}]{QuadricIntrinsic2026}節の中央二重平方族について、
canonical compact 圏と任意の Noetherian 基底変更を含む全相対条件を確認する。
重み付き射影スタック上の Clifford 代数として heart を表し、
その対角加群の perfectness を幾何ファイバーから持ち上げる。
これにより表示スタック $\mathfrak N^{\mathrm{pres}}$ から
内在的 $\mathfrak H_2$ への実際の射を得る。
全ての基底上平坦 coherent 対象が compact であるという追加条件は用いない。
任意の noncubic 特殊ファイバーを持つ相対族の正規形や、
cubic を横断する開近傍の効果性とは区別する。

\subsection{中央有限代数に対する fiberwise perfectness の持上げ}

\begin{lemma}\label{lem:nc-fiber-perfect}
$R\to D$ を Noetherian 可換環の準同型、$H$ を中心に $D$ の作用する
有限 $D$-代数、$M$ を有限 $D$-加群である右 $H$-加群とする。
$H,M$ は $R$ 上平坦と仮定する。
$\mathfrak q\in\operatorname{Spec} D$、$\mathfrak p=\mathfrak q\cap R$ とする。
対応する fiber の点において
\[
 M\otimes_R\kappa(\mathfrak p)
 \quad\text{が}\quad
 H\otimes_R\kappa(\mathfrak p)\text{ 上 perfect}
\]
ならば、$M$ は $\mathfrak q$ のある近傍で $H$ 上 perfect である。
\end{lemma}

\begin{proof}
まず $R$ を $R_{\mathfrak p}$、$D$ を $D_{\mathfrak q}$ に替える。
さらに $D$ の henselization へ忠実平坦に移る。
中心に沿う平坦 tensor は $R$-Tor を計算する複体と可換するので、
$H,M$ の $R$-平坦性はこれらの変更後も保たれる。
$D$ は Noetherian henselian 局所環、極大イデアルを $\mathfrak n$ と書く。
$R$ の極大イデアル $\mathfrak pR_{\mathfrak p}$ を以下 $\mathfrak m$ と書けば、
$\mathfrak mD\subset\mathfrak n$ である。

有限 $D$-代数 $H$ は semiperfect である。
この標準事実は非可換 $H$ にも適用される。
実際、$\mathfrak nH\subset\operatorname{Jac}(H)$ であり、$H/\mathfrak nH$ は有限次元代数である。
同代数の idempotent の lift $e\in H$ を取り、有限 $D$-生成系上の左乗法行列に
determinant trick を用いれば、$e$ を消す monic 多項式で、剰余が
$t^a(t-1)^b$ となるものを得る。
Hensel 分解と中国剰余定理により $e$ の多項式として idempotent を持ち上げられる。
corner 代数へ帰納的に適用して完全直交系を持ち上げると、
有限生成 $H$-加群は有限生成 projective cover を持つ。

有限自由右 $H$-加群による $M$ の部分解消を、fiber 上の射影次元以上の長さ $n$ まで取る。
$H$ は Noetherian（有限 $D$-加群）なので全 syzygy は有限である。
各段で商が $R$-平坦であることから、帰納的に全 syzygy も $R$-平坦となる。
従ってこの部分解消の fiber は完全であり、$n$ 番目の syzygy $K$ に対し
\[
 \overline K=K/\mathfrak mK
\]
は $\overline H=H/\mathfrak mH$ 上 projective である。
fiber の perfectness は幾何的剰余体への忠実平坦拡大でも検出できる。

最小 projective cover $P\twoheadrightarrow K$ を取り、核を $L$ とする。
$\mathfrak mH\subset\operatorname{Jac}(H)$ なので
\[
 \operatorname{Jac}(\overline H)=\operatorname{Jac}(H)/\mathfrak mH.
\]
したがって $\overline P\twoheadrightarrow\overline K$ も最小 projective cover である。
$\overline K$ 自身が projective なのでこれは同型となる。
一方 $K$ の $R$-平坦性から、$0\to L\to P\to K\to0$ を
$\kappa(\mathfrak p)$ へ tensor した列は左端でも完全である。
従って $L/\mathfrak mL=0$。
$L$ は有限 $D$-加群で $\mathfrak mD\subset\mathfrak n$ だから、Nakayama により $L=0$。
従って $K$ は projective であり、$M$ は有限 projective 解消を持つ。

有限表示加群の projectivity は中心に沿う忠実平坦拡大から降下するので、
henselization の前の $D_{\mathfrak q}$ 上でも $K$ は projective である。
有限表示と splitting の有限個の係数を広げれば、同じ性質が $\mathfrak q$ の近傍で成立する。
\end{proof}

\begin{remark}
この補題は、$R$ 上の平坦性を省いた fiber 判定ではない。
特に有限自由解消の syzygy の $R$-平坦性と、最後の
$\operatorname{Tor}_1^R(K,\kappa(\mathfrak p))=0$ の両方を使う。
基底の被約性、正則性、有限 Krull 次元は用いない。
\end{remark}

\subsection{noncubic heart の weighted stack 表示}

$k$ を標数零の体、$R$ を可換 Noetherian $k$-代数とする。
判別式が可逆な二元二次形式 $f_1,f_2\in R[u,v]_2$ を取り、
\[
 B=R[u,v]\langle c,d\rangle/(c^2-f_1,d^2-f_2),
 \qquad h=cd+dc,
 \qquad C=R[u,v,h],\quad\deg(u,v,h)=(1,1,2)
\]
と置く。$u,v$ は中心にある。
補題~\cite[\texttt{lem:noncubic-relative-clifford-center}]{QuadricIntrinsic2026} より
\[
 B=C\oplus Cc\oplus Cd\oplus C(cd)
\]
であり、この表示は全ての基底変更と両立する。
$X=\PP_R(1,1,2)=[(\operatorname{Spec} C\setminus V(C_+))/\mathbb G_m]$ とし、
$\mathcal B$ を $B$ の対応する $X$ 上の代数層とする。
\[
 \mathcal B\simeq\OO_X\oplus\OO_X(-1)^{\oplus2}\oplus\OO_X(-2)
 \quad\text{as }\OO_X\text{-modules}.
 \tag{NCF-1}\label{eq:nc-clifford-sheaf}
\]

\begin{proposition}\label{prop:nc-qgr-weighted}
構造対象を保つ完全同値
\[
 \operatorname{QGr} B\simeq\operatorname{QCoh}(X,\mathcal B),\qquad
 \qgr B\simeq\coh(X,\mathcal B)
\]
がある。これは任意の Noetherian 基底変更における module-object 比較と整合する。
coherent 対象間の Hom は有限 $R$-加群であり、任意の平坦基底変更と可換する。
対象・射・完全列について fppf descent が成立する。
\end{proposition}

\begin{proof}
$B$ は Noetherian 中央環 $C$ 上有限なので左右 Noetherian であり、
これは任意の Noetherian $R$-代数への引戻しでも成立する。
有限個の斉次 $C$-基底の次数は $0,1,1,2$ なので、
右 $B$-加群の元について $B_+$-torsion と $C_+$-torsion は一致する。
例えば $mC_{\ge r}=0$ なら $mB_{\ge r+2}=0$ である。
中心元 $u,v,h$ の局所化による sheafification の核もこの torsion である。
従って graded 加群の開部分への局所化と $\mathbb G_m$ 同変 descent が第一の同値を与える。
両辺の Noetherian 対象を取れば第二の同値を得る。

$X$ は $R$ 上 smooth proper な tame Deligne--Mumford stack である。
具体的に $u\ne0$、$v\ne0$ の二つの affine chart と、
$h\ne0$ の chart $[\mathbb A^2_R/\mu_2]$ で覆われる。
$2$ は可逆であり、$\mu_2$ と $\mathbb G_m$ の不変部分を取る操作は完全である。
したがって上の sheafification、局所化、同変 descent は任意の基底変更と整合する。
さらに affine 射 $X_{R'}\to X_R$ の下で、$R'$ 作用を持つ
$\mathcal B$-加群はちょうど $\mathcal B_{R'}$-加群である。
これは通常の引戻しが $R'\otimes_R-$ である module-object 基底変更である。

coherent $\mathcal B$-加群 $E,F$ に対して、
$\mathcal Hom_{\mathcal B}(E,F)$ は coherent $\OO_X$-加群である。
その大域切断の有限性は weighted projective Serre 有限性から従う。
この有限性は、有限 graded $C$-加群の偶数・奇数成分を
通常の次数一生成環 $C^{(2)}$ 上の有限加群として扱うことでも得られる。
平坦基底変更では、有限表示性により内部 Hom が基底変更し、
上の有限 chart の \v{C}ech 複体と完全な不変部分によって大域 $H^0$ も基底変更する。
fppf descent は各 affine chart の加群の忠実平坦降下を同変に貼り合わせればよい。
\end{proof}

\begin{lemma}[weighted cohomology]\label{lem:nc-weighted-cohomology}
任意の $j\in\mathbb Z$ について $\mathbf R\Gamma(X,\OO_X(j))$ は
有限 projective $R$-複体であり、任意の基底変更と可換する。
\[
 H^0(X,\OO_X(j))=C_j,\qquad H^1(X,\OO_X(j))=0,
 \qquad H^2(X,\OO_X(j))=C_{-j-4}^{\vee},
\]
ただし負次数の $C$ は零とする。
従って $\mathbf R\Gamma(X,\mathcal B(j))$ も同じ普遍性を持つ。
特に $n\ge0$ なら
\[
 \mathbf R\Gamma(X,\mathcal B(n))=B_n[0].
 \tag{NCF-2}\label{eq:nc-full-orbit-cohomology}
\]
\end{lemma}

\begin{proof}
$u,v,h$ の局所化の graded degree $j$ 成分を取る有限 \v{C}ech 複体を用いる。
Laurent 単項式の各指数ベクトルごとに分けると、
三指数が全て非負の場合に $H^0$、全て負の場合に $H^2$ が残り、
その他は係数 $\pm1$ の simplex 複体となって分裂完全である。
$H^2$ の基底を $u^{-a-1}v^{-b-1}h^{-c-1}$ と書けば、重みは $-a-b-2c-4$ である。
これが表示された双対を与える。
固定重みでは残る単項式は有限個で、収縮は $R$ に依存しない。
従って任意の非平坦な基底変更でも成立する。
\eqref{eq:nc-clifford-sheaf} へ適用し、$n\ge0$ では最小 twist が $-2$ となることを使う。
\end{proof}

\subsection{canonical compact 圏と相対 smooth proper 性}

\begin{proposition}\label{prop:nc-compact-generator}
\[
 G=\bigoplus_{i=0}^3\mathcal B(i)
\]
は $D(\operatorname{QCoh}(X,\mathcal B))$ の compact generator である。
その dg 自己準同型代数を $E_R=\mathbf R\End(G)$ とすると、
$E_R$ は $R$ 上 perfect であり、任意の Noetherian $R$-代数 $R'$ に対して
\[
 E_R\otimes_R^{\mathbf L}R'\simeq E_{R'},\qquad
 D(\operatorname{QGr} B)^c\simeq\operatorname{Perf}(E_R)
\]
となる。この同定は実際の module-object 引戻しの導来関手に対して
\[
 D(\operatorname{QGr} B)^c\otimes_{\operatorname{Perf} R}\operatorname{Perf} R'
 \simeq D(\operatorname{QGr} B_{R'})^c
 \tag{NCF-3}\label{eq:nc-perfect-basechange}
\]
を与える。compact 対象は $D^b(\qgr B)$ に属する。
\end{proposition}

\begin{proof}
忘却関手は直和と可換し、
\[
 \mathbf R\Hom_{\mathcal B}(\mathcal B(i),M)
 =\mathbf R\Gamma(X,M(-i))
\]
は有限 \v{C}ech 複体と完全な不変部分で計算できる。
従って全ての $\mathcal B(i)$ は compact である。
全整数の $\OO_X(i)$ は $D(\operatorname{QCoh} X)$ の生成系である。
実際、全ての $\mathbf R\Gamma(X,N(i))$ が零なら、
$u$ または $v$ の乗法に沿う filtered colimit を取ると二つの affine chart 上での
導来切断が零となる。
$h$ の乗法では偶数・奇数の二つの colimit を取れば、
$[\mathbb A^2/\mu_2]$ 上の二つの character 成分がともに零となる。
有限 \v{C}ech 複体はこれらの colimit と可換するため、$N$ は全 chart で零である。
この議論は複体を上下有界と仮定しない。
随伴により全ての $\mathcal B(i)$ も生成系となる。
さらに $u,v,h$ が $X$ 上共通零点を持たないことから、
\[
 0\to\mathcal B(n-4)\to
 \mathcal B(n-3)^{\oplus2}\oplus\mathcal B(n-2)\to
 \mathcal B(n-1)^{\oplus2}\oplus\mathcal B(n-2)\to
 \mathcal B(n)\to0
 \tag{NCF-4}\label{eq:nc-weighted-koszul}
\]
が完全である。これは中心の weighted Koszul 列を
局所自由な $\mathcal B$ へ tensor したものである。
上下の帰納法により任意の $\mathcal B(n)$ は $G$ の thick closure に属する。
従って $G$ は compact generator である。

各 dg Hom 成分は
$\mathbf R\Gamma(X,\mathcal B(j-i))$ なので前補題を適用できる。
積も元の $B$ の積と \v{C}ech の積から構成され、基底変更と整合する。
これで $E_R$ の proper 性と dg 代数としての基底変更を得る。
$G$ の各項は基底上平坦であり、実際の導来引戻しで $G_{R'}$ へ移る。
よって標準 dg Morita 比較の下でその引戻しは
$-\otimes_{E_R}^{\mathbf L}E_{R'}$ となり、\eqref{eq:nc-perfect-basechange} を得る。
最後に $G$ は coherent な heart の対象であり、その有限 cones と直和因子は
$D^b(\qgr B)$ に入る。
また canonical compacts は、局所的に有限生成 projective $\mathcal B$-加群の
有限複体で表される対象にちょうど等しい。
一方向は $G$ が局所自由であることから従う。
逆向きには、有限個の上の chart 上での有限 projective 表示を用いると、
内部導来 Hom が直和と可換し、一様に有限な cohomological amplitude を持つ。
有限 \v{C}ech 複体と完全な $\mu_2$ 不変部分を取ってもこの性質を保つので、
その対象は global に compact である。
同じ証明は $X\times_R X$ とその有限局所自由代数にも適用できる。
\end{proof}

\begin{remark}
$G$ を tilting と主張していない。
例えば負の差に \eqref{eq:nc-clifford-sheaf} を適用すると $H^2$ が現れる。
有限な通常の End 代数に置き換えず、$E_R$ は dg 代数として扱う。
\end{remark}

\begin{proposition}\label{prop:nc-relative-smooth}
上の canonical compact 圏は $R$ 上 smooth proper である。
この結論は任意の Noetherian 基底変更後に成立する。
\end{proposition}

\begin{proof}
proper 性は前命題による。
$Y=X\times_R X$ 上で
\[
 \mathcal H=p_1^*\mathcal B^{\mathrm{op}}\otimes_{\OO_Y}p_2^*\mathcal B,
 \qquad M=\Delta_*\mathcal B
\]
と置く。$\mathcal H$ は $\OO_Y$ 上有限局所自由で、$M$ は coherent である。
$X/R$ の diagonal $\Delta$ は有限であり、基底変更に対して pushforward は
通常の加群の制限として計算される。
従って $M$ も $R$ 上平坦である。
ここで stack の diagonal を closed immersion と仮定していない。

まず幾何ファイバー上の $M$ の perfectness を確認する。
既存の体上の逆構成
定理~\cite[\texttt{thm:ore-anticommuting-central}]{QuadricIntrinsic2026} により、
非退化な $f_1,f_2$ の全ての幾何ファイバーは $\mathcal H_2$ の対象である。
定理~\cite[\texttt{thm:normal-helix}]{QuadricIntrinsic2026} の四項 foundation を $T$ とする。
その canonical compact 性と非有界導来圏での生成性は、
定理~\cite[\texttt{thm:artin-foundation-lift}]{QuadricIntrinsic2026} を Artin 基底 $A=k$ に適用した場合である。
同定理の証明は局所 Noetherian 性・体上の大域次元二・DDL の compactness 条件を確認している。
従って $T$ はここで必要な非有界導来圏の compact tilting generator である。
対応する有限代数 $\Lambda$ は相対 normalized bar 解消により smooth である。
この入力から幾何ファイバー上の対角加群が perfect と分かることを、次のように具体化する。
$T$ とその左双対 $T^\vee=\mathbf R\mathcal Hom_{\mathcal B}(T,\mathcal B)$ の外積は、
右 $\mathcal H$-加群の compact generator を与える。
実際、前命題の $G$ は $T$ の thick closure に、逆に $T$ は $G$ の thick closure に属する。
左右の有限 cones と外積により、$T^\vee\boxtimes T$ は
$(\mathcal B^{\mathrm{op}}(-i))\boxtimes\mathcal B(j)$ の生成する圏を生成する。
後者は \eqref{eq:nc-weighted-koszul} を両因子へ適用すれば全ての compacts を生成する。
この外積の dg End は $\Lambda^{\mathrm{op}}\otimes_k\Lambda$ であり、
$M$ は Morita 表示で対角双加群 $\Lambda$ に対応する。
従って $\Lambda$ の有限双加群解消は $M$ の perfectness を与える。
これは幾何ファイバー上の主張であって、族の foundation を先に仮定していない。

$Y$ の任意の affine étale chart $\operatorname{Spec} D\to Y$ へ移る。
$D$ は Noetherian で、$\mathcal H$ は有限 $D$-代数、$M$ は有限 $D$-加群となる。
双方の $R$-平坦性は étale pullback で保たれる。
幾何ファイバーの perfectness は剰余体へ忠実平坦降下するので、
補題~\ref{lem:nc-fiber-perfect} が適用できる。
従って $M$ は $\mathcal H$ 上局所 perfect であり、étale descent により $Y$ 上 perfect である。
重み二の chart では $\mu_2$ が étale かつ linearly reductive なので、同じ議論がそのまま降下する。

最後に、$G^\vee\boxtimes G$ による積上の compact Morita 表示を取る。
その dg End は $E_R^{\mathrm{op}}\otimes_R^{\mathbf L}E_R$ であり、
対角 $M$ は $E_R$ の対角双加群に対応する。
従ってこの双加群は perfect であり、$E_R$ は $R$ 上 smooth である。
同じ入力は任意の Noetherian $R'$ 上で成立し、
また対角の有限解消は導来基底変更でも保存される。
\end{proof}

\subsection{同じ根・全軌道・相対 AS4 被覆}

\begin{proposition}\label{prop:nc-root-ample-orbit}
étale 局所に、$B$ の次数付き自己同型 $\sigma$ で
\[
 \sigma^2=\iota,\qquad
 \iota(u)=u,\quad\iota(v)=v,\quad\iota(c)=-c,\quad\iota(d)=-d
\]
となるものを選べる。
canonical compact 圏の相対 Serre 関手と、heart の完全自己同値は
\[
 S=(-)_\iota(-2)[2],\qquad s=(-)_\sigma(1),\qquad s^2=S^{-1}[2]
\]
を満たす。$(\mathcal B,s)$ は AZ ample であり、
右軌道積の全軌道環は $B^\sigma$ に等しい。
これらは任意の Noetherian 基底変更と整合する。
\end{proposition}

\begin{proof}
補題~\cite[\texttt{lem:noncubic-local-exchange}]{QuadricIntrinsic2026} により、étale 局所に二形式を交換する
対合 $L$ と単元 $\eta,\gamma$ を
\[
 L^2=1,\quad Lf_1=\eta f_2,\quad Lf_2=\eta^{-1}f_1,\quad\gamma^2=\eta
\]
と選び、$\sigma|_{Ru\oplus Rv}=L$、$\sigma(c)=\gamma d$、
$\sigma(d)=-\gamma^{-1}c$ と置ける。
判別式が可逆な二根を étale 局所で順序付け、第一根対を $\{0,\infty\}$、
第二根対を $\{a,1/b\}$ とすれば
$1-ab$ は単元である。
行列 $\left(\begin{smallmatrix}1&-a\\b&-1\end{smallmatrix}\right)$ の平方は
$(1-ab)I$ なので、単元の平方根を二度付加することで上の選択を行える。
これは二根対の一致点を含む。

命題~\cite[\texttt{prop:noncubic-relative-serre-on-perfect}]{QuadricIntrinsic2026} の Clifford trace は
\[
 \mathcal Hom_{\OO_X}(\mathcal B,\OO_X(-4))
 \simeq\mathcal B_\iota(-2)
\]
を与える。
前節で canonical compacts がまさに $\mathcal B$ 上 perfect な対象であることを確認したので、
この trace の相対双対性は求める canonical compact 圏の Serre 双対性である。
双対性は weighted \v{C}ech の residue pairing でも直接確認でき、
補題~\ref{lem:nc-weighted-cohomology} と同様に任意の基底変更と整合する。
平方の式は $\sigma^2=\iota$ から従う。

$s^n\mathcal B$ は右 $\mathcal B$-加群として $\mathcal B(n)$ と同型である。
AZ の生成条件は、有限生成 graded $B$-加群の高い truncation の有限生成性から従う。
全射の coherent kernel $K$ に対して
\[
 \Ext^1(\mathcal B(-n),K)=H^1(X,K(n))=0\quad(n\gg0)
\]
が成り立つ。
最後の消滅は中央 $C$ に関する Serre 消滅で、
偶数・奇数の二つの有限 $C^{(2)}$-加群に通常の Serre 消滅を適用すればよい。
これが AZ の Hom 全射条件を与える。

全軌道の次数 $n\ge0$ は \eqref{eq:nc-full-orbit-cohomology} により $B_n$ である。
$a\in B_m$ は $\mathcal B\to s^m\mathcal B$ の射
$x\mapsto a\sigma^m(x)$ に対応する。
従って $a*b=s^n(a)\circ b$ という右規約の積は
\[
 a*b=a\sigma^m(b)\qquad(a\in B_m,b\in B_n)
\]
となり、通常の Zhang twist $B^\sigma$ の積に一致する。
全次数・任意基底変更で同じ計算を用いる。
\end{proof}

\begin{proposition}[全相対 AS4 条件]\label{prop:nc-full-relative-as4}
前命題の選択の下で
\[
 A'=R[u,v]\langle c,d\rangle/
 \bigl((c^2-f_1)-\eta(d^2-f_2)\bigr),\qquad
 g'=(c^2-f_1)+\eta(d^2-f_2)
\]
と置く。
$A=(A')^\sigma$ と $g'$ は、定義~\ref{def:relative-h2}の相対 quadratic AS4 被覆の全条件を満たし、
\[
 A/(g')=B^\sigma
\]
である。特に全ての Noetherian 基底変更後も $A$ は左右 Noetherian であり、
左右・双加群 Koszul 解消および正則正規切断は任意の底変更後も完全である。
\end{proposition}

\begin{proof}
$x=c+\gamma d$、$y=c-\gamma d$ とすると $2\gamma$ は単元であり、
\[
 \begin{gathered}
 A'=R[u,v,x][y;\tau,\partial],\qquad
 \tau(u)=u,\quad\tau(v)=v,\quad\tau(x)=-x,\\
 \partial(u)=\partial(v)=0,\qquad\partial(x)=2(f_1-\eta f_2).
 \end{gathered}
\]
従って $u^av^bx^iy^j$ は $R$-基底で、
全次数は階数 $\binom{n+3}{3}$ の有限自由加群である。
可逆自己同型による Ore 拡大の左右 Hilbert 基底定理により、
$R$ が Noetherian なら $A'$ も左右 Noetherian である。
各 Noetherian $R'$ について同じ Ore 表示を持つため普遍 Noetherian 性が成立する。
Zhang twist は有限生成 graded 加群と graded Noetherian 性を保つ。
一般のイデアルは最高内部次数の初期項を用いるので、$A$ 自身も左右 Noetherian となる。

二次関係は monic な PBW 書換えであり、関係加群は $V^{\otimes2}$ の階数六の直和因子である。
既存の体上の Ore 計算は、全幾何ファイバーが quadratic Koszul AS 正則
型 $(1,4,6,4,1)$ であることを与える。
ここで相対結論を fiber 次元一定だけから導かない。
$A'_1,A'_2,A'_3$ は PBW 基底により既に自由なので、
補題~\ref{lem:relative-koszul-degree-three} の相対 Koszul 定理を適用できる。
二次双対の全次数は基底変更する有限 projective 加群となり、
階数は $1,4,6,4,1$、次数五以上は零である。
双対の積の Frobenius pairing は有限 projective 加群間の射であり、
全幾何ファイバーで同型なので基底上で同型となる。
左右 Koszul 複体は各内部次数で有限 projective 複体となり、
その fiber の完全性により局所分裂完全である。
双加群 Koszul 複体にも同じ判定を適用できる。
従ってそれらの完全性は任意の非平坦基底変更後も保たれ、
相対 AS 双対性も Frobenius pairing によって従う。
同じ議論は $A$ の Zhang twist についても成立する。

最後に
\[
 g'=\tfrac12(x^2+y^2)-(f_1+\eta f_2)
\]
は $A'$ の中央にある。
$y$ に関する最高係数が $1/2$ なので、$R$ に零因子があっても左右乗法は単射である。
その商の基底は $u^av^b$ と $c,d$ の交互語との積であり、全次数自由である。
また $\sigma(g')=g'$、$\sigma$ は $A'$ 上にも延びる。
twist 後は
\[
 g'*_\sigma a=g'\sigma^2(a),\qquad a*_\sigma g'=ag'
\]
なので $g'$ は正則正規であり、normalizer は $\sigma^2$ である。
元の正則乗法と自己同型の合成であるため、任意底変更後も正則性を保つ。
商は表示から $B^\sigma$ である。
\end{proof}

\subsection{表示スタックからの実際の相対族と残る境界問題}

\begin{theorem}\label{thm:noncubic-relative-effectivity}
第~\cite[\texttt{sec:noncubic-stack}]{QuadricIntrinsic2026}節の表示スタック
\[
 \mathfrak N^{\mathrm{pres}}=[P/G],\qquad
 P=Q^{\mathrm{nd}}\times Q^{\mathrm{nd}}
\]
上の降下した反 Serre 環 $\mathcal R$ は、任意の有限型 affine test $T$ 上で
\[
 (\mathcal C_T,\mathcal O_T)=(\qgr\mathcal R_T,\pi\mathcal R_T)
\]
という実際の相対 $\mathfrak H_2$ 族を定める。
従って自然な射
\[
 q_{\mathrm{nc}}:\mathfrak N^{\mathrm{pres}}\longrightarrow\mathfrak H_2
\]
が存在し、その全ての幾何ファイバーは noncubic である。
ここでは平方根・AS4 被覆は局所的 witness であり、出力の標識に加えない。
\end{theorem}

\begin{proof}
命題~\cite[\texttt{prop:noncubic-antiserre-descent}]{QuadricIntrinsic2026} により $\mathcal R$ は表示スタック上に実際に降下する。
affine 局所に行列式線束を自明化し、前節の $\sigma$ を選ぶと
$\mathcal R=(B^\sigma)^{(2)}$ である。
命題~\cite[\texttt{prop:noncubic-relative-heart-veronese}]{QuadricIntrinsic2026} の Veronese 比較は構造対象を保ち、
任意基底変更に整合する。
従って以下の全条件はこの fppf 局所表示で確認すればよい。

命題~\ref{prop:nc-qgr-weighted} が noetherian Abel 系、module-object 基底変更、
coherent Hom の有限性と平坦基底変更、fppf descent を与える。
命題~\ref{prop:nc-compact-generator} と命題~\ref{prop:nc-relative-smooth} が
canonical compact 部分の smooth proper 性、$D^b(\mathcal C_T)$ への包含、
任意導来基底変更を与える。
標識 $\mathcal B$ は $R$ 上平坦で、補題~\ref{lem:nc-weighted-cohomology} により
$\mathbf R\Hom(\mathcal B,\mathcal B)=R[0]$ である。
従って (RQ1) が成立する。

全幾何ファイバーでは、既存の noncubic 逆構成により Gabriel--Krull の
rank one torsion-free 条件を含む体上の全 $\mathcal H_2$ 条件が成立する。
fiber 圏は上の module-object 基底変更と一致するから (RQ2) が成立する。
命題~\ref{prop:nc-root-ample-orbit} は同じ局所 $s$ に対する相対 Serre 平方と AZ 豊富性を与える。
命題~\ref{prop:nc-full-relative-as4} はその全軌道環 $B^\sigma$ の相対 AS4 被覆を与える。
従って (RQ3), (RQ4) も成立する。

$\mathcal R$ 自体とその加群の降下は選んだ $\sigma$ に依存しない。
一方 $s$ と被覆の存在は相対 $\mathfrak H_2$ の定義により fppf 局所条件である。
よってこれらの局所的な検証は、表示スタック上の一つの相対系へ貼り合わさる。
また AZ 生成系は基底上平坦なので、補題~\ref{lem:relative-h2-flat-center} と同じ co-effaceability により
Lowen--Van den Bergh の flat Abel 性も従う。
\end{proof}

\begin{remark}[noncubic 境界との区別]
$q_{\mathrm{nc}}$ の存在だけでは、noncubic 点を含む開近傍の atlas は得られない。
特に任意の Artin 変形は、閉 fiber が noncubic でも
隣接乗法の四次 minor が非零の冪零元となり得る。
従って二形式表示の次元一のパラメータ族で全ての横断変形を捉えたことにはならない。
任意の相対 $\mathfrak H_2$ 族の noncubic 表示への局所還元、
この射の相対全忠実性・monomorphism 性、および cubic を横断する近傍の効果性は別途証明が必要である。
ここでの対角加群の議論は、明示的な中央有限な Clifford order を持つ族に適用したものであり、
一般の root-family が同じ中央有限表示を持つとは仮定していない。
\end{remark}

% !TeX root = ../nc-quadric-moduli.tex
% 2026-11-22 / v1.48（Clifford PBWとfpqc閉部分降下を分離）
\section{noncubic 境界を横断する有限型効果性}
\label{sec:noncubic-finite-type}

局所環上の構成を一つの有限型近傍へ広げる。
有限 foundation の Lagrangian 表示に全幾何点を認識する有限開条件を課し、
その同じ基底の全 fiber の geometric helix を用いて全 Hom を先に平坦化する。
二つの cubic 開部分と非被約 noncubic 閉部分から Noetherian 性を貼り合わせ、
局所中央被覆を有限型開へ延長する。
一般の Koszul 性の開性や、無限個の開集合の交叉の開性は用いない。
以下の相対結果とその組合せの新規性は未確認である。

\subsection{Lagrangian 開部分の全幾何点の認識}

以下では標数零の代数閉体上で考える。
前節の Grassmann 表示で、第二 summand の $D_L$ による基底変更を行い、
$H=U\oplus U^*$、$U=k^2\otimes k^2\simeq\operatorname{Mat}_2(k)$ と書く。
交代形式は二 summand 間の標準評価から作るものとする。
この基底変更は二次元因子上の同型なので分解テンソルの集合を保つ。
$\Sigma_U\subset\mathbb P(U)$ と $\Sigma_{U^*}\subset\mathbb P(U^*)$ を
二つの Segre 二次曲面とする。

\begin{lemma}[両階数三の Lagrangian 平面の直接正規形]
\label{lem:lagrangian-double-rank-three-recognition}
$L\subset U\oplus U^*$ を Lagrangian 四平面とする。
二つの射影の階数がともに三で、$L\cap U$ と $L\cap U^*$ の各生成元が
二次正方行列として可逆であると仮定する。
このとき $\operatorname{GL}_2\times\operatorname{GL}_2$ の $U$ 上の tensor 作用と
$U^*$ 上の反傾作用、および $L$ の基底変更により、$L$ は
\begin{equation}\label{eq:lagrangian-rank-three-normal-form}
\begin{aligned}
 L=\bigl\langle&(-E_{21},\alpha E_{12}+\gamma E_{21}),\quad
 (E_{12},\beta E_{12}-\alpha E_{21}),\\
 &(-D_0,0),\quad (0,I_2)\bigr\rangle,
 \qquad \alpha^2+\beta\gamma=1.
\end{aligned}
\end{equation}
となる。ここで $D_0=\operatorname{diag}(1,-1)$ である。
対応する有限 foundation は、中央二重平方環
\[
 k[u,v]\langle c,d\rangle/
 \bigl(c^2-uv,\ d^2-(\alpha u+\beta v)(\gamma u-\alpha v)\bigr)
\]
の既知の normal foundation と同型である。
この認識は、候補 foundation 自身の domain 性や $\mathcal H_2$ 所属を仮定しない。
\end{lemma}
\begin{proof}
交点線を可逆行列 $X\in U$、$Y\in U^*$ で生成する。
Lagrangian 条件から $\operatorname{tr}(X^tY)=0$ である。
作用を $X\mapsto gXh^t$、$Y\mapsto g^{-t}Yh^{-1}$ と書くと、
$X^tY$ は $h$ により共役変換される。
この行列は可逆かつ trace 零であり、標数零では異なる固有値 $\lambda,-\lambda$ を持つ。
従って基底変更と交点線の生成元の定数倍により
\[
 L\cap U=kD_0,\qquad L\cap U^*=kI_2
\]
とできる。
二射影の像は、それぞれ反対の交点線の annihilator なので、
$\langle D_0,E_{12},E_{21}\rangle$ と
$\langle I_2,E_{12},E_{21}\rangle$ である。
交点線の成分を引けば、残る二次元部分は
$\langle E_{12},E_{21}\rangle$ からその双対への graph となる。
Lagrangian 条件によりその行列は対称であり、交点が指定した二線だけであることから可逆である。
これを $Q=(q_{ij})$ と書く。

$U$ の scalar $t$ と $U^*$ の scalar $t^{-1}$ により $Q$ は $t^{-2}Q$ に移る。
代数閉性を用いて $t^4=-\det Q$ と選べば $\det Q=-1$ と正規化できる。
$\alpha=-q_{12}$、$\beta=q_{11}$、$\gamma=-q_{22}$ と置くと
$\alpha^2+\beta\gamma=1$ であり、graph の二基底と交点二線は
\eqref{eq:lagrangian-rank-three-normal-form} となる。
第二 summand の $D_L$ による基底変更を戻すと、これは
式~\ref{eq:nchh-plane} の四列の行列そのものである。
その式は中央二重平方環の二つの行列分解から直接計算した有限 foundation の表示である。
二形式のうち $f_1=uv$ は非退化である。
$f_2=(\alpha u+\beta v)(\gamma u-\alpha v)$ の二つの線形因子は
\[
 \det\begin{pmatrix}\alpha&\beta\\\gamma&-\alpha\end{pmatrix}
 =-(\alpha^2+\beta\gamma)=-1
\]
により一次独立であるため、$f_2$ も非退化である。
従って第~\ref{sec:relative-noncubic-families}節の逆構成によって
実際の $\mathcal H_2$ 対象を与える。
この比較は有限代数の積を保ち、元の構造頂点も保つ。
\end{proof}

\begin{proposition}[非被約閉基底上の両階数三正規形]
\label{prop:lagrangian-relative-double-rank-three}
$R$ を標数零の体上の可換環、$L\subset U_R\oplus U_R^*$ を
Lagrangian な階数四の部分束とする。
二射影の三次小行列式がそれぞれ単位イデアルを生成し、その四次行列式が
$R$ の元としてともに零であると仮定する。
さらに二つの交点線束の局所生成元は、二次正方行列として可逆とする。
このとき $\operatorname{Spec}R$ の \texorpdfstring{\'etale}{etale} 被覆上で、
前補題の式~\eqref{eq:lagrangian-rank-three-normal-form} が成立する。
従ってその有限 foundation は同じ基底上の明示中央二重平方族の
有限 foundation と同型である。基底の被約性は不要である。
\end{proposition}
\begin{proof}
三次小行列式を可逆化し、行基本変形と列基本変形を行うと、射影行列は
$\operatorname{diag}(I_3,a)$ の形となる。
四次行列式が実際に零なので $a=0$ であり、各 kernel、image、cokernel は
所定階数の局所自由部分束または商束となる。
特に $L\cap U_R$ と $L\cap U_R^*$ は線束である。
Zariski 局所にその可逆行列としての生成元を $X,Y$ と選ぶ。
等方性から $T=X^tY$ は trace 零であり、$\det T$ は可逆である。
Cayley--Hamilton により $T^2=-\det(T)I_2$ である。

$\lambda^2=-\det T$ を満たす可逆元 $\lambda$ を有限 \texorpdfstring{\'etale}{etale}
被覆により加える。その上で
\[
 e_\pm=\frac12(I_2\pm\lambda^{-1}T)
\]
は互いに直交する階数一の idempotent であり、$R^2$ を二つの固有線束へ分解する。
各固有線束の局所基底を選び、$hTh^{-1}=\lambda D_0$ とする。
$g=D_0h^{-t}X^{-1}$ と置くと、前補題の tensor 作用は
\[
 gXh^t=D_0,\qquad g^{-t}Yh^{-1}=\lambda I_2
\]
を与える。$Y$ の生成元を $\lambda^{-1}$ 倍すれば交点線は $RD_0$ と $RI_2$ となる。

両射影の像は反対の交点線の annihilator に一致する。
これらの束はそれぞれ
$RD_0\oplus\langle E_{12},E_{21}\rangle_R$ と
$RI_2\oplus\langle E_{12},E_{21}\rangle_R$ である。
$L$ から二交点線を割った階数二束は、両側の
$\langle E_{12},E_{21}\rangle_R$ へ同型に射影する。
従ってこれは可逆行列 $Q$ の graph であり、等方性により $Q=Q^t$ である。
次に $t^4=-\det Q$ の有限 \texorpdfstring{\'etale}{etale} 被覆へ移り、
$U_R,U_R^*$ の scalar を $t,t^{-1}$ 倍すると $\det Q=-1$ となる。
この被覆が \texorpdfstring{\'etale}{etale} であることには、$4t^3$ が可逆であることだけを使う。
$\alpha=-q_{12}$、$\beta=q_{11}$、$\gamma=-q_{22}$ と置けば、
$\alpha^2+\beta\gamma=1$ と式~\eqref{eq:lagrangian-rank-three-normal-form} が
$R$ 上の恒等式として成立する。
式~\ref{eq:nchh-plane} の比較もそのまま相対的に成立する。
\end{proof}

\paragraph{Cross reshuffling.}
$M\in\operatorname{Hom}(U,U^*)$ を四因子テンソルと見て、その係数を
$M_{ck,ab}$ と書く。cross reshuffling を
\[
 \operatorname{Cr}(M)_{ka,bc}=M_{ck,ab}
\]
で定める。これは中間の隣接二因子を入力とする縮約を表す四次正方行列である。
$M=y\otimes\nu$ が階数一で、$y$ と $\nu$ がそれぞれ二次正方行列として可逆なら、
$\operatorname{Cr}(M)$ は行・列の並替を除いて二つの可逆二次行列の Kronecker 積であり、可逆である。

\begin{proposition}[noncubic 点を含む全幾何点認識の開部分]
\label{prop:lagrangian-geometric-recognition-open}
$L_0$ を前補題の条件を満たす点とし、有限 foundation の Grassmann chart 内にあるとする。
$\operatorname{LG}(4,H)$ における $L_0$ の開近傍 $W$ を取れて、その全幾何点の
有限 foundation は $\mathcal H_2$ 対象の normal foundation となる。
具体的に、$L$ の局所基底を列行列
\[
 \begin{pmatrix}M\\N\end{pmatrix},\qquad M,N\in\operatorname{Mat}_4
\]
で表し、次の有限開条件を課せばよい。
\begin{enumerate}
\item $\mathbb P(L\cap U)\cap\Sigma_U=\varnothing$、
      $\mathbb P(L\cap U^*)\cap\Sigma_{U^*}=\varnothing$。
\item $\operatorname{Cr}(N\operatorname{adj}M)$ と
      $\operatorname{Cr}(M\operatorname{adj}N)$ が可逆である。
\end{enumerate}
二つの射影がともに階数三なら前補題の noncubic 型であり、
少なくとも一方の階数が四なら一方の ruling foundation が geometric cubic 型である。
この命題は幾何点での有限 foundation の認識であり、相対 heart の存在を結論に先取りしない。
\end{proposition}
\begin{proof}
(1) の失敗は、$L$ が固定された Segre 二次曲面上の点を含むという
射影 incidence の像である。従って (1) は開条件である。
$\mathbb P^3$ の任意の直線は代数閉体上で Segre 二次曲面と交わるので、
(1) により両方の交点の次元は高々一、両射影の階数は三以上である。

$L$ の基底を $T$ により変えると
\[
 N\operatorname{adj}M\longmapsto
 BT\operatorname{adj}(AT)=(\det T)N\operatorname{adj}M
\]
である。従って (2) も基底に依存しない開条件である。
指定点 $L_0$ では $M$ の階数が三なので、$\operatorname{adj}M$ の像は
$\ker M$、その行空間は $\operatorname{ann}(\operatorname{im}M)$ である。
$(M;N)$ の単射性により $N\ker M$ は非零で、$L_0\cap U^*$ を与える。
また Lagrangian 条件により
$\operatorname{ann}(\operatorname{im}M)=L_0\cap U^*$ である。
従って $N\operatorname{adj}M$ は、二次正方行列として非退化な二元の階数一外積である。
上の Kronecker 積の観察により、その cross reshuffling は可逆である。
$M\operatorname{adj}N$ についても同様である。
従って (1), (2) はともに $L_0$ で成立する。

$M$ が可逆な幾何点を取る。元の $F$ は同型であり、
その cubic quintuple の四因子テンソルは、二次元因子上の同型と定数倍を除いて
\[
 K=BA^{-1}
\]
である。Lagrangian 条件 $M^tB=N^tA$ により $K$ は対称である。
一組の隣接二因子の縮約が零となることは、$K$ の核が非零の分解テンソルを含むことに等しい。
この核は $L\cap U$ なので (1) に反する。
反対の隣接二因子については $K^t=K$ によって同じ結論を得る。
残る二組の隣接二因子は $\operatorname{Cr}(K)$ とその転置で表され、
\[
 \operatorname{Cr}(K)
 = (\det M)^{-1}\operatorname{Cr}(N\operatorname{adj}M)
\]
は (2) により可逆である。
従って四組全ての隣接縮約が非零であり、geometric 条件を満たす。
命題~\ref{prop:cubic-foundation-open} の cubic 逆構成により、
元の foundation は実際の $\mathcal H_2$ 対象から来る。

$N$ が可逆な場合は、命題~\ref{prop:relative-ruling-involution} の二短完全列で
ruling を交換する。その第一合成は $j^t$ に相当し、したがって可逆である。
新しい cubic tensor は二次元因子の同型と隣接対を保つ因子順の反転を除いて $AB^{-1}$ であり、
同じ議論と第二の cross 条件が geometric 性を与える。
この因子順は次のように直接確認できる。
元の四因子を $(V_0,V_1,V_2,K^*)$ とし、
$F:V_0\otimes V_1\to D$、$j:D^*\to V_2\otimes K^*$ と書く。
ruling 交換後の第一合成は $j^t$、新しい関係束の包含は $F$ の curry であり、
四因子は $(K,V_2^*,V_1^*,V_0^*)$ となる。
従って新しいテンソルの係数は
\[
 w'_{kcba}=(F^\vee j^{-1})_{ab,ck}.
\]
標準交代形式への基底変更を入れた射影行列は
$M=-F^\vee$、$N=D_Lj$ なので、$D_L^2=I_4$ を使うと
\[
 F^\vee j^{-1}=-AB^{-1}D_L.
\]
$D_L$ は一つの二次元因子だけの基底変更であり、
$(a,b,c,k)$ から $(k,c,b,a)$ への反転は隣接対の集合を保つ。
また $M^tB=N^tA$ から $AB^{-1}=(AB^{-1})^t$ が従うので、
両 summand を交換した前段の判定がそのまま適用できる。
ruling 交換は導来圏内の指定構造頂点を保つ mutation である。
認識した cubic $\mathcal H_2$ 対象の二つ目の normal foundation が元の foundation なので、
元の点も $\mathcal H_2$ に由来する。
両射影が階数三の場合は前補題を適用する。
\end{proof}

\begin{corollary}[cubic 二開部分と noncubic 閉部分]
\label{cor:lagrangian-open-closed-recognition}
前命題の $W$ において、二つの行列式を零とする閉部分スキームを
$D=V(\det M,\det N)$ と置く。
後の補題~\ref{lem:lagrangian-boundary-transverse} により、
$k$ 上のこの $D$ 自身は smooth（従って reduced）である。
ただし非被約な任意の基底へ引き戻した閉部分スキームは
非被約となり得る。以下の相対正規形では各行列式が
\emph{基底環の元として零}であることを要求し、
幾何 fiber で階数三となるだけの条件に置き換えない。
$D$ 上では命題~\ref{prop:lagrangian-relative-double-rank-three} の正規形が
\texorpdfstring{\'etale}{etale} 局所に成立し、対応する有限 foundation は
相対中央二重平方族のものに一致する。
開補集合は $D(\det M)$ と $D(\det N)$ で覆われ、
前者は元の ruling の geometric cubic 表示、後者は交換した ruling の
geometric cubic 表示を持つ。
\end{corollary}
\begin{proof}
$D$ 上では各三次小行列式が単位イデアルを生成し、
交点線の determinant も幾何点で零とならないので局所的に可逆である。
命題~\ref{prop:lagrangian-relative-double-rank-three} を適用する。
二つの cubic 開部分については前命題の geometric 判定が全幾何点で成立し、
geometric locus は射影 incidence の像の補集合として定義された開部分なので、
その判定は任意の基底変更と両立する。
\end{proof}

\begin{remark}[cross reshuffling の添字規約と行列式]
\label{rem:lg-cross-reshuffling-sign}
式 $\operatorname{Cr}(M)_{ka,bc}=M_{ck,ab}$ では
行添字を $(k,a)$、列添字を $(b,c)$ とし、
それぞれ辞書式に $(00,01,10,11)$ と並べる。
前の二射影とも階数三となる正規形
\eqref{eq:lagrangian-rank-three-normal-form} を
この規約で四行四列に書くと、係数の厳密な多項式計算により
\[
 \det\operatorname{Cr}(N\operatorname{adj}M)=-1,\qquad
 \det\operatorname{Cr}(M\operatorname{adj}N)
       =-(\alpha^2+\beta\gamma)^4=-1
\]
となる。
一般の rank-one テンソル $Y\otimes\nu$ では
$\det\operatorname{Cr}(Y\otimes\nu)=-(\det Y\,\det\nu)^2$ である。
列添字の $(b,c)$ と $(c,b)$ を交換した別の reshuffling を用いると
四列のうち $01$ と $10$ が入れ替わり、行列式の符号が反転する。
両方の規約で可逆性の条件は同じだが、
本文の添字規約と検算コードは上記の符号で一致させる。
これは有限型認識命題の開条件の非消滅を補強するものであり、
同命題の全相対認識をこの行列式計算だけで証明するものではない。
\end{remark}

\begin{example}[幾何 fiber が noncubic でも閉部分を経由しない一次変形]
\label{ex:lg-nilpotent-transverse-determinant}
$R=k[\varepsilon]/(\varepsilon^2)$ とし、標準交代形式を持つ
$U_R\oplus U_R^*$ 内の Lagrangian 四平面を、
その基底の二射影行列
\[
 M_\varepsilon=
 \begin{pmatrix}
 0&0&-1&-\varepsilon/2\\
 0&1&0&0\\
 -1&0&0&0\\
 0&0&1&-\varepsilon/2
 \end{pmatrix},\qquad
 N_\varepsilon=
 \begin{pmatrix}
 0&0&0&1\\
 1&0&0&0\\
 0&-1&0&0\\
 0&0&0&1
 \end{pmatrix}
\]
で定める。直接の行列積により
$M_\varepsilon^{t}N_\varepsilon=N_\varepsilon^{t}M_\varepsilon$
が成立する。剰余体への還元では両射影の階数は三、
各交点線は可逆二次正方行列で生成される。
従って上の認識開部分の $R$-点となる。
一方
\[
 \det M_\varepsilon=\varepsilon\ne0,\qquad
 \det N_\varepsilon=0.
\]
さらに本文の cross reshuffling 規約では
\[
 \det\operatorname{Cr}(N_\varepsilon\operatorname{adj}M_\varepsilon)
 =\varepsilon^2-1=-1,\qquad
 \det\operatorname{Cr}(M_\varepsilon\operatorname{adj}N_\varepsilon)=-1
\]
が $R$ 上で成立する。
唯一の幾何 fiber は noncubic であるが、この族は
閉部分 $V(\det M,\det N)$ を経由しない。
したがって命題~\ref{prop:lagrangian-relative-double-rank-three} の
\emph{scheme-theoretic} な行列式零条件は本質的である。
この検算は
\path{computations/verify_lagrangian_recognition.py} で再現できる。
\end{example}

\subsection{全幾何 fiber の有限 foundation 認識からの一様化}

\begin{lemma}[normal helix の幾何 fiber における BW 条件]
\label{lem:normal-foundation-BW-fiber}
$K$ を標数零の代数閉体、$(\mathcal C,\mathcal O)\in\HH(K)$
を強い構造対象付き二次曲面とし、
論文Aの定理~\cite[\texttt{thm:normal-helix}]{QuadricIntrinsic2026}
から得られる normal foundation を $E_i$ とする。
$D_{ij}=\operatorname{Hom}(E_{-j},E_{-i})$、
$P_i=e_iD$、$S_i$ を頂点単純加群とする。
このとき各 $S_i$ は頂点 $i,\ldots,i+4$ だけに支持される
長さ三の有限 projective 解消を持ち、その最終項は
$P_{i+4}$ 一個である。さらに
\[
 \operatorname{Ext}_{D}^{q}(S_i,P_j)=
 \begin{cases}
 K,&q=3,\ j=i+4,\\
 0,&\text{otherwise}.
 \end{cases}
 \tag{BW-normal}\label{eq:normal-foundation-BW-fiber}
\]
従って五頂点bar補題の (BW-fiber) は
normal foundation として認識された各幾何 fiber で成立する。
\end{lemma}
\begin{proof}
論文Aの系~
\cite[\texttt{cor:normal-cubic-defect}]{QuadricIntrinsic2026}
は、二つのMCM rulingの距離二の積の余次元 $d_i\in\{0,1\}$ を用いて、
右単純加群の極小解消を
\[
 0\longrightarrow P_{i+4}
 \longrightarrow P_{i+3}^{2}\oplus P_{i+2}^{d_{i+2}}
 \longrightarrow P_{i+1}^{2}\oplus P_{i+2}^{d_i}
 \longrightarrow P_i\longrightarrow S_i\longrightarrow0
\]
と記述する。ここで $d_{i+2}=d_i$ であり、
cubicの場合には $d_i=0$、noncubicの場合には
偶数始点で $d_i=1$、奇数始点で $d_i=0$ となる
（添字の始点の選択を固定した場合）。
従って支持・解消長・最終項の条件が得られる。
また論文Aの normal helix定理が与える
三次元・Gorenstein parameter四の左右 AS 正則
$\mathbb Z$-代数について、AS正則性の定義に含まれる単純加群の双対条件
\cite[Definition 4.2, (ASR1)--(ASR2)]{MN2025}
を適用すると、右 projective 列に対する
Ext は \eqref{eq:normal-foundation-BW-fiber} となる。
この引用はAS正則性からASF正則性への含意
（同論文Theorem 4.8）と区別する。
これは単純加群の解消とAS双対性を合わせた主張であり、
単にprojective resolutionの最終項が一個であることだけから
全Ext消滅を導いていない。
\end{proof}

\begin{theorem}[normal foundation の全 fiber 認識による全軌道の一様化]
\label{thm:uniform-orbit-from-normal-fibers}
$R$ を本稿の基礎体 $k$ 上の可換 Noetherian 代数、$\Lambda$ を有限 projective $R$-代数とし、
$\operatorname{Perf}\Lambda$ は相対 smooth proper とする。
$F_0,F_1,F_2,F_3$ を相対 full strong exceptional foundation とし、
$\operatorname{End}(\bigoplus_{a=0}^3F_a)\simeq\Lambda$ の同定を固定する。
実際の可逆 perfect Morita 双加群による自己同値 $\mathsf s$ と同型
\[
 \mathsf s^2\simeq S^{-1}[2],\qquad
 \mathsf sF_0\simeq F_2,\qquad \mathsf sF_1\simeq F_3
\]
を固定する。ここで $S$ は canonical な相対 Serre 関手である。
各幾何点 $x$ において、$\operatorname{Perf}\Lambda_x$ とある体上の
$\HH$ 対象 $(\mathcal C_x,\mathcal O_x)$ の導来圏との同値があり、
$(F_0,F_1,F_2,F_3)_x$ を定理~\cite[\texttt{thm:normal-helix}]{QuadricIntrinsic2026} の
normal helix の最初の四項に送ると仮定する。
この仮定は各 fiber の有限 foundation に関するものであり、
$R$ 上の heart や相対正規形を先に仮定しない。

\[
 E_{2a}=\mathsf s^aF_0,\qquad E_{2a+1}=\mathsf s^aF_1,
 \qquad C_{ij}=\operatorname{Hom}_{\Lambda}(E_{-j},E_{-i})\quad(i\le j)
\]
と置き、$i>j$ では $C_{ij}=0$ とする。この同じ $R$ 上で次が成立する。
\begin{enumerate}
\item 全 $i\le j$ に対して $\mathbf R\operatorname{Hom}(E_i,E_j)$ は
  次数零の有限 projective $R$-加群に集中し、その階数は
  \[
  h_{2n}=(n+1)^2,\qquad h_{2n+1}=(n+1)(n+2)
  \]
  である。全 Hom と合成は任意の導来基底変更と整合する。
\item 全軌道環
  \[
  B=\bigoplus_{n\ge0}\operatorname{Hom}(F_0,\mathsf s^nF_0)
  \]
  は一次生成・二次・相対 Koszul であり、全 $B_n$ は階数 $(n+1)^2$ の
  有限 projective 加群である。二次関係は階数七の直和因子となり、
  双対 $B^!$ の全次数、左右および双加群 Koszul 解消は任意基底変更と両立する。
\item $C$ の各頂点単純加群には、距離四以内の五頂点 bar 解消があり、
  これは任意基底変更後にも完全である。
  任意の可換 $R$-代数 $R'$ に対し、canonical な比較
  \[
  D(\operatorname{QGr}B_{R'})\simeq D(\operatorname{Mod}\Lambda_{R'}),
  \qquad D(\operatorname{QGr}B_{R'})^c\simeq\operatorname{Perf}\Lambda_{R'}
  \]
  があり、$\pi B$ を $F_0$ に、$(1)$ を最初に選んだ実際の $\mathsf s$ に送る。
  この比較とその Serre 平方同型は、非平坦なものを含む導来基底変更と整合する。
\end{enumerate}
さらに $B$ が左右 Noetherian なら、任意の有限型可換 $R$-代数 $R'$ 上で、
$\operatorname{qgr}B_{R'}$ は Noetherian heart、module-object 基底変更と fppf descent、
一般 coherent Hom の有限性と平坦基底変更、全 $\chi_q$、AZ 豊富性を持つ。
構造対象は平坦かつ exceptional であり、全幾何 fiber は仮定した $\HH$ の heart と
構造対象を保って exact 同値である。
同じ根の相対 AS4 被覆の存在は、この定理の結論とは別に確認する。
\end{theorem}

\begin{proof}
各 fiber の列を同定する。
$E_{i+4}=S^{-1}E_i[2]$ であり、最初の四項は指定された $F_0,\ldots,F_3$ である。
体上の normal helix も同じ四項と同じ反 Serre 四周期を持つから、
全整数 $i$ で両列が一致する。
ここで選び直した別の根との比較は必要ない。
最初に与えた $\mathsf s$ がこの列の二項移動を実現する。

$\Lambda$ の相対 smooth proper 性により、各
$K_{ij}=\mathbf R\operatorname{Hom}_{\Lambda}(E_i,E_j)$ は $R$ 上 perfect である。
normal helix の geometric 性から、全幾何 fiber の $K_{ij}$ は
$i\le j$ のとき階数 $h_{j-i}$ の次数零成分だけを持つ。
perfect 複体の fiber 判定
\cite[Tag 0BDK, Lemma 36.31.3]{StacksPerfectFibers}
によれば、その所定階数・所定次数への集中を表す開部分は全幾何点を含むので
$\operatorname{Spec}R$ 全体である。
従って各 $K_{ij}$ 自体が指定階数の vector bundle となる。
この操作は $i,j$ ごとに新たな開を選ぶものではない。
既に同じ基底の全 fiber に認識仮定があり、その全体上に判定を適用している。
有限 projective 性から、非平坦な基底変更でも通常の tensor が導来 tensor と一致する。
これで (1) を得る。

体上の normal helix では、$C$ は距離高々二の成分で生成され、
奇数頂点からの正距離の射は直後の偶数頂点を経由する。
対応する積の余核は各成分で有限生成であり、全幾何 fiber で零なので零である。
従ってこれらの生成条件は $R$ 上で成立する。
同様に、$B_1\otimes B_n\to B_{n+1}$ は全射となる。
二次関係 $J=\ker(B_1^{\otimes2}\to B_2)$ は、$B_2$ の projectivity により
階数七の直和因子である。
$T_R(B_1)/(J)\to B$ は各次数で有限生成加群から projective 加群への全射となる。
この全射は $R$ 上分裂する。
各 fiber の偶数 $\mathbb Z$-代数は既知の $\HH$ 切断環のものと同じであり、
二次で Koszul である。
実際の二周期の選択が異なっても、頂点単純加群の線形解消は同じ偶数
$\mathbb Z$-代数に属するので、この二つの性質は変わらない。
具体的には、周期を指定して得る次数付き環の次数付き加群圏は、各次数成分を
対応する頂点に置くことで、この偶数 $\mathbb Z$-代数の加群圏と同値になる。
生成次数と解消の線形性もこの対応で保たれる。
従って上の全射の kernel の全 fiber は零であり、有限生成性と Nakayama により零である。
これで $B$ の二次提示を得る。

ここまでに $B_1,B_2,B_3$ の有限 projective 性と、全幾何 fiber の Koszul 性を
互いに独立に確認した。
各局所環で Polishchuk--Positselski の局所 Koszul 変形定理
\cite[Chapter 6, Theorem 7.3]{PolishchukPositselski2005} を適用できる。
補題~\ref{lem:relative-koszul-degree-three} と同じ局所貼合せにより、
全双対成分の projectivity、相対 Koszul 性、任意底変更と両立する解消を得る。
一つの特殊点の Koszul 性を開全体へ外挿していないことに注意する。

全ての幾何 fiber を実際の $\HH$ の normal foundation に
同定するという定理の仮定に、補題~
\ref{lem:normal-foundation-BW-fiber} を適用すると、
五頂点barの支持・終項とAS双対の (BW-fiber) が一括して得られる。
この仮定は、有限foundationが単にstrong exceptionalであるという
条件だけからは従わない。
従って補題~\ref{lem:five-window-bar-lifting} の全幾何 fiber 版を
$R$ 上で適用できる。
全 $C_{ij}$ の projectivity は既に (1) で得ているため、この適用は循環しない。
極小 Betti 数が cubic と noncubic の間で変わっても、五頂点の非極小 bar 複体は共通である。
命題~\ref{prop:five-window-canonical-tails} により、全基底変更後の
canonical torsion 局所化、四項 tilting、compact 圏の比較を得る。

比較が指定の全軌道を保つことも確認する。
$P_i\mapsto E_{-i}$ とする cocontinuous derived functor を取る。
前向きの derived Hom が次数零に集中しているため、向き付けられた dg 部分圏の
truncation と $H^0$ の zigzag によりこの functor を構成できる。
五頂点 bar 解消の像は有限個の perfect $\Lambda$-対象からなる複体である。
各幾何 fiber では、normal helix の AZ 比較の下で単純 torsion 加群の解消の像となるので零である。
従ってその totalization は基底上の perfect fiber 判定で零となる。
二項移動と整合するため、二つの parity のこの比較だけで全 $i$ を扱える。
導来 torsion kernel は $\operatorname{loc}\{S_i\}$ なので、functor は tails へ降りる。
四項 tilting generator の End 代数を元の $\Lambda$ と同定するから、これは Morita 同値である。
奇数始点の生成条件から偶数頂点への restriction は
$\operatorname{QGr}C\simeq\operatorname{QGr}B$ を与える。
全 $P_i$ の像が指定された $E_{-i}$ であり、二項移動が指定された $\mathsf s$ なので、
次数移動および実際の平方同型も対応する。
全構成は有限 projective 係数の解消と導来 tensor から得られ、任意基底変更と整合する。

最後に $B$ が左右 Noetherian とする。
有限型中心基底変更は有限個の中心変数の付加と商であるため、全 $B_{R'}$ も左右 Noetherian である。
補題~\ref{lem:relative-cubic-coherent-cohomology} と
補題~\ref{lem:relative-cubic-global-chi} の証明は、Noetherian 生成系と
式~\ref{eq:five-window-tails-cohomology} の有限幅の cohomology 公式を用いるので、
この五頂点の場合にも同じ証明が適用できる。
特に coherent 対象間の平坦基底変更は、
四項tiltingによる $D^b(\operatorname{mod}\Lambda)$ 表示を経由し、
補題~\ref{lem:finite-algebra-flat-hom-basechange} から従う。
ここでは非平坦変更については導来圏の比較を行うのであって、
一般 coherent 対象間の通常の Hom の基底変更は要求しない。
module-object 基底変更は補題~\ref{lem:qgr-module-object-basechange} による。
各 fiber では元の normal helix が豊富であり、その偶数部分も元の $\HH$ の
切断列であるから、得られた heart と指定した $(\mathcal C_x,\mathcal O_x)$ は
構造対象を保って exact 同値となる。
\end{proof}

\begin{remark}[有限型認識で必要な入力の順序]
この定理では、有限 foundation の全幾何 fiber を既知の normal または cubic 型へ
認識する有限型開条件が、全次数の一様性を供給している。
各次数の平坦 locus を無限に交叉する必要はない。
ここでの大域的な結論は有限 projective 性であり、無限個の Hom 束を
一つの開上で同時に自由化するという、不要な強い主張は含めていない。
また $B$ の Koszul 性を使って全 forward Hom を証明してはいない。
全 fiber の geometric helix から Hom を先に平坦化し、その後に相対 Koszul 原理を使う。
Noetherian 性と同じ根の AS4 被覆は、この順序の後に別途供給すべき入力である。
\end{remark}

\subsection{閉部分と二つの cubic 開部分からの Noetherian 性}

\begin{lemma}[標準 noncubic 族の相対 normal tilting]
\label{lem:nc-relative-normal-tilting}
$R$ を可換 Noetherian $k$-代数とし、
$f_1=uv$、$f_2=(\alpha u+\beta v)(\gamma u-\alpha v)$、
$\alpha^2+\beta\gamma=1$ とする。
\[
 B_0=R[u,v]\langle c,d\rangle/(c^2-f_1,d^2-f_2),\qquad
 Z=R[u,v,h],\quad h=cd+dc
\]
と置く。式~\ref{eq:nchh-plane} の行列分解の cokernel を $M_c,M_d$ とすれば、
両者は $Z$-加群として
\[
 M_c\simeq Z^{\oplus2}\oplus Z(-1)^{\oplus2},\qquad
 M_d\simeq Z^{\oplus2}\oplus Z(-1)^{\oplus2}
\]
であり、これらの表示は任意の基底変更と両立する。
特に四項
\[
 T=(\pi B_0,\pi M_c,\pi B_0(1),\pi M_d(1))
\]
は相対 full strong exceptional foundation で、その End 代数は
式~\ref{eq:nchh-plane} の有限 foundation 代数に一致する。
対応する相対 Morita 同値は四対象、合成、相対 Serre 関手を保ち、
任意の Noetherian 基底変更と整合する。
\end{lemma}
\begin{proof}
$M_c$ の右生成元を行 $e=(e_1,e_2)$ と書く。
$e\Phi_c=0$ は
\[
 e c=e C_c,\qquad C_c=\begin{pmatrix}0&u\\v&0\end{pmatrix},
 \qquad C_c^2=f_1I_2
\]
に等しい。
$B_0$ は $Z$ 上の基底 $1,c,d,cd$ を持つので、この関係による reduction で
$e_1,e_2,e_1d,e_2d$ が $M_c$ を $Z$ 上生成する。
線形独立性を、作用を明示して確認する。
基底行 $\mathcal E=(e_1,e_2,e_1d,e_2d)$ に対する係数を列に置くと、
右乗法 $c,d$ はそれぞれ次の行列で表される。
\[
 C=\begin{pmatrix}C_c&hI_2\\0&-C_c\end{pmatrix},\qquad
 D=\begin{pmatrix}0&f_2I_2\\I_2&0\end{pmatrix}.
\]
すなわち $(\mathcal E z)c=\mathcal E Cz$、
$(\mathcal E z)d=\mathcal E Dz$ である。
$C^2=f_1I_4$、$D^2=f_2I_4$、$CD+DC=hI_4$ だから、
これらは $Z^{\oplus2}\oplus Z(-1)^{\oplus2}$ 上の右 $B_0$-作用を定める。
最初の二基底は $e\Phi_c=0$ を満たすので、この自由加群への
$M_c$ からの全射を得る。
他方、前記の reduction は逆向きの $Z$-線形全射を与える。
両合成は指定四基底上で恒等なので、求める同型である。
$M_d$ も $c,d$ を交換し、
$C_c$ を $\left(\begin{smallmatrix}0&\alpha u+\beta v\\
\gamma u-\alpha v&0\end{smallmatrix}\right)$ に替えれば同じである。
この証明は係数環上の恒等式なので、任意基底変更と両立する。

命題~\ref{prop:nc-qgr-weighted} の weighted stack
$X=\mathbb P_R(1,1,2)$ と Clifford order $\mathcal B_0$ を用いる。
上の二加群は中心 $Z$ 上有限自由なので、対応する sheaf は $R$ 上平坦である。
全幾何 fiber の $\pi M_c,\pi M_d$ は既知の noncubic $\HH$ 対象の
normal foundation に属し、canonical compact 対象である。
中心の affine \texorpdfstring{\'etale}{etale} chart 上で
補題~\ref{lem:nc-fiber-perfect} を適用すると、両 sheaf は $\mathcal B_0$ 上局所 perfect となる。
命題~\ref{prop:nc-compact-generator} によって、それらは global にも compact である。
この compactness と平坦性は、任意の Noetherian 基底変更に保たれる。

命題~\ref{prop:nc-relative-smooth} の相対 smooth proper 性により、
四項の間の全 derived Hom は $R$ 上 perfect で、導来基底変更と両立する。
各幾何 fiber で四項は normal foundation なので、perfect fiber 判定により、
その全 Hom は所定階数の次数零 vector bundle、逆向きの Hom は零となる。
行列分解からの射は、式~\ref{eq:nchh-plane} の有限代数から
$\Gamma=\operatorname{End}(\bigoplus T_i)$ への代数準同型を与える。
各幾何 fiber で同型であり、双方は所定階数の有限 projective 加群なので、
これは $R$ 上でも同型である。
特に $\Gamma$ は有限 foundation の normalized bar 解消により相対 smooth である。

full 性を別に確認する。
$G=\bigoplus_{j=0}^3\mathcal B_0(j)$ を既知の compact generator とし、
$T_\Sigma=\bigoplus T_i$ と置く。
相対 Morita 随伴
\[
 \mathsf F=(-)\otimes_\Gamma^{\mathbf L}T_\Sigma,
 \qquad \mathsf G=\mathbf R\operatorname{Hom}(T_\Sigma,-)
\]
を取る。$\mathbf R\operatorname{End}(T_\Sigma)=\Gamma[0]$ より
$\mathsf F$ は $\operatorname{Perf}\Gamma$ 上 fully faithful である。
$\mathsf G(G)$ は $R$ 上 perfect であり、$\Gamma$ の相対 smooth 性から
$\Gamma$ 上も perfect である。
従って counit の cone
\[
 K=\operatorname{Cone}(\mathsf F\mathsf G(G)\longrightarrow G)
\]
は compact である。各幾何 fiber では $T_\Sigma$ が tilting generator なので $K$ は零となる。
$\mathbf R\operatorname{Hom}(G,K)$ は $R$ 上 perfect で全幾何 fiber が零だから零であり、
$G$ の生成性から $K=0$ である。
従って $G$ は $T_\Sigma$ の thick 閉包に属し、$T_\Sigma$ も compact generator となる。
これで相対 Morita 同値を得る。
全構成は perfect 対象の評価と tensor から成り、基底変更と整合する。
Serre 関手との比較は相対双対 functor の Yoneda 一意性による。
\end{proof}

\paragraph{四周期からの Hom $\mathbb Z$-代数の比較。}
標識付き四項 foundation の同型から得た相対 Morita 同値を $F$、
両側の反 Serre 関手を $t,t_0$ と書く。
上の比較は自然同型 $\eta:Ft\simeq t_0F$ を与える。
最初の四項の同型 $\varphi_a:F(E_a)\simeq E_a^0$ ($0\le a<4$) と、
選んだ実際の平方同型による四周期を用い、
$i=a+4n$ に対して $\eta$ の反復から
$\varphi_i:F(E_i)\simeq E_i^0$ を定める。負の $n$ には逆同型を使う。
\[
 f\longmapsto\varphi_jF(f)\varphi_i^{-1}
 \qquad(f:E_i\longrightarrow E_j)
\]
は合成を保つ全 Hom $\mathbb Z$-代数の同型であり、任意基底変更と両立する。
この比較に、二つの根そのものを同一視する仮定は要らない。
異なる $\varphi_i$ の選択は頂点ごとの共役を与えるだけなので、
偶数制限の row と column の Noetherian 性、および線型解消の存在は変わらない。



\begin{lemma}[中央 Clifford 閉部分の普遍 Noetherian 性と忠実平坦降下]
\label{lem:closed-clifford-fpqc-noetherian}
$R_0$ を可換 Noetherian 環、$f_1,f_2\in R_0[u,v]_2$ とし、
\[
 B_0=R_0[u,v]\langle c,d\rangle/(c^2-f_1,d^2-f_2)
\]
と置く。任意の可換 Noetherian $R_0$-代数 $T$ について
$B_0\otimes_{R_0}T$ は左右 Noetherian である。
さらに $D$ を正の向きの $R_0$-線形 $\mathbb Z$-代数とし、
全 $D_{ij}$ は有限 projective で任意基底変更と可換すると仮定する。
有限個の affine \texorpdfstring{\'etale}{etale} 射
$\operatorname{Spec}R_\nu\to\operatorname{Spec}R_0$ が被覆し、
その上で $D$ の全 indexed Hom と合成が
上の型の標準中央 Clifford 族の indexed Hom と同型になるなら、
任意の可換 Noetherian $R_0$-代数 $T$ 上で
$D_T$ の全 row と column は Noetherian となる。
\end{lemma}
\begin{proof}
補題~\ref{lem:noncubic-relative-clifford-center} の
任意基底上の PBW 表示により
\[
 B_0\otimes_{R_0}T
 \cong T[u,v,h]\oplus T[u,v,h]c
                  \oplus T[u,v,h]d\oplus T[u,v,h]cd.
\]
$T$ が Noetherian なので中心多項式環 $T[u,v,h]$ は
Noetherian であり、その有限加群である上の非可換環も
左右 Noetherian である。
標準族の indexed Hom $\mathbb Z$-代数はその次数付き環の
自由シフトの Hom を記述するので、各 row と column は
対応する次数付き自由加群と同じ Noetherian 性を持つ。
これは任意の Noetherian $T$ でも成立する。

$R_*=\prod_\nu R_\nu$ と置く。有限個の flat \'{e}tale
射の和が全射なので $R_0\to R_*$ は忠実平坦である。
任意の Noetherian $R_0$-代数 $T$ についても
$T\to T\otimes_{R_0}R_*$ は忠実平坦となり、
後者は有限型 \'{e}tale $T$-代数の有限直積として
Noetherian である。
指定した indexed Hom の同型は任意の $T$ への基底変更後も
全成分と合成について保たれる。
従って各 row と column はこの忠実平坦な被覆の上で
Noetherian である。

右 row $P_i$ の任意の部分加群 $N\subset P_i$ に対し、
$N\otimes_T(T\otimes_{R_0}R_*)$ は有限個の頂点同次元で
生成される。それらの生成元をテンソルの有限和で表し、
現れた有限個の $N$ の頂点同次元が生成する部分加群を
$N_0$ とする。忠実平坦性から
$(N/N_0)\otimes_T(T\otimes_{R_0}R_*)=0$ は $N=N_0$
を意味する。従って $P_i$ は Noetherian。
左 column にも同じ議論を適用する。
この方法では被約な幾何点集合での同型だけでは足りず、
\emph{閉部分の環 $R_0$ 上の実際の相対同型}を入力とする。
\end{proof}

\begin{lemma}[二つの cubic 開部分と noncubic 閉部分からの普遍 Noetherian 性]
\label{lem:lg-two-rulings-universal-noetherian}
$R$ を可換 Noetherian 環とし、有限 LG 認識開上の foundation と
実際の Morita 根 $s$ を取る。四項を $E_0,E_1,E_2,E_3$ と書き、
指定された同型
\[
 sE_0\simeq E_2,\qquad sE_1\simeq E_3,
 \qquad s^2\simeq S^{-1}[2]
\]
を用いて $E_{2n}=s^nE_0$, $E_{2n+1}=s^nE_1$ と置く。
全ての前向き導来 Hom は有限 projective 加群に次数零で集中し、
任意の底変更と両立することを既に確認したとする。
このとき同じ根による偶数軌道環
\[
 C_n=\operatorname{Hom}(E_0,s^nE_0),\qquad
 a*b=s^{\deg b}(a)\circ b
\]
は、任意の可換 Noetherian $R$-代数 $T$ に対して
$C_T=C\otimes_RT$ が左右 Noetherian となる。

ここで有限 LG 認識開について用いる入力は次の三つである。
二つの射影行列を $M,N$ として
\[
 a=\det M,\quad b=\det N,\quad I=(a,b),
 \qquad Z=\operatorname{Spec}(R/I)
\]
と置く。
\begin{enumerate}
\item $D(a)$ では選んだ ruling の helix が geometric cubic 型である。
\item $D(b)$ では反対の ruling の helix が geometric cubic 型である。
\item $Z$ 上では、非被約構造を含めて、foundation が \'etale 局所的に
  判別式の可逆な二形式による標準 noncubic foundation と同型である。
\end{enumerate}
これらは追加の全次数条件ではなく、有限 LG 認識で確認する条件である。
\end{lemma}

\begin{proof}
まず二つ目の cubic 開部分で、偶数軌道環が変わらないことを明示する。
$U=\operatorname{Hom}(E_2,E_3)$ とし、相対評価射から perfect 対象
\[
 Y=\operatorname{Cone}\bigl(U\otimes_RE_2\longrightarrow E_3\bigr)[-1]
\]
を定める。各幾何 fiber の正規 helix では
式~\cite[\texttt{eq:normal-ruling-evaluation}]{QuadricIntrinsic2026} が
\[
 0\longrightarrow Y\longrightarrow E_2^{\oplus2}
 \longrightarrow E_3\longrightarrow0
\]
を与え、$Y$ は $E_1$ と反対の初期 ruling である。
実際、同式の記号で $E_3=\pi X_1(1)$ と書けば
核は $\pi p(X_1)$ であり、$X_0=pc(X_1)$ より
$p(X_1)=c(X_0)$ となる。
従って反対側の列は
\[
 E'_{2n}=E_{2n},\qquad E'_{2n+1}=s^nY
\]
と取れる。評価射と cone の形成は導来底変更と両立する。
この列の前向き Hom も、全幾何 fiber での正規 helix の消滅から
同じ基底上で有限 projective に次数零で集中する。

特に ruling の交換は偶数対象を全く変えない。
偶数対象間の Hom、その合成、および指定された $s$ による二周期同型も
変わらない。従って $D(b)$ で反対側の cubic helix から得る偶数環は
同じ $C_b$ である。単に二つの環の fiber を比較しているのではない。

各 cubic 開部分では、全 Hom の projectivity と底変更整合性が既にある。
隣接生成性、階数二の cubic relation、可逆な四次 syzygy、およびその縮約同型は、
各固定次数における有限 projective 加群の写像として定義できる。
幾何 fiber で geometric cubic の提示と標準解消になるので、
それらの余核と解消のホモロジーは零である。
例えば関係部分加群は、既に projective と分かっている Hom 成分への
全射の核として得るため、projectivity と底変更をここで循環して仮定しない。
これによって定理~\ref{thm:relative-geometric-cubic-noetherian} の
相対入力が二つの開部分で全て成立する。
系~\ref{cor:relative-periodic-even-noetherian} により
$C_a$ と $C_b$ は、任意の Noetherian 基底変更後にも左右 Noetherian である。

次に閉部分 $Z$ を扱う。
仮定 (3) の \'etale 局所標準化は、三次 minor の可逆性と
$\det M=\det N=0$ の実際の等式から得られる。
従ってこれは被約点集合上の標準化ではなく、$R/I$ 上の foundation の同型である。
標準 noncubic 族の係数環を $R'$ と書く。
補題~\ref{lem:noncubic-relative-clifford-center} は、
冪零元を含む任意の係数環上でも
$B_0=R'[u,v]\langle c,d\rangle/(c^2-f_1,d^2-f_2)$ が
中心 $R'[u,v,h]$ 上階数四の自由代数となることを保証する。
ここでは $h=cd+dc$ であり、
monic PBW の四つの重なりの検証は係数の被約性に依存しない。

標準族の根と選んだ $s$ を同一視する必要もない。
補題~\ref{lem:nc-relative-normal-tilting} の相対 tilting が
最初の四対象・合成・相対 Serre 関手を同定する。
両者の四周期 $E_{i+4}\simeq S^{-1}E_i[2]$ を使えば
全 indexed Hom $\mathbb Z$-代数が頂点ごとの共役で一致し、
偶数部分に制限しても同じである。
そこで補題~\ref{lem:closed-clifford-fpqc-noetherian} を、
この \'{e}tale 局所的な \emph{相対} 同型に適用する。
結果として $C/IC$ は左右 Noetherian であり、
任意の Noetherian $R/I$-代数への基底変更後も左右 Noetherian となる。
忠実平坦降下の対象は各 row と column であり、
点ごとの Noetherian 性をそのまま降下したものではない。

最後に定理~\ref{thm:patch-noetherian} を、環 $C$ と中心イデアル
$I=(a,b)$ に適用する。
各 $C_n$ は有限 $R$-加群であり、$C/IC,C_a,C_b$ の左右 Noetherian 性を
確認したので、$C$ も左右 Noetherian である。
この段階は、各 stalk の Noetherian 性から全体へ推論するものではない。
実際には $I$-進完成と二つの開部分の有限忠実平坦被覆による貼合せである。

任意の可換 Noetherian $R$-代数 $T$ に対しては、$I_T=(a,b)T$ と置く。
商と局所化は底変更と可換するので
\[
 C_T/I_TC_T\cong(C/IC)\otimes_{R/I}(T/I_T),
 \qquad (C_T)_a=C_a\otimes_{R_a}T_a,
 \qquad (C_T)_b=C_b\otimes_{R_b}T_b
\]
である。右辺は上で証明した普遍性により左右 Noetherian であり、
$(C_T)_n$ は有限 $T$-加群である。
同じ貼合せ定理を $T,I_T,C_T$ に適用して結論を得る。
\end{proof}

\subsection{同じ根に対する中央 AS4 被覆の延長}

\begin{lemma}[局所中央被覆の有限型近傍への延長]
\label{lem:spread-local-central-hull}
$R$ を可換 Noetherian 環、$\mathfrak p\in\operatorname{Spec}R$ とする。
$B=T_R(V)/(J)$ は連結二次 $R$-代数であり、
$V,J$ はそれぞれ階数四、七の有限 projective 加群、
$J\subset V^{\otimes2}$ は直和因子とする。
$B$ は相対 Koszul であり、全ての $B_n$ は階数 $(n+1)^2$ の
有限 projective 加群として任意の底変更と両立すると仮定する。
$E=B^!$ と置く。

局所環 $R_{\mathfrak p}$ 上で、中央正則元
$\phi_{\mathfrak p}\in(E_{\mathfrak p})_2$ があり、
$P_{\mathfrak p}=E_{\mathfrak p}/(\phi_{\mathfrak p})$ は
階数列 $(1,4,6,4,1)$ の有限自由 Frobenius 代数であるとする。
このとき $\mathfrak p\in D(f)$ となる一つの主開部分へ制限すれば、
$\phi_{\mathfrak p}$ は中央元 $\phi\in(E_f)_2$ に延長し、次が成り立つ。
\begin{enumerate}
\item $\phi$ は全次数で正則であり、この性質は任意の底変更に保たれる。
\item $P=E_f/(\phi)$ は階数列 $(1,4,6,4,1)$ の相対 Koszul
  Frobenius 代数であり、$A=P^!$ の全次数は有限 projective である。
\item $A$ は相対 quadratic Koszul AS4 型の代数であり、
  ある中央正則二次元 $g$ によって $A/(g)\cong B_f$ と書ける。
  この表示、正則性、左右および双加群 Koszul 解消は任意の底変更と両立する。
\end{enumerate}
さらに任意の可換 Noetherian $R_f$-代数 $T$ に対して $B_T$ が左右 Noetherian
ならば、$A_T$ も左右 Noetherian である。
この主張には、$D(f)$ の全ての noncubic fiber において指定した根を
別の ``良い根'' と比較する仮定は要らない。
\end{lemma}

\begin{proof}
相対 Koszul 双対性により、$E$ の各次数は有限 projective であり、
任意の底変更と両立する。また全ての幾何 fiber $E_{\bar s}$ は Koszul である。
まず $\phi_{\mathfrak p}$ を有限生成加群 $E_2$ の局所化の元として表し、
一つの主開部分上の元 $\phi$ へ延長する。
一次生成加群 $E_1$ の有限個の生成元との交換子は $E_3$ の元であり、
$R_{\mathfrak p}$ 上で零である。従ってさらに一つの主開部分へ制限して
これら全てを零とできる。$E$ は一次生成なので $\phi$ は中央である。

$P=E/(\phi)$ と置く。各 $P_i$ は有限表示 $R$-加群である。
$0\leq i\leq4$ に対する指定階数での局所自由性と $P_5=0$ は、
有限個の条件として $\mathfrak p$ の一つの開近傍へ広がる。
同様に、$P_4$ を自明化したときの Frobenius pairings
$P_i\otimes P_{4-i}\to P_4$ の非退化性も、有限個の行列式を
可逆化することによって同じ近傍で成立させられる。
$P$ は一次生成なので $P_5=0$ は $P_n=0$ ($n\geq5$) を含意する。
補題~\ref{lem:spread-central-regular-line} を適用すると、さらに一度縮めて
$\phi$ の全次数正則性と任意底変更後の正則性が得られる。

ここで Koszul 性を新たに開条件として使う必要はない。
各幾何 fiber において $E_{\bar s}$ は既に Koszul であり、
$\phi_{\bar s}$ は中央正則二次元である。
正則正規二次元による Koszul 代数の商は Koszul だから、
全ての $P_{\bar s}$ が Koszul となる。
$P_1,P_2,P_3$ の局所自由性は既に確認したので、
相対 Koszul 変形原理
（補題~\ref{lem:relative-koszul-degree-three}）を適用できる。
従って $P$ は相対 Koszul であり、$A=P^!$ の全次数も有限 projective、
その階数は $\binom{n+3}{3}$ である。
Frobenius pairings と相対 Koszul 双対性が相対 AS4 型を与える。

最後に $\phi$ を $J=E_2^\vee$ 上の全射としてみて $H=\ker\phi$ と置く。
局所的に $J/H$ を生成する $g$ を選べば、$A=T(V)/(H)$、$A/(g)=B$ となる。
ここでは $\phi$ を一つの大域的元として延長したので $J/H\cong R$ であり、
$g$ も大域的に選べる。
$K_3=(J\otimes V)\cap(V\otimes J)=E_3^\vee$ 上の二つの縮約
$\phi\otimes1,1\otimes\phi:K_3\to V$ は、$\phi$ の中央性により等しい。
$\phi E_1\to E_3$ は直和因子への単射だから、この共通写像は全射である。
従って任意の $v\in V$ に対して $A_3$ 内で $gv=vg$ となり、$g$ は中央である。
短完全列の候補
\[
 A_{n-2}\xrightarrow{\,g\,}A_n\longrightarrow B_n\longrightarrow0
\]
で、$B_n$ が projective なので右側の核は projective である。
$h_B=(1-t^2)h_A$ よりその階数は $A_{n-2}$ の階数に等しく、
この核への全射は同階数の有限 projective 加群間の同型となる。
従って左側も単射であり、列は split exact な $R$-加群列である。
任意底変更後にも $g$ は正則となる。
Koszul 解消の底変更整合性は相対 Koszul 双対性に含まれる。
Noetherian 性の最後の主張は、中央正則元による次数有界な $g$-進
フィルトレーションから $\operatorname{gr}_gA_T\cong B_T[g]$ を取り、
左右の有限生成性を持ち上げれば従う。
\end{proof}

\begin{remark}[一般の Koszul 開性との区別]
上の証明は、指定点の Frobenius 代数 $P_{\mathfrak p}$ が Koszul だから
近傍全体でも Koszul になる、という推論を使っていない。
先に全幾何 fiber の helix を H2 モデルとして認識し、
その同じ even Hom $\mathbb Z$-代数の線形解消から $B_{\bar s}$ の
Koszul 性を得る。この全 fiber の入力と、中央元の一様正則性を経由して
$P_{\bar s}$ の Koszul 性を導く点が重要である。
一般の二次代数では、三次成分の次元を固定したパラメータ空間における
$n$-Koszul locus は開であり、Koszul locus はこれらの可算共通部分である
\cite[Chapter 6, Theorem 2.1 and Corollary 2.2]{PolishchukPositselski2005}。
また同書 Chapter 6, \S6 は、Koszul locus が constructible でない族を与える。
これらを本研究の特別な中央延長に対する否定結果と解するべきではないが、
局所 Koszul 変形原理だけから単一の有限型開近傍での Koszul 性を
推論できるわけではない。本証明は、全幾何 fiber の明示的認識という
別の有限型入力を使って、この問題を回避する。
\end{remark}

\subsection{noncubic 点の有限型近傍での全相対認識}

\begin{theorem}[noncubic 境界を横断する有限型効果性]
\label{thm:noncubic-finite-type-effectivity}
任意の noncubic $\HH$ 対象の normal foundation に対応する点 $x$ と、
定理~\ref{thm:noncubic-hh-trivial-root} の良い根を取る。
有限 foundation の smooth affine chart $(S,x)$ の有限型 affine
\texorpdfstring{\'etale}{etale} 近傍 $(U,u)\to(S,x)$ 上に、
構造対象だけを指定した実際の相対族
\[
 q_U:U\longrightarrow\mathfrak H_2
\]
が存在する。この族の canonical compact 圏は普遍 foundation 代数
$\Lambda_U$ の perfect 圏であり、四項窓は普遍 projective に対応する。

同じ $U$ 上で、選んだ実際の Morita 根 $\mathsf s$ は heart の exact 自己同値となり、
全軌道環 $B$ は二次・相対 Koszul・全次数有限 projective である。
任意の Noetherian 基底変更後に $B$ は左右 Noetherian であり、
同じ根に対する中央正則二次被覆
\[
 A/(g)\simeq B,
 \qquad \rank A_n=\binom{n+3}{3}
\]
を持つ。$A$ は相対 quadratic Koszul AS4 型で、同じ普遍 Noetherian 性を持つ。
全 Hom、canonical tails の比較、相対 Serre 平方同型、被覆の正則性と
左右・双加群 Koszul 解消は任意の導来基底変更と整合する。
根・根同型・被覆はこの族を構成する補助選択であり、$q_U$ の標識には含めない。
\end{theorem}

\begin{proof}
最初に有限型の近傍を選ぶ。
命題~\ref{prop:noncubic-lagrangian-etale-slice} の representable な
\texorpdfstring{\'etale}{etale} 射を $S$ へ引き戻す。
$G^+$ は smooth なので、その torsor は指定幾何点の
\texorpdfstring{\'etale}{etale} 近傍で切断を持つ。
従ってさらに affine \texorpdfstring{\'etale}{etale} 近傍へ移り、
実際の Lagrangian 四平面とその射影行列を持つものとできる。
命題~\ref{prop:lagrangian-geometric-recognition-open} の有限開条件を課す。
その条件は指定点で成立するので、この制限は $x$ を失わない。
系~\ref{cor:noncubic-etale-morita-root} の近傍との共通
\texorpdfstring{\'etale}{etale} 細分を取り、実際の Morita 根、平方同型、
式~\ref{eq:noncubic-etale-root-window} の二対象比較を同時に固定する。
必要なら $u$ の affine 開へ縮める。

この同じ近傍の全幾何 fiber の四項 foundation は、
命題~\ref{prop:lagrangian-geometric-recognition-open} により
既知の $\HH$ の normal foundation である。
定理~\ref{thm:uniform-orbit-from-normal-fibers} を適用して、
全 forward Hom の有限 projective 性、全軌道の二次提示と相対 Koszul 性、
五頂点 bar 解消と canonical tails 比較を得る。
この段階では全次数の平坦開を交叉していない。
全幾何 fiber に同じ有限次元の認識を先に与えたことを用いている。

系~\ref{cor:lagrangian-open-closed-recognition} は
$V(\det M,\det N)$ の非被約構造を含む相対正規形と、二つの cubic 開部分を与える。
補題~\ref{lem:lg-two-rulings-universal-noetherian} により、
全軌道環の任意の Noetherian 基底変更後の左右 Noetherian 性が従う。
従って定理~\ref{thm:uniform-orbit-from-normal-fibers} の最後の主張も適用できる。
一般 coherent Hom の有限性と平坦基底変更、module-object 基底変更、
fppf descent、AZ 豊富性と全 $\chi_q$ を同じ近傍上で得る。
各 fiber の heart は、元の normal helix の AZ 再構成によって
既知の $\HH$ の heart と構造対象を保って exact 同値である。
特に Gabriel--Krull の rank-one torsion-free 条件も満たす。
これは導来同値だけから heart を認識する推論ではない。

$S$ は smooth で $U\to S$ は \texorpdfstring{\'etale}{etale} なので、
$R=\mathcal O_{U,u}$ は正則局所環である。
指定点の既存の局所中央被覆定理
\ref{thm:noncubic-local-central-hull} を適用すると、
$E_R=B_R^!$ の中央正則二次元 $\phi_R$ と
階数列 $(1,4,6,4,1)$ の Frobenius 商を得る。
ここで使う全軌道と根は既に同じものとして比較済みである。
また geometric cubic locus は、既約な有限 foundation 表示空間の稠密開である
（命題~\ref{prop:cubic-foundation-open}）。
$S$ はその非空開部分で、$U\to S$ は \texorpdfstring{\'etale}{etale} だから、
$U$ の各成分は $S$ に dominant であり、その一般点はこの cubic 開部分に入る。
従って同定理の generic cubic 条件も満たされる。
補題~\ref{lem:spread-local-central-hull} により、一つの開近傍へ縮めて
中央 AS4 被覆を全次数・任意底変更について延長できる。
$B$ の普遍 Noetherian 性から被覆 $A$ のそれも従う。

最後に $\operatorname{qgr}B$ と $\pi B$ を取る。
次数移動は exact であり、canonical 導来比較の下で最初に選んだ
$\mathsf s$ と実際の平方同型に対応する。
四項 tilting は相対 smooth proper な $\Lambda_U$ との Morita 同値を与え、
canonical compact、相対 Serre、非平坦を含む導来基底変更の条件を満たす。
上で確認した heart、fiber、AZ、全軌道、同じ根の AS4 被覆と合わせ、
定義~\ref{def:relative-h2} の全公理が成立する。
従って $q_U$ は、幾何点の表示だけでなく実際の相対 $\mathfrak H_2$ 族である。
\end{proof}

% 相対効果性の補助結果と普遍 cubic 表示の入力は分類論文
% \cite{QuadricIntrinsic2026} の対応する節から引用する。
% !TeX root = ../nc-quadric-surface-moduli.tex
% 2026-11-22 / v1.44（pointed Artin smooth bridgeとnoncubic chart次元を明確化）
\section{内在的な全モジュライの有限型代数性}
\label{sec:global-h2-smooth-atlas}

本節では、全ての noncubic 点で得た有限型の効果的近傍を用いる。
入力は実際の相対 $\mathfrak H_2$ 族であり、全幾何 fiber の分類だけではない。
これらの近傍から指定された pointed 同型を保つ smooth atlas を作り、
既存の noncubic 表示族の準コンパクト性を使って有限個の chart に縮約する。
有限 foundation スタックと内在的モジュライの位相を先に同一視する必要はない。

\subsection{noncubic 点を含む実際の smooth chart}

\begin{corollary}[効果的 noncubic 近傍の smooth atlas]
\label{cor:noncubic-effective-smooth-neighborhood}
$x\in\mathfrak H_2(k)$ を noncubic 点とする。
有限 foundation の標準表示空間 $S_0=\mathbb A^{64}$ の開部分 $S$ と、
$x$ の normal foundation を表す点 $s_x\in S(k)$ を取る。
定理~\ref{thm:noncubic-finite-type-effectivity} で構成した有限型 affine étale 近傍
\[
 a_x:(U_x,u_x)\longrightarrow(S,s_x)
\]
および実際の族 $q_x:U_x\to\mathfrak H_2$ を取る。
全軌道と canonical tails の比較により、この族には次の性質を持つ列 $(E_i)$ がある。
各 $E_i$ は平坦な noetherian compact 対象で、$E_0=\OO$ であり、
各四項窓は full strong exceptional collection、
全前向き derived Hom は次数零の有限 projective 加群である。
さらに $E_{i+4}=S^{-1}E_i[2]$ であり、最初の四項窓の End 代数は
頂点・積・任意基底変更を保って $a_x^*\Lambda$ と同定される。

このとき $q_x$ はスキームで表現可能な smooth 射である。
その fppf 像 $\mathfrak V_x\subset\mathfrak H_2$ は $x$ を含む開部分スタックであり、
有限型で滑らかな Artin スタックである。
\end{corollary}

\begin{proof}
定理~\ref{thm:atlas-from-effective-foundation} の仮定を確認する。
tails の row は次数成分ごとに基底上平坦であり、商関手は tensor と整合するので、
各 $E_i$ は平坦である。Noetherian 性と compact 性は実際の相対族の構成による。
四項の完全性・strong 性と前向き Hom の消滅は、五頂点 bar 比較により任意基底変更上で成立する。
full exceptional collection の Serre rotation は mutation による helix の延長と一致するため、
$E_{i+4}=S^{-1}E_i[2]$ は De Deken--Lowen の geometric $(4,3)$-helix 条件を与える。
最後の End 代数の同定も全軌道比較の一部である。
従って同定理を適用できる。

指定された同型に関する点を明記する。
small extension $A'\twoheadrightarrow A$ と
\[
 u_A\in U_x(A),\qquad c_{A'}\in\mathfrak H_2(A'),\qquad
 \alpha_A:q_x(u_A)\xrightarrow{\sim}c_{A'}|_A
\]
を取ると、$\alpha_A$ により四項窓とその全 Hom の基底を還元先へ移す。
De Deken--Lowen の持上げを基底 $A$ に直接適用し、
例外対象の $\Ext^1,\Ext^2$ の消滅と $A'^\times\to A^\times$ の全射性により、
$E'_0=\OO_{c_{A'}}$ とその指定された pointing を保つ。
五つの Hom 空間の基底も指定還元を延長して選ぶと、実際の合成から $S(A')$ の点が得られる。
$a_x$ の全点での étale 性により、これは $u_A$ を保つ $u_{A'}\in U_x(A')$ へ持ち上がる。

両方の四項 tilting 窓の End 代数をこの基底で同定する。
得られる標準 dg Morita 同値の還元は、対象だけでなく全ての射上で $D(\alpha_A)$ と一致する。
補題~\ref{lem:artin-two-heart-comparison}、または反 Serre 再構成の全忠実性を使えば、
これは $\alpha_A$ 自体を延長する pointed exact 同型となる。
module-object 基底変更により、その同型は全有限型 test 上の相対系にも延長する。
この議論は根や被覆を同型の標識へ加えず、指定した pointed 同型をそのまま保つ。

これらの指定された Artin 同型の持上げと、
$\Phi:\mathfrak H_2\to\mathfrak P$ の全忠実性は
補題~\ref{lem:pointed-artin-lifting-smooth-bridge} の全仮定を満たす。
従って $q_x$ はスキームで表現可能な smooth 射で、
その開像は全ての pointed 同型を relation に含めた
smooth groupoid により表される。
この推論は有限型スキーム上の Tag 02HY の Artin 試験に帰着し、
現時点で $\mathfrak H_2$ 全体が Artin スタックであるとは仮定しない。
\end{proof}

\subsection{有限個の chart による全体の表示}

\begin{theorem}[全 $\mathfrak H_2$ の有限型代数性と滑らかさ]
\label{thm:global-h2-smooth-finite-type}
局所有限型 $k$-スキームの site 上の内在的スタック $\mathfrak H_2$ は、
$k$ 上有限型で滑らかな純次元三の Artin スタックの制限により表される。
その対角は affine かつ有限表示である。
有限型の滑らかなスキーム $U$ と representable smooth 全射
\[
 q:U\longrightarrow\mathfrak H_2
\]
を取れ、全ての pointed exact 同値を用いる relation scheme
\[
 R=U\times_{\mathfrak H_2}U
   \simeq U\times_{\mathfrak P}U
\]
に対して $\mathfrak H_2\simeq[U/R]$ である。
ここで選ばれた foundation、Serre 根、平方同型、AS4 被覆を保存することは、射の条件に加えない。
\end{theorem}

\begin{proof}
まず有限個の noncubic 近傍で足りることを示す。
定理~\ref{thm:noncubic-relative-effectivity} による実際の射を、
表示スタックの atlas へ引き戻すと
\[
 p:P=Q^{\mathrm{nd}}\times Q^{\mathrm{nd}}
   \longrightarrow\mathfrak N^{\mathrm{pres}}
   \longrightarrow\mathfrak H_2
\]
を得る。$P$ は有限型の準コンパクトな $k$-スキームである。
命題~\ref{prop:noncubic-presentation-stack} により、全 noncubic $k$-対象はその像に入る。
各 noncubic 点 $x$ に前系の開部分 $\mathfrak V_x$ を取る。
引戻し $p^{-1}(\mathfrak V_x)$ は $P$ の開部分スキームであり、これらは $P$ の全閉点を覆う。
有限型 $k$-スキームは Jacobson なので、全体も覆う。
従って準コンパクト性により有限個
\[
 p^{-1}(\mathfrak V_{x_1}),\ldots,p^{-1}(\mathfrak V_{x_m})
\]
で $P$ を覆える。この段階で使ったのは $P\to\mathfrak H_2$ の存在と
各 $\mathfrak V_x$ の内在的な開性だけであり、$\mathfrak N^{\mathrm{pres}}\to\mathfrak H_2$
の full faithfulness や、任意の相対 noncubic 族の正規形は使わない。

定理~\ref{thm:cubic-smooth-open} の有限型 smooth atlas $U_{\mathrm{cub}}$ と合わせ、
\[
 U=U_{\mathrm{cub}}\amalg U_{x_1}\amalg\cdots\amalg U_{x_m}
\]
と置く。有限和なので $U$ は有限型で滑らかである。
各成分の射は representable smooth なので $q:U\to\mathfrak H_2$ も同様であり、
その開像は全ての $k$-対象を含む。

全ての相対族を覆うことも確認する。
有限型 affine $T\to\mathfrak H_2$ に対して、$U\times_{\mathfrak H_2}T\to T$ の像は開である。
$T$ の各閉点の剰余体は $k$ であり、その fiber は上の分類と有限被覆によって像に入る。
従ってこの開像は全閉点を含み、Jacobson 性により $T$ 全体である。
ゆえにこの射は smooth 全射であり、fppf 局所に section を持つ。
任意の局所有限型 $T$ については有限型 affine 開被覆へ制限すればよい。
これで $q$ が元の site 上の全対象を fppf に覆うことが示された。

相対反 Serre 認識の全忠実性から、relation は表示した $U\times_{\mathfrak P}U$ と等しい。
$\mathfrak P$ の対角が affine 有限表示なので、$R$ は有限型スキームである。
両射影 $R\rightrightarrows U$ は $q$ の基底変更として smooth であり、
\cite[Tags 04X0, 04TK]{StacksSmoothAtlas} により $[U/R]$ は通常の全スキーム上の Artin スタックとなる。
その元の site への制限は $\mathfrak H_2$ と同値である。
有限型 atlas $U$ によって有限型性が従い、$U$ の滑らかさによって $k$ 上の滑らかさが従う。
対角の性質は命題~\ref{prop:relative-h2-algebraicity-reduction} による。

最後に、系~\ref{cor:cubic-formal-h2-abelian} と系~\ref{cor:noncubic-all-artin-h2} により、
各 $k$-点の形式変形 groupoid は pointed Abel 変形と同値である。
その接空間は $\operatorname{HH}^2(\Lambda)$、無限小自己同型空間は
$\operatorname{HH}^1(\Lambda)$ であり、
命題~\ref{prop:foundation-hochschild} によって両次元の差は三である。
滑らかな Artin スタックの局所次元はこの差に等しいから、全閉点で三となる。
有限型 smooth atlas の各成分は閉点を持ち、次元は局所一定なので、全体は純次元三である。
\end{proof}

\begin{corollary}[noncubic 有効 chart の相対次元]
\label{cor:noncubic-atlas-relative-dimension}
上の $q_x:U_x\to\mathfrak H_2$ は、像上で相対次元
$64-3=61$ の smooth 射である。
その relation scheme
$R_x=U_x\times_{\mathfrak H_2}U_x$ は
$k$ 上 smooth な純次元 $64+61=125$ のスキームである。
\end{corollary}
\begin{proof}
$U_x$ は $S\subset\mathbb A^{64}$ の有限型 étale 近傍なので、
その各点における次元は64である。
定理~\ref{thm:global-h2-smooth-finite-type} により
$\mathfrak H_2$ は純次元3でsmoothである。
smoothなスタック射の相対次元公式より $q_x$ は相対次元61。
$R_x\to U_x$ は $q_x$ の基底変更であるため、
相対次元61の smooth 射で、総次元は125となる。
これはsmooth atlasのfiber次元の整合性を表すものであり、
非cubic点のモジュライに新たな変形方向を加えたものではない。
\end{proof}

\subsection{普遍 cubic 族の降下と内在的な既約性}

\begin{lemma}[根を標識にしない普遍非線型 cubic 族]
\label{lem:universal-nonlinear-cubic-h2-family}
$S_{\mathrm e}=\mathscr G\setminus\mathscr G_{\mathrm{lin}}$ を、
第~\ref{sec:relative-picard-roots}節の非線型 geometric 五つ組の開部分とする。
普遍 $\mathbb Z$-代数 $D$ と $C_{ij}=D_{-j,-i}$ から、実際の相対族
\[
 q_{\mathrm e}:S_{\mathrm e}\longrightarrow\mathfrak V_{\mathrm{cub}}
\]
が得られる。この射は representable smooth であり、その開像は既約である。
開像の補集合の幾何点は、存在するとすれば可換二次曲面の pointed 圏 $Q$ だけである。
\end{lemma}

\begin{proof}
系~\ref{cor:finite-type-root-chart} の étale 被覆
$\coprod_\alpha T_\alpha\to S_{\mathrm e}$ を取る。
各 $T_\alpha$ には実際の二周期があり、定理~\ref{thm:geometric-cubic-relative-effectivity} により
$q_\alpha:T_\alpha\to\mathfrak H_2$ が構成されている。
その構成における偶数 restriction は、任意の基底変更上で
\[
 \operatorname{QGr}B_\alpha\simeq\operatorname{QGr}C_{T_\alpha}
\]
という構造対象を保つ exact 同値を与える。
右辺の $C$ は根を選ぶ前の普遍代数であり、二周期の選択に依存しない。
従って $T_\alpha\times_{S_{\mathrm e}}T_\beta$ 上では、両族は同じ
$\operatorname{QGr}C$ を通じて pointed exact 同値となる。
三重の重なりの cocycle は、この同じ圏を通る比較の結合性による。
coherent 対象と全有限型 module-object 基底変更にも整合するので、
$\mathfrak H_2$ の fppf descent により $q_{\mathrm e}$ が得られる。
根と AS4 被覆は局所的 witness であり、それら自体を重なりで同定する必要はない。

各 $q_\alpha$ は系~\ref{cor:quintuple-effective-smooth-chart} により smooth である。
smoothness は始域の étale 被覆で判定できるから $q_{\mathrm e}$ も smooth である。
$S_{\mathrm e}$ は射影空間の非空開部分で既約なので、その連続像は既約であり、
smooth 性によって像は開である。
全非線型 cubic 表示が像に入ることは、定理~\ref{thm:cubic-smooth-open} の逆構成による。
線型 geometric cubic 表示は全て同じ pointed $Q$ を与えるので、最後の主張が従う。
\end{proof}

\begin{proposition}[全内在的モジュライの既約性]
\label{prop:global-h2-irreducible}
$\mathfrak V_{\mathrm{cub}}$ は既約であり、$\mathfrak H_2$ において稠密である。
従って $\mathfrak H_2$ は既約かつ連結である。
\end{proposition}

\begin{proof}
前補題の既約な開像を $\mathfrak V_{\mathrm e}$ と書く。
残り得る点 $Q$ をその閉包に入れる。
系~\ref{cor:quadric-effective-smooth-neighborhood} の $U_Q$ で、指定点 $u_Q$ を含む
既約成分 $U_Q^0$ に制限する。
$U_Q$ は smooth なので、その既約成分は開である。
$U_Q^0\to S\subset S_0=\mathbb A^{64}$ は étale であり、像は非空開部分である。

線型 geometric cubic 表示の四項 foundation は、可換二次曲面の
\[
 (\OO,\OO(1,0),\OO(1,1),\OO(2,1))
\]
および ruling を交換した列から得られる。
従ってその $S_0$ 上の基底付き表示は、有限個の $\mathcal G/K$ 軌道に含まれる。
$\dim(\mathcal G/K)=61<64$ なので、この有限和の閉包は $S_0$ の真の閉部分である。
さらに命題~\ref{prop:cubic-foundation-open} の非空な geometric cubic 開部分と交差する。
この二つの稠密開条件を同時に課すと、$U_Q^0$ の稠密開部分の foundation は非線型 cubic であり、
その $q_Q$ による像は $\mathfrak V_{\mathrm e}$ に入る。
実際の射 $q_Q$ の連続性から、$Q$ は $\mathfrak V_{\mathrm e}$ の閉包に入る。
前補題の像が最初から $Q$ を含む場合にも同じ結論である。
従って $\mathfrak V_{\mathrm{cub}}$ は、この一つの既約な開部分の閉包であり、既約である。

任意の noncubic $k$-点 $x$ に対し、
定理~\ref{thm:noncubic-finite-type-effectivity} の $U_x$ のうち指定点を含む既約成分を取る。
これも $S_0$ の非空開部分への étale 射を持つ。
命題~\ref{prop:cubic-foundation-open} の geometric cubic foundation は
$S_0$ の非空開部分なので、$U_x$ の稠密開部分は cubic である。
その実際の族 $q_x$ は全幾何 fiber の foundation と heart を保つため、
この稠密開部分は $\mathfrak V_{\mathrm{cub}}$ へ送られる。
従って $x$ は $\mathfrak V_{\mathrm{cub}}$ の閉包に入る。

以上で全 $k$-点がこの閉包に入った。
もし補集合が非空なら、有限型 smooth atlas への逆像は非空開部分であり、
$k$-点を持つ。この像も $k$-点なので矛盾する。
従って $\mathfrak V_{\mathrm{cub}}$ は $\mathfrak H_2$ において稠密である。
既約部分空間の閉包は既約なので結論を得る。
この証明で用いた連続写像は全て実際の相対 $\mathfrak H_2$ 族の射であり、
finite foundation の幾何点像の位相を内在位相と同一視する操作はない。
\end{proof}

\subsection{noncubic 境界の滑らかさと余次元}

\begin{lemma}[二つの射影行列式の横断性]
\label{lem:lagrangian-boundary-transverse}
$W\subset\operatorname{LG}(4,U\oplus U^*)$ を
命題~\ref{prop:lagrangian-geometric-recognition-open} の認識開部分とし、
tautological 四平面束の二射影を
\[
 M:\mathscr L\longrightarrow U\otimes\OO_W,
 \qquad N:\mathscr L\longrightarrow U^*\otimes\OO_W
\]
と書く。
\[
 D=V(\det M,\det N)\subset W
\]
は $k$ 上 smooth な純次元八のスキームであり、$D\hookrightarrow W$ は
余次元二の正則閉埋込みである。
$D$ 上の交点線束を
\[
 \ell=\mathscr L\cap U=\ker N,\qquad
 m=\mathscr L\cap U^*=\ker M
\]
とすると、その法束は
\[
 N_{D/W}\simeq(m^*)^{\otimes2}\oplus(\ell^*)^{\otimes2}
 \tag{NC-boundary-normal}\label{eq:noncubic-boundary-normal}
\]
と同定される。
\end{lemma}

\begin{proof}
認識開条件により、$D$ の全幾何点で $M,N$ はともに階数三である。
三次小行列式を可逆化して行列を掃き出すと、射影は $\operatorname{diag}(I_3,a)$ と書ける。
$D$ 上ではその四次行列式が実際に零なので $a=0$ である。
従って非被約構造を先に除かなくても kernel と cokernel は線束であり、
$\ell,m$ は $\mathscr L$ の線部分束となる。
二 summand の交わりが零なので $\ell\oplus m\subset\mathscr L$ は階数二の部分束である。

幾何点 $L\in D$ において、標準交代形式から
\[
 T_L\operatorname{LG}\simeq\operatorname{Sym}^2L^*
\]
を得る。
具体的には一次変形 $\varphi:L\to(U\oplus U^*)/L$ に対し、
持上げ $\widetilde\varphi$ を用いた
$q(v,w)=\omega(v,\widetilde\varphi(w))$ が対称形式となる。
$q$ は持上げの選択に依存しない。

$\operatorname{im}M=m^\perp\subset U$ なので、
$\operatorname{coker}M\simeq m^*$ である。
行列式を微分すると、階数三の射 $M$ に対するその一次変分は
$\ker M=m$ から $\operatorname{coker}M=m^*$ への誘導射で計算される。
局所的には、これは $\operatorname{diag}(I_3,0)$ の右下成分の微分にほかならない。
従って $\det M$ の微分は、符号を含む単元倍を除いて
\[
 \operatorname{Sym}^2L^*\longrightarrow(m^*)^{\otimes2},
 \qquad q\longmapsto q|_{m\times m}
\]
である。同様に $\det N$ の微分は $q|_{\ell\times\ell}$ である。
二線は独立なので、この二つの restriction を合わせた射は全射である。
例えば $\ell,m$ の生成元を $L$ の基底の最初の二元とすれば、
二つの写像は対称行列の二つの独立な対角係数を取る。

行列式線束を局所自明化すれば、$D$ は滑らかな十次元スキーム $W$ 上の
二つの関数で定義され、その微分は全幾何点で独立である。
Jacobian 判定により $D$ は smooth で余次元二、両関数は正則列となる。
特にこの証明は $D$ の reducedness も同時に与える。

最後に行列式線束を intrinsic に読むと、kernel と cokernel の完全列から
\[
 (\det U\otimes\det\mathscr L^*)|_D
 \simeq\det(\operatorname{coker}M)\otimes\det(\ker M)^*
 \simeq(m^*)^{\otimes2}
\]
であり、$N$ 側は $(\ell^*)^{\otimes2}$ である。
上の微分の全射とその kernel の接束への同定から
\eqref{eq:noncubic-boundary-normal} を得る。
\end{proof}

\begin{theorem}[内在的 noncubic 境界は滑らかな余次元二の閉部分]
\label{thm:noncubic-boundary-smooth-codim-two}
定理~\ref{thm:global-h2-smooth-finite-type} の有限型 Artin スタックにおいて、
\[
 \mathfrak N=
 \bigl(\mathfrak H_2\setminus\mathfrak V_{\mathrm{cub}}\bigr)_{\mathrm{red}}
\]
を reduced induced closed substack とする。
このとき $\mathfrak N$ は有限型で滑らかな既約 Artin スタックであり、純次元一である。
閉埋込み $\mathfrak N\hookrightarrow\mathfrak H_2$ は余次元二の正則閉埋込みである。
具体的に、各 noncubic 点を含む適当な smooth scheme chart では、
そのイデアルは実際の二つの射影行列式 $(\det M,\det N)$ で生成される。
\end{theorem}

\begin{proof}
\emph{境界の定義。}
$\mathfrak V_{\mathrm{cub}}$ は内在的な開部分であり、$\mathfrak H_2$ は有限型なので、
その閉補集合には一意な reduced induced closed substack がある
\cite[Tag 0509, Lemma 100.10.1]{StacksReducedBoundary}。
これを $\mathfrak N$ とする。

\emph{quotient chart での余次元二。}
指定した noncubic 点に対し、良い根に対応する Lagrangian 表示を取る。
命題~\ref{prop:noncubic-lagrangian-etale-slice} の étale locus と認識開条件に制限し、
$G^+$-不変開 $W$ と
\[
 \mathfrak L=[W/G^+]\longrightarrow\mathfrak F
\]
を得る。ここで $\mathfrak F$ は該当する finite foundation スタックの開部分であり、
この射は representable étale である。
$G^+$ は二つの summand を保つ tensor 作用と反傾作用で作用するので、
射影の行列式は局所的に単元倍される。
従ってイデアル $(\det M,\det N)$ と $D$ は $G^+$-不変である。
前補題と smooth atlas $W\to[W/G^+]$ により
\[
 \mathfrak D=[D/G^+]\hookrightarrow\mathfrak L
\]
は余次元二の正則閉埋込みであり、$\mathfrak D$ は smooth である。
その次元は $8-\dim G^+=8-7=1$ である。
ここで正則性の降下は Koszul-regular immersion の fpqc locality を用い、
局所 Noetherian な表示上で通常の regular immersion と同値であることを使う
\cite[Tags 0694, 063L]{StacksReducedBoundary}。

\emph{両方の smooth 性を保つ共通 chart。}
定理~\ref{thm:noncubic-finite-type-effectivity} と
系~\ref{cor:noncubic-effective-smooth-neighborhood} の
$q_0:U_0\to\mathfrak H_2$ を取る。
$U_0\to S$ は étale、$S\to\mathfrak F$ は foundation の smooth atlas なので、
$U_0\to\mathfrak F$ は smooth である。
そこで
\[
 V=U_0\times_{\mathfrak F}\mathfrak L
\]
を取り、指定点を含む étale scheme chart $U\to V$ に移る。
このとき $V\to U_0$ は étale、$V\to\mathfrak L$ は smooth であるから、
\[
 q:U\longrightarrow\mathfrak H_2,
 \qquad a:U\longrightarrow\mathfrak L
\]
はともに smooth である。
必要ならさらに étale に細分して $G^+$-torsor を自明化し、
tautological 束の基底を選べば、実際の行列 $M,N$ を $U$ 上で取れる。
この細分後も二つの smooth 性は保たれる。

\[
 D_U=U\times_{\mathfrak L}\mathfrak D
\]
と置く。これは $U$ 上で $V(\det M,\det N)$ と書ける閉部分スキームである。
$a$ は smooth なので、前段の正則閉埋込みを flat に引き戻している。
従って $D_U\hookrightarrow U$ は余次元二の正則閉埋込みとなり、
$D_U\to\mathfrak D$ が smooth であることから $D_U$ 自身も $k$ 上 smooth である。
この議論は、torsor の切断から得る lift $U\to W$ の smooth 性や flat 性を仮定しない。

\emph{内在的境界との scheme-theoretic な同定。}
一方、$\mathfrak N_U=U\times_{\mathfrak H_2}\mathfrak N$ は $U$ の閉部分スキームである。
射 $\mathfrak N_U\to\mathfrak N$ は $q$ の基底変更として smooth であり、
$\mathfrak N$ は reduced なので $\mathfrak N_U$ も reduced である。
ここでは reducedness が smooth topology に局所的であることを使う
\cite[Tag 034E]{StacksReducedBoundary}。

認識開条件の下では、両射影が階数三の幾何点は noncubic normal foundation を与え、
一方が階数四なら一方の ruling foundation が geometric cubic となる。
実際の相対族はその fiber heart と構造対象を保つため、
系~\ref{cor:lagrangian-open-closed-recognition} と体上の分類により
\[
 |\mathfrak N_U|=|D_U|
\]
である。
両者は同じスキーム $U$ の reduced closed subscheme なので、
reduced induced 構造の一意性から
\[
 \mathfrak N_U=D_U
\]
が closed subscheme として成立する。
幾何点の一致から非被約な二つの部分スキームを無条件に同一視しているのではなく、
前補題と smooth pullback によって双方の reducedness を確認した後の同定である。

全 noncubic 点でこの共通 chart を取れるので、これらの $D_U$ は $\mathfrak N$ を smooth に覆う。
従って $\mathfrak N$ は $k$ 上 smooth であり、
$\mathfrak N\hookrightarrow\mathfrak H_2$ は smooth 局所で余次元二の正則閉埋込みである。
$\mathfrak H_2$ は純次元三なので $\mathfrak N$ は純次元一となる。

\emph{既約性。}
定理~\ref{thm:noncubic-relative-effectivity} の実際の族を
\[
 P=Q^{\mathrm{nd}}\times Q^{\mathrm{nd}}\longrightarrow\mathfrak H_2
\]
へ引き戻す。$P$ は $\mathbb A^6$ の非空開部分であり、reduced かつ既約である。
その全幾何 fiber は noncubic なので、この射は $\mathfrak N$ を経由する。
この因数分解には、reduced な始域から閉部分への set-theoretic な包含が
scheme-theoretic な因数分解を与えることを使う
\cite[Tag 050B, Lemma 100.10.3]{StacksReducedBoundary}。
体上の noncubic 分類を各代数閉拡大体上で適用すると、$P\to\mathfrak N$ は全幾何点を覆う。
従って $|\mathfrak N|$ は既約空間 $|P|$ の連続像であり、既約である。
この結論に、$\mathfrak N^{\mathrm{pres}}\to\mathfrak N$ の相対全忠実性や同値は用いない。
\end{proof}

\begin{remark}[非被約 test 上で境界に属する条件]
$\mathfrak N$ は reduced induced closed substack である。
非被約な test $T\to\mathfrak H_2$ が $\mathfrak N$ を経由することは、
全幾何 fiber が noncubic であることだけからは従わない。
上の smooth chart では、前者は $\det M=\det N=0$ が
$\OO_T$ の元として実際に成立することを要求する。
行列式が非零の冪零元となる横断的な Artin 変形は、閉 fiber が noncubic でも境界には属さない。
また本定理は reduced boundary の正則埋込みを証明するものであり、
それ以外の候補 functor に非被約構造を無条件に与える主張ではない。
\end{remark}

\subsection{framed 実現空間と ambient 内の非開性}

\begin{corollary}[framed 実現空間のスキーム性]
\label{cor:global-xas-smooth-scheme}
\[
 X_{\mathrm{as}}=X_3\times_{\mathfrak P}\mathfrak H_2
\]
は有限型で滑らかな既約・純次元八十三のスキームにより表される。
$X_{\mathrm{as}}\to X_3$ は有限型の monomorphism で、separated かつ unramified、quasi-finite である。
\end{corollary}

\begin{proof}
全スキーム上では前定理の $[U/R]$ を用いる。
$R=U\times_{\mathfrak P}U$ は全ての同型を記録するため、
$[U/R]\to\mathfrak P$ は全スキーム上でも全忠実である。
実際、対象を $U$ へ fppf 局所に持ち上げると Isom はまさに $R$ であり、
この一致は降下する。従って $X_{\mathrm{as}}\to X_3$ は monomorphism である。

命題~\ref{prop:relative-h2-algebraicity-reduction} により、$X_{\mathrm{as}}$ は有限型代数空間である。
また $X_{\mathrm{as}}\to\mathfrak H_2$ は $\operatorname{PGL}_9$-torsor なので smooth で相対次元八十である。
従って $X_{\mathrm{as}}$ は $k$ 上 smooth で純次元 $3+80=83$ である。
命題~\ref{prop:global-h2-irreducible} により $\mathfrak H_2$ は既約で、
この torsor は開全射かつ幾何 fiber が既約な $\operatorname{PGL}_9$ である。
従って $X_{\mathrm{as}}$ も既約である。
実際、二つの非空開部分の像は底の既約性により交わり、
その fiber 内の二つの非空開部分は fiber の既約性により交わる。

$X_{\mathrm{as}}$ と $X_3$ はともに有限型 $k$-空間なので、両者の間の射は有限型である。
局所有限型 monomorphism は separated、unramified、locally quasi-finite かつスキームで表現可能である
\cite[Tag 05W6, Lemma 76.14.8]{StacksRepresentableMonomorphism}。
特に $X_3$ がスキームであるため $X_{\mathrm{as}}$ 自身もスキームである。
同じ結論は separated locally quasi-finite な代数空間からスキームへの射に関する
\cite[Tag 03XX, Proposition 67.50.2]{StacksRepresentableMonomorphism} からも従う。
有限型性は準コンパクト性を含むので、射は quasi-finite である。
\end{proof}

\begin{proposition}[canonical 実現像は ambient 全体で開ではない]
\label{prop:global-h2-not-ambient-open}
$\Phi:\mathfrak H_2\to\mathfrak P=[X_3/\operatorname{PGL}_9]$ は、
noncubic の二つの根対が一致する点で open immersion ではない。
従って $X_{\mathrm{as}}\to X_3$ も、その点の任意の frame の近傍で open immersion ではない。
\end{proposition}

\begin{proof}
一致点 $x$ において、系~\ref{cor:noncubic-all-artin-h2} は
\[
 \dim T_x\mathfrak H_2=4,\qquad
 \dim\operatorname{HH}^1(\Lambda)=1
\]
を与える。ここで $T_x$ は閉 fiber の同定を保つ一次変形の同型類の空間である。
一方、系~\ref{cor:artin-polarization-retraction} の分裂完全列は
\[
 0\longrightarrow\operatorname{HH}^1(\Lambda)
 \longrightarrow T_x\mathfrak P
 \longrightarrow\operatorname{HH}^2(\Lambda)\longrightarrow0
\]
を与え、$\dim T_x\mathfrak P=5$ である。
$\Phi$ は canonical 反 Serre 環を取る射なので、その微分の像は
この分裂における $\operatorname{HH}^2$ の四次元の像である。
従って一次変形の射は全射ではなく、$\Phi$ はこの点で open immersion にならない。
開 immersion なら形式変形 groupoid を同一にすることに反するからである。

framed 空間については、その任意の frame $J$ で
\[
 \dim T_JX_{\mathrm{as}}=83,\qquad
 \dim T_JX_3=5+(80-1)=84
\]
となる。後者の $80-1$ は $\operatorname{PGL}_9$ の軌道接空間の次元である。
open immersion ならこの二つの接空間は同型になるので、framed 射についても結論を得る。
\end{proof}

\begin{remark}[表現可能性から自動的には出ない比較]
上で構成したのは内在的な族のスタックそのものの有限型 smooth atlas と、
framed 部分層のスキームによる表現である。
これによって $X_{\mathrm{as}}$ を $X_3$ の開部分と同一視することはできず、
実際に前命題が一致根対でその同一視を否定する。
また有限 foundation の ruling 商 $\mathfrak Q_{\mathrm{fnd}}$ との相対同値や、
その表示からの内在位相の比較は、この atlas の構成だけでは結論しない。
全スキーム上の $[U/R]$ は元の有限型 site 上の $\mathfrak H_2$ を延長するが、
任意の Noetherian 基底上で素朴に同じ Abel 公理を課した関手との一致は別の連続性の問題である。
\end{remark}

\subsection{偏極を変えた関係式の開近傍}
\label{subsec:twisted-polarization-neighborhood}

$V=k^9$、$G=\operatorname{PGL}(V)$ と置く。以下の $J$ は
$X_{\mathrm{as}}$ の点と、それが $V^{\otimes2}$ 内に定める関係部分束を表す。
系~\ref{cor:global-xas-smooth-scheme} の monomorphism により、後者を固定すれば
前者の選択は高々一つである。有限型のスキーム
\[
 D=\{(J,g)\in X_{\mathrm{as}}\times G:
                      (g\otimes g)J=J\}
\]
を取る。等式は Grassmannian の対角の逆像であり、スキームとして課す。
射影行列を線型行列に持ち上げても、スカラーは部分束を変えない。
$D$ は $X_{\mathrm{as}}$ 上で pointed 自己同値を表す群スキームである
（命題~\ref{prop:relative-h2-stack-recognition}）。特に平坦性は要求しない。

\begin{theorem}[次数移動を記録する開埋込み]
\label{thm:polarization-inertia-open}
次の $G$-同変射は開埋込みである。
\[
 \theta:D\longrightarrow X_3,\qquad
 (J,g)\longmapsto (g^{-1}\otimes 1)J .
\]
始域で $G$ は関係式の枠を変え、$g$ には共役で作用する。
像を $W$ とすると $[W/G]$ は $\HHstack$ の慣性スタックに同型である。
$X_{\mathrm{as}}\to X_3$ は $W$ の中の恒等切断 $g=1$ に一致し、
有限表示の局所閉埋込みである。
\end{theorem}
\begin{proof}
\emph{写像と逆構成の一意性。}
枠と次数の線束を局所自明化する。$R=T(V)/(J)$ の次数付き自己同型として
$g$ を持ち上げると、積 $a*b=g^{\deg b}(a)b$ の右 Zhang twist は同じ
次数成分を持ち、関係部分束は $(g^{-1}\otimes1)J$ である。
実際、長さ $n$ の語の積は順に $g^{n-1},\ldots,g,1$ を掛けてから元の積を
取るので、二次関係の全挿入もこの同型で対応する。
スカラーの取り方によらず $X_3$ への射が得られる。
右加群の明示同値は命題~\ref{prop:antiserre-polarization-family} の
式~\eqref{eq:polarization-right-zhang} の証明と同じで、任意の係数変更と可換する
\cite[Theorem 3.1]{Zhang1996}。

圏の言葉では、$(\cC,\OO)$ の pointed 自己同値 $u$ が反 Serre 環に
誘導する射影行列が $g$ である。$t=S^{-1}[2]$ と $u$ の可換同型は
Serre 双対の自然性から得られる。$u\OO\simeq\OO$ を使うと、
$F=ut$ の軌道対象 $F^n\OO$ は $t^n\OO$ と同定される。
右規約の積 $F^{\deg b}(a)\circ b$ をこの同定で移したものが上記の積である。
負の軌道対象も一致し、従って AZ の二条件は $t$ から $F$ へ移る。
全軌道環の tails は、構造対象と表示された移動を含めて
$(\cC,\OO,F)$ を再現する。
逆にこの triple から $t$ を内在的に回収すると $u=Ft^{-1}$ が決まる。
$u\OO\simeq\OO$ の選択はスカラーだけ異なり、同じ pointed 射の同型類を与える。
また、移動を比較する自然同型の違いは
補題~\ref{lem:relative-h2-flat-center} によりスカラーで、
次数 $n$ にはその $n$ 乗だけ作用する。従って、もとの次数一の
\emph{射影枠}も回収できる。具体的には
$\Hom(\OO,F\OO)\simeq\Hom(\OO,t\OO)$ に枠を移す。
線束 twist のある局所自明化の重なりでもこの射影同定は変わらない。
以上で同じ framed twist から二つの異なる $(J,g)$ は得られない。
特に任意の代数閉拡大体上で単射である。

\emph{非恒等 twist の Artin 近傍。}
$k$ 点 $(J_0,g_0)$ において、剰余体 $k$ の局所 Artin 環 $A$ 上の
任意の framed 関係式変形 $J'_A$ を取り、
$B_A=T(V_A)/(J'_A)$、$F_A=(1)$ と書く。
閉 fiber $B_0$ は Koszul である。実際、命題~\ref{prop:all-h2-antiserre-koszul}
の環の単純加群の線型自由解消を twist すれば、同じ生成次数の自由解消を得る。
左右の次数付き noetherian 性も twist で保たれる。
$X_3$ の次数三の平坦性と補題~\ref{lem:relative-koszul-degree-three} から
$B_A$ は全次数自由である。補題~\ref{lem:artin-noetherian-lift} と
命題~\ref{prop:artin-relations-flat-abelian} の torsion 商の証明を適用する。
同証明が使うのは全次数平坦性、閉 fiber の noetherian 性、
および次数一生成性であり、これらはいま確認した。
よって $(\qgr B_A,\pi B_A)$ は元の $(\cC_0,\OO_0)$ の flat pointed Abel 変形である。
系~\ref{cor:cubic-formal-h2-abelian}、\ref{cor:noncubic-all-artin-h2}
によって、これは Artin 基底上で実際の相対 $\HHstack$ 族になり、
指定した還元同型も保たれる。

この族の反 Serre を $t_A$ とする。$u_A=F_A t_A^{-1}$ は閉 fiber で
$u_0$ となる。$u_A\OO_A$ と $\OO_A$ は同じ例外対象の平坦持上げである。
$\Ext^1(\OO_0,\OO_0)=0$ と小拡大による帰納法で閉 fiber の指定同型を持ち上げられる。
得られた pointed 自己同値は加群対象への基底変更で系全体に延びる。
従って $F_A=u_A t_A$ の軌道対象は $t_A$ のものと同定され、
その全軌道 Hom は自由で基底変更と可換する。
自然射 $B_A\to\bigoplus_{n\ge0}\Hom(\OO_A,F_A^n\OO_A)$ は
閉 fiber で上の twist の同型である。各次数が同じ階数の自由加群なので
Nakayama によって全次数で同型となる。
ゆえに次数一の枠を移した canonical 関係式と $u_A$ の行列が
$(J_A,g_A)\in D(A)$ を与え、$\theta(J_A,g_A)=J'_A$ である。

この逆構成は Artin 基底変更と整合する。実際、平坦な $B_A$ の
torsion 商の係数還元は、商の右随伴を用いる
命題~\ref{prop:artin-relations-flat-abelian} の議論を
$A\to A/I$ に対して適用すれば、$\operatorname{QGr}(B_A/IB_A)$ である。
$F_A$ と自由 row の還元もこれに従う。圏と $t_A$ の比較は相対認識の
全忠実性による。逆に $D(A)$ から始めれば、twist の pointed tails 同値と
射影枠の回収により元の点に戻る。特に小拡大で還元先の $D$ の点を
あらかじめ指定していても、その点を保つ一意の持上げとなる。

\emph{開埋込みと切断。}
両スキームは代数閉体 $k$ 上有限型なので、局所 Artin 小拡大による
smoothness 試験 \cite[Tag 02HY]{StacksSmoothAtlas} の有限型閉点の剰余体は $k$ である。
上の持上げにより $\theta$ は smooth で、双対数での一意性により
閉点で相対接空間は零である。従って全域でエタールである。
幾何点上の単射性と \cite[Tag 025G]{StacksBoundaryComparison}
により開埋込みを得る。$[D/G]$ は $\HHstack=[X_{\mathrm{as}}/G]$
の慣性であり、$D$ の恒等切断は $G$ の単位元の逆像だから閉である。
\end{proof}

\begin{corollary}[任意の関係式族への適用]
\label{cor:x3-foundation-neighborhood}
局所有限型基底 $S$ と任意の $j:S\to X_3$ に対し、$S_W=j^{-1}W$
は、その幾何 fiber が構造対象を保つ自己同値による反 Serre 環の
Zhang twist となる点の開集合である。$S_W$ の fppf 局所で階数二十四の
normal foundation $\Lambda$ と、表示された次数移動を実現する
可逆 perfect $\Lambda$-双加群を構成できる。canonical 実現の部分は
$S_W$ 上で $g=1$ を課して得られ、任意の基底変更に整合する。
\end{corollary}
\begin{proof}
開埋込みの逆を $S_W$ 上で取ると、唯一の canonical 族とその自己同値が得られる。
定理~\ref{thm:global-h2-smooth-finite-type} の cubic・noncubic の
smooth atlas を引き戻し、その相対 tilting 窓の和 $T$ と $\Lambda=\End(T)$ を取る。
twist の tails 同値で移した $F=ut$ に対して
$P=\mathbf R\Hom(T,F(T))$ と置く。左作用は $F$ が $T$ の射に及ぼす作用、
右作用は前合成である。標準の導来増強で $F$ は余極限を保ち、
自由生成対象上の値とこの作用から関手は $-\otimes_\Lambda^{\mathbf L}P$
となる。$F^{-1}$ が逆双加群を与える。$P$ は基底上 perfect であり、
$\Lambda^e$ の相対 smooth 性により双加群としても perfect である
（$\Lambda^e$ の有限対角解消に $P$ を tensor する）。
相対 compact Hom の基底変更からこの表示は引戻しと整合する。
次数の正規化による基底線束は局所自明化してこの記述を適用する。
恒等切断を引き戻すと所要の閉条件を得る。行列代表を取れば、
非対角成分と対角成分相互の差をすべて零にする方程式である。
非被約な基底ではこれらの式が基底環で消えることを要求する。
\end{proof}

\begin{example}[$W$ の外の三次表示]
九生成元 $x_i$ に対して、二文字 $x_ix_j$ は $1\le j\le d_i$ のときだけ残し、
その他を零とする。ただし $(d_1,\ldots,d_9)=(0,3,4,0,3,3,4,4,4)$ とする。
長さ二の語は $25$ 個である。$x_i$ で始まる長さ三の語数は
$\sum_{j\le d_i}d_j$ で、$i=1,4$ で零、残りの七箇所で七である。
同じ漸化式により長さ四の語数は $98$ となる。
従って $X_3$ の点を与えるが、$W$ 上の四次成分の階数は $81$ である。
この例では $x_2,x_3$ 以外を零とする商が二生成自由代数になるので、
右 noetherian 性も成り立たない。
\end{example}


\section{Noncubic 境界の二形式表示と内在的認識}\label{sec:moduli-boundary}
$V$ を二次元ベクトル空間とし、
\[
 P=(\operatorname{Sym}^2V\setminus V(\operatorname{disc}))^2,\quad
 \widetilde G=(\operatorname{GL}(V)\times\mathbb G_m^2)\rtimes S_2,\quad
 D=\{(\lambda I_V,\lambda,\lambda)\},\quad G=\widetilde G/D.
\]
これらは第~\ref{sec:noncubic-stack}節の二形式表示スタック
$\mathfrak N^{\mathrm{pres}}=[P/G]$ を与える。
表示商の滑らかさ・次元は命題~\ref{prop:noncubic-presentation-stack}、
内在的被約境界との相対同値は
定理~\ref{thm:noncubic-intrinsic-presentation-equivalence} に示す。
二つの表示が与える幾何点の同値類が一致することと、
非被約基底上の相対群oidの同値を区別する。

% !TeX root = ../nc-quadric-moduli.tex
% 2026-11-22 / v0.1（非可換二次曲面の相対モジュライ原稿へ証明移管）
\subsection{二つの二元二次形式の不変式と good moduli}
\label{subsec:moduli-noncubic-invariants}

本節の $P$、$G$ は第~\ref{sec:moduli-boundary}節の二形式表示であり、
$\mathcal D=\mathbb P(\operatorname{Sym}^2V)\setminus V(\Delta)$ を
非退化二元二次形式の射影パラメータ空間とする。
二形式 $f=a u^2+buv+c v^2$、$g=a'u^2+b'uv+c'v^2$ に対し
\[
 \Delta(f)=b^2-4ac,\qquad D(f,g)=bb'-2ac'-2a'c,\qquad
 \chi(f,g)=\frac{D(f,g)^2}{\Delta(f)\Delta(g)}
\]
と置く。$\chi$ は二形式の基底変換・スカラー変換・交換に不変であり、
$P$ 上の正則関数である。
$\mathfrak N^{\mathrm{pres}}=[P/G]$ と書く。

\begin{proposition}[表示スタックの良モジュライ]
\label{prop:moduli-noncubic-good-moduli}
標数零の代数閉体 $k$ 上で
\[
 k[P]^G=k[\chi].
\]
従って $\chi:\mathfrak N^{\mathrm{pres}}\to\mathbb A^1_\chi$ は
Alper の意味の good moduli space である。
\end{proposition}
\begin{proof}
まず核の二次元トーラスの不変式を取ると、形式の代表の定数倍が除かれ、
\[
 k[P]^G
 =k[\mathcal D^2]^{\operatorname{PGL}_2\times S_2}
\]
となる。トーラスの作用が逆二乗であることは、標数零では不変式環を変えない。
$\operatorname{PGL}_2$ は $\mathcal D$ に推移的に作用するため、
第一形式を $uv$ と固定すれば、順序付き二形式の不変式環は
$uv$ の射影安定化群 $H=\mathbb G_m\rtimes C_2$ による
$\mathcal D$ の不変式環に同定される。
これは斉次空間 $\mathcal D=\operatorname{PGL}_2/H$ の faithfully flat な局所自明化で
不変関数を降下させることによる。

第二形式を $[a:b:c]$ と書く。
$\mathcal D$ は $\mathbb P^2$ の判別式非消滅開集合であり、アフィンで
\[
 k[\mathcal D]=\bigl(k[a,b,c,\Delta^{-1}]\bigr)_0,
 \qquad \Delta=b^2-4ac.
\]
残るトーラスの作用は射影座標で
$(a,b,c)\mapsto(t a,b,t^{-1}c)$ と書ける。
次数零のトーラス不変単項式は
$a^i b^j c^i\Delta^{-m}$、$2i+j=2m$ の形なので、$j$ は偶数である。
従って不変式環は
\[
 x=\frac{b^2}{\Delta},\qquad y=\frac{ac}{\Delta}
\]
で生成され、$x-4y=1$ により $k[x]$ である。
残る $C_2$ は $a,c$ を交換して $x$ を保つ。
第一形式が $uv$ なら $D(uv,f)=b$、$\Delta(uv)=1$ なので $x=\chi$ である。
二形式を交換する $S_2$ も対称な $\chi$ を保つため、全不変式環は $k[\chi]$ となる。

$P$ はアフィン、$G$ は標数零で線型簡約なので、商スタックから不変式環の
スペクトルへの射は good moduli space である
\cite[Example 8.3, Theorem 13.2]{Alper2009}。
\end{proof}

この計算は二形式の\emph{表示商}についての主張である。
表示商と内在的被約 noncubic 境界との相対同値、および
全内在的 $\mathfrak H_2$ における粗モジュライ不存在の証明は、
相対認識・自己同型の比較が必要なので別に扱う。

% !TeX root = ../nc-quadric-moduli.tex
% 2026-11-22 / v0.1（表示商の特殊fiber証明を論文Bへ移管）
\subsection{良モジュライ射の平坦性と特殊ファイバー}
\label{subsec:moduli-noncubic-fibers}

本節は表示商 $\mathfrak N^{\mathrm{pres}}=[P/G]$ そのものの計算であり、
その内在的被約境界との相対同値は分類原稿
\cite{QuadricIntrinsic2026} の結果を別途入力とする。
$Q^{\mathrm{nd}}=\operatorname{Sym}^2 V\setminus V(\operatorname{disc})$、
$P=Q^{\mathrm{nd}}\times Q^{\mathrm{nd}}$ とする。
\begin{proposition}[平坦な良モジュライ射と特殊 fiber]
\label{prop:moduli-noncubic-good-moduli-fibers}
$\chi$ は平坦かつ有限表示である。$\xi\in k$ に対し
$\mathfrak N^{\mathrm{pres}}_\xi=\chi^{-1}(\xi)$ を scheme-theoretic fiber とすると、次が成り立つ。
\begin{enumerate}
\item $\xi\ne0,1$ の fiber は smooth である。
\item $\xi=0$ の fiber は smooth な被約 Cartier divisor の二重構造である。
\item $\xi=1$ の fiber は reduced である。その幾何点は、一根共有の類と二根対一致の類の二つである。
一根共有の類は smooth な開部分をなし、その閉包は一致類を含む。
一致類の近傍で、smooth chart 上の fiber 方程式は $ac=0$ となる。
\end{enumerate}
従って $\chi$ が smooth でない点は、$\xi=0$ の全点と $\xi=1$ の一致類に限られる。
\end{proposition}

\begin{proof}
まず固定した第一形式への帰着を smooth な被覆として実現する。
$\operatorname{GL}(V)$ は $Q^{\mathrm{nd}}$ に推移的に作用し、
$uv$ の安定化群は $O(uv)\simeq\mathbb G_m\rtimes C_2$ で smooth である。
従って軌道射 $\operatorname{GL}(V)\to Q^{\mathrm{nd}}$ は smooth 全射である。
例えば単位元における微分は、二つの非対角成分から $u^2,v^2$ を、
対角成分の和から $uv$ を得るため全射であり、群作用により全点で同様である。
その基底変更に同型な射
\[
 \operatorname{GL}(V)\times Q^{\mathrm{nd}}\longrightarrow P,
 \qquad (A,g)\longmapsto
 \bigl(\operatorname{Sym}^2(A)(uv),\operatorname{Sym}^2(A)g\bigr)
\]
も smooth 全射である。
さらに $P\to[P/G]$ は smooth 全射なので、これらの合成上で fiber の性質を調べられる。

$g=au^2+buv+cv^2$、$\Delta=b^2-4ac$ と書けば、この被覆上で
\[
 \chi=\frac{b^2}{\Delta},\qquad \Delta\in\OO^\times.
\]
従って fiber の方程式は
\begin{equation}\label{eq:moduli-noncubic-fiber-quadric}
 (1-\xi)b^2+4\xi ac=0,\qquad \Delta\ne0
\end{equation}
である。$\xi\ne0,1$ なら、この二次形式は三変数の非退化二次形式であり、
その唯一の affine 特異点である原点は $\Delta\ne0$ により除かれる。
従って fiber は smooth である。

$\xi=0$ では式は $b^2=0$ となる。
被約部分 $b=0$ 上では $ac$ が可逆なので、これは smooth な Cartier divisor の二重構造である。
この divisor は $P$ 上では判別式の極化 $D(f_1,f_2)=0$ で与えられる。
$D$ は $G$ の半不変式であり、分母 $\Delta(f_1)\Delta(f_2)$ は可逆であるため、
商スタック上でも $\operatorname{div}(\chi)=2\mathfrak E_0$ と書ける。
ここで $\mathfrak E_0$ は smooth な被約 Cartier divisor である。

$\xi=1$ では $ac=0$ となり、$\Delta\ne0$ は $b$ が可逆であることと同値である。
これは二つの座標軸の通常交差と可逆な $b$ の積なので reduced である。
$(a,c)\ne(0,0)$ なら二形式は一根だけを共有し、fiber は smooth である。
$a=c=0$ は二形式が比例する場合であり、fiber はそこで交差特異点を持つ。
$b$ を一に正規化した横断座標では
\begin{equation}\label{eq:moduli-noncubic-chi-coincident-node}
 \chi-1=\frac{4ac}{1-4ac}.
\end{equation}
$a$ 軸または $c$ 軸の原点への特殊化が、一根共有の類から一致類への特殊化である。
二つの軸は $u,v$ の交換で同じ軌道へ移るので、幾何点の類は記述した二つで尽くされる。

平坦性は $P$ 上で確認できる。
$P$ は整 affine スキームであり、$\chi$ は非定数なので $k[\chi]\to k[P]$ は単射である。
従って $k[P]$ は PID $k[\chi]$ 上 torsion-free、従って flat である。
有限表示性は有限型 $k$-スキーム・スタックの間の射であることから従う。
平坦性は smooth 全射 $P\to\mathfrak N^{\mathrm{pres}}$ から降下する。
最後の smooth locus の記述は、平坦有限表示性と上で判定した幾何 fiber の smooth locus から従う。
\end{proof}

\begin{remark}[resultant による共有根の方程式]
二つの二元二次形式に対し、標準的な Sylvester resultant を用いると
\[
 D(f_1,f_2)^2-\Delta(f_1)\Delta(f_2)=4\operatorname{Res}(f_1,f_2).
\]
従って $\chi=1$ は共有根の resultant divisor である。
二つの形式自身は非退化であるが、共有根が二つになる一致類では、この divisor の局所方程式が $ac=0$ となる。
\end{remark}

% !TeX root = ../nc-quadric-surface-moduli.tex
% 2026-11-22 / v1.48（AHR形式模型の二次不変式と接微分を明示）
\section{noncubic 境界の内在的商表示とモジュライ写像}
\label{sec:noncubic-boundary-comparison}

標数零の代数閉体 $k$ 上で考える。
本節では、二形式の表示スタックと前節の内在的被約境界を、相対的な対象・射まで比較する。
幾何点の全単射だけでなく、境界の smooth chart 上の実際の相対族の比較を用いる。
その結果、表示側の不変式と慣性計算が内在的境界へ移る。
以下の相対同定とその組合せについて、包括的な新規性は未確認である。

定理~\ref{thm:noncubic-relative-effectivity} の射は、有限型 atlas $P$ と
有限型作用 groupoid $G\times P\rightrightarrows P$ 上の実際の族と同型から作られる。
従ってその降下データは、定理~\ref{thm:global-h2-smooth-finite-type} の
通常の代数スタック $[U/R]$ への射も定める。
以下で代数スタックの一般判定を使う際はこの射を用い、元の有限型 site へ制限する。
任意の Noetherian 基底に素朴に同じ Abel 公理を課す関手との一致は仮定しない。

\subsection{幾何点の慣性から相対的な忠実性へ}

\begin{lemma}[標数零における相対慣性の消滅]
\label{lem:charzero-geometric-inertia-faithful}
$k$ を標数零の体、$f:\mathcal X\to\mathcal Y$ を
有限型の代数的スタックの間の射とし、両対角は affine かつ有限表示とする。
任意の代数閉拡大体 $K/k$ と $x\in\mathcal X(K)$ に対して
\[
 \operatorname{Aut}_{\mathcal X(K)}(x)
 \longrightarrow\operatorname{Aut}_{\mathcal Y(K)}(f(x))
\]
が単射ならば、$f$ は代数空間で表現可能である。
この仮定は幾何点における抽象群の単射でよい。
\end{lemma}

\begin{proof}
有限型 smooth atlas $V\to\mathcal X$ を取り、対応する対象を $x_V$ と書く。
二つの自己同型群スキームの射の核を
\[
 K_V=\ker\bigl(
 \underline{\operatorname{Aut}}_{\mathcal X}(x_V)
 \longrightarrow
 \underline{\operatorname{Aut}}_{\mathcal Y}(f(x_V))\bigr)
\]
とする。対角の仮定と $V$ の Noetherian 性から、$K_V\to V$ は
affine かつ有限表示の群スキームである。
各幾何点 $\bar v\to V$ の上で、その fiber はただ一つの
幾何点を持つ有限型群スキームである。
Cartier の定理により標数零の有限型群スキームは smooth なので、
この fiber は $\operatorname{Spec}\kappa(\bar v)$ に等しい
\cite[Tag 0BF6, Lemma 39.8.2]{StacksBoundaryComparison}。

相対微分は fiber への基底変更と整合する。
従って有限型加群 $\Omega_{K_V/V}$ は全ての幾何点で零となり、
Nakayama の補題により零である。
ゆえに $K_V\to V$ は unramified である
\cite[Tag 02G3, Lemma 29.36.2]{StacksBoundaryComparison}。
その単位切断 $e:V\to K_V$ は unramified 射の切断だから開埋込みである
\cite[Tag 024T, Proposition 41.6.1]{StacksBoundaryComparison}。
全ての幾何 fiber が単位点だけからなるので、$e$ は全射でもある。
従って $e$ は同型であり、$K_V$ は自明である。

これは相対慣性 $\mathcal I_{\mathcal X/\mathcal Y}$ が atlas $V$ 上で自明であることを意味する。
fppf 降下により相対慣性は $\mathcal X$ 上で自明となり、
相対慣性による表現可能性判定を適用できる
\cite[Tag 04YY, Lemma 101.6.2]{StacksBoundaryComparison}。
\end{proof}

\begin{remark}[標数零の仮定が必要な理由]
\label{rem:charzero-inertia-faithful-necessity}
上の補題を正標数へそのまま拡張することはできない。
$\operatorname{char}k=p>0$ のとき
$B\alpha_p\to\operatorname{Spec}k$ は、代数閉拡大体
$K/k$ 上の自己同型の\emph{抽象群}については
$\alpha_p(K)=\{0\}$ なので単射条件を満たす。
しかし相対慣性群スキームは非自明な $\alpha_p$ であり、
Stacks Project \cite[Tag 04YY]{StacksBoundaryComparison} の判定から
この射は代数空間で表現可能ではない。
標数零では有限型群スキームの幾何 fiber が smooth になるため、
この障害を排除できる。
\end{remark}

\begin{lemma}[幾何点の自己同型群のスキームとしての比較]
\label{lem:nc-geometric-inertia-scheme-iso}
自然な射
\[
 f:\mathfrak N^{\mathrm{pres}}\longrightarrow\mathfrak N
\]
は代数空間で表現可能であり、各幾何点で自己同型群スキームの同型を誘導する。
\end{lemma}

\begin{proof}
まず $\mathfrak N^{\mathrm{pres}}=[P/G]$ は smooth、従って reduced である。
定理~\ref{thm:noncubic-relative-effectivity} の射は全幾何点を
noncubic 部分へ送るので、定理~\ref{thm:noncubic-boundary-smooth-codim-two}
の reduced closed substack $\mathfrak N$ を経由する。
この因数分解は reduced atlas $P$ 上で確認して降下させればよい。

命題~\ref{prop:noncubic-presentation-stack} の慣性計算は、
この射が誘導する幾何点の自己同型群の写像を同型と同定している。
実際、$G$ の元は $B_f$ の次数付き同型を与え、その tails 関手が
相対族の構成で用いる射である。次数スカラーは、$\mathcal O$ 上で恒等となる
自然同型によって消え、他の核はない。
したがって前補題より $f$ は表現可能である。

幾何点 $p$ の自己同型群スキーム間の射を $a:G_p\to H_{f(p)}$ と書く。
両群は affine 有限型であり、Cartier の定理によって smooth である。
そのスキームとしての核は前補題の議論により自明なので、
$a$ は群スキームの monomorphism であり、その接写像は単射である。
一方、$a$ は幾何点で全射である。
各 fiber は核の移送で零次元だから両群の次元は等しい。
従って単位元における接写像は同型であり、
smooth な $k$-スキームの間の Jacobian 判定により $a$ は単位元で étale である。
群の平行移動により全ての点で étale となる。
étale な monomorphism は開埋込みであり、幾何点で全射なので同型である
\cite[Tag 025G, Theorem 41.14.1]{StacksBoundaryComparison}。
\end{proof}

\subsection{二形式表示と内在的境界の相対比較}

\begin{lemma}[境界の smooth chart 上での実際の二形式表示]
\label{lem:noncubic-boundary-chart-factorization}
任意の noncubic 幾何点の近傍で、定理
\ref{thm:noncubic-boundary-smooth-codim-two} の smooth chart
\[
 q_D:D_U\longrightarrow\mathfrak N
\]
のある \texorpdfstring{\'etale}{etale} 被覆 $D'_U\to D_U$ と射
$h:D'_U\to P$ を取れて、実際の相対族の同型
\[
 q_D|_{D'_U}\simeq q_{\mathrm{nc}}\circ h
 \tag{NC-chart-factor}\label{eq:noncubic-boundary-chart-factorization}
\]
が成立する。この同型は構造対象を保ち、任意の有限型基底変更と整合する。
根・根同型・AS4 被覆を保存する条件は加えない。
\end{lemma}
\begin{proof}
共通 chart 上で二つの行列式は実際に零である。
命題~\ref{prop:lagrangian-relative-double-rank-three} を適用し、
必要な \texorpdfstring{\'etale}{etale} 被覆上で有限 foundation を
\[
 f_1=uv,\qquad f_2=(\alpha u+\beta v)(\gamma u-\alpha v),
 \qquad\alpha^2+\beta\gamma=1
\]
の標準 foundation と同定する。両形式の判別式は単元なので、これは
$h:D'_U\to P$ を定める。

この段階で得たものは有限代数の同型であり、次に相対 Abel 圏を比較する。
補題~\ref{lem:nc-relative-normal-tilting} により、二形式族の標準四項は
相対 tilting foundation であり、その End 代数は上の有限 foundation に一致する。
一方、$q_D$ の canonical compact 圏も同じ有限代数の perfect 圏に同値である。
従って両側の間には、構造対象 $\mathcal O$ と四項を保つ相対 Morita 同値 $F$ がある。
この同値は perfect 双加群により与えられ、任意の基底変更に整合する。

相対 Serre 双対の Yoneda 一意性により、反 Serre 関手
$t=S^{-1}[2]$ と $t_f=S_f^{-1}[2]$ に対して自然同型
$Ft\simeq t_fF$ を得る。その反復と構造対象の同型を用いると
\[
 \bigoplus_{n\ge0}\operatorname{Hom}(\mathcal O,t^n\mathcal O)
 \simeq
 \bigoplus_{n\ge0}\operatorname{Hom}(\mathcal O_f,t_f^n\mathcal O_f)
\]
が積を保つ graded 環の同型となる。積との整合性は、同じ自然同型の反復を
用いて各次数を比較することから従う。各 Hom は次数零の有限 projective 加群であり、
比較は全有限型基底変更に整合する。

相対反 Serre 再構成（命題~\ref{prop:relative-h2-antiserre-reconstruction}）と
その全忠実性（命題~\ref{prop:relative-h2-stack-recognition}）によって、
この graded 同型は両相対 Abel 系の pointed exact 同値を与える。
特に canonical compact 圏だけでなく、coherent heart と module-object 基底変更の
系が一致する。これが~\eqref{eq:noncubic-boundary-chart-factorization} である。
\end{proof}

\begin{lemma}[smooth 同次元スタックの接判定]
\label{lem:smooth-stack-etale-tangent-comparison}
$k$ を標数零の代数閉体とし、$f:\mathcal X\to\mathcal Y$ を
$k$ 上 smooth、有限型、同じ純次元の Artin スタックの間の射とする。
両対角は affine 有限表示とする。全ての代数閉拡大体 $K/k$ と
$x\in\mathcal X(K)$ について、次を仮定する。
\begin{enumerate}
\item $\operatorname{Aut}_{\mathcal X}(x)(K)\to
  \operatorname{Aut}_{\mathcal Y}(f(x))(K)$ は同型である。
\item 同型類の一次変形の接空間の射
  $T^0_x\mathcal X\to T^0_{f(x)}\mathcal Y$ は全射である。
\end{enumerate}
このとき $f$ は algebraic space で表現可能な \texorpdfstring{\'etale}{etale} 射である。
\end{lemma}
\begin{proof}
補題~\ref{lem:charzero-geometric-inertia-faithful} により $f$ は表現可能である。
各幾何点の慣性群スキームの同型は、補題~\ref{lem:nc-geometric-inertia-scheme-iso}
の群スキームの議論と同じく、Cartier の定理、核の消滅、同次元性から従う。

smooth なスタックでは
\[
 \dim T^0_x\mathcal X
  =\dim_x\mathcal X+\dim\operatorname{Lie}\operatorname{Aut}(x)
\]
である。前段の慣性群の同型と同次元性から、仮定の全射は同型となる。
相対一次変形の完全列
\[
 \operatorname{Lie}\operatorname{Aut}(x)
 \longrightarrow\operatorname{Lie}\operatorname{Aut}(f(x))
 \longrightarrow T_x(\mathcal X/\mathcal Y)
 \longrightarrow T^0_x\mathcal X
 \longrightarrow T^0_{f(x)}\mathcal Y
\]
によって $T_x(\mathcal X/\mathcal Y)=0$ である。
この列は、$f(x)$ の自明な一次変形との同型を付けた $x$ の一次変形を考えれば得られる。
その同型の変更が二つの無限小自己同型空間の余核を与える。

最後に smooth scheme atlas $V\to\mathcal Y$ を取り、
$Z=\mathcal X\times_{\mathcal Y}V$ とする。表現可能性により $Z$ は代数空間で、
$Z\to\mathcal X$ は smooth なので $Z$ 自身も $k$ 上 smooth である。
各点で $Z$ と $V$ の次元は等しく、相対接空間は零である。
両者の \texorpdfstring{\'etale}{etale} scheme chart 上の Jacobian 判定により
$Z\to V$ は \texorpdfstring{\'etale}{etale} となる。
従って $f$ も \texorpdfstring{\'etale}{etale} である。
\end{proof}

\begin{theorem}[二形式表示は内在的被約境界を表す]
\label{thm:noncubic-intrinsic-presentation-equivalence}
二形式の実際の相対族は、局所有限型 $k$-スキームの site 上で
\[
 q_{\mathrm{nc}}:\mathfrak N^{\mathrm{pres}}=[P/G]
  \xrightarrow{\ \sim\ }\mathfrak N
\]
という同値を与える。ここで
$\mathfrak N=(\mathfrak H_2\setminus\mathfrak V_{\mathrm{cub}})_{\mathrm{red}}$
は内在的な被約境界である。従って任意の非被約 test も含め、対象と
pointed exact 同型の両方がこの表示商で記述される。
\end{theorem}
\begin{proof}
表示スタックは smooth、従って reduced であり、その全幾何点は noncubic である。
従って既存の実際の射 $q_{\mathrm{nc}}$ は一意に $\mathfrak N$ を経由する。
この因数分解は atlas $P$ 上で reduced source の性質を使って確認し、降下させる。
両スタックは有限型、smooth、純次元一で、対角は affine 有限表示である。

幾何点の同型類と自己同型群の対応は
命題~\ref{prop:noncubic-presentation-stack} による。
接空間の全射性を確認する。
任意の幾何点 $p\in P(K)$ とその像 $x$ に、
補題~\ref{lem:noncubic-boundary-chart-factorization} の
$q_D:D'_U\to\mathfrak N$ と $h:D'_U\to P$ を取る。
$x$ の上の点 $u$ における $h(u)$ は $p$ と同じ pointed 圏を与えるので、
命題~\ref{prop:noncubic-presentation-stack} により $G(K)$ の同じ軌道に入る。
$h$ を定数の $G(K)$ 元で移せば、$h(u)=p$ としてよい。
\eqref{eq:noncubic-boundary-chart-factorization} と $q_D$ の smooth 性から
\[
 T_uD'_U\longrightarrow T^0_p\mathfrak N^{\mathrm{pres}}
 \longrightarrow T^0_x\mathfrak N
\]
の合成は全射であり、第二の射も全射である。
比較同型の $u$ 上の値を替えると接空間上で可逆な自己同型を合成するだけなので、
この結論は比較の選択に依存しない。
前補題により $q_{\mathrm{nc}}$ は表現可能な \texorpdfstring{\'etale}{etale} 射である。

最後に、全代数閉体上で幾何点 groupoid が同値である。
従って表現可能な $q_{\mathrm{nc}}$ の各幾何 fiber は一点であり、
\texorpdfstring{\'etale}{etale} 性と合わせて全射な開埋込み、すなわち同型となる。
これは \cite[Tag 0DUD, Lemma 101.45.7]{StacksBoundaryComparison} の判定とも一致する。
代数スタックとしての同型から元の site 上での同値が従い、非被約基底も含まれる。
\end{proof}

\begin{remark}[一致根対の二つの接方向と良モジュライの局所不変式]
\label{rem:noncubic-coincident-chi-tangent-loss}
根対が異なる点では無限小慣性は零であり、両境界スタックの接次元は一である。
一致根対では無限小慣性は一次元であり、接次元は二である。
第一形式を $uv$、第二形式の $uv$ 係数を一とする
局所表示
\[
 f_2=a u^2+uv+c v^2,\qquad \Delta(f_2)=1-4ac\ne0
\]
を使うと、$(a,c)=(0,0)$ における二つの接方向が得られる。
連結慣性トーラスは
$(a,c)\mapsto(t^2a,t^{-2}c)$ と作用し、
二因子の交換と二形式の交換も積 $ac$ を保つ
（命題~\ref{prop:noncubic-coincident-gerbe-tangent}）。
したがってこの二変数の\emph{形式的な接表示}に対する
不変式環は
\[
 k[[a,c]]^{\mathbb G_m}=k[[ac]]
\]
であり、全慣性の有限成分も $ac$ に自明に作用する。
ここで主張したのは表示チャート上の接表現の不変式である。
内在的スタック全体の形式局所模型は、後の
定理~\ref{thm:noncubic-coincident-formal-pair} に従う。

表示上の良モジュライ座標は
\[
 \chi=\frac1{1-4ac},\qquad
 \chi-1=\frac{4ac}{1-4ac}
\]
なので、$ac$ は $\chi-1$ の形式局所座標でもある。
厳密に
\[
 \left.\frac{\partial\chi}{\partial a}\right|_{(0,0)}
 =\left.\frac{\partial\chi}{\partial c}\right|_{(0,0)}=0,
 \qquad
 \left.\frac{\partial^2\chi}{\partial a\,\partial c}
 \right|_{(0,0)}=4
\]
が成立する。この一次微分の消滅は、境界の二つの接方向が
消えることを意味しない。
前定理の接空間全射性は、$\chi$ の微分ではなく、
smooth 境界チャートの相対\emph{二形式族の因数分解}から得ている。
これらの導関数と重みは
\path{computations/verify_noncubic_boundary.py} の有限記号計算と整合する。
\end{remark}

\subsection{境界の良モジュライ空間と全体の粗モジュライ障害}

\begin{lemma}[定値な穿孔直線の特殊化]
\label{lem:constant-punctured-line-space}
$k$ を代数閉体、$Z$ を任意の $k$-代数空間とし、
$a:\mathbb A^1_k\to Z$ を射とする。
$a|_{\mathbb G_m}$ が $z:\operatorname{Spec}k\to Z$ を経由する定値射なら、
$a(0)=z$ である。$Z$ に準分離性や有限型性は仮定しない。
\end{lemma}
\begin{proof}
スキームからの \texorpdfstring{\'etale}{etale} chart $U\to Z$ と
$a(0)$ 上の $k$-点 $u_0\in U(k)$ を選び、$u_0$ の affine 近傍に制限する。
代数空間の対角はスキームで表現可能なので
\[
 W=\mathbb A^1_k\times_Z U
\]
はスキームであり、$\mathbb A^1_k$ 上 \texorpdfstring{\'etale}{etale} である。
$(0,u_0)$ を含む既約な affine 開近傍 $W_0\subset W$ を取る。
これは滑らかな曲線であり、$W_0^\times=W_0\times_{\mathbb A^1}\mathbb G_m$
は非空な稠密既約開部分である。

仮定により、$W_0^\times\to U$ は
$U_z=U\times_Z\operatorname{Spec}k$ を経由する。
$U_z$ は $k$ 上 \texorpdfstring{\'etale}{etale} なスキームなので
$k$-点の離散和である。従って既約性により $W_0^\times$ の像は
ただ一つの $k$-点 $u_1\in U(k)$ に含まれる。
この点は affine スキーム $U$ の閉点である。
連続性と $W_0^\times$ の稠密性から、$W_0$ 全体の像も $\{u_1\}$ に含まれ、
$u_0=u_1$ となる。これを $Z$ へ送れば $a(0)=z$ を得る。
ここで使った閉点性はスキーム $U$ 上のものだけであり、
任意の代数空間 $Z$ の $k$-点が閉であるとは仮定していない。
\end{proof}

\begin{corollary}[全 $\mathfrak H_2$ の粗モジュライ代数空間の不存在]
\label{cor:global-h2-no-coarse-space}
一根を共有する根対に対応する点を $x_{\mathrm{sh}}$、
二つの根対が一致する点を $x_{\mathrm{eq}}$ とする。
任意の代数空間 $Z$ と射 $\psi:\mathfrak H_2\to Z$ に対して
\[
 \psi(x_{\mathrm{sh}})=\psi(x_{\mathrm{eq}})
\]
である。従って $\mathfrak H_2$ は、幾何点の同型類を全単射に分類する
粗モジュライ代数空間を持たない。
同じ結論は $\mathfrak N^{\mathrm{pres}}$ および $\mathfrak N$ にも成り立つ。
\end{corollary}
\begin{proof}
$\mathbb A^1_t$ 上の実際の表示族
\[
 f_1=uv,\qquad f_2=u(v-tu)
\]
を使う。両判別式は一であり、これは $P$ への射を定める。
$t=0$ の fiber は $x_{\mathrm{eq}}$、$t\ne0$ の fiber は $x_{\mathrm{sh}}$ である。
穿孔直線上で同型類が一定であるだけでなく、族自身が fppf 局所的に定値となることを確認する。
$t=r^2$ の有限 \texorpdfstring{\'etale}{etale} 被覆を取り、
\[
 A_r=\operatorname{diag}(r^2,1),\qquad a_1=a_2=r
\]
と置く。$G$ の作用式を使うと、$(A_r,r,r)$ は
$(uv,u(v-u))$ を $(uv,u(v-r^2u))$ へ送る。
よってこの族は穿孔直線の被覆上で一定の対象と同型である。

任意の $\psi$ と合成して得る $a:\mathbb A^1\to Z$ は、
穿孔直線上では fppf 局所的に値 $\psi(x_{\mathrm{sh}})$ の定値射である。
代数空間は fppf sheaf なので、これは穿孔直線上の定値射そのものである。
前補題から $\psi(x_{\mathrm{eq}})=\psi(x_{\mathrm{sh}})$ が従う。
体上の分類により二点は異なる pointed exact 同値類であるから、
幾何点の同型類に対して単射となる射は存在しない。
表示スタックについても同じ族と議論を使える。
境界については同値定理、または実際の表示族の $\mathfrak N$ への因数分解を使う。
\end{proof}

\begin{corollary}[内在的境界の良モジュライ空間]
\label{cor:intrinsic-boundary-good-moduli}
定理~\ref{thm:noncubic-intrinsic-presentation-equivalence} の同値の下で、
命題~\ref{prop:moduli-noncubic-good-moduli} の不変式は
\[
 \chi:\mathfrak N\longrightarrow\mathbb A^1_\chi
\]
という good moduli space を与える。
一根共有と根対一致はともに $\chi=1$ へ写る。
この結論は境界の good moduli space を記述するものであり、
$\mathfrak H_2$ 全体の good moduli space の存在はこの系からは導かない。
大域構成は、本稿を入力とする別稿C
\cite[\texttt{thm:qa-global-h2-quasiprojective-good-moduli}]{QuiverAvatarModuli2026}に置く。
\end{corollary}
\begin{proof}
命題~\ref{prop:moduli-noncubic-good-moduli} の
$k[P]^G=k[\chi]$ と linearly reductive 群による affine 商の結果を、
表示スタックと内在的境界の同値に沿って移せばよい。
\end{proof}


\subsection{良モジュライ射の平坦性と二つの特殊 fiber}

以下では $k$ は標数零の代数閉体とする。
定理~\ref{thm:noncubic-intrinsic-presentation-equivalence} により $\mathfrak N=[P/G]$ と同定し、
良モジュライ射 $\chi:\mathfrak N\to\mathbb A^1$ を用いる。
以下の計算は表示商の smooth atlas 上で行い、この同値に沿って内在的境界へ移す。

\begin{proposition}[平坦な良モジュライ射と特殊 fiber]
\label{prop:noncubic-good-moduli-fibers}
$\chi$ は平坦かつ有限表示である。$\xi\in k$ に対し
$\mathfrak N_\xi=\chi^{-1}(\xi)$ を scheme-theoretic fiber とすると、次が成り立つ。
\begin{enumerate}
\item $\xi\ne0,1$ の fiber は smooth である。
\item $\xi=0$ の fiber は smooth な被約 Cartier divisor の二重構造である。
\item $\xi=1$ の fiber は reduced である。その幾何点は、一根共有の類と二根対一致の類の二つである。
一根共有の類は smooth な開部分をなし、その閉包は一致類を含む。
一致類の近傍で、smooth chart 上の fiber 方程式は $ac=0$ となる。
\end{enumerate}
従って $\chi$ が smooth でない点は、$\xi=0$ の全点と $\xi=1$ の一致類に限られる。
\end{proposition}
\begin{proof}
二形式表示商の良モジュライ射の平坦性、各ファイバーの既約・被約性と局所方程式については
本稿の命題~\ref{prop:moduli-noncubic-good-moduli-fibers} に
二形式表示商での全証明を与えた。定理~\ref{thm:noncubic-intrinsic-presentation-equivalence} の
相対スタック同値により、表示商上の命題が内在的境界に移る。
\end{proof}


\subsection{全配置の慣性と一致類の接表現}

以下で $D_8$ は位数八の二面体群を表す。
$c,d$ の独立の符号変更を $e_1,e_2$、交換を $\tau$ とする。
これらは $D_8=\langle e_1,e_2,\tau\rangle$ を生成し、その中心の非自明元は $e_1e_2$ である。

\begin{proposition}[特殊配置も含む全慣性]
\label{prop:noncubic-all-stabilizers}
内在的境界の幾何点の慣性は次の通りである。
\begin{enumerate}
\item 四根が相異なり $\chi\ne0,1$ なら $D_8\times C_2$ で、位数は十六である。
\item $\chi=0$ では位数三十二であり、根配置の慣性 $D_8$ を商とする完全列
\[
 1\longrightarrow\mu_2^2\longrightarrow H_0\longrightarrow D_8\longrightarrow1
\]
を持つ。この完全列は分裂しない。
\item 一根共有では $D_8$ である。
\item 二根対一致では
\begin{equation}\label{eq:noncubic-coincident-full-stabilizer}
 H_{\mathrm{eq}}\simeq
 \bigl((\mathbb G_m\rtimes C_2)\times D_8\bigr)/\mu_2,
\end{equation}
である。$C_2$ は $\mathbb G_m$ に逆元で作用し、割る中央 $\mu_2$ は
$(-1,e_1e_2)$ で生成される。この群の単位成分は $\mathbb G_m$ で、連結成分は八個である。
\end{enumerate}
従って Deligne--Mumford locus は、一致類を除いた開部分である。
\end{proposition}

\begin{proof}
四根が相異なる場合には、根対を $\{0,\infty\}$ と $\{1,\lambda\}$ と書ける。
順序付き四点の交比に対する $S_4$ の作用の核は Klein 四群 $V_4$ である。
二対への分割を保つ順列群は位数八の $D_8$ であり、$V_4$ を含む。
その商の非自明元は $\lambda\mapsto\lambda^{-1}$ と作用するため、
根配置の自己同型群が $V_4$ より大きくなるのは $\lambda=-1$ の場合に限られる。
$\lambda=1$ は四根相異の条件で除かれている。
$\chi=(1+\lambda)^2/(1-\lambda)^2$ なので、この例外はちょうど $\chi=0$ である。
一般慣性の既存の生成元と関係は $\lambda\ne0,1$ の全範囲で定義されるので、
$\lambda\ne-1$ では慣性は $D_8\times C_2$ となる。
特に別の pairing に関する harmonic 値や equianharmonic 値からは、追加の慣性は生じない。

$\lambda=-1$ とし、$\iota^2=-1$ を選ぶ。$f_1=uv$, $f_2=u^2-v^2$ に対して、
列 $(u,v,c,d)^{\mathsf t}$ への置換行列
\[
 R=\operatorname{diag}\!\left(\begin{pmatrix}0&-1\\1&0\end{pmatrix},\iota,\iota\right),
 \quad
 T=\operatorname{diag}\!\left(\begin{pmatrix}-1&1\\1&1\end{pmatrix},
              \begin{pmatrix}0&\iota\\-2\iota&0\end{pmatrix}\right),
 \quad S=\operatorname{diag}(1,-1,\iota,1)
\]
を取る。$R,T$ は既存の一般慣性を生成する変換であり、$S$ はその外の根配置の対称性を持ち上げる。
これらと $e_1,e_2$ は $H_0$ を生成する。
根配置の群が位数八、核が $\mu_2^2$ なので位数は三十二である。
さらに $S^2=e_1$ であり、$S$ は $e_1,e_2$ と可換である。
従って $S$ の根配置への像の任意の持上げは二乗が $e_1$ となり、位数二ではない。
根配置のこの元は位数二なので、群拡大に section は存在しない。
一根共有の場合の $D_8$ は既存の生成元の計算による。

一致類では $f_1=f_2=uv$ と選ぶ。
$O(uv)\simeq\mathbb G_m\rtimes C_2$ は $W=ku\oplus kv$ に作用し、
$D_8$ は $kc\oplus kd$ に符号付き置換行列として作用する。
二作用を block diagonal に並べて射影化する写像
\[
 O(uv)\times D_8\longrightarrow\operatorname{Aut}_{\mathrm{gr}}(B_{uv,uv})/D
\]
を考える。
右辺の任意の元は、$uv$ を定数倍する $A\in\operatorname{GL}(W)$ と、
同じ定数の平方根を係数に持つ $c,d$ の monomial 変換で表される。
その係数の一つを用いて全体を次数スカラーで割ると、$A$ は $O(uv)$ に入り、
$c,d$ の係数は $\pm1$ となる。従ってこの写像は fppf 全射である。
核は block diagonal 行列がスカラーになる条件で決まり、
$(I_2,I_2)$ と $(-I_2,-I_2)$ の二元だけである。
これで式~\eqref{eq:noncubic-coincident-full-stabilizer} を得る。
$O(uv)\times D_8$ は十六個の連結成分を持ち、核の非自明元はその単位成分に入らないので、
商の連結成分数は八である。単位成分にはこの核との交わりがなく、$\mathbb G_m$ そのものである。

標数零の有限型群スキームは smooth である。
従って一致類以外で慣性の Lie 代数は零、そこでスタックは Deligne--Mumford である。
一致類では一次元の Lie 代数があるので Deligne--Mumford ではない。
\end{proof}

\begin{proposition}[一致類の閉 gerbe と接表現]
\label{prop:noncubic-coincident-gerbe-tangent}
一致類の reduced closed substack は $BH_{\mathrm{eq}}$ と同値であり、
$\mathfrak N$ への埋込みは余次元二の正則閉埋込みである。
その法空間、すなわち一致点での一次変形空間は二次元で、
$T(t):(u,v,c,d)\mapsto(tu,t^{-1}v,c,d)$ の作用の重みは $2,-2$ である。
$u,v$ の交換はこの二つの重み空間を交換し、$\tau$ は両者に $-1$ で作用し、
$e_1,e_2$ は恒等に作用する。
\end{proposition}

\begin{proof}
$P$ 上で二形式が比例する locus は
\[
 \{(f,\rho f): f\in Q^{\mathrm{nd}},\ \rho\in\mathbb G_m\}
 \simeq Q^{\mathrm{nd}}\times\mathbb G_m
\]
であり、smooth な四次元閉部分である。
閉性は二つの係数ベクトルを並べた $2\times3$ 行列の二次小行列式で与えられる。
$P$ は smooth な六次元スキームなので、この包含は余次元二の正則閉埋込みである。
$G$ はこの locus に推移的に作用し、その安定化群が $H_{\mathrm{eq}}$ であるため、
商は $BH_{\mathrm{eq}}$ となる。正則埋込みは商へ降下する。

$(uv,uv)$ において、二形式の個別のスカラー変化と共通の
$\operatorname{GL}(W)$ による変化を除くと、接空間は
第二形式と第一形式の差の $u^2,v^2$ 係数で表される。
これを $(a,c)$ と書く。
$T(t)$ は $(a,c)\mapsto(t^2a,t^{-2}c)$ と作用し、
$u,v$ の交換は $(a,c)\mapsto(c,a)$ と作用する。
二形式の交換は差を反転するので $\tau$ は $-1$ となる。
符号変更 $e_1,e_2$ は形式自体に作用しないので接空間でも恒等である。
これは商に用いた中央 $\mu_2$ とも整合する。
\end{proof}


\begin{proposition}[一致根対における連結慣性の接表現]
\label{prop:noncubic-coincident-torus-weights}
$x\in\mathfrak N$ を一致根対の点とし、
$T=\operatorname{Aut}(x)^\circ\simeq\mathbb G_m$ とする。
指標 $\varepsilon$ は、中央二重平方表示における
\[
 u\longmapsto z u,\qquad v\longmapsto z^{-1}v,
 \qquad c\longmapsto c,\qquad d\longmapsto d
\]
の原始指標として取る。その逆指標に取り替えても以下の表示は変わらない。
このとき $T$-表現として
\[
 \begin{split}
 \operatorname{HH}^1(\Lambda_x)&\simeq k,\\
 T^0_x\mathfrak H_2\simeq\operatorname{HH}^2(\Lambda_x)
   &\simeq k_{2\varepsilon}\oplus k_{-2\varepsilon}\oplus k^{\oplus2},\\
 T^0_x\mathfrak N&\simeq k_{2\varepsilon}\oplus k_{-2\varepsilon},\\
 N_{\mathfrak N/\mathfrak H_2,x}&\simeq k^{\oplus2}.
 \end{split}
\]
特に二つの横断方向に対する連結慣性の作用は自明である。
非連結成分を含む全作用は命題~\ref{prop:noncubic-coincident-full-normal-representation} で計算する。
\end{proposition}
\begin{proof}
$f_1=uv$, $f_2=-uv$ の良い対角根を使う。
第~\ref{sec:noncubic-hh-roots}節の Grassmann 模型において
\[
 (g_A,g_B,g_C,g_K)=(T_z,T_z^{-1},T_z,T_z),
 \qquad T_z=\operatorname{diag}(z,1)
\]
が安定化トーラスである。その $H$ 上の行列は
\[
 \rho_z=\operatorname{diag}
 (1,z^{-1},z,1,\ 1,z,z^{-1},1).
\]
式~\ref{eq:nchh-plane} の係数行列を $L$ とすれば
\[
 \rho_zL=L\operatorname{diag}(z,z^{-1},1,1),
 \qquad\rho_z^tJ^{-1}\rho_z=J^{-1}.
\]
従って交代形式は正確に保たれ、conformal 係数による指標の補正はない。
tautological 四平面の指標は $(\varepsilon,-\varepsilon,0,0)$ である。
このトーラスは有効群への写像の核を持たず、$D=\operatorname{Hom}(E_0,E_2)$
の projective 表現で原始指標が現れるため、上の $z$ は本文の内在的トーラスの
パラメータ、またはその逆に一致する。非自明な冪による isogeny ではない。

Lagrangian 接空間 $\operatorname{Sym}^2L^*$ の指標分解は
\[
 k_{2\varepsilon}\oplus k_{-2\varepsilon}
 \oplus k_\varepsilon^{\oplus2}\oplus k_{-\varepsilon}^{\oplus2}
 \oplus k^{\oplus4}.
\]
$\mathfrak g^+$ は二つの $\mathfrak{gl}_2$ から自明な一次元中心を除いたものであり、
その分解は
\[
 \mathfrak g^+
 \simeq k_\varepsilon^{\oplus2}\oplus k_{-\varepsilon}^{\oplus2}\oplus k^{\oplus3}.
\]
軌道写像の核 $\operatorname{HH}^1(\Lambda_x)$ は安定化トーラスの Lie 代数であり、
トーラスの共役作用は自明である。従って軌道接空間は
\[
 k_\varepsilon^{\oplus2}\oplus k_{-\varepsilon}^{\oplus2}\oplus k^{\oplus2}
\]
となる。Lagrangian chart の \texorpdfstring{\'etale}{etale} 性から、
この余核は内在的な全 $\operatorname{HH}^2$ に一致し、表示した分解を得る。

交点二線 $\ell=L\cap U$, $m=L\cap U^*$ は上の四基底の最後の二つで生成され、
両者に対するトーラス作用は自明である。
二つの行列式の微分はそれぞれ $q|_{m\times m}$ と $q|_{\ell\times\ell}$ であり、
上で確認した交代形式の不変性により、この同定はトーラス同変である。
従って境界の法表現は自明な二次元表現となる。
接空間の完全列とトーラスの線型簡約性から、境界の接表現は
$k_{2\varepsilon}\oplus k_{-2\varepsilon}$ となる。
\end{proof}

\subsection{Deligne--Mumford 開部分の粗モジュライ空間}

$\mathfrak N^{\mathrm{DM}}$ を、一致した二根対の閉 gerbe を除いた開部分とする。
相対境界同値の下で、これは $[P^\circ/G]$ と書ける。
ここで $P^\circ\subset P$ は二形式が比例しない開部分である。

\begin{theorem}[粗空間は存在するが良モジュライ射にはならない]
\label{thm:noncubic-dm-coarse}
制限された射
\[
 \chi_{\mathrm{DM}}:\mathfrak N^{\mathrm{DM}}\longrightarrow\mathbb A^1_\chi
\]
は粗モジュライ空間である。
すなわち、任意の代数空間への射に対して普遍であり、任意の代数閉拡大体上で
幾何点の同型類から $\mathbb A^1$ の点への全単射を与える。
この二つの性質は、$\mathbb A^1_\chi$ 上の任意の基底変更後にも成立する。
しかし $\chi_{\mathrm{DM}}$ は cohomologically affine ではなく、従って good moduli space ではない。
また $\mathfrak N^{\mathrm{DM}}$ の inertia 射は proper ではなく、特に finite ではない。
$\chi_{\mathrm{DM}}$ 自身も separated ではない。
\end{theorem}

\begin{proof}
\emph{順序付き根配置への帰着。}
$\mathcal D=\mathbb P(\operatorname{Sym}^2W)\setminus V(\operatorname{disc})$ とし、
$Y=(\mathcal D\times\mathcal D)\setminus\operatorname{Diag}(\mathcal D)$ と置く。
命題~\ref{prop:noncubic-sign-gerbe} の gerbe を制限して
\[
 \mathfrak N^{\mathrm{DM}}\longrightarrow
 \mathfrak R^{\mathrm{DM}}=[Y/(\operatorname{PGL}_2\times S_2)]
\]
を得る。gerbe から代数空間への射は、その基底からの射と同じである。
実際、fppf 局所に gerbe の対象を選ぶと、異なる選択はさらに局所的に同型なので、
代数空間への像は一致し、一意に降下する。
従ってまず根配置の商を扱えばよい。

順序付き二対の商を $\mathfrak R_{\mathrm{ord}}^{\mathrm{DM}}=[Y/\operatorname{PGL}_2]$ とする。
$\operatorname{PGL}_2$ の $\mathcal D$ 上の推移性を使うと
\[
 \mathfrak R_{\mathrm{ord}}^{\mathrm{DM}}
 \simeq[Z/H],\qquad
 Z=\mathcal D\setminus\{[uv]\},\quad
 H=\operatorname{Stab}_{\operatorname{PGL}_2}([uv])
   \simeq\mathbb G_m\rtimes C_2.
\]
この同値は、第一根対を $[uv]$ とすることを構造群の $H$ への reduction と読むことによる。
第一根対の二根は étale 局所に分離できるため、この reduction は相対的にも fppf 局所に取れる。

\emph{一つの affine chart の有限商。}
$Z$ の座標を $[a:b:c]$、判別式を $\Delta=b^2-4ac$ とする。
$Z$ は二つの開部分 $Z_a=\{a\ne0\}$ と $Z_c=\{c\ne0\}$ で覆われ、
$H$ の反射 $J:u\leftrightarrow v$ はこれらを交換する。
$Z_a$ 上で $a=1$ と正規化すると
\[
 Z_a=\operatorname{Spec}k[b,c,(b^2-4c)^{-1}].
\]
$H$ のトーラス $T$ は、パラメータを必要なら逆に取ることで
$(b,c)\mapsto(tb,t^2c)$ と作用する。
$\delta^2=b^2-4c$ を満たす可逆元を加えた有限 étale $\mu_2$-torsor は
\[
 \widetilde Z_a\simeq\mathbb A^1_x\times\mathbb G_{m,\delta},
 \qquad b=\delta x,\quad c=\frac{\delta^2(x^2-1)}4
\]
である。$T$ は $\delta$ のみに通常の乗法で作用する。
被覆の deck involution は $(x,\delta)\mapsto(-x,-\delta)$ であるから、
\begin{equation}\label{eq:noncubic-dm-affine-slice}
 [Z_a/T]\simeq[\mathbb A^1_x/\mu_2],\qquad \chi=x^2.
\end{equation}
右辺の商スタックは有限 étale affine atlas を持つ。
従って不変式環 $k[x]^{\mu_2}=k[x^2]$ による射は、\emph{任意の}代数空間への射に対する
categorical moduli space である
\cite[Tags 0DUP, 0DUR; Lemmas 106.13.4 and 106.13.6]{StacksDMCoarse}。
ここでは対象の代数空間に quasi-separatedness を課していない。

\emph{二つの chart の降下。}
任意の代数空間 $M$ と $H$-不変な射 $f:Z\to M$ を取る。
$Z_a$ への制限は $T$-不変なので、前段により一意な
$h:\mathbb A^1_\chi\to M$ が存在し、$f|_{Z_a}=h\chi$ となる。
$J$ は $Z_a,Z_c$ を交換し、$\chi$ を保つ。
従って $J$-不変性から $f|_{Z_c}=h\chi$ でもあり、$f=h\chi$ が $Z$ 全体で成立する。
一意性も $Z_a$ への制限から従う。
よって $\mathfrak R_{\mathrm{ord}}^{\mathrm{DM}}\to\mathbb A^1_\chi$ は categorical である。
二対の交換も $\chi$ を保つため、この普遍性は
$\mathfrak R^{\mathrm{DM}}$ へ降下し、最初の gerbe の議論で
$\mathfrak N^{\mathrm{DM}}$ へも降下する。

任意の幾何点の $H$-軌道は、$J$ を使えば $Z_a$ と交わる。
式~\eqref{eq:noncubic-dm-affine-slice} の幾何点では、$x^2$ の値ごとに一つの $\mu_2$-軌道がある。
従って $\chi$ の各値に、ちょうど一つの幾何点の同型類が対応する。
$\chi=1$ に対応するのは一根共有の類であり、一致類はこの開部分から除かれている。
gerbe と二対交換も既知の分類によりこの同型類集合を変えない。
これで粗モジュライ空間の二条件が示された。

\emph{任意の基底変更。}
$T=\operatorname{Spec}A\to\mathbb A^1_\chi$ が $\chi\mapsto s\in A$ に対応するとする。
式~\eqref{eq:noncubic-dm-affine-slice} の基底変更は
\[
 [\operatorname{Spec}A[x]/(x^2-s)\,/\mu_2]
\]
であり、その不変式環は $A$ である。
実際、$A[x]/(x^2-s)=A\oplus Ax$ で、$\mu_2$ は第二項に $-1$ で作用し、$2$ は可逆である。
再び同じ有限 affine atlas に対する categorical quotient の定理を適用できる。
前段の二 chart と反射による降下もそのまま基底変更されるので、
$\mathfrak N^{\mathrm{DM}}\times_{\mathbb A^1}T\to T$ は categorical である。
一般のスキーム・代数空間 $T$ については affine étale 被覆上で構成し、一意性により降下すればよい。
幾何点については代数閉体上の先の議論がそのまま使える。

\emph{good moduli space ではないこと。}
$P\setminus P^\circ$ は、比例する二形式の非空な smooth 閉部分で、$P$ 内の余次元は二である。
$P$ は normal affine スキームなので
\[
 \Gamma(P^\circ,\OO)=\Gamma(P,\OO)
\]
となる。これは normal domain がその高さ一の局所環の共通部分に等しいことから従う。
もし $P^\circ$ が affine なら、開埋込み $P^\circ\hookrightarrow P$ がこの同じ環の恒等写像から
誘導される同型となるため矛盾する。従って $P^\circ$ は affine ではない。
$G$ は affine なので atlas $P^\circ\to[P^\circ/G]$ は affine である。
仮に $\chi_{\mathrm{DM}}$ が cohomologically affine なら、その合成
\[
 P^\circ\longrightarrow\mathfrak N^{\mathrm{DM}}
 \longrightarrow\mathbb A^1\longrightarrow\operatorname{Spec}k
\]
も cohomologically affine となる。
Serre の affineness criterion により $P^\circ$ は affine となって矛盾する
\cite[Propositions 3.3 and 3.9]{Alper2009}。

\emph{inertia の非 proper 性。}
$R=k[[t]]$、$K=k((t))$ とし、$P^\circ(R)$ の族
\[
 f_1=uv,\qquad f_2=(u-v)(u-t^2v)
\]
を取る。判別式は $1$ と $(1-t^2)^2$ で可逆であり、
特殊 fiber も一根共有の類なので全体が $P^\circ$ に属する。
一般 fiber には次数付き自己同型を表す射影行列
\[
 A_t=\operatorname{diag}\!\left(
 \begin{pmatrix}0&t^2\\1&0\end{pmatrix},t,t\right)
 \in G(K)\subset\operatorname{PGL}_4(K)
\]
がある。実際、$A_t(f_i)=t^2f_i$、$c,d\mapsto tc,td$ である。
しかし $A_t$ の射影類は $\operatorname{PGL}_4(R)$ に延長できない。
$R$ 上の線束は自明なので、延長できるならある $\lambda\in K^\times$ に対し
$\lambda A_t\in\operatorname{GL}_4(R)$ となる。
$A_t$ には成分 $1$ があるので、全成分が $R$ に属するためには $v(\lambda)\ge0$ を要求するが、
$\det A_t=-t^4$ より $v(\det(\lambda A_t))=4v(\lambda)+4>0$ となり可逆性に矛盾する。
この議論は任意の DVR 拡大後にも同じである。

$G\subset\operatorname{PGL}_4$ は既知の閉埋込みであり、
$P^\circ\to[P^\circ/G]$ への inertia の引戻しは実際の stabilizer 群スキームである。
従ってこの一般点の inertia section は $R$ 上へ延長しない。
proper 射の付値判定に反するので inertia 射は proper ではない。
もし $\chi_{\mathrm{DM}}$ が separated ならその対角は proper であり、
その基底変更である inertia も proper となるため、最後の主張も従う。
\end{proof}

\begin{remark}
全境界では $\chi=1$ に共有根類と一致類の二つがあり、粗モジュライ空間は存在しない。
一致類を除くと、上の粗モジュライ空間が存在する一方、good moduli space の性質は失われる。
この違いは、有限慣性という語を各幾何点の安定化群が有限であるという意味に用いるだけでは捉えられない。
この開部分は Deligne--Mumford だが、その inertia \emph{射}は finite ではない。
\end{remark}

\subsection{一致類の全法表現と形式局所模型}

\begin{proposition}[非連結成分を含む一致類の接表現]
\label{prop:noncubic-coincident-full-normal-representation}
一致類を $f_1=f_2=uv$ で表し、
式~\eqref{eq:noncubic-coincident-full-stabilizer} の群を $H=H_{\mathrm{eq}}$ と書く。
$T(z):(u,v,c,d)\mapsto(zu,z^{-1}v,c,d)$、
$r:(u,v)\mapsto(v,u)$、$e_1:c\mapsto-c$、$e_2:d\mapsto-d$、
$\tau:(c,d)\mapsto(d,c)$ を生成元として取る。
境界の接表現を $V_{\mathrm b}$、その法表現を $V_{\mathrm n}$ とすると、
適当な座標によりその全作用は次の表となる。
表の $a,c$ は変形座標であり、次数一生成元 $c$ とは区別する。
\[
\begin{array}{c|c|c|c}
 &V_{\mathrm b}\ (a,c)&V_{\mathrm n}\ (x,y)&\operatorname{HH}^1\\\hline
 T(z)&(z^2a,z^{-2}c)&(x,y)&1\\
 e_1,e_2&(a,c)&(-x,-y)&1\\
 r&(c,a)&(-x,-y)&-1\\
 \tau&(-a,-c)&(y,x)&1
\end{array}
\]
原始トーラス指標を逆に取る場合は $a,c$ の順序を交換する。
また連結トーラスの重み分解による $H$-表現の直和分解
\[
 \operatorname{HH}^2(\Lambda_x)\simeq V_{\mathrm b}\oplus V_{\mathrm n}
\]
が存在する。
\end{proposition}
\begin{proof}
境界の接作用は命題~\ref{prop:noncubic-coincident-gerbe-tangent} で計算済みである。
$\operatorname{HH}^1$ は $H$ の Lie 代数なので、共役作用から表の最後の列を得る。
法表現を、有限成分で保存されるとは限らない Lagrangian slice 上だけでなく、
全 Grassmann 模型上で計算する。

行列分解を
\[
 \Phi_c=\begin{pmatrix}c&u\\-v&-c\end{pmatrix},\qquad
 \Phi_d=\begin{pmatrix}d&u\\-v&-d\end{pmatrix}
\]
とし、$L=\operatorname{im}(-\Phi_c^t,\Phi_d)\subset H_0$
をこの正規形の四平面とする。連結トーラスの計算で用いた $f_2=-uv$ の表示とは、
$d$ の非零スカラー変更で対応する。ここで $H_0$ は八次元の線型空間を表し、安定化群 $H$ と区別する。
$D=\operatorname{diag}(1,-1)$、$S=\left(\begin{smallmatrix}0&1\\1&0\end{smallmatrix}\right)$ と置く。
係数列の順序を $(u,v,c,d)$ とすれば、四列は
\[
 (-E_{21},E_{12}),\quad(E_{12},-E_{21}),\quad(-D,0),\quad(0,D)
\]
である。

二つの符号変更と根交換は、二 summand の行列 $(X,Y)$ にそれぞれ
\[
 \rho_{e_1}(X,Y)=(-DXD,Y),\qquad
 \rho_{e_2}(X,Y)=(X,-DYD),\qquad
 \rho_r(X,Y)=(-SXS,-SYS)
\]
と作用する。
これは例えば $\Phi_c(-c)=-D\Phi_cD$ と
$\Phi_c(v,u)=-S\Phi_c(u,v)S$、および $d$ 側の同じ恒等式から得られる。
四平面の基底への誘導作用はそれぞれ
\[
 R_{e_1}=\operatorname{diag}(1,1,-1,1),\quad
 R_{e_2}=\operatorname{diag}(1,1,1,-1),\quad
 R_r=P_{uv},
 \qquad\rho_gL=LR_g
\]
である。$P_{uv}$ は最初の二列を交換する行列である。

ruling 交換が実際の自己同値の変形作用であることを確認する。
$c\leftrightarrow d$ は標準 foundation
$(\mathcal O,M_c,\mathcal O(1),M_d(1))$ を
$(\mathcal O,M_d,\mathcal O(1),M_c(1))$ へ送る。
ここで行列分解の周期列から
\[
 0\longrightarrow M_d\longrightarrow B(1)^{\oplus2}
   \longrightarrow M_d(1)\longrightarrow0
\]
および $c$ 側の同じ列が得られる。
この列は、まず $k[u,v,d]/(d^2-uv)$ 上の行列分解で得られ、
$B$ が交代語の PBW 基底によりこの部分環上で左自由であることを使って移せる。
従って $M_d$ は実際の mutation
$\operatorname{Cone}(\operatorname{Hom}(E_2,E_3)\otimes E_2\to E_3)[-1]$
であり、第四項も同様に $M_c(1)$ である。
命題~\ref{prop:ruling-involution-mutation} の相対構成は任意の Artin foundation 変形にも適用できる。
そのため閉点の行列を固定するというだけでなく、当該 pointed 自己同値による
全変形 groupoid の作用を、この mutation と閉点での指定同型で計算できる。

その作用を命題~\ref{prop:relative-ruling-involution} の二つの短完全列から読む。
新しい $F'$ は $K\hookrightarrow C\otimes D$ の縮約なので、
$(F')^*$ は $j$ の二因子を交換した $j^t$ である。
また新しい $G'$ の核は $\operatorname{im}\alpha_F\simeq A$ なので、
その curry 写像 $j'$ は $F^*$ の二因子を交換した $(F^*)^t$ となる。
従って旧四平面を $(-F^*,j)$ と書くと、新しい第一成分は $-j^t$、
第二成分は $(F^*)^t$ である。
従って全体の共通符号を除いて、交換は
\[
 \rho_\tau(X,Y)=(Y^t,X^t)
\]
となる。一致点では $\Phi_c,\Phi_d$ の交換が実際の graded 自己同型 $c\leftrightarrow d$
に対応し、$P_{cd}$ を最後の二列の交換行列として
\[
 \rho_\tau L=-LP_{cd}
\]
である。この計算は実際の mutation の二写像表示に従うものであり、
foundation 商全体と内在的 $\mathfrak H_2$ の相対同値を仮定していない。

二つの射影行列を $M,N$ とし、四平面の一次変分 $\dot L$ に対して
\[
 \nu(\dot L)=
 \bigl(\operatorname{tr}(\operatorname{adj}M\,\dot M),
       \operatorname{tr}(\operatorname{adj}N\,\dot N)\bigr)
\]
と置く。これは orbit 接空間を消し、Lagrangian slice 上では前節の独立な
二つの行列式微分に一致する。従って $\operatorname{HH}^2$ から法空間への
全射を与える。その kernel は二次元であり、境界の接空間に一致する。
ここでは全 Grassmann 接空間の次元十六、orbit 次元十二、境界の接次元二を使う。

各 block の行列式は一、$\det R_{e_1}=\det R_{e_2}=\det R_r=-1$ なので、
基底を戻す $\dot L\mapsto\rho_g\dot L R_g^{-1}$ の下で、
二つの射影行列式の微分はともに $-1$ 倍される。
$\tau$ では二 summand を交換し、二次正方行列の転置が四次元空間上で
行列式 $-1$ を持つ。他方、$\det(-P_{cd})=-1$ である。
この二つの符号が打ち消し合い、法方向への作用は $(x,y)\mapsto(y,x)$ となる。
連結トーラスの法作用は命題~\ref{prop:noncubic-coincident-torus-weights} で自明と示した。
これで法表現の表が得られる。

$H$ はトーラスと有限群の積の有限中心商なので、標数零では線型簡約である。
ここではさらに、連結トーラスの非零重み部分が境界の接空間に一致し、
零重み部分から法空間への写像が同型である。
$H$ はその単位成分を正規化するので、この二つの部分はともに $H$-不変である。
従ってこの重み分解は $H$-同変な直和分解を与え、表により全 $\operatorname{HH}^2$ の表現が決まる。
中央の $(-1,e_1e_2)$ は表の全ての表現で自明に作用するため、
表は式~\eqref{eq:noncubic-coincident-full-stabilizer} の商群上に正しく降下する。
\end{proof}

\begin{theorem}[一致類に沿う形式的な対の線型化]
\label{thm:noncubic-coincident-formal-pair}
$\mathcal G=BH_{\mathrm{eq}}\subset\mathfrak N\subset\mathfrak H_2$ を
一致類の閉 gerbe とし、$V=V_{\mathrm b}\oplus V_{\mathrm n}$ を前命題の全接表現とする。
$\mathcal G$ に沿う形式完備化の対は、非標準的に
\[
 \begin{aligned}
 &(\widehat{\mathfrak H}_{2,\mathcal G},\widehat{\mathfrak N}_{\mathcal G})\\
 &\qquad\simeq
 \left(
 [\operatorname{Spf}k[[a,c,x,y]]/H_{\mathrm{eq}}],
 [\operatorname{Spf}k[[a,c]]/H_{\mathrm{eq}}]
 \right).
 \end{aligned}
\]
となる。右の閉包含は $x=y=0$ で与えられ、群作用は前命題の線型作用である。
また $\mathfrak H_2$ と $[V/H_{\mathrm{eq}}]$ は、
指定した gerbe の近傍に共通の表現可能な \texorpdfstring{\'etale}{etale} 近傍を持つ。
\end{theorem}
\begin{proof}
命題~\ref{prop:noncubic-coincident-gerbe-tangent} より、$\mathcal G$ は
$\mathfrak N$ の閉 substack である。$\mathfrak N$ 自身が $\mathfrak H_2$ の閉 substack
なので、一致類は $\mathfrak H_2$ の閉点である。
$\mathfrak H_2$ は有限型 smooth Artin スタックで、対角は affine 有限表示である。
前命題の $H_{\mathrm{eq}}$ は線型簡約である。
従って Alper--Hall--Rydh の定理
\cite[Theorems 1.1, 1.2]{AlperHallRydh2020} の全仮定が成立する。
同論文の閉点の normal representation は、smooth な安定化群の場合、
ここでの一次変形空間 $V=\operatorname{HH}^2$ に一致する。
Theorem 1.2 によって、$[V/H_{\mathrm{eq}}]$ と $\mathfrak H_2$ の共通
\texorpdfstring{\'etale}{etale} 近傍を得る。
affine diagonal の仮定によって、$\mathfrak H_2$ への射は affine、特に表現可能に取れる。

境界も同時に線型化できることを別に示す。
同定理の affine slice $W$ の固定点を $w$ とする。
$W$ は $w$ の近傍で smooth であり、境界の引戻し $D\subset W$ も smooth で、
$H_{\mathrm{eq}}$-不変な余次元二の正則閉部分である。
ここで $\mathfrak m$ と $I$ は affine slice の座標環 $k[W]$ における
$w$ と $D$ のイデアルとする。両方とも有理 $H_{\mathrm{eq}}$-加群なので、
有限次元表現からの射の同変持上げに線型簡約性を適用できる。
接空間の双対の完全列
\[
 0\longrightarrow (I+\mathfrak m^2)/\mathfrak m^2
 \longrightarrow\mathfrak m/\mathfrak m^2
 \longrightarrow\mathfrak m/(I+\mathfrak m^2)\longrightarrow0
\]
は $H_{\mathrm{eq}}$-同変に分裂する。
線型簡約性により、法方向の余接表現を $I$ の元へ、
境界方向の余接表現を $\mathfrak m$ の元へ、それぞれ同変に持ち上げられる。
選んだ四元は regular system of parameters をなし、
完備局所環の同変同型
\[
 \widehat{\mathcal O}_{W,w}\simeq k[[a,c,x,y]]
\]
を与える。最初に $I$ から選んだ二元を $x,y$ とすれば、
正則埋込みと Nakayama の補題から $I\widehat{\mathcal O}_{W,w}=(x,y)$ である。
従って形式的な対が同時に線型化される。

slice から $\mathfrak H_2$ への射は表現可能な \texorpdfstring{\'etale}{etale} 射であり、
指定 gerbe 上では同型である。従って全ての無限小近傍を同型に写し、
上の形式局所環の対を群で割ることで主張を得る。
ここでは接表現の計算だけから形式模型を推測せず、
\texorpdfstring{\'etale}{etale} slice と線型簡約性による同変座標の構成を使っている。
\end{proof}

\begin{corollary}[形式不変式環と smooth な局所商]
\label{cor:noncubic-coincident-polynomial-invariants}
前定理の全接表現について
\[
 k[V]^{H_{\mathrm{eq}}}
   =k\bigl[ac,(x+y)^2,(x-y)^2\bigr]
\]
は三変数の多項式環である。従って形式不変式環は
\[
 k[[a,c,x,y]]^{H_{\mathrm{eq}}}
   \simeq k[[w,p,q]],
 \qquad w=ac,\quad p=(x+y)^2,\quad q=(x-y)^2
\]
となり、境界への制限は $p=q=0$ で与えられる。
Alper--Hall--Rydh の共通 \texorpdfstring{\'etale}{etale} 近傍は、
原点近傍で smooth な三次元スキームを good moduli space として持つ。
この局所計算と、別稿Cの明示GIT開による全体のgood moduli構成を区別する。
\end{corollary}
\begin{proof}
まずトーラスの不変式を取ると
\[
 k[V]^{\mathbb G_m}=k[w,x,y],\qquad w=ac.
\]
$e_1$ は $(x,y)$ を同時に反転し、$r$ はこの環上では $e_1$ と同じ作用を持つ。
$\tau$ は $w$ を保ち、$x,y$ を交換する。
$\xi=x+y$, $\eta=x-y$ と置くと、残る有限群の像は
$\xi,\eta$ の独立の符号変更群である。
従って不変式環は $k[w,\xi^2,\eta^2]$ であり、三つの生成元は代数的に独立である。
形式冪級数については、トーラス不変な単項式の $a,c$ の指数が等しいことから、
まず不変式環は $k[[w,x,y]]$ となる。
ここで $w$ の重み付き次数を二とする filtration と通常の極大イデアル進 filtration は cofinal である。
次に $\xi,\eta$ の独立の符号変更による不変性は両指数が偶数であることと同値なので、
$k[[w,\xi^2,\eta^2]]$ を得る。
境界は $x=y=0$ なので制限の主張も従う。

\cite[Theorem 1.2]{AlperHallRydh2020} の共通近傍は
\[
 [W/H_{\mathrm{eq}}]
   =[V/H_{\mathrm{eq}}]\times_{V/\!/H_{\mathrm{eq}}}U
\]
と取れ、$U\to V/\!/H_{\mathrm{eq}}=\mathbb A^3$ は affine
\texorpdfstring{\'etale}{etale} である。
線型簡約群の affine 商の good moduli morphism をこの
\texorpdfstring{\'etale}{etale} 射で引き戻すと、$[W/H_{\mathrm{eq}}]\to U$ が得られる。
従ってこの共通近傍の good moduli space は原点近傍で smooth である。
$[W/H_{\mathrm{eq}}]\to\mathfrak H_2$ が近傍全体で慣性を同型に写すとは仮定せず、
good moduli space をそのまま $\mathfrak H_2$ の開部分へ降下させていない。
\end{proof}

\begin{remark}[形式良モジュライ射の二次的分岐]
\label{rem:coincident-local-good-quotient-quadratic}
一致点で全接表現を
$V=V_{\mathrm b}\oplus V_{\mathrm n}=\langle a,c,x,y\rangle$
と書き、四変数に同じ次数一を与える。
前系の三つの不変式
\[
 w=ac,\qquad p=(x+y)^2,\qquad q=(x-y)^2
\]
はいずれも次数二であり、非零の一次不変式は存在しない。
従ってこの線型化模型からその三次元良モジュライ空間への
射の原点における一次微分は\emph{零}である。
一方、商の座標環は三元多項式環であり、その Hilbert 級数は
\[
 \sum_{n\ge0}\dim k[V]^{H_{\mathrm{eq}}}_n\,z^n
 =\frac{1}{(1-z^2)^3}.
\]
これは商の原点が滑らかな三次元点であることと矛盾しない。
接空間の微分が零でも、最初の非自明な商座標が二次に現れるためである。
境界では $p=q=0$、$\chi-1=4w/(1-4w)$ となり、
局所商の境界像は滑らかな一次元の座標軸となる。
ただしこの良モジュライ射は前定理の\emph{共通 étale 近傍}上のもので、
$\mathfrak H_2$ の Zariski 開部分への降下を意味しない。
全安定化群の生成元が $w,p,q$ を保つことと、
三生成元の Jacobian の generic rank が三であることを
\path{computations/verify_noncubic_boundary.py} で厳密検算する。
\end{remark}

\begin{corollary}[一致根対での偏極付き形式近傍の二成分]
\label{cor:coincident-polarized-inertia-branches}
$\mathfrak P=[X_3/\operatorname{PGL}_9]$ を、一致根対の canonical 偏極が与える
閉 gerbe（開 $[W/G]$ 内）に沿って完備化する。その形式スタックは非標準的に
\[
 \left[\operatorname{Spf}
     \bigl(k[[a,c,x,y,z]]/(za,zc)\bigr)/H_{\mathrm{eq}}\right]
\]
である。$H_{\mathrm{eq}}$ は最初の四座標には
命題~\ref{prop:noncubic-coincident-full-normal-representation} の表で作用する。
$z$ には $r$ が $-1$、他の表示生成元が $1$ で作用する。
canonical 偏極の成分は $z=0$、もう一つの成分は $a=c=0$ である。
\end{corollary}
\begin{proof}
定理~\ref{thm:polarization-inertia-open} により、自己同型が閉 fiber で恒等となる
慣性を計算すればよい。定理~\ref{thm:noncubic-coincident-formal-pair} の
商表示を Artin test に適用すると、その対象は局所的に
変形点 $v=(a,c,x,y)$ と $hv=v$ を満たす $h\in H_{\mathrm{eq}}$ の組、
射は $v$ に作用し $h$ を共役する元である。
恒等に還元する $h$ は連結トーラスにあり、標数零の形式座標
$z=\log T$ を使う（$T$ はトーラスのパラメータ）。
固定条件は $(\exp(2z)-1)a=(\exp(-2z)-1)c=0$ で、
$x,y$ には条件がない。二つの係数はいずれも $z$ の単元倍である。
$r$ はトーラスを反転し、他の生成元はこれと可換なので、$z$ の作用も従う。
最後に $(za,zc)=(z)\cap(a,c)$ は相異なる二つの素イデアルの交わりであり、
両成分の共通部分は $\operatorname{Spf} k[[x,y]]$ の商になる。
\end{proof}

% !TeX root = ../nc-quadric-surface-moduli.tex
% 2026-11-21 / v1.39
\section{構造対象の忘却と楕円型の有限・無限ファイバー}
\label{sec:unpointed-elliptic-fiber}

全ての Abel 圏を一つの代数スタックに集めることは仮定しない。
ここでは、$\HH$ の同型類から、その基礎 Abel 圏の exact 同値類への
集合論的な忘却写像を考える。固定 $\cC$ でのこの写像のファイバーは、
許容される $\OO$ の同型類を
$\operatorname{Aut}^{\mathrm{ex}}(\cC)$ で割った集合である。
固定 $\cC$ の上のスタックの \emph{2-fiber} はこの商ではなく、
許容 $\OO$ の同型群oidであることに注意する。

\begin{lemma}[許容構造対象の数値類と一意性]
\label{lem:forgetting-numerical-rigidity}
$(\cC,\OO)\in\HH$ を固定し、
定理~\cite[\texttt{thm:normal-helix}]{QuadricIntrinsic2026} の foundation によって
$K_0(\cC)\simeq K_0(Q)$、$Q=\PP^1\times\PP^1$ と同定する。
$\cC$ 内の Gabriel--Krull rank-one torsion-free exceptional 対象 $L$ は
ある $a,b\in\mathbb Z$ に対して
\[
 \operatorname{ch}[L]
 =1+aH_1+bH_2+abH_1H_2
 =\operatorname{ch}[\OO_Q(a,b)]
\]
を満たす。この数値類を指定したとき、該当する $L$ の同型類は
高々一つである。特に、固定 $\cC$ 上の許容構造対象の
同型類の集合は高々可算である。
\end{lemma}
\begin{proof}
命題~\cite[\texttt{prop:first-orbit-object}]{QuadricIntrinsic2026} の Euler 形式の計算と
$r=1$ という rank 条件を使う。
$\operatorname{ch}[L]=r+aH_1+bH_2+cH_1H_2$ と書くと
\[
 1=\chi(L,L)=r^2+2rc-2ab=1+2c-2ab.
\]
従って $c=ab$。上の $K_0$ は捩れを持たず、
この Chern 指標で整数類が一意に定まる。

同じ類を持つ二つの対象 $L,L'$ を取る。
$t=S^{-1}[2]$ と置く。もし $\Ext^2(L,L')\ne0$ なら
Serre 双対性により $\Hom(L',t^{-1}L)\ne0$ であり、
rank-one torsion-free 性から $L'\hookrightarrow t^{-1}L$ が得られる。
命題~\cite[\texttt{prop:first-orbit-object}]{QuadricIntrinsic2026} の証明で用いた AZ 豊富性と
自由被覆・負切断消滅を任意の coherent 対象に適用すると、
$n\gg0$ で
$\dim\Hom(\OO,t^n M)=\chi(\OO,t^nM)$ が成り立つ。
両辺に $\Hom(\OO,t^n-)$ を施すと、Riemann--Roch により
\[
 (2n+a+1)(2n+b+1)
 \le(2n+a-1)(2n+b-1)
\]
となり、矛盾する。従って $\Ext^2(L,L')=0$。
$\chi(L,L')=1$ から $\Hom(L,L')\ne0$ であり、同様に単射
$L\hookrightarrow L'$ を得る。余核は $K_0$ 類が零なので
Hilbert 多項式も零である。
$t$ の AZ 豊富性により非零 coherent tails 対象には
十分高い次数に非零切断が存在するから、余核は零である。
\end{proof}

滑らかな楕円曲線 $E$ と次数二の線束 $L_0,L_1,L_2$ に対し、
線束の類を加法的に $\ell_i=[L_i]\in\operatorname{Pic}^2(E)$ と書く。
差を
\[
 \eta=\ell_1-\ell_0,\qquad \delta=\ell_2-\ell_0
 \quad\in\operatorname{Pic}^0(E)
\]
とする。$(E,L_0,L_1,L_2)$ が VdB 許容であることは、
この滑らかな場合には $\eta\ne0,\eta\ne\delta,\delta\ne0$
と同値である。実際、相異なる次数二の線束による二つの
二重被覆 $E\to\PP^1$ は異なる被覆対合を持つので、
積写像は閉埋込みとなる。

\begin{lemma}[楕円四つ組の二つの反射]
\label{lem:elliptic-two-reflections}
許容楕円四つ組の helix の添字を一つ進める操作 $T$ は
\[
 T:(\eta,\delta)\longmapsto(\delta-\eta,\delta)
\]
を誘導する。また、
\[
 \eta\ne2\delta
 \tag{R-open}\label{eq:elliptic-ruling-open}
\]
の下で、命題~\ref{prop:ruling-involution-mutation} による
もう一方の ruling の foundation も cubic 型であり、
対応する四つ組は
\[
 R:(L_0,L_1,L_2)\longmapsto
 (L_0L_1L_2^{-1},L_2,L_1).
\]
従って
\[
 R:(\eta,\delta)\longmapsto(2\delta-\eta,\delta),
 \qquad RT:(\eta,\delta)\longmapsto(\eta+\delta,\delta).
\]
二つの操作は、それぞれ underlying Abel 圏の exact 同値類を保つ。
$R$ は $\OO$ を保ち、$T$ は一般に $\OO$ を変更する。
\end{lemma}
\begin{proof}
helix の関係
$L_{i+3}=L_i^{-1}L_{i+1}L_{i+2}$ により、
$\ell_{i+2}-\ell_i=\delta$ は $i$ に依存しない。
従って $(L_1,L_2,L_3)$ の差は
$(\delta-\eta,\delta)$ である。
添字の変更は同じ $\mathbb Z$-代数の reindexing なので、
基礎 tails Abel 圏を変えない。

後半は ruling 交換の二写像を楕円曲線の切断で計算する。
$V_i=H^0(E,L_i)$、$W_2=H^0(E,L_0L_1)$、
$W_3=H^0(E,L_0L_1L_2)$ とすると、
元の foundation の二写像は
\[
 V_1\otimes V_0\longrightarrow W_2,\qquad
 V_2\otimes W_2\longrightarrow W_3
\]
という切断の積である。
次数二の線束の評価列を $L_0L_1$ で捩ると、
\[
 \ker(V_2\otimes W_2\longrightarrow W_3)
 \simeq H^0(E,L_0L_1L_2^{-1})
\]
を得る。ここでは $\deg(L_0L_1L_2^{-1})=2$ による
$H^1$ 消滅を用いた。
一方、$L_1$ の評価列を $L_0L_1$ で捩ると、
$V_0\to V_1^\vee\otimes W_2$ の余核は
$H^0(E,L_0L_1^2)$ に同定される。
この同定では二次元空間の
$V_1^\vee\simeq V_1\otimes(\det V_1)^{-1}$ に由来する
一次元のスカラー因子を許す。
従って交換後の二写像は、
$(L_0L_1L_2^{-1},L_2,L_1)$ から得る
隣接二つの切断乗法と一致する。

新しい隣接次数二線束が相異なる条件は
$\eta\ne2\delta$ および $\eta\ne\delta$ であり、
残る非周期性条件は $\delta\ne0$ である。
ゆえに式~\eqref{eq:elliptic-ruling-open} の下で
新しい四つ組も VdB 許容である。
命題~\ref{prop:ruling-involution-mutation} と
定理~\cite[\texttt{thm:normal-helix}]{QuadricIntrinsic2026} による qgr の再構成により、
これは元の $(\cC,\OO)$ のもう一方の normalized foundation である。
差の計算は直接に
$\eta'=2\delta-\eta,\ \delta'=\delta$ を与える。
\end{proof}

\begin{theorem}[楕円型の全標識と忘却ファイバー]
\label{thm:elliptic-marking-lattice-fiber}
標数零の代数閉体上の滑らかな楕円許容四つ組
$U=(E,L_0,L_1,L_2)$ に対し、
\begin{equation}\label{eq:elliptic-nonresonance}
 \delta\ne0,\qquad \eta\notin\mathbb Z\delta
\end{equation}
を仮定する。$\cC=\qgr D(U)$ と置き、元の normalized foundation により
$K_0(\cC)$ を $K_0(Q)$ と同定する。
\begin{enumerate}
\item 各 $(a,b)\in\mathbb Z^2$ に対し、
$\operatorname{ch}L(a,b)=\operatorname{ch}\OO_Q(a,b)$ を持つ
rank-one torsion-free exceptional 対象 $L(a,b)$ が一意に存在する。
これらがそのような対象を尽くし、全ての $(\cC,L(a,b))$ が $\HH$ に属する。
\item 許容構造対象の同型類を $\operatorname{Aut}^{\mathrm{ex}}(\cC)$ で割った
忘却ファイバーは $\mathbb Z/{\sim}$ と同定される。
整数 $n$ は $L(n,0)$ に対応し、同値関係は次の有限な幾何学的判定で与えられる：
\begin{equation}\label{eq:elliptic-exact-fiber-relation}
 n\sim m
 \quad\Longleftrightarrow\quad
 \exists\epsilon\in\Aut(E,0),\quad
 \epsilon(\delta)=\delta,\quad
 \epsilon(\eta+n\delta)\in
 \{\eta+m\delta,\ 2\delta-\eta-m\delta\}.
\end{equation}
原点の選択はこの条件を変えない。
\item $\delta$ が非捩れなら、このファイバーは可算無限である。
$\delta$ が位数 $r$ の捩れ元なら、ファイバーは有限であり、
\eqref{eq:elliptic-exact-fiber-relation} を $\mathbb Z/r\mathbb Z$ に課した商である。
特にその濃度は $r$ 以下である。
\end{enumerate}
本定理は幾何点上の対象と同値類についての主張であり、
相対スタック射の離散性や étale 被覆性を結論しない。
\end{theorem}
\begin{proof}
まず対象を全て構成する。補題~\ref{lem:elliptic-two-reflections} の操作により
$U_n=(RT)^nU$ は具体的に
\[
 U_n=(E,L_0,L_1\otimes(L_2L_0^{-1})^{\otimes n},L_2)
\]
と書け、その差は $(\eta+n\delta,\delta)$ である。
仮定~\eqref{eq:elliptic-nonresonance} は、正負いずれの反復でも
$\eta+n\delta$ が $0,\delta,2\delta$ になることを排除する。
途中の $T$ による差 $\delta-\eta-n\delta$ に対しても同じである。
従って全ての四つ組が許容であり、各 $U_n$ から得る pointed 対象を
同じ underlying Abel 圏 $\cC$ へ移せる。

この構成の $K_0$ 上の効果を確認する。
元の foundation の数値類は
\[
 (\OO_Q,\OO_Q(1,0),\OO_Q(1,1),\OO_Q(2,1))
\]
であり、四つ先の項は $\OO_Q(2,2)$ に対応する。
$T$ は構造対象を第一 ruling へ移し、二つの座標軸を交換する。
$R$ の二つの評価列は、foundation を $(O,X,M,Z)$ と書くと
\[
 [Y]=2[M]-[Z],\qquad [Z']=2[M]-[X]
\]
を与える。従って $R$ は構造対象を固定し、二つの軸を交換する。
特に $RT$ による foundation の数値類は
\[
 (\OO_Q(1,0),\OO_Q(2,0),\OO_Q(2,1),\OO_Q(3,1))
\]
である。帰納的に、$U_n$ の構造対象は
$\operatorname{ch}\OO_Q(n,0)$ の対象 $L(n,0)$ となる。
逆操作でも同じ計算が成り立つ。
これらは定理~\cite[\texttt{thm:all-admissible-inverse}]{QuadricIntrinsic2026} により許容構造対象である。

元の cubic helix の適合二周期から来る根 $s$ を選ぶ。
この根は全 helix 項を二つずつ移す。
四項の基底と $s^2=t$ によれば、その $K_0$ 作用は
$\OO_Q(1,1)$ の tensor 作用に一致する。
ゆえに
\[
 L(a,b)=s^b L(a-b,0)
\]
は指定数値類を持つ許容構造対象である。
補題~\ref{lem:forgetting-numerical-rigidity} により、
全ての rank-one torsion-free exceptional 対象の数値類はこの形であり、
同じ類の対象は一意なので (1) が従う。
また自己同値 $s$ による対角移動を除くと、
全ての忘却ファイバーの元を $L(n,0)$ で代表できる。

次にその重複を決める。
系~\ref{cor:two-intrinsic-foundations} により、pointed exact 同値は
二つの内在的 foundation の非順序対を保つ。
$U_n$ の二つの foundation は $U_n,RU_n$ であり、
仮定により両者は geometric cubic 型である。
Van den Bergh の四つ組の復元
\cite[arXiv v3, Theorem 5.5.10]{VdB2011} を用いると、
$L(n,0)$ と $L(m,0)$ が pointed 同値となる必要十分条件は、
$U_n$ が $U_m$ または $RU_m$ と曲線同型を介して同型になることである。
ここでは ordered foundation から full helix を反 Serre 移動で復元する。
逆に四つ組の同型は添字零を保つ tails 同値を与える。

楕円曲線の自己同型を translation と $\epsilon\in\Aut(E,0)$ に分ける。
translation は $\operatorname{Pic}^0(E)$ 上の差には作用しないので、
必要条件はちょうど \eqref{eq:elliptic-exact-fiber-relation} である。
十分性も成り立つ。差が一致した後は、次数二線束 $L_0$ の差を
translation で合わせれば三つの線束が全て一致する。
この translation の存在は、代数閉体上の
$[2]:E\to E$ の全射性による。従って (2) が従う。

$\delta$ が非捩れなら、$\epsilon(\delta)=\delta$ となる
原点を保つ自己同型は恒等に限る。
非自明な $\epsilon$ に対して $\epsilon-1$ は非零 isogeny であり、
その核の元は捩れだからである。
従って同値関係は $n=m$、または
\[
 2\eta=(2-n-m)\delta
\]
だけである。第二式が可能でも、ある整数 $q$ による反射
$n\mapsto 2-q-n$ の二元以下の軌道を同一視するだけなので、商は可算無限である。
$\delta$ の位数が $r$ なら、まず $n\equiv m\pmod r$ が同値を与え、
残る条件はこの有限巡回群上の \eqref{eq:elliptic-exact-fiber-relation} となる。
これで (3) を得る。
\end{proof}

\begin{corollary}[一般楕円型における忘却の可算無限ファイバー]
\label{thm:generic-elliptic-unpointed-fiber}
滑らかな楕円許容四つ組で $\delta$ が非捩れ、
$\eta\notin\mathbb Z\delta$ なら、忘却の同型類ファイバーは可算無限である。
さらに $2\eta\notin\mathbb Z\delta$ なら、整数 $n$ による上記の代表に重複はない。
特に幾何学的一般点ではこの最後の条件を満たす。
\end{corollary}
\begin{proof}
前定理と、\eqref{eq:elliptic-exact-fiber-relation} の非捩れの場合の記述による。
幾何学的一般点では独立な Picard パラメータに対し
全ての等式 $r\delta=0$ と $2\eta=q\delta$ を避ける。
\end{proof}

\begin{corollary}[余分な楕円自己同型がない捩れの場合の濃度]
\label{cor:elliptic-torsion-fiber-cardinality}
前定理の仮定に加え $\Aut(E,0)=\{1,-1\}$ とし、
$\delta$ の位数を $r\ge2$ とする。
$2\eta\notin\langle\delta\rangle$ なら忘却ファイバーは正確に $r$ 個である。
$2\eta=q\delta$ なら、その個数は
\[
 \frac{r+\#\{n\in\mathbb Z/r\mathbb Z:2n=2-q\}}2
 =\begin{cases}
   (r+1)/2,&r\text{ が奇数},\\
   r/2+1,&r\text{ が偶数で }q\text{ が偶数},\\
   r/2,&r\text{ が偶数で }q\text{ が奇数}.
  \end{cases}
\]
余分な自己同型を持つ楕円曲線には、前定理の有限群を含む判定式を使う。
\end{corollary}
\begin{proof}
$r>2$ なら $\delta$ を固定する $\epsilon$ は恒等だけなので、
重複がある場合は反射 $n\mapsto2-q-n$ による商である。
$r=2$ には $\epsilon=-1$ も許されるが、
$2\delta=0$ を \eqref{eq:elliptic-exact-fiber-relation} に代入すると
恒等と同じ二つの条件しか生じない。
反射の軌道数は固定点数を用いる表示式となり、合同式 $2n=2-q$ を解けばよい。
\end{proof}

\begin{remark}[既存研究と残る問題]
上の定理は条件~\eqref{eq:elliptic-nonresonance} の下で全標識と
幾何点ファイバーの同値判定を与えるが、
自然な $\mathbb Z$-torsor 構造や相対ファイバーの代数性は与えない。
Kitamura--Okawa の2022年 Oberwolfach 報告\cite{OkawaKitamura2022OWR}は
cubic $\mathbb Z$-代数の無標識同値に $D_\infty$ が現れることと
rank-one torsion-free line bundle の $\mathbb Z^2$ 分類を述べる。
同報告で詳細論文は当時準備中とされているので、
本節の証明はその未照合の全分類定理に依存させず、
本稿の二つの foundation と楕円評価列から構成した。
\end{remark}


\section{残る比較問題}\label{sec:moduli-remaining}
$X_{\mathrm{as}}\to X_3$ の局所閉性・偏極認識方程式、
foundation stack との全相対 groupoid 同値、
任意 Noetherian 基底上での素朴な Abel 公理の延長は、
既証明の有限型代数性と区別する。
標識を忘れた場合について定理~\ref{thm:unmarked-diagonal} が決めるのは
分類射の標的の対角の性質である。非準分離な無標識モジュライの構成は残る。

\bibliographystyle{alpha}
\bibliography{references}
\end{document}
